\documentclass[lang=cn,11pt,scheme=chinese,thmcnt=section,usesamecnt]{elegantbook}
\usepackage{unicode-math}
\setmathfont{NewCMMath-Regular.otf}
\usepackage{imakeidx}
\usepackage{longtable}
% elegantbook 的 note 环境以 \endnote 作结束码，与 endnotes 宏包冲突；先保存再恢复。
\let\elegantnoteend\endnote
\usepackage{endnotes}
\let\latexendnote\endnote
\let\endnote\elegantnoteend
\makeindex[title=中文索引,intoc,columns=2]
\indexsetup{othercode={\raggedcolumns\renewcommand{\indexspace}{\par\vskip2pt plus1pt minus1pt}}}
\counterwithout{section}{chapter}
\setcounter{tocdepth}{1}
\setcounter{secnumdepth}{2}
\hypersetup{bookmarksdepth=2,pdftitle={Fukaya 范畴与 Picard–Lefschetz 理论},pdfauthor={Paul Seidel}}
\renewcommand{\thechapter}{\Roman{chapter}}
\renewcommand{\thesection}{\arabic{section}}
\renewcommand{\thesubsection}{(\arabic{section}\alph{subsection})}
\counterwithout{equation}{chapter}
\counterwithin{equation}{section}
\renewcommand{\theequation}{\thesection.\arabic{equation}}
\makeatletter
% Place page destinations in the shipout foreground so figures cannot discard them.
\let\Hy@EveryPageAnchor\relax
\AddToHook{shipout/foreground}{%
  \ifHy@pageanchor
    \xdef\@currentHpage{page.\thepage}%
    \put(0,-18){\hyper@anchorstart{page.\thepage}\hyper@anchorend}%
  \fi
}
\renewcommand{\theELEGANT@samecnt}{\thesection.\arabic{ELEGANT@samecnt}}
\AtBeginEnvironment{note}{\refstepcounter{ELEGANT@samecnt}}
\renewcommand{\notename}{笔记~\theELEGANT@samecnt}
\AtBeginEnvironment{assumption}{\refstepcounter{ELEGANT@samecnt}}
\renewcommand{\assumptionname}{假设~\theELEGANT@samecnt}
\AtBeginEnvironment{example}{\stepcounter{ELEGANT@samecnt}}
\AtBeginEnvironment{exercise}{\stepcounter{ELEGANT@samecnt}}
\AtBeginEnvironment{problem}{\stepcounter{ELEGANT@samecnt}}
\renewcommand{\theexam}{\theELEGANT@samecnt}
\renewcommand{\theexer}{\theELEGANT@samecnt}
\renewcommand{\theprob}{\theELEGANT@samecnt}
\newcommand{\StatementNumber}[1]{\setcounter{ELEGANT@samecnt}{\numexpr#1-1\relax}}
\makeatother
\newcommand{\SourcePage}[1]{\phantomsection\label{source:#1}\ignorespaces}
\newcommand{\sr}[2]{\hyperref[#1]{#2}}
\newcommand{\secref}[1]{\hyperref[sec:#1]{小节~(#1)}}
\newcommand{\stref}[2]{\hyperref[st:#2]{#1~#2}}
\newcommand{\eref}[1]{\hyperref[eq:#1]{(#1)}}
\newcommand{\bcite}[1]{[\hyperlink{bib:#1}{#1}]}
\newcounter{citetag}
\ExplSyntaxOn
\cs_new_protected:Npn \book_cite:n #1
  {
    \stepcounter{citetag}
    \phantomsection\label{cite:\arabic{citetag}}
    \clist_if_exist:cF {g_book_bib_#1_clist}{\clist_new:c {g_book_bib_#1_clist}}
    \clist_gput_right:cx {g_book_bib_#1_clist}{cite:\arabic{citetag}}
    [\hyperlink{bib:#1}{#1}]
  }
\cs_new_protected:Npn \book_backref:n #1
  {
    \group_begin:
    \seq_clear:N \l_tmpa_seq
    \clist_if_exist:cT {g_book_bib_#1_clist}
      {
        \clist_map_inline:cn {g_book_bib_#1_clist}
          {
            \tl_set:Nx \l_tmpa_tl {\getpagerefnumber{##1}}
            \exp_args:NNV \seq_if_in:NnF \l_tmpa_seq \l_tmpa_tl
              {
                \seq_put_right:NV \l_tmpa_seq \l_tmpa_tl
                \allowbreak\enspace\hyperref[##1]{\pageref*{##1}}
              }
          }
      }
    \group_end:
  }
\renewcommand{\bcite}[1]{\book_cite:n{#1}}
\newcommand{\BibBackref}[1]{{\footnotesize\book_backref:n{#1}}}
\ExplSyntaxOff
\newcommand{\Ai}{A_\infty}
\newcommand{\A}{\mathcal A}
\newcommand{\B}{\mathcal B}
\newcommand{\C}{\mathcal C}
\newcommand{\F}{\mathcal F}
\newcommand{\G}{\mathcal G}
\newcommand{\M}{\mathcal M}
\newcommand{\N}{\mathcal N}
\newcommand{\Y}{\mathcal Y}
\newcommand{\LG}{\mathcal L}
\newcommand{\RG}{\mathcal R}
\newcommand{\sgnmark}{\maltese}
\newcommand{\Q}{\mathcal Q}
\newcommand{\K}{\mathbb K}
\newcommand{\Z}{\mathbb Z}
\newcommand{\R}{\mathbb R}
\newcommand{\CC}{\mathbb C}
\DeclareMathOperator{\Ob}{Ob}
\DeclareMathOperator{\Hom}{Hom}
\DeclareMathOperator{\Tw}{Tw}
\DeclareMathOperator{\Cone}{Cone}
\DeclareMathOperator{\Id}{Id}
\DeclareMathOperator{\id}{id}
\DeclareMathOperator{\im}{im}
\DeclareMathOperator{\HH}{HH}
\DeclareMathOperator{\End}{End}
\newcommand{\nufun}{\operatorname{nu\text{-}fun}}
\newcommand{\modcat}{\operatorname{mod}}
\newcommand{\Ch}{\operatorname{Ch}}
\newcommand{\dg}{\mathrm{dg}}
\newcommand{\CP}{\mathbb{CP}}
\newcommand{\RP}{\mathbb{RP}}
\newcommand{\br}{\sharp}
\newcommand{\eps}{\epsilon}
\newcommand{\ACone}{\mathcal C\!\mathit{one}}
\newcommand{\ModShift}{\mathit{SS}}
\setCJKfamilyfont{zhkai}[BoldFont=SimHei,ItalicFont=KaiTi,BoldItalicFont=SimHei]{KaiTi}
\setCJKfamilyfont{zhfs}[BoldFont=SimHei,ItalicFont=FangSong,BoldItalicFont=SimHei]{FangSong}
\allowdisplaybreaks[2]
\emergencystretch=2em
\usepackage{float}
\newcommand{\SourceDiagram}[2]{\begin{figure}[H]\centering\includegraphics[width=#2\paperwidth]{images/#1.pdf}\end{figure}}
\newcommand{\SourceEquationDiagram}[3]{\begin{figure}[H]\begin{equation}\label{eq:#3}\includegraphics[width=#2\paperwidth]{images/#1.pdf}\end{equation}\end{figure}}

\counterwithout{figure}{chapter}
\counterwithin{figure}{section}
\renewcommand{\thefigure}{\thesection.\arabic{figure}}
\newcommand{\SourceFigure}[3]{\begin{figure}[H]\centering\includegraphics[width=#2\paperwidth]{images/#1.pdf}\caption{}\label{fig:#3}\end{figure}}

\newcommand{\Efamily}{\widetilde{\mathcal E}}
\newcommand{\Sfamily}{\widetilde{\mathcal S}}

% ---- 译者注体例（2026-11 校订版）----
\newif\ifshownotes
\shownotestrue
\newcommand{\znote}[2]{\ifshownotes\latexendnote{\textbf{【原书 p.~#1】}\ #2}\fi}
\renewcommand{\notesname}{译者注}
\renewcommand{\enoteformat}{%
  \rightskip=0pt \leftskip=3.0em \parindent=-3.0em \parskip=4pt plus 1pt
  \leavevmode\makebox[3.0em][l]{\normalfont\bfseries\theenmark.}\ignorespaces}
\renewcommand{\enoteheading}{%
  \cleardoublepage
  \chapter*{译者注}%
  \markboth{译者注}{译者注}%
  \addcontentsline{toc}{chapter}{译者注}%
  \noindent\small 本附录汇集全书译者注。方括号内为原书页码；每条注末给出依据或来源。}%

\title{Fukaya 范畴与 Picard–Lefschetz 理论}
\subtitle{中文译注本}
\author{Paul Seidel}
\date{}
\cover{images/cover.png}
\begin{document}
\maketitle
\frontmatter
\hypersetup{pageanchor=true}
\chapter*{前言}
\markboth{前言}{前言}
\addcontentsline{toc}{chapter}{前言}
\SourcePage{5}
Fukaya 范畴素有难以入门之名。这既是因为学习它需要大量背景知识——来自同调代数、辛几何和几何分析——也是因为它的基本定义本身颇为复杂。本书希望为有意了解 Fukaya 范畴、并可能将其用于自身工作的研究生和研究人员提供一份参考资料。我尽量把注意力集中在一组相当基础的课题上，力求尽可能精确地加以叙述，并补上某些早期文献中的空缺。这使本书的风格颇为简峭（因此，深入研读本书时，最好同时阅读一些讨论应用的论文）。本书的第二个目标，是介绍若干此前尚未发表的结果；这些结果把 Fukaya 范畴与 Lefschetz 纤维化理论联系起来。这一目标在最后几节中逐渐占据主导地位，正文也随之逐步转变为一部研究专著。

我广泛借鉴了许多人的工作，其中首先要提到的是 Fukaya、Kontsevich 和 Donaldson。Fukaya 的奠基性贡献，当然是把 $\Ai$ 结构引入辛几何。在代数方面，他率先使用了 Yoneda 嵌入的 $\Ai$ 版本，而我们将在全书中系统地采用这一工具。此外，若干几何思想，例如 Pin 结构的作用，以及借助带参数模空间构造 $\Ai$ 同态的方法，都来自 Fukaya、Oh、Ohta 和 Ono 的工作。Kontsevich 引入了 $\Ai$ 范畴的导出范畴，并建立了其中的大部分理论，特别是正合三角的内在刻画。他还猜想了 Dehn 扭转与扭转函子之间的关系，而这正是本书的主要结果之一。最后，在与 Barannikov 的合作中，他提出了 Lefschetz 纤维化的 Fukaya 范畴的一种构造；我们采用的定义表面上有所不同，但想来应当与之等价。Donaldson 的影响同样贯穿始终。除了在 Lefschetz 束方面的开创性工作之外，他还引入了匹配圈，并提出以此为出发点建立 Floer 上同调的组合公式；这一设想在本书中确实得到了部分实现。其他数学家也作出了重要贡献。例如，我们对 Picard–Lefschetz 理论的部分叙述反映了 Auroux 的观点。Smith 的一个结果——四维 Lefschetz 束中的消失圈必然充满纤维——对于提出这些圈可能“分裂生成”Fukaya 范畴的设想起了关键作用。此外，Fukaya–Smith 关于余切丛的工作，为有关 Lefschetz 纤维化的一些更大胆的想法提供了很好的试验场。我们处理横截性问题的方法，源于与 Lazzarini 的数次交谈。最后，Abouzaid 的建议大大改进了关于辛嵌入的讨论。

\SourcePage{6}
在机构方面，我要感谢 ETH Zürich 邀请我讲授作为本书基础的课程；感谢 NSF 提供经费支持（项目编号 DMS-9729992 和 DMS-0405516）；还要感谢 EMS Publishing House 接受书稿，并悉心推动本书出版。

\chapter*{译校说明}
\markboth{译校说明}{译校说明}
\addcontentsline{toc}{chapter}{译校说明}

\noindent\small
\textbf{底本。}本书译自 Paul Seidel 的讲义 \emph{Fukaya Categories and Picard--Lefschetz Theory}（欧洲数学会，2008；作者公开的 PDF 版本，共 334 页）。翻译底本的原书页码与 PDF 页码相差 8：PDF 第 9 页对应原书正文第 1 页。书中脚注与译者注所称“原书 p.~$N$”均指原书\textbf{印刷页码}，与 PDF 阅读器的页码不同。

\textbf{初版与本次校订。}初版（修订版 01）据英文原文译出，止于第 20 节 20.7 末，共 243 页，未加译者注。本次校订把全书 82 个章节文件按原书页序分成 12 段，逐段与英文原文对照通读，公式、交换图与符号约定另据原书页面图像逐式核对；遇原书排印可疑处，先核对页面图像，再决定改译文还是存疑。校订只改真正的误译、漏译、术语不统一与公式/指标错误，不作风格性改写；原书自身的笔误保留原译的正确写法，并在译者注与《原书笔误登记》中说明，不回改。

\textbf{译者注。}全书译者注统一排在书末“译者注”附录，正文只在相应句末加数字上标，不改动正文语义。每条注以“【原书 p.~$N$】”开头，先给页码，再给注文，注末给依据。依据分三级：A 级为原书前后文、作者他文、经典教材与原始论文；B 级为 nLab、Stacks Project、EMS 等可靠百科或数据库，均经实际打开核对并注明访问日期；此类注只作概念与背景说明，凡涉及具体定理归属或数值者，另举原书前后文、经典教材或原始论文为据（A 级），不单凭百科立论。无外部依据、仅由上下文推得者，明确写作“译者按（上下文推断，无外部依据）”。凡未能核实者不入注，宁缺勿滥。全书共 187 条译者注，A 级 125 条、B 级 51 条、C 级 11 条。

\textbf{术语。}全书术语以书末《术语表》为准。较初版有两处统一：Lefschetz thimble 由“指套”改为“顶针”（thimble 为顶针，finger cot 才是指套）；operad 由“运算胚”改为“操作子”。前者是初版沿用的日文译法，后者“胚”字在中文数学中另有 germ 的含义，均易致误，故改。

\textbf{使用建议。}初次阅读可先读“引言”与“读书报告”，把握定理 A、B 与推论 C 三条主线；再按章节顺序推进。译者注按需查阅，不必先读。书中少数记号（如用 \S{} 形符号表示 Koszul 符号）与原书一致，遇符号歧义时以注文说明为准。

\textbf{免责。}译者注是理解辅助，不是原文的一部分；凡注文与原文冲突，以原文为准。译文与注文中的疏漏由译者负责，欢迎指正。
\normalsize

\tableofcontents
\mainmatter
\chapter*{引言}
\markboth{引言}{引言}
\addcontentsline{toc}{chapter}{引言}
\label{intro}
\SourcePage{9}
\noindent\textbf{初步印象。}Fukaya 范畴是与一个辛流形相联系的代数结构。它编码了 Lagrangian 子流形的信息、任意两个此类子流形的相交方式，以及伪全纯圆盘（或多边形）的信息。把如此丰富的数据集中在一个对象之中，使 Fukaya 范畴引人入胜，却也让人有些望而生畏，从辛拓扑的角度来看尤其如此。在继续之前，我们想对这一点稍作展开。下面是一段面向非专家读者的非正式讨论，目的只是让读者对整体情形有一个大致印象；其中有不少并不严格的陈述，讨论范围也比本书其余部分略宽。

固定一个辛流形 $M^{2n}$。粗略地说，Fukaya 范畴的对象是 Lagrangian 子流形 $L^n\subset M$。此类子流形组成的空间是无限维的（局部模型为 $L$ 上闭 $1$ 形式组成的空间）。不过，Hamilton 同痕的 Lagrangian 子流形在 Fukaya 范畴中给出同构的对象，这就把局部自由度降到了第一 Betti 数 $b_1(L)$。事实上，在某些简单情形中，按 Hamilton 同痕分类的问题已经完全解决：例如曲面上的曲线，此时问题本质上属于拓扑；还有较不显然的例子，如 $S^2\times S^2$ 中的 Lagrangian 球面\bcite{75}，以及少数相关例子。另有一些一般性的现象也有助于这一问题：由于存在文献\bcite{60}意义下的障碍，并非所有 Lagrangian 子流形都会实际出现在 Fukaya 范畴中；有些虽会出现，却是平凡的，即零对象；而 Fukaya 范畴中的同构关系也可能比 Lagrangian 同痕更弱。第一点的一个例子，只需考察 $S^2$ 中的圆周：只有当其补集中的两个半球面积相等时，所需的障碍抵消才会发生；更有意思的例子见\bcite{31}、\bcite{32}。第二种现象发生在实点集 $\RP^n\subset\CP^n$（$n>1$）上，条件是它具有 Spin 结构（$n=4k-1$），并且 Fukaya 范畴采用 $\mathbb Q$ 系数。对于第三点，$\CP^2$ 中某些 Lagrangian 环面是候选例子，但据作者所知，其中尚无任何一个得到严格验证。有了这些现象，对于辛拓扑相对简单的流形，原则上计算整个 Fukaya 范畴就成为可行的目标。上面引用的论文，以及其他一些论文，例如\bcite{132}、\bcite{102}、\bcite{101}、\bcite{59}，都可视为朝这一方向推进的、更广泛的持续研究的一部分。然而，对于任何复杂程度稍高的辛流形（例如 $\CP^3$ 中次数至少为 $4$ 的复超曲面），直接分类其 Fukaya 范畴中的对象看来仍遥不可及。

也许只是我们把目标定得太高了。毕竟，即使在纯代数背景下，例如考虑有限维代数上的模，能够实现完全分类的情形也只占极少数\bcite{41}。
\SourcePage{10}
不过，两者有一个显著区别：对于两个明确给出的模，存在一个显然的算法，可以判定它们是否同构。更一般地，给定有限个模，很容易确定它们之间的态射以及态射的复合映射。相比之下，在 Fukaya 范畴中，态射是 Floer 上链群，复合映射则是这些群上的运算；二者都通过适当的伪全纯映射模空间来定义。在某些情形中，可以用几何方法处理这些空间：当 $M$ 是余切丛，而 Lagrangian 子流形 $L=\operatorname{graph}(df)$ 是截面时，一种“绝热极限”方法把问题化为有限维 Morse 理论中的问题\bcite{58}（这一方法的变体也已在若干相关情形中取得成功）；当 $M$ 是定义在实数域上的代数簇，而我们只考虑其实点集 $L=M_{\R}$ 时，可以运用代数几何工具；最后，当 $M$ 是曲面时，对伪全纯曲线的研究在很大程度上成为组合问题。但在这些情形之外，Floer 理论的计算历来依赖逐个例子处理的方法，无法事先保证成功。本书试图应对这一问题。希望本书与其他进展一道，能够使 Fukaya 范畴成为更易于运用的工具。

\noindent\textbf{研究策略。}鉴于上述困难，显然若想有所进展，就必须作出一些大幅简化。首先，我们只考虑正合辛流形 $M$，即辛形式 $\omega$ 是某个 $1$ 形式 $\theta$ 的外微分的流形。这当然意味着 $M$ 不可能是闭流形；不过，我们施加标准的凸性条件，使有限面积的伪全纯曲线感受不到这一缺憾。仿射代数簇提供了丰富的例子。Lagrangian 子流形 $L\subset M$ 中最简单的一类是正合的那些，即 $\theta|L$ 是某个函数的微分；我们将只处理这一类，事先排除其他所有情形。这一限制带来许多有用的技术性结论：例如，正合 Lagrangian 子流形不会有障碍，而且每一个都给出 Fukaya 范畴中的非零对象，其自同态环就是经典的上同调环（没有量子修正）。此外，正合性更根本的好处，是消除了 Fukaya 范畴对辛面积参数（或 Novikov 环参数）的依赖。镜像对称表明，这种依赖往往由非平凡的超越函数给出，而\bcite{108}、\bcite{57}等具体计算也证实了这一点。在我们的情形中，该参数被设为零（或无穷大，视解释而定）；这使问题从计数几何的角度看不那么有趣，却更便于计算。还应提到，最终可以尝试把仿射簇的 Fukaya 范畴作为研究其射影完备化的 Fukaya 范畴的踏脚石\bcite{129}。

\SourcePage{11}
为了使问题更接近经典同调代数，我们还将施加一个 Calabi–Yau 型条件（确切地说，是 $2c_1(M)=0$），并对 Lagrangian 子流形作出相应假设，从而使所有 Floer 群都具有 $\Z$ 分次。由此得到的设置可概括如下。Fukaya 范畴 $\F(M)$ 的对象称为\emph{正合 Lagrangian 膜}，是三元组 $L^{\br}=(L,\alpha^{\br},P^{\br})$，其中 $L$ 是 $M$ 的闭正合 Lagrangian 子流形，$\alpha^{\br}$ 是 $L$ 上称为\emph{分次}的实值函数，而 $P^{\br}$ 是 $L$ 上的一个 Pin 结构。态射空间是 Floer 上链群 $CF^*(L_0^{\br},L_1^{\br})$；它们是有限维分次向量空间（系数取在某个固定但任意的域 $\K$ 中；Pin 结构实际上只在 $\operatorname{char}(\K)\ne2$ 时才必需）。这些空间带有多重线性的复合映射 $\mu^d$，$d\ge1$，并由此构成一个 $\Ai$ 结构。\znote{3}{这里 $\mu^d$ 并非彼此独立：$d=1$ 是微分，$d=2$ 是同伦意义下的乘法，$d=3$ 度量乘法偏离结合律的程度，更高阶项度量下一层的偏差，全部项合起来满足二次的 $\Ai$ 结合律方程。用操作子（operad）的话说，$\Ai$ 操作子是结合操作子的一个 dg 分辨，$\Ai$ 代数就是其上的代数；Fukaya 情形之所以自然出现这种只在同伦意义下成立的结合律，是因为伪全纯曲线模空间只在边界处退化成低阶模空间。[来源：nLab \textquotedblleft A-infinity-algebra\textquotedblright，https:/\allowbreak /\allowbreak ncatlab.org/\allowbreak nlab/\allowbreak show/\allowbreak A-infinity-algebra（2026-11-15 访问）。]}定义需要作出辅助选择，但所得结果在拟同构意义下与这些选择无关。

我们的第二个主要想法来自 Kontsevich 的工作，即对 Fukaya 范畴作形式扩张。这些都是纯代数构造，可以应用于任意 $\Ai$ 范畴 $\A$。第一步是引入链复形的类似物，称为\emph{扭曲复形}；它们组成另一个包含 $\A$ 的 $\Ai$ 范畴 $\Tw\A$。\znote{3}{扭曲复形是链复形的推广：在 dg 情形，一个扭曲复形由一族除有限个外为零的分次对象 $\{E_i\}$ 连同次数为 $i-j+1$ 的态射 $q_{ij}:E_i\to E_j$ 组成，并满足 $dq_{ij}+\sum_k q_{kj}q_{ik}=0$；$q$ 就是所谓\textquotedblleft 扭\textquotedblright，只保留其中一项时便是普通链复形与它的微分。只留一项时 $q$ 给出映射锥，故 $\Tw\A$ 自动含有锥；而令 $q=0$ 便把 $\A$ 嵌入其中。[来源：nLab \textquotedblleft twisted complex\textquotedblright，https:/\allowbreak /\allowbreak ncatlab.org/\allowbreak nlab/\allowbreak show/\allowbreak twisted+\allowbreak complex（2026-11-15 访问）。]}这一范畴含有态射的映射锥，因此具有我们所谓的三角化的 $\Ai$ 范畴结构；其上同调层次的范畴 $H^0(\Tw\A)$，也称为\emph{导出范畴} $D(\A)$，是经典意义下的三角范畴。第二个有用步骤，是取导出范畴的\emph{分裂闭包}（或 Karoubi 完备化）$D^{\pi}(\A)$，把所有幂等自同态的像作为形式直和项加入进来；在 $\Ai$ 层次上也有相应构造，记为 $\Pi(\Tw\A)$。\znote{3}{Karoubi 完备化（幂等完备化）把所有幂等态射 $e:X\to X$ 形式地分裂，得到对象 $(X,e)$ 以及 $X\to(X,e)\to X$；它是使一切幂等态射都可分裂的最小扩张，对任意具备该性质的范畴满足普遍性。$\Pi(\Tw\A)$ 正是这一构造：$\Tw\A$ 只给出锥，再分裂幂等才得到直和项，于是三角闭包与幂等闭包一并完成。$\Tw\A$ 与 $\Pi(\Tw\A)$ 的差别正在于此。[来源：nLab \textquotedblleft Karoubi envelope\textquotedblright，https:/\allowbreak /\allowbreak ncatlab.org/\allowbreak nlab/\allowbreak show/\allowbreak Karoubi+\allowbreak envelope（2026-11-15 访问）。]}每次扩张都会使范畴更容易描述：这一说法表面上似乎矛盾，但许多读者会很清楚它为什么成立。作为说明，考虑\emph{分裂生成元}的概念：若从某一组对象出发，反复取映射锥并分裂出直和项，可以构造出 $\Pi(\Tw\A)$ 的所有对象，就称这组对象分裂生成 $\Pi(\Tw\A)$。只要 $\B\subset\Pi(\Tw\A)$ 是由这样一组分裂生成元构成的全 $\Ai$ 子范畴，$\Pi(\Tw\B)$ 就与 $\Pi(\Tw\A)$ 拟等价。因此，如果只关心 $\Pi(\Tw\A)$，实际上确定一个这样的 $\B$ 就够了。要记住，分裂生成是一个很弱的性质；例如，它并不意味着 Grothendieck 群 $K_0(D^{\pi}(\A))$ 由这些分裂生成元的类所张成。因此，期待找到有限个 Lagrangian 子流形，分裂生成 $\Pi(\Tw\F(M))$，并非不合理；我们将看到，在有利的情形中确实可以做到这一点。

谈到 Fukaya 范畴时使用“计算”一词，意味着必须以某种组合方式编码其基础几何结构，即辛流形和 Lagrangian 子流形。为此，我们将系统地运用 Picard–Lefschetz 理论。一般而言，这一方法在辛几何中的重要性在很大程度上来自 Donaldson 的工作\bcite{38}；这里我们只使用该理论最初等的部分，即单值变换的辛性质，以及 Lefschetz 顶针和消失圈的 Lagrangian 性质。
\SourcePage{12}
为了利用这些性质，我们将依照 Kontsevich 的一个想法，引入 Lefschetz 纤维化的 Fukaya 范畴；从某种意义上说，它介于总空间的通常 Fukaya 范畴与光滑纤维的 Fukaya 范畴之间。随后，我们将以此建立一套按维数归纳进行计算的机制；这套机制最初以猜想的形式提出于\bcite{126}。

\noindent\textbf{计算 Fukaya 范畴。}为简化陈述，我们只考虑代数簇的情形。具体地，设 $X\subset\CC^N$ 是一个 $(n+1)$ 维仿射代数簇，假设它光滑，并且在无穷远处也光滑。后一条件意味着它的闭包 $\overline X\subset\CP^N$ 仍然光滑，且与无穷远超平面 $\CP^{N-1}=\CP^N\setminus\CC^N$ 横截相交。给 $X$ 配备 Fubini–Study Kähler 形式的限制，它便成为正合辛流形。此外，我们假设典范丛 $K_{\overline X}$ 对某个 $d\in\Z$ 同构于 $\mathcal O_{\overline X}(-d)$，从而 $X$ 是 Calabi–Yau 的。一般的仿射超平面截面 $X^1=X\cap\CC^{N-1}$ 也继承所有这些性质。考虑系数取在某个域 $\K$ 中的 Fukaya 范畴 $\F(X)$ 和 $\F(X^1)$。我们先讨论 $X^1$，因为在这一情形中得到的结果略强一些。

\begingroup\renewcommand{\theoremname}{定理 A}
\begin{theorem*}\phantomsection\label{st:A}
假设 $d\ne2$、$n\ge1$ 且 $\operatorname{char}(\K)\ne2$。那么，$\Pi(\Tw\F(X^1))$ 可以由对 $X$ 应用 Picard–Lefschetz 理论所得的组合数据计算出来。\znote{4}{定理 A 的假设在原书印刷页 4 作 $n\ge1$；紧接其后的解释段（印刷页 5）却说 $n>1$ 排除了最低的非平凡维数，并称 $n=1$ 的情形“只是为方便起见”才排除、可用类似方法处理。可见 $n>1$ 是证明叙述中的便利假设，不是定理本身的限制；第 19 节的归纳证明也按 $n\ge1$ 陈述。译文从原书定理陈述，取 $n\ge1$，不将 $n=1$ 排除在定理之外。[来源：原书 p.~4 定理 A、p.~5 说明与 §19 证明；属原书表述不统一，无外部依据。]}
\end{theorem*}\endgroup

这一陈述需要作些解释。首先，对某个 $\Ai$ 范畴 $\A$，称 $\Pi(\Tw\A)$ 可计算，含义如下：可以明确写出另一个在很强意义下有限的 $\Ai$ 范畴 $\B$——其对象只有有限个，态射组成有限维向量空间，且当阶数 $d\gg0$ 足够大时，复合映射 $\mu^d_{\B}$ 全部为零——使 $\Pi(\Tw\B)$ 与 $\Pi(\Tw\A)$ 成为拟等价的 $\Ai$ 范畴。特别地，这意味着 $D^{\pi}(\B)\cong D^{\pi}(\A)$，作为经典三角范畴同构。注意，若合理选择系数域 $\K$（例如有限域或 $\mathbb Q$），$\B$ 确实只包含有限量的信息。可以将这些信息输入计算机程序，通过明确构造扭曲复形及其幂等自同态，生成 $D^{\pi}(\B)$ 中对象的同构类的完整无限列表，并同时给出态射、复合映射和正合三角。当然，仅凭这一点本身，并不能解决有关 $D^{\pi}(\B)$ 的任何重大结构问题。至于“Picard–Lefschetz 理论所产生的组合数据”的含义，精确解释留待\sr{section:19}{第 19 节}；大致想法是考察从 $X$ 到 $\CC$ 和 $\CC^2$ 的线性映射，并用辫单值变换来编码这些映射的分支数据。同样的过程可以应用于 $X^1$，并反复进行，直到维数降至零。给定定义 $X$ 的明确方程，可以借助消元理论，以算法方式提取这些数据
\SourcePage{13}
（不过，除了最简单的例子之外，计算复杂度通常都会高到难以承受）。

还需要简要考察该定理的限制。条件 $n>1$ 排除了最低的非平凡维数，即 $X^1$ 为 Riemann 曲面的情形。这实际上只是为了方便：Riemann 曲面的情形可用相似、事实上还略简单的方法处理，但需要作一些特别调整，这会使讨论复杂化。条件 $\operatorname{char}(\K)\ne2$ 是我们证明方法中的技术性产物，因为证明使用了二重分歧覆盖和 $\Z/2$ 作用；很可能存在其他方法，可以取消这一限制。另一个要求 $d\ne2$ 则略为本质一些，我们把讨论留到后面（见\stref{笔记}{19.8}）。

如前所述，对于 $\F(X)$ 本身，我们得到的结果略弱一些。这并不令人意外：例如，目前没有已知方法可以判定 $X$ 是否含有闭正合 Lagrangian 子流形，从而判定 $\F(X)$ 及其导出范畴是否非平凡。与此相对，在 $X^1$ 的情形中，可以把消失圈作为对象的自然来源。我们所能得到的最好结果如下。

\begingroup\renewcommand{\theoremname}{定理 B}
\begin{theorem*}\phantomsection\label{st:B}
假设 $\operatorname{char}(\K)\ne2$。那么，由对 $X$ 应用 Picard–Lefschetz 理论所得的组合数据，可以构造一个 $\Ai$ 范畴，其中包含一个与 $\F(X)$ 拟等价的全子范畴。
\end{theorem*}\endgroup

这里的 $\Ai$ 范畴具有 $\Tw\B$ 的形式，其中 $\B$ 在与前面相同的意义下是有限的。事实上，按构造 $\B$ 将是有向的，这带来几个值得注意的结论。首先，$\Tw\B$ 已经是分裂闭的，因此它包含的不仅是 $\F(X)$，而是整个 $\Pi(\Tw\F(X))$。其次，有向 $\Ai$ 范畴与态射空间有限维的 dg 范畴拟同构。由于扭曲复形范畴继承这一性质，便有如下推论。

\begingroup\renewcommand{\corollaryname}{推论 C}
\begin{corollary*}\phantomsection\label{st:C}
在\sr{st:B}{定理 B} 的条件下，$\F(X)$ 与一个态射空间有限维的 dg 范畴拟同构。
\end{corollary*}\endgroup

\sr{st:B}{定理 B} 在实际中的用途，取决于我们对映射 $\F(X)\longrightarrow\Tw\B$ 的几何意义了解得有多充分。对于 $X$ 中最简单的一类 Lagrangian 子流形，即能够表示为某个超平面截面 Lefschetz 束的匹配圈的 Lagrangian 球面，这些问题很容易处理，因而确实可以得到计算 Floer 上同调群的算法（严格地说，出于与\sr{st:A}{定理 A} 相同的原因，这里应排除 $n=1$ 的情形）。对于更一般的 Lagrangian 子流形，至少就目前而言，这一陈述还主要具有理论性质。

上述结果的证明将在\sr{section:19}{第 19 节}给出；证明涉及本书各处引入的大部分材料。作为一个自然的副产物，我们建立了消失圈的 Hurwitz 移动与有向 Fukaya 范畴的变异之间的一般关系。
\SourcePage{14}
事实上，我们将给出这一关系的两个证明（\stref{定理}{17.20}和\stref{定理}{18.24}；文献\bcite{127}中最初设想的证明与这两个证明都不相同，迄今仍未发表）。另一个有趣结论，是存在一个计算 Lefschetz 纤维化总空间中 Lagrangian 子流形的 Floer 上同调的谱序列（\stref{推论}{18.27}）。当然，与任何谱序列一样，这只是把问题分解为可见的部分，即起始项，以及隐藏的高阶信息，即微分。在这里，微分表达了这些 Lagrangian 子流形如何由 Lefschetz 顶针构成；“构成”一词应作代数解释，即 Fukaya 范畴中的 Postnikov 分解。

除了结尾处的一处例外，本书并不讨论例子或应用。所幸已有文献弥补了这一不足（若想获得一些线索，可参看\sr{chapter:III}{第三章}开头）。此外，任何认真研究 Fukaya 范畴的读者，最终都必须研读奠基性的文献\bcite{60}；其中讨论的许多几何与代数问题，远远超出了本书建立的框架。

\chapter{\texorpdfstring{$A_\infty$}{A∞} 范畴}
\label{chapter:I}
\SourcePage{15}
本章介绍 Fukaya 范畴背后的同调代数。我们先快速回顾 $\Ai$ 结构。这绝不是对该理论的完整阐述。基本上，唯一详细讨论的方面是恒等态射，因为它对 Fukaya 范畴而言可能颇为棘手。此外，我们有意使形式体系保持具体。这意味着我们不使用操作子（operad）或非交换几何，尽管这两种语言中的任何一种，都可以帮助理清错综复杂的公式。\znote{7}{操作子（operad）是用\textquotedblleft $n$ 元运算的合成法则\textquotedblright 编码代数结构的语言：$\Ai$ 操作子是结合操作子的一个 dg 分辨，其上的代数恰为 $\Ai$ 代数，$\mu^d$ 就是 $d$ 元运算。本书不用这套语言，但如\stref{命题}{1.12}那样的递归构造，用操作子说就是\textquotedblleft 把分辨搬到范畴上\textquotedblright，公式里的树形求和因此有了统一解释。[来源：nLab \textquotedblleft A-infinity-algebra\textquotedblright，https:/\allowbreak /\allowbreak ncatlab.org/\allowbreak nlab/\allowbreak show/\allowbreak A-infinity-algebra（2026-11-15 访问）。]}关于这些材料，还有风格各异的其他参考文献。这里给出几点建议：一些较早的基础论文，例如\bcite{70}、\bcite{62}，即使重点与本书稍有不同，仍然值得一读。在较新的文献中，Keller 的工作\bcite{82}、\bcite{83}、\bcite{84}既对基础材料作了易于理解的说明，也讨论了它们在代数和表示论问题中的应用。Lefèvre 的学位论文\bcite{95}采用同伦代数语言，从相近的观点作了全面阐述。Fukaya 及其合作者写过多份专门面向辛几何的论述\bcite{51}、\bcite{52}、\bcite{53}、\bcite{58}、\bcite{60}、\bcite{56}、\bcite{55}。最后三篇文献的一般性远远超过我们的需要，但愿意顾及这一点的读者会发现，许多论证可以简化，适用于我们较为有限的框架。最后，Kontsevich 的论文和讲义\bcite{90}、\bcite{88}、\bcite{89}中也蕴藏着丰富的见解。

本章其余部分的总主题，是 $\Ai$ 范畴中对象的构造与运用。我们将讨论映射锥和正合三角，以及将对象分裂成直和项的问题。还会考察一种更具体的构造，即所谓一个对象沿另一个对象的扭转。此前，这一构造在文献中出现于两个不同背景。其一是例外集与变异的概念；它最初是为在代数几何中构造刚性向量丛而产生的：参看\bcite{67}和论文集\bcite{117}。这一理论逐步推广到一般三角范畴（不同方面见\bcite{65}、\bcite{26}、\bcite{27}、\bcite{28}）；文献\bcite{66}是一篇有用的新近综述。扭转构造第二次、也是较近的一次出现，与球面对象有关；在这一背景下，它给出导出范畴的自同构\bcite{134}。向 $\Ai$ 范畴的推广由 Kontsevich 开创，我们将广泛借鉴他的思想。除了建立一般形式体系外，我们还将推出若干结论，特别是\stref{推论}{5.8}和\stref{命题}{5.17}；后面把这些材料与辛几何及 Lefschetz 纤维化联系起来时，它们会派上用场。本章最后以一段略带技术性的讨论收尾，简要介绍具有 $\Z/2$ 对称性的 $\Ai$ 范畴。

\section{定义}\label{section:1}
\SourcePage{16}
\subsection{范畴}\label{sec:1a}
\begin{definition*}
固定系数域 $\K$。我们假定所有范畴都是 $\K$ 上的线性小范畴。一个\emph{非含幺 $\Ai$ 范畴} $\A$ 由对象集 $\Ob\A$、对每对对象给出的分次向量空间 $\hom_{\A}(X_0,X_1)$，以及对每个阶数 $d\ge1$ 给出的复合映射组成：
\begin{equation}\label{eq:1.1}
\mu^d_{\A}:\hom_{\A}(X_{d-1},X_d)\otimes\cdots\otimes\hom_{\A}(X_0,X_1)\longrightarrow\hom_{\A}(X_0,X_d)[2-d].
\end{equation}
照例，$[k]$ 表示将向量空间的分次向下平移 $k\in\Z$。映射\eref{1.1}必须满足以下二次的 $\Ai$ 结合律方程：
\begin{equation}\label{eq:1.2}
\sum_{m,n}(-1)^{\sgnmark_n}\mu^{d-m+1}_{\A}\bigl(a_d,\ldots,a_{n+m+1},\mu^m_{\A}(a_{n+m},\ldots,a_{n+1}),a_n,\ldots,a_1\bigr)=0,
\end{equation}
其中对所有可能的项求和，即 $1\le m\le d$、$0\le n\le d-m$，符号由 $\sgnmark_n=|a_1|+\cdots+|a_n|-n$ 决定。同一和式还会出现在许多其他符号表达式中，因此从现在起，我们将它规定为通用记号。\znote{8}{记号 $\sgnmark_n$ 是纯粹的 Koszul 记账：把元素逐个移过前面的态射，每跨过一个次数为 $e$ 的态射便得因子 $(-1)^{e}$；但 $\mu^d$ 的输入已带 $[2-d]$ 的次数平移，实际出现的因而是 $|a_i|-1$ 而非 $|a_i|$，$n$ 项累加便得 $|a_1|+\cdots+|a_n|-n$。后文凡出现同一记号都可直接代入此式，不必重新推导。[译者按（上下文推断，无外部依据）：由\eref{1.1}的 $[2-d]$ 平移与 Koszul 符号规则相加即得。]}
\end{definition*}

\begin{definition*}
对于任意非含幺 $\Ai$ 范畴 $\A$，可以赋予它一个\emph{上同调范畴} $H(\A)$。其对象与 $\A$ 相同；态射空间为上同调群 $H(\hom_{\A}(X_0,X_1),\mu^1_{\A})$；复合则定义为
\begin{equation}\label{eq:1.3}
[a_2]\cdot[a_1]=(-1)^{|a_1|}[\mu^2_{\A}(a_2,a_1)].
\end{equation}
除了可能缺少恒等态射外，这是一个通常的线性分次范畴。\znote{8}{\textquotedblleft 非含幺\textquotedblright 只是说定义里不把恒等态射列为数据，并不禁止它存在。原书后面区分严格含幺（单位元在链层次上与全部 $\mu^d$ 相容）与上同调含幺（单位元只需在 $H$ 中起作用）等情形，并在讲完恒等态射之后约定（§2h）一律假定 c-含幺。阅读后文时务必先确认当前处在哪种假设下，否则\textquotedblleft 同构\textquotedblright、\textquotedblleft 拟等价\textquotedblright 的说法会失准。[来源：Seidel, Fukaya categories and Picard–Lefschetz theory, §2h（术语约定）及 §2 对恒等态射的讨论。]}还有一个变体，即只保留零次上同调的子范畴 $H^0(\A)\subset H(\A)$。同样地，如果一个非含幺 $\Ai$ 范畴 $\A$ 满足对所有 $d>2$ 都有 $\mu^d_{\A}=0$，那么它与一个非含幺\emph{微分分次范畴}（dg 范畴）是同一回事。更精确地，相关 dg 范畴的微分为 $\partial a=(-1)^{|a|}\mu^1_{\A}(a)$，乘法为 $a_2a_1=(-1)^{|a_1|}\mu^2_{\A}(a_2,a_1)$，与\eref{1.3}一致。
\end{definition*}

\begin{definition*}
非含幺 $\Ai$ 范畴的\emph{反范畴}定义如下：令 $\Ob\A^{\mathrm{opp}}=\Ob\A$、$\hom_{\A^{\mathrm{opp}}}(X_0,X_1)=\hom_{\A}(X_1,X_0)$，并将复合映射中各项的顺序反转：
\[
\mu^d_{\A^{\mathrm{opp}}}(a_d,\ldots,a_1)=(-1)^{\sgnmark_d}\mu^d_{\A}(a_1,\ldots,a_d).
\]
在上同调层次上，这给出通常的反范畴（不过不带 Koszul 符号），即 $H(\A^{\mathrm{opp}})=H(\A)^{\mathrm{opp}}$。\znote{8}{反范畴公式里的 $(-1)^{\sgnmark_d}$ 只是把被反转的输入顺序重新记账所需的 Koszul 符号，并非新的数学内容：到了上同调层次，$\mu^2_{\A^{\mathrm{opp}}}$ 恰给出反向复合，符号自动抵消。后文凡见 $\A^{\mathrm{opp}}$（例如 $\nufun(\A^{\mathrm{opp}},\Ch)$）都指这一约定，与取对偶无关。[译者按（上下文推断，无外部依据）：把\eref{1.2}用于 $\A^{\mathrm{opp}}$，再与 $H(\A)^{\mathrm{opp}}$ 的复合式逐项比对即得。]}
\end{definition*}

\begin{note}\label{st:1.1}
关于术语和约定，补充几点说明：(i) “非含幺”其实应当理解为“不要求含幺”。(ii) 假定范畴是小范畴，偶尔会造成形式上的麻烦，例如在定义扭曲复形时；我们把必要的集合论处理留给读者。(iii) 使用没有恒等态射的范畴时，需要小心：
\SourcePage{17}
例如，此时没有合理的对象同构概念；全忠实函子仍有意义，范畴的同构也有意义，但范畴的等价却没有意义，至少不能直接按通常方式定义。(iv) 我们的符号遵循“条形约定”，这是代数学家最常使用的约定；\eref{1.1}中递减的编号，试图调和态射复合的标准约定与圆周上标记点同样标准的循环次序，见\secref{9c}。
\end{note}

\begin{note}\label{st:1.2}
阶数 $d>2$ 的复合不是链映射。不过，它们仍以 Massey 积的形式，部分体现在上同调层次上。设有态射 $[a_k]\in\Hom_{H(\A)}(X_{k-1},X_k)$，$k=1,2,3$，并且两个乘积 $[a_3]\cdot[a_2]$ 与 $[a_2]\cdot[a_1]$ 都为零。选择链层次的代表元 $a_k$，再取 $h_k\in\hom_{\A}(X_{k-1},X_{k+1})$，$k=1,2$，使 $\mu^1_{\A}(h_k)=\mu^2_{\A}(a_{k+1},a_k)$。那么
\[
c=\mu^3_{\A}(a_3,a_2,a_1)-\mu^2_{\A}(h_2,a_1)-\mu^2_{\A}(a_3,h_1)\in\hom_{\A}(X_0,X_3)
\]
是次数为 $|a_1|+|a_2|+|a_3|-1$ 的上闭链，而类
\begin{equation}\label{eq:1.4}
[c]\in\frac{\Hom_{H(\A)}(X_0,X_3)}{[a_3]\cdot\Hom_{H(\A)}(X_0,X_2)+\Hom_{H(\A)}(X_1,X_3)\cdot[a_1]}
\end{equation}
只依赖于最初给定的 $[a_k]$。通常记作 $\langle[a_3],[a_2],[a_1]\rangle=(-1)^{|a_2|}[c]$。
\end{note}

\subsection{函子}\label{sec:1b}
\begin{definition*}
两个非含幺 $\Ai$ 范畴 $\A,\B$ 之间的\emph{非含幺 $\Ai$ 函子}，由映射 $\F:\Ob\A\to\Ob\B$ 以及对所有 $d\ge1$ 给出的多重线性映射组成：
\begin{equation}\label{eq:1.5}
\F^d:\hom_{\A}(X_{d-1},X_d)\otimes\cdots\otimes\hom_{\A}(X_0,X_1)\longrightarrow\hom_{\B}(\F X_0,\F X_d)[1-d],
\end{equation}
其中 $X_0,\ldots,X_d\in\Ob\A$，并满足以下多项式方程：
\begin{equation}\label{eq:1.6}
\begin{aligned}
&\sum_r\sum_{s_1,\ldots,s_r}\mu^r_{\B}\bigl(\F^{s_r}(a_d,\ldots,a_{d-s_r+1}),\ldots,\F^{s_1}(a_{s_1},\ldots,a_1)\bigr)\\
&\quad=\sum_{m,n}(-1)^{\sgnmark_n}\F^{d-m+1}\bigl(a_d,\ldots,a_{n+m+1},\mu^m_{\A}(a_{n+m},\ldots,a_{n+1}),a_n,\ldots,a_1\bigr).
\end{aligned}
\end{equation}
左侧对所有 $r\ge1$ 以及分拆 $s_1+\cdots+s_r=d$ 求和。
\end{definition*}

\begin{definition*}
它诱导一个通常的线性分次非含幺函子 $H(\F):H(\A)\to H(\B)$，在态射上的作用为 $[a]\mapsto[\F^1(a)]$。若 $H(\F)$ 是范畴同构，就称 $\F$ 为\emph{拟同构}；若 $H(\F)$ 全忠实，就称 $\F$ 为\emph{上同调全忠实}的。
\end{definition*}

非含幺 $\Ai$ 函子可按如下方式复合：
\begin{equation}\label{eq:1.7}
\begin{aligned}
(\G\circ\F)^d(a_d,\ldots,a_1)
=\sum_r\sum_{s_1,\ldots,s_r}\G^r\bigl(\F^{s_r}(a_d,\ldots,a_{d-s_r+1}),\ldots,\F^{s_1}(a_{s_1},\ldots,a_1)\bigr),
\end{aligned}
\end{equation}
\SourcePage{18}
其中求和范围与\eref{1.6}相同。复合严格满足结合律，而恒等函子——$\F X=X$、$\F^1=\id$（作为各态射空间的自同态），所有高阶项为零——是复合的单位元。

\subsection{形式微分同胚}\label{sec:1c}
\begin{definition*}
下面从一个稍有不同的角度看待某一特殊类的 $\Ai$ 函子。设有对象集 $\Ob\A$，并对每对对象给定分次向量空间 $\hom_{\A}(X_0,X_1)$，但暂时还没有任何复合映射。从 $\A$ 到自身的\emph{形式微分同胚} $\Phi$，是任意一列形如\eref{1.5}的映射，唯一的条件是每个 $\Phi^1$ 都是 $\hom_{\A}(X_0,X_1)$ 的线性自同构。公式\eref{1.7}使这些形式微分同胚组成一个群。现在设映射 $\mu^{\bullet}_{\A}$ 使 $\A$ 成为非含幺 $\Ai$ 范畴。给定 $\Phi$，可以递归求解\eref{1.6}，证明存在唯一的另一组 $\Ai$ 结构 $\mu^{\bullet}_{\widetilde{\A}}$（其基础向量空间相同），使 $\Phi:\A\to\widetilde{\A}$ 成为非含幺 $\Ai$ 函子。这就在所有可能的 $\mu^{\bullet}_{\A}$ 组成的集合上定义了形式微分同胚群的作用；我们记 $\widetilde{\A}=\Phi_*\A$。
\end{definition*}

\begin{note}\label{st:1.3}
“形式微分同胚”这一名称来自非交换几何；在那里，可以把 $\Ai$ 结构看作分次非交换仿射空间上次数为 $-1$ 的向量场，而上述作用就是坐标变换。
\end{note}

\subsection{自然变换}\label{sec:1d}
\begin{definition*}
非含幺 $\Ai$ 函子 $\F:\A\to\B$ 本身构成一个非含幺 $\Ai$ 范畴 $\Q=\nufun(\A,\B)$ 的对象。该范畴态射链空间中的元素 $T\in\hom^g_{\Q}(\F_0,\F_1)$，是一列 $(T^0,T^1,\ldots)$，其中每个 $T^d$ 都是对所有 $(X_0,\ldots,X_d)$ 给出的一族多重线性映射
\[
\hom_{\A}(X_{d-1},X_d)\otimes\cdots\otimes\hom_{\A}(X_0,X_1)\longrightarrow\hom_{\B}(\F_0X_0,\F_1X_d)[g-d].
\]
特别地，$T^0$ 是一族 $\hom^g_{\B}(\F_0X,\F_1X)$ 中的元素，每个 $\A$ 的对象 $X$ 对应一个。沿用 Fukaya 的术语，称这样的 $T$ 为从 $\F_0$ 到 $\F_1$ 的、次数为 $g=|T|$ 的\emph{预自然变换}。边界算子为
\end{definition*}

\begin{equation}\label{eq:1.8}
\begin{aligned}
&\mu^1_{\Q}(T)^d(a_d,\ldots,a_1)\\
&=\sum_{r,i}\sum_{s_1,\ldots,s_r}(-1)^{\dagger}\mu^r_{\B}\bigl(
\F_1^{s_r}(a_d,\ldots,a_{d-s_r+1}),\ldots,\F_1^{s_{i+1}}(\ldots,a_{s_1+\cdots+s_i+1}),\\
&\hspace{4em}T^{s_i}(a_{s_1+\cdots+s_i},\ldots,a_{s_1+\cdots+s_{i-1}+1}),
\F_0^{s_{i-1}}(a_{s_1+\cdots+s_{i-1}},\ldots),\ldots,\F_0^{s_1}(a_{s_1},\ldots,a_1)\bigr)\\
&\quad-\sum_{m,n}(-1)^{\sgnmark_n+|T|-1}T^{d-m+1}\bigl(a_d,\ldots,a_{n+m+1},
\mu^m_{\A}(a_{n+m},\ldots,a_{n+1}),a_n,\ldots,a_1\bigr).
\end{aligned}
\end{equation}
第一重求和取遍 $1\le i\le r$ 以及分拆 $s_1+\cdots+s_r=d$，其中 $s_i$ 可以为零；并且
\[
\dagger=(|T|-1)(|a_1|+\cdots+|a_{s_1+\cdots+s_{i-1}}|-s_1-\cdots-s_{i-1}).
\]
\SourcePage{19}
\begin{definition*}
上闭链 $\mu^1_{\Q}(T)=0$ 才是真正的自然变换，而上边界给出了它们之间的链同伦概念。
\end{definition*}

对于 $T_1\in\hom_{\Q}(\F_0,\F_1)$、$T_2\in\hom_{\Q}(\F_1,\F_2)$，复合公式要简单得多：
\begin{equation}\label{eq:1.9}
\begin{aligned}
&\mu^2_{\Q}(T_2,T_1)^d(a_d,\ldots,a_1)\\
&=\sum_{r,i,j}\sum_{s_1,\ldots,s_r}(-1)^{\circ}\mu^r_{\B}\bigl(
\F_2^{s_r}(a_d,\ldots,a_{d-s_r+1}),\ldots,\F_2^{s_{j+1}}(\ldots),\\
&\qquad T_2^{s_j}(a_{s_1+\cdots+s_j},\ldots,a_{s_1+\cdots+s_{j-1}+1}),
\F_1^{s_{j-1}}(\ldots),\ldots,\F_1^{s_{i+1}}(\ldots),\\
&\qquad T_1^{s_i}(a_{s_1+\cdots+s_i},\ldots,a_{s_1+\cdots+s_{i-1}+1}),
\F_0^{s_{i-1}}(\ldots),\ldots,\F_0^{s_1}(a_{s_1},\ldots,a_1)\bigr),
\end{aligned}
\end{equation}
其中
\[
\circ=\sum_{k=1}^{s_1+\cdots+s_{j-1}}(|T_2|-1)(|a_k|-1)
+\sum_{k=1}^{s_1+\cdots+s_{i-1}}(|T_1|-1)(|a_k|-1).
\]
注意，这里不再出现 $\mu^{\bullet}_{\A}$。$d>2$ 时的 $\mu^d_{\Q}$ 的公式，完全遵循\eref{1.9}的同一模式。\znote{11}{两个符号 $\dagger$ 与 $\circ$ 都是纯记账：在\eref{1.8}里要把 $T$ 所在的块移到最右端，须跨过其右侧共 $s_1+\cdots+s_{i-1}$ 个输入，每个输入贡献因子 $(-1)^{(|T|-1)(|a_k|-1)}$，累乘即得 $\dagger$；\eref{1.9}要先后把 $T_2$ 与 $T_1$ 各移一次，于是符号是两段求和。看清这一点，就不必每次重算。[译者按（上下文推断，无外部依据）：由 $T$ 的块次数 $|T|-1$ 与各输入的次数平移逐项相消得出。]}

给定 $\Ai$ 函子 $\F_0,\F_1:\A\to\B$ 之间的自然变换 $T$，考虑元素 $[T_X^0]\in\Hom_{H(\B)}(\F_0X,\F_1X)$。对于所有 $[a]\in\Hom_{H(\A)}(X_0,X_1)$，它们满足自然性条件
\[
[T^0_{X_1}]\cdot[\F_0^1(a)]=(-1)^{|a|\cdot|T|}[\F_1^1(a)]\cdot[T^0_{X_0}].
\]
因此，它们组成上同调函子 $H(\F_0),H(\F_1)$ 之间的自然变换，记作 $H(T)$。更形式化地，用
\[
\operatorname{Nu\text{-}Fun}(H(\A),H(\B))
\]
表示由非含幺线性分次函子 $H(\A)\to H(\B)$ 及其自然变换组成的范畴。那么，赋值 $\F\mapsto H(\F)$、$T\mapsto H(T)$ 定义了非含幺函子
\begin{equation}\label{eq:1.10}
H(\nufun(\A,\B))\longrightarrow\operatorname{Nu\text{-}Fun}(H(\A),H(\B)).
\end{equation}

\clearpage
\subsection{复合函子}\label{sec:1e}
与固定的 $\G:\A\to\B$ 作左复合和右复合，对任意非含幺 $\Ai$ 范畴 $\C$，定义了两个非含幺 $\Ai$ 函子
\[
\begin{aligned}
\LG_{\G}&:\nufun(\C,\A)\longrightarrow\nufun(\C,\B),\\
\RG_{\G}&:\nufun(\B,\C)\longrightarrow\nufun(\A,\C).
\end{aligned}
\]
它们在对象上的作用已作过说明。在预自然变换上的一阶映射，对于 $T\in\hom_{\nufun(\C,\A)}(\F_0,\F_1)$ 为
\begin{equation}\label{eq:1.11}
\begin{aligned}
&(\LG^1_{\G}T)^d(a_d,\ldots,a_1)\\
&=\sum_{r,i}\sum_{s_1,\ldots,s_r}(-1)^{\dagger}\G^r\bigl(
\F_1^{s_r}(a_d,\ldots,a_{d-s_r+1}),\ldots,\F_1^{s_{i+1}}(\ldots),\\
&\qquad T^{s_i}(a_{s_1+\cdots+s_i},\ldots,a_{s_1+\cdots+s_{i-1}+1}),
\F_0^{s_{i-1}}(\ldots),\ldots,\F_0^{s_1}(a_{s_1},\ldots,a_1)\bigr),
\end{aligned}
\end{equation}
\SourcePage{20}
其中 $\dagger$ 与\eref{1.8}相同；而对于 $T\in\hom_{\nufun(\B,\C)}(\F_0,\F_1)$，有
\[
\begin{aligned}
(\RG^1_{\G}T)^d(a_d,\ldots,a_1)
=\sum_r\sum_{s_1,\ldots,s_r}T^r\bigl(\G^{s_r}(a_d,\ldots,a_{d-s_r+1}),\ldots,\G^{s_1}(a_{s_1},\ldots,a_1)\bigr).
\end{aligned}
\]
高阶项 $\RG^d_{\G}$，$d\ge2$，全部为零。$\LG_{\G}$ 则并非如此，但其相应公式是\eref{1.11}的直接推广，例如见\bcite{55}，公式 (8.42)。特别地，其中不出现 $\Ai$ 范畴 $\A,\B,\C$ 的结构映射 $\mu$。复合函子本身也可复合，其行为与预期一致：
\begin{equation}\label{eq:1.12}
\RG_{\G_1}\circ\RG_{\G_2}=\RG_{\G_2\circ\G_1},\qquad
\LG_{\G_1}\circ\LG_{\G_2}=\LG_{\G_1\circ\G_2},\qquad
\LG_{\G_1}\circ\RG_{\G_2}=\RG_{\G_2}\circ\LG_{\G_1}.
\end{equation}
下面还有一个稍不显然的关系；它推广了 Gerstenhaber 关于 Hochschild 上同调上的 Yoneda 积具有分次交换性的观察。

\begin{lemma}\label{st:1.4}
设 $\Q=\nufun(\A,\A)$。假设有函子 $\G_1,\G_2\in\Ob\Q$ 和自然变换 $T_1\in\hom_{\Q}(\G_1,\Id_{\A})$、$T_2\in\hom_{\Q}(\G_2,\Id_{\A})$。那么
\[
[\mu^2_{\Q}(T_1,\LG^1_{\G_1}(T_2))]
=(-1)^{\ddagger}[\mu^2_{\Q}(T_2,\RG^1_{\G_2}(T_1))]
\in H(\hom_{\Q}(\G_1\circ\G_2,\Id_{\A})),
\]
其中 $\ddagger=(|T_1|-1)(|T_2|-1)+1$。
\end{lemma}
\begin{proof}
定义 $U\in\hom^{|T_1|+|T_2|-1}_{\Q}(\G_1\circ\G_2,\Id_{\A})$ 为
\[
\begin{aligned}
&U^d(a_d,\ldots,a_1)\\
&=\sum_{m,n,r}\sum_{s_1,\ldots,s_r}(-1)^{(|T_2|-1)\sgnmark_n}
T_1^{r+d-n-m+1}\bigl(a_d,\ldots,a_{n+m+1},\\
&\qquad T_2^m(a_{n+m},\ldots,a_{n+1}),
\G_2^{s_r}(a_n,\ldots,a_{n-s_r+1}),\ldots,\G_2^{s_1}(a_{s_1},\ldots,a_1)\bigr).
\end{aligned}
\]
直接计算可知，$\mu^1_{\Q}(U)$ 就是本引理所陈述等式两侧的差。
\end{proof}

\begin{note}\label{st:1.5}
与固定函子作复合，当然只是完整情形的一部分。为求简明，我们就讨论到这里，尽管这会使后面讨论拟等价时略有不便。
\end{note}

\clearpage
\subsection{长度滤过}\label{sec:1f}
\begin{definition*}
设 $\G_0,\G_1$ 是 $\Q=\nufun(\A,\B)$ 的对象。\emph{长度滤过}是 $\hom_{\Q}(\G_0,\G_1)$ 的一个完备递减滤过 $F^{\bullet}$：其 $F^r$ 项由满足 $T^0=\cdots=T^{r-1}=0$ 的预自然变换 $T$ 组成。
\end{definition*}

\SourcePage{21}
令 $A=H(\A)$、$B=H(\B)$、$G_k=H(\G_k)$，则相应谱序列的起始项为
\begin{equation}\label{eq:1.13}
E_1^{rs}=\prod_{X_0,\ldots,X_r}\Hom^s_{\K}\bigl(\Hom_A(X_{r-1},X_r)\otimes\cdots\otimes\Hom_A(X_0,X_1),\Hom_B(G_0X_0,G_1X_r)\bigr).
\end{equation}
即使该谱序列一般未必收敛，仍可将其用于比较论证。在最初等的层次上，由于\eref{1.13}只涉及上同调数据，有以下结论。

\begin{lemma}\label{st:1.6}
取 $\G_0,\G_1,\G_2\in\Ob\Q$，并设 $T$ 是从 $\G_0$ 到 $\G_1$ 的自然变换。假设对每个 $X\in\Ob\A$，与 $[T^0]$ 作右复合给出同构 $\Hom_{H(\B)}(\G_1X,\G_2X)\to\Hom_{H(\B)}(\G_0X,\G_2X)$。那么，在 $H(\Q)$ 中与 $T$ 作右复合给出同构 $\Hom_{H(\Q)}(\G_1,\G_2)\to\Hom_{H(\Q)}(\G_0,\G_2)$。（左复合也有相应结论。）
\end{lemma}

\begin{lemma}\label{st:1.7}
若 $\G:\A\to\B$ 是上同调全忠实的非含幺 $\Ai$ 函子，那么对任意 $\C$，$\LG_{\G}:\nufun(\C,\A)\to\nufun(\C,\B)$ 也具有这一性质。（对于 $\RG_{\G}:\nufun(\B,\C)\to\nufun(\A,\C)$ 的相应陈述一般不成立，但例如在 $\G$ 为拟同构时成立。）
\end{lemma}

由\eref{1.8}可见，\eref{1.13}上的微分 $\partial=\partial^{rs}$ 为
\[
\begin{aligned}
(\partial T)(a_{r+1},\ldots,a_1)
&=(-1)^{(r+s)(|a_1|-1)+1}T(a_{r+1},\ldots,a_2)G_0(a_1)\\
&\quad+(-1)^{\sgnmark_r+(r+s)}G_1(a_{r+1})T(a_r,\ldots,a_1)\\
&\quad+\sum_n(-1)^{\sgnmark_n+1+(r+s-1)}T(a_{r+1},\ldots,a_{n+2}a_{n+1},\ldots,a_1).
\end{aligned}
\]
（这里的约定与通常使用的约定略有不同：$T$ 是 $E_1^{rs}$ 的元素，因此作用在上同调层次上；同样，$a_i$ 是 $A$ 中的态射。）

\begin{definition*}
我们把这一双分次复形记作 $CC^{r+s}(A,B)^s=E_1^{rs}$，把它的 $\partial$ 上同调记作 $HH^{r+s}(A,B)^s=E_2^{rs}$，并分别称为 $A$ 以 $B$ 为系数的\emph{Hochschild 复形}和\emph{Hochschild 上同调}（它们依赖于函子 $G_0,G_1$，尽管记号并未显示这一点）。第一个非平凡列为
\[
E_2^{0s}=HH^s(A,B)^s=\Hom^s_{\operatorname{Nu\text{-}Fun}(A,B)}(G_0,G_1),
\]
而边缘同态 $H^s(\hom_{\Q}(\G_0,\G_1))\to E_2^{0s}$ 为 $T\mapsto H(T)$。因此，其他各列的 Hochschild 上同调群可以视为\eref{1.10}全忠实的障碍。
\end{definition*}

\SourcePage{22}
\begin{note}\label{st:1.8}
Hochschild 上同调的经典定义针对一个 $\K$ 代数 $A$ 及一个双模 $M$，给出分次群 $HH(A,M)$。这一概念可直接推广到分次代数与分次双模；此时 Hochschild 上同调成为双分次的。作为一个特殊情形，设有两个代数同态 $G_0,G_1:A\to B$。那么 $B$ 成为 $A$ 双模（左侧通过 $G_1$ 作用，右侧通过 $G_0$ 作用），因而得到 Hochschild 上同调群 $HH(A,B)$。上面考察的群，就是把这一概念从代数（可以视为只有一个对象的范畴）直接推广到一般范畴的结果，只是我们对双分次分量的标记略有些不合常规（下面将在\stref{笔记}{2.7}中，从更抽象的同调代数观点继续这一讨论）。

为避免术语混淆，要指出，在 $\Ai$ 世界中也有 Hochschild 上同调理论，即与一个 $\Ai$ 代数及其上的一个 $\Ai$ 双模相联系的分次群 $HH(\A,\M)$。同样地，如果 $\G_0,\G_1:\A\to\B$ 是 $\Ai$ 代数同态，可以借助它们使 $\B$ 成为 $\A$ 上的 $\Ai$ 双模，从而定义 $HH(\A,\B)$。把这一概念推广到 $\Ai$ 范畴，所得恰好就是我们的谱序列所趋向的上同调 $H(\hom_{\Q}(\G_0,\G_1))$。
\end{note}

\subsection{构造函子}\label{sec:1g}
接下来几段讨论 $\Ai$ 范畴与 $\Ai$ 函子的构造，它们属于同伦代数，而非同调代数。作为第一个问题，设有非含幺 $\Ai$ 范畴 $\A,\B$，其基础上同调范畴为 $A,B$，并给定一个非含幺函子 $G:A\to B$。考虑与函子对 $(G,G)$ 相联系的 Hochschild 上同调 $HH(A,B)$。

\begin{lemma}\label{st:1.9}
若对所有 $r\ge3$ 都有 $HH^2(A,B)^{2-r}=0$，那么 $G$ 可以由一个非含幺 $\Ai$ 函子 $\G:\A\to\B$ 实现，即 $H(\G)=G$。
\end{lemma}
\begin{proof}
这是一个障碍理论论证，可看作谱序列的非线性类似物。设对某个 $r\ge3$，已经定义 $\G^1,\ldots,\G^{r-1}$，使 $\Ai$ 函子方程\eref{1.6}在不超过 $r-1$ 阶时成立（当 $r=3$ 时，由于 $G$ 是函子，总能做到这一点）。任意暂取一个 $\G^r$。由此得到的 $r$ 阶 $\Ai$ 函子方程两侧之差，是一族映射
\begin{equation}\label{eq:1.14}
\theta:\hom_{\A}(X_{r-1},X_r)\otimes\cdots\otimes\hom_{\A}(X_0,X_1)\longrightarrow\hom_{\B}(\G(X_0),\G(X_r))[2-r].
\end{equation}
直接计算表明，$\theta$ 在上同调上诱导一个映射 $T$；计算从略。此外，$T$ 是 $CC^2(A,B)^{2-r}$ 中的上闭链。由假设，它在 Hochschild 复形中可写成 $\partial S$。选择一族链层次的映射 $\sigma$，使其类似于\eref{1.14}，并代表 $S$；
\SourcePage{23}
然后把已给定的 $\G^{r-1}$ 改为 $\G^{r-1}-\sigma$。这不影响 $r-1$ 阶的 $\Ai$ 函子方程，却将 $T$ 改为 $T-\partial S$，而后者按假设为零。因此，$r$ 阶 $\Ai$ 函子方程现在在上同调层次上成立。适当选择新的 $\G^r$，便可使它在链层次上也成立，这就完成了归纳步骤。
\end{proof}

我们还需要一个相对版本。
\begin{lemma}\label{st:1.10}
设 $\A,\B$ 是非含幺 $\Ai$ 范畴，$\widetilde{\A}\subset\A$ 是全子范畴，相应的上同调范畴为 $A,B,\widetilde A$。设 $G:A\to B$ 是通常的非含幺函子，$\widetilde G$ 是它在 $\widetilde A$ 上的限制。假设由此得到的 Hochschild 上同调群之间的自然限制映射
\[
HH^{r+s}(A,B)^s\longrightarrow HH^{r+s}(\widetilde A,B)^s
\]
当 $r+s=2$、$r\ge3$ 时是单射，当 $r+s=1$、$r\ge2$ 时是满射。那么，给定任意满足 $H(\widetilde{\G})=\widetilde G$ 的非含幺 $\Ai$ 函子 $\widetilde{\G}:\widetilde{\A}\to\B$，都可以将其扩张为满足 $H(\G)=G$ 的 $\G:\A\to\B$。
\end{lemma}
\begin{proof}
为此，使用由 $CC(A,B)$ 中限制到 $CC(\widetilde A,B)$ 后为零的链组成的子复形。假设保证这一子复形在 $r\ge3$ 的双次数 $(2,2-r)$ 上无上同调。作此修改后，证明与前面完全相同。
\end{proof}

\clearpage
\subsection{同伦}\label{sec:1h}
\begin{definition*}
考虑 $\Q=\nufun(\A,\B)$ 的对象集中的两个非含幺 $\Ai$ 函子 $\F_0,\F_1$，假设它们在对象上的作用相同。令 $D=\F_0-\F_1\in\hom^1_{\Q}(\F_0,\F_1)$ 为由 $D^0=0$、$D^d=\F_0^d-\F_1^d$（$d>0$）定义的预自然变换。它是真正的自然变换，即 $\mu^1_{\Q}(D)=0$。如果存在还满足 $T^0=0$ 的 $T\in\hom^0_{\Q}(\F_0,\F_1)$，使 $D=\mu^1_{\Q}(T)$，就称 $\F_0,\F_1$ \emph{同伦}。至少从表面上看，这一名称是合理的，因为同伦函子在上同调上诱导相同的函子：$H(\F_0)=H(\F_1)$。\znote{15}{要求 $T^0=0$ 并非技术洁癖：同伦的差 $D=\F_0-\F_1$ 在零阶已为零，若允许 $T^0\neq0$，则 $\mu^1_{\Q}(T)$ 的零阶项会混入 $\F_0^1-\F_1^1$ 的修正量，$H(\F_0)=H(\F_1)$ 便不再自动成立。把 $T^0$ 归零后，$\mu^1_{\Q}(D)=0$ 逐阶给出 $\F_0^1-\F_1^1$ 是上边界，两个函子在上同调上因而一致。[来源：Seidel, Fukaya categories and Picard–Lefschetz theory, §2c（同伦与同构）；并见\stref{笔记}{1.11}。]}与固定函子作左复合、右复合，都保持同伦。这意味着，若 $T$ 是从 $\F_0$ 到 $\F_1$ 的同伦，那么 $\RG^1_{\G}(T)$ 是 $\F_0\circ\G$ 与 $\F_1\circ\G$ 之间的同伦，而 $\LG^1_{\G}(T)$ 是 $\G\circ\F_0$ 与 $\G\circ\F_1$ 之间的同伦。
\end{definition*}

同伦是等价关系。原则上这不难验证，但计算可能显得有些混乱，原因如下：由于考察的函子 $\F_k$ 在对象上的作用相同，对于所有 $k,l$，分次向量空间 $\hom_{\Q}(\F_k,\F_l)$ 都相同，而我们需要把预自然变换从一个空间移到另一个空间中。为避免混淆，记 $T_{[kl]}$ 表示把 $T$ 视为 $\hom_{\Q}(\F_k,\F_l)$ 中的元素。采用这一记号，基本观察是：对于任意三个函子 $\F_j,\F_k,\F_l$，
\begin{equation}\label{eq:1.15}
\begin{aligned}
\mu^1_{\Q}(T_{[jl]})
&=\mu^1_{\Q}(T_{[jk]})-\mu^2_{\Q}(\F_k-\F_l,T_{[jk]})\\
&=\mu^1_{\Q}(T_{[kl]})+\mu^2_{\Q}(T_{[kl]},\F_j-\F_k).
\end{aligned}
\end{equation}
\SourcePage{24}
为证明传递性，设 $T_1$ 是从 $\F_0$ 到 $\F_1$ 的同伦，$T_2$ 是从 $\F_1$ 到 $\F_2$ 的同伦。由\eref{1.15}可见
\begin{equation}\label{eq:1.16}
T=(T_1)_{[02]}+(T_2)_{[02]}+\mu^2_{\Q}(T_2,T_1)
\end{equation}
满足
\[
\begin{aligned}
\mu^1_{\Q}(T)
&=\F_1-\F_2-\mu^2_{\Q}(\F_1-\F_2,T_1)\\
&\quad+\F_0-\F_1+\mu^2_{\Q}(T_2,\F_0-\F_1)+\mu^1_{\Q}\mu^2_{\Q}(T_2,T_1)
=\F_0-\F_2,
\end{aligned}
\]
因此给出从 $\F_0$ 到 $\F_2$ 的同伦。对称性的证明也类似：设 $T_1$ 是从 $\F_0$ 到 $\F_1$ 的同伦，考虑映射
\begin{equation}\label{eq:1.17}
\begin{aligned}
\phi &: \hom_{\Q}(\F_1,\F_0)\longrightarrow\hom_{\Q}(\F_1,\F_1),
&\phi(T)&=(-1)^{|T|}\mu^2_{\Q}(T_1,T)+T_{[11]},\\
\psi &: \hom_{\Q}(\F_1,\F_1)\longrightarrow\hom_{\Q}(\F_0,\F_1),
&\psi(T)&=\mu^2_{\Q}(T,T_1)+T_{[01]}.
\end{aligned}
\end{equation}
由于各式中的最后一项，它们是同构；再应用\eref{1.15}，可见它们还是链映射。最后，再次利用\eref{1.15}，得到
\[
\begin{aligned}
\phi((\F_0-\F_1)_{[10]})&=\mu^1_{\Q}((T_1)_{[11]}),\\
\psi(\mu^1_{\Q}((T_1)_{[11]}))&=\mu^1_{\Q}\mu^2_{\Q}((T_1)_{[11]},T_1)+(\F_0-\F_1).
\end{aligned}
\]
由此可知，$-T_1-\mu^2_{\Q}((T_1)_{[11]},T_1)$ 在 $\psi\circ\phi$ 下的原像，是从 $\F_1$ 到 $\F_0$ 的链同伦。

\begin{note}\label{st:1.11}
当 $\B$ 具有严格恒等态射时，上述计算及一般的同伦概念会变得更清楚一些；见后面的\secref{2c}。这里先给出一个来自有理同伦论的不同解释，参看\bcite{68}，第 X 章。考虑如下 dg 代数：$I=\K u_0\oplus\K u_1\oplus\K h$，其中 $|u_0|=|u_1|=0$、$|h|=1$；微分和乘法为（省略因次数原因而为零的式子）$\partial u_0=h=-\partial u_1$，以及
\[
\begin{gathered}
u_0u_0=u_0,\qquad u_1u_1=u_1,\qquad u_0u_1=u_1u_0=0,\\
u_1h=h,\qquad hu_1=u_0h=0,\qquad hu_0=h.
\end{gathered}
\]
$H(I)\cong\K$，生成元为 $u_0+u_1$，故投影到 $u_k$ 所生成子空间的两个映射 $f_k:I\to\K$（$k=0,1$）都是拟同构。若 $\B$ 是非含幺 $\Ai$ 代数，很容易定义张量积 $I\otimes\B$：在基础分次向量空间的张量积上，令
\[
\mu^1(i\otimes b)=i\otimes\mu^1_{\B}(b)+(-1)^{|i|+|b|}\partial(i)\otimes b,
\]
而对于 $d\ge2$，令
\[
\mu^d(i_d\otimes b_d,\ldots,i_1\otimes b_1)
=(-1)^{\triangleleft}i_d\cdots i_1\otimes\mu^d_{\B}(b_d,\ldots,b_1),
\qquad\triangleleft=\sum_{j>k}(|b_j|-1)|i_k|.
\]
简单计算表明，$T$ 是两个 $\Ai$ 同态 $\F_0,\F_1:\A\to\B$ 之间的同伦，当且仅当
\[
\begin{aligned}
\G(a_d,\ldots,a_1)
&=u_0\otimes\F_0(a_d,\ldots,a_1)+u_1\otimes\F_1(a_d,\ldots,a_1)\\
&\quad+(-1)^{\sgnmark_d}h\otimes T(a_d,\ldots,a_1)
\end{aligned}
\]
\SourcePage{25}
定义了一个 $\Ai$ 同态 $\G:\A\to I\otimes\B$。这一情形可归结为下面的交换图。
\SourceDiagram{diagram-1h}{0.24715358}
作出显然的修改后，同样的叙述也适用于 $\Ai$ 范畴与函子。
\end{note}

\subsection{同调扰动理论}\label{sec:1i}
设 $\B$ 是非含幺 $\Ai$ 范畴。假设对每对对象 $(X_0,X_1)$，有如下图：
\SourceDiagram{diagram-1i}{0.40292265}
其中 $(\hom_{\A}(X_0,X_1),\mu^1_{\A})$ 是某个向量空间链复形；$\F^1,\G^1$ 是链映射；$T^1$ 的次数为 $-1$，且满足 $\mu^1_{\B}T^1+T^1\mu^1_{\B}=\F^1\G^1-\id$。\emph{扰动引理}即以下陈述。\znote{17}{这正是经典的同伦扰动引理在 $\Ai$ 范畴上的形式：若 $f:X\to Y$、$\nabla:Y\to X$ 为链映射，$\Phi$ 为满足 $f\nabla=\id$、$\nabla f=\id+d\Phi+\Phi d$ 的收缩同伦，则扰动后的微分可由显式公式 $\Phi^{\tau}=\sum_{r\ge0}\Phi(1+d^{\tau}\Phi)^{r}$ 给出。原书的 $\F,\G,T$ 分别对应 $f,\nabla,\Phi$，\eref{1.18}即该显式公式的展开。[来源：nLab \textquotedblleft homological perturbation theory\textquotedblright，https://\allowbreak ncatlab.org/\allowbreak nlab/\allowbreak show/\allowbreak homological+\allowbreak perturbation+\allowbreak theory（2026-11-15 访问）。]}

\begin{proposition}\label{st:1.12}
可以构造一个非含幺 $\Ai$ 范畴 $\A$，其对象集为 $\Ob\A=\Ob\B$，态射空间为 $\hom_{\A}(X_0,X_1)$，一阶结构映射为 $\mu^1_{\A}$；同时得到非含幺 $\Ai$ 函子 $\F:\A\to\B$、$\G:\B\to\A$，它们在对象上均为恒等，一阶部分是给定的映射 $\F^1,\G^1$；还得到 $\F\circ\G$ 与 $\Id_{\B}$ 之间的同伦 $T$，其起始项是给定的 $T^1$。
\end{proposition}

所有这些结构都由明确公式定义，通常写成递归形式。例如，
\begin{equation}\label{eq:1.18}
\begin{aligned}
\F^d(a_d,\ldots,a_1)
&=\sum_r\sum_{s_1,\ldots,s_r}T^1\bigl(\mu^r_{\B}(\F^{s_r}(a_d,\ldots,a_{d-s_r+1}),\ldots,\F^{s_1}(a_{s_1},\ldots,a_1))\bigr),\\
\mu^d_{\A}(a_d,\ldots,a_1)
&=\sum_r\sum_{s_1,\ldots,s_r}\G^1\bigl(\mu^r_{\B}(\F^{s_r}(a_d,\ldots,a_{d-s_r+1}),\ldots,\F^{s_1}(a_{s_1},\ldots,a_1))\bigr),
\end{aligned}
\end{equation}
其中两个和式都对 $r\ge2$ 的分拆 $s_1+\cdots+s_r=d$ 求和。很容易验证，$\mu^{\bullet}_{\A}$ 确实是 $\Ai$ 结构，$\F$ 确实是 $\Ai$ 函子。\znote{17}{把 $\hom_{\A}$ 取成 $\hom_{\B}$ 的上同调（\stref{命题}{1.12}中\textquotedblleft 分裂\textquotedblright 的极端情形），即得 Kadeishvili 定理：dg 代数的上同调自然带有 $\Ai$ 代数结构，且与原代数弱等价。这一转移定理正是镜像对称中\textquotedblleft Floer 上同调仍带 $\Ai$ 结构\textquotedblright 的来源，也是本书后面以上同调范畴为出发点的依据。[来源：nLab \textquotedblleft A-infinity-algebra\textquotedblright（Kadeishvili–Merkulov 定理一节），https:/\allowbreak /\allowbreak ncatlab.org/\allowbreak nlab/\allowbreak show/\allowbreak A-infinity-algebra（2026-11-15 访问）。]}
\SourcePage{26}
为了使计算更具几何意味，可以把递归定义“展开”，写成涉及树的求和的直接表达式（更精确地，是\secref{9d}意义下的稳定树）。$\G$ 和 $T$ 的构造也可以这样处理，但相应公式略为复杂一些。

\begin{note}\label{st:1.13}
一个典型应用如下。给定非含幺 $\Ai$ 范畴 $\B$，把每个复形 $(\hom_{\B}(X_0,X_1),\mu^1_{\B})$ 分裂为微分为零的直和项与一个无上同调的补项。取 $\hom_{\A}(X_0,X_1)$ 为前一个直和项，$\F^1$ 为包含映射，$\G^1$ 为投影。为无上同调的补项选择一个收缩同伦，就得到所需的 $T^1$。于是\stref{命题}{1.12}构造出满足 $\mu^1_{\A}=0$、且与 $\B$ 拟同构的 $\A$。在这一特殊情形中，$\G\circ\F$ 也同伦于 $\Id_{\A}$。为说明这一点，先注意，由于 $\F\circ\G$ 同伦于 $\Id_{\B}$，$\G\circ\F\circ\G$ 同伦于 $\G$。\stref{引理}{1.7}说明
\[
\RG^1_{\G}:\hom_{\nufun(\A,\A)}(\Id_{\A},\G\circ\F)
\longrightarrow\hom_{\nufun(\B,\A)}(\G,\G\circ\F\circ\G)
\]
是链复形的拟同构。此外，改为考察零阶项为零的预自然变换组成的子复形，这一结论仍然成立。显然，$\RG^1_{\G}$ 将 $\Id_{\A}-\G\circ\F$ 映为 $\G-\G\circ\F\circ\G$，因此这就证明了所需的第二个同伦的存在性。
\end{note}

\begin{corollary}\label{st:1.14}
非含幺 $\Ai$ 范畴之间的任意拟同构，都具有同伦意义下的逆。
\end{corollary}
\begin{proof}
设给定的拟同构为 $\mathcal K:\B\to\widetilde{\B}$。应用扰动引理，构造由非含幺 $\Ai$ 范畴及拟同构组成的如下图：
\SourceDiagram{diagram-1-14}{0.22638437}
其中 $\mu^1_{\A}=\mu^1_{\widetilde{\A}}=0$。复合 $\mathcal H=\widetilde{\G}\circ\mathcal K\circ\F$ 是两个一阶复合为零的 $\Ai$ 范畴之间的拟同构，因此必然是形式微分同胚，从而相对于复合具有严格的逆 $\mathcal H^{-1}$。考虑 $\mathcal I=\F\circ\mathcal H^{-1}\circ\widetilde{\G}$。按构造，$\mathcal K$ 同伦于 $\widetilde{\F}\circ\mathcal H\circ\G=\widetilde{\F}\circ\widetilde{\G}\circ\mathcal K\circ\F\circ\G$。因此 $\mathcal K\circ\mathcal I$ 同伦于 $\widetilde{\F}\circ\mathcal H\circ\G\circ\F\circ\mathcal H^{-1}\circ\widetilde{\G}$，而后者又同伦于恒等函子。另一个同伦以同样方式构造。
\end{proof}

\SourcePage{27}
\begin{note}\label{st:1.15}
即使不计拓扑学中的先驱结果\bcite{137}，\stref{命题}{1.12}也有很长的历史。早期特殊情形见\bcite{80}和\bcite{69}；经过此后种种推广与简化，最终得到这里陈述的版本，取自\bcite{98}。证明、更多参考文献，以及一幅鲤鱼图，都可参看该文。至少在 $\operatorname{char}(\K)=0$ 时，可以由控制 $\Ai$ 结构的操作子的“余纤维性”推出抽象扰动引理的存在，见\bcite{99}。\stref{推论}{1.14}最早的出处似乎是\bcite{81}、\bcite{111}。对于 $L_\infty$ 代数，略弱一些的类似结果出现于\bcite{88}，第 4.5 节；其证明与上面基本相同。若想了解体现\stref{引理}{1.9}思想的障碍理论方法，可参看\bcite{60}，附录 A5。
\end{note}

\clearpage
\subsection{\texorpdfstring{$A_\infty$}{A∞} 模}\label{sec:1j}
\begin{definition*}
设 $\Ch$ 是 $\K$ 向量空间链复形组成的 dg 范畴，并将其视为 $\Ai$ 范畴。对于固定的非含幺 $\Ai$ 范畴 $\A$，将 $\A$ 上非含幺（右）$\Ai$ 模组成的非含幺 $\Ai$ 范畴定义为
\[
\operatorname{nu\text{-}mod}(\A)=\nufun(\A^{\mathrm{opp}},\Ch).
\]
\end{definition*}

\begin{definition*}
有必要具体写出细节。$\Q=\operatorname{nu\text{-}mod}(\A)$ 的对象 $\M$，由对每个 $X\in\Ob\A$ 给出的分次向量空间 $\M(X)$，以及对所有 $d\ge1$ 给出的映射组成：
\[
\mu^d_{\M}:\M(X_{d-1})\otimes\hom_{\A}(X_{d-2},X_{d-1})\otimes\cdots\otimes\hom_{\A}(X_0,X_1)
\longrightarrow\M(X_0)[2-d].
\]
若从 $\Ai$ 函子的角度理解，$\mu^1_{\M}$ 就是像对象 $\M(X_0)\in\Ob\Ch$ 上的微分，而对 $d\ge2$，$\mu^d_{\M}$ 是函子在态射上的作用；更精确地，是把自变量顺序反转后的 $\M^{d-1}$。方程\eref{1.6}变为
\begin{equation}\label{eq:1.19}
\begin{aligned}
&\sum_n(-1)^{\sgnmark_n}\mu^{n+1}_{\M}\bigl(\mu^{d-n}_{\M}(b,a_{d-1},\ldots,a_{n+1}),a_n,\ldots,a_1\bigr)\\
&\quad+\sum_{m,n}(-1)^{\sgnmark_n}\mu^{d-m+1}_{\M}\bigl(b,a_{d-1},\ldots,a_{n+m+1},\mu^m_{\A}(a_{n+m},\ldots,a_{n+1}),a_n,\ldots,a_1\bigr)=0,
\end{aligned}
\end{equation}
其中第二个和式满足 $n+m<d$。\znote{19}{把右模写成 $\nufun(\A^{\mathrm{opp}},\Ch)$，是把\textquotedblleft 模 $=$ 单对象范畴上的函子\textquotedblright 这一经典对应推广到 $\Ai$ 情形：$\Ai$ 代数就是只有一个对象的 $\Ai$ 范畴，其上的模便是定义在反范畴上、取值于链复形范畴的 $\Ai$ 函子。也正因如此，\eref{1.19}中 $\M(X_{d-1})$ 排在最左、$a$ 的指标递减，与 $\hom_{\A}$ 的复合顺序正好相反。[来源：nLab \textquotedblleft A-infinity-category\textquotedblright，https:/\allowbreak /\allowbreak ncatlab.org/\allowbreak nlab/\allowbreak show/\allowbreak A-infinity-category（2026-11-15 访问）。]}
\end{definition*}

\begin{definition*}
再设 $\M_0,\M_1$ 为两个非含幺 $\Ai$ 模。$\hom_{\Q}(\M_0,\M_1)$ 中的态射称为\emph{预模同态}，写成一列映射
\[
t^d:\M_0(X_{d-1})\otimes\hom_{\A}(X_{d-2},X_{d-1})\otimes\cdots\otimes\hom_{\A}(X_0,X_1)
\longrightarrow\M_1(X_0)[|t|-d+1],\qquad d\ge1.
\]
它与预自然变换 $T$ 的一般表述之间，通过显然的自变量调换相联系，同时作一个为方便而选取的符号调整：
\begin{equation}\label{eq:1.20}
t^d(b,a_{d-1},\ldots,a_1)=(-1)^{\S} T^{d-1}(a_1,\ldots,a_{d-1})(b),
\end{equation}
\SourcePage{28}
其中 $T^{d-1}(a_1,\ldots,a_{d-1})\in\hom_{\Ch}(\M_0(X_{d-1}),\M_1(X_0))$，且 $\S=(|T|-1)|b|+|T|(|T|-1)/2$。真正的模同态是满足 $\mu^1_{\Q}(t)=0$ 的那些 $t$。
\end{definition*}

为看清其具体含义，写出 $\Q$ 中的复合映射：
\begin{equation}\label{eq:1.21}
\begin{aligned}
&(\mu^1_{\Q}t)^d(b,a_{d-1},\ldots,a_1)\\
&=\sum_n(-1)^{\ddagger}\mu^{n+1}_{\M_1}\bigl(t^{d-n}(b,a_{d-1},\ldots,a_{n+1}),a_n,\ldots,a_1\bigr)\\
&\quad+\sum_n(-1)^{\ddagger} t^{n+1}\bigl(\mu^{d-n}_{\M_0}(b,a_{d-1},\ldots,a_{n+1}),a_n,\ldots,a_1\bigr)\\
&\quad+\sum_{m,n}(-1)^{\ddagger} t^{d-m+1}\bigl(b,a_{d-1},\ldots,a_{n+m+1},\mu^m_{\A}(a_{n+m},\ldots,a_{n+1}),a_n,\ldots,a_1\bigr);\\[2pt]
&(\mu^2_{\Q}(t_2,t_1))^d(b,a_{d-1},\ldots,a_1)\\
&=\sum_n(-1)^{\ddagger} t_2^{n+1}\bigl(t_1^{d-n}(b,a_{d-1},\ldots,a_{n+1}),a_n,\ldots,a_1\bigr);\\
&\mu^d_{\Q}=0\qquad\text{对所有 }d\ge3.
\end{aligned}
\end{equation}
这里 $\ddagger=|a_{n+1}|+\cdots+|a_{d-1}|+|b|-d+n+1$。$\mu^3_{\Q},\mu^4_{\Q},\ldots$ 为零，是一个非正式但有用的一般原则的实例：$\nufun(\A,\B)$ 继承的是 $\B$ 的性质，而非 $\A$ 的性质。

在上同调层次上，任意 $\M$ 都给出一个非含幺 $H(\A)$ 模 $H(\M)$；更形式化地说，它是从 $H(\A)$ 到分次向量空间范畴的非含幺反变函子。它由 $\M(X)$ 关于 $\partial(b)=(-1)^{|b|}\mu^1_{\M}(b)$ 的上同调，以及 $b\cdot a=(-1)^{|a|}\mu^2_{\M}(b,a)$ 所诱导的模结构组成。同样，给定 $\Ai$ 模同态 $t$，得到诱导映射 $H(t):H(\M_0(X))\to H(\M_1(X))$，$[b]\mapsto[(-1)^{|b|}t^1(b)]$；它们组成通常意义下的模同态。下面是\stref{引理}{1.6}的一个特殊情形。

\begin{lemma}\label{st:1.16}
设 $t:\M_0\to\M_1$ 是 $\Ai$ 模同态，并且对每个 $X\in\Ob\A$，映射 $H(t):H(\M_0(X))\to H(\M_1(X))$ 都是同构。那么，对于任意 $\N$，与 $t$ 作右复合及左复合，分别诱导拟同构
\[
\hom_{\Q}(\M_1,\N)\longrightarrow\hom_{\Q}(\M_0,\N),\qquad
\hom_{\Q}(\N,\M_0)\longrightarrow\hom_{\Q}(\N,\M_1).
\]
\end{lemma}

特别地，如果非含幺 $\Ai$ 模 $\M$ 对所有 $X$ 都满足 $H(\M(X),\mu^1_{\M})=0$，那么对所有 $\N$ 都有 $H(\hom_{\Q}(\M,\N))=H(\hom_{\Q}(\N,\M))=0$。借用\bcite{136}、\bcite{21}中的术语，可以说在 $\Ai$ 模的世界中，每个对象同时都是 K-投射的和 K-内射的。

\subsection{拉回}\label{sec:1k}
\begin{definition*}
任意非含幺 $\Ai$ 函子 $\G:\A\to\B$ 都诱导一个拉回函子
\begin{equation}\label{eq:1.22}
\G^*=\RG_{\G^{\mathrm{opp}}}:\operatorname{nu\text{-}mod}(\B)\longrightarrow\operatorname{nu\text{-}mod}(\A).
\end{equation}
\SourcePage{29}
与任意右复合一样，它的高阶项为零，因此是 dg 范畴之间的非含幺 dg 函子。具体地，
\[
\begin{aligned}
(\G^*\M)(X)&=\M(\G(X)),\\
\mu^d_{\G^*\M}(b,a_{d-1},\ldots,a_1)
&=\sum_r\sum_{s_1,\ldots,s_r}\mu^{r+1}_{\M}\bigl(b,\G^{s_r}(a_{d-1},\ldots,a_{d-s_r}),\ldots,\G^{s_1}(a_{s_1},\ldots,a_1)\bigr),\\
(\G^*t)^d(b,a_{d-1},\ldots,a_1)
&=\sum_r\sum_{s_1,\ldots,s_r}t^{r+1}\bigl(b,\G^{s_r}(a_{d-1},\ldots,a_{d-s_r}),\ldots,\G^{s_1}(a_{s_1},\ldots,a_1)\bigr).
\end{aligned}
\]
\end{definition*}

\subsection{Yoneda 嵌入}\label{sec:1l}
\begin{definition*}
对 $\A$ 的每个对象 $Y$，都有相应的非含幺 $\Ai$ 模 $\Y$，定义为 $\Y(X)=\hom_{\A}(X,Y)$、$\mu^d_{\Y}=\mu^d_{\A}$。这是一个典范的非含幺 $\Ai$ 函子 $\mathcal I=\mathcal I_{\A}:\A\to\Q=\operatorname{nu\text{-}mod}(\A)$ 的一部分，该函子称为\emph{Yoneda 嵌入}。对于 $c\in\hom_{\A}(Y_0,Y_1)$，$\mathcal I^1(c)\in\hom_{\Q}(\Y_0,\Y_1)$ 是如下预模同态：
\[
\begin{aligned}
\Y_0(X_{d-1})\otimes\hom_{\A}(X_{d-2},X_{d-1})\otimes\cdots\otimes\hom_{\A}(X_0,X_1)&\longrightarrow\Y_1(X_0),\\
(b,a_{d-1},\ldots,a_1)&\longmapsto\mu^{d+1}_{\A}(c,b,a_{d-1},\ldots,a_1).
\end{aligned}
\]
高阶项 $\mathcal I^d$，$d\ge2$，由这一公式的显然类似式定义。\znote{21}{Yoneda 嵌入给出的模 $\Y=\hom_{\A}(\cdot,Y)$ 是\textquotedblleft 自由模\textquotedblright：$\Ai$ 版 Yoneda 引理说明 $\hom_{\Q}(\Y,\M)$ 的上同调就是 $H(\M(Y))$，\eref{1.25}中的 $\lambda$ 正是这一对应的显式形式。一个常见推论是每个线性 $\Ai$ 范畴都 $\Ai$ 等价于某个 dg 范畴；但该等价只在抽象层次上成立，对 Fukaya 范畴直接使用并不现实。[来源：nLab \textquotedblleft A-infinity-category\textquotedblright，https:/\allowbreak /\allowbreak ncatlab.org/\allowbreak nlab/\allowbreak show/\allowbreak A-infinity-category（2026-11-15 访问）。]}
\end{definition*}

Yoneda 嵌入与拉回在如下意义下相容。设 $\F:\A\to\B$ 为非含幺 $\Ai$ 函子，则存在典范自然变换
\begin{equation}\label{eq:1.23}
T:\mathcal I_{\A}\longrightarrow\F^*\circ\mathcal I_{\B}\circ\F.
\end{equation}
其一阶项 $T^1$ 对每个 $Y$ 给出一个模同态，而这个模同态本身是一列映射
\begin{equation}\label{eq:1.24}
\begin{aligned}
(\mathcal I_{\A}(Y))(X_{d-1})\otimes\hom_{\A}(X_{d-2},X_{d-1})\otimes\cdots\otimes\hom_{\A}(X_0,X_1)\\
\longrightarrow (\F^*\circ\mathcal I_{\B}\circ\F(Y))(X_0).
\end{aligned}
\end{equation}
按定义，$(\mathcal I_{\A}(Y))(X_{d-1})=\hom_{\A}(X_{d-1},Y)$，且 $(\F^*\circ\mathcal I_{\B}\circ\F)(Y)(X_0)=\hom_{\B}(\F(X_0),\F(Y))$；据此，\eref{1.24}由 $(b,a_{d-1},\ldots,a_1)\mapsto(-1)^{\sgnmark_{d-1}+|b|}\F^d(b,a_{d-1},\ldots,a_1)$ 给出。类似的表达式定义 $T^r$，$r\ge2$。

从另一个方向，可以把 $\mathcal I^1$ 推广为链映射
\begin{equation}\label{eq:1.25}
\begin{aligned}
\lambda=\lambda_{\M}&:\M(Y)\longrightarrow\hom_{\Q}(\Y,\M),\\
\lambda(c)^d(b,a_{d-1},\ldots,a_1)&=\mu^{d+1}_{\M}(c,b,a_{d-1},\ldots,a_1),
\end{aligned}
\end{equation}
其中 $\M$ 是任意非含幺 $\Ai$ 模。在上同调层次上，这一映射相对于两侧的乘法具有以下容易验证的性质：
\begin{equation}\label{eq:1.26}
[\lambda_{\M}(\mu^2_{\M}(c,b))]
=[\mu^2_{\Q}(\lambda_{\M}(c),\mathcal I^1(b))]\in\Hom_{H(\Q)}(\Y_0,\M),
\end{equation}
\SourcePage{30}
其中 $[c]\in H(\M(Y_1))$、$[b]\in\Hom_{H(\A)}(Y_0,Y_1)$；以及
\begin{equation}\label{eq:1.27}
[\mu^2_{\Q}(t,\lambda_{\M_0}(c))]
=[\lambda_{\M_1}(t^1(c))]\in\Hom_{H(\Q)}(\Y,\M_1),
\end{equation}
其中 $[t]\in\Hom_{H(\Q)}(\M_0,\M_1)$、$[c]\in H(\M_0(Y))$。注意，如果只考虑来自 Yoneda 嵌入的模，并相应地令 $t$ 为 $\A$ 中某个态射的像，那么这两个等式都归结为 $\mathcal I$ 的 $\Ai$ 函子方程的特殊情形。

\begin{note}\label{st:1.17}
映射 $\lambda$ 可以视为更广泛构造的一部分，这也解释了上面列出的性质。把 $\Q$ 的 Yoneda 嵌入与沿 $\A$ 的 Yoneda 嵌入作拉回相结合，得到非含幺 dg 函子，称为\emph{线性化函子}：
\[
\LG:\Q\xrightarrow{\mathcal I_{\Q}}\operatorname{nu\text{-}mod}(\Q)
\xrightarrow{\mathcal I_{\A}^*}\Q.
\]
（这一名称来自如下事实：每个模同态在 $\LG$ 下的像，只有一阶线性项可能非零。）这个函子带有典范自然变换 $L:\Id_{\Q}\to\LG$。在一阶，它对每个 $\Ai$ 模 $\M$ 给出模同态 $l=L^1_{\M}:\M\to\M_{\mathrm{lin}}=\LG(\M)$；后者的起始项由对所有 $Y$ 给出的线性映射 $l^1:\M(Y)\to\M_{\mathrm{lin}}(Y)$ 组成。按定义，$\M_{\mathrm{lin}}(Y)=\hom_{\Q}(\Y,\M)$，其中 $\Y$ 是 $Y$ 在 Yoneda 嵌入下的像，而该线性映射与我们的 $\lambda$ 只差一个符号 $(-1)^{|c|}$。$l$ 是 $\Ai$ 模同态这一事实蕴含\eref{1.26}；$L$ 是自然变换这一事实蕴含\eref{1.27}。
\end{note}

\section{恒等态射与等价}\label{section:2}
\subsection{定义}\label{sec:2a}
设 $\A$ 是非含幺 $\Ai$ 范畴。“$\A$ 具有恒等态射”有几种不同含义。

\begin{definition*}
若对每个对象 $X$，都有一个必然唯一的 $e_X\in\hom^0_{\A}(X,X)$，满足
\begin{equation}\label{eq:2.1}
\begin{gathered}
\mu^1_{\A}(e_X)=0;\\
(-1)^{|a|}\mu^2_{\A}(e_{X_1},a)=a=\mu^2_{\A}(a,e_{X_0})\qquad\text{对 }a\in\hom_{\A}(X_0,X_1);\\
\mu^d_{\A}(a_{d-1},\ldots,a_{n+1},e_{X_n},a_n,\ldots,a_1)=0\\
\text{对 }d>2,\ a_k\in\hom_{\A}(X_{k-1},X_k)\text{ 及任意 }0\le n<d,
\end{gathered}
\end{equation}
则称 $\A$ \emph{严格含幺}。
\end{definition*}

\begin{definition*}
若 $H(\A)$ 含幺，即对每个对象 $X$ 都有恒等态射，从而是通常意义下的分次线性范畴，则称 $\A$ \emph{上同调含幺}，或简称 \emph{c-含幺}。
\end{definition*}

\SourcePage{31}
\begin{definition*}
\emph{同伦含幺 $\Ai$ 范畴}由一个集合 $\Ob\A$、分次向量空间 $\hom_{\A}(X_0,X_1)$，以及对所有 $d+i_0+\cdots+i_d>0$ 给出的多重线性映射组成：
\begin{equation}\label{eq:2.2}
\begin{aligned}
\mu^{d,(i_d,\ldots,i_0)}_{\A}:\hom_{\A}(X_{d-1},X_d)\otimes\cdots\otimes\hom_{\A}(X_0,X_1)\\
\longrightarrow\hom_{\A}(X_0,X_d)\bigl[2-d-2\textstyle\sum_k i_k\bigr].
\end{aligned}
\end{equation}
这些映射须满足某些推广的结合律方程；当 $i_0=\cdots=i_d=0$ 时，它们归结为\eref{1.2}。最先出现的几个新方程是
\begin{equation}\label{eq:2.3}
\begin{gathered}
\mu^1_{\A}(\mu^{0,(1)}_{\A})=0,\\
(-1)^{|a|-1}\mu^2_{\A}(\mu^{0,(1)}_{\A},a)
+\mu^1_{\A}(\mu^{1,(1,0)}_{\A}(a))
+\mu^{1,(1,0)}_{\A}(\mu^1_{\A}(a))=a,\\
\mu^2_{\A}(a,\mu^{0,(1)}_{\A})
+\mu^1_{\A}(\mu^{1,(0,1)}_{\A}(a))
+(-1)^{|a|-1}\mu^{1,(0,1)}_{\A}(\mu^1_{\A}(a))=-a,\\
\mu^{1,(1,0)}_{\A}(\mu^{0,(1)}_{\A})
+\mu^{1,(0,1)}_{\A}(\mu^{0,(1)}_{\A})
+\mu^1_{\A}(\mu^{0,(2)}_{\A})=0.
\end{gathered}
\end{equation}

\end{definition*}
除去通常的符号因素，前面的两条乘法关系说明，与上闭链 $e_X=-\mu^{0,(1)}_{\A}$ 作左乘或右乘，都链同伦于恒等映射。利用这两个同伦中的任意一个，可见 $\mu^2_{\A}(e_X,e_X)$ 与 $e_X$ 只差一个上边界。这两个上边界之差是上闭链，而\eref{2.3}的最后一部分表明，它本身在上同调中也是平凡的。完整的方程组可理解为把这一层级继续推进，依次给出越来越高阶的同伦；细节见\bcite{55}，第 5 节。

任意严格含幺 $\Ai$ 范畴都是同伦含幺的：令 $\mu^{0,(1)}_{\A}=-e_X$，并把\eref{2.2}中所有其他满足 $i_0+\cdots+i_d>0$ 的映射取为零即可。此外，如前面的讨论所示，同伦含幺蕴含 c-含幺。\znote{23}{这三种“含幺”其实处在不同层次：严格含幺要求单位元在链层面就满足结合律，同伦含幺只要求它以链同伦的方式充当单位元，而上同调含幺只要求 $H(\A)$ 含有单位元。几何中自然出现的 Fukaya 范畴通常只是上同调含幺的，故后面以它为默认框架，必要时再用形式微分同胚换成严格含幺的模型。[来源：译者按（上下文推断，无外部依据）：依据本页 (2.1)–(2.3) 及下页对三种概念差别的讨论]}

\begin{definition*}
还应提到函子的相应概念。若 $\A,\B$ 严格含幺，则称 $\Ai$ 函子 $\F:\A\to\B$ \emph{严格含幺}，如果对所有对象 $X$ 都有 $\F^1(e_X)=e_{\F(X)}$，并且对任意 $d\ge2$ 及任意 $n$，都有
\[
\F^d(a_{d-1},\ldots,a_{n+1},e_{X_n},a_n,\ldots,a_1)=0.
\]
若 $\A,\B$ 为 c-含幺，则称 $\Ai$ 函子 $\F:\A\to\B$ 为 \emph{c-含幺}，只需 $H(\F)$ 含幺。同伦含幺版本见\bcite{55}。
\end{definition*}

严格含幺在许多方面很方便，因此在文献中广泛使用，尽管它看起来可能与同伦代数结构的一般理念不太协调。c-含幺乍看也可能令人怀疑，因为它走向另一个极端，把问题完全推到上同调层次。就我们预期的应用而言，Fukaya 范畴并不自然带有严格单位元，但很容易看出它是 c-含幺的。事实上，可以通过几何方式为它构造同伦单位元，利用附加的模空间定义\eref{2.2}中所需的、满足 $i_0+\cdots+i_d>0$ 的映射，见\bcite{60}，§20。所有这些似乎都表明同伦含幺是最恰当的概念，然而下面的讨论将出现一个或许出人意料的转折：
\SourcePage{32}
考察相关障碍后发现，存在一般的代数步骤，可以给任意 c-含幺 $\Ai$ 范畴配备同伦单位元。此外，由\bcite{55}中的推导可见，任意同伦含幺 $\Ai$ 范畴都与一个严格含幺的范畴拟同构。这意味着，三种概念之间的区别在于偏好和便利，而非实质。我们选择主要在 c-含幺的框架中工作。不过，作为证明一些基本性质的技术工具，仍将给出把 c-含幺 $\Ai$ 范畴变为严格含幺范畴的直接方法。\znote{24}{把只在上同调层面成立的单位元提升为链层面的同伦单位元，是同调扰动理论的典型应用：先在同调上选好代表，再逐阶修正高阶映射以抵消障碍，最后由收敛的无穷级数得到新结构。这里的形式微分同胚正是这种“扰动”，引理 2.1 的递归构造则是它的具体实现。[来源：nLab "homological perturbation theory"，https:/\allowbreak /\allowbreak ncatlab.org/\allowbreak nlab/\allowbreak show/\allowbreak homological+\allowbreak perturbation+\allowbreak theory（2026-11-15 访问）]}

\begin{lemma}\label{st:2.1}
设 $\A$ 是上同调含幺 $\Ai$ 范畴，并且对每个对象 $X$，或者复形 $\hom_{\A}(X,X)$ 具有非零上同调，或者对所有 $Y$ 都有 $\hom_{\A}(X,Y)=\hom_{\A}(Y,X)=0$。那么存在满足 $\Phi^1=\Id$ 的形式微分同胚 $\Phi$，使修改后的 $\Ai$ 结构 $\Phi_*\A$ 严格含幺。\znote{24}{原书引理 2.1 没有第二个条件，但该条件不可省：若范畴含零对象 $X$，则 $\hom_{\A}(X,X)$ 上同调为零，代表恒等态射的闭链 $e_X$ 只能是 $0$，而严格含幺要求 $\mu^2(e_X,a)=a$、$\mu^2(a,e_X)=a$，矛盾；故原陈述在含零对象时失效。译文补上的条件正是为排除这一退化情形：当 $X$ 与一切对象的双向 hom 复形都为零时，$e_X=0$ 已使严格含幺方程空洞地成立。此系对原书陈述的必要修补，不是原文内容。[来源：原书 p.~24 引理 2.1；反例由本书 \S1 的 $\Ai$ 范畴定义直接验证，无外部依据。]}
\end{lemma}
\begin{proof}
第一步如下。对每个对象 $X$，选择一个上闭链 $e_X\in\hom^0_{\A}(X,X)$，代表 $H(\A)$ 中的恒等态射。对于任意三个对象，可以找到一个链映射
\[
\mu^2_{\widetilde{\A}}:\hom_{\A}(X_1,X_2)\otimes\hom_{\A}(X_0,X_1)\longrightarrow\hom_{\A}(X_0,X_2),
\]
使它代表 $H(\A)$ 中的复合，并满足\eref{2.1}的第二个条件。取一个形式微分同胚，令 $\Phi^1=\Id$，并令每个 $\Phi^2$ 都是 $\mu^2_{\A}$ 与 $\mu^2_{\widetilde{\A}}$ 之间的链同伦。那么 $\widetilde{\A}=\Phi_*\A$ 的前两个复合映射就是 $\mu^1_{\widetilde{\A}}=\mu^1_{\A}$ 和选定的 $\mu^2_{\widetilde{\A}}$。

从这里起，构造递归进行。设非含幺 $\Ai$ 范畴 $\A$ 带有元素 $e_X\in\hom^0_{\A}(X,X)$，满足\eref{2.1}的前两个等式。考虑对这样的 $\A$ 给出的一族条件 $(U_{d,n})$，其中 $d>2$、$0\le n\le d$，且只考察阶数 $i>2$：
\[
\begin{aligned}
(U_{d,n})\quad&\mu^i_{\A}(a_{i-1},\ldots,a_{j+1},e_{X_j},a_j,\ldots,a_1)\text{ 在以下两种情形为零：}\\
&i<d\text{ 且 }j\text{ 任意；以及 }i=d\text{ 且 }j<n.
\end{aligned}
\]
纳入 $n=0$ 的情形纯粹是为了记号方便，因为显然 $(U_{d,0})=(U_{d-1,d-1})$。设 $\A$ 已对某个 $n<d$ 满足 $(U_{d,n})$。取一个形式微分同胚 $\Phi$，使 $\Phi^1=\Id$，对 $2\le k<d-1$ 有 $\Phi^k=0$，并令
\begin{equation}\label{eq:2.4}
\begin{aligned}
\Phi^{d-1}(a_{d-1},\ldots,a_1)&=(-1)^{\sgnmark_n}\mu^d_{\A}(a_{d-1},\ldots,a_{n+1},e_{X_n},a_n,\ldots,a_1),\\
\Phi^d(a_d,\ldots,a_1)&=(-1)^{\sgnmark_n}\mu^{d+1}_{\A}(a_d,\ldots,a_{n+1},e_{X_n},a_n,\ldots,a_1).
\end{aligned}
\end{equation}
由定义显然可见，在低于 $d-1$ 阶时，$\widetilde{\A}=\Phi_*\A$ 的复合映射与 $\A$ 的相同。利用 $\Ai$ 结合律方程及归纳假设，作一个略复杂的计算，可知 $d-1$ 阶时也相同。现在考察 $d$ 阶；为简化公式，记
\SourcePage{33}
$(b_1,\ldots,b_{d+1})=(a_1,\ldots,a_n,e_{X_n},a_{n+1},\ldots,a_d)$。可得
\begin{equation}\label{eq:2.5}
\begin{aligned}
&\mu^d_{\widetilde{\A}}(a_d,\ldots,a_1)\\
&=\mu^d_{\A}(a_d,\ldots,a_1)
-(-1)^{\S}\mu^2_{\A}\bigl(\mu^d_{\A}(a_d,\ldots,e_{X_{n+1}},a_{n+1},\ldots,a_2),a_1\bigr)\\
&\quad-(-1)^{\sgnmark_{n+1}}\sum_{j<n}(-1)^{\dagger}
\mu^d_{\A}\bigl(a_d,\ldots,e_{X_{n+1}},a_{n+1},\ldots,\mu^2_{\A}(a_{j+2},a_{j+1}),\ldots,a_1\bigr)\\
&\quad-(-1)^{\sgnmark_n}\sum_{i,j}(-1)^{\ddagger}
\mu^{d+2-i}_{\A}\bigl(b_{d+1},\ldots,\mu^i_{\A}(b_{j+i},\ldots,b_{j+1}),b_j,\ldots,b_1\bigr),
\end{aligned}
\end{equation}
其中符号为 $\S=|a_2|+\cdots+|a_{n+1}|-n$、$\dagger=|a_1|+\cdots+|a_j|-j$、$\ddagger=|b_1|+\cdots+|b_j|-j$；最后一个和式取遍以下四种可能：$i=1$ 且 $j\ne n$；$i=2$ 且 $j>n$；$i=d$ 且 $j=0$；最后是 $i=d+1$。注意，按归纳假设，最后一个和式中的表达式在 $i=1$、$j=n$ 时，以及 $3\le i\le d-1$ 时，也都为零。

先设 $n=0$，此时\eref{2.5}右侧第三个表达式为空和。去掉它，并对第四个表达式应用 $\Ai$ 结合律方程，最终得到
\[
\mu^d_{\widetilde{\A}}(a_d,\ldots,a_1)
=\mu^d_{\A}(a_d,\ldots,a_1)-\mu^2_{\A}\bigl(\mu^d_{\A}(a_d,\ldots,a_2,e_{X_1}),a_1\bigr).
\]
当 $a_1$ 也是恒等态射时，此式显然为零，表明 $\widetilde{\A}$ 具有性质 $(U_{d,1})$。接着考虑 $n>0$ 的情形。此时，\eref{2.5}最后一个表达式中，令 $i=2$ 并对所有 $j\le n$ 求和所得的和也为零；$i=d$、$j=1$ 时的项同样为零。像前面一样利用 $\Ai$ 结合律方程，可证明整个和式为零。消去这一和式之后，很容易验证 $\widetilde{\A}$ 具有性质 $(U_{d,n+1})$。

反复进行这一过程，得到一列形式微分同胚 $\Phi_1,\Phi_2,\ldots$，使得当 $k\to\infty$ 时，$(\Phi_k\circ\cdots\circ\Phi_1)_*\A$ 满足越来越多的严格含幺条件。$\Phi_k$ 本身越来越接近恒等映射；也就是说，当 $k\to\infty$ 时，$\Phi_k$ 的第一个非平凡项的阶数越来越高。因此，无穷复合 $\Phi=\cdots\circ\Phi_3\circ\Phi_2\circ\Phi_1$ 是有意义的，且 $\Phi_*\A$ 严格含幺，正如所要证明的那样。\znote{25}{形式微分同胚保持对象不变，只逐阶修改高阶复合映射，因此它是 $\A$ 与 $\Phi_*\A$ 之间在对象上为恒等的 $\Ai$ 同构。证明的每一轮都把第一个非平凡项的阶数推高，故无穷复合在每一阶上只有有限多个非零项，收敛性由此得到保证。[来源：nLab "homological perturbation theory"，https:/\allowbreak /\allowbreak ncatlab.org/\allowbreak nlab/\allowbreak show/\allowbreak homological+\allowbreak perturbation+\allowbreak theory（2026-11-15 访问）]}
\end{proof}

\begin{note}\label{st:2.2}
\stref{引理}{2.1}由 Lefèvre\bcite{95}，定理 3.2.1.1，与作者各自独立得到。还有一个更具概念性的证明，基于同调扰动理论，以及含幺代数的完整 Hochschild 复形与约化 Hochschild 复形拟同构这一事实（表达式\eref{2.4}受到这些复形之间经典同伦等价的启发）。$\Ai$ 函子、$\Ai$ 模等也有类似结果，分别见\bcite{95}，定理 3.2.2.1 和定理 3.3.1.2。
\end{note}

\SourcePage{34}
\subsection{函子范畴}\label{sec:2b}
若 $\A$ 是严格含幺 $\Ai$ 范畴，那么对任意非含幺 $\Ai$ 范畴 $\C$，$\nufun(\C,\A)$ 仍然严格含幺。具体地，严格恒等态射是自然变换 $E_{\F}$，定义为 $E^0_{\F}=e_{\F(X)}\in\hom^0_{\A}(\F(X),\F(X))$，以及对所有 $d>0$ 令 $E^d_{\F}=0$。还有一个稍不显然的结论。

\begin{lemma}\label{st:2.3}
若 $\A$ 为 c-含幺，则 $\Q=\nufun(\C,\A)$ 也为 c-含幺。此外，若 $E_{\F}\in\hom^0_{\Q}(\F,\F)$ 是代表 $H(\Q)$ 中恒等态射的任意链，那么对每个 $X\in\Ob\C$，$E^0_{\F}\in\hom^0_{\A}(\F(X),\F(X))$ 代表 $H(\A)$ 中的恒等态射。
\end{lemma}
\begin{proof}
可将其归结为严格含幺的情形。先利用\stref{笔记}{1.13}取 $\A$ 的最小模型，再应用\stref{引理}{2.1}严格化，得到与 $\A$ 拟同构的严格含幺范畴 $\widetilde{\A}$，以及一个在对象上为恒等的拟同构 $\Phi:\A\to\widetilde{\A}$。通过诱导的上同调同构，将 $H(\widetilde{\A})$ 与 $H(\A)$ 识别。令 $\widetilde{\Q}=\nufun(\C,\widetilde{\A})$。考虑与 $\Phi$ 作左复合的 $\Ai$ 函子，并记 $\widetilde{\F}=\Phi\circ\F$。在自然变换层次上，得到交换图
\SourceEquationDiagram{diagram-2-6}{0.55038403}{2.6}
这里 $C=H(\C)$、$A=H(\A)=H(\widetilde{\A})$、$F=H(\F)=H(\widetilde{\F})$ 是上同调层次的结构，向下的箭头则是\eref{1.10}中的映射，把 $\Ai$ 函子的自然变换送到相应的上同调自然变换。\stref{引理}{1.7}说明 $H(\LG^1_{\Phi})$ 是同构。选择一个链 $E_{\F}$，使 $[\LG^1_{\Phi}(E_{\F})]$ 为严格单位元的类 $[E_{\widetilde{\F}}]$。由 $H(\LG^1_{\Phi})$ 的函子性，$[E_{\F}]$ 是 $H(\Q)$ 中的恒等态射。在\eref{2.6}中追图可知，$H(E_{\widetilde{\F}})=H(E_{\F})$ 是 $\operatorname{Nu\text{-}Fun}(C,A)$ 中的恒等自然变换。这就表明，对每个 $X$，在上述识别下 $[E^0_{\widetilde{\F}}]=[E^0_{\F}]$ 是 $A$ 中的恒等态射，正如所要证明的那样。
\end{proof}

\begin{lemma}\label{st:2.4}
设 $\G:\A\to\B$ 是 c-含幺 $\Ai$ 函子。那么对任意非含幺 $\Ai$ 范畴 $\C$，复合函子 $\LG_{\G}:\nufun(\C,\A)\to\nufun(\C,\B)$ 为 c-含幺。（右复合也有类似结果：设 $\G:\A\to\B$ 是非含幺 $\Ai$ 范畴之间的非含幺 $\Ai$ 函子。那么，对任意 c-含幺 $\Ai$ 范畴 $\C$，复合函子 $\RG_{\G}:\nufun(\B,\C)\to\nufun(\A,\C)$ 为 c-含幺。）
\end{lemma}
\begin{proof}
设 $E_{\F}\in\hom^0_{\nufun(\C,\A)}(\F,\F)$ 是代表 $H(\nufun(\C,\A))$ 中恒等态射的链。取其在复合下的像：
\begin{equation}\label{eq:2.7}
[\LG^1_{\G}(E_{\F})]\in\Hom^0_{H(\nufun(\C,\B))}(\G\circ\F,\G\circ\F).
\end{equation}
我们已经看到，对每个 $X$，$E^0_{\F}\in\hom^0_{\A}(\F(X),\F(X))$ 代表上同调恒等态射，因此 $\LG^1_{\G}(E_{\F})^0=\G^1(E^0_{\F})\in\hom^0_{\B}(\G(\F(X)),\G(\F(X)))$ 也代表上同调恒等态射。
\SourcePage{35}
利用\stref{引理}{1.6}及 $\nufun(\C,\B)$ 的 c-含幺性，可知\eref{2.7}是可逆态射。由于 $\LG_{\G}$ 的函子性，它还是幂等的，故必等于恒等态射（后面会反复使用这一类论证）。对于 $\RG_{\G}$，证明相同。
\end{proof}

\subsection{同伦与同构}\label{sec:2c}
设 $\A$ 是严格含幺 $\Ai$ 范畴，$\C$ 是非含幺 $\Ai$ 范畴，而 $\F_0,\F_1:\C\to\A$ 是在对象上作用相同的非含幺 $\Ai$ 函子。设 $T$ 是从 $\F_0$ 到 $\F_1$ 的零次预自然变换，满足 $T^0=0$。用 $S$ 表示修改后的预自然变换，定义为对所有 $X$ 令 $S^0=e_{\F_0(X)}$，且当 $d>0$ 时令 $S^d=T^d$。直接计算表明，$T$ 是 $\F_0$ 与 $\F_1$ 之间的同伦，当且仅当 $S$ 是真正的自然变换。在这一背景下，原本看似神秘的同伦复合表达式\eref{1.16}，可以看作相关自然变换的通常复合；\eref{1.17}也有类似解释。

利用\stref{引理}{1.6}及 $\Ai$ 函子范畴的严格含幺性，可知两个同伦函子作为 $H^0(\nufun(\C,\A))$ 的对象是同构的。我们还需要 c-含幺背景下的类似结果。

\begin{lemma}\label{st:2.5}
设 $\A$ 是 c-含幺 $\Ai$ 范畴，$\C$ 是非含幺 $\Ai$ 范畴。若两个非含幺 $\Ai$ 函子 $\F_0,\F_1:\C\to\A$ 同伦，那么它们作为 $H^0(\nufun(\C,\A))$ 的对象也是同构的。
\end{lemma}
\begin{proof}
沿用\stref{引理}{2.3}证明中的记号。由于左复合保持同伦，$\widetilde{\F}_k=\Phi\circ\F_k$ 同伦，因此由前面的观察，它们作为 $H^0(\widetilde{\Q})$ 的对象同构。另一方面，由\stref{引理}{1.7}，$\LG_{\Phi}$ 上同调全忠实，所以 $\F_k$ 本身必然同构。
\end{proof}

\subsection{Morita 不变性}\label{sec:2d}
设 $A,B$ 是线性分次含幺范畴，$F_0,F_1:A\to B$ 是线性分次含幺函子。设子范畴 $\widetilde A\subset A$ 的包含函子为等价。

\begin{lemma}\label{st:2.6}
限制映射给出 Hochschild 上同调群的同构 $HH(A,B)\cong HH(\widetilde A,B)$。
\end{lemma}
\begin{proof}
设 $\rho:CC(A,B)\to CC(\widetilde A,B)$ 是链层次的限制映射。对任意 $X\in\Ob A$，选择一个 $\widetilde X\in\Ob\widetilde A$，以及同构 $u_X:X\to\widetilde X$。对于已属于 $\widetilde A$ 的对象 $X$，可以取 $\widetilde X=X$，并令 $u_X$ 为恒等态射。\emph{Dennis 余迹}是如下链映射：
\[
\begin{aligned}
\chi &:CC^{r+s}(\widetilde A,B)^s\longrightarrow CC^{r+s}(A,B)^s,\\
(\chi T)^r(a_r,\ldots,a_1)
&=F_1(u_{X_r})^{-1}T^r(u_{X_r}a_ru_{X_{r-1}}^{-1},\ldots,u_{X_1}a_1u_{X_0}^{-1})F_0(u_{X_0}).
\end{aligned}
\]
\SourcePage{36}
它是同伦等价，因为 $\rho\circ\chi=\id$，且 $\id-\chi\circ\rho=\partial\eta+\eta\partial$，其中
\[
\begin{aligned}
(\eta T)(a_{r-1},\ldots,a_1)
&=\sum_n(-1)^{\sgnmark_n+(r+s)}F_1(u_{X_{r-1}})^{-1}\\
&\qquad\cdot T(u_{X_{r-1}}a_{r-1}u_{X_{r-2}}^{-1},\ldots,u_{X_{n+1}}a_{n+1}u_{X_n}^{-1},u_{X_n},a_n,\ldots,a_1).
\end{aligned}
\]
\end{proof}

\begin{note}\label{st:2.7}
也可以采用以下更抽象的途径。考虑分次 $A$ 双模组成的Abel 范畴 $\operatorname{Bimod}^0(A)$；更精确地说，其对象是从 $(A^{\mathrm{opp}},A)$ 到分次向量空间范畴的双函子（双模同态就是零次自然变换）。任意一对函子 $(G_0,G_1)$ 都给出一个双模，由向量空间 $\Hom_B(G_0(X_0),G_1(X_1))$ 组成。略微滥用记号，把这个双模直接记作 $B$，而把考察函子对 $(\Id_A,\Id_A)$ 得到的相应对象记作 $A$。Hochschild 复形可以解释为来自最后这个双模在 $\operatorname{Bimod}^0(A)$ 中的一个投射分解，因此
\begin{equation}\label{eq:2.8}
HH(A,B)\cong\operatorname{Ext}_{\operatorname{Bimod}^0(A)}(A,B).
\end{equation}
这里考察所有阶数的 Ext 群，并把 $B$ 的给定分次作所有可能的平移，所以右侧是双分次的。由初等理由，将双函子从 $A$ 限制到 $\widetilde A$，给出范畴等价 $\operatorname{Bimod}^0(A)\cong\operatorname{Bimod}^0(\widetilde A)$，从而得到 Ext 群的同构。

现在，这与代数的 Hochschild 同调和上同调的经典 Morita 不变性之间的关系应当已经清楚，例如参看\bcite{143}，定理 9.5.6。事实上，我们的情形类似于矩阵代数给出的 Morita 不变性的基本例子，上面的初等论证沿用了\bcite{96}，定理 1.2.4 的证明。
\end{note}

\subsection{拟等价}\label{sec:2e}
\begin{definition*}
还需要补充一点记号和术语：给定两个 c-含幺 $\Ai$ 范畴 $\A,\B$，用 $\operatorname{fun}(\A,\B)$ 表示 $\nufun(\A,\B)$ 中由 c-含幺 $\Ai$ 函子组成的全子范畴。若 c-含幺 $\Ai$ 函子 $\F:\A\to\B$ 的基础上同调函子为等价，则称 $\F$ 为\emph{拟等价}。
\end{definition*}

设 $\A$ 是 c-含幺 $\Ai$ 范畴，$\widetilde{\A}\subset\A$ 为全 $\Ai$ 子范畴，且包含函子为拟等价。
\begin{lemma}\label{st:2.8}
对于任意 c-含幺 $\Ai$ 范畴 $\B$，将 c-含幺 $\Ai$ 函子及其预自然变换作限制，给出拟等价 $\Q=\operatorname{fun}(\A,\B)\to\widetilde{\Q}=\operatorname{fun}(\widetilde{\A},\B)$。
\end{lemma}
\begin{proof}
考虑任意两个 c-含幺 $\Ai$ 函子 $\G_k$（$k=0,1$），以及它们在 $\widetilde{\A}$ 上的限制 $\widetilde{\G}_k$。预自然变换上的限制映射
\begin{equation}\label{eq:2.9}
\hom_{\Q}(\G_0,\G_1)\longrightarrow\hom_{\widetilde{\Q}}(\widetilde{\G}_0,\widetilde{\G}_1)
\end{equation}
\SourcePage{37}
显然与长度滤过相容，因此诱导相应谱序列之间的映射。在 $E_2$ 层次上，它归结为相关上同调函子的 Hochschild 上同调上的限制映射；我们刚刚已经证明，这个映射是同构。因此\eref{2.9}是拟同构，也就是说，限制函子上同调全忠实。最后，给定任意 c-含幺 $\Ai$ 函子 $\widetilde{\G}:\widetilde{\A}\to\B$，可以把其基础上同调函子扩张为 $G:H(\A)\to H(\B)$，并且这一扩张在自然同构意义下唯一；随后\stref{引理}{1.10}允许构造一个扩张 $\widetilde{\G}$、满足 $H(\G)=G$ 的 $\G:\A\to\B$，因而它为 c-含幺。这证明限制函子在对象上为满射，故为拟等价。
\end{proof}

特别地，扩张步骤说明存在 c-含幺 $\Ai$ 函子 $\mathcal P:\A\to\widetilde{\A}$，满足 $\mathcal P|\widetilde{\A}=\Id_{\widetilde{\A}}$。若用 $\mathcal K:\widetilde{\A}\to\A$ 表示包含函子，则 $\mathcal P\circ\mathcal K=\Id_{\widetilde{\A}}$。$\mathcal P$ 是拟等价，特别地上同调全忠实，因此由\stref{引理}{1.7}，$\LG_{\mathcal P}$ 也上同调全忠实。由于 $\LG_{\mathcal P}(\mathcal K\circ\mathcal P)=\LG_{\mathcal P}(\Id_{\A})$，可知 $\mathcal K\circ\mathcal P$ 与 $\Id_{\A}$ 作为 $H^0(\operatorname{fun}(\A,\A))$ 的对象同构。\znote{29}{在链复形层面，诱导同调同构的链映射称为拟同构；它一般不可逆，“存在从 $C$ 到 $D$ 的拟同构”甚至不是对称关系。拟等价要求上同调函子本身为等价，因而天然可逆，\stref{定理}{2.9}正是这一点的定量表述。[来源：nLab "quasi-isomorphism"，https:/\allowbreak /\allowbreak ncatlab.org/\allowbreak nlab/\allowbreak show/\allowbreak quasi-isomorphism（2026-11-15 访问）]}

\begin{theorem}\label{st:2.9}
设 $\F:\A\to\B$ 为拟等价。那么存在拟等价 $\G:\B\to\A$，使得 $\G\circ\F\cong\Id_{\A}$ 在 $H^0(\operatorname{fun}(\A,\A))$ 中成立，且 $\F\circ\G\cong\Id_{\B}$ 在 $H^0(\operatorname{fun}(\B,\B))$ 中成立。
\end{theorem}
\begin{proof}
存在全 $\Ai$ 子范畴 $\widetilde{\A}\subset\A$、$\widetilde{\B}\subset\B$，使包含函子为拟等价，并使 $\F$ 限制为从 $\widetilde{\A}$ 到 $\widetilde{\B}$ 的拟同构 $\widetilde{\F}$。如前所述，两个包含函子都能在各自的函子范畴中取到同构意义下的逆。由\stref{推论}{1.14}，$\widetilde{\F}$ 本身具有同伦意义下的逆；再由\stref{引理}{2.5}，同伦蕴含函子范畴中的同构。
\end{proof}

\begin{lemma}\label{st:2.10}
若 $\G:\A\to\B$ 为拟等价，那么对于任意 c-含幺 $\Ai$ 范畴 $\C$，$\RG_{\G}:\operatorname{fun}(\B,\C)\to\operatorname{fun}(\A,\C)$ 也为拟等价。（$\LG_{\G}$ 也有相应结论；事实上，后一情形还存在适用于非含幺 $\C$ 及范畴 $\nufun$ 的版本。）
\end{lemma}
\begin{proof}
首先需要知道 $\RG_{\G}$ 上同调全忠实。若 $\G$ 是拟同构，这是\stref{引理}{1.7}的一部分；若 $\G$ 是全子范畴的包含函子，这是\stref{引理}{2.8}的一部分。与前面的讨论一样，这两个特殊情形推出一般结论。接着，\stref{定理}{2.9}给出一个 $\F:\B\to\A$，使 $\F\circ\G\cong\Id_{\A}$。由\stref{引理}{2.4}，对任意 $\mathcal H\in\Ob\operatorname{fun}(\A,\C)$，有
\[
\RG_{\G}(\mathcal H\circ\F)=\LG_{\mathcal H}(\F\circ\G)
\cong\LG_{\mathcal H}(\Id_{\A})=\mathcal H,
\]
这表明 $\RG_{\G}$ 在对象的同构类上为满射。$\LG_{\G}$ 的论证类似，但略为简单。
\end{proof}

\begin{note}\label{st:2.11}
\stref{定理}{2.9}的严格含幺版本，似乎最早由 Kontsevich 提出（未发表）。它出现于\bcite{55}，定理 8.6，以及\bcite{95}，定理 9.2.04。这里讨论的 c-含幺版本也见于\bcite{97}，定理 8.8。
\end{note}

\SourcePage{38}
\subsection{\texorpdfstring{$A_\infty$}{A∞} 模}\label{sec:2f}
设 $\A$ 是 c-含幺 $\Ai$ 范畴。由于 $\Ch$ 严格含幺，$\Q=\operatorname{nu\text{-}mod}(\A)=\nufun(\A^{\mathrm{opp}},\Ch)$ 也严格含幺。注意，由于\eref{1.20}中引入的符号，$\M$ 的严格恒等态射是如下模同态 $t=e_{\M}$：
\begin{equation}\label{eq:2.10}
t^1(b)=(-1)^{|b|}b,\qquad t^d=0\quad\text{对所有 }d\geq2.
\end{equation}
\begin{definition*}
若 $\Ai$ 模 $\M$ 的基础上同调模含幺，就称 $\M$ 为 \emph{c-含幺}。这意味着：若 $e_X\in\hom^0_{\A}(X,X)$ 代表上同调单位元，则链映射 $\mu^2_{\M}(\cdot,e_X):\M(X)\to\M(X)$ 在上同调上诱导恒等映射。用 $\modcat(\A)$ 表示由 c-含幺 $\Ai$ 模组成的全子范畴。
\end{definition*}
若 $\G:\A\to\B$ 为 c-含幺 $\Ai$ 函子，可以把拉回\eref{1.22}限制到这些子范畴，得到
\begin{equation}\label{eq:2.11}
\G^*:\modcat(\B)\longrightarrow\modcat(\A),
\end{equation}
它是严格含幺 dg 函子。作为\stref{引理}{2.10}的特殊情形，若 $\G$ 为拟等价，那么\eref{2.11}也为拟等价。

\subsection{Yoneda 嵌入}\label{sec:2g}
由定义，对 c-含幺的 $\A$，Yoneda 嵌入 $\mathcal I=\mathcal I_{\A}$ 的像落在子范畴 $\modcat(\A)\subset\operatorname{nu\text{-}mod}(\A)$ 中。
\begin{lemma}\label{st:2.12}
对任意 c-含幺 $\Ai$ 模 $\M$ 及 $\A$ 的任意对象 $Y$，\eref{1.25}中的映射 $\lambda$ 是拟同构。
\end{lemma}
\begin{proof}
由定义，$\lambda$ 的映射锥平移 $-1$ 后，是链复形
\begin{equation}\label{eq:2.12}
\left(\M(Y)\oplus\hom_{\Q}(\Y,\M)[-1],\begin{pmatrix}\mu^1_{\M}&0\\\lambda&-\mu^1_{\Q}\end{pmatrix}\right).
\end{equation}
先取子复形 $\hom_{\Q}(\Y,\M)[-1]$，再接上它通常的长度滤过，以此给这个复形赋予滤过。记 $A=H(\A)$、$M=H(\M)$，所得谱序列满足
\[
E_1^{rs}=\begin{cases}
M^s(Y),&r=0,\\[3pt]
\displaystyle\prod_{X_0,\ldots,X_{r-1}}\Hom^s_{\K}\!\left(\begin{aligned}&\Hom_A(X_{r-1},Y)\otimes\Hom_A(X_{r-2},X_{r-1})\\[-2pt]&\qquad\otimes\cdots\otimes\Hom_A(X_0,X_1),\,M(X_0)\end{aligned}\right),&r>0.
\end{cases}
\]
由\eref{1.25}和\eref{1.21}可见，微分 $\partial=\partial^{rs}:E_1^{rs}\to E_1^{r+1,s}$ 在 $r=0$ 时由 $(\partial c)(b)=(-1)^{|b|}cb$ 给出，而在 $r>0$ 时由
\[
\begin{aligned}
(\partial t)(b,a_r,\ldots,a_1)
&=(-1)^{\sgnmark_r+|b|}t(b,a_r,\ldots,a_2)a_1\\
&\quad+(-1)^{|b|}t(ba_r,a_{r-1},\ldots,a_1)\\
&\quad+\sum_n(-1)^{\triangle} t(b,a_r,\ldots,a_{n+2}a_{n+1},\ldots,a_1)
\end{aligned}
\]
\SourcePage{39}
给出，其中 $\triangle=|a_{n+2}|+\cdots+|a_r|+|b|-r+n+1$。这是标准的条形分解。当 $A$ 具有恒等态射，且这些态射在 $M$ 上作为恒等映射作用时，它具有收缩同伦
\[
(\kappa t)(b,a_{r-1},\ldots,a_1)=t(e_Y,b,a_{r-1},\ldots,a_1).
\]
因此\eref{2.12}无上同调，从而得到所需结论。
\end{proof}

专门考察来自 Yoneda 嵌入的 $\Ai$ 模 $\M$，此时 $\lambda=\mathcal I^1$，因此：
\begin{corollary}\label{st:2.13}
$\mathcal I$ 上同调全忠实（从而也 c-含幺）。
\end{corollary}
考察 $\mathcal I$ 的像，还可得到：
\begin{corollary}\label{st:2.14}
任意 c-含幺 $\Ai$ 范畴，都以典范方式拟同构于一个严格含幺 dg 范畴。
\end{corollary}

现在考察 Yoneda 嵌入对于 c-含幺 $\Ai$ 范畴之间的上同调全忠实 $\Ai$ 函子 $\F:\A\to\B$ 的自然性。对任意这样的函子 $\F$，由\stref{引理}{1.16}，\eref{1.23}中自然变换 $T$ 的一阶项是 c-含幺 $\Ai$ 模的同构。结合\stref{引理}{1.6}，可知 $T$ 本身在 $\Ai$ 函子的基础上同调范畴中是同构。换言之，图
\SourceEquationDiagram{diagram-2-13}{0.28246124}{2.13}
在 $H^0(\operatorname{fun}(\A,\modcat(\A)))$ 中在典范同构意义下交换。

\begin{note}\label{st:2.15}
仿照\stref{笔记}{2.7}的思路，可以把\stref{引理}{2.12}的证明重新表述得更抽象一些。设 $\operatorname{Mod}^0(A)$ 为从 $A$ 到分次向量空间的反变（含幺）函子组成的Abel 范畴。若 $\M_0,\M_1$ 为 c-含幺 $\Ai$ 模，它们的基础上同调模 $M_0,M_1$ 就是 $\operatorname{Mod}^0(A)$ 的对象。借助条形分解，由 $\hom_{\Q}(\M_0,\M_1)$ 的长度滤过得到的谱序列，其 $E_2$ 项可解释为
\begin{equation}\label{eq:2.14}
E_2=\operatorname{Ext}_{\operatorname{Mod}^0(A)}(M_0,M_1),
\end{equation}
右侧的双分次方式与\eref{2.8}相同。与 Yoneda $\Ai$ 模关联的上同调模，是 $\operatorname{Mod}^0(A)$ 的投射对象。因此，若 $\M_0=\Y$，则\eref{2.14}只有一列非零，而这一列可与 $M_1(Y)$ 等同。所以 $\M_1(Y)$ 拟同构于 $\hom_{\Q}(\Y,\M_1)$。
\end{note}

\SourcePage{40}
\subsection{术语约定}\label{sec:2h}
关于恒等态射的问题已经讨论完毕；今后我们将不再在术语中强调这一点，而采用如下约定：\emph{除非另作说明，以下所有 $\Ai$ 范畴、$\Ai$ 函子和 $\Ai$ 模均假定为 c-含幺。}

\section{正合三角}\label{sec:3}
\subsection{基本思路}\label{sec:3a}
在任何一种范畴中构造对象时，常见的策略都是：先用所期望的表示函子刻画所需对象，再研究存在性（可表示性）问题。在我们的情形中，适当的形式体系是 $\Ai$ 模与 Yoneda 嵌入。固定一个 $\Ai$ 范畴 $\A$，记 $\Q=\modcat(\A)$。设 $\M$ 是一个 $\Ai$ 模。假设存在 $Y\in\Ob\A$，其 Yoneda 嵌入下的像 $\Y$ 在 $H^0(\Q)$ 中与 $\M$ 同构，并固定该范畴中的一个同构 $[t]:\Y\to\M$。
\begin{definition*}
称二元组 $(Y,[t])$ \emph{拟表示} $\M$。
\end{definition*}
\stref{推论}{2.13}表明，若 $(\widetilde Y,[\tilde t])$ 是另一个这样的二元组，那么在 $H^0(\A)$ 中存在一个优选的同构 $[a]:Y\to\widetilde Y$，满足 $[\tilde t]\cdot[\mathcal I^1(a)]=[t]$。\znote{32}{“拟表示”把通常的表示概念在同伦意义下放宽：并不要求 $\hom_{\Q}(\Y,\M)$ 与 $\M(Y)$ 同构，只要求 Yoneda 像 $\Y$ 与 $\M$ 在 $H^0(\Q)$ 中同构。这样在作映射锥、平移等构造时，每个对象只需在“相差一个同构”的意义下被表示，正是构造三角化包络所需要的灵活性。[来源：译者按（上下文推断，无外部依据）：对照\stref{引理}{3.1}与\stref{推论}{2.13}]}略微滥用术语，有时我们会直接称 $Y$ 为拟表示 $\M$ 的对象，但仍把 $[t]$ 的选择视为数据的一部分。利用\eref{1.25}中的链映射 $\lambda=\lambda_{\M}:\M(Y)\to\hom_{\Q}(\Y,\M)$，可以给出更具体的刻画：
\begin{lemma}\label{st:3.1}
$(Y,[t])$ 拟表示 $\M$，当且仅当如下条件成立：存在 $c\in\M^0(Y)$，满足 $\mu^1_{\M}(c)=0$ 以及 $[t]=[\lambda(c)]$，并且对所有 $X\in\Ob\A$，映射
\begin{equation}\label{eq:3.1}
\hom_{\A}(X,Y)\longrightarrow\M(X),\qquad b\longmapsto(-1)^{|b|}\mu^2_{\M}(c,b)
\end{equation}
都是拟同构。
\end{lemma}
\begin{proof}
由于 $\mu^1_{\M}(c)=0$，$\lambda(c)$ 是一个 $\Ai$ 模同态 $\Y\to\M$，其一阶项除符号外就是\eref{3.1}。\stref{引理}{1.16}表明 $[t]=[\lambda(c)]$ 在 $H^0(\Q)$ 中是同构，因此 $(Y,[t])$ 拟表示 $\M$。反过来，若已有拟表示二元组 $(Y,[t])$，那么由\stref{引理}{2.12}，存在 $c$ 满足 $\mu^1_{\M}(c)=0$ 及 $[t]=[\lambda(c)]$。由于 $[\lambda(c)]$ 在 $H^0(\Q)$ 中是同构，它必然在基础上同调模之间诱导同构，这就意味着\eref{3.1}对所有 $X$ 都是拟同构。
\end{proof}

\subsection{直和}\label{sec:3b}
下面是这种方法的一个初等例子。给定 $\Ai$ 模 $\M_0,\M_1$，把相应的分次向量空间取直和 $\M_0(X)\oplus\M_1(X)$（$X\in\Ob\A$），并赋予结构映射 $\mu^d_{\M_0}\oplus\mu^d_{\M_1}$，即可定义它们的直和 $\M_0\oplus\M_1$。现在，设 $Y_0,Y_1$ 为 $\A$ 的对象，$\Y_0,\Y_1$
\SourcePage{41}
为其 Yoneda 嵌入下的像。
\begin{definition*}
把任意一个拟表示 $\Ai$ 模 $\Y_0\oplus\Y_1$ 的对象记为 $Y_0\oplus Y_1$。
\end{definition*}
由定义，对每个 $X$，它都带有典范同构
\[
\Hom_{H(\A)}(X,Y_0\oplus Y_1)\cong\Hom_{H(\A)}(X,Y_0)\oplus\Hom_{H(\A)}(X,Y_1).
\]

\subsection{张量积}\label{sec:3c}
同样在这个基本层次上，设 $Z=(Z,\partial_Z)$ 是 $\K$ 向量空间的链复形，$\M$ 是 $\Ai$ 模。
\begin{definition*}
张量积模 $Z\otimes\M$ 定义为
\begin{equation}\label{eq:3.2}
\begin{aligned}
(Z\otimes\M)(X)&=Z\otimes\M(X),\\
\mu^1_{Z\otimes\M}(z\otimes b)&=(-1)^{|b|-1}\partial_Z(z)\otimes b+z\otimes\mu^1_{\M}(b),\\
\mu^d_{Z\otimes\M}(z\otimes b,a_{d-1},\ldots,a_1)&=z\otimes\mu^d_{\M}(b,a_{d-1},\ldots,a_1)\quad(d\geq2).
\end{aligned}
\end{equation}

\end{definition*}
可以使 $Z\otimes-$ 成为 $\Q$ 到自身的 $\Ai$ 函子。在一阶上，它把 $t\in\hom_{\Q}(\M_0,\M_1)$ 映到 $\id_Z\otimes t\in\hom_{\Q}(Z\otimes\M_0,Z\otimes\M_1)$，后者由下式定义：
\[
(\id_Z\otimes t)^d(z\otimes b,a_{d-1},\ldots,a_1)=(-1)^{|z|(|t|-1)}z\otimes t^d(b,a_{d-1},\ldots,a_1);
\]
所有高阶项均置为零。给定 $\A$ 的对象 $Y$，其 Yoneda 像为 $\Y$；把任意一个拟表示 $Z\otimes\Y$ 的对象记为 $Z\otimes Y$。在上同调层次上，有
\[
\begin{aligned}
\Hom_{H(\A)}(X,Z\otimes Y)&\cong\Hom_{H(\Q)}(X,Z\otimes\Y)\cong H((Z\otimes\Y)(X))\\
&\cong H(Z)\otimes H(\Y(X))\cong H(Z)\otimes\Hom_{H(\A)}(X,Y),
\end{aligned}
\]
其中的同构分别来自 Yoneda 嵌入、映射 $\lambda$ 以及通常的 Künneth 公式。

\begin{note}\label{st:3.2}
给定链复形 $(Z,\partial_Z)$，令 $\widetilde Z=H(Z)$ 且 $\partial_{\widetilde Z}=0$，再选择一个线性映射 $\psi:\widetilde Z\to Z$，为每个上同调类选取链代表。由\stref{引理}{1.16}，所得嵌入 $\widetilde Z\otimes\M\hookrightarrow Z\otimes\M$ 在 $H^0(\Q)$ 中诱导同构（而且这个同构不依赖于 $\psi$ 的选择）。这说明，把张量积运算限制在 $Z$ 只是一个微分为零的分次 $\K$ 向量空间的情形，并无本质上的一般性损失。
\end{note}

\subsection{平移}\label{sec:3d}
\begin{definition*}
最常用的张量积例子，是 $Z$ 为集中在某个次数 $-\sigma$ 的一维向量空间 $\K[\sigma]$ 的情形。所得对象分别记为 $\K[\sigma]\otimes\M=\ModShift^{\sigma}\M$ 与 $\K[\sigma]\otimes Y=S^{\sigma} Y$，称为 $\M$ 或 $Y$ 的
\SourcePage{42}
$\sigma$ 次\emph{平移}。当 $\sigma=1$ 时省略指标，只写 $\ModShift\M$、$SY$。
\end{definition*}
这些平移如果存在，就带有典范同构
\begin{align}
\Hom_{H(\A)}(Y_0,S^{\sigma} Y_1)&\cong\Hom_{H(\A)}(Y_0,Y_1)[\sigma],\label{eq:3.3}\\
\Hom_{H(\A)}(Y_0,Y_1)&\cong\Hom_{H(\A)}(S^{\sigma} Y_0,S^{\sigma} Y_1),\label{eq:3.4}\\
\Hom_{H(\A)}(S^{\sigma} Y_0,Y_1)[\sigma]&\cong\Hom_{H(\A)}(Y_0,Y_1).\label{eq:3.5}
\end{align}
第一个来自定义，第二个来自 $\ModShift^{\sigma}=\K[\sigma]\otimes-$ 的函子性。第三个则先对二元组 $(S^{\sigma} Y_0,Y_1)$ 应用\eref{3.3}，再对 $(Y_0,Y_1)$ 反向（从右到左）应用\eref{3.4}而得。\znote{34}{这类符号差异在三角范畴中是常态：一个三角 $(f,g,h)$ 换成 $(f,g,-h)$ 一般不再是殊三角，尽管它仍诱导出长正合列；而 $(f,-g,-h)$ 因与原三角同构而仍是殊三角。因此使用平移同构时必须固定符号约定，否则不同来源的公式会相差 $(-1)^{\sigma}$。[来源：nLab "triangulated category"，https:/\allowbreak /\allowbreak ncatlab.org/\allowbreak nlab/\allowbreak show/\allowbreak triangulated+\allowbreak category（2026-11-15 访问）]}

\clearpage
\begin{note}\label{st:3.3}
得到\eref{3.5}的另一种途径，是从 $(S^{\sigma} Y_0,Y_1)$ 出发，应用\eref{3.4}得到 $(Y_0,S^{-\sigma}Y_1)$，然后反向使用\eref{3.3}回到 $(Y_0,Y_1)$。遗憾的是，这与前面的定义相差一个符号 $(-1)^{\sigma}$，因此使用这些同构时需要稍加留意。
\end{note}

假设 $\Ai$ 范畴 $\A$ 满足：对所有 $Y$，$SY$ 都存在。那么可以构造一个平移 $\Ai$ 函子 $S:\A\to\A$，它在 $H^0(\operatorname{fun}(\A,\A))$ 中的同构意义下唯一，并使图
\SourceDiagram{diagram-3d}{0.13915370}
在 $H^0(\operatorname{fun}(\A,\Q))$ 中在同构意义下交换。具体地，考虑由所有形如 $SY$ 的对象组成的全 $\Ai$ 子范畴 $\widetilde{\A}\subset\A$；相应地，令 $\widetilde{\Q}\subset\Q$ 为由如下 $\Ai$ 模组成的子范畴：它们在 $H^0(\Q)$ 中同构于某个 $\ModShift\Y=\ModShift\mathcal I(Y)$。由构造，Yoneda 嵌入的限制是拟等价 $\widetilde{\A}\to\widetilde{\Q}$，故由\stref{定理}{2.9}可取其逆。给定一个这样的逆 $\mathcal K$，令 $S=\mathcal K\circ\ModShift\circ\mathcal I$ 即可。这里以平移为例作了说明，但同样的论证适用于 $\Ai$ 模的其他函子性构造。\znote{34}{这里构造的 $S$ 就是经典三角范畴中的平移函子：代数情形常记 $(-)[1]$，也写作 $S$ 或 $\Sigma$，它在三角范畴中应当是范畴的自等价。本节的论证说明，只要每个对象的平移都存在，就能把这一运算提升为 $\Ai$ 函子，从而在 $H^0$ 上得到真正的自等价。[来源：nLab "triangulated category"，https:/\allowbreak /\allowbreak ncatlab.org/\allowbreak nlab/\allowbreak show/\allowbreak triangulated+\allowbreak category（2026-11-15 访问）]}

\begin{note}\label{st:3.4}
任意平移 $S^{\sigma} Y$（$\sigma\in\Z$）和直和 $Y_0\oplus Y_1$ 的存在性，蕴含对于上同调有限维的复形 $(Z,\partial_Z)$，张量积 $Z\otimes Y$ 也存在。具体地，在 $\widetilde Z=H(Z)$ 中选取一组基 $\{[z^i]\}$，令 $\sigma^i=-|z^i|$，则
\[
\bigoplus_{i\in I}S^{\sigma^i}Y
\]
在 Yoneda 嵌入下映到 $\bigoplus_{i\in I}\ModShift^{\sigma^i}\Y\cong\widetilde Z\otimes\Y$。由\stref{笔记}{3.2}，它也是 $Z\otimes\Y$ 的拟表示对象。
\end{note}
\SourcePage{43}
\subsection{锥}\label{sec:3e}
设 $Y_0,Y_1$ 为 $\A$ 的对象，$c\in\hom^0_{\A}(Y_0,Y_1)$ 为零次上闭链（$\mu^1_{\A}(c)=0$）。
\begin{definition*}
$c$ 的\emph{抽象映射锥}是如下定义的 $\Ai$ 模 $\C=\ACone(c)$：
\begin{equation}\label{eq:3.6}
\begin{aligned}
\C(X)&=\hom_{\A}(X,Y_0)[1]\oplus\hom_{\A}(X,Y_1),\\
\mu^d_{\C}((b_0,b_1),a_{d-1},\ldots,a_1)
&=\bigl(\mu^d_{\A}(b_0,a_{d-1},\ldots,a_1),\\[-2pt]
&\qquad\mu^d_{\A}(b_1,a_{d-1},\ldots,a_1)+\mu^{d+1}_{\A}(c,b_0,a_{d-1},\ldots,a_1)\bigr).
\end{aligned}
\end{equation}

\end{definition*}
直接计算即可验证\eref{1.19}。至于 $\C$ 的 c-含幺性，需要检验：若 $e_X\in\hom^0_{\A}(X,X)$ 代表 $H(\A)$ 中的恒等态射，则 $(b_0,b_1)\mapsto\mu^2_{\C}((b_0,b_1),e_X)$ 在 $H(\C(X),\mu^1_{\C})$ 上诱导的自同态是恒等映射。我们已经知道它是幂等的；考察 $\C(X)$ 中由 $b_0=0$ 给出的子复形及其商，又可知它可逆，因此它就是恒等映射。
\begin{definition*}
把拟表示抽象映射锥的对象记为 $\Cone(c)$，简称 $c$ 的\emph{映射锥}。
\end{definition*}
与通常的情形一样，这样的对象如果存在，就在 $H^0(\A)$ 中的典范同构意义下唯一。\znote{35}{抽象映射锥以链复形的映射锥为原型：对链映射 $f:X\to Y$，$\Cone(f)$ 的第 $n$ 项是 $Y_n\oplus X_{n-1}$，其微分的左上、右下两块分别是两个复形的微分，右上块为 $f$，左下块为零，且右下块带负号。把 $X$ 换成 $\hom_{\A}(X,Y_0)$、$f$ 换成与 $c$ 复合，即得\eref{3.6}，其中 $c$ 前的符号也源于这一约定。[来源：nLab "mapping cone"，https:/\allowbreak /\allowbreak ncatlab.org/\allowbreak nlab/\allowbreak show/\allowbreak mapping+\allowbreak cone（2026-11-15 访问）]}

考察在一个固定上同调类内改变上闭链 $c$ 的选择时，$\ACone(c)$ 如何变化，是很有启发性的。若另取代表 $\tilde c=c+\mu^1_{\A}(h)$，则
\begin{equation}\label{eq:3.7}
\begin{aligned}
t^1(b_0,b_1)&=\bigl((-1)^{|b_0|-1}b_0,\,(-1)^{|b_1|}b_1+\mu^2_{\A}(h,b_0)\bigr),\\
t^d((b_0,b_1),a_{d-1},\ldots,a_1)&=\bigl(0,\mu^{d+1}_{\A}(h,b_0,a_{d-1},\ldots,a_1)\bigr)\quad(d\geq2)
\end{aligned}
\end{equation}
定义了模同态 $\ACone(c)\to\ACone(\tilde c)$。由\stref{引理}{1.16}，它在 $H^0(\Q)$ 中是同构（这里及下文类似情形中，$|b_0|$ 指 $b_0$ 在 $\hom_{\A}(X,Y_0)$ 中的次数，即施加平移 $[1]$ 之前的次数）。因此，上同调类 $[c]$ 在该范畴中的非典范同构意义下确定 $\ACone(c)$；之所以非典范，是因为\eref{3.7}需要选择 $h$。由定义，对作为 $H^0(\A)$ 对象的 $\Cone(c)$ 也成立同样的结论。

\begin{note}\label{st:3.5}
直和、平移与映射锥之间有一个初等关系：平凡态射 $c=0\in\hom_{\A}(Y_0,Y_1)$ 的锥，拟表示 $\ModShift\Y_0\oplus\Y_1$。例如，假设 $\A$ 对映射锥封闭，也就是说，对所有 $c$，$\Cone(c)$ 都存在。先取一个辅助对象 $Z$，以及代表恒等态射的链 $e_Z$，则 $\ACone(e_Z)$ 同构于零模。于是，对任意 $Y$，零映射 $Y\to\Cone(e_Z)$ 的锥其实就是 $SY$，这意味着 $\A$ 对平移 $S^{\sigma}$（$\sigma>0$）封闭（显然也可推出，对所有 $Y_0,Y_1$，$SY_0\oplus SY_1$ 都存在）。
\end{note}

\subsection{正合三角}\label{sec:3f}
任意抽象映射锥 $\C=\ACone(c)$ 都带有典范模同态 $\iota\in\hom^0_{\Q}(\Y_1,\C)$、$\pi\in\hom^1_{\Q}(\C,\Y_0)$；用\eref{3.6}中的记号，
\SourcePage{44}
它们由 $\iota^1(b_1)=(0,(-1)^{|b_1|}b_1)$、$\pi^1(b_0,b_1)=(-1)^{|b_0|-1}b_0$ 给出，高阶项为零。这些态射在 $H(\Q)$ 中组成三角图：
\SourceEquationDiagram{diagram-3-8}{0.21807669}{3.8}
箭头上的标记 $[1]$ 提醒我们，相应态射的次数为一。为便于后文引用，我们还要给出这些映射的另一些代表：
\begin{lemma}\label{st:3.6}
设 $e_{Y_0},e_{Y_1}$ 是代表 $H(\A)$ 中恒等态射的链。那么由
\[
\begin{aligned}
\tilde\iota^d(b,a_{d-1},\ldots,a_1)&=\bigl(0,\mu^{d+1}_{\A}(e_{Y_1},b,a_{d-1},\ldots,a_1)\bigr),\\
\tilde\pi^d((b_0,b_1),a_{d-1},\ldots,a_1)&=-\mu^{d+1}_{\A}(e_{Y_0},b_0,a_{d-1},\ldots,a_1)
\end{aligned}
\]
给出的模同态 $\tilde\iota\in\hom^0_{\Q}(\Y_1,\C)$、$\tilde\pi\in\hom^1_{\Q}(\C,\Y_0)$，分别与 $\iota$、$\pi$ 上同调。
\end{lemma}
\begin{proof}
已知 Yoneda 嵌入上同调含幺。因此存在 $h\in\hom^{-1}_{\Q}(\Y_1,\Y_1)$，使 $\mu^1_{\Q}(h)=e_{\Y_1}-\mathcal I^1(e_{Y_1})$，即 $\Q$ 中的严格恒等态射与来自 $\A$ 中恒等态射代表的链之差。把它作为第二个分量，第一个分量置零，就得到 $k=(0,h)\in\hom^{-1}_{\Q}(\Y_1,\C)$，显然满足 $\mu^1_{\Q}(k)=\iota-\tilde\iota$。对 $\pi-\tilde\pi$ 的论证相同。
\end{proof}

\begin{definition*}
把 $H(\A)$ 中任意如下形式的图
\SourceEquationDiagram{diagram-3-9}{0.21807669}{3.9}
称为\emph{正合三角}，如果它在 Yoneda 嵌入下变为与\eref{3.8}同构的图。
\end{definition*}
这意味着：令 $c=c_1$，则在 $H^0(\Q)$ 中存在同构 $[t]:\Y_2\to\C=\ACone(c)$，满足 $[\pi]\circ[t]=[\mathcal I^1(c_3)]$ 以及 $[\iota]=[t]\circ[\mathcal I^1(c_2)]$。根据前面关于 $\ACone(c)$ 如何依赖链 $c$ 的选择的讨论，这里的选择并不重要：若对某个 $c$ 的选择能构造同构 $[t]$，那么对任何其他选择也都能做到。沿用\stref{引理}{3.1}的一般思路，可以得到更具体的正合性判据：

\SourcePage{45}
\begin{lemma}\label{st:3.7}
三角\eref{3.9}正合，当且仅当如下条件成立。存在
\[
h_1\in\hom^0_{\A}(Y_1,Y_0),\qquad h_2\in\hom^0_{\A}(Y_2,Y_1),\qquad k\in\hom^{-1}_{\A}(Y_1,Y_1),
\]
满足
\begin{equation}\label{eq:3.10}
\begin{aligned}
\mu^1_{\A}(h_1)&=\mu^2_{\A}(c_3,c_2),\\
\mu^1_{\A}(h_2)&=-\mu^2_{\A}(c_1,c_3),\\
\mu^1_{\A}(k)&=-\mu^2_{\A}(c_1,h_1)+\mu^2_{\A}(h_2,c_2)+\mu^3_{\A}(c_1,c_3,c_2)-e_{Y_1},
\end{aligned}
\end{equation}
其中 $e_{Y_1}$ 是恒等态射的链代表；此外，对每个对象 $X$，如下链复形必须无上同调：
\begin{equation}\label{eq:3.11}
\begin{gathered}
\hom_{\A}(X,Y_2)[1]\oplus\hom_{\A}(X,Y_0)[1]\oplus\hom_{\A}(X,Y_1),\\
\partial=\begin{pmatrix}
\mu^1_{\A}&0&0\\
\mu^2_{\A}(c_3,\cdot)&\mu^1_{\A}&0\\
\mu^2_{\A}(h_2,\cdot)+\mu^3_{\A}(c_1,c_3,\cdot)&\mu^2_{\A}(c_1,\cdot)&\mu^1_{\A}
\end{pmatrix}.
\end{gathered}
\end{equation}
\end{lemma}
\begin{proof}
假设上述所有数据都存在。令 $\C=\ACone(c_1)$，则 $b=(c_3,h_2)$ 是 $\C(Y_2)$ 中的上闭链，从而给出模同态 $t=\lambda_{\C}(b)\in\hom^0_{\Q}(\Y_2,\C)$。复形\eref{3.11}就是相关链映射\eref{3.1}的映射锥，因此其无上同调性保证 $[t]$ 在 $H^0(\Q)$ 中是同构。由\eref{1.26}及 $\C$ 的定义可知，$[\mu^2_{\Q}(t,\mathcal I^1(c_2))]=[\lambda_{\C}(\mu^2_{\C}(b,c_2))]$ 是
\[
\bigl(\mu^2_{\A}(c_3,c_2),\,\mu^2_{\A}(h_2,c_2)+\mu^3_{\A}(c_1,c_3,c_2)\bigr)\in\C(Y_1)
\]
在 $\lambda_{\C}$ 下的像。由假设\eref{3.10}，它与 $(0,e_{Y_1})$ 上同调，而后者在 $\lambda_{\C}$ 下的像，就是\stref{引理}{3.6}中的模同态 $\tilde\iota$。因此
\[
[\iota]=[\mu^2_{\Q}(t,\mathcal I^1(c_2))].
\]
接着，由\eref{1.27}可见，在 $\Hom^1_{H(\Q)}(\Y_2,\Y_0)$ 中，
\[
[\mu^2_{\Q}(\pi,t)]=[\lambda_{\Y_0}(\pi^1(b))]=[\lambda_{\Y_0}(c_3)]=[\mathcal I^1(c_3)].
\]
这样就验证了正合三角所需的全部条件。\znote{37}{引理 3.7 的条件正对应经典三角范畴的公理：\eref{3.10}的前两式记录“三角中相继复合为零”，而\eref{3.11}无上同调则等价于“用 Hom 作用后得到长正合列”。后者是三角结构最核心的性质，在通常的公理体系中单列为一条引理。[来源：nLab "triangulated category"，https:/\allowbreak /\allowbreak ncatlab.org/\allowbreak nlab/\allowbreak show/\allowbreak triangulated+\allowbreak category（2026-11-15 访问）]}

反方向的证明，实质上就是把上述论证倒过来。任意态射 $[t]:\Y_2\to\C$ 都经由 $\lambda_{\C}$ 来自某个元素 $b\in\C(Y_2)$。由于 $[\mathcal I^1(\pi^1(b))]=[\mu^2_{\Q}(\pi,t)]=[\mathcal I^1(c_3)]$，$b$ 的第一个分量与 $c_3$ 上同调。加上一个适当的上边界后，可以假设 $b=(c_3,h_2)$，其中 $h_2$ 为某个元素。接着，
\[
\begin{aligned}
{}[\lambda_{\C}(0,e_{Y_1})]&=[\tilde\iota]=[\mu^2_{\Q}(t,\mathcal I^1(c_2))]\\
&=[\lambda_{\C}(\mu^2_{\C}(b,c_2))].
\end{aligned}
\]
因为 $\lambda_{\C}$ 是拟同构，$(0,e_{Y_1})$ 与 $\bigl(\mu^2_{\A}(c_3,c_2),\mu^2_{\A}(h_2,c_2)+\mu^3_{\A}(c_1,c_3,c_2)\bigr)$ 上同调。选择一个余链，其微分为这两个上闭链之差，即得到剩余数据 $(k,h_1)$。
\end{proof}

\SourcePage{46}
\subsection{另一种刻画}\label{sec:3g}
虽然拟表示对象这一概念为各种构造提供了统一的方法，但用它推导具体性质时，可能会显得繁琐。因此，接下来关于正合三角的讨论暂时不给出证明。稍后改用更具体的观点时，我们会补上这些证明。

考虑对象为 $Z_0,Z_1,Z_2$ 的严格含幺 $\Ai$ 范畴 $\mathcal D$，其中
\[
\begin{aligned}
\hom_{\mathcal D}(Z_k,Z_k)&=\K\cdot e_{Z_k},\\
\hom_{\mathcal D}(Z_0,Z_1)&=\K\cdot x_1,\quad |x_1|=0,\\
\hom_{\mathcal D}(Z_1,Z_2)&=\K\cdot x_2,\quad |x_2|=0,\\
\hom_{\mathcal D}(Z_2,Z_0)&=\K\cdot x_3,\quad |x_3|=1,\\
\hom_{\mathcal D}(Z_2,Z_1)&=\hom_{\mathcal D}(Z_1,Z_0)=\hom_{\mathcal D}(Z_0,Z_2)=0;
\end{aligned}
\]
除严格单位元所要求的运算外，仅有的非零复合映射为
\[
\mu^3_{\mathcal D}(x_3,x_2,x_1)=e_{Z_0},\quad
\mu^3_{\mathcal D}(x_1,x_3,x_2)=e_{Z_1},\quad
\mu^3_{\mathcal D}(x_2,x_1,x_3)=e_{Z_2}.
\]
\begin{proposition}\label{st:3.8}
$H(\A)$ 中的三角\eref{3.9}正合，当且仅当存在 $\Ai$ 函子 $\F:\mathcal D\to\A$，使 $\F(Z_k)=Y_k$，并在上同调上满足 $[\F^1(x_k)]=[c_k]$。
\end{proposition}
将它与\stref{引理}{3.7}比较很有启发性。假设存在如上所述的 $\Ai$ 函子 $\F$。选择 $d_k$，使 $\mu^1_{\A}(d_k)=\F^1(x_k)-c_k$。那么
\begin{equation}\label{eq:3.12}
\begin{aligned}
h_1&=-\F^2(x_3,x_2)+\mu^2_{\A}(\F^1(x_3),d_2)-\mu^2_{\A}(d_3,c_2),\\
h_2&=\F^2(x_1,x_3)-\mu^2_{\A}(\F^1(x_1),d_3)-\mu^2_{\A}(d_1,c_3),\\
k&=-\mu^3_{\A}(d_1,c_3,c_2)-\mu^3_{\A}(\F^1(x_1),d_3,c_2)+\mu^3_{\A}(\F^1(x_1),\F^1(x_3),d_2)\\
&\quad-\mu^2_{\A}(\F^2(x_1,x_3),d_2)\\
&\quad+\mu^2_{\A}\bigl(d_1,\mu^2_{\A}(\F^1(x_3),d_2)-\mu^2_{\A}(d_3,c_2)-\F^2(x_3,x_2)\bigr)
\end{aligned}
\end{equation}
满足\eref{3.10}中的全部等式，其中 $e_{Y_1}=\F^1(e_{Z_1})$。$\F$ 的其他分量可用来构造复形\eref{3.11}的收缩同伦（这里不把它写出），这给出了证明\stref{命题}{3.8}中一个方向的一种方法。稍后给出的证明在原则上与之相似，但会有更好的理论基础，从而减少对\eref{3.12}这类显式公式的依赖。
\begin{corollary}\label{st:3.9}
若 $\F:\A\to\B$ 为任意 $\Ai$ 函子，则 $H(\F)$ 把 $H(\A)$ 中的正合三角映为 $H(\B)$ 中的正合三角。
\end{corollary}
\begin{corollary}\label{st:3.10}
设 $\F:\A\to\B$ 为上同调全忠实 $\Ai$ 函子。假设 $H(\B)$ 中有一个正合三角，组成它的对象都落在 $\F$ 的像中。那么它的原像是 $H(\A)$ 中的正合三角。
\end{corollary}

\begin{note}\label{st:3.11}
\stref{命题}{3.8}出自 Kontsevich。我们的 $\mathcal D=\mathcal D_3$ 实际上属于一族 $\Ai$ 范畴 $\mathcal D_n$（$n\geq1$），后者有 $n$ 个对象 $Z_0,\ldots,Z_{n-1}$。
\SourcePage{47}
存在满足 $\F(Z_0)=Y_0$ 的 $\Ai$ 函子 $\F:\mathcal D_1\to\A$，当且仅当 $Y_0$ 是 $H(\A)$ 中的零对象。类似地，存在满足 $\F(Z_k)=Y_k$ 的 $\F:\mathcal D_2\to\A$，当且仅当 $Y_0,Y_1$ 在 $H^0(\A)$ 中同构（有关讨论参看\bcite{95}，第 193 页）。因此，正合三角可视为把同构概念从两个对象推广到三个对象。
\end{note}

\begin{note}\label{st:3.12}
还有一个值得注意的推论。取 $\mathcal D$ 中由对象 $Z_k$ 组成的基本正合三角。按\stref{笔记}{1.2}定义的任意三个相继映射的 Massey 积，除一个符号外就是恒等态射：
\begin{equation}\label{eq:3.13}
\begin{aligned}
\langle[x_3],[x_2],[x_1]\rangle&=[e_{Z_0}],\\
\langle[x_1],[x_3],[x_2]\rangle&=[-e_{Z_1}],\\
\langle[x_2],[x_1],[x_3]\rangle&=[e_{Z_2}].
\end{aligned}
\end{equation}
由 Massey 积对于 $\Ai$ 函子的函子性，任意正合三角都满足同样的结论。例如，取一个正合三角，把其中三个态射之一乘以某个 $\K\setminus\{0,1\}$ 中的常数。一般而言，所得三角将不满足\eref{3.13}，因而不再正合。这里说“一般而言”，是因为 Massey 积取值的商群\eref{1.4}可能为零；在这种情形下，修改后的三角仍可能正合。这在经典的三角范畴框架中也成立，例如参看\bcite{20}，推论 1.1.10(ii)。
\end{note}

\subsection{三角化的 \texorpdfstring{$A_\infty$}{A∞} 范畴}\label{sec:3h}
\begin{definition*}
若一个 $\Ai$ 范畴非空（$\Ob\A\neq\varnothing$），且满足下列条件，则称它为\emph{三角化的}：
\begin{itemize}
\item $H^0(\A)$ 中每个态射 $[c_1]$ 都能补全为正合三角\eref{3.9}；
\item 对每个对象 $Y$，存在 $\widetilde Y$，使 $S\widetilde Y\cong Y$ 在 $H^0(\A)$ 中成立。
\end{itemize}

\end{definition*}
这里需要略作解释。由\stref{笔记}{3.5}，第一个条件意味着每个对象 $Y$ 都有平移 $SY$，因此第二个条件才有意义。用更具函子性的语言来说，第一个条件允许构造平移函子 $S:\A\to\A$（参看\secref{3d}），第二个条件则断言它是拟等价。\znote{39}{三角化的两条要求恰好是经典公理 TR2 与平移可逆性的范畴化：第一条保证任意态射都能补全为殊三角，第二条保证平移在两个方向上都有定义。对 $H^0(\A)$ 而言平移是自等价，所以第二条自动成立，真正的内容全在第一条。[来源：nLab "triangulated category"，https:/\allowbreak /\allowbreak ncatlab.org/\allowbreak nlab/\allowbreak show/\allowbreak triangulated+\allowbreak category（2026-11-15 访问）]}
\begin{example}\label{st:3.13}
$\Ai$ 范畴 $\Q=\modcat(\A)$ 是三角化的。
\end{example}
\begin{proposition}\label{st:3.14}
若 $\A$ 是三角化的 $\Ai$ 范畴，则 $H^0(\A)$ 是经典意义下的三角范畴。此外，若 $\A,\B$ 是三角化的 $\Ai$ 范畴，$\F:\A\to\B$ 是 $\Ai$ 函子，则 $H^0(\F)$ 是正合函子。
\end{proposition}
回顾经典意义下的三角范畴\bcite{141}、\bcite{61}：它是一个加性范畴，配以一个“平移”自等价 $T$ 以及一类“殊三角”，后者是如下形式的周期图：
\[
\cdots\longrightarrow Y_0\longrightarrow Y_1\longrightarrow Y_2\longrightarrow TY_0\longrightarrow TY_1\longrightarrow TY_2\longrightarrow\cdots.
\]
\SourcePage{48}
在我们的情形中，这个平移由平移函子诱导，即 $T=H^0(S)$；殊三角就是正合三角，其中借助\eref{3.3}把 $[c_3]$ 转化为态射 $Y_2\to SY_0$。\znote{40}{经典定义要求加性范畴、平移自等价与一族满足 (TR1)–(TR4) 的三角，其中 (TR4) 即八面体公理。本书取 $T=H^0(S)$，并把正合三角规定为殊三角，于是 $\Ai$ 层面的正合性与经典三角结构被衔接起来，后面各条推论都可以回到这套公理去检验。[来源：nLab "triangulated category"，https:/\allowbreak /\allowbreak ncatlab.org/\allowbreak nlab/\allowbreak show/\allowbreak triangulated+\allowbreak category（2026-11-15 访问）]}

\clearpage
\subsection{三角的性质}\label{sec:3i}
\stref{命题}{3.14}表明，三角化 $\Ai$ 范畴中的正合三角类满足三角结构的公理性质。事实上，这些性质经过适当表述后，也适用于一般的 $\Ai$ 范畴。为了方便不太熟悉同调代数的读者，下面列出其中一部分：
\begin{corollary}\label{st:3.15}
设\eref{3.9}是 $\Ai$ 范畴 $\A$ 中的正合三角。那么其中任意两个相继态射的复合为零；而且，固定对象 $X$，从左侧或右侧取 Hom，可得到向量空间的长正合序列
\SourceDiagram{diagram-3-15}{0.61061473}
\end{corollary}
反过来不成立，即上述推论中的性质不足以刻画正合三角（利用\stref{笔记}{3.12}，很容易构造反例）。\znote{40}{由殊三角取 Hom 得到的两条长正合列，是三角结构最主要的推论：一条关于协变 Hom，一条关于反变 Hom，与拓扑中纤维序列给出的同伦群长正合列完全平行。反过来，仅凭长正合列并不能判定一个三角是殊三角，这正是本段强调反向不成立的意思。[来源：nLab "triangulated category"，https:/\allowbreak /\allowbreak ncatlab.org/\allowbreak nlab/\allowbreak show/\allowbreak triangulated+\allowbreak category（2026-11-15 访问）]}
\begin{corollary}\label{st:3.16}
取两个正合三角
\SourceEquationDiagram{diagram-3-14}{0.38007652}{3.14}
以及映射 $[b_0]\in\Hom^0_{H(\A)}(Y_0,\widetilde Y_0)$、$[b_1]\in\Hom^0_{H(\A)}(Y_1,\widetilde Y_1)$，使其满足
\[
[\mu^2_{\A}(\tilde c_1,b_0)]=[\mu^2_{\A}(b_1,c_1)].
\]
那么存在一个（不唯一的）态射 $[b_2]$，将图补全为\eref{3.14}所示的交换图。此外，若 $[b_0],[b_1]$ 是同构，那么 $[b_2]$ 也是同构。
\end{corollary}

\SourcePage{49}
\begin{corollary}\label{st:3.17}
假设已有正合三角\eref{3.9}，且 $Y_0$ 存在平移 $SY_0$。那么旋转后的三角
\SourceDiagram{diagram-3-17}{0.23053822}
仍然正合。
\end{corollary}
在上一推论的陈述中，我们隐含地使用了\eref{3.3}，把 $[c_3]$ 视为到 $SY_0$ 的态射。对 $[c_1]$，类似地使用\eref{3.5}，其中非平凡的部分是应用平移函子 $S$，故采用记号 $[S(c_1)]$。

\begin{note}\label{st:3.18}
刚才的推论有一个简单而有用的结论。假设正合三角\eref{3.9}中的映射 $[c_1]$ 为零。由于抽象映射锥的同构类不依赖于上闭链的选择，\stref{笔记}{3.5}中的结论仍成立，即在 $H^0(\A)$ 中有 $Y_2\cong SY_0\oplus Y_1$。通过旋转正合三角，当其他任意一个映射为零时，也能得到类似结论。完整列表为
\[
\begin{aligned}
{}[c_1]=0&\quad\Longrightarrow\quad Y_2\cong SY_0\oplus Y_1,\\
[c_2]=0&\quad\Longrightarrow\quad Y_0\cong Y_1\oplus S^{-1}Y_2,\\
[c_3]=0&\quad\Longrightarrow\quad Y_1\cong Y_2\oplus Y_0.
\end{aligned}
\]
\end{note}
\begin{corollary}\label{st:3.19}
给定正合三角
\SourceDiagram{diagram-3-19a}{0.62515318}
还存在正合三角
\SourceDiagram{diagram-3-19b}{0.13084601}
\end{corollary}
习惯上把三角范畴的公理列为 (TR1)–(TR4)\bcite{61}，第 5 章。\stref{推论}{3.17}对应于 (TR2)；\stref{推论}{3.15}与\stref{推论}{3.16}是 (TR1)–(TR3) 的推论；\stref{推论}{3.19}则是八面体公理 (TR4) 的一个不完整版本。

\SourcePage{50}
\subsection{生成元}\label{sec:3j}
设 $\B$ 是三角化的 $\Ai$ 范畴，$\A\subset\B$ 是非空的全 $\Ai$ 子范畴。
\begin{definition*}
考虑满足如下条件的最小全子范畴 $\widetilde{\B}\subset\B$：它包含 $\A$，对同构封闭（即若 $X_0\cong X_1$ 在 $H^0(\B)$ 中成立，且 $X_0\in\Ob\widetilde{\B}$，则 $X_1\in\Ob\widetilde{\B}$），且自身三角化。称它为 $\B$ 中\emph{由 $\A$ 生成的三角化子范畴}；等价地，也可以说由 $\Ob\A\subset\Ob\B$ 生成。若 $\widetilde{\B}=\B$，则称 $\A$ \emph{生成} $\B$，或称其对象为 $\B$ 的\emph{生成元}。
\end{definition*}
具体地，从 $\Ob\A$ 出发，构造所有可能的映射锥及平移（向上与向下），再把这个过程重复任意有限多次，就得到 $\Ob\widetilde{\B}$。这样得到的子范畴是三角化的，因为作为\stref{推论}{3.9}与\stref{推论}{3.10}的特殊情形，一个三角在 $H(\widetilde{\B})$ 中正合，当且仅当它在 $H(\B)$ 中正合。

\begin{definition*}
设 $\A$ 是非空 $\Ai$ 范畴。一个\emph{三角化包络}是二元组 $(\B,\F)$，其中 $\B$ 为三角化的 $\Ai$ 范畴，$\F:\A\to\B$ 为上同调全忠实函子，且 $\F$ 的像中的对象生成 $\B$。
\end{definition*}
\begin{lemma}\label{st:3.20}
三角化包络总是存在的。此外，若 $(\B,\F)$ 与 $(\widetilde{\B},\widetilde{\F})$ 是 $\A$ 的两个这样的包络，则存在拟等价 $\G:\B\to\widetilde{\B}$，使 $\G\circ\F$ 在 $H^0(\operatorname{fun}(\A,\widetilde{\B}))$ 中同构于 $\widetilde{\F}$。
\end{lemma}
特别地，三角范畴 $H^0(\B)$ 在等价意义下不依赖于包络的选择；事实上，由\stref{命题}{3.14}的后一部分，这是正合等价。
\begin{definition*}
称这个范畴为 $\A$ 的\emph{导出范畴}，记作 $D(\A)$。
\end{definition*}
这个定义的优点，是允许采用不同的三角化包络给出不同实现。不过，对大多数应用而言，这种一般性带来的妨碍多于帮助。因此，下面不再保持中立，而选定一种具体的包络，即扭曲复形的 $\Ai$ 范畴 $\Tw\A$。它的定义仿照加性范畴上的链复形所组成的 dg 范畴。在这个范畴中，映射锥运算有显式实现，因而容易研究；我们将以此为出发点，证明上述所有抽象性质。
\begin{note}\label{st:3.21}
至少还有一种典范的三角化包络构造，它基于\stref{例题}{3.13}：在 $\Q=\modcat(\A)$ 中，取由 Yoneda 嵌入的像生成的全子范畴即可。
\end{note}
\begin{note}\label{st:3.22}
这里对三角结构采用的“链层次”或“内在”观点，最早出现于\bcite{28}。该文还引入了扭曲复形，至少讨论了 dg 范畴的情形；推广到 $\Ai$ 范畴的工作则见于\bcite{90}。
\end{note}

\subsection{加性扩张}\label{sec:3k}
构造 $\Tw\A$ 的第一步是加性扩张 $\Sigma\A$。设 $\A$ 是非含幺 $\Ai$ 范畴。
\begin{definition*}
$\Sigma\A$ 的对象是三元组
\SourcePage{51}
$X=(I,\{X^i\},\{V^i\})$，更直观地写作形式直和
\begin{equation}\label{eq:3.15}
X=\bigoplus_{i\in I}V^i\otimes X^i,
\end{equation}
其中 $I$ 是有限集，$\{X^i\}_{i\in I}$ 是 $\A$ 的一族对象，$\{V^i\}_{i\in I}$ 是相应的一族有限维分次向量空间，称为\emph{重数空间}。两个这样的直和之间的态射，由其组成直和项之间的态射适当地张量相应向量空间之间的线性映射构成：
\begin{equation}\label{eq:3.16}
\begin{aligned}
\hom_{\Sigma\A}\left(\bigoplus_{i\in I_0}V_0^i\otimes X_0^i,\,\bigoplus_{j\in I_1}V_1^j\otimes X_1^j\right)
&=\bigoplus_{ij}\hom_{\K}(V_0^i,V_1^j)\otimes\hom_{\A}(X_0^i,X_1^j),
\end{aligned}
\end{equation}
并带有自然的分次。
\end{definition*}
把态射写作矩阵 $a=(a^{ji})$，行列指标分别为 $j\in I_1$、$i\in I_0$；每个矩阵元形如
\[
a^{ji}=\sum_k\phi^{jik}\otimes x^{jik}\in\hom_{\K}(V_0^i,V_1^j)\otimes\hom_{\A}(X_0^i,X_1^j).
\]
写公式时，常只考虑求和仅含一项的情形 $a^{ji}=\phi^{ji}\otimes x^{ji}$，并默认按线性扩张。有了这一约定，即可定义 $\Sigma\A$ 上的复合映射。它把 $\A$ 的复合与线性映射的通常复合结合起来，正如\stref{笔记}{1.11}中的张量积：
\begin{equation}\label{eq:3.17}
\begin{aligned}
\mu^d_{\Sigma\A}(a_d,\ldots,a_1)^{i_d,i_0}
&=\sum_{i_1,\ldots,i_{d-1}}(-1)^{\triangleleft}\phi_d^{i_d,i_{d-1}}\cdots\phi_1^{i_1,i_0}\\[-2pt]
&\qquad\otimes\mu^d_{\A}(x_d^{i_d,i_{d-1}},\ldots,x_1^{i_1,i_0}),
\end{aligned}
\end{equation}
其中 $\triangleleft=\sum_{p<q}|\phi_p^{i_p,i_{p-1}}|\cdot(|x_q^{i_q,i_{q-1}}|-1)$，$a_k=(a_k^{ji})=(\phi_k^{ji}\otimes x_k^{ji})$。从现在起固定一个单元素集 $\{*\}$，把 $\A$ 的对象等同于三元组 $(\{*\},\{X^*=X\},\{V^*=\K\})$，其中 $\K$ 集中在零次。这样就把 $\A$ 嵌入为 $\Sigma\A$ 的全 $\Ai$ 子范畴。

\subsection{扭曲复形}\label{sec:3l}
\begin{definition*}
$\A$ 中的一个\emph{预扭曲复形}是二元组 $(X,\delta_X)$，其中 $X\in\Ob\Sigma\A$，而微分（有时也称联络）为 $\delta_X\in\hom^1_{\Sigma\A}(X,X)$。展开来写，就是
\begin{equation}\label{eq:3.18}
\left(X=\bigoplus_{i\in I}V^i\otimes X^i,\quad\delta_X=(\delta_X^{ji})\right),
\end{equation}
其中 $\delta_X^{ji}=\sum_k\phi^{jik}\otimes x^{jik}$，$\phi^{jik}\in\Hom_{\K}(V^i,V^j)$、$x^{jik}\in\hom_{\A}(X^i,X^j)$，且 $|\phi^{jik}|+|x^{jik}|=1$。
\end{definition*}
出于技术上的原因，还需要引入（简单的）子复形 $\widetilde X\subset X$ 与商复形 $X/\widetilde X$ 的概念。
\begin{definition*}
子复形的定义是：对每个 $i\in I$，取分次子空间 $\widetilde V^i\subset V^i$，使所有 $\phi^{jik}$ 都把 $\widetilde V^i$ 映入 $\widetilde V^j$，然后取 $\bigoplus_i\widetilde V^i\otimes X^i$，配以相应限制的微分。
\SourcePage{52}
商复形以相同方式构造，使用向量空间 $V^i/\widetilde V^i$ 以及 $\phi^{jik}$ 在它们之间诱导的线性映射。
\end{definition*}
\begin{definition*}
\emph{扭曲复形}是具有如下两个性质的预扭曲复形 $(X,\delta_X)$。首先，$\delta_X$ \emph{严格下三角}，即存在有限递减的子复形滤过
\[
X=F^0X\supset F^1X\supset\cdots\supset F^nX=0,
\]
使商 $F^kX/F^{k+1}X$ 上的诱导微分为零；滤过并不是 $(X,\delta_X)$ 结构的一部分，只要求它存在。其次，满足\emph{广义 Maurer–Cartan 方程}
\begin{equation}\label{eq:3.19}
\sum_{r=1}^{\infty}\mu^r_{\Sigma\A}(\delta_X,\ldots,\delta_X)=0.
\end{equation}

\end{definition*}
这个式子有意义，因为结合下三角条件与定义\eref{3.17}可见，\eref{3.19}中除有限多项外，其余各项均为零。常常省略记号中的微分，直接写作 $X\in\Ob\Tw\A$。

扭曲复形组成非含幺 $\Ai$ 范畴 $\Tw\A$ 的对象。态射空间与之前相同，即 $\hom_{\Tw\A}(X_0,X_1)=\hom_{\Sigma\A}(X_0,X_1)$，但所有复合都由微分的贡献作了变形：
\begin{equation}\label{eq:3.20}
\begin{aligned}
\mu^d_{\Tw\A}(a_d,\ldots,a_1)
&=\sum_{i_0,\ldots,i_d}\mu^{d+i_0+\cdots+i_d}_{\Sigma\A}\bigl(\overbrace{\delta_{X_d},\ldots,\delta_{X_d}}^{i_d},a_d,\\[-2pt]
&\qquad\underbrace{\delta_{X_{d-1}},\ldots,\delta_{X_{d-1}}}_{i_{d-1}},a_{d-1},\ldots,a_1,\underbrace{\delta_{X_0},\ldots,\delta_{X_0}}_{i_0}\bigr),
\end{aligned}
\end{equation}
其中对所有 $i_0,\ldots,i_d\geq0$ 求和。同样，微分的下三角性保证，插入足够多个 $\delta_{X_k}$ 后，复合就为零，因此\eref{3.20}是有限和。$\Tw\A$ 中结合律等式\eref{1.2}的成立，恰好归结为广义 Maurer–Cartan 方程\eref{3.19}。还要注意，微分为零的扭曲复形所组成的全子范畴可以与 $\Sigma\A$ 等同。

设 $X_0,X_1$ 是配有上述滤过的扭曲复形。这些滤过诱导链复形 $\hom_{\Tw\A}(X_0,X_1)$ 的滤过，从而给出收敛到其上同调的谱序列，起始项为
\begin{equation}\label{eq:3.21}
E_1^{rs}=\bigoplus_{j-i=r}\Hom^{r+s}_{H(\Sigma\A)}(F^iX_0/F^{i+1}X_0,\,F^jX_1/F^{j+1}X_1).
\end{equation}
该谱序列与 $H(\Tw\A)$ 中的复合相容。这意味着，对任意三个扭曲复形 $X_0,X_1,X_2$，映射 $\mu^2_{\Tw\A}$ 保持滤过，从而在相应谱序列上诱导乘积。

\SourcePage{53}
\subsection{函子性}\label{sec:3m}
设 $\G:\A\to\B$ 为非含幺 $\Ai$ 函子。诱导函子 $\Sigma\G:\Sigma\A\to\Sigma\B$ 定义为
\begin{equation}\label{eq:3.22}
\begin{aligned}
(\Sigma\G)\left(\bigoplus_{i\in I}V^i\otimes X^i\right)&=\bigoplus_{i\in I}V^i\otimes\G(X^i),\\
(\Sigma\G)^d(a_d,\ldots,a_1)^{i_d,i_0}
&=\sum_{i_1,\ldots,i_{d-1}}(-1)^{\triangleleft}\phi_d^{i_d,i_{d-1}}\cdots\phi_1^{i_1,i_0}\\[-2pt]
&\qquad\otimes\G^d(x_d^{i_d,i_{d-1}},\ldots,x_1^{i_1,i_0}),
\end{aligned}
\end{equation}
其中 $\triangleleft$ 与 $a_k$ 同\eref{3.17}。此外，任意 $T\in\hom_{\nufun(\A,\B)}(\G_0,\G_1)$ 都诱导 $\Sigma^1T\in\hom_{\nufun(\Sigma\A,\Sigma\B)}(\Sigma\G_0,\Sigma\G_1)$，同样由下式给出：
\begin{equation}\label{eq:3.23}
\begin{aligned}
(\Sigma^1T)^d(a_d,\ldots,a_1)^{i_d,i_0}
&=\sum_{i_1,\ldots,i_{d-1}}(-1)^{\triangleleft}\phi_d^{i_d,i_{d-1}}\cdots\phi_1^{i_1,i_0}\\[-2pt]
&\qquad\otimes T^d(x_d^{i_d,i_{d-1}},\ldots,x_1^{i_1,i_0}).
\end{aligned}
\end{equation}
公式\eref{3.22}、\eref{3.23}连同为零的高阶项，定义一个非含幺 $\Ai$ 函子 $\Sigma:\nufun(\A,\B)\to\nufun(\Sigma\A,\Sigma\B)$。

扭曲复形的情形完全类似。有一个诱导的非含幺 $\Ai$ 函子 $\Tw\G:\Tw\A\to\Tw\B$，
\[
\begin{aligned}
(\Tw\G)(X,\delta_X)&=\left(\Sigma\G(X),\sum_{e\geq1}(\Sigma\G)^e(\delta_X,\ldots,\delta_X)\right),\\[2pt]
(\Tw\G)^d(a_d,\ldots,a_1)
&=\sum_{i_0,\ldots,i_d}(\Sigma\G)^{d+i_0+\cdots+i_d}\bigl(\overbrace{\delta_{X_d},\ldots,\delta_{X_d}}^{i_d},a_d,\\[-2pt]
&\qquad\underbrace{\delta_{X_{d-1}},\ldots,\delta_{X_{d-1}}}_{i_{d-1}},a_{d-1},\ldots,a_1,\underbrace{\delta_{X_0},\ldots,\delta_{X_0}}_{i_0}\bigr).
\end{aligned}
\]
注意，若 $X$ 有子复形滤过 $F^kX$，使微分严格下三角，那么 $\Sigma\G(X)$ 由 $\Sigma\G(F^kX)$ 诱导的滤过也有同样性质，故它确实是扭曲复形。对 $T\in\hom_{\nufun(\A,\B)}(\G_0,\G_1)$，定义 $\Tw^1T\in\hom_{\nufun(\Tw\A,\Tw\B)}(\Tw\G_0,\Tw\G_1)$ 如下：
\begin{equation}\label{eq:3.24}
\begin{aligned}
(\Tw^1T)^d(a_d,\ldots,a_1)
&=\sum_{i_0,\ldots,i_d}(\Sigma^1T)^{d+i_0+\cdots+i_d}\bigl(\overbrace{\delta_{X_d},\ldots,\delta_{X_d}}^{i_d},a_d,\\[-2pt]
&\qquad\underbrace{\delta_{X_{d-1}},\ldots,\delta_{X_{d-1}}}_{i_{d-1}},a_{d-1},\ldots,a_1,\underbrace{\delta_{X_0},\ldots,\delta_{X_0}}_{i_0}\bigr).
\end{aligned}
\end{equation}
与之前相同，这些数据合在一起构成非含幺 $\Ai$ 函子 $\Tw:\nufun(\A,\B)\to\nufun(\Tw\A,\Tw\B)$，并且与左、右复合相容。这意味着，对固定的 $\G:\A\to\B$，有
\begin{equation}\label{eq:3.25}
\begin{aligned}
\LG_{\Tw\G}\circ\Tw&=\Tw\circ\LG_{\G}:\nufun(\C,\A)\longrightarrow\nufun(\Tw\C,\Tw\B),\\
\RG_{\Tw\G}\circ\Tw&=\Tw\circ\RG_{\G}:\nufun(\B,\C)\longrightarrow\nufun(\Tw\A,\Tw\C).
\end{aligned}
\end{equation}

\SourcePage{54}
\begin{lemma}\label{st:3.23}
若 $\G:\A\to\B$ 上同调全忠实，那么 $\Tw\G:\Tw\A\to\Tw\B$ 也上同调全忠实。
\end{lemma}
\begin{proof}
对 $(\Tw\G)^1$ 在谱序列\eref{3.21}上的作用使用比较论证，就把结论归结为 $\Sigma\G$ 的相应结论，而后者显然成立。
\end{proof}

\subsection{单位元与拟等价}\label{sec:3n}
若 $\A$ 是严格含幺 $\Ai$ 范畴，那么 $\Tw\A$ 也严格含幺。具体地，扭曲复形\eref{3.18}的严格恒等态射是“单位矩阵” $(\phi^{ji}\otimes x^{ji})$，其中 $\phi^{ii}\otimes x^{ii}=\id_{V^i}\otimes e_{X^i}$，而所有非对角矩阵元均为零。此外，若 $\G:\A\to\B$ 是严格含幺 $\Ai$ 函子，那么 $\Tw\G$ 也严格含幺。最后，对严格含幺的 $\A$ 及任意 $\C$，
\begin{equation}\label{eq:3.26}
\Tw:\nufun(\C,\A)\longrightarrow\nufun(\Tw\C,\Tw\A)
\end{equation}
是严格含幺 $\Ai$ 函子。只需回顾\secref{2b}中 $E_{\G}$ 的表达式，并代入\eref{3.24}，即可看出这一点。其上同调含幺版本则不那么显然：
\begin{lemma}\label{st:3.24}
\textup{(i)} 若 $\A$ 为 c-含幺，那么 $\Tw\A$ 也为 c-含幺。

\textup{(ii)} 假设 $\A,\B$ 都为 c-含幺，且 $\G:\A\to\B$ 是 c-含幺 $\Ai$ 函子。那么 $\Tw\G$ 也为 c-含幺。

\textup{(iii)} 对 c-含幺的 $\A$ 及任意 $\C$，\eref{3.26}是 c-含幺 $\Ai$ 函子。
\end{lemma}
\begin{proof}
(i) 如\stref{引理}{2.3}的证明，先取最小模型，再将其严格化，得到严格含幺的 $\widetilde{\A}$ 及在对象上为恒等映射的拟同构 $\Phi:\A\to\widetilde{\A}$。由\stref{引理}{3.23}，诱导的非含幺 $\Ai$ 函子 $\Tw\Phi:\Tw\A\to\Tw\widetilde{\A}$ 上同调全忠实。后者严格含幺；其恒等态射可经由这一全忠实函子在上同调上唯一拉回，故前者必然上同调含幺。

(ii) 固定 $X\in\Tw\A$ 及适当的滤过 $F^kX$。在保滤过的自同态子复形上，$(\Tw\Phi)^1$ 仍是拟同构；这由\eref{3.21}的同一个有限滤过比较论证得出。因此可以选择代表 $X$ 的上同调恒等态射的保滤过上闭链 $e_X$，使其在各相继商上诱导上同调恒等态射。由定义，$(\Tw\G)^1(e_X)$ 对 $\Tw\G(F^kX)$ 给出的像滤过是下三角的，且在商 $\Tw\G(F^kX)/\Tw\G(F^{k+1}X)$ 上诱导的自同态是恒等态射的链代表，因为 $\G$ 上同调含幺。利用\eref{3.21}可知，以 $[(\Tw\G)^1(e_X)]$ 作左乘和右乘，在 $\Hom_{H(\Tw\B)}(\Tw\G(X),\Tw\G(X))$ 上都是自同构。因此该元素本身可逆。它还是幂等元，因为它是恒等态射在一个函子下的像，所以它等于恒等态射。这证明 $\Tw\G$ 上同调含幺。

\SourcePage{55}
(iii) 固定非含幺 $\Ai$ 函子 $\G:\C\to\A$，令 $\widetilde{\G}=\Phi\circ\G$，其中 $\Phi$ 如上。作为\eref{3.25}的特殊情形，有交换图
\[
\begin{array}{ccc}
\hom_{\nufun(\C,\A)}(\G,\G)&\xrightarrow{\ \LG^1_{\Phi}\ }&\hom_{\nufun(\C,\widetilde{\A})}(\widetilde{\G},\widetilde{\G})\\[3pt]
\Big\downarrow\scriptstyle\Tw^1&&\Big\downarrow\scriptstyle\Tw^1\\[3pt]
\hom_{\nufun(\Tw\C,\Tw\A)}(\Tw\G,\Tw\G)&\xrightarrow{\ \LG^1_{\Tw\Phi}\ }&\hom_{\nufun(\Tw\C,\Tw\widetilde{\A})}(\Tw\widetilde{\G},\Tw\widetilde{\G}).
\end{array}
\]
由 (ii) 及\stref{引理}{2.4}，水平箭头在上同调层次上都把恒等态射映到恒等态射。右侧竖直箭头处有严格单位元，故它也如此。下方水平箭头由\stref{引理}{3.23}与\stref{引理}{1.7}在上同调上为同构，因此交换性给出所需结论。
\end{proof}

\begin{lemma}\label{st:3.25}
设 $\G:\A\to\B$ 是 c-含幺 $\Ai$ 范畴之间的拟等价。那么 $\Tw\G:\Tw\A\to\Tw\B$ 也是拟等价。
\end{lemma}
这是\stref{定理}{2.9}及上一引理的直接推论：先在相应函子范畴中的同构意义下取 $\G$ 的逆，再考察诱导的函子及自然变换即可。另一种较少依赖形式体系、更多借助正合三角的方法，见后文\stref{引理}{3.36}的证明。

\subsection{基本运算}\label{sec:3o}
从这里起，恢复之前的通常假设：所涉及的一切都上同调含幺。

$\Sigma\A$ 对象的形式允许定义直和运算 $Y_0\oplus Y_1$：形式上，把指标集及其相应的族作不交并即可。这个运算显然延伸到扭曲复形。此外，$Y_0\oplus Y_1$ 在 Yoneda 嵌入下的像，恰好就是\secref{3b}中的 $\Ai$ 模直和 $\Y_0\oplus\Y_1$。

同样自然地，可以对有限维分次向量空间 $Z$ 与 $\Sigma\A$ 中对象 $Y=\bigoplus_{i\in I}W^i\otimes Y^i$ 取张量积 $Z\otimes Y$。它有相同的指标集 $I$ 与对象 $Y^i$，但将每个 $W^i$ 替换为 $Z\otimes W^i$。给定 $a=(a^{ji})=(\phi^{ji}\otimes x^{ji})\in\hom_{\Sigma\A}(Y_0,Y_1)$，把每个线性映射 $\phi^{ji}\in\Hom_{\K}(W_0^i,W_1^j)$ 替换为分次张量积 $\id_Z\otimes\phi^{ji}\in\Hom_{\K}(Z\otimes W_0^i,Z\otimes W_1^j)$，即可定义 $\id_Z\otimes a\in\hom_{\Sigma\A}(Z\otimes Y_0,Z\otimes Y_1)$；这里
\[
(\id_Z\otimes\phi^{ji})(z\otimes w)=(-1)^{|\phi^{ji}|\cdot|z|}z\otimes\phi^{ji}(w).
\]
这给出一个高阶项为零的非含幺 $\Ai$ 函子 $Z\otimes-:\Sigma\A\to\Sigma\A$。扭曲复形也有类似构造，它把 $(Y,\delta_Y)$ 映到 $(Z\otimes Y,\delta_{Z\otimes Y}=\id_Z\otimes\delta_Y)$，并以与之前相同的方式作用于态射。在 Yoneda 嵌入下，它成为\secref{3c}中的 $\modcat(\Tw\A)$ 上的函子 $Z\otimes-$，这里限制在 $Z$ 有限维且 $\partial_Z=0$ 的情形。

\SourcePage{56}
当 $Z=\K[\sigma]$ 时，记 $S^{\sigma} Y=Z\otimes Y$、$S^{\sigma}(a)=\id_Z\otimes a$。把扭曲复形上的运算显式写出，也许会有所帮助：
\[
\begin{gathered}
\left(Y=\bigoplus_iW^i\otimes Y^i,\quad\delta_Y=\left(\sum_k\phi^{jik}\otimes x^{jik}\right)_{ji}\right)\\
\Big\downarrow\\
\left(S^{\sigma} Y=\bigoplus_iW^i[\sigma]\otimes Y^i,\quad
\delta_{S^{\sigma} Y}=S^{\sigma}(\delta_Y)=\left(\sum_k(-1)^{\sigma|\phi^{jik}|}\phi^{jik}\otimes x^{jik}\right)_{ji}\right).
\end{gathered}
\]
将前面的讨论专门用于此情形，即可看出：$S^{\sigma}$ 作为阶数 $\geq2$ 的项均为零的 $\Ai$ 函子，确实表示\secref{3d}中的抽象平移函子 $\ModShift^{\sigma}$。事实上，在链层次上有自然的等同
\[
\hom_{\Tw\A}(Y_0,SY_1)\cong\hom_{\Tw\A}(Y_0,Y_1)[1],
\]
它诱导\eref{3.3}。

作为补充，现在简要考察 $\A$ 严格含幺的情形。此时可以如下定义有限维向量空间链复形 $(Z,\partial_Z)$ 与扭曲复形 $Y=(\bigoplus_{i\in I}W^i\otimes Y^i,\delta_Y)$ 的张量积 $Z\otimes Y$。首先把微分改为 $\tilde\partial_Z(z)=(-1)^{|z|-1}\partial_Z(z)$，然后令
\begin{equation}\label{eq:3.27}
Z\otimes Y=\left(\bigoplus_{i\in I}(Z\otimes W^i)\otimes Y^i,\quad\id_Z\otimes\delta_Y+\tilde\partial_Z\otimes e_Y\right),
\end{equation}
其中微分的第二项定义为对角矩阵，其矩阵元为 $(\tilde\partial_Z\otimes\id_{W^i})\otimes e_{Y^i}$。直接计算可验证广义 Maurer–Cartan 方程；这里关键地用到了 $e_{Y^i}$ 是严格单位元。显然，$\partial_Z$ 相对于 $Z$ 的按次数滤过严格下三角。将这个滤过与 $Y$ 的滤过结合，就得到\eref{3.27}的滤过，从而完成验证：所得对象确实是扭曲复形。

\subsection{映射锥}\label{sec:3p}
设 $c\in\hom_{\Tw\A}(Y_0,Y_1)$ 是零次上闭链。
\begin{definition*}
映射锥 $C=\Cone(c)$ 是扭曲复形
\begin{equation}\label{eq:3.28}
\left(C=SY_0\oplus Y_1,\quad\delta_C=\begin{pmatrix}S(\delta_{Y_0})&0\\-S(c)&\delta_{Y_1}\end{pmatrix}\right).
\end{equation}
\end{definition*}
这里先用平移函子把 $c$ 变为元素 $S(c)\in\hom^0_{\Tw\A}(SY_0,SY_1)$，再以直接的方式把该空间与 $\hom^1_{\Tw\A}(SY_0,Y_1)$ 等同；比较\eref{3.5}。若每个 $Y_k$ 都有子复形滤过 $F^iY_k$，使微分严格下三角，那么映射锥的组合滤过——先取 $F^iSY_0\oplus Y_1$，再接上 $0\oplus F^iY_1$——对 $C$ 也具有同样的性质。广义 Maurer–Cartan 方程等价于各 $Y_k$ 的对应方程，以及 $c$ 是 $\Tw\A$ 中上闭链这一事实。$C$ 在 Yoneda 嵌入下的像可以与抽象映射锥 $\C=\ACone(c)$ 等同，只需利用
\[
\begin{aligned}
\hom_{\Tw\A}(X,C)&=\hom_{\Tw\A}(X,SY_0)\oplus\hom_{\Tw\A}(X,Y_1)\\
&\cong\hom_{\Tw\A}(X,Y_0)[1]\oplus\hom_{\Tw\A}(X,Y_1)=\C(X).
\end{aligned}
\]
\SourcePage{57}
综上，我们已经看到，$\Tw\A$ 含有表示下列运算的优选对象：直和；与有限维分次向量空间的张量积，特别是平移；在严格含幺情形下，还包括与有限维链复形的张量积；以及最重要的映射锥。以后处理扭曲复形时，自然都会采用这些代表。\znote{49}{这些优选代表使扭曲复形范畴成为“链层面可计算”的三角范畴：直和、张量与平移都有具体模型，映射锥则由\eref{3.28}给出。值得注意的是，在一般的三角范畴中映射锥并非函子性构造，本书用扭曲复形把它具体化，正是为了避免这一缺陷。[来源：nLab "triangulated category"，https:/\allowbreak /\allowbreak ncatlab.org/\allowbreak nlab/\allowbreak show/\allowbreak triangulated+\allowbreak category（2026-11-15 访问）]}

\begin{note}\label{st:3.26}
还可以换一种方式定义加性扩张：只允许一维重数空间，更具体地说，只允许 $\K$ 的平移。此时对象 $X$ 写作直和 $\bigoplus_iS^{-\sigma^i}X^i$，其中 $X^i\in\Ob\A$，$\sigma^i\in\Z$。用显然的同构，并对每个系数 $a^{ji}$ 加上额外的符号 $(-1)^{(|a^{ji}|-1)\sigma_0^i}$，来作如下等同：
\begin{equation}\label{eq:3.29}
\hom_{\Sigma\A}\left(\bigoplus_iS^{-\sigma_0^i}X_0^i,\,\bigoplus_jS^{-\sigma_1^j}X_1^j\right)
=\bigoplus_{i,j}\hom_{\A}(X_0^i,X_1^j)[\sigma_1^j-\sigma_0^i].
\end{equation}
这里 $|a^{ji}|$ 指施加任何平移之前，在 $\hom_{\A}(X_0^i,X_1^j)$ 中的次数。相对于这一等同，复合变为
\[
\mu^d_{\Sigma\A}(a_d,\ldots,a_1)^{i_d,i_0}
=(-1)^{\sigma_0^{i_0}}\sum_{i_1,\ldots,i_{d-1}}\mu^d_{\A}(a_d^{i_d,i_{d-1}},\ldots,a_1^{i_1,i_0}),
\]
其中 $a_k\in\hom_{\Sigma\A}(X_{k-1},X_k)$，$\sigma_0^{i_0}$ 是 $X_0$ 的第 $i_0$ 个直和项的平移量。平移函子 $S$ 将态射 $a\in\hom^k_{\Sigma\A}(X_0,X_1)$ 乘以 $(-1)^{k-1}$。若沿同样思路继续定义一个版本的 $\Tw\A$，则 $C=\Cone(c)$ 上的微分为
\[
\delta_C=\begin{pmatrix}\delta_{Y_0}&0\\c&\delta_{Y_1}\end{pmatrix};
\]
\eref{3.28}中的平移运算与符号已经吸收到等同\eref{3.29}中。在拟等价意义下，允许一般的重数空间，还是只允许 $\K$ 的平移，当然并无区别。较一般的方法，其优点是无需选择基就能进行张量积运算 $Z\otimes-$，而且符号也更系统，尽管稍微复杂一些；不过，总体差别很小。
\end{note}


\clearpage
\subsection{\texorpdfstring{$\Tw\A$ 中的正合三角}{TwA 中的正合三角}}\label{sec:3q}
\eref{3.28}中的扭曲复形 $C$ 带有典范态射
\begin{equation}\label{eq:3.30}
\begin{aligned}
\relax[i]&\in\Hom^0_{H(\Tw\A)}(Y_1,C),&i&=\begin{pmatrix}0\\e_{Y_1}\end{pmatrix},\\
[p]&\in\Hom^1_{H(\Tw\A)}(C,Y_0),&p&=(S(e_{Y_0}),0),
\end{aligned}
\end{equation}
\SourcePage{58}
其中 $e_{Y_k}$ 是 $Y_k$ 的恒等态射的链代表（容易看出，具体选择无关紧要）。在 Yoneda 嵌入下，它们成为\stref{引理}{3.6}中的 $\Ai$ 模同态 $\tilde\iota$ 与 $\tilde\pi$。由此以及正合三角和三角化 $\Ai$ 范畴的定义，立即得到：
\begin{lemma}\label{st:3.27}
$H^0(\Tw\A)$ 中的三角\eref{3.9}正合，当且仅当存在 $H^0(\Tw\A)$ 中的同构 $[b]:Y_2\to\Cone(c_1)$，使得
\[
[\mu^2_{\Tw\A}(p,b)]=[c_3],\qquad [i]=[\mu^2_{\Tw\A}(b,c_2)].
\]
\end{lemma}
\begin{lemma}\label{st:3.28}
$\Tw\A$ 是三角化的 $\Ai$ 范畴。
\end{lemma}
构造\eref{3.28}以加性范畴上的链复形范畴中的经典映射锥为原型，在大多数方面也具有相同的性质。特别地，证明其相应同伦范畴为三角范畴的标准论证\bcite{61}（第 IV 章 §2），可以毫无困难地推广到我们的情形，从而得到：
\begin{proposition}\label{st:3.29}
赋予 $H^0(\Tw\A)$ 平移函子 $T=H^0(S)$，并像\stref{命题}{3.14}那样，把正合三角规定为殊三角。这样，$H^0(\Tw\A)$ 就成为经典意义下的三角范畴。
\end{proposition}
由 $\Ai$ 函子 $\F$ 诱导的函子 $\Tw\F:\Tw\A\to\Tw\B$ 与 $S$ 严格交换。此外，可以作等同 $\Tw\F(\Cone(c))\cong\Cone((\Tw\F)^1(c))$，且这一等同与两边相应的态射\eref{3.30}相容。因此：
\begin{lemma}\label{st:3.30}
$H(\Tw\F):H(\Tw\A)\to H(\Tw\B)$ 将正合三角映为正合三角；因而 $H^0(\Tw\F):H^0(\Tw\A)\to H^0(\Tw\B)$ 是经典三角范畴之间的正合函子。
\end{lemma}

\subsection{一般的正合三角}\label{sec:3r}
现在来推导\secref{3g}至\secref{3j}中未经证明而陈述的一般 $\Ai$ 范畴的正合三角性质。主要工具是：
\begin{lemma}\label{st:3.31}
$H(\A)$ 中的三角正合，当且仅当它在嵌入 $H(\A)\hookrightarrow H(\Tw\A)$ 下的像正合。
\end{lemma}
\begin{proof}
取 $\A$ 中的对象 $Y_k$ 和态射 $c_k$，使它们构成一个三角。\stref{引理}{3.7}给出的正合性充要条件，要求存在满足关系\eref{3.10}的某些其他态射，并要求相应的链复形\eref{3.11}无上同调。对于前一个条件，把三角视为位于 $H(\A)$ 还是 $H(\Tw\A)$ 中没有区别，因为嵌入 $\A\hookrightarrow\Tw\A$ 在链层次上就是全忠实的。
\SourcePage{59}
对于后一个条件，若\eref{3.11}对所有 $X\in\Ob\Tw\A$ 都无上同调，那么对所有 $X\in\Ob\A$ 当然也如此。反过来，给定 $X\in\Ob\Tw\A$，可以利用它的子复形递降滤过 $F^iX$，在
\[
\hom_{\Tw\A}(X,Y_2)[1]\oplus\hom_{\Tw\A}(X,Y_0)[1]\oplus\hom_{\Tw\A}(X,Y_1)
\]
上定义一个递增滤过，使之与\eref{3.11}中的微分相容。该滤过的每一个相伴分次部分，都是对 $\A$ 中某些对象取复形\eref{3.11}、再作平移后所得的有限直和。因此，若这些复形均无上同调，与整个 $X$ 相应的复形也无上同调。
\end{proof}
由此很容易得到\stref{推论}{3.15}至\stref{推论}{3.19}所述的正合三角性质：根据\stref{命题}{3.29}，它们在 $H(\Tw\A)$ 中成立，因此在全子范畴 $H(\A)$ 中也成立。

\begin{proof}[\stref{推论}{3.9}的证明]
设 $H(\A)$ 中有正合三角\eref{3.9}，并有 $\Ai$ 函子 $\F:\A\to\B$。由\stref{引理}{3.31}，该三角在 $H(\Tw\A)$ 中也正合。由\stref{引理}{3.30}，它在 $\Tw\F$ 下的像是 $H(\Tw\B)$ 中的正合三角。而 $\Tw\F$ 在 $\A\subset\Tw\A$ 上的限制就是原来的 $\F$；再用一次\stref{引理}{3.31}，即得结论。
\end{proof}

\begin{proof}[\stref{命题}{3.8}的证明（同时也直接证明\stref{推论}{3.10}）]
\stref{引理}{3.7}的一个平凡特例（取 $h_1=h_2=k=0$）表明，$\mathcal D$ 中由对象 $Z_k$ 与态射 $x_k$ 构成的原型三角正合。于是由\stref{推论}{3.9}，对于任意 $\Ai$ 函子 $\F:\mathcal D\to\A$，对象 $Y_k=\F(Z_k)$ 与态射 $[c_k]=[\F^1(x_k)]$ 都构成一个正合三角。

为证明逆向结论，先假设 $\A$ 严格含幺，从而 $\Tw\A$ 也严格含幺。取零次上闭链 $c\in\hom^0_{\A}(Y_0,Y_1)$。考虑非全 $\Ai$ 子范畴 $\widetilde{\mathcal D}\subset\Tw\A$：其对象为 $\widetilde Z_0=Y_0$、$\widetilde Z_1=Y_1$、$\widetilde Z_2=\Cone(c)$，态射空间由 $e_{Y_0},e_{Y_1},c$ 的所有可能出现形式生成：
\begin{align*}
\hom_{\widetilde{\mathcal D}}(\widetilde Z_0,\widetilde Z_0)&=\K e_{Y_0},\\
\hom_{\widetilde{\mathcal D}}(\widetilde Z_1,\widetilde Z_1)&=\K e_{Y_1},\\
\hom_{\widetilde{\mathcal D}}(\widetilde Z_2,\widetilde Z_2)&=\left\{\begin{pmatrix}\lambda_{00}S(e_{Y_0})&0\\\lambda_{10}S(c)&\lambda_{11}e_{Y_1}\end{pmatrix}\right\},\\
\hom_{\widetilde{\mathcal D}}(\widetilde Z_0,\widetilde Z_1)&=\K c,\\
\hom_{\widetilde{\mathcal D}}(\widetilde Z_1,\widetilde Z_2)&=\K\begin{pmatrix}0\\e_{Y_1}\end{pmatrix},\\
\hom_{\widetilde{\mathcal D}}(\widetilde Z_2,\widetilde Z_0)&=\K(S(e_{Y_0}),0),\\
\hom_{\widetilde{\mathcal D}}(\widetilde Z_1,\widetilde Z_0)&=0,\\
\hom_{\widetilde{\mathcal D}}(\widetilde Z_2,\widetilde Z_1)&=\{(\nu_{10}S(c),\nu_{11}e_{Y_1})\},\\
\hom_{\widetilde{\mathcal D}}(\widetilde Z_0,\widetilde Z_2)&=\left\{\begin{pmatrix}\sigma_{00}e_{Y_0}\\\sigma_{10}c\end{pmatrix}\right\}.
\end{align*}
\SourcePage{60}
其中 $\lambda,\nu,\sigma$ 均在 $\K$ 中取值。对于 $d>2$，复合 $\mu^d_{\widetilde{\mathcal D}}$ 为零，因此 $\widetilde{\mathcal D}$ 是 dg 范畴。现在在态射空间 $\hom_{\widetilde{\mathcal D}}(\widetilde Z_i,\widetilde Z_j)$ 上引入自同态 $T^1$。相对于上面给出的基，它们为
\[
\begin{aligned}
T^1(\lambda_{00},\lambda_{10},\lambda_{11})&=(0,0,-\lambda_{10}),\\
T^1(\nu_{10},\nu_{11})&=(0,-\nu_{10}),\\
T^1(\sigma_{00},\sigma_{10})&=(-\sigma_{10},0),
\end{aligned}
\]
其余情形均为零。于是 $\id+\mu^1_{\widetilde{\mathcal D}}T^1+T^1\mu^1_{\widetilde{\mathcal D}}$ 是到一个代表 $\mu^1_{\widetilde{\mathcal D}}$ 上同调的子空间的投影。按照\stref{笔记}{1.13}所述的方法，应用扰动引理，便得到一个拟同构的 $\Ai$ 范畴。\znote{52}{扰动引理是这套证明的技术核心：若一个复形经收缩同伦等价于较小的复形，那么对它施加的“小扰动”可以转移到小复形上，并给出显式的级数公式。命题 3.8 的证明正是先用 $T^1$ 把 $\widetilde{\mathcal D}$ 的微分扰动为 $\mu^1_{\widetilde{\mathcal D}}$，再由扰动引理读出所需的 $\Ai$ 范畴。[来源：nLab "homological perturbation theory"，https:/\allowbreak /\allowbreak ncatlab.org/\allowbreak nlab/\allowbreak show/\allowbreak homological+\allowbreak perturbation+\allowbreak theory（2026-11-15 访问）]}显式计算表明，它恰好就是\stref{命题}{3.8}中的范畴 $\mathcal D$（注意，出于次数的原因，阶数大于三的复合都为零，故计算到此为止）。因此得到 $\Ai$ 函子 $\G:\mathcal D\to\widetilde{\mathcal D}\subset\Tw\A$，满足
\[
\G(Z_k)=\widetilde Z_k,\quad[\G^1(x_1)]=[c],\quad[\G^1(x_2)]=[i],\quad[\G^1(x_3)]=[p].
\]

现在设 $H(\A)$ 中有正合三角\eref{3.9}。令 $\C\subset\A$ 为由三角的三个对象组成的全子范畴；令 $\widetilde{\C}\subset\Tw\A$ 为包含这些对象在 $\A\hookrightarrow\Tw\A$ 下的像以及 $\Cone(c)$ 的全子范畴，其中 $c=c_1$。由\stref{引理}{3.31}与\stref{引理}{3.27}，显然的嵌入 $\C\hookrightarrow\widetilde{\C}$ 是拟等价，故由\stref{定理}{2.9}可以找到逆向拟等价 $\mathcal K:\widetilde{\C}\to\C$。这个逆在对象 $Y_0,Y_1$ 上为恒等，将 $\Cone(c)$ 映为 $Y_2$，并将 $[i],[p]$ 分别映为 $[c_2],[c_3]$。于是 $\F=\mathcal K\circ\G:\mathcal D\to\C\subset\A$ 就是具有所需性质的 $\Ai$ 函子。

还需去掉严格含幺的假设。设 $\A$ 仅为 c-含幺，并考虑其中的某个正合三角。先取最小模型，再用\stref{引理}{2.1}将其严格化，得到严格含幺范畴 $\widetilde{\A}$ 及拟等价 $\Phi:\A\to\widetilde{\A}$。由\stref{推论}{3.9}，所给三角在 $\Phi$ 下的像仍正合。对这个像应用前面的讨论，得到具有所需性质的 $\Ai$ 函子 $\widetilde{\F}:\mathcal D\to\widetilde{\A}$。再取\stref{定理}{2.9}给出的逆向拟等价 $\Psi$，令 $\F=\Psi\circ\widetilde{\F}$，并利用其与原对象及态射的上同调同构作识别，即得结论。
\end{proof}

\begin{lemma}\label{st:3.32}
$\Tw\A$ 由其全子范畴 $\A$ 生成。
\end{lemma}
\begin{proof}
取扭曲复形 $Y$ 及其子复形滤过 $F^iY$。由于 $\delta_Y$ 严格下三角，对每个 $i$，它都诱导映射 $d^i\in\hom^1_{\Tw\A}(F^iY/F^{i+1}Y,F^{i+1}Y)$。若将其视为从 $S^{-1}(F^iY/F^{i+1}Y)$ 到 $F^{i+1}Y$ 的零次态射，那么相应映射锥同构于 $F^iY$，故有正合三角
\SourceEquationDiagram{diagram-3-31}{0.31361505}{3.31}
\SourcePage{61}
由定义，$F^iY/F^{i+1}Y$ 的微分为零，因而同构于 $\A$ 中对象的若干平移的直和。利用这一点及\eref{3.31}，可见 $Y$ 能够由 $\A$ 中的对象反复施加锥与平移构造出来，这正是生成子范畴的定义所要求的。
\end{proof}

\begin{lemma}\label{st:3.33}
$\Ai$ 范畴 $\A$ 三角化，当且仅当嵌入 $\A\hookrightarrow\Tw\A$ 是拟等价。
\end{lemma}
\begin{proof}
设该嵌入为拟等价。$\A$ 中任意零次上闭链的抽象映射锥，都由 $\Tw\A$ 的某个对象拟表示，故由假设也由 $\A$ 的某个对象拟表示。任一方向的平移也具有同样性质；由定义，这就是说 $\A$ 是三角化的 $\Ai$ 范畴。

反之，设 $\A$ 三角化。由\stref{引理}{3.32}，$H^0(\Tw\A)$ 的任意对象，在同构意义下，都可以从 $H^0(\A)$ 的对象出发，反复取映射锥与任一方向的平移构造出来。不过，由于 $\A$ 三角化，所有这些运算都不会离开 $H^0(\Tw\A)$ 中的如下全子范畴：其对象同构于 $H^0(\A)$ 的某个对象。因此该嵌入是拟等价。
\end{proof}
由此结合\stref{命题}{3.29}与\stref{引理}{3.31}，可以推出\stref{命题}{3.14}。具体地，$H^0(\Tw\A)$ 是经典意义下的三角范畴；而 $H^0(\A)$ 是与之等价的全子范畴，具有相同的平移函子及相应的殊三角类，因此它本身也是三角范畴。\znote{53}{此处 $H^0(\A)$ 作为 $H^0(\Tw\A)$ 的全子范畴，与后者有相同的平移函子，殊三角类也相同（由\stref{引理}{3.31}），所以“本身也是三角范畴”并非形式推论：等价的全子范畴只要在三角运算下封闭，就继承三角结构。[来源：nLab "triangulated category"，https:/\allowbreak /\allowbreak ncatlab.org/\allowbreak nlab/\allowbreak show/\allowbreak triangulated+\allowbreak category（2026-11-15 访问）]}

\begin{lemma}\label{st:3.34}
设 $\F:\A\to\B$ 是上同调全忠实的 $\Ai$ 函子。假设 $\B$ 是三角化的 $\Ai$ 范畴，且 $\F$ 的像中的对象生成 $\B$。那么存在拟等价 $\widetilde{\F}:\Tw\A\to\B$，它在 $\A\subset\Tw\A$ 上的限制，在 $H^0(\operatorname{fun}(\A,\B))$ 中与 $\F$ 同构。
\end{lemma}
\begin{proof}
由假设，$\B\hookrightarrow\Tw\B$ 是拟等价。取\stref{定理}{2.9}给出的逆向拟等价 $\mathcal K:\Tw\B\to\B$，定义 $\widetilde{\F}=\mathcal K\circ\Tw\F$。它在 $\A$ 上的限制显然满足要求，而且至少是上同调全忠实的。由假设，$H^0(\B)$ 的任意对象，在同构意义下，都可以从形如 $\F(X)$ 的对象出发，反复取映射锥与任一方向的平移构造出来。但 $\widetilde{\F}$ 本身保持正合三角和平移，见\stref{推论}{3.9}，故所有这些运算都不会离开同构于 $\widetilde{\F}$ 的像中对象的那一类对象。因此 $\widetilde{\F}$ 是拟等价。
\end{proof}
这一结论结合\stref{引理}{3.32}，蕴含\stref{引理}{3.20}。

\SourcePage{62}
\subsection{\texorpdfstring{$\Ai$ 模的映射锥}{A∞ 模的映射锥}}\label{sec:3s}
现在考虑\stref{例题}{3.13}。设 $t:\M_0\to\M_1$ 是 $\Ai$ 模之间的零次同态。
\begin{definition*}
定义 $\C=\ACone(t)$ 为如下 $\Ai$ 模：
\begin{equation}\label{eq:3.32}
\begin{aligned}
\C(X)&=\M_0(X)[1]\oplus\M_1(X),\\
\mu^d_{\C}((b_0,b_1),a_{d-1},\ldots,a_1)
&=\bigl(\mu^d_{\M_0}(b_0,a_{d-1},\ldots,a_1),\\[-2pt]
&\qquad\mu^d_{\M_1}(b_1,a_{d-1},\ldots,a_1)+t^d(b_0,a_{d-1},\ldots,a_1)\bigr).
\end{aligned}
\end{equation}
\end{definition*}
这显然推广了\eref{3.6}；后者对应 $\M_k=\mathcal I(Y_k)$、$t=\mathcal I^1(c)$ 均来自 Yoneda 嵌入的情形。存在次数分别为零与一的典范模同态 $\iota:\M_1\to\C$、$\pi:\C\to\M_0$，其定义与刚才所说的特例完全相同。
\begin{lemma}\label{st:3.35}
$H^0(\Q)$ 中的下列三角正合：
\SourceDiagram{diagram-3-35}{0.21392285}
\end{lemma}
\begin{proof}
考虑全嵌入 $\Q\hookrightarrow\Tw\Q$。在后一个范畴中，有\eref{3.28}定义的典范锥 $\Cone(t)$。存在自然给出的态射
\begin{equation}\label{eq:3.33}
\C\longrightarrow\Cone(t),
\end{equation}
它对所有 $\M\in\Ob\Tw\Q$ 都诱导同构 $\hom_{\Tw\Q}(\M,\C)\to\hom_{\Tw\Q}(\M,\Cone(t))$（事实上，只需对 $\M\in\Ob\Q$ 检查这一点，因为一般情形可由一个直接的滤过论证推出）。这尤其意味着\eref{3.33}是同构。此外，很容易验证，它把 $\Cone(t)$ 与其组成部分之间的自然态射 $i,p$，同上面定义的 $\iota,\pi$ 联系起来。再用\stref{引理}{3.27}与\stref{引理}{3.31}，即得结论。
\end{proof}
这表明，$H^0(\Q)$ 中的任意态射都能补成一个正合三角。任一方向的平移模的存在性是显然的，因此 $\Q$ 三角化。\znote{54}{$\Q=\modcat(\A)$ 的三角化说明：$\Ai$ 模范畴可以充当 $\A$ 的三角化包络，这与导出范畴由链复形范畴对拟同构作局部化得到是同一类构造。区别在于这里用拟表示对象代替形式局部化，构造是具体的、可计算的。[来源：nLab "triangulated category"，https:/\allowbreak /\allowbreak ncatlab.org/\allowbreak nlab/\allowbreak show/\allowbreak triangulated+\allowbreak category（2026-11-15 访问）]}

为使讨论完整，我们回到\stref{笔记}{3.21}提出的要点。由于 $\Q$ 三角化，可以利用\stref{引理}{3.34}构造上同调全忠实函子 $\Tw\A\to\Q$。还有一种更具体的构造，参见\bcite{82}第 7 节：将 $\A$ 与 $\Tw\A$ 的 Yoneda 嵌入组合，得到 $\Ai$ 函子
\[
\widetilde{\mathcal I}:\Tw\A\xrightarrow{\mathcal I_{\Tw\A}}\modcat(\Tw\A)\xrightarrow{\mathcal I_{\A}^{*}}\modcat(\A)=\Q,
\]
其在 $\A\subset\Tw\A$ 上的限制显然等于 $\mathcal I_{\A}$。
\SourcePage{63}
\begin{lemma}\label{st:3.36}
$\widetilde{\mathcal I}$ 上同调全忠实。
\end{lemma}
\begin{proof}
考虑 $\Tw\A$ 中满足如下性质的对象 $Y$ 所成的子集：对所有 $X\in\Ob\A$，
\begin{equation}\label{eq:3.34}
H(\widetilde{\mathcal I}):\Hom_{H(\Tw\A)}(X,Y)\longrightarrow\Hom_{H(\Q)}(\widetilde{\mathcal I}(X),\widetilde{\mathcal I}(Y))
\end{equation}
都是同构。由于 $\widetilde{\mathcal I}$ 延伸了 $\A$ 的 Yoneda 嵌入，当 $Y$ 是 $\A$ 的对象时，这显然成立。\stref{推论}{3.9}说明，$\widetilde{\mathcal I}$ 将正合三角映为正合三角（事实上，可以直接验证：它严格地将锥映为锥，而不只是同构意义下如此）。利用\stref{推论}{3.15}中的相应长正合列，可以证明满足条件的 $Y$ 的类对映射锥封闭。再由平移相容性及\stref{引理}{3.32}，可知所有 $Y$ 都有这一性质。现在修改\eref{3.34}的要求，使双射性对所有 $X\in\Ob\Tw\A$ 都成立，再重复上述论证即可。
\end{proof}

\section{幂等元}\label{sec:4}
\subsection{背景}\label{sec:4a}
先回顾范畴论中的一些初等概念。设 $A$ 是 $\K$ 上的线性范畴，并设某个对象 $Y$ 有幂等自同态 $p\in\Hom_A(Y,Y)$。
\begin{definition*}
对象 $Z$ 连同包含态射与收缩态射 $k,r$，满足
\begin{equation}\label{eq:4.1}
Z\mathrel{\mathop{\rightleftarrows}\limits^{k}_{r}}Y,\qquad rk=e_Z,\quad kr=p,
\end{equation}
称为 $p$ 的\emph{像}。
\end{definition*}
称其为“这个”像是有理由的：给定两个三元组 $(Z,k,r)$、$(Z',k',r')$，两者之间有典范同构 $Z\cong Z'$。
\begin{definition*}
若 $A$ 包含所有幂等自同态的像，就称 $A$ \emph{分裂闭}。
\end{definition*}
特别地，若 $A$ 分裂闭且为加性范畴，那么任意幂等自同态都给出直和分解 $Y=Z\oplus Z^{\perp}$，其中两个直和项分别是 $p$ 与互补投影 $p^{\perp}=e_Y-p$ 的像。
\begin{definition*}
范畴 $A$ 的\emph{分裂闭包}是二元组 $(B,F)$，其中 $B$ 是分裂闭范畴，$F:A\to B$ 是全忠实嵌入，并且 $B$ 的任意对象都是 $A$ 中某个幂等自同态的像。
\end{definition*}
对任意 $A$，都有一个平凡的分裂闭包构造，称为 \emph{Karoubi 完备化} $PA$。\znote{55}{Karoubi 完备化在幂等完备化理论中是标准构造：$PA$ 的对象可等价地写成满足 $p^2=p$ 的二元组 $(Y,p)$，两对象间的态射 $\phi$ 须满足 $\phi=\phi p_0=p_1\phi$，这正是本节公式的另一种写法。嵌入 $A\to PA$ 全忠实，且 $PA$ 的每个对象都是 $A$ 中某个对象的收缩（直和项）。在一切幂等自同态都可分裂的范畴中，$PA$ 是接收 $A$ 的泛对象，故任意分裂闭包都由此得到唯一的等价扩张。[来源：nLab "Karoubi envelope"，https:/\allowbreak /\allowbreak ncatlab.org/\allowbreak nlab/\allowbreak show/\allowbreak Karoubi+\allowbreak envelope（2026-11-15 访问）。]}$PA$ 的对象是二元组 $(Y,p)$，态射空间则是相应的直和项
\[
\Hom_{PA}((Y_0,p_0),(Y_1,p_1))=p_1\Hom_A(Y_0,Y_1)p_0.
\]
嵌入 $A\to PA$ 为 $Y\mapsto(Y,e_Y)$。若 $(B,F)$ 是 $A$ 的任意分裂闭包，则 $F$ 典范地延伸为等价 $PA\to B$。这就证明，任意两个分裂闭包都是等价的，并且这一等价与 $A$
\SourcePage{64}
的有关嵌入相容。当然，在实际工作中，分裂闭包的一般概念多少有些多余，因为通常直接使用 $PA$。\znote{56}{“$A$ 上的三角结构诱导 $PA$ 上的三角结构”是一条经典的稳定性结论：三角范畴的幂等完备化仍是三角范畴，完备化函子 $A\to PA$ 是三角函子，正合三角在完备化后仍正合。证明的关键与本节引理 4.1 相同：把态射沿幂等投影作直和项分解，再验证补全三角的第三个对象仍落在完备化之内。也正因如此，后文才可以在 $PA$ 或 $\Pi\A$ 中自由地取直和项。[来源：nLab "Karoubi envelope"，https:/\allowbreak /\allowbreak ncatlab.org/\allowbreak nlab/\allowbreak show/\allowbreak Karoubi+\allowbreak envelope（2026-11-15 访问）；并见 Seidel, Fukaya categories and Picard–Lefschetz theory, §4a。]}

现在转向与我们的目的尤其相关的三角范畴。可以证明，$A$ 上的三角结构诱导 $PA$ 上的三角结构\bcite{18}。尽管我们不需要这一结果，下面一个相关但更简单的观察将在后文用到：
\begin{lemma}\label{st:4.1}
设 $B$ 是分裂闭三角范畴，$A\subset B$ 是三角全子范畴（即非空，并且对平移与正合三角封闭）。令 $\widetilde A\subset B$ 为如下全子范畴：其对象是所有同构于 $A$ 中某个对象的直和项的对象，故 $\widetilde A\cong PA$。那么 $\widetilde A$ 仍是 $B$ 的三角全子范畴。
\end{lemma}
\begin{proof}
取 $\widetilde A$ 中的态射 $\tilde c:\widetilde Y_0\to\widetilde Y_1$。由定义，每个 $\widetilde Y_j$ 都同构于 $A$ 中某个对象 $Y_j$ 的直和项。将相应的包含态射、收缩态射分别记为 $k_j,r_j$。将 $\tilde c$ 与 $c=k_1\tilde c r_0$ 各自补成正合三角，第三个对象分别记为 $\widetilde Y_2,Y_2$。存在态射 $r_2,k_2$，补全下列交换图：
\SourceDiagram{diagram-4-1}{0.57323015}
其中 $T$ 为平移函子。由于 $j=0,1$ 时 $r_jk_j=\id_{\widetilde Y_j}$，所以 $r_2k_2$ 是 $\widetilde Y_2$ 的自同构（这是三角结构公理的标准推论；若 $A=H^0(\A)$ 来自一个三角化的 $\Ai$ 范畴，我们已将其陈述为\stref{推论}{3.16}）。因此，$\widetilde Y_2$ 是 $Y_2$ 的幂等自同态 $p=k_2(r_2k_2)^{-1}r_2$ 的像，故按定义属于 $\widetilde A$。
\end{proof}

\subsection{幂等元与\texorpdfstring{$\Ai$}{A∞} 模}\label{sec:4b}
以显然的方式，把基域 $\K$ 视为只有一个对象 $*$ 的 $\Ai$ 范畴。当然，它的上同调范畴仍是 $\K$，只是此时将其视为普通范畴。我们在记号上不区分这两者；这应当只会引起极少的混淆。

设 $\A$ 是 $\Ai$ 范畴，$A=H(\A)$ 是其上同调范畴，并令 $\Q=\modcat(\A)$ 为 $\Ai$ 模范畴。
\begin{definition*}
对象 $Y$ 的\emph{同伦幂等元}定义为非含幺 $\Ai$ 函子 $\wp:\K\to\A$，满足 $\wp(*)=Y$。\znote{56}{把基域 $\K$ 视为单对象 $\Ai$ 范畴时，非含幺 $\Ai$ 函子 $\wp:\K\to\A$ 只由作用在唯一对象上的一列元素决定，其相容性方程恰为 (4.2)。在 $\Ai$ 代数的语言中，结构映射 $m_n$ 的次数为 $n-2$（同调次数约定），而 $\wp^d$ 的次数为 $1-d$，两者相差一个平移，这正是函子性方程能逐阶写出的原因。当 $\wp^d=0$（$d\geq2$）时，(4.2) 退化为 $\wp^1\circ\wp^1=\wp^1$，即普通幂等元。[来源：nLab "A-infinity-category"，https:/\allowbreak /\allowbreak ncatlab.org/\allowbreak nlab/\allowbreak show/\allowbreak A-infinity-category（2026-11-15 访问）。]}具体地，它由元素 $\wp^d\in\hom^{1-d}_{\A}(Y,Y)$（$d\geq1$）给出，满足一系列方程
\begin{equation}\label{eq:4.2}
\sum_r\sum_{s_1,\ldots,s_r}\mu^r_{\A}(\wp^{s_1},\ldots,\wp^{s_r})=
\begin{cases}\wp^{d-1},&d\text{ 为偶数},\\0,&d\text{ 为奇数},\end{cases}
\end{equation}
\SourcePage{65}
其中对分拆 $s_1+\cdots+s_r=d$ 求和。
\end{definition*}
\begin{lemma}\label{st:4.2}
设 $p\in\Hom_A(Y,Y)$ 是上同调范畴中的幂等自同态。那么存在同伦幂等元 $\wp$，使 $[\wp^1]=p$。
\end{lemma}
\begin{proof}
$p$ 定义了非含幺函子 $\K\to A$。由\stref{引理}{1.9}，将它提升为 $\Ai$ 函子的障碍位于 $\HH^2(\K,A)^{2-r}$ 中，其中 $r>2$。基础 Hochschild 复形为
\[
\begin{gathered}
CC^{r+s}(\K,A)^s=\Hom_A^s(Y,Y),\\
\partial^{r,s}(T)=(-1)^s
\begin{cases}pT-Tp,&r\text{ 为偶数},\\pT-Tp^{\perp},&r\text{ 为奇数}.\end{cases}
\end{gathered}
\]
将 $\Hom_A(Y,Y)$ 按照与 $p$ 作左复合、右复合的像与核进行分解，就可直接看出，唯一非零的上同调群为
\[
\HH^s(\K,A)^s=p\Hom_A^s(Y,Y)p\oplus p^{\perp}\Hom_A^s(Y,Y)p^{\perp}.
\]
\end{proof}
\begin{note}\label{st:4.3}
除了直接计算 Hochschild 复形，也可以使用一些基本的同调代数。先向 $\K$ 添加一个单位元，使所讨论的函子变为含幺的；这样，\stref{笔记}{2.7}中的相应双模范畴就是半单的，因此所有高阶 $\operatorname{Ext}$ 群都为零。
\end{note}
设 $\wp$ 是同伦幂等元。与之相应，我们定义一个 $\Ai$ 模 $\mathcal Z$，称为 $\wp$ 的\emph{抽象像}：
\begin{definition*}

\[
\mathcal Z(X)=\hom_{\A}(X,Y)[q]
\]
是以 $\hom_{\A}(X,Y)$ 为系数、以分次变量 $q$（$|q|=-1$）为变量的形式多项式空间。用 $\delta_q$ 表示正规化的形式微分运算：它将 $q^0$ 映为零，将其他幂 $q^k$ 映为 $q^{k-1}$。用这一记号，微分为
\begin{equation}\label{eq:4.3}
\begin{aligned}
\mu^1_{\mathcal Z}(b(q))={}&\sum_r\sum_{s_2,\ldots,s_r}\delta_q^{s_2+\cdots+s_r}\mu^r_{\A}(\wp^{s_r},\ldots,\wp^{s_2},b(q))\\
&+(-1)^{|b|}\frac{b(q)-b(-q)}{2q},
\end{aligned}
\end{equation}
其中最后一项是先将幂级数反对称化，再施加 $\delta_q$；它实际上不涉及乘以 $\frac12$，因此在 $\K$ 的特征为二时也有意义。对于 $d\geq2$，高阶结构映射为
\[
\begin{aligned}
&\mu^d_{\mathcal Z}(b(q),a_{d-1},\ldots,a_1)\\
&\quad=\sum_r\sum_{s_{d+1},\ldots,s_r}\delta_q^{s_{d+1}+\cdots+s_r}
\mu^r_{\A}(\wp^{s_r},\ldots,\wp^{s_{d+1}},b(q),a_{d-1},\ldots,a_1).
\end{aligned}
\]
\end{definition*}
\SourcePage{66}
\begin{lemma}\label{st:4.4}
基础上同调模 $H(\mathcal Z)$ 可以如下描述：$H(\mathcal Z)(X)\cong p\Hom_A(X,Y)$，其中 $p=[\wp^1]\in\Hom_A(Y,Y)$；模结构 $H(\mathcal Z)(X_1)\otimes\Hom_A(X_0,X_1)\to H(\mathcal Z)(X_0)$ 来自通常的复合
\[
p\Hom_A(X_1,Y)\otimes\Hom_A(X_0,X_1)\longrightarrow p\Hom_A(X_0,Y).
\]
\end{lemma}
\begin{proof}
$\mathcal Z(X)$ 有一个自然的递增滤过，即 $F_r=\hom_{\A}(X,Y)q^0\oplus\cdots\oplus\hom_{\A}(X,Y)q^{r-1}$。我们通常使用递降滤过；为统一起见，将 $r$ 换成 $-r$，就得到一个递降滤过。相应谱序列满足
\[
E_1^{-r,s}=H^{s-r}(F_{r+1}/F_r)=
\begin{cases}\Hom_A^s(X,Y),&r\geq0,\\0,&\text{其余情形}.\end{cases}
\]
微分由\eref{4.3}中恰含一个 $\delta_q$ 或 $1/q$ 的项诱导，因此为
\[
\partial^{-r,s}(b)=(-1)^{r+s}
\begin{cases}p^{\perp}b,&r>0\text{ 且为奇数},\\pb,&r>0\text{ 且为偶数},\\0,&\text{其余情形}.\end{cases}
\]
所以谱序列在 $E_2$ 页退化，$H(\mathcal Z)(X)$ 同构于唯一非零的列 $E_2^{0\bullet}=p\Hom_A(X,Y)$。事实上，谱序列的边缘同态给出这两个群之间的典范同构；利用这一同构，所需的模结构相容性就是显然的。
\end{proof}
\begin{lemma}\label{st:4.5}
对任意 $\Ai$ 模 $\M$，有同构
\[
\Hom_{H(\Q)}(\mathcal Z,\M)\cong H(\M)(Y)p,
\]
并且它对于以模同态作左复合而言是自然的。
\end{lemma}
\begin{proof}
考虑\stref{笔记}{2.15}中的 Abel 范畴 $\operatorname{Mod}^0(A)$。\stref{引理}{4.4}说明，作为该范畴的对象，$H(\mathcal Z)$ 是 $\Hom_A(-,Y)$ 的直和项；后者为投射对象，故前者也是。结合\eref{2.14}，这意味着 $\hom_{\Q}(\mathcal Z,\M)$ 上长度滤过所诱导的谱序列在 $E_2$ 页退化，而且映射
\[
\Hom_{H(\Q)}(\mathcal Z,\M)\longrightarrow\Hom_{\operatorname{Mod}^0(A)}(H(\mathcal Z),H(\M)[\bullet])
\]
是同构。由经典 Yoneda 引理及\stref{引理}{4.4}，后一个分次群自然同构于 $H(\M)(Y)p$。
\end{proof}
考虑将 Yoneda 嵌入应用于 $p$ 所得的 $H^0(\Q)$ 中的幂等自同态。由\stref{引理}{4.5}直接可知，$\mathcal Z$ 是这个自同态在\eref{4.1}意义下的像。特别地，在 $H^0(\Q)$ 中的同构意义下，抽象像只依赖于 $p$，而不依赖于 $\wp$ 的选择。
\SourcePage{67}
\begin{note}\label{st:4.6}
$\mathcal Z$ 的定义受到\bcite{24}中一个观察的启发：在适当的复形范畴中，利用无限分解的技巧可以得到直和项的存在性（感谢 François Loeser 向我指出这一点）。
\end{note}

\subsection{分裂闭的\texorpdfstring{$\Ai$}{A∞} 范畴}\label{sec:4c}
\begin{definition*}
若任意同伦幂等元的抽象像都由 $\A$ 的某个对象拟表示，就称 $\A$ \emph{分裂闭}。$\Ai$ 范畴 $\A$ 的\emph{分裂闭包}是二元组 $(\B,\F)$，其中 $\B$ 为分裂闭 $\Ai$ 范畴，$\F:\A\to\B$ 为上同调全忠实函子，且在 $H^0(\B)$ 中，每个对象都是 $H^0(\A)$ 中某个对象的幂等自同态的像。
\end{definition*}
表述这些概念时，我们沿用了\sr{sec:3}{第 3 节}的一般模式，但其实可以大幅简化它们的含义。首先，$\A$ 分裂闭，当且仅当 $H^0(\A)$ 具有这一性质（按最初在上同调层次上的含义）。事实上，设 $H^0(\A)$ 分裂闭，并取同伦幂等元 $\wp$，满足 $[\wp^1]=p$。由假设，$p$ 在 $H^0(\A)$ 中有像 $Z$。利用\stref{引理}{4.5}，很容易证明 $Z$ 的 Yoneda 模在 $H^0(\Q)$ 中同构于 $\wp$ 的抽象像，因此给出了一个拟表示对象。逆向蕴含类似，但还需使用\stref{引理}{4.2}。直接推论是：$(\B,\F)$ 为 $\A$ 的分裂闭包，当且仅当 $(H^0(\B),H^0(\F))$ 为 $H^0(\A)$ 的分裂闭包。
\begin{lemma}\label{st:4.7}
每个 $\Ai$ 范畴都有分裂闭包。此外，若 $(\B,\F)$ 与 $(\widetilde{\B},\widetilde{\F})$ 是 $\A$ 的两个分裂闭包，则存在拟等价 $\G:\B\to\widetilde{\B}$，使 $\G\circ\F$ 在 $H^0(\operatorname{fun}(\A,\widetilde{\B}))$ 中与 $\widetilde{\F}$ 同构。
\end{lemma}
\begin{proof}
用 $\Pi\A$ 表示 $\Q=\modcat(\A)$ 的如下全 $\Ai$ 子范畴：其对象是在 $H^0(\Q)$ 中同构于某个同伦幂等元 $\wp$ 的抽象像 $\mathcal Z$ 的所有对象。我们断言，$\Pi\A$ 包含所有与 $\A$ 的对象 $Y$ 相应的 Yoneda $\Ai$ 模 $\Y$。为说明这一点，取同伦幂等元，使 $[\wp^1]$ 为 $Y$ 的恒等自同态（由\stref{引理}{4.2}，它存在）；由\stref{引理}{4.5}，$\wp$ 的抽象像同构于 $Y$ 的 Yoneda 模。于是二元组 $(\Pi\A,\mathcal I_{\A}:\A\to\Pi\A\subset\Q)$ 就是 $\A$ 的分裂闭包。事实上，如前所述，只需考察上同调范畴；在那里，同样结合\stref{引理}{4.2}与\stref{引理}{4.5}就得到所需结论。

对于第二个结论，设 $(\B,\F)$ 是 $\A$ 的任意分裂闭包。利用 $\B$ 的 Yoneda 嵌入与 $\Ai$ 模的拉回，定义 $\Ai$ 函子
\[
\G:\B\xrightarrow{\mathcal I_{\B}}\modcat(\B)\xrightarrow{\F^*}\Q.
\]
由于 $\F$ 上同调全忠实，$\G\circ\F$ 同构于 $\mathcal I_{\A}$，见\eref{2.13}。这意味着，将 $\G$ 限制在形如 $\F(Y)$ 的对象上时，它上同调全忠实；但由定义，$H^0(\B)$ 的每个对象都是某个 $\F(Y)$ 的幂等
\SourcePage{68}
自同态的像，因此 $\G$ 在整个 $\B$ 上都上同调全忠实。类似论证表明，$\G$ 的像包含于 $\Pi\A$，而且 $\G$ 实际上是拟等价 $\B\to\Pi\A$。当然，既然已有任意分裂闭包与 $\Pi\A$ 之间的拟等价，就可以利用\stref{定理}{2.9}构造任意两个这样的闭包之间的拟等价。
\end{proof}
\begin{lemma}\label{st:4.8}
三角化的 $\Ai$ 范畴的分裂闭包仍然三角化。
\end{lemma}
\begin{proof}
设 $\A$ 为三角化的 $\Ai$ 范畴。由\stref{例题}{3.13}，$\Q=\modcat(\A)$ 三角化；经 Yoneda 嵌入，$H^0(\A)$ 被识别为 $H^0(\Q)$ 的一个经典意义下的三角子范畴。此外，$H^0(\Pi\A)$ 是由同构于 $H^0(\A)$ 中某个对象的直和项的对象组成的子范畴，因此由\stref{引理}{4.1}，它仍是 $H^0(\Q)$ 的三角子范畴。换言之，若 $H(\Q)$ 中有正合三角，其中两个对象属于 $\Pi\A$，那么第三个也属于 $\Pi\A$。由\stref{推论}{3.10}，这意味着 $\Pi\A$ 本身是三角化的 $\Ai$ 范畴。
\end{proof}
\begin{definition*}
设 $\A$ 是任意 $\Ai$ 范畴；先取它的三角化包络，再取后者的分裂闭包，将结果记为 $\B$。例如，两次都采用标准选择，就得到 $\B=\Pi(\Tw\A)$。称 $H^0(\B)$ 为\emph{分裂闭导出范畴}，记为 $D^{\pi}(\A)$。
\end{definition*}
与\secref{3j}的讨论相对应，还应提到一个相关概念。设 $\B$ 是分裂闭、三角化的 $\Ai$ 范畴，$\A\subset\B$ 为非空的全 $\Ai$ 子范畴。考虑包含 $\A$、对 $H^0(\B)$ 中的同构封闭、且本身三角化并分裂闭的最小全子范畴 $\widetilde{\B}\subset\B$。
\begin{definition*}
称它为 $\B$ 中由 $\A$ \emph{分裂生成}的子范畴。\znote{60}{分裂生成在 §5 中反复出现：推论 5.7 断言，若某族对象使所有 $X$ 的逐次扭转都为零，则该族既生成 $\A$、又使 $\A$ 分裂闭；推论 5.8 则处理 $T_{Y_1}\cdots T_{Y_m}X\cong S^{\sigma}X$（$\sigma\ne0$）的情形，结论是 $Y_k$ 分裂生成 $\A$。直观上，分裂闭包把“由映射锥生成”与“取直和项”两件事一并完成，这正是 Fukaya 范畴中由球面等对象生成整个范畴时出现的现象。[来源：并见 Seidel, Fukaya categories and Picard–Lefschetz theory, §5e（推论 5.7 与 5.8）。]}
\end{definition*}
具体地，从 $\Ob\A$ 出发，取映射锥和平移，任意多次地重复这一过程，最后取先前构造的对象的所有幂等自同态的像，就得到 $\Ob\widetilde{\B}$。若 $\widetilde{\B}=\B$，就称 $\A$ \emph{分裂生成} $\B$，或称 $\A$ 的对象是 $\B$ 的\emph{分裂生成元}。例如，$\Pi(\Tw\A)$ 由 $\A$ 的对象分裂生成。由\stref{引理}{3.20}与\stref{引理}{4.7}中的唯一性结论，立即得到：
\begin{corollary}\label{st:4.9}
设 $\B$ 是分裂闭、三角化的 $\Ai$ 范畴，$\A\subset\B$ 为分裂生成 $\B$ 的全子范畴。那么存在拟等价 $\Pi(\Tw\A)\to\B$，它诱导三角范畴的等价 $H^0(\B)\cong D^{\pi}(\A)$。
\end{corollary}
\begin{note}\label{st:4.10}
分裂闭包 $\Pi\A$ 有一点在美学上不够理想：它并不把 $\A$ 本身包含为全子范畴（当然，它包含一个与 $\A$ 拟同构的子范畴）。现在概述另一种构造，它是无界扭曲复形的一种变体。这一构造更
\SourcePage{69}
具体，也更接近最初的想法；但至少这里给出的版本需要严格含幺。

设 $\A$ 是严格含幺的 $\Ai$ 范畴，并假设它有零对象 $O$，使 $\hom_{\A}(O,X)$ 与 $\hom_{\A}(X,O)$ 都为零群。如下定义另一个 $\Ai$ 范畴 $\Sigma^-\A$。其对象是形式无限序列 $X=(X^0,X^{-1},X^{-2},\ldots)$，其中 $X^i\in\Ob\A$。态射 $a\in\hom^k_{\Sigma^-\A}(X_0,X_1)$ 是无限矩阵
\[
a=(a^{ji}),\qquad a^{ji}\in\hom^{k+j-i}_{\A}(X_0^i,X_1^j),
\]
要求它“本质上下三角”：存在某个 $N$（依赖于 $a$），使 $j<i-N$ 时 $a^{ji}=0$。复合映射通过矩阵乘法定义，与加性扩张的情形相同。注意，每个 $X^i$ 部分都被视为平移了 $i$，这会产生一些符号。事实上，为简化下面的公式，我们采用\stref{笔记}{3.26}中修改后的符号约定。

考虑二元组 $(X,\delta_X)$，其中 $X\in\Ob\Sigma^-\A$，$\delta_X\in\hom^1_{\Sigma^-\A}(X,X)$ 满足\eref{3.19}，并且严格下三角，即 $j\leq i$ 时 $\delta_X^{ji}=0$。定义 $\Ai$ 范畴 $\Tw^-\A$，其对象为这样的二元组，态射的结构映射则像\eref{3.20}那样用微分进行变形。可以把 $\A$ 与 $\Tw^-\A$ 中由序列 $(X^0,X^{-1}=O,X^{-2}=O,\ldots)$ 组成的全 $\Ai$ 子范畴等同。

给定 $Y\in\Ob\A$ 的同伦幂等元 $\wp$，如下定义 $\Tw^-\A$ 的对象 $(Z,\delta_Z)$：
\begin{equation}\label{eq:4.4}
\begin{gathered}
Z^i=Y\quad\text{对所有 }i,\\
\delta_Z^{ji}=\begin{cases}
\wp^1,&i\text{ 为偶数且 }j=i+1\leq0,\\
\wp^1-e_Y,&i\text{ 为奇数且 }j=i+1\leq0,\\
\wp^{j-i},&i+1<j\leq0,\\
0,&\text{其余情形}.
\end{cases}
\end{gathered}
\end{equation}
广义 Maurer–Cartan 方程等价于\eref{4.2}。定义
\[
\begin{aligned}
r&\in\hom^0_{\Tw^-\A}(Y,Z),&r^{00}&=e_Y,\\
k&\in\hom^0_{\Tw^-\A}(Z,Y),&k^{0i}&=\wp^{1-i},\\
h&\in\hom^{-1}_{\Tw^-\A}(Z,Z),&h^{i-1,i}&=(-1)^i e_Y,
\end{aligned}
\]
除列出的系数外，其余系数都为零。于是
\[
\begin{aligned}
\mu^1_{\Tw^-\A}(r)&=0,&\mu^1_{\Tw^-\A}(k)&=0,\\
\mu^2_{\Tw^-\A}(k,r)&=\wp^1,&\mu^2_{\Tw^-\A}(r,k)-\mu^1_{\Tw^-\A}(h)&=e_Z.
\end{aligned}
\]
\SourcePage{70}
这意味着，在上同调范畴中，$Z$ 是 $Y$ 的幂等自同态 $[\wp^1]$ 的像。因此，若取 $\Tw^-\A$ 中由 $\A$ 的所有对象以及所有\eref{4.4}类型的对象组成的全 $\Ai$ 子范畴，就得到 $\A$ 的一个分裂闭包。为说明它与前一种方法的联系，注意与 $Z$ 相应的 Yoneda $\A$ 模恰好就是抽象像 $\mathcal Z$。
\end{note}

\section{扭转}\label{sec:5}
\subsection{抽象运算}\label{sec:5a}
固定 $\Ai$ 范畴 $\A$ 及其对象 $Y$。给定 $\A$ 上的 $\Ai$ 模 $\M$，可以如下定义另一个模 $\C=\mathcal T_Y\M$：
\begin{definition*}

\begin{equation}\label{eq:5.1}
\begin{aligned}
\C(X)&=(\M(Y)\otimes\hom_{\A}(X,Y))[1]\oplus\M(X),\\
\mu^1_{\C}(c_0\otimes b,c_1)
&=\bigl((-1)^{|b|-1}\mu^1_{\M}(c_0)\otimes b+c_0\otimes\mu^1_{\A}(b),\\[-2pt]
&\qquad\mu^1_{\M}(c_1)+\mu^2_{\M}(c_0,b)\bigr),\\
\mu^d_{\C}((c_0\otimes b,c_1),a_{d-1},\ldots,a_1)
&=\bigl(c_0\otimes\mu^d_{\A}(b,a_{d-1},\ldots,a_1),\\[-2pt]
&\qquad\mu^d_{\M}(c_1,a_{d-1},\ldots,a_1)\\[-2pt]
&\qquad+\mu^{d+1}_{\M}(c_0,b,a_{d-1},\ldots,a_1)\bigr)\quad(d\geq2).
\end{aligned}
\end{equation}
\end{definition*}
$\C$ 的 c-含幺性可以用与\eref{3.6}相同的方法证明。

映射 $\M\mapsto\mathcal T_Y\M$ 是 $\Ai$ 函子 $\mathcal T_Y:\Q\to\Q$ 的一部分；这里仍令 $\Q=\modcat(\A)$，称此函子为\emph{沿 $Y$ 的抽象扭转函子}。预模同态 $t\in\hom_{\Q}(\M_0,\M_1)$ 的像为
\[
\begin{aligned}
\tilde t&=\mathcal T_Y^1(t)\in\hom_{\Q}(\mathcal T_Y\M_0,\mathcal T_Y\M_1),\\
\tilde t^1(c_0\otimes b,c_1)&=\bigl((-1)^{|b|-1}t^1(c_0)\otimes b,\ t^1(c_1)+t^2(c_0,b)\bigr),\\
\tilde t^d((c_0\otimes b,c_1),a_{d-1},\ldots,a_1)
&=\bigl(0,\ t^d(c_1,a_{d-1},\ldots,a_1)\\[-2pt]
&\qquad+t^{d+1}(c_0,b,a_{d-1},\ldots,a_1)\bigr)\quad(d\geq2).
\end{aligned}
\]
高阶项 $\mathcal T_Y^d$（$d\geq2$）为零，因此 $\mathcal T_Y$ 实际上是 dg 函子。考察\eref{2.10}可见，它还是严格含幺的。

可以将抽象扭转的构造分成两步。首先，取\secref{3c}意义下的张量积 $\M(Y)\otimes\Y$。它带有\emph{抽象求值同态}
\begin{equation}\label{eq:5.2}
\begin{aligned}
\phi&:\M(Y)\otimes\Y\longrightarrow\M,\\
\phi^d(c_0\otimes b,a_{d-1},\ldots,a_1)&=\mu^{d+1}_{\M}(c_0,b,a_{d-1},\ldots,a_1).
\end{aligned}
\end{equation}
而 $\mathcal T_Y\M=\ACone(\phi)$ 就是它在\eref{3.32}意义下的锥。特别地，有相应模同态 $\iota,\pi$，与 $\phi$ 一起构成
\SourcePage{71}
$H(\Q)$ 中的正合三角
\SourceEquationDiagram{diagram-5-3}{0.30946121}{5.3}
\begin{note}\label{st:5.1}
扭转运算很容易推广到有多个扭转对象的情形。例如，给定两个对象 $Y_+,Y_-$ 和 $\Ai$ 模 $\M$，定义 $\C=\mathcal T_{Y_-,Y_+}\M$ 为
\[
\begin{aligned}
\C(X)&=\bigl(\M(Y_-)\otimes\hom_{\A}(X,Y_-)\\[-2pt]
&\qquad\oplus\M(Y_+)\otimes\hom_{\A}(X,Y_+)\bigr)[1]\oplus\M(X),\\
\mu^1_{\C}(c_{0,-}\otimes b_-,c_{0,+}\otimes b_+,c_1)
&=\bigl((-1)^{|b_-|-1}\mu^1_{\M}(c_{0,-})\otimes b_-+c_{0,-}\otimes\mu^1_{\A}(b_-)\\[-2pt]
&\qquad+(-1)^{|b_+|-1}\mu^1_{\M}(c_{0,+})\otimes b_++c_{0,+}\otimes\mu^1_{\A}(b_+),\\[-2pt]
&\qquad\mu^1_{\M}(c_1)+\mu^2_{\M}(c_{0,-},b_-)+\mu^2_{\M}(c_{0,+},b_+)\bigr),
\end{aligned}
\]
高阶项也以类似方式定义。若 $\hom_{\A}(Y_-,Y_+)=0$，那么这与 $\mathcal T_{Y_-}(\mathcal T_{Y_+}\M)$ 是同一个 $\Ai$ 模。若 $\hom_{\A}(Y_-,Y_+)$ 仅仅无上同调，则存在一个模到另一个模的显然嵌入，其商无上同调，因此两种运算仍给出 $H^0(\Q)$ 中同构的对象。
\end{note}

\subsection{Yoneda 模的扭转}\label{sec:5b}
\begin{definition*}
给定 $\A$ 的两个对象 $Y_0,Y_1$，用 $T_{Y_0}(Y_1)$ 表示任意一个拟表示 $\mathcal T_{Y_0}(\Y_1)$ 的对象，并称其为\emph{沿 $Y_0$ 扭转 $Y_1$} 的结果。
\end{definition*}
若它存在，就带有 $H^0(\A)$ 中的典范态射 $[i]:Y_1\to T_{Y_0}(Y_1)$，后者在 Yoneda 嵌入下成为 $[\iota]$。

现在假设 $H(\A)$ 对有限直和及平移 $S^{\sigma}$（$\sigma\in\Z$）封闭；并假设 $\Hom_{H(\A)}(Y_0,Y_1)$ 为有限维分次向量空间。那么由\stref{笔记}{3.4}，可以构造张量积 $\Hom_{H(\A)}(Y_0,Y_1)\otimes Y_0$，并且它带有到 $Y_1$ 的典范求值态射 $[f]$，即 $[\phi]$ 在 Yoneda 嵌入下的原像。具体地，若选取 $\Hom_{H(\A)}(Y_0,Y_1)$ 的一组基 $[b^i]$，令 $\sigma^i=-|b^i|$，并把张量积看作 $Y_0$ 的若干平移的直和，那么 $[f]$ 在第 $i$ 个直和项上的限制，就是 $[b^i]$ 在\eref{3.5}中的同构
\[
\Hom^{-\sigma^i}_{H(\A)}(Y_0,Y_1)\cong\Hom^0_{H(\A)}(S^{\sigma^i}Y_0,Y_1)
\]
\SourcePage{72}
下的像。若 $T_{Y_0}(Y_1)$ 在 $\A$ 中存在，它位于三角
\SourceEquationDiagram{diagram-5-4}{0.42369185}{5.4}
之中，其中 $[p]$ 在 Yoneda 嵌入下对应于 $[\pi]$。应用\stref{推论}{3.10}可知，\eref{5.4}正合。\znote{64}{三角 (5.4) 是扭转到扭曲复形情形的具体化：$T_{Y_0}Y_1$ 由态射 $f:\hom_{\A}(Y_0,Y_1)\otimes Y_0\to Y_1$ 的锥给出，$[i]$ 与 $[p]$ 分别是锥的标准态射与投影。在 Bondal–Kapranov 的语言中，它正是由对象 $Y_1$（次数 $0$）与 $Y_0$（次数 $-|b^i|+1$，$i$ 取遍一组基）组成的扭曲复形，其微分由参数化的态射 $\hom(Y_0,Y_1)\otimes Y_0\to Y_1$ 给出。[来源：nLab "twisted complex"，https:/\allowbreak /\allowbreak ncatlab.org/\allowbreak nlab/\allowbreak show/\allowbreak twisted+\allowbreak complex（2026-11-15 访问）。]}

最后，考虑 $\A$ 三角化且 c-有限（上同调有限）的情形。
\begin{definition*}
所谓 c-有限，是指对于任意两个对象 $Y_0,Y_1$，空间 $\Hom_{H(\A)}(Y_0,Y_1)$ 都有限维。
\end{definition*}
$\A$ 的三角化性质蕴含它对直和与平移封闭（\stref{笔记}{3.5}）。更重要的是，求值态射 $[f]$ 在 $\A$ 中有映射锥，并位于相应的标准正合三角中。由\stref{推论}{3.16}，这个锥在 Yoneda 嵌入下的像同构于 $\mathcal T_{Y_0}(\Y_1)$。因此，所有对象 $T_{Y_0}(Y_1)$ 都在 $\A$ 中存在。利用与\secref{3d}相同的论证，可以使扭转成为函子性的。也就是说，对每个 $Y\in\Ob\A$，存在\emph{扭转函子} $T_Y:\A\to\A$，它在 $H^0(\operatorname{fun}(\A,\A))$ 中的同构意义下唯一，并且使图
\SourceEquationDiagram{diagram-5-5}{0.18484595}{5.5}
在 $H^0(\operatorname{fun}(\A,\Q))$ 中在同构意义下交换。

一个重要特例是：$\B=\Tw\A$ 是 c-有限 $\Ai$ 范畴上的扭曲复形范畴。回顾\stref{引理}{3.36}，Yoneda 嵌入 $\A\to\Q=\modcat(\A)$ 自然延伸为上同调全忠实函子 $\widetilde{\mathcal I}:\B\to\Q$。固定 $Y_0\in\Ob\A$、$Y_1\in\Ob\B$。由定义，Yoneda 嵌入将 $T_{Y_0}(Y_1)\in\Ob\B$ 映为在 $\modcat(\B)$ 中定义的抽象扭转 $\mathcal T_{Y_0}(\Y_1)$。此外，抽象扭转的构造与限制 $\modcat(\B)\to\Q$ 相容。因此，$\B$ 与 $\Q$ 中的扭转概念相容：
\begin{equation}\label{eq:5.6}
\widetilde{\mathcal I}(T_{Y_0}(Y_1))\cong\mathcal T_{Y_0}(\widetilde{\mathcal I}(Y_1))\quad\text{在 }H^0(\Q)\text{ 中}.
\end{equation}
也可以由\eref{5.4}与\stref{推论}{3.16}推出这一点，但上面的论证给出了一个典范同构。这不仅稍微更精确，还可以使该同构具有\eref{5.5}意义下的函子性。
\SourcePage{73}
\begin{note}\label{st:5.2}
将上述讨论推广到\stref{笔记}{5.1}的情形是直接的。所得对象 $T_{Y_{0,-},Y_{0,+}}(Y_1)$ 位于正合三角
\SourceDiagram{diagram-5-2}{0.55453786}
之中。
\end{note}

\subsection{判定拟表示对象}\label{sec:5c}
给定 $\Ai$ 范畴 $\A$ 的对象 $Y_0,Y_1,Y_2$，若有一个显式判据，能够判断何时可以取 $Y_2$ 为 $\mathcal T_{Y_0}\Y_1$ 的拟表示对象，将会很有用。可以以\stref{引理}{3.1}为基础，因为它指明了定义态射 $\Y_2\to\mathcal T_{Y_0}\Y_1$ 所需的数据；不过，我们更愿意考虑反方向的映射。这一选择以及下面表述中的技术细节，都由预期的几何应用决定（见\sr{sec:17}{第 17 节}）。\znote{65}{之所以考察“从 $\mathcal T_{Y_0}\Y_1$ 到 $\Y_2$ 的映射”而不是反方向，是因为几何应用中真正给出的是 $Y_2$ 与 $T_{Y_0}Y_1$ 拟同构的一个显式代表：清空循环与伪全纯圆盘的计数自然给出一次态射 $c:Y_1\to Y_2$ 以及把 $\hom(Y_0,Y_1)$ 提升一阶的映射 $k$，即数据 $(c,k)$。若反过来写，同样的几何数据就要摊成无限多个方程。[来源：译者按（上下文推断，无外部依据）：由本节 (5.7) 的方程形状与 §17 的几何来源推得。]}
\begin{lemma}\label{st:5.3}
\textup{(i)} 取二元组 $(c,k)$，其中
\[
c\in\hom^0_{\A}(Y_1,Y_2),\qquad k:\hom_{\A}(Y_0,Y_1)\longrightarrow\hom_{\A}(Y_0,Y_2)[-1],
\]
满足
\begin{equation}\label{eq:5.7}
\begin{aligned}
\mu^1_{\A}(c)&=0,\\
\mu^1_{\A}(k(c_0))-k(\mu^1_{\A}(c_0))+\mu^2_{\A}(c,c_0)&=0\quad\text{对所有 }c_0\in\hom_{\A}(Y_0,Y_1).
\end{aligned}
\end{equation}
它确定模同态 $t:\mathcal T_{Y_0}(\Y_1)\to\Y_2$，满足
\[
[\mu^2_{\Q}(t,\iota)]=[\mathcal I^1(c)]\in\Hom_{H(\Q)}(\Y_1,\Y_2).
\]
\textup{(ii)} 给定 \textup{(i)} 中的二元组 $(c,k)$，以及任意
\[
\gamma\in\hom_{\A}(Y_1,Y_2)[-1],\qquad
\kappa:\hom_{\A}(Y_0,Y_1)\longrightarrow\hom_{\A}(Y_0,Y_2)[-2],
\]
令
\[
\begin{aligned}
\tilde c&=c+\mu^1_{\A}(\gamma),\\
\tilde k(c_0)&=k(c_0)+(-1)^{|c_0|-1}\bigl(\mu^1_{\A}(\kappa(c_0))-\kappa(\mu^1_{\A}(c_0))+\mu^2_{\A}(\gamma,c_0)\bigr),
\end{aligned}
\]
便得到\eref{5.7}的一个新解；与之相应的同态 $\tilde t$ 与原来的 $t$ 上同调。

\SourcePage{74}
\textup{(iii)} 设 $(c,k)$ 如 \textup{(i)}。此外，假设对每个 $X\in\Ob\A$，给定映射
\[
\begin{aligned}
d&:\hom_{\A}(X,Y_1)\longrightarrow\hom_{\A}(X,Y_2),\\
l&:\hom_{\A}(Y_0,Y_1)\otimes\hom_{\A}(X,Y_0)\longrightarrow\hom_{\A}(X,Y_2)[-1],\\
\xi&:\hom_{\A}(X,Y_1)\longrightarrow\hom_{\A}(X,Y_2)[-1],\\
\psi&:\hom_{\A}(Y_0,Y_1)\otimes\hom_{\A}(X,Y_0)\longrightarrow\hom_{\A}(X,Y_2)[-2],
\end{aligned}
\]
对所有 $c_0,b,c_1$ 满足
\begin{align*}
&\mu^1_{\A}(d(c_1))+d(\mu^1_{\A}(c_1))=0,\\
&d(\mu^2_{\A}(c_0,b))+\mu^1_{\A}(l(c_0,b))+(-1)^{|b|-1}l(\mu^1_{\A}(c_0),b)+l(c_0,\mu^1_{\A}(b))=0,\\
&\xi(\mu^1_{\A}(c_1))-\mu^1_{\A}(\xi(c_1))+d(c_1)-\mu^2_{\A}(c,c_1)=0,\\
&(-1)^{|b|-1}\psi(\mu^1_{\A}(c_0),b)+\psi(c_0,\mu^1_{\A}(b))-\mu^1_{\A}(\psi(c_0,b))\\
&\qquad+\xi(\mu^2_{\A}(c_0,b))+l(c_0,b)-\mu^2_{\A}(k(c_0),b)-\mu^3_{\A}(c,c_0,b)=0,
\end{align*}
并且下列链复形无上同调：
\begin{equation}\label{eq:5.8}
\begin{aligned}
D={}&(\hom_{\A}(Y_0,Y_1)\otimes\hom_{\A}(X,Y_0))[2]\\[-2pt]
&\quad\oplus\hom_{\A}(X,Y_1)[1]\oplus\hom_{\A}(X,Y_2),\\
\partial(c_0\otimes b,c_1,c_2)={}&\bigl((-1)^{|b|-1}\mu^1_{\A}(c_0)\otimes b+c_0\otimes\mu^1_{\A}(b),\\[-2pt]
&\quad\mu^1_{\A}(c_1)+\mu^2_{\A}(c_0,b),\ \mu^1_{\A}(c_2)+d(c_1)+l(c_0,b)\bigr).
\end{aligned}
\end{equation}
那么由 $(c,k)$ 定义的同态 $t$ 在 $H^0(\Q)$ 中诱导同构。事实上，只需对组成拟等价全子范畴 $\widetilde{\A}\subset\A$ 的一组对象验证这一条件，它就自动对其余所有对象成立。
\end{lemma}
\begin{proof}
(i) 直接计算表明，按下式定义 $t$，它便具有所有所需性质：
\[
\begin{aligned}
&t^d((c_0\otimes b,c_1),a_{d-1},\ldots,a_1)\\
&\quad=\mu^{d+1}_{\A}(c,c_1,a_{d-1},\ldots,a_1)+\mu^{d+1}_{\A}(k(c_0),b,a_{d-1},\ldots,a_1)\\
&\qquad+\mu^{d+2}_{\A}(c,c_0,b,a_{d-1},\ldots,a_1).
\end{aligned}
\]
(ii) 仿照 (i)，以 $(\gamma,\kappa)$ 代替 $(c,k)$，定义 $u\in\hom^{-1}_{\Q}(\mathcal T_{Y_0}\Y_1,\Y_2)$，便有 $\mu^1_{\Q}(u)=\tilde t-t$。

(iii) 要使 $[t]$ 为同构，必须使 $t^1$ 的映射锥无上同调。具体地，这些锥是链复形
\[
\begin{aligned}
\widetilde D={}&(\hom_{\A}(Y_0,Y_1)\otimes\hom_{\A}(X,Y_0))[2]\\[-2pt]
&\quad\oplus\hom_{\A}(X,Y_1)[1]\oplus\hom_{\A}(X,Y_2),\\
\tilde\partial(c_0\otimes b,c_1,c_2)={}&\bigl((-1)^{|b|-1}\mu^1_{\A}(c_0)\otimes b+c_0\otimes\mu^1_{\A}(b),\\[-2pt]
&\quad\mu^1_{\A}(c_1)+\mu^2_{\A}(c_0,b),\\[-2pt]
&\quad\mu^1_{\A}(c_2)+\mu^2_{\A}(c,c_1)+\mu^2_{\A}(k(c_0),b)+\mu^3_{\A}(c,c_0,b)\bigr).
\end{aligned}
\]
\SourcePage{75}
利用 $(\xi,\psi)$ 可以构造同构 $D\to\widetilde D$，即
\[
(c_0\otimes b,c_1,c_2)\longmapsto(c_0\otimes b,c_1,c_2+\xi(c_1)+\psi(c_0\otimes b)).
\]
所以\eref{5.8}无上同调蕴含所需结论。最后一项断言是显然的，因为限制函子 $\Q=\modcat(\A)\to\modcat(\widetilde{\A})$ 是拟等价（\stref{引理}{2.8}）。
\end{proof}

\begin{lemma}\label{st:5.4}
固定\stref{引理}{5.3}\textup{(i)} 中的 $(c,k)$。假设此外存在
\[
\begin{aligned}
g&=\sum_jg_2^j\otimes g_1^j\in(\hom_{\A}(Y_0,Y_1)\otimes\hom_{\A}(Y_2,Y_0))^1,\\
p&=\sum_jp_2^j\otimes p_1^j\in(\hom_{\A}(Y_0,Y_1)\otimes\hom_{\A}(Y_1,Y_0))^0,\\
q&=\sum_jq_2^j\otimes q_1^j\in\bigl(\hom_{\A}(Y_0,Y_1)\otimes\Hom_{\K}(\hom_{\A}(Y_0,Y_1),\hom_{\A}(Y_0,Y_0))\bigr)^{-1},
\end{aligned}
\]
满足
\[
\begin{aligned}
&\sum_j\bigl((-1)^{|g_1^j|-1}\mu^1_{\A}(g_2^j)\otimes g_1^j+g_2^j\otimes\mu^1_{\A}(g_1^j)\bigr)=0,\\
&\sum_jg_2^j\otimes\mu^2_{\A}(g_1^j,c)\\
&\quad=\sum_j(-1)^{|p_1^j|-1}\mu^1_{\A}(p_2^j)\otimes p_1^j+\sum_jp_2^j\otimes\mu^1_{\A}(p_1^j),\\
&\sum_jg_2^j\otimes\bigl(\mu^2_{\A}(g_1^j,k(\cdot))+\mu^3_{\A}(g_1^j,c,\cdot)\bigr)\\
&\quad=\sum_jq_2^j\otimes\mu^1_{\A}(q_1^j(\cdot))-\sum_jq_2^j\otimes q_1^j(\mu^1_{\A}(\cdot))\\
&\qquad+\sum_j(-1)^{|q_1^j|+|\cdot|-1}\mu^1_{\A}(q_2^j)\otimes q_1^j(\cdot)+\sum_jp_2^j\otimes\mu^2_{\A}(p_1^j,\cdot)\\
&\qquad+\sum_j b^j\otimes(e_{Y_0}\otimes b^{j,\vee}),
\end{aligned}
\]
其中，$e_{Y_0}$ 是 $Y_0$ 的恒等态射的代表，$b^j$ 为 $\hom_{\A}(Y_0,Y_1)$ 的一组基，$b^{j,\vee}$ 为其对偶基，而 $(e_{Y_0}\otimes b^{j,\vee})\in\Hom_{\K}(\hom_{\A}(Y_0,Y_1),\hom_{\A}(Y_0,Y_0))$ 为将 $b^j$ 映为 $e_{Y_0}$、将所有其他基元素映为零的线性映射。那么，$g$ 给出一次模同态 $u:\Y_2\to\hom_{\A}(Y_0,Y_1)\otimes\Y_0$，满足 $[\mu^2_{\Q}(u,t)]=[\pi]$，其中 $[\pi]$ 是\eref{5.3}中的映射。
\end{lemma}
\begin{proof}
定义预模同态
\[
\begin{aligned}
u&:\Y_2\longrightarrow\hom_{\A}(Y_0,Y_1)\otimes\Y_0,\\
u^d(b,a_{d-1},\ldots,a_1)&=\sum_jg_2^j\otimes\mu^{d+1}_{\A}(g_1^j,b,a_{d-1},\ldots,a_1),
\end{aligned}
\]
以及
\SourcePage{76}
\[
\begin{aligned}
v&:\mathcal T_{Y_0}(\Y_1)\longrightarrow\hom_{\A}(Y_0,Y_1)\otimes\Y_0,\\
v^d((c_0\otimes b,c_1),a_{d-1},\ldots,a_1)
&=\sum_jp_2^j\otimes\mu^{d+1}_{\A}(p_1^j,c_1,a_{d-1},\ldots,a_1)\\
&\quad+\sum_jq_2^j\otimes\mu^{d+1}_{\A}(q_1^j(c_0),b,a_{d-1},\ldots,a_1)\\
&\quad+\sum_jp_2^j\otimes\mu^{d+2}_{\A}(p_1^j,c_0,b,a_{d-1},\ldots,a_1)\\
&\quad-\sum_jg_2^j\otimes\mu^{d+2}_{\A}(g_1^j,c,c_1,a_{d-1},\ldots,a_1)\\
&\quad-\sum_jg_2^j\otimes\mu^{d+2}_{\A}(g_1^j,k(c_0),b,a_{d-1},\ldots,a_1)\\
&\quad-\sum_jg_2^j\otimes\mu^{d+3}_{\A}(g_1^j,c,c_0,b,a_{d-1},\ldots,a_1),
\end{aligned}
\]
两者次数分别为一与零。利用假设可算得 $\mu^1_{\Q}(u)=0$，并且 $w=\mu^2_{\Q}(u,t)-\mu^1_{\Q}(v)$ 为
\[
w^d((c_0\otimes b,c_1),a_{d-1},\ldots,a_1)=(-1)^{|c_0|}c_0\otimes\mu^{d+1}_{\A}(e_{Y_0},b,a_{d-1},\ldots,a_1).
\]
由与\stref{引理}{3.6}相同的论证，它与 $\pi$ 上同调。
\end{proof}

\begin{note}\label{st:5.5}
设处于\stref{引理}{5.3}的 (i)、(iii) 所述情形；另假设 $\A$ 对有限直和与平移封闭，且 $\Hom_{H(\A)}(Y_0,Y_1)$ 有限维，从而 $\Hom_{H(\A)}(Y_0,Y_1)\otimes Y_0$ 有意义。由构造，在 $H(\A)$ 中有下列正合三角：
\SourceEquationDiagram{diagram-5-9}{0.39046112}{5.9}
从这个角度看，扭转过程退居幕后：只需利用\stref{引理}{5.3}给出的一组条件，就能判定具有\eref{5.9}这种特殊形式的三角正合。不过要注意，三角中的第三个态射仅被间接刻画：它在 Yoneda 嵌入下的像是 $[\pi]\circ[t]^{-1}$。这时就需要\stref{引理}{5.4}。假设我们有该引理指定的额外数据，并考虑 $H(\A)$ 中的态射
\[
[g]=\sum_j[g_2^j]\otimes[g_1^j]:Y_2\longrightarrow\Hom_{H(\A)}(Y_0,Y_1)\otimes Y_0.
\]
在 Yoneda 嵌入下，它映为模同态 $[u]$，由构造满足 $[u]\circ[t]=[\pi]$。因此，\eref{5.9}中的第三个映射恰好就是 $[g]$。
\end{note}

\subsection{函子性}\label{sec:5d}
设 $\F:\A\to\B$ 为 $\Ai$ 范畴之间的上同调全忠实函子。取 $Y_0,Y_1\in\Ob\A$。
\begin{lemma}\label{st:5.6}
设两个扭转 $T_{Y_0}Y_1$ 与 $T_{\F(Y_0)}\F(Y_1)$ 都存在。那么在 $H^0(\B)$ 中有典范同构
\begin{equation}\label{eq:5.10}
\F(T_{Y_0}Y_1)\cong T_{\F(Y_0)}\F(Y_1).
\end{equation}
\end{lemma}
\SourcePage{77}
\begin{proof}
先给出一个较简单的初步论证。由\stref{推论}{3.16}，正合三角\eref{5.4}在 $H^0(\A)$ 中的同构意义下刻画 $T_{Y_0}(Y_1)$。由此结合\stref{推论}{3.9}，很容易推出\eref{5.10}中的两个对象同构。不过，这并未给出同构的典范选择；要得到更强的结论，必须稍微更仔细地论证。

沿用\eref{1.23}的思路，对任意 $Y\in\Ob\A$ 与 $\M\in\Ob(\modcat(\B))$，都有典范模同态
\[
s:\mathcal T_Y(\F^*\M)\longrightarrow\F^*(\mathcal T_{\F(Y)}\M).
\]
具体地，$s^1$ 将 $(c_0\otimes b,c_1)$ 映为 $((-1)^{|c_0|+|b|-1}c_0\otimes\F^1(b),(-1)^{|c_1|}c_1)$；对于 $d\geq2$，$s^d$ 将 $((c_0\otimes b,c_1),a_{d-1},\ldots,a_1)$ 映为 $((-1)^{\triangle}c_0\otimes\F^d(b,a_{d-1},\ldots,a_1),0)$，其中 $\triangle=|a_1|+\cdots+|a_{d-1}|+|b|+|c_0|-d$。考察 $d=1$ 就可见，$[s]$ 是 $H^0(\modcat(\A))$ 中的同构。\znote{69}{$s$ 比较的是“先扭转再沿 $\F$ 拉回”与“先沿 $\F$ 拉回再扭转”两个模，它在 $d=1$ 时已是 $H^0$ 中的同构，原因在于 $\F$ 上同调全忠实：$s^1$ 的两个分量分别由 $\F^1$ 与恒等给出，而 $\hom$ 空间上的同构正是全忠实性所保证的。高阶项 $s^d$ 的符号 $(-1)^{\triangle}$ 只是把 $\F^d$ 逐个移过前面的元素所生的 Koszul 符号，不影响同构性。[来源：并见 Seidel, Fukaya categories and Picard–Lefschetz theory, §5d（引理 5.6 的证明）。]}对 $Y=Y_0$、$\M=(\mathcal I_{\B}\circ\F)(Y_1)$ 应用这一点，得到
\begin{equation}\label{eq:5.11}
\mathcal T_{Y_0}(\F^*\circ\mathcal I_{\B}\circ\F)(Y_1)
\cong\F^*\bigl(\mathcal T_{\F(Y_0)}(\mathcal I_{\B}\circ\F)(Y_1)\bigr).
\end{equation}
由\eref{2.13}，\eref{5.11}的左边典范同构于 $\mathcal T_{Y_0}(\mathcal I_{\A}(Y_1))$；由 $T_{Y_0}$ 的定义，它典范同构于 $\mathcal I_{\A}(T_{Y_0}(Y_1))$，从而也典范同构于 $(\F^*\circ\mathcal I_{\B}\circ\F)(T_{Y_0}Y_1)$。对右边应用 $T_{\F(Y_0)}$ 的定义，可知它典范同构于 $(\F^*\circ\mathcal I_{\B})(T_{\F(Y_0)}\F(Y_1))$。所以，\eref{5.10}中的两个对象在施加 $\F^*\circ\mathcal I_{\B}$ 后成为典范同构的对象。再利用\eref{2.13}，可知 $\F^*\circ\mathcal I_{\B}$ 在 $\F$ 的像上上同调全忠实，因此在 $H^0(\B)$ 中同构于某个 $\F(X)$ 的对象所成的稍大对象类上，也上同调全忠实。\eref{5.10}的两边都属于这一类，故得到它们之间的典范同构。
\end{proof}
这样构造的同构关于 $Y_1$ 是自然的（事实上，再多作一些工作，就能将它延伸为函子范畴 $H^0(\operatorname{fun}(\A,\B))$ 中 $\F\circ T_{Y_0}$ 与 $T_{\F(Y_0)}\circ\F$ 之间的典范同构）。它也与典范映射 $Y_1\to T_{Y_0}Y_1$、$\F(Y_1)\to T_{\F(Y_0)}\F(Y_1)$ 相容。若进一步假设这些 $\Ai$ 范畴对有限直和与平移封闭，且有关的 $\Hom$ 空间有限维，那么正合三角\eref{5.4}在 $H(\F)$ 下的像，典范同构于对象 $\F(Y_0),\F(Y_1)$ 的相应三角。

\subsection{扭转的复合}\label{sec:5e}
取 c-有限、三角化的 $\Ai$ 范畴 $\A$，固定对象 $Y_1,\ldots,Y_m$。从任意对象 $X=X_m$ 出发，依次进行扭转，令 $X_{k-1}=T_{Y_k}X_k$。由定义，有正合三角
\SourceDiagram{diagram-5e}{0.2990766}
\SourcePage{78}
其中 $Z_k=\Hom_{H(\A)}(Y_k,X_k)$。反复应用\stref{推论}{3.19}，得到合成的正合三角
\SourceEquationDiagram{diagram-5-12}{0.30946121}{5.12}
其中对象 $Y$ 通过一系列正合三角由 $Z_k\otimes Y_k$ 构造而成，因此属于 $\A$ 中由 $Y_k$ 生成的三角化 $\Ai$ 子范畴。现在假设 $[i_1]\circ\cdots\circ[i_m]=0$。\znote{70}{在三角范畴中，正合三角 $X\to Y\to Z\to TX$ 分裂（即 $Y\cong X\oplus Z$）当且仅当中段复合 $X\to Y\to Z$ 为零。这里 $[i_1]\circ\cdots\circ[i_m]$ 正是合成三角 (5.12) 中从 $X$ 出发的那一段，故它为零时 (5.12) 分裂，得到 (5.13)；其中 $[-1]$ 来自第三项与锥的平移约定。[来源：nLab "triangulated category"，https:/\allowbreak /\allowbreak ncatlab.org/\allowbreak nlab/\allowbreak show/\allowbreak triangulated+\allowbreak category（2026-11-15 访问）；并见 Seidel, Fukaya categories and Picard–Lefschetz theory, 笔记 3.18。]}那么\eref{5.12}按\stref{笔记}{3.18}所述方式分裂，故
\begin{equation}\label{eq:5.13}
X\oplus T_{Y_1}\cdots T_{Y_m}(X)[-1]\cong Y.
\end{equation}
显然，这个观察在具体情形中是否有用，取决于能否证明 $[i_k]$ 的复合为零。下面给出两个特别简单的情形：
\begin{corollary}\label{st:5.7}
设 $\A$ 是 c-有限、三角化的 $\Ai$ 范畴。假设存在 $Y_1,\ldots,Y_m$，使对所有对象 $X$，
\begin{equation}\label{eq:5.14}
T_{Y_1}\cdots T_{Y_m}X\cong0
\end{equation}
都是 $H^0(\A)$ 的零对象（等价地，$T_{Y_1}\cdots T_{Y_m}X$ 在该范畴中的自同态群为零）。那么 $Y_k$ 生成 $\A$，而且这个范畴分裂闭。
\end{corollary}
\begin{proof}
第一项断言直接由\eref{5.13}得到。对于第二项，考虑分裂闭包 $\Pi\A$；它显然仍 c-有限，并由\stref{引理}{4.8}三角化。由扭转的定义可见，$T_{Y_k}$ 保持直和，因此\eref{5.14}对 $\Pi\A$ 的所有对象 $X$ 都成立。于是 $Y_k$ 也生成 $\Pi\A$，这意味着嵌入 $\A\hookrightarrow\Pi\A$ 是拟等价。
\end{proof}
\begin{corollary}\label{st:5.8}
设 $\A$ 是 c-有限、三角化且分裂闭的 $\Ai$ 范畴。假设存在 $Y_1,\ldots,Y_m$ 与 $\sigma\ne0$，使得对所有 $X$，在 $H^0(\A)$ 中有
\[
T_{Y_1}\cdots T_{Y_m}(X)\cong S^{\sigma}X.
\]
那么 $Y_k$ 分裂生成 $\A$。
\end{corollary}
\begin{proof}
有限维假设意味着，对每个 $X$，存在 $s>0$，使 $\Hom^{s\sigma}_{H(\A)}(X,X)=0$。将给定对象族 $Y_k$ 与自身连接 $s$ 次。在所得正合三角\eref{5.12}中，态射
\[
X\longrightarrow(T_{Y_1}\cdots T_{Y_m})^s(X)\cong S^{s\sigma}X
\]
必然为零，因此由\eref{5.13}，$X$ 由 $Y_k$ 分裂生成。
\end{proof}

\SourcePage{79}
\subsection{扭曲复形}\label{sec:5f}
设 $\B$ 是严格含幺的 $\Ai$ 范畴，其链层次的 $\hom$ 空间有限维，并取 $\A=\Tw\B$。在此情形下，扭转函子有显式代表。具体地，给定 $Y_0,Y_1\in\Ob\A$，取张量积 $\hom_{\A}(Y_0,Y_1)\otimes Y_0$，其中如\secref{3o}所述，将第一个因子视为边界算子为 $\partial=\mu^1_{\A}$ 的链复形。由定义，
\begin{equation}\label{eq:5.15}
\hom_{\A}(\hom_{\A}(Y_0,Y_1)\otimes Y_0,Y_1)=\hom_{\A}(Y_0,Y_1)^{\vee}\otimes\hom_{\A}(Y_0,Y_1).
\end{equation}
这个态射群上的 $\mu^1_{\A}$ 算子为 $(c^{\vee}\otimes b)\mapsto(-1)^{|b|}\partial^{\vee}(c^{\vee})\otimes b+c^{\vee}\otimes\partial(b)$，其中 $\partial^{\vee}$ 为通常意义下的对偶微分。特别地，若把\eref{5.15}与 $\End_{\K}(\hom_{\A}(Y_0,Y_1))$ 等同，使 $c^{\vee}\otimes b$ 对应于自同态 $a\mapsto\langle c^{\vee},a\rangle b$，那么恒等自同态就是一个上闭链。将\eref{5.15}中相应的元素记为 $f$，并取\secref{3p}定义的映射锥 $\Cone(f)$。在 Yoneda 嵌入下，$\hom_{\A}(Y_0,Y_1)\otimes Y_0$ 成为 $\hom_{\A}(Y_0,Y_1)\otimes\Y_0$；$\mathcal I^1(f)$ 是\eref{5.2}定义的抽象求值映射 $\phi$；因此 $\mathcal I(\Cone(f))=\mathcal T_{Y_0}(\Y_1)$。\znote{71}{这里把锥 $\Cone(f)$ 显式写成扭曲复形，是扭曲复形一般定义的直接应用：对象为 $\hom_{\A}(Y_0,Y_1)\otimes Y_0$ 与 $Y_1$ 的直和，前者按基 $b^i$ 分解成若干 $Y_0$ 的平移，微分为 (5.18)，其下三角形状来自 $Y_0$ 与 $Y_1$ 的次序。一般框架中两对象扭曲复形的条件 $dq_{ij}+\sum_k q_{kj}q_{ik}=0$ 在这里恰好就是 $f$ 与 $\mu^1_{\A}$ 的相容性。[来源：nLab "twisted complex"，https:/\allowbreak /\allowbreak ncatlab.org/\allowbreak nlab/\allowbreak show/\allowbreak twisted+\allowbreak complex（2026-11-15 访问）。]}这就是说，在这种 $\Ai$ 范畴 $\A$ 中，可以取
\begin{equation}\label{eq:5.16}
T_{Y_0}Y_1=\Cone(f).
\end{equation}
\begin{note}\label{st:5.9}
把这一点写得更显式一些可能很有帮助，至少在 $Y_0,Y_1$ 是 $\B$ 的对象这一简单情形下如此。选取 $\hom_{\B}(Y_0,Y_1)$ 的一组基 $\{b^i\}$，指标为 $i\in I$，并写作 $\mu^1_{\B}(b^i)=\sum_j\beta^{ji}b^j$。那么 $T_{Y_0}(Y_1)$ 是扭曲复形
\begin{equation}\label{eq:5.17}
Y=\bigoplus_{i\in I}S^{-|b^i|+1}Y_0\oplus Y_1,
\end{equation}
其微分 $\delta_Y=(\phi^{ji}\otimes a^{ji})$ 中，$\phi$ 是 $\K$ 的各个平移之间的恒等映射，而
\begin{equation}\label{eq:5.18}
a^{ji}=\begin{cases}
(-1)^{\sigma^i+1}\beta^{ji}e_{Y_0},&i,j\in I,\\
(-1)^{\sigma^i}b^i,&i\in I,\ j=*,\\
0,&\text{其余情形}.
\end{cases}
\end{equation}
这里 $\sigma^i=|b^i|-1$，$*$ 是为\eref{5.17}中的 $Y_1$ 直和项添加的行、列所用的指标（或者，按照\stref{笔记}{3.26}修改等同\eref{3.29}后，\eref{5.18}两行中的符号 $(-1)^{\sigma^i}$ 都消失，分别剩下 $-\beta^{ji}e_{Y_0}$ 与 $b^i$）。
\end{note}
典范代表\eref{5.16}的存在具有许多实际优点。例如，可以直接写出公式来定义 $\Ai$ 函子 $T_{Y_0}$，而无须采用\secref{5b}中的间接构造。
\SourcePage{80}
\begin{note}\label{st:5.10}
设 $\widetilde{\A}$ 是 c-有限、三角化的 $\Ai$ 范畴。应用扰动引理，将其替换为 $\mu^1=0$ 的拟同构范畴，后者在链层次的 $\hom$ 空间因此有限维。再利用\stref{引理}{2.1}使之严格含幺。将结果记为 $\B$，并令 $\A=\Tw\B$。由构造，$\B$ 与 $\widetilde{\A}$ 拟同构，因此也是三角化的 $\Ai$ 范畴；这又意味着 $\A$ 与 $\widetilde{\A}$ 拟等价。利用这一点，常常可以把 c-有限、三角化 $\Ai$ 范畴的一般情形，归结为刚才讨论的技术上更简单的情形。
\end{note}

\subsection{伴随扭转}\label{sec:5g}
扭转运算有一个对偶，更准确地说是伴随运算。原则上它同样重要，但本书很少用到，因此这里只给出定义。取 $\Ai$ 范畴 $\A$ 的对象 $Y_0,Y_1$，如下定义 $\Ai$ 模 $\C^{\vee}$：
\begin{definition*}
\begin{equation}\label{eq:5.19}
\begin{aligned}
\C^{\vee}(X)&=\hom_{\A}(X,Y_1)\oplus\Hom_{\K}(\hom_{\A}(Y_1,Y_0),\hom_{\A}(X,Y_0))[-1],\\
\mu^1_{\C^{\vee}}(c,\beta)&=\bigl(\mu^1_{\A}(c),\ \mu^1_{\A}\circ\beta+(-1)^{|\beta|-1}\beta\circ\mu^1_{\A}+\mu^2_{\A}(\cdot,c)\bigr),\\
\mu^d_{\C^{\vee}}((c,\beta),a_{d-1},\ldots,a_1)
&=\bigl(\mu^d_{\A}(c,a_{d-1},\ldots,a_1),\\[-2pt]
&\qquad\mu^d_{\A}(\beta(\cdot),a_{d-1},\ldots,a_1)+\mu^{d+1}_{\A}(\cdot,c,a_{d-1},\ldots,a_1)\bigr).
\end{aligned}
\end{equation}
任意一个拟表示 $\C^{\vee}$ 的对象，都称为\emph{沿 $Y_0$ 对 $Y_1$ 作伴随扭转}的结果，记为 $T_{Y_0}^{\vee}Y_1$。
\end{definition*}
现在假设 $\A$ c-有限且三角化。那么，对所有 $Y_0,Y_1$，这样的拟表示对象都存在。它位于正合三角
\SourceEquationDiagram{diagram-5-20}{0.43615338}{5.20}
之中，其中的“对偶求值映射” $[f^{\vee}]$ 同样是所有可能的态射 $Y_1\to Y_0$ 的组合。

\subsection{球面对象}\label{sec:5h}
设 $\A$ 是 c-有限、三角化的 $\Ai$ 范畴。沿用最初由 Kontsevich 提出的想法，作如下定义：
\begin{definition*}
若 $Y_0\in\Ob\A$ 满足以下条件，就称它为维数 $n\in\Z$ 的\emph{球面对象}。
\begin{itemize}
\item （Poincaré 对偶性）存在一个线性映射 $\int:\Hom^n_{H(\A)}(Y_0,Y_0)\to\K$，使对所有 $Y_1$，配对
\[
\Hom^{n-\bullet}_{H(\A)}(Y_1,Y_0)\otimes\Hom^{\bullet}_{H(\A)}(Y_0,Y_1)
\xrightarrow{\text{复合}}\Hom^n_{H(\A)}(Y_0,Y_0)\xrightarrow{\int}\K
\]
非退化。
\SourcePage{81}
\item （球面自同态环）$Y_0$ 的自同态空间满足下列短正合列
\[
0\longrightarrow\K\xrightarrow{[e_{Y_0}]}\Hom_{H(\A)}(Y_0,Y_0)\xrightarrow{\int}\K\longrightarrow0.
\]
\end{itemize}
\end{definition*}
若 $n\ne0$，最后一个条件等价于 $\Hom_{H(\A)}(Y_0,Y_0)\cong\K[t]/t^2$，其中 $t$ 的次数为 $n$。当 $n>0$ 时，这恰好是 $S^n$ 的通常上同调环，这也解释了所用术语。后面将看到，Lagrangian 球面（以及同调球面）在 Fukaya 范畴中成为球面对象。\znote{73}{“球面对象”的术语来自拓扑：当维数 $n>0$ 时，自同态环 $\K[t]/t^2$（$\deg t=n$）与球面 $S^n$ 的上同调环 $\HH^*(S^n)$ 一致，而 Poincaré 对偶条件则模仿了上同调中的相交配对。在 Fukaya 范畴中，精确 Lagrangian 球面 $\mathbb{S}^n\subset\mathbb{C}^{n+1}$ 的 Floer 上同调同构于 $\HH^*(S^n)$，配对由 Lagrangian 相交数给出，因此它确实成为球面对象。[来源：译者按（上下文推断，无外部依据）：由本节定义与 $\HH^*(S^n)$ 的经典计算推得。]}
\begin{lemma}\label{st:5.11}
若 $Y_0$ 是球面对象，则 $T_{Y_0}$ 是从 $\A$ 到自身的拟等价。
\end{lemma}
\begin{proof}
给定 $\A$ 的两个对象 $Y_k$（$k=1,2$），令 $C$ 为映射 $\mathcal T_{Y_0}^1:\hom_{\Q}(\Y_1,\Y_2)\to\hom_{\Q}(\mathcal T_{Y_0}\Y_1,\mathcal T_{Y_0}\Y_2)$ 的映射锥（按通常链复形的意义）。它带有三步递降滤过 $F^{\bullet}C$，其相伴分次空间为
\[
\begin{aligned}
C/F^1C&=\hom_{\Q}(\Y_1,\hom_{\A}(Y_0,Y_2)\otimes\Y_0)[1],\\
F^1C/F^2C&=\hom_{\Q}(\Y_1,\Y_2)[1]\oplus\hom_{\Q}(\Y_1,\Y_2)\\[-2pt]
&\qquad\oplus\hom_{\Q}(\hom_{\A}(Y_0,Y_1)\otimes\Y_0,\hom_{\A}(Y_0,Y_2)\otimes\Y_0),\\
F^2C&=\hom_{\Q}(\hom_{\A}(Y_0,Y_1)\otimes\Y_0,\Y_2)[-1].
\end{aligned}
\]
由于 $F^1C/F^2C$ 层次上的抵消，所得谱序列具有以下非零项：
\[
E_1^{rs}=\begin{cases}
(\Hom_{H(\A)}(Y_0,Y_2)\otimes\Hom_{H(\A)}(Y_1,Y_0))^{r+s+1},&r=0,\\
\bigl(\Hom_{H(\A)}(Y_0,Y_2)\otimes\Hom_{H(\A)}(Y_0,Y_0)\\
\qquad\otimes\Hom_{H(\A)}(Y_0,Y_1)^{\vee}\bigr)^{r+s},&r=1,\\
(\Hom_{H(\A)}(Y_0,Y_2)\otimes\Hom_{H(\A)}(Y_0,Y_1)^{\vee})^{r+s-1},&r=2.
\end{cases}
\]
$E_1$ 页的微分 $\partial^{rs}$ 由 $[\phi_k]$ 诱导，也就是说，它由 $H(\A)$ 中的复合映射或其对偶组成。例如，$\partial^{1\bullet}$ 满射，因为
\SourceDiagram{diagram-5-11a}{0.55661479}
\SourcePage{82}
是恒等映射。类似地，$\partial^{0\bullet}$ 单射，因为与 Poincaré 映射复合后得到同构：
\SourceDiagram{diagram-5-11b}{0.55661479}
再作维数计算，就可见谱序列的 $E_2$ 页为零。由构造，这意味着 $T_{Y_0}$ 上同调全忠实。\znote{74}{证明的策略是：把 $T_{Y_0}$ 在 $\hom$ 空间上的作用看成对映射锥 $C$ 作三步滤过，$E_1$ 页由三项组成，分别涉及 $\Hom(Y_0,Y_2)$、$\Hom(Y_0,Y_2)\otimes\Hom(Y_0,Y_0)\otimes\Hom(Y_0,Y_1)^{\vee}$ 与 $\Hom(Y_0,Y_2)\otimes\Hom(Y_0,Y_1)^{\vee}$。Poincaré 对偶使 $E_1$ 页的两个微分一满一单，配合有限维性作维数计数即得 $E_2=0$；等价地，$T_{Y_0}$ 在同调层次上全忠实。[来源：并见 Seidel, Fukaya categories and Picard–Lefschetz theory, §5h（引理 5.11 的证明）。]}

还需证明它在对象的同构类上满射。用 $\B\subset\A$ 表示由同构于某个 $T_{Y_0}(Y_1)$ 的对象组成的全 $\Ai$ 子范畴。由\eref{5.4}与 $Y_0$ 的球面性可得
\[
T_{Y_0}(S^{n-1}Y_0)\cong Y_0,
\]
因此 $Y_0$ 属于 $\B$。利用这一点，再用\eref{5.4}，可见 $\B$ 生成 $\A$。另一方面，\stref{推论}{3.9}表明，$\B$ 本身是三角化的 $\Ai$ 范畴，故包含函子 $\B\hookrightarrow\A$ 必为拟等价。
\end{proof}

\begin{note}\label{st:5.12}
事实上，在 $H^0(\operatorname{fun}(\A,\A))$ 中的同构意义下，伴随扭转 $T_{Y_0}^{\vee}$ 是 $T_{Y_0}$ 的逆。这里不证明这一点（经典导出范畴背景下的平行定理，见\bcite{134}命题 2.10），但要说明一个部分结果，即在 $H^0(\A)$ 中有
\begin{equation}\label{eq:5.21}
T_{Y_0}T_{Y_0}^{\vee}(Y_1)\cong Y_1.
\end{equation}
为此，转到模范畴 $\Q=\modcat(\A)$，取\eref{5.19}中的 $\C^{\vee}$。$\C^{\vee}(Y_0)$ 的上同调位于长正合列
\begin{equation}\label{eq:5.22}
\begin{gathered}
\vdots\\
\Hom_{H(\A)}(Y_0,Y_1)[-1]\\
\big\downarrow\\
\Hom_{H(\A)}(Y_1,Y_0)^{\vee}\otimes\Hom_{H(\A)}(Y_0,Y_0)[-1]\\
\big\downarrow\\
H(\C^{\vee}(Y_0))\\
\vdots
\end{gathered}
\end{equation}
之中，其中第一个同态又是一个对偶复合映射。由 Poincaré 对偶性，经过投影 $\int:\Hom_{H(\A)}(Y_0,Y_0)\to\K$ 后，它成为同构。利用
\SourcePage{83}
$Y_0$ 的自同态环的球面性，可知\eref{5.22}中第二个映射在子空间 $\Hom_{H(\A)}(Y_1,Y_0)^{\vee}\otimes\K\cdot[e_{Y_0}]$ 上的限制为同构。因此，将 $\mathcal T_{Y_0}\C^{\vee}$ 构造为抽象求值映射的锥时，可以把 $\C^{\vee}(Y_0)$ 替换为拟同构的子复形 $\hom_{\A}(Y_1,Y_0)^{\vee}=\Hom_{\K}(\hom_{\A}(Y_1,Y_0),\K\cdot e_{Y_0})$。替换后得到的 $\Ai$ 模为
\[
\begin{aligned}
\M(X)&=(\hom_{\A}(Y_1,Y_0)^{\vee}\otimes\hom_{\A}(X,Y_0))\oplus\hom_{\A}(X,Y_1)\\[-2pt]
&\qquad\oplus\Hom_{\K}(\hom_{\A}(Y_1,Y_0),\hom_{\A}(X,Y_0))[-1],\\
\mu^1_{\M}(\gamma\otimes b,c,\beta)
&=\bigl((-1)^{|b|+|\gamma|}(\gamma\circ\mu^1_{\A})\otimes b+\gamma\otimes\mu^1_{\A}(b),\ \mu^1_{\A}(c),\\[-2pt]
&\qquad\mu^1_{\A}\circ\beta+(-1)^{|\beta|-1}\beta\circ\mu^1_{\A}\\[-2pt]
&\qquad+\mu^2_{\A}(\cdot,c)+\mu^2_{\A}(e_{Y_0},b)\otimes\gamma(\cdot)\bigr),\\
\mu^d_{\M}((\gamma\otimes b,c,\beta),a_{d-1},\ldots,a_1)
&=\bigl(\gamma\otimes\mu^d_{\A}(b,a_{d-1},\ldots,a_1),\ \mu^d_{\A}(c,a_{d-1},\ldots,a_1),\\[-2pt]
&\qquad\mu^d_{\A}(\beta(\cdot),a_{d-1},\ldots,a_1)+\mu^{d+1}_{\A}(\cdot,c,a_{d-1},\ldots,a_1)\\[-2pt]
&\qquad+\mu^{d+1}_{\A}(e_{Y_0},b,a_{d-1},\ldots,a_1)\otimes\gamma(\cdot)\bigr).
\end{aligned}
\]
存在模同态 $\M\to\Y_1$，其一阶项为 $(\gamma\otimes b,c,\beta)\mapsto(-1)^{|c|}c$，高阶项为零。它的核是 $\M(X)$ 中满足 $c=0$ 的子复形；利用 $e_{Y_0}$ 在上同调层次上是恒等态射这一事实，容易看出该核无上同调。因此，在 $H^0(\Q)$ 中有同构 $\mathcal T_{Y_0}\C^{\vee}\cong\M\cong\Y_1$。利用定义回到 $H^0(\A)$，就得到\eref{5.21}。
\end{note}

\subsection{例外集}\label{sec:5i}
在下面的整个讨论中，$\A$ 始终是 c-有限、三角化的 $\Ai$ 范畴。
\begin{definition*}
若 $\A$ 的有限有序对象族 $Y=(Y_1,\ldots,Y_m)$ 满足以下两个条件，就称它为\emph{例外集}。
\begin{itemize}
\item 每个 $Y_k$ 都是\emph{例外对象}；按定义，这意味着其自同态空间由恒等态射张成：
\[
\Hom_{H(\A)}(Y_k,Y_k)=\K\cdot[e_{Y_k}].
\]
\item 在上同调层次上，不同 $Y_k$ 之间的态射具有方向性，即
\[
\Hom_{H(\A)}(Y_k,Y_l)=0\quad(k>l).
\]
\end{itemize}
若此外 $Y_k$ 在\secref{3j}意义下生成 $\A$，就称 $Y$ 为\emph{满例外集}。\znote{76}{“例外对象”一词来自 Bondal 的工作：在三角范畴中，若对象 $E$ 满足 $\Hom(E,E)=\K$ 且 $\Hom(E,E[i])=0$（$i\ne0$），就称 $E$ 例外；两两例外、且态射只有单向的一族例外对象称为例外集，若它还生成整个范畴，就称为满的。对满例外集，任意对象都可沿这族对象作有限阶梯式的 Postnikov 分解，这正是引理 5.13 与 (5.23) 的内容。[来源：nLab "exceptional collection"，https:/\allowbreak /\allowbreak ncatlab.org/\allowbreak nlab/\allowbreak show/\allowbreak exceptional+\allowbreak collection（2026-11-15 访问）。]}
\end{definition*}
\begin{lemma}\label{st:5.13}
设 $Y$ 为满例外集。那么，对任意 $X$，在 $H^0(\A)$ 中有同构
\[
T_{Y_1}\cdots T_{Y_m}(X)\cong0.
\]
\end{lemma}
\SourcePage{84}
\begin{proof}
如\secref{5e}，令 $X_m=X$、$X_{k-1}=T_{Y_k}(X_k)$。目标是证明，对所有 $X$ 都有 $X_0\cong0$。先考虑 $X=Y_l$ 的情形。只要 $k>l$，就有 $\Hom_{H(\A)}(Y_k,X)=0$，因此由\eref{5.4}，$X_{k-1}\cong X_k\cong\cdots\cong X$。下一步中，三角化简为
\SourceDiagram{diagram-5-13}{0.27830739}
因此 $X_{l-1}\cong\Cone(e_{Y_l})\cong0$。此后，所有较低的 $X_k$（$k<l$）也立即可见为零。

$\F=T_{Y_1}\cdots T_{Y_m}$ 是 $\A$ 到自身的 $\Ai$ 函子，因此由\stref{推论}{3.9}保持正合三角。这意味着，若两个对象 $X_0,X_1$ 的像 $\F(X_0),\F(X_1)$ 均为零，那么它们之间任意态射的映射锥 $X_2=\Cone(X_0\to X_1)$ 的像也为零。此外，$\F$ 在同构意义下也与平移相容，这更为显然。由于 $Y_k$ 生成 $\A$，所需结论随之成立。
\end{proof}
这个论证实际上给出了任意对象 $X$ 到 $Y_k$ 的若干平移的显式 \emph{Postnikov 分解}（在一般三角范畴中的例外集框架下，这是众所周知的，参见\bcite{66}§2）。具体地，仍令 $Z_k=\Hom_{H(\A)}(Y_k,X_k)$，有如下阶梯图，其左侧为正合三角，右侧为交换三角：
\SourceEquationDiagram{diagram-5-23}{0.28246123}{5.23}
\SourcePage{85}
\begin{note}\label{st:5.14}
由\stref{推论}{5.7}可知，若 $\A$ 有满例外集，则它分裂闭。事实上，还有稍强的结论：若 $\Q^{\mathrm{fin}}\subset\Q=\modcat(\A)$ 为如下全子范畴，其对象是对每个 $X\in\Ob\A$ 都有 $H(\M(X))$ 有限维的 $\Ai$ 模 $\M$，那么 Yoneda 嵌入给出拟等价 $\A\hookrightarrow\Q^{\mathrm{fin}}$。经典框架下的相应定理是\bcite{27}第 530 页的推论：若 $\mathcal T$ 是具有满例外集的三角范畴，则从 $\mathcal T$ 到有限维 $\K$ 向量空间范畴的任意上同调函子都可表示（受三角范畴理论的局限，这里的证明比我们所处理情形的证明困难得多）。
\end{note}

\subsection{变异}\label{sec:5j}
可以通过以下变换修改给定例外集 $Y=(Y_1,\ldots,Y_m)$：
\begin{itemize}
\item 将每个 $Y_k$ 替换为在 $H^0(\A)$ 中与之同构的对象 $\widetilde Y_k$，形成所得对象族 $\widetilde Y$。
\item 给定 $(\sigma_1,\ldots,\sigma_m)\in\Z^m$，取 $S^{\sigma}Y=(S^{\sigma_1}Y_1,\ldots,S^{\sigma_m}Y_m)$，其中 $S$ 为 $\A$ 上的平移函子。
\item 最后，第 $k$ 个\emph{基本变异}（$1\leq k<m$）为运算 $Y\mapsto U=\mathscr L_kY$，其中 $j\ne k,k+1$ 时 $U_j=Y_j$，其余两个对象为
\begin{equation}\label{eq:5.24}
U_k=T_{Y_k}Y_{k+1},\qquad U_{k+1}=Y_k.
\end{equation}
\end{itemize}
暂且在例外集的同构类层次上讨论，即把第一类变换视为平凡变换。那么，第二类与第三类变换都可逆。第三类变换的逆为 $Y\mapsto V=\mathscr L_k^{-1}Y$，其中仍在 $j\ne k,k+1$ 时取 $V_j=Y_j$，而
\[
V_k=Y_{k+1},\qquad V_{k+1}=T_{Y_{k+1}}^{\vee}Y_k.
\]
此外，还有辫关系
\[
\begin{aligned}
\mathscr L_j\mathscr L_kY&\cong\mathscr L_k\mathscr L_jY&&(|j-k|\geq2),\\
\mathscr L_j\mathscr L_{j+1}\mathscr L_jY&\cong\mathscr L_{j+1}\mathscr L_j\mathscr L_{j+1}Y.
\end{aligned}
\]
由此得到群 $\Z^m\rtimes\operatorname{Br}_m$ 在长度为 $m$ 的例外集的同构类集合上的作用。子群 $\Z^m$ 通过平移作用；辫群 $\operatorname{Br}_m$ 的标准生成元 $\beta_1,\ldots,\beta_{m-1}$ 通过基本变异作用；半直积则通过对称群 $\operatorname{Sym}_m$ 的置换表示形成，这里把它视为 $\operatorname{Br}_m$ 的商群。此外，若一个例外集是满的，那么其 $\Z^m\rtimes\operatorname{Br}_m$ 轨道上的所有例外集也都是满的。\znote{78}{平移给出 $\Z^m$，基本变异 $\mathscr L_k$ 给出辫群 $\operatorname{Br}_m$ 的标准生成元，上两式正是辫关系的两种情形（相距较远的生成元交换、相邻生成元满足辫关系）。这套作用在 Bondal 关于螺旋与 Koszul 代数的工作中已经出现：螺旋由反复作变异得到，而变异不改变导出范畴；稍后 §5n 会把这一点重述为有向 $\Ai$ 范畴上的作用。[来源：nLab "exceptional collection"，https:/\allowbreak /\allowbreak ncatlab.org/\allowbreak nlab/\allowbreak show/\allowbreak exceptional+\allowbreak collection（2026-11-15 访问）。]}
\SourcePage{86}
证明只需对正合三角作简单论证，可以在三角范畴框架中完成，而无须涉及 $\Ai$ 结构；例如见\bcite{66}命题 2.4.2。需要强调，$\mathscr L_k$ 并非函子：它们作用于 $\A$ 内的对象族，却不作用于这些对象之间的态射。

\subsection{Koszul 对偶集}\label{sec:5k}
设 $Y=(Y_1,\ldots,Y_m)$ 是 $\A$ 中的满例外集。对于 $k\leq l$，记
\[
Y_{k,l}=T_{Y_k}T_{Y_{k+1}}\cdots T_{Y_{l-1}}(Y_l),
\]
最简单的情形为 $Y_{k,k}=Y_k$。
\begin{definition*}
\emph{右 Koszul 对偶集} $Y^!$ 定义为由对象 $Y^!_{m+1-k}=Y_{1,k}$（$1\leq k\leq m$）组成的对象族。
\end{definition*}
很容易看出，它仍是满例外集，并由原来的 $Y$ 沿“整体半扭转”
\begin{equation}\label{eq:5.25}
\Delta^{1/2}=\beta_{m-1}(\beta_{m-2}\beta_{m-1})\cdots(\beta_2\cdots\beta_{m-1})(\beta_1\cdots\beta_{m-1})\in\operatorname{Br}_m
\end{equation}
变异而得（辫群中的复合从右到左进行，因此\eref{5.25}对应的变异从 $\mathscr L_{m-1}$ 开始，接着依次施加 $\mathscr L_{m-2}$ 直到 $\mathscr L_1$，如此继续）。复合有关正合三角\eref{5.4}中的映射 $[i]$，便对所有 $j<k\leq l$ 得到 $H^0(\A)$ 中的典范态射
\[
[i_{j,k,l}]:Y_{k,l}\longrightarrow Y_{k-1,l}\longrightarrow\cdots\longrightarrow Y_{j,l}.
\]
令 $[i_{k,k,l}]$ 为恒等态射，就将其延伸到 $j=k$ 的情形。特别地，对每个 $k$，有典范态射 $[i_{1,k,k}]:Y_k\to Y_{1,k}=Y^!_{m+1-k}$。现在讨论 Koszul 对偶集的一些性质。这些性质基本上都来自三角\eref{5.4}诱导的长正合列。为简洁起见，用 $\Hom$ 代替 $\Hom_{H(\A)}$；同时省略通常括在上同调层次态射外面的方括号，即以 $i_{j,k,l}$ 表示 $[i_{j,k,l}]$。
\begin{lemma}\label{st:5.15}
\textup{(i)} 当 $k=l$ 时，$\Hom(Y_l,Y^!_{m+1-k})$ 为一维，由 $i_{1,k,k}$ 生成；当 $k\ne l$ 时为零。

\textup{(ii)} $\Hom(X,Y^!_{m+1-k})\cong\Hom(Y_k,X_k)^{\vee}$，其中如\eref{5.23}，$X_k=T_{Y_{k+1}}\cdots T_{Y_m}(X)$。

\textup{(iii)} $\Hom(Y^!_{m-k},Y^!_{m+1-k})\cong\Hom(Y_k,Y_{k+1})^{\vee}[-1]$。
\end{lemma}
\begin{proof}
省略直接可得的 (i)。为由它推出更一般的 (ii)，从两端遍历分解\eref{5.23}，并施加 $\Hom(-,Y^!_{m+1-k})$，得到
\[
\begin{aligned}
\Hom(X,Y^!_{m+1-k})&\cong\Hom(X_{m-1},Y^!_{m+1-k})\cong\cdots\cong\Hom(X_k,Y^!_{m+1-k}),\\
0&\cong\Hom(X_0,Y^!_{m+1-k})\cong\cdots\cong\Hom(X_{k-1},Y^!_{m+1-k}).
\end{aligned}
\]
\SourcePage{87}
余下的正合三角联系 $X_k$ 与 $X_{k-1}$，由此得到
\[
\Hom(X_k,Y^!_{m+1-k})\cong\Hom(\Hom(Y_k,X_k)\otimes Y_k,Y^!_{m+1-k})\cong\Hom(Y_k,X_k)^{\vee}.
\]
对于 (iii)，在 (ii) 中取 $X=Y^!_{m-k}$。此时 $X_k$ 是映射 $i_{1,k+1,k+1}:Y_{k+1}\to Y^!_{m-k}$ 的锥，因此 $\Hom(Y_k,X_k)\cong\Hom(Y_k,Y_{k+1})[1]$。
\end{proof}
\begin{lemma}\label{st:5.16}
取\eref{5.23}中的一次态射 $Z_k\otimes Y_k\to Z_{k+1}\otimes Y_{k+1}$，利用上面 \textup{(ii)} 的同构，将其视为映射
\begin{equation}\label{eq:5.26}
\Hom(X,Y^!_{m+1-k})^{\vee}\otimes Y_k\longrightarrow\Hom(X,Y^!_{m-k})^{\vee}\otimes Y_{k+1}.
\end{equation}
利用 \textup{(iii)}，它可以表述如下：取 $\Hom(Y_k,Y_{k+1})$ 与 $\Hom(Y^!_{m-k},Y^!_{m+1-k})[1]$ 的对偶基 $c_l,c_l^{\vee}$。对每个 $l$，先取与 $c_l^{\vee}$ 作乘法的线性映射 $\Hom(X,Y^!_{m-k})\to\Hom(X,Y^!_{m+1-k})$，将其对偶化，再与 $c_l$ 组合，就得到\eref{5.26}两侧空间之间的映射。最后对所有 $l$ 求和，即得所需结果。\znote{79}{用对偶基 $c_l,c_l^{\vee}$ 描述映射，是把两件事配对：$u_1\cdot c_l$ 把 $\Hom(X,Y^!_{m-k})$ 的元素与 $c_l\in\Hom(Y_k,Y_{k+1})$ 复合，而 $c_l^{\vee}\cdot u_0$ 把 $c_l^{\vee}$ 作用到 $\Hom(X,Y^!_{m+1-k})$ 上。两项互为转置，故求和后得到 (5.26) 两侧空间之间的恰当映射；谱序列微分 (5.29) 中的符号 $(-1)^{|c_l|(|u_0|+|u_1|)}$ 则来自把 $c_l^{\vee}$ 移过 $u_0$ 时交换两个奇次元素所生的 Koszul 符号。[来源：并见 Seidel, Fukaya categories and Picard–Lefschetz theory, §5k–5l（引理 5.15、5.16）。]}
\end{lemma}
\begin{proof}[证明概要]
先假设 $X$ 只是某个态射 $Z_k\otimes Y_k\to Z_{k+1}\otimes Y_{k+1}$ 的锥。在这一情形下，前一引理的所有同构都容易写出，可以通过显式计算验证结论（可以使用三角范畴的标准工具抽象计算，也可以用扭曲复形计算）。现在考虑一般对象 $X$。取典范映射 $X\to X_{k-1}$，将其补成正合三角 $X'\to X\to X_{k-1}$。再应用标准正合三角，可见 $X'\to X$ 在与 $X$ 相应的\eref{5.26}两侧空间及与 $X'$ 相应的对应空间之间诱导同构，且与两侧之间的映射交换。类似地，若令 $X''=X'_{k+1}$，则映射 $X'\to X''$ 也在\eref{5.26}上诱导同构。因此，只需对 $X''$ 证明交换性。但它满足 $T_{Y_k}T_{Y_{k+1}}X''=X'_{k-1}=0$，因此是上述类型的锥。
\end{proof}

\subsection{Beilinson 型谱序列}\label{sec:5l}
设对象 $X$ 具有 Postnikov 分解\eref{5.23}。令 $K$ 为上同调函子，即从 $H(\A)$ 到分次 $\K$ 向量空间的反变函子，要求它与有限直和及平移相容，并将正合三角映为长正合列；精确定义可见\bcite{61}、\bcite{143}。存在收敛到 $K(X)$ 的谱序列，其起始项为
\begin{equation}\label{eq:5.27}
E_1^{rs}=K^{r+s}(Z_{m-r}\otimes Y_{m-r})\cong(Z_{m-r}^{\vee}\otimes K(Y_{m-r}))^{r+s}.
\end{equation}
微分 $\partial_1^{r\bullet}$ 是\eref{5.23}中态射 $Z_{m-r-1}\otimes Y_{m-r-1}\to Z_{m-r}\otimes Y_{m-r}$ 在 $K$ 下的像。为建立这一谱序列，可以采用
\SourcePage{88}
Atiyah–Hirzebruch 谱序列的标准定义\bcite{138}（第 15 章），分别用 Postnikov 分解与上同调函子代替胞腔分解与广义同调理论；也可以采用正合偶的表述\bcite{61}（第 263 页）。将这一一般构造用于当前情形，得到：
\begin{proposition}\label{st:5.17}
对 $H(\A)$ 的任意一对对象 $U_0,U_1$，存在收敛到 $\Hom(U_0,U_1)$ 的谱序列，其起始项为
\begin{equation}\label{eq:5.28}
E_1^{rs}=(\Hom(U_0,Y^!_{r+1})\otimes\Hom(Y_{m-r},U_1))^{r+s}.
\end{equation}
按\stref{引理}{5.16}选取对偶基，其中 $k=m-r-1$。那么微分 $\partial^{r\bullet}:E_1^{r\bullet}\to E_1^{r+1,\bullet}$ 为
\begin{equation}\label{eq:5.29}
\partial(u_0\otimes u_1)=\sum_l(-1)^{|c_l|(|u_0|+|u_1|)}(c_l^{\vee}\cdot u_0)\otimes(u_1\cdot c_l).
\end{equation}
\end{proposition}
\begin{proof}
取 $X=U_0$、$K=\Hom(-,U_1)$。利用同构 $Z_{m-r}^{\vee}=\Hom(Y_{m-r},X_{m-r})^{\vee}\cong\Hom(X,Y^!_{r+1})$，起始项\eref{5.27}就成为\eref{5.28}，而\stref{引理}{5.16}给出公式\eref{5.29}（严格地说，符号问题还有些含糊，因为我们并未完整说明谱序列如何建立；不过此处并不需要完全精确）。\znote{80}{这一构造是把拓扑中的 Atiyah–Hirzebruch 谱序列移植到三角范畴：把空间换成对象 $X$、把胞腔分解换成 Postnikov 分解 (5.23)、把广义同调换成上同调函子 $K$，$E_1$ 项就是各层“胞腔”的上同调，微分由粘合映射诱导。经典情形是 Beilinson 对射影空间上结构层与对角线的分解，它给出 $\mathbb{P}^n$ 上相干层上同调的谱序列。[来源：并见 Seidel, Fukaya categories and Picard–Lefschetz theory, §5l 与笔记 5.18 所引文献。]}
\end{proof}
\begin{note}\label{st:5.18}
谱序列\eref{5.28}最初的例子来自 Beilinson 对射影空间上对角线的分解，见\bcite{19}或\bcite{104}第 245 页。它在\bcite{67}中被推广到凝聚层的例外集；从这里推广到一般三角范畴相当直接。完整论述见\bcite{66}，特别是其中的 2.7.3 与 3.3.2 节（注意，由于使用左 Koszul 对偶而非右对偶，他们的谱序列 $E_1$ 项与我们的略有不同）。与 Koszul 对偶性的联系在\bcite{26}、\bcite{25}中有解释；另见下面的\secref{5o}。
\end{note}

\clearpage
\subsection{有向\texorpdfstring{$\Ai$}{A∞} 范畴}\label{sec:5m}
到目前为止，我们采用了外在的观点，在一个固定的三角化 $\Ai$ 范畴内部研究例外集及其变异。另一种办法是只考虑一类特殊的 $\Ai$ 范畴：它们的对象本身在导出范畴中构成例外集。在这种内在的背景下，变异就是将这样的范畴变成另一个具有等价导出范畴的范畴的操作。
\begin{definition*}
\emph{有向 $\Ai$ 范畴}是一个具有 $m$ 个对象的严格含幺 $\Ai$ 范畴 $\A^{\to}$，连同这些对象的一个排序 $(Y_1,\ldots,Y_m)$，使
\[
\hom_{\A^{\to}}(Y_i,Y_j)=\begin{cases}
0,&i>j,\\\K\cdot e_{Y_i},&i=j,\\\K\text{ 上的有限维空间},&i<j.
\end{cases}
\]

\end{definition*}
\SourcePage{89}
由于严格含幺，只有下列复合映射不是由单位元条件直接确定的：
\begin{equation}\label{eq:5.30}
\begin{aligned}
\mu^d_{\A^{\to}}:\hom_{\A^{\to}}(Y_{i_{d-1}},Y_{i_d})\otimes\cdots\otimes\hom_{\A^{\to}}(Y_{i_0},Y_{i_1})\\
\longrightarrow\hom_{\A^{\to}}(Y_{i_0},Y_{i_d})[2-d],\qquad i_0<\cdots<i_d.
\end{aligned}
\end{equation}
所以 $\A^{\to}$ 可以由有限多个 $\K$ 中的元素完全描述，即\eref{5.30}的矩阵系数。令 $\A$ 为 $\A^{\to}$ 的三角化包络。由定义，$Y_k$ 在上同调全忠实函子 $\A^{\to}\to\A$ 下的像构成一个例外集。由构造，$\A$ 由 $\A^{\to}$ 的对象生成，因此这个例外集是满的。
\begin{note}\label{st:5.19}
更具体地，取 $\A=\Tw\A^{\to}$。有向性使扭曲复形的定义在技术上可以作若干小的简化。首先，任意这样的复形都可以写成
\begin{equation}\label{eq:5.31}
\left(X=\bigoplus_{i=1}^mV^i\otimes Y_i,\quad\delta_X=(\delta_X^{ji})\right),
\end{equation}
其中 $V^1,\ldots,V^m$ 为有限维分次向量空间，而 $\delta_X^{ji}\in\Hom_{\K}(V^i,V^j)\otimes\hom_{\A}(Y_i,Y_j)$ 为一次元素。对角项形如 $\delta_X^{ii}=\partial_{V^i}\otimes e_{Y_i}$，其中 $(V^i,\partial_{V^i})$ 是向量空间的链复形。

其次，$\delta_X$ 的下三角条件自动满足，因此在这一特殊情形下，可以从扭曲复形的定义中删去。具体地，可以先用 $\bigoplus_{i\geq k}V^i\otimes Y_i$ 对\eref{5.31}进行滤过，再用首个链复形 $(V^k,\partial_{V^k})$ 的次数滤过加以细化。这样得到 $X$ 的一个滤过，相对于它，$\delta_X$ 满足所需的下三角条件。

最后，任意扭曲复形 $X$ 都在 $H^0(\A)$ 中同构于某个微分的对角项全为零的 $\widetilde X$，即 $\delta_{\widetilde X}^{ii}=0$。事实上，Postnikov 分解\eref{5.23}说明，在同构意义下，可以从 $X_0=0$ 出发，通过一系列映射锥
\begin{equation}\label{eq:5.32}
X_k\cong\Cone(S^{-1}X_{k-1}\longrightarrow Z_k\otimes Y_k)
\end{equation}
重建 $X$。假设已有与 $X_{k-1}$ 同构的对象 $\widetilde X_{k-1}$，它是只包含 $Y_1,\ldots,Y_{k-1}$ 直和项、且微分对角项为零的扭曲复形。令 $\widetilde X_k$ 为与\eref{5.32}相应的态射 $S^{-1}\widetilde X_{k-1}\to Z_k\otimes Y_k$ 的锥。它同构于 $X_k$，并具有与 $\widetilde X_{k-1}$ 相同的性质；反复进行这一构造，最终就得到所需的 $\widetilde X=\widetilde X_m$。
\end{note}

\SourcePage{90}
\subsection{有向子范畴}\label{sec:5n}
暂时只考虑所有 $\hom$ 空间都有限维的 $\Ai$ 范畴，这可以简化一些技术细节。设 $Y=(Y_1,\ldots,Y_m)$ 是这样的范畴 $\A$ 中的对象族，未必是例外集。
\begin{definition*}
相应的\emph{有向 $\Ai$ 子范畴} $\A^{\to}=\A^{\to}(Y)$ 定义为如下唯一的有向 $\Ai$ 范畴：其对象为 $Y_k$，在 $i<j$ 时 $\hom_{\A^{\to}}(Y_i,Y_j)=\hom_{\A}(Y_i,Y_j)$，而非平凡复合映射\eref{5.30}与 $\A$ 中的对应映射相同。
\end{definition*}
这一构造具有函子性：设 $\F:\A\to\B$ 为 $\Ai$ 函子，令 $\B^{\to}$ 为与像对象族 $\F(Y)$ 相应的有向 $\Ai$ 子范畴。则存在唯一的严格含幺 $\Ai$ 函子 $\F^{\to}:\A^{\to}\to\B^{\to}$，其在对象的严格递增链上的限制等于 $\F$。若 $\F$ 上同调全忠实，则 $\F^{\to}$ 为拟同构。

需要提醒一点：若 $\A$ 严格含幺，就可以把 $\A^{\to}$ 与 $\A$ 中由 $i<j$ 时的 $\hom_{\A}(Y_i,Y_j)$ 以及各 $k$ 对应的子空间 $\K\cdot e_{Y_k}$ 组成的子范畴等同。不过，一般的 c-含幺情形下，$\A^{\to}$ 严格说来并非子范畴。对此有如下补救：
\begin{lemma}\label{st:5.20}
设 $Y$ 为 $\A$ 中的对象族，$\A^{\to}$ 为相应的有向 $\Ai$ 子范畴。那么存在在对象上为恒等映射的 $\Ai$ 函子 $\mathcal H:\A^{\to}\to\A$，对所有 $i<j$ 诱导同构
\[
\Hom_{H(\A^{\to})}(Y_i,Y_j)\cong\Hom_{H(\A)}(Y_i,Y_j).
\]
\end{lemma}
\begin{proof}
先取最小模型，再用\stref{引理}{2.1}将其严格化，得到严格含幺的 $\widetilde{\A}$ 及在对象上为恒等映射的拟等价 $\Phi:\A\to\widetilde{\A}$。令 $\widetilde{\A}^{\to}$ 为与同一对象族 $Y$ 相应的有向子范畴。严格含幺性给出包含函子 $\widetilde{\mathcal H}:\widetilde{\A}^{\to}\to\widetilde{\A}$，同时有诱导拟同构 $\Phi^{\to}:\A^{\to}\to\widetilde{\A}^{\to}$。取在对象上为恒等映射的逆向拟等价 $\Psi:\widetilde{\A}\to\A$，令 $\mathcal H=\Psi\circ\widetilde{\mathcal H}\circ\Phi^{\to}$ 即可。
\end{proof}
\begin{lemma}\label{st:5.21}
设 $Y,Y'$ 是 $\A$ 中的两个对象族，并且对每个 $j$，$Y_j$ 与 $Y'_j$ 在 $H^0(\A)$ 中同构。那么相应的有向子范畴 $\A^{\to}(Y)$、$\A^{\to}(Y')$ 拟同构，且拟同构可由保持对象顺序的严格含幺 $\Ai$ 函子给出。
\end{lemma}
\begin{proof}
令 $\B\subset\A$ 为只包含对象 $Y_j$ 的全 $\Ai$ 子范畴，$\B'$ 为只包含对象 $Y'_j$ 的相应子范畴，$\C$ 为同时包含两类对象的子范畴。由假设，$\B$ 与 $\B'$ 到 $\C$ 的包含都是拟等价。因此，利用\stref{定理}{2.9}可以构造拟等价 $\B\to\B'$。由构造很容易看出，它将每个 $Y_j$ 映到 $Y'_j$，故实际上是拟同构。再转到有向子范畴即可。
\end{proof}
\begin{lemma}\label{st:5.22}
设 $Y$ 是三角化 $\Ai$ 范畴 $\A$ 中的满例外集，$\A^{\to}$ 为相应的有向 $\Ai$ 子范畴，$\mathcal H:\A^{\to}\to\A$ 为\stref{引理}{5.20}中的
\SourcePage{91}
$\Ai$ 函子。那么存在拟等价 $\G:\Tw\A^{\to}\to\A$，其在 $\A^{\to}\subset\Tw\A^{\to}$ 上的限制，在 $H^0(\operatorname{fun}(\A^{\to},\A))$ 中同构于 $\mathcal H$。特别地，$D(\A^{\to})$ 与 $H^0(\A)$ 作为三角范畴等价。
\end{lemma}
\begin{proof}
$Y$ 是例外集这一事实意味着 $\mathcal H$ 上同调全忠实。由于 $\A$ 三角化，且由 $Y$ 的对象生成，它实际上就是 $\A^{\to}$ 的三角化包络。$\Tw\A^{\to}$ 也一样；现在应用\stref{引理}{3.20}即可。
\end{proof}
一个应用如下。取有向 $\Ai$ 范畴 $\A^{\to}$，以及 $\A=\Tw\A^{\to}$ 中由 $\A^{\to}$ 本身的对象构成的满例外集。按\secref{5j}施加变异来变换这个例外集，令 $\B^{\to}$ 为 $\A$ 中与变异后对象族相应的有向 $\Ai$ 子范畴。\stref{引理}{5.22}给出拟等价 $\Tw\B^{\to}\to\Tw\A^{\to}$，从而给出三角范畴的等价 $D(\B^{\to})\cong D(\A^{\to})$。正如 Kontsevich 最初观察到的，这定义了群 $\Z^m\rtimes\operatorname{Br}_m$ 在具有 $m$ 个对象的有向 $\Ai$ 范畴的拟同构类集合上的作用，同一轨道中的所有范畴都具有等价的导出范畴（实际上，中心元素 $(1,\ldots,1)\in\Z^m\subset\Z^m\rtimes\operatorname{Br}_m$ 作用平凡，因此这一作用下降到商群 $\Z^{m-1}\rtimes\operatorname{Br}_m$）。

我们还需要\stref{引理}{5.22}对非例外对象族的一个版本：
\begin{lemma}\label{st:5.23}
设 $Y$ 是 $\A$ 中的对象族，每个对象都是球面的。在 $\Tw\A$ 中，考虑由与\eref{5.24}相同的公式定义的对象族 $U=\mathscr L_kY$。用 $\B^{\to},\C^{\to}$ 表示与这两个对象族相应的有向 $\Ai$ 子范畴。那么 $\Tw\B^{\to}$ 与 $\Tw\C^{\to}$ 拟等价。
\end{lemma}
\begin{proof}
根据\stref{引理}{5.6}，可以把 $\Ai$ 范畴 $\A$ 替换为任意拟等价范畴，并把 $Y,U$ 替换为新范畴中的对应对象族。因此，可以假设 $\A$ 严格含幺且 $\mu^1_{\A}=0$；这纯粹是为了方便。

稍微区分一下记号，用 $\widetilde Y_k$ 表示 $\B^{\to}$ 中对应于 $\A$ 中 $Y_k$ 的对象。这些对象在 $\B=\Tw\B^{\to}$ 中构成满例外集 $\widetilde Y$。令 $\widetilde U=\mathscr L_k\widetilde Y$ 为这个例外集的第 $k$ 个基本变异。包含函子 $\mathcal H:\B^{\to}\to\A$ 诱导 $\Ai$ 函子 $\Tw\mathcal H:\B\to\Tw\A$，满足
\begin{equation}\label{eq:5.33}
\Tw\mathcal H(T_{\widetilde Y_i}\widetilde Y_j)=T_{Y_i}Y_j\quad(i<j).
\end{equation}
左边的扭转在 $\B$ 内进行，右边是在 $\Tw\A$ 中的相应构造。两边都采用\secref{5f}的显式公式，于是等式归结为 $\hom_{\B^{\to}}(\widetilde Y_i,\widetilde Y_j)=\hom_{\A}(Y_i,Y_j)$。由定义，\eref{5.33}尤其意味着 $\Tw\mathcal H(\widetilde U)=U$。下一步是证明映射
\begin{equation}\label{eq:5.34}
(\Tw\mathcal H)^1:\hom_{\B}(\widetilde U_i,\widetilde U_j)\longrightarrow\hom_{\Tw\A}(U_i,U_j)\quad(i<j)
\end{equation}
\SourcePage{92}
为拟同构。只有 $i$ 或 $j$ 等于 $k$ 时，这一点才非平凡。大多数情形可由以下观察处理：一边由 $p<q$ 时的 $\hom_{\B^{\to}}(\widetilde Y_p,\widetilde Y_q)$ 构成，另一边由相应的 $\hom_{\A}(Y_p,Y_q)$ 构成；依 $\B^{\to}$ 的定义，它们相同。唯一剩下的是 $i=k,j=k+1$ 的情形，此时\eref{5.34}中的空间为
\[
\begin{aligned}
\hom_{\B}(T_{\widetilde Y_k}\widetilde Y_{k+1},\widetilde Y_k)&=\hom_{\A}(Y_k,Y_{k+1})^{\vee}[-1],\\
\hom_{\Tw\A}(T_{Y_k}Y_{k+1},Y_k)&=\hom_{\A}(Y_{k+1},Y_k)\\[-2pt]
&\qquad\oplus\hom_{\A}(Y_k,Y_{k+1})^{\vee}[-1]\otimes\hom_{\A}(Y_k,Y_k).
\end{aligned}
\]
由于 $\mu^1_{\A}=0$，第一个空间的微分为零；第二个空间上的微分则是将复合部分对偶化所得的映射
\[
(\mu^2_{\A})':\hom_{\A}(Y_{k+1},Y_k)\longrightarrow\hom_{\A}(Y_k,Y_{k+1})^{\vee}\otimes\hom_{\A}(Y_k,Y_k).
\]
球面性意味着
\[
(\id\otimes\textstyle\int)\circ(\mu^2_{\A})':\hom_{\A}(Y_{k+1},Y_k)\longrightarrow\hom_{\A}(Y_k,Y_{k+1})^{\vee}[-n_k]
\]
是同构，其中 $n_k$ 为球面对象 $Y_k$ 的维数。因此 $(\mu^2_{\A})'$ 单射；再作维数计算，可知其像是与 $\ker(\id\otimes\int)=\hom_{\A}(Y_k,Y_{k+1})^{\vee}\otimes\K e_{Y_k}$ 互补的子空间。而\eref{5.34}中的映射恰好为 $(\Tw\mathcal H)^1=(0,\id\otimes e_{Y_k})$，因此是拟同构。

令 $\widetilde{\C}^{\to}$ 为 $\B$ 中与对象族 $\widetilde U$ 相应的有向 $\Ai$ 子范畴。它自然嵌入 $\B$；将此嵌入与 $\Tw\mathcal H$ 组合，得到 $\Ai$ 函子 $\G:\widetilde{\C}^{\to}\to\Tw\A$。这个函子实际上严格含幺，其像落在有向 $\Ai$ 子范畴 $\C^{\to}\subset\Tw\A$ 中；由上面的计算，它给出 $\widetilde{\C}^{\to}$ 与 $\C^{\to}$ 之间的拟同构。因此，$\Tw\C^{\to}$ 与 $\Tw\widetilde{\C}^{\to}$ 拟等价；而根据对变异的一般讨论，后一个范畴又与 $\Tw\B^{\to}$ 拟等价。
\end{proof}
要将有向子范畴的构造推广到 $\A$ 仅 c-有限的情形，可以这样做：利用扰动引理，转到一个拟同构的 $\Ai$ 范畴 $\widetilde{\A}$，使实际的 $\hom_{\widetilde{\A}}$ 空间有限维，再对这个范畴应用上述构造。将结果仍记为 $\A^{\to}$；它仍带有\stref{引理}{5.20}中的 $\Ai$ 函子 $\mathcal H:\A^{\to}\to\A$。缺点是 $\A^{\to}$ 不再典范，因为它依赖于 $\widetilde{\A}$ 的选择；但至少在拟同构意义下唯一，这对于大多数用途已经足够。特别地，\stref{引理}{5.22}与\stref{引理}{5.23}仍然成立。

\clearpage
\subsection{再论 Koszul 对偶集}\label{sec:5o}
设 $\A^{\to}$ 是有向 $\Ai$ 范畴，$\A=\Tw\A^{\to}$。回顾\stref{引理}{3.36}，存在典范的上同调全忠实 $\Ai$ 函子 $\widetilde{\mathcal I}:\A\to\Q=\modcat(\A^{\to})$。由定义，$Y^!_{m+1-k}\in\Ob\A$ 在此函子下的像是 $\Ai$ 模 $\Y^!_{m+1-k}=\mathcal T_{Y_1}\cdots\mathcal T_{Y_{k-1}}(\Y_k)$。
\SourcePage{93}
这样的模的具体形式很容易写出。基础分次向量空间为
\begin{equation}\label{eq:5.35}
\begin{aligned}
\Y^!_{m+1-k}(Y_l)=\bigoplus_r\bigoplus_{i_1,\ldots,i_r}
\bigl(&\hom_{\A}(Y_{i_r},Y_k)\otimes\hom_{\A}(Y_{i_{r-1}},Y_{i_r})\otimes\cdots\\[-2pt]
&\otimes\hom_{\A}(Y_{i_1},Y_{i_2})\otimes\hom_{\A}(Y_l,Y_{i_1})\bigr)[r],
\end{aligned}
\end{equation}
其中对 $r\geq0$、$l\leq i_1<i_2<\cdots<i_r<k$ 求和；特别地，$r=0$ 的直和项就是 $\hom_{\A}(Y_l,Y_k)$。其次，
\begin{equation}\label{eq:5.36}
\begin{aligned}
&\mu^1_{\Y^!_{m+1-k}}(b_{r+1}\otimes\cdots\otimes b_1)\\
&\quad=\sum_{p,n}(-1)^{\dagger}b_{r+1}\otimes\cdots\otimes b_{n+p+1}
\otimes\mu^p_{\A}(b_{n+p},\ldots,b_{n+1})\otimes b_n\otimes\cdots\otimes b_1,
\end{aligned}
\end{equation}
其中 $\dagger=|b_1|+\cdots+|b_n|-n$；对于 $d\geq2$，
\[
\begin{aligned}
&\mu^d_{\Y^!_{m+1-k}}(b_{r+1}\otimes\cdots\otimes b_1,a_{d-1},\ldots,a_1)\\
&\quad=\sum_{n\geq1}b_{r+1}\otimes\cdots\otimes b_{n+1}\otimes\mu^{d-1+n}_{\A}(b_n,\ldots,b_1,a_{d-1},\ldots,a_1).
\end{aligned}
\]
\begin{definition*}
现在考虑\emph{简单 $\Ai$ 模} $\mathit{SS}_k$：令 $\mathit{SS}_k(Y_k)=\K$，对所有其他 $l$ 令 $\mathit{SS}_k(Y_l)=0$，并取 $\mu^2_{\mathit{SS}_k}(\cdot,e_{Y_k})=\Id_{\K}$，其余结构映射均为零。
\end{definition*}
存在典范模同态
\[
t:\Y^!_{m+1-k}\longrightarrow\mathit{SS}_k.
\]
一阶部分 $t^1$ 为从 $\Y^!_{m+1-k}(Y_k)=\hom_{\A}(Y_k,Y_k)=\K\cdot e_{Y_k}$ 到 $\mathit{SS}_k(Y_k)$ 的恒等映射；高阶映射 $t^d$（$d\geq2$）恒为零。\znote{85}{简单模 $\mathit{SS}_k$ 是有向 $\Ai$ 范畴 $\A^{\to}$ 的模范畴中最小的对象：它只在 $Y_k$ 处取值 $\K$，其余为零，且 $\mu^2(\cdot,e_{Y_k})$ 为恒等。引理 5.24 说 Koszul 对偶对象 $Y^!_{m+1-k}$ 的 Yoneda 模同构于 $\mathit{SS}_k$，因此 Koszul 对偶集恰好给出一族简单模；这也是后文把 $\Q$ 与“每个对象都由简单模经迭代锥得到”这一结构联系起来的出发点。[来源：并见 Seidel, Fukaya categories and Picard–Lefschetz theory, §5o（引理 5.24）与 §5p。]}
\begin{lemma}\label{st:5.24}
$[t]$ 是 $H^0(\Q)$ 中的同构。
\end{lemma}
\begin{proof}
复形\eref{5.35}具有显然的递增长度滤过 $F^r$。将 $r$ 换成 $-r$，使其变为递降滤过，考察相应的谱序列，它收敛到 $H(\Y^!_{m+1-k}(Y_l))$。起始项为
\[
E_1^{-r,s}=\left(\bigoplus_{i_1,\ldots,i_r}\Hom_{H(\A)}(Y_{i_r},Y_k)\otimes\cdots\otimes\Hom_{H(\A)}(Y_l,Y_{i_1})\right)^s,
\]
其中 $i_k$ 如\eref{5.35}，微分为
\[
\partial^{-r,s}(b_{r+1}\otimes\cdots\otimes b_1)=\sum_n(-1)^{\dagger+|b_{n+1}|}b_{r+1}\otimes\cdots\otimes b_{n+2}b_{n+1}\otimes\cdots\otimes b_1,
\]
\SourcePage{94}
其中 $\dagger$ 如\eref{5.36}。这正是一个约化条形分解；更准确地说，若对所有 $l$ 取直和，所得就是分次 $H(\A)$ 模的 Abel 范畴 $\operatorname{Mod}^0(H(\A))$ 中，简单对象 $H(\mathit{SS}_k)$ 的标准投射分解。因此，当 $k=l$ 时，谱序列的 $E_2$ 项为一维，由 $[e_{Y_k}]$ 张成；其余情形为零。与 $\mathit{SS}_k(Y_l)$ 的平凡滤过作比较很直接，结论是对任意 $l$，$t^1:\Y^!_{m+1-k}(Y_l)\to\mathit{SS}_k(Y_l)$ 都在上同调上诱导同构。再由\stref{引理}{1.16}，得到所需结果。
\end{proof}
\begin{note}\label{st:5.25}
概述一下这个观点如何给出谱序列\eref{5.28}的另一种构造。任意 $\mathcal U\in\Ob\Q$ 都带有自然的子模递降滤过
\[
(F^r\mathcal U)(Y_k)=\begin{cases}\mathcal U(Y_k),&k\leq m-r,\\0,&\text{其余情形}.\end{cases}
\]
分次部分为 $F^r\mathcal U/F^{r+1}\mathcal U\cong\mathcal U(Y_{m-r})\otimes\mathit{SS}_{m-r}$，其中将 $\mathcal U(Y_{m-r})$ 视为向量空间的链复形。现在取 $\mathcal U=\mathcal U_1$，考虑 $\hom_{\Q}(\mathcal U_0,\mathcal U_1)$ 上的诱导滤过。所得谱序列的起始项为
\begin{equation}\label{eq:5.37}
E_1^{rs}=\bigl(H(\hom_{\Q}(\mathcal U_0,\mathit{SS}_{m-r}))\otimes H(\mathcal U_1(Y_{m-r}))\bigr)^{r+s}.
\end{equation}
现在假设 $\mathcal U_0,\mathcal U_1$ 是扭曲复形在上同调全忠实嵌入 $\widetilde{\mathcal I}:\A\to\Q$ 下的像。利用\stref{引理}{2.12}与\stref{引理}{5.24}，就可以将\eref{5.37}化为\eref{5.28}的形式。
\end{note}

\subsection{含幺\texorpdfstring{$\Ai$}{A∞} 模}\label{sec:5p}
令 $\mathcal P=\modcat(\A^{\to})^{\mathrm{fin}}$ 为有限维、严格含幺 $\Ai$ 模组成的 $\Ai$ 范畴。有限维是指 $\M(Y_k)$ 是总维数有限的分次向量空间；严格含幺是指 $\mu^2_{\M}(b,e_{Y_k})=b$，并且在 $d\geq3$ 时，只要某个 $a_j$ 是单位元，就有 $\mu^d_{\M}(b,a_{d-1},\ldots,a_1)=0$；此外，只允许具有类似消失性质的预模同态。由于有向性，后一个条件意味着态射空间 $\hom_{\mathcal P}(\M_0,\M_1)$ 有限维。Yoneda 嵌入 $\mathcal I$ 及其在扭曲复形上的延伸 $\widetilde{\mathcal I}$，其像自然落在 $\mathcal P$ 中。另一方面，嵌入 $\mathcal P\to\Q$ 上同调全忠实；考察长度滤过所给谱序列的 $E_1$ 项即可看出：对 $\Q$ 与 $\mathcal P$，它们分别是\eref{2.14}中用于计算模 $\operatorname{Ext}$ 群的标准条形复形的非约化版本与约化版本。
\begin{corollary}\label{st:5.26}
$\mathcal P$ 是 $\A^{\to}$ 的三角化包络。
\end{corollary}
\SourcePage{95}
\begin{proof}
由上述观察，$\mathcal I:\A^{\to}\to\mathcal P$ 是上同调全忠实的 $\Ai$ 函子。此外，$\mathcal P$ 对\eref{3.32}意义下的映射锥封闭，故用与\stref{引理}{3.35}相同的论证，可知它三角化。还需证明 $\mathcal I$ 的像生成 $\mathcal P$，等价地，证明 $\widetilde{\mathcal I}:\Tw(\A^{\to})\to\mathcal P$ 是拟等价。但由\stref{引理}{5.24}，简单模位于 $\widetilde{\mathcal I}$ 的像中；另一方面，任意模都容易表示为由简单模形成的迭代锥。
\end{proof}
\begin{corollary}\label{st:5.27}
任意有向 $\Ai$ 范畴都拟同构于另一个仍然有向、而且实际上为 dg 范畴的范畴，即 $d>2$ 时 $\mu^d=0$。
\end{corollary}
\begin{proof}
$\mathcal I(Y_k)$ 在 $\mathcal P$ 中构成例外集。取相应的有向 $\Ai$ 子范畴；由于 $\mathcal P$ 严格含幺且 $\hom$ 空间有限维，这一步很直接。高阶复合映射为零是 $\Ai$ 模范畴的一般性质。
\end{proof}
当然，我们已经知道，任意 $\Ai$ 范畴都可以变换为一个拟同构的 dg 范畴（\stref{推论}{2.14}）。上述结果的要点在于：有向情形下，可以在保持 $\hom$ 空间有限维的同时做到这一点；在更一般的背景下，这是不可能的。\znote{87}{关键在于：有向 $\Ai$ 范畴的模范畴 $\mathcal P$ 既严格含幺、又只有有限维 $\hom$ 空间，故在 $\mathcal P$ 内部作形式消去（把高阶结构映射逐步杀掉）不会破坏有限维性。对一般的 $\Ai$ 范畴，把 $\mu^d$（$d>2$）消去通常要引入无穷维的 $\hom$ 空间，因此推论 5.27 的有向性假设不能去掉。[来源：并见 Seidel, Fukaya categories and Picard–Lefschetz theory, §5p（推论 5.27）与推论 2.14。]}

\section{\texorpdfstring{$\Z/2$}{Z/2} 作用}\label{sec:6}
\subsection{基本概念}\label{sec:6a}
本节讨论带有最简单意义下的对合的 $\Ai$ 范畴。
\begin{definition*}
$\Ai$ 范畴 $\A$ 上的\emph{朴素 $\Z/2$ 作用}是一个 $\Ai$ 函子 $\epsilon_{\A}:\A\to\A$，其阶数不小于二的项为零，并满足 $\epsilon_{\A}\circ\epsilon_{\A}=\Id_{\A}$。
\end{definition*}
事实上，它只由对象上的对合与相应的态射映射 $\epsilon_{\A}:\hom_{\A}(X_0,X_1)\to\hom_{\A}(\epsilon_{\A}X_0,\epsilon_{\A}X_1)$ 组成，满足
\[
\mu^d_{\A}(\epsilon_{\A}(a_d),\ldots,\epsilon_{\A}(a_1))=\epsilon_{\A}(\mu^d_{\A}(a_d,\ldots,a_1)).
\]
\begin{definition*}
若 $\epsilon_{\A}X=X$，就称 $X$ 为 $\epsilon_{\A}$-\emph{不动对象}。\emph{不动 $\Ai$ 子范畴} $\A^{\mathrm{fix}}$ 是由这些对象组成的全 $\Ai$ 子范畴。\emph{不变 $\Ai$ 子范畴} $\A^{\mathrm{inv}}$ 具有与 $\A^{\mathrm{fix}}$ 相同的对象，但其态射为不变部分 $\hom_{\A^{\mathrm{inv}}}(X_0,X_1)=\hom_{\A}(X_0,X_1)^{\Z/2}$，复合映射为相应的限制。
\end{definition*}
我们需要取不变部分这一过程具有良好的同调性质，因此从现在起假设 $\operatorname{char}(\K)\ne2$。

后面会多次遇到以下初等问题。令 $A=H^0(\A)$，$\epsilon_A$ 为 $\epsilon_{\A}$ 所诱导的 $\Z/2$ 作用。取两个 $\epsilon_{\A}$-不动对象 $X_0,X_1$，假设它们在 $A$ 中同构。何时可以把这个同构
\SourcePage{96}
选为 $\Z/2$-不变的？障碍是类
\begin{equation}\label{eq:6.1}
[y=\epsilon_A(x^{-1})\cdot x]\in\{y\in\operatorname{Aut}_A(X_0):y^{-1}=\epsilon_A(y)\}/\operatorname{Aut}_A(X_0),
\end{equation}
其中 $x$ 为任意一个从 $X_0$ 到 $X_1$ 的同构，而 $\operatorname{Aut}_A(X_0)=\Hom_A^0(X_0,X_0)^{\times}$ 通过扭曲共轭作用于自身，即 $z$ 将 $y$ 映为 $\epsilon_A(z)^{-1}yz$。\znote{88}{这个障碍类可以化归为纯代数验算：若 $x:X_0\to X_1$ 为同构，令 $y=\epsilon_A(x^{-1})x$，则 $y^{-1}=x^{-1}\epsilon_A(x)=\epsilon_A(y)$，故 $y$ 必落在所示集合中；若换用另一同构 $xz$（$z\in\operatorname{Aut}_A(X_0)$），$y$ 就变成 $\epsilon_A(z)^{-1}yz$。而 $y$ 为单位元恰等价于 $\epsilon_A(x)=x$。因此这个商集退化为单点，当且仅当存在 $\Z/2$-不变的同构；扭曲共轭正是把「换一个 $x$」造成的差别除掉的等价关系。[来源：译者按（上下文推断，无外部依据）：由 $y$ 的定义直接验算 $y^{-1}=\epsilon_A(y)$，并代入 $x\mapsto xz$ 得到扭曲共轭公式。]}
\begin{definition*}
最简单的情形是 $X_0$ \emph{自由}：这意味着 $\Hom_A(X_0,X_0)\cong R\oplus R$ 是同一个分次代数 $R$ 的两个副本的直和，且 $\epsilon_A$ 通过交换两个直和项作用。
\end{definition*}
在这种情形下，\eref{6.1}中的集合由 $y=(r,r^{-1})$ 这样的元素组成，每个元素都可以通过以 $z=(r^{-1},e)$ 作扭曲共轭而化为平凡元素。因此：
\begin{lemma}\label{st:6.1}
设有两个 $\epsilon_{\A}$-不动对象，其中一个已知是自由的。那么，若它们在 $H^0(\A)$ 中同构，它们也在 $H^0(\A^{\mathrm{inv}})$ 中同构。
\end{lemma}

\subsection{等变扭曲复形}\label{sec:6b}
首先需要讨论加性扩张的等变版本 $\Sigma^{\mathrm{eq}}\A$。其对象是 $\Sigma\A$ 的对象，即写成\eref{3.15}中形式和的三元组 $(I,\{X^i\},\{V^i\})$，再加上下列数据：指标集 $I$ 上的对合 $i\mapsto\underline i$，满足 $X^{\underline i}=\epsilon_{\A}(X^i)$；以及向量空间的分次映射 $\epsilon_{V^i}:V^i\to V^{\underline i}$，对所有 $i$ 满足 $\epsilon_{V^{\underline i}}\circ\epsilon_{V^i}=\Id$。态射空间与复合映射的定义与非等变版本相同。每个态射空间带有诱导的对合；用\eref{3.16}的记号，它由以下两个映射的张量积得到：
\[
\epsilon_{V_1^j}\circ-\circ\epsilon_{V_0^{\underline i}}:
\Hom_{\K}(V_0^i,V_1^j)\longrightarrow\Hom_{\K}(V_0^{\underline i},V_1^{\underline j})
\]
以及
\[
\epsilon_{\A}:\hom_{\A}(X_0^i,X_1^j)\longrightarrow\hom_{\A}(X_0^{\underline i},X_1^{\underline j}).
\]
\begin{definition*}
\emph{等变扭曲复形}是二元组 $(X,\delta_X)$，其中 $X\in\Ob\Sigma^{\mathrm{eq}}\A$，微分 $\delta_X$ 属于不变部分 $\hom^1_{\Sigma^{\mathrm{eq}}\A}(X,X)^{\Z/2}$，并满足通常的条件：存在滤过（此时应由等变子复形 $F^iX$ 组成），以及 Maurer–Cartan 方程。
\end{definition*}
态射空间沿用 $\Sigma^{\mathrm{eq}}\A$ 中的定义，复合映射按\eref{3.20}定义，所得为 $\Ai$ 范畴 $\Tw^{\mathrm{eq}}\A$，它带有在对象上平凡的朴素 $\Z/2$ 作用。

最简单的等变扭曲复形是 $\A$ 中的单个 $\epsilon_{\A}$-不动对象 $X$（更精确地说，取固定的单元素集 $I=\{*\}$、一维向量空间 $V^*=\K$，其映射 $\epsilon_{V^*}=\Id$ 为平凡映射，并取 $X^*=X$）。一个稍不平凡的例子是 $\A$ 中任意对象与其像的直和 $X\oplus\epsilon_{\A}X$；微分仍为零，但对合 $i\mapsto\underline i$ 交换两个直和项。注意，第一类对象给出 $\A^{\mathrm{fix}}$ 到 $\Tw^{\mathrm{eq}}\A$ 的全嵌入，与 $\Z/2$ 作用相容。更一般地，$\A^{\mathrm{inv}}$ 中的任意扭曲复形都可以
\SourcePage{97}
视为等变扭曲复形（集 $I$ 与向量空间 $V^i$ 上的对合均平凡）；这样，$\Tw(\A^{\mathrm{inv}})$ 成为 $(\Tw^{\mathrm{eq}}\A)^{\mathrm{inv}}$ 的全 $\Ai$ 子范畴。最后，直接忘掉对合，便定义了全忠实 $\Ai$ 函子 $\Tw^{\mathrm{eq}}\A\to\Tw\A$。

再说明一些熟悉运算的等变版本。直和的定义是显然的。其次，给定带 $\Z/2$ 作用 $\epsilon_Z$ 的有限维分次向量空间 $Z$ 和等变扭曲复形 $X$，可以取张量积 $Z\otimes V^i$，赋予组合的 $\Z/2$ 作用 $\epsilon_Z\otimes\epsilon_{V^i}$，并取微分 $\delta_Y=\id_Z\otimes\delta_X$，从而形成等变张量积 $Y=Z\otimes X$。这包括平移运算 $S$，其定义照常是取带平凡 $\Z/2$ 作用的 $Z=\K[1]$。另一个有趣情形是 $\Z/2$ 的非平凡特征标，即 $Z=\K$、$\epsilon_Z=-\id_Z$。这时张量积只是将所有 $\epsilon_{V^i}$ 的符号反转，把这一过程记为 $X\mapsto\check X$。一般的张量积 $Z\otimes X$ 本质上是 $X$ 与 $\check X$ 的若干平移的有限和，$Z$ 的每个不变或反不变基元素各对应一项。最后，若 $c\in\hom_{\Tw\A}(Y_0,Y_1)$ 是不变的零次上闭链，那么赋予 $SY_0\oplus Y_1$ 微分\eref{3.28}所定义的映射锥，仍是等变扭曲复形。这说明 $(\Tw^{\mathrm{eq}}\A)^{\mathrm{inv}}$ 三角化，而 $\Tw^{\mathrm{eq}}\A$ 本身并不三角化。\znote{89}{成败取决于特征标：等变扭曲复形的每个对象都带有 $\Z/2$ 在其上的作用，可以平凡，也可以按 $-\id$ 作用（后者即 $\check X$）。要在 $SY_0\oplus Y_1$ 上写出对合，必须让两块的符号一致，这只有当 $c$ 落在不变部分时才做得到；对一般的 $c$，锥没有自然的等变结构。反过来，一旦取不变部分，$Y_0,Y_1,c$ 全都不变，锥便自动留在 $(\Tw^{\mathrm{eq}}\A)^{\mathrm{inv}}$ 里——这就是「只有取不变部分之后才三角化」的确切含义。[来源：原书 §6b 本段关于锥的构造与 (3.28) 的锥定义，并见本段关于 $\check X$ 的讨论。]}

\begin{lemma}\label{st:6.2}
设 $Y,Y_0,Y_1\in\Ob\A$ 均为 $\epsilon_{\A}$-不动对象；另假设 $Y$ 自由，且 $\Hom_{H(\A)}(Y_0,Y_1)$ 上的 $\Z/2$ 作用平凡。若在 $D(\A)=H^0(\Tw\A)$ 中，$Y$ 同构于某个态射 $Y_0\to Y_1$ 的锥，那么同样的结论在 $D(\A^{\mathrm{inv}})$ 中成立。
\end{lemma}
\begin{proof}
由假设，$\Hom_{H(\A)}(Y_0,Y_1)$ 是 $\Z/2$-不变的，因此可以在不改变锥的同构类的情况下，假设有关上闭链 $c\in\hom^0_{\A}(Y_0,Y_1)$ 本身 $\Z/2$-不变。这时可以把 $\Cone(c)$ 构造为 $\Tw^{\mathrm{eq}}\A$ 的对象。事实上，它位于嵌入 $\Tw(\A^{\mathrm{inv}})\subset\Tw^{\mathrm{eq}}\A$ 的像中。由于遗忘函子 $H^0(\Tw^{\mathrm{eq}}\A)\to H^0(\Tw\A)$ 全忠实，就在 $H^0(\Tw^{\mathrm{eq}}\A)$ 中得到同构 $Y\cong\Cone(c)$。再由于 $Y$ 自由，利用\stref{引理}{6.1}可以使这个同构不变；这就意味着它来自 $H^0(\Tw(\A^{\mathrm{inv}}))$ 中的同构。
\end{proof}
\begin{lemma}\label{st:6.3}
设 $\A$ c-有限，$X$ 为 $\epsilon_{\A}$-不动对象。假设 $X$ 拟表示某个对象 $Y$ 与其 $\epsilon_{\A}$ 像的直和 $Y\oplus\epsilon_{\A}Y$，并且 $\Hom_{H(\A)}(Y,\epsilon_{\A}Y)=0$。再假设在 $D(\A)$ 中有同构
\begin{equation}\label{eq:6.2}
T_{Y_1}\cdots T_{Y_m}(Y)\cong S(\epsilon_{\A}Y).
\end{equation}
这里 $Y_1,\ldots,Y_m$ 是 $\A$ 的 $\epsilon_{\A}$-不动对象，并且对所有 $j<k$，$\Hom_{H(\A)}(Y_j,Y_k)$ 上的 $\Z/2$ 作用平凡；$S$ 为平移。那么 $X$ 属于 $\Tw(\A^{\mathrm{inv}})$ 中由 $Y_k$ 生成的三角化子范畴。
\end{lemma}
\SourcePage{98}
\begin{proof}
最清楚的情形是 $\A$ 严格含幺且 $\mu^1_{\A}=0$（这尤其意味着 $\hom_{\A}$ 空间本身有限维）。迭代\eref{5.16}，可知\eref{6.2}左边的对象具有形式
\begin{equation}\label{eq:6.3}
\begin{aligned}
C=Y\oplus\bigoplus\bigl(&\hom_{\A}(Y_{j_r},Y)[1]\otimes\hom_{\A}(Y_{j_{r-1}},Y_{j_r})[1]\otimes\cdots\\[-2pt]
&\otimes\hom_{\A}(Y_{j_1},Y_{j_2})[1]\otimes Y_{j_1}\bigr),
\end{aligned}
\end{equation}
其中对所有 $r\geq1$ 与 $j_1<j_2<\cdots<j_r$ 取直和。无须显式写出完整微分 $\delta_C$，只需了解它的一般形式，这很容易由构造推出。首先：
\begin{itemize}
\item $\delta_C$ 严格减小张量积表达式的长度。特别地，它在第一个直和项 $Y$ 上为零。
\end{itemize}
因此，$C$ 可以写为态射 $U\to Y$ 的锥，其中 $(U,\delta_U)\in\Ob\Tw\A$ 包含\eref{6.3}中除首项之外的所有项，整体反向平移一次，并带有微分的相应部分。其次：
\begin{itemize}
\item 只考虑满足 $j_1\geq j$ 的直和项所给的滤过与 $\delta_U$ 相容。此外，相对于这个滤过的对角项，即相伴分次空间上的诱导微分，形如 $\partial_{Z_j}\otimes e_{Y_j}$，其中 $\partial_{Z_j}$ 是重数空间
\[
Z_j=\bigoplus_{\substack{r\geq1\\j=j_1<j_2<\cdots<j_r}}
\bigl(\hom_{\A}(Y_{j_r},Y)\otimes\hom_{\A}(Y_{j_{r-1}},Y_{j_r})\otimes\cdots\otimes\hom_{\A}(Y_j,Y_{j_2})\bigr)[r-1]
\]
上的某个微分；$r=1$ 时该项理解为 $\hom_{\A}(Y_j,Y)$。仔细考察会发现 $\partial_{Z_j}$ 是某种条形微分，但我们不用这一事实。
\end{itemize}
结合 $j<k$ 时 $\hom_{\A}(Y_j,Y_k)$ 完全 $\Z/2$-不变的假设，可知 $\delta_U$ 的每个系数，都是某个向量空间同态与 $\A$ 中某个 $\epsilon_{\A}$-不变态射的张量积。因此，$U$ 属于子范畴 $\Tw(\A^{\mathrm{inv}})\subset\Tw\A$；再看上述描述，可知它实际上属于 $\Tw(\A^{\mathrm{inv}})$ 中由 $Y_k$ 生成的三角化子范畴。根据\eref{6.2}，在 $H(\Tw\A)$ 中有形如
\SourceDiagram{diagram-6-3}{0.24299974}
的正合三角。由假设，在 $H(\A)$ 中，从 $Y$ 到 $\epsilon_{\A}Y$ 没有任何次数的非零态射，因此该正合三角分裂，即在 $H^0(\Tw\A)$ 中有
\begin{equation}\label{eq:6.4}
X\cong Y\oplus\epsilon_{\A}Y\cong U.
\end{equation}
\eref{6.4}中的所有对象都是等变扭曲复形，且中间的对象自由。用\stref{引理}{6.2}证明结尾的论证，便得到 $H^0(\Tw(\A^{\mathrm{inv}}))$ 中所需的同构。

\SourcePage{99}
为将结论推广到一般 $\A$，去掉上面引入的技术限制，作如下处理。选择\stref{笔记}{1.13}中的数据，但要保证复形 $\hom_{\A}(X_0,X_1)$ 上的分裂与同伦和 $\Z/2$ 作用相容。只有 $X_0,X_1$ 都为 $\epsilon_{\A}$-不动对象时，才可能有困难；在这些情形下，只需将复形分成不变与反不变直和项，分别处理后再合起来。于是，\stref{扰动引理}{1.12}之后的公式\eref{1.18}保持原有的对称性，即得到等变拟同构 $\F:\A\to\B$，其中 $\B$ 带有朴素 $\Z/2$ 作用且 $\mu^1_{\B}=0$。接着，像\stref{引理}{2.1}那样使 $\B$ 严格含幺。该引理证明的第一步是平凡的：由于 $\mu^1_{\B}=0$，上同调单位元具有唯一的代表，且自动 $\Z/2$-不变，可以取 $\Phi^2=0$。构造的其余部分不涉及选择，见\eref{2.4}，因而自然与对合相容，给出等变形式微分同胚 $\Phi$，使 $\widetilde{\B}=\Phi_*\B$ 严格含幺。取 $\G=\Phi\circ\F$，考虑 $X,Y,Y_k$ 在其下的像。显然，在 $H^0(\widetilde{\B})$ 中有
\[
\begin{gathered}
\G X\cong\G Y\oplus\G(\epsilon_{\A}Y)\cong\G Y\oplus\epsilon_{\widetilde{\B}}(\G Y),\\
\Hom^*_{H(\widetilde{\B})}(\G Y,\epsilon_{\widetilde{\B}}(\G Y))=0,\\
\Hom^*_{H(\widetilde{\B})}(\G Y_i,\G Y_j)\text{ 在 }i<j\text{ 时为 }\Z/2\text{-不变的}.
\end{gathered}
\]
由\stref{引理}{5.6}，\eref{6.2}的对应式也成立。因此处于上面考虑的特殊情形，即 $\G X$ 属于 $\Tw(\widetilde{\B}^{\mathrm{inv}})$ 中由 $\G Y_k$ 生成的三角化子范畴。但不变部分 $\G^{\mathrm{inv}}:\A^{\mathrm{inv}}\to\widetilde{\B}^{\mathrm{inv}}$ 也是拟同构，因此由\stref{引理}{3.25}诱导导出范畴的等价，所需结论随之成立。
\end{proof}

\subsection{\texorpdfstring{$\Ai$}{A∞} 模}\label{sec:6c}
与上面对扭曲复形的讨论相对应，也可以引入 $\Ai$ 模的等变版本。
\begin{definition*}
$\A$ 上的\emph{等变 $\Ai$ 模}是 $\Ai$ 模 $\M$，连同模同态 $\epsilon_{\M}:\M\to\epsilon_{\A}^*\M$，要求其阶数不小于二的项为零，并且 $\epsilon_{\M}$ 与 $\epsilon_{\A}^*(\epsilon_{\M})$ 以任一顺序复合，均为恒等模自同态。\znote{91}{记号 $\epsilon_{\A}^*\M$ 需要说明：沿对合 $\epsilon_{\A}$ 拉回，把 $\A$-模 $\M$ 变成另一个 $\A$-模，其分量集合仍是 $\M(X)$，但结构映射改为经由 $\M(\epsilon_{\A}X)$ 给出。因此所谓模同态 $\epsilon_{\M}:\M\to\epsilon_{\A}^*\M$，逐分量就是下一句写出的那族映射 $\M(X)\to\M(\epsilon_{\A}X)$，并要求与 $\mu^d_{\M}$ 相容；写成模同态而非逐分量数据，好处是「$\epsilon_{\M}$ 给出等变结构」这句话可以直接用模范畴的语言陈述，不必每次重复相容性公式。[来源：原书 §6c 本段及紧接的分量公式，并见 §6b 中 $\Sigma^{\mathrm{eq}}\A$ 的同类构造。]}
\end{definition*}
具体地，附加数据由一族映射 $\epsilon_{\M}:\M(X)\to\M(\epsilon_{\A}(X))$ 组成，$X$ 与 $\epsilon_{\A}(X)$ 对应的映射互为逆，并且以如下方式与 $\epsilon_{\A}$ 相容：
\[
\mu^d_{\M}(\epsilon_{\M}(b),\epsilon_{\A}(a_{d-1}),\ldots,\epsilon_{\A}(a_1))=\epsilon_{\M}(\mu^d_{\M}(b,a_{d-1},\ldots,a_1)).
\]
态射及其复合按非等变情形定义，将所得 $\Ai$ 范畴记为 $\Q=\modcat^{\mathrm{eq}}(\A)$。它带有诱导的朴素 $\Z/2$
\SourcePage{100}
作用，在对象上平凡；具体地，对 $t\in\hom_{\Q}(\M_0,\M_1)$，令
\[
(\epsilon_{\Q}t)^d(b,a_{d-1},\ldots,a_1)=\epsilon_{\M_1}\bigl(t^d(\epsilon_{\M_0}(b),\epsilon_{\A}(a_{d-1}),\ldots,\epsilon_{\A}(a_1))\bigr).
\]
若 $Y$ 是 $\A$ 的 $\epsilon_{\A}$-不动对象，相应 Yoneda 模 $\Y$ 带有自然的等变结构，即在 $\Y(X)=\hom_{\A}(X,Y)$ 上取 $\epsilon_{\Y}=\epsilon_{\A}$。这定义了上同调全忠实 $\Ai$ 函子 $\mathcal I_{\A}^{\mathrm{eq}}:\A^{\mathrm{fix}}\to\Q$，称为\emph{等变 Yoneda 嵌入}；它与 $\Z/2$ 作用严格交换（利用\stref{引理}{3.36}的等变版本，可以将它延伸为上同调全忠实函子 $\Tw^{\mathrm{eq}}\A\to\Q$）。另一方面，若 $\M$ 为等变 $\Ai$ 模，则只对 $\epsilon_{\A}$-不动对象 $X$ 取空间 $\M(X)$ 的 $\Z/2$-不变部分，就得到 $\A^{\mathrm{inv}}$ 上的 $\Ai$ 模 $\M^{\mathrm{inv}}$。连同对 $\Z/2$-不变自然变换的相应限制，这定义了\emph{不变部分 $\Ai$ 函子} $\Q^{\mathrm{inv}}\to\modcat(\A^{\mathrm{inv}})$。最后，直接忘掉映射 $\epsilon_{\M}$，就得到全忠实 $\Ai$ 函子 $\Q\to\modcat(\A)$。

给定带 $\Z/2$ 作用的向量空间链复形 $Z$ 与等变模 $\M$，张量积 $Z\otimes\M$ 仍然等变。其次，给定两个等变模之间的 $\Z/2$-不变同态，\eref{3.32}意义下的锥仍为等变模。最后，若 $Y$ 是 $\epsilon_{\A}$-不动对象，$\M$ 为等变模，那么扭转模 $\mathcal T_Y\M$ 自然等变，而且在 $H(\Q^{\mathrm{inv}})$ 中有正合三角\eref{5.3}。
\begin{lemma}\label{st:6.4}
设 $Y_0,Y_1,Y_2$ 是 $\epsilon_{\A}$-不动对象，且 $\Hom_{H(\A)}(Y_0,Y_1)$ 上的 $\Z/2$ 作用平凡。假设在 $H^0(\Q)$ 中有
\begin{equation}\label{eq:6.5}
\mathcal T_{Y_0}(\Y_1)\cong\Y_2,
\end{equation}
其中同构属于该范畴态射群的 $\Z/2$-不变部分。那么在 $H^0(\modcat(\A^{\mathrm{inv}}))$ 中有与\eref{6.5}平行的同构，这里的 Yoneda 模与扭转都相对于 $\A^{\mathrm{inv}}$ 而取。最后，若 $\A$ c-有限，则在 $D(\A^{\mathrm{inv}})$ 中有 $T_{Y_0}(Y_1)\cong Y_2$。
\end{lemma}
\begin{proof}
为明确区分各个范畴，将等变 Yoneda $\Ai$ 模记为 $\mathcal I_{\A}^{\mathrm{eq}}(Y)$，将 $\modcat(\A^{\mathrm{inv}})$ 中的对应模记为 $\mathcal I_{\A^{\mathrm{inv}}}(Y)$。应用不变部分函子 $\Q^{\mathrm{inv}}\to\modcat(\A^{\mathrm{inv}})$。它将 $\mathcal I_{\A}^{\mathrm{eq}}(Y_2)$ 映为 $\mathcal I_{\A^{\mathrm{inv}}}(Y_2)$。另一方面，$\mathcal T_{Y_0}(\mathcal I_{\A}^{\mathrm{eq}}(Y_1))$ 的像是 $\A^{\mathrm{inv}}$ 上的 $\Ai$ 模 $\C$，满足
\[
\C(X)=\bigl((\hom_{\A}(Y_0,Y_1)\otimes\hom_{\A}(X,Y_0))[1]\bigr)^{\Z/2}\oplus\hom_{\A}(X,Y_1)^{\Z/2},
\]
并带有通常的结构映射\eref{5.1}。由于 $\hom_{\A}(Y_0,Y_1)$ 的反不变部分无上同调，将第一个直和项替换为 $(\hom_{\A}(Y_0,Y_1)^{\Z/2}\otimes\hom_{\A}(X,Y_0)^{\Z/2})[1]$，就得到一个拟同构子模；由\stref{引理}{1.16}，它在 $H^0(\modcat(\A^{\mathrm{inv}}))$ 中是同构对象。
\SourcePage{101}
但按定义，这个子模恰好就是 $\mathcal T_{Y_0}(\mathcal I_{\A^{\mathrm{inv}}}(Y_1))$。结合\eref{5.6}，最后一项断言立即成立。
\end{proof}
\begin{note}\label{st:6.5}
最后作一些略为宽泛的一般评论。对 $\Ai$ 范畴上群作用采取朴素处理，与同伦代数结构的思想并不一致；但对于有限群以及特征为零、或特征 $p\gg0$ 足够大的基域 $\K$，它仍然相当有效。\secref{10b}会重新考虑一般情形；不过，那里提出的定义虽然比这里的更精细，也至多只是初步近似。\znote{93}{「朴素」在此有确切含义：把 $\Z/2$ 作用实现为范畴自身的严格对合，即 $\epsilon_{\A}\circ\epsilon_{\A}=\Id$ 逐项成立；而真正等变的结构只应在同伦意义下满足这一条，即带有更高阶的同伦相干数据。两者在上链层次上给出的不动对象并不相同，这正是注 6.5 所说「不协调」的来源。本章只在特征为零或特征 $p\gg0$ 时使用朴素版本，因为那时取不变部分是正合函子（Maschke 定理），不会丢失信息。[来源：原书注 6.5 与 §10b、§14b，并见本章 6(a) 关于 $\operatorname{char}(\K)\neq2$ 的假设。]}

我们使用不动对象的方式也有些朴素。在恰当的范畴表述中，应当改为考虑“等变对象”：它们只在指定的同构意义下被群作用固定。其上同调层次的版本将在后面简要出现（\secref{14b}），但定义适当的上链层次概念并不直接。从形式上说，这种方法的优点在于：可以分别把 $\Tw^{\mathrm{eq}}\A$ 与 $\modcat^{\mathrm{eq}}(\A)$ 定义为 $\Tw\A$ 与 $\modcat(\A)$ 上诱导作用的不动对象子范畴。不过，这里采用的专门定义已经足以满足我们的目的。
\end{note}

\chapter{Fukaya 范畴}\label{chapter:II}
\SourcePage{103}
本章的主要目标，是在正合辛流形这一相对简单的特殊情形下，利用 Gromov 与 Floer 开创的经典分析方法，说明 Fukaya $\Ai$ 范畴的定义。可以说，这种方法如今已经过时，由 Fukaya–Oh–Ohta–Ono 的奠基性工作\bcite{60}、\bcite{56}所取代，后者允许在完全一般的情形下定义 Fukaya 范畴。不过，一般定义在分析与代数两方面都相当复杂；对于这里讨论的情形，其中许多精细技术并无必要。

为了方便非专门研究这一领域的读者，先从正合 Lagrangian 子流形的 Floer 上同调群的原始定义\bcite{44}开始。自 20 世纪 90 年代初以来，人们已知道 Floer 群上有一些运算，以适当一类具有边界标记点的 Riemann 曲面为指标，使之成为开弦拓扑量子场论（TQFT；这一想法似乎源于 Donaldson；讨论它或其闭弦对应物的一些论文有\bcite{119}、\bcite{106}、\bcite{135}、\bcite{131}）。在允许复结构变化的 Riemann 曲面族时，就出现了最初在\bcite{52}中提到的 Fukaya $\Ai$ 范畴（实际上，它只使用带边界标记点的圆盘，因此应视为某个更复杂理论的第一部分，而该理论仍有很大部分尚未探索）。这里所谓经典方法，是用零阶非齐次项扰动伪全纯映射方程，使解的模空间变得光滑，从而避免使用虚基本链。像\bcite{116}一样，非齐次项随 Riemann 曲面的复结构变化；我们将花一些篇幅说明精确要求。一旦澄清这一点，基本分析结果便没有特别困难或新颖之处，因此只给出尽量简短的证明概要。我们将更多强调分次与符号问题，因为两者对于应用都很重要。对于分次，采用 Kontsevich 的分次 Lagrangian 子流形形式体系\bcite{90}、\bcite{125}；与 Floer、Viterbo 较早提出的相对分次相比，它保留了稍多的信息。研究 Floer 理论中符号问题的一般框架建立于\bcite{48}，Lagrangian Floer 上同调的情形则在\bcite{60}中处理；但我们仍认为值得重新整理其中一部分理论。\znote{95}{两种分次的差别可以这样看：Floer--Viterbo 的相对分次只记录相位差落在 $\Z$ 里的部分，因此只能对每一对 $(L_0,L_1)$ 给出 $\Z$-分次，且需先任选一个基准；Kontsevich 的分次则给每个 $L$ 附加整体相位提升 $\theta_L:L\to\R$，于是分次是绝对的，平移 $[1]$ 在单个对象上也有定义。代价是这类提升并不总存在，通常要在 $2c_1(M)=0$ 之类的条件下工作；一旦存在，符号与次数的计算都可以逐点进行，这正是本节说它「保留了稍多信息」的意思。[来源：原书 §8 引言提到的 Kontsevich 分次 Lagrangian 形式体系；并见原书第三章关于分次与符号的讨论。]}

本章还讨论带对合的辛流形的等变 Fukaya 范畴，以及曲面的 Fukaya 范畴所具有的部分组合性质。后一主题已经在多种背景下出现：椭圆曲线的同调镜像对称\bcite{90}、\bcite{108}、\bcite{107}，曲面上曲线的 Floer 同调\bcite{135}、\bcite{35}，以及 Legendrian 纽结理论\bcite{30}、\bcite{43}。我们着重讨论将分析方法与组合方法联系起来的正则性论证，因为文献对此论述较少。

\SourcePage{104}
\section{一点辛几何}\label{sec:7}
\subsection{带角的正合辛流形}\label{sec:7a}
我们的基本设置涉及带角流形；相较于带边流形，其优点是这一类对乘积封闭，在讨论纤维化时很方便（见\sr{chapter:III}{第三章}）。先回顾有关术语。在带角流形中，考虑固定余维数 $k\geq1$ 的角点集合。这个集合的连通分支称为\emph{连通边界层}；这样的层在整个流形中的闭包称为\emph{连通边界面}。始终要求这些面没有自交；这意味着每个余维数为 $l$ 的角点，恰好属于余维数 $1\leq k\leq l$ 的 $2^l-1$ 个不同的连通边界面。去掉连通性要求，定义\emph{边界面}为同一余维数、两两不交的连通边界面的并；\emph{边界层}则为同一余维数、闭包两两不交的连通边界层的不交并。\znote{96}{$2^l-1$ 来自局部模型：在余维 $l$ 的角点附近，带角流形微分同胚于 $\{x\in\R^n:x^{n-l+1},\dots,x^n\geq0\}$ 中原点的一个邻域，此时恰有 $l$ 个坐标为零。过该点的每个连通边界层由「哪几个坐标为零、其余为正」指定，故有 $2^l-1$ 个非空选择，各层的闭包就是过该点的边界面；「无自交」的假设排除了更复杂的相遇方式。特别地，余维一的角点（$l=1$）恰属于一个边界面，这与带边流形的直觉一致。[来源：nLab, "manifold with boundary", https:/\allowbreak /\allowbreak ncatlab.org/\allowbreak nlab/\allowbreak show/\allowbreak manifold+\allowbreak with+\allowbreak boundary（2026-11-15 访问）。]}

\begin{definition*}
\emph{带角的正合辛流形} $(M,\varpi_M,\theta_M,I_M)$ 是紧致光滑带角流形 $M$，赋予辛形式 $\varpi_M$、满足 $d\theta_M=\varpi_M$ 的一形式 $\theta_M$，以及与 $\varpi_M$ 相容的殆复结构 $I_M$。要求满足以下两个凸性条件：由
\[
\varpi_M(\cdot,X_{\theta_M})=\theta_M
\]
定义的\emph{负 Liouville 向量场}沿 $M$ 的所有边界面都严格指向内部；边界还必须\emph{弱 $I_M$-凸}，即 $I_M$-全纯曲线不能接触 $\partial M$，除非完全包含在边界中。
\end{definition*}
这个定义是为兼顾各种需要而作的折中。缺点是两类凸性并未很好地联系起来，偶尔会使处理变得笨拙。当讨论只集中于一个流形时，省略 $\varpi_M,\theta_M,I_M$ 中的下标。

\begin{definition*}
带角正合辛流形的\emph{同构}是微分同胚 $\phi:M\to N$，满足 $\phi^*\theta_N=\theta_M+d(\text{某个函数})$，因而也满足 $\phi^*\varpi_N=\varpi_M$。
\end{definition*}
在所有这类同构构成的空间 $\operatorname{Iso}(M,N)$ 上赋予 $C^{\infty}$ 拓扑。$M=N$ 时，它成为拓扑群 $\operatorname{Aut}(M)$。注意，由于正合性假设，$\operatorname{Aut}(M)$ 中的任意同痕自动是 Hamilton 同痕。还将使用稍受限制、但技术上更简单的自同构概念：令 $\operatorname{Aut}(M,\partial M)$ 为在 $\partial M$ 附近等于恒等映射、且满足 $\phi^*\theta=\theta+d(\text{某个在 }\partial M\text{ 附近为零的函数})$ 的微分同胚 $\phi$ 所成的群。\znote{96}{这一句含一步验算。设 $\phi_t$ 是 $\operatorname{Aut}(M)$ 中的同痕，$\phi_t^*\theta=\theta+df_t$；对 $t$ 求导得 $\mathcal L_{X_t}\theta=d\dot f_t$，再用 Cartan 公式 $\mathcal L_{X_t}\theta=d(\theta(X_t))+\iota_{X_t}\varpi$，便得 $\iota_{X_t}\varpi=d(\dot f_t-\theta(X_t))$，故 $X_t$ 是 Hamilton 向量场，其 Hamilton 函数为 $\theta(X_t)-\dot f_t$。关键在右端自动恰当：$\varpi$ 自身正合，所以 $\iota_{X_t}\varpi$ 的 de Rham 类必为零，不存在 $H^1(M;\R)$ 型的障碍。反过来，$\varpi$ 正合也使每个 Hamilton 向量场的流都留在 $\operatorname{Aut}(M)$ 里：沿流 $\phi_t^*\varpi=\varpi$ 与 $\phi_t^*\theta=\theta+df_t$ 同时成立，$f_t$ 由 $\iota_{X_t}\varpi=dH_t$ 沿轨道积分给出。于是 $\operatorname{Aut}(M)$ 中的同痕与 Hamilton 同痕是同一批东西，$\operatorname{Iso}(M,N)$ 上取的 $C^\infty$ 拓扑也只反映 Hamilton 方向的变化。[来源：nLab, "Hamiltonian vector field", https:/\allowbreak /\allowbreak ncatlab.org/\allowbreak nlab/\allowbreak show/\allowbreak Hamiltonian+\allowbreak vector+\allowbreak field（2026-11-15 访问）。]}

\begin{definition*}
若 $\theta|_L$ 为正合一形式，就称 Lagrangian 子流形 $L\subset M$ \emph{正合}；除非另有说明，始终假设这些子流形是闭的，且与 $\partial M$ 不交。
\end{definition*}
这一类子流形在 $\operatorname{Aut}(M)$ 中的同痕以及负 Liouville 向量场的流下保持不变。\znote{96}{正合性的作用在此：$\varpi|_L=d(\theta|_L)$ 恒成立，故 $\dim L=\tfrac12\dim M$ 且 $\theta|_L$ 闭的子流形已是 Lagrangian；要求 $\theta|_L$ 整体恰当，即 $[\theta|_L]=0\in H^1(L;\R)$，是更强的条件，等价于 $L$ 上任何闭曲线的辛面积为零。有了整体原函数 $h_L$，$\phi^*\theta=\theta+df$ 直接把正合性搬过去：$\phi^*(\theta|_{\phi(L)})=d(h_L+f|_L)$，故同痕像仍正合；负 Liouville 流的情形同理。这正是这一类子流形能作为本章「对象」的前提：$\theta|_L$ 的整体原函数 $h_L$ 使 $L$ 上闭曲线的辛面积为零，作用量泛函因而无需 Novikov 系数即可定义。[来源：nLab, "lagrangian submanifold", https:/\allowbreak /\allowbreak ncatlab.org/\allowbreak nlab/\allowbreak show/\allowbreak Lagrangian+\allowbreak submanifold（2026-11-15 访问）。]}

\begin{definition*}
带角正合辛流形的\emph{变形}，就是固定流形 $M$ 上的一参数数据族 $(\varpi^s,\theta^s,I^s)$（$s\in[0,1]$），对每个 $s$ 都
\SourcePage{105}
满足上述条件。
\end{definition*}
这一背景下有 Moser 引理的一个较弱对应物，可表述如下。设 $M,N$ 是带角正合辛流形。
\begin{definition*}
若对某个 $c\in\R$ 有 $\rho^*\theta_N=e^{-c}\theta_M+d(\text{某个函数})$，就称嵌入 $\rho:M\to N$ 为\emph{正合共形辛嵌入}。
\end{definition*}
\begin{lemma}\label{st:7.1}
若两个带角正合辛流形可以相互变形，则存在从其中一个到另一个的正合共形辛嵌入。
\end{lemma}
\begin{proof}
取一个变形，像标准 Moser 论证那样，考察由 $\partial\theta^s/\partial s$ 得到的依赖于 $s$ 的向量场。它未必沿边界面指向内部，但可以通过加上 $X_{\theta^s}$ 的一个足够大的倍数来补救。积分所得向量场，即得所需嵌入。
\end{proof}
还可以说得更精确些。在所有正合共形辛嵌入的空间中，考察刚才构造的嵌入 $\rho$ 的同痕类。它不依赖于任何选择；在保持端点不变的变形同伦下不变；并且与变形的复合相容。特别地，若 $\tilde\rho$ 是将变形反向进行所得到的嵌入，那么 $\tilde\rho\circ\rho$ 与恒等映射同痕。

\subsection{接触型边界的情形}\label{sec:7b}
为补充上面的讨论，现在考虑一类更熟悉的带边辛流形。一些基础参考文献为\bcite{42}、\bcite{100}。
\begin{definition*}
\emph{具有接触型边界的正合辛流形}是一个紧致带边而不带角的流形 $M$，赋予辛形式 $\varpi_M$ 与一形式 $\theta_M$，满足 $d\theta_M=\varpi_M$，且 $X_{\theta_M}$ 沿 $\partial M$ 指向内部。
\end{definition*}
一个基本构造，是以典范方式接上一个无限锥，将这样的 $M$ 扩大为非紧致正合辛流形：
\begin{equation}\label{eq:7.1}
\begin{aligned}
(\widehat M,\varpi_{\widehat M},\theta_{\widehat M})
={}&(M,\varpi_M,\theta_M)\cup_{\partial M}\\[-2pt]
&\bigl(\R^+\times\partial M,\ d(e^r(\theta_M|_{\partial M})),\ e^r(\theta_M|_{\partial M})\bigr),
\end{aligned}
\end{equation}
其中 $r$ 为锥的径向变量；全书约定 $\R^{\pm}\subset\R$ 是闭的，即都包含零。“接缝”处的可微结构，利用两边 Liouville 流给出的 $\partial M$ 的领状邻域来定义。因此，$M$ 上的负 Liouville 场 $X_{\theta_M}$ 与 $\R^+\times\partial M$ 上的 $-\partial_r$ 拼合为 $\widehat M$ 上的负 Liouville 场。类似地，在 $\partial M$ 的某个邻域 $U\subset M$ 中，存在唯一的函数芽 $h_M:U\to(0,1]$，满足
\begin{equation}\label{eq:7.2}
h_M^{-1}(1)=\partial M,\qquad X_{\theta_M}\cdot h_M=-h_M.
\end{equation}
在锥上令 $h_{\widehat M}(r,y)=e^r$，就可将它光滑延伸为 $U\cup(\R^+\times\partial M)\subset\widehat M$ 上的函数 $h_{\widehat M}$。\znote{97}{$h_M$ 可读作「到边界的 Liouville 距离」：$X_{\theta_M}\cdot h_M=-h_M$ 说明沿负 Liouville 流 $\phi_t$（$\dot\phi_t=X_{\theta_M}$）有 $h_M(\phi_t x)=e^{-t}h_M(x)$，而 $h_M^{-1}(1)=\partial M$ 把它归一化。正因如此，锥上取 $h_{\widehat M}=e^r$ 才与 $\theta_{\widehat M}=e^r(\theta_M|_{\partial M})$ 配对，负 Liouville 场拼成 $-\partial_r$。芽的唯一性也随之清楚：$h_M$ 沿每条轨道被端点值唯一确定，故它在 $\partial M$ 的邻域内没有选择余地。顺带一提，后面引理 7.5 证明中取的 $M^{\epsilon}=M\setminus h_M^{-1}((e^{-\epsilon},1])$ 正是用 $h_M$ 把边界推里一点，而 $h_M$ 在 $1$ 附近的水平集给出 $\partial M$ 的领状邻域。[来源：原书 (7.1)--(7.2) 及本节关于 $h_{\widehat M}$ 延拓的论述。]}

\SourcePage{106}
\begin{definition*}
两个\eref{7.1}类型的流形之间，适当的同构概念是微分同胚 $\hat\phi:\widehat M\to\widehat N$，满足 $\hat\phi^*\theta_{\widehat N}-\theta_{\widehat M}=d(\text{某个紧支集函数})$。
\end{definition*}
特别地，这样的映射在某个紧子集之外保持 Liouville 向量场 $\partial_r$。因此，对某个 $C\gg0$，它在 $[C,\infty)\times\partial M\subset\widehat M$ 上形如
\[
\hat\phi(r,y)=(r-w(y),\psi(y)),
\]
其中 $\psi:\partial M\to\partial N$ 为接触微分同胚，$w\in C^{\infty}(\partial M,(-\infty,C])$ 满足 $\psi^*(\theta_N|_{\partial N})=e^w(\theta_M|_{\partial M})$。称 $\hat\phi$ \emph{在无穷远附近为接触型}。所有这样的同构构成的空间 $\operatorname{Iso}(\widehat M,\widehat N)$ 带有自然拓扑；它允许同痕改变无穷远处的接触微分同胚，因此这些同痕未必具有紧支集。照常记 $\operatorname{Aut}(\widehat M)=\operatorname{Iso}(\widehat M,\widehat M)$。以下 Moser 引理的版本是标准事实：
\begin{lemma}\label{st:7.2}
设 $(\varpi_M^s,\theta_M^s)$（$s\in[0,1]$）是 $M$ 上具有接触型边界的正合辛结构的光滑族。那么，由此得到的非紧致流形 $\widehat M^0,\widehat M^1$ 之间，存在在无穷远附近为接触型的正合辛同构。
\end{lemma}
设 $M$ 为具有接触型边界的正合辛流形。沿用一个众所周知的方法，例如见\bcite{76}，考虑与 $\varpi_M$ 相容的殆复结构 $J$，要求它在 $\partial M$ 的某个邻域中在 $X_{\theta_M}$ 的流下不变，并满足
\begin{equation}\label{eq:7.3}
dh_M\circ J=-\theta_M,
\end{equation}
其中 $h_M$ 如\eref{7.2}。称这样的殆复结构\emph{在边界附近为接触型}。给定带复结构 $I_S$ 的 Riemann 曲面 $S$ 与 $J$-全纯映射 $u:S\to M$，考虑函数 $\rho=h_M\circ u:u^{-1}(U)\to(0,1]$。由\eref{7.3}，它满足
\begin{equation}\label{eq:7.4}
\begin{aligned}
\Delta\rho&=-d(d\rho\circ I_S)=-d(dh_M\circ J\circ du)\\
&=d(\theta_M\circ du)=u^*\varpi_M\geq0.
\end{aligned}
\end{equation}
由极大值原理，若 $\rho$ 在 $S$ 的某个内点达到其最大可能值一，则它在那里必然局部为常数。因此，任取在边界附近为接触型的 $I_M=J$，就能使 $\partial M$ 弱凸，从而使 $M$ 成为带角正合辛流形。

我们需要这一基本凸性论证的几个变体。首先，它在以下意义下适用于殆复结构族：
\begin{lemma}\label{st:7.3}
设 $(J^z)_{z\in S}$ 是 $M$ 上与 $\varpi_M$ 相容的殆复结构族，以某个 Riemann 曲面 $S$ 为参数。假设每个 $J^z$ 都在 $\partial M$ 附近为接触型。那么，$S\times M$ 上由 $J(z,x)=I_S(z)\times J^z(x)$ 给出的殆复结构使 $S\times\partial M$ 弱凸。
\end{lemma}
\SourcePage{107}
其次，可以考察完备化\eref{7.1}中的伪全纯曲线。$M$ 上每个在边界附近为接触型的殆复结构，都有到 $\widehat M$ 的典范延伸 $\hat J$，在锥上对于 $r$ 的平移不变。它满足 $dh_{\widehat M}\circ\hat J=-\theta_{\widehat M}$，因此用与前面相同的计算，得到：
\begin{lemma}\label{st:7.4}
设 $S$ 是紧致连通带边 Riemann 曲面，$u:S\to\widehat M$ 是 $\hat J$-全纯映射，满足 $u(\partial S)\subset M$。那么也有 $u(S)\subset M$。
\end{lemma}
极大值原理论证还有积分形式；它既可以给出前一引理的另一种证明，也能得到稍为一般的结果。具体地，设 $N$ 为任意正合辛流形，$M$ 为具有接触型边界的同维正合辛流形，$\iota:M\hookrightarrow\operatorname{int}(N)$ 为满足 $\omega_M=\iota^*\omega_N$、$\theta_M=\iota^*\theta_N$ 的嵌入。为简洁起见，把它写作包含。取与 $\omega_N$ 相容的殆复结构 $J$，其在 $M$ 上的限制在边界附近为接触型。令 $S$ 为紧致连通带边 Riemann 曲面，$u:S\to N\setminus\operatorname{int}(M)$ 为满足 $u(\partial S)\subset\partial M$ 的 $J$-全纯映射。显然，
\begin{equation}\label{eq:7.5}
\begin{aligned}
0\leq\int_Su^*\omega_N&=\int_{\partial S}u^*\theta_M=\int_{\partial S}u^*(-dh_M\circ J)\\
&=\int_{\partial S}-dh_M\circ du\circ I_S.
\end{aligned}
\end{equation}
若 $\xi$ 是沿 $\partial S$、对于边界定向为正的切向量，则 $I_S\xi$ 指向内部，而 $du(I_S\xi)$ 不指向 $M$ 的内部，因此 $dh_M(du(I_S\xi))\geq0$。除非 $u$ 为常值，否则这与\eref{7.5}矛盾。作为一个简单应用，得到：
\begin{lemma}\label{st:7.5}
设 $S$ 是紧致连通带边 Riemann 曲面，$u:S\to N$ 是 $J$-全纯映射，满足 $u(\partial S)\subset\operatorname{int}(M)$。那么也有 $u(S)\subset\operatorname{int}(M)$。
\end{lemma}
\begin{proof}
对于某个 $\epsilon>0$，以稍小的子集 $M^{\epsilon}=M\setminus h_M^{-1}((e^{-\epsilon},1])$ 替换 $M$。适当选择 $\epsilon$，可以使 $J|_{M^{\epsilon}}$ 在边界附近仍为接触型，使 $u(\partial S)$ 位于 $\operatorname{int}(M^{\epsilon})$ 中，并使 $u$ 与 $\partial M^{\epsilon}$ 横截。因此 $S^{\epsilon}=u^{-1}(N\setminus\operatorname{int}(M^{\epsilon}))$ 仍为紧致带边 Riemann 曲面，其边界映到 $\partial M^{\epsilon}$。对其每个连通分支应用前面的论证，得到 $u(S^{\epsilon})\subset M^{\epsilon}$，从而 $u(S)\subset\operatorname{int}(M)$。
\end{proof}
最后，给定任意带角正合辛流形 $N$，可以如下\emph{磨圆角点}。令 $F_1,\ldots,F_d$ 为 $N$ 的余维一边界面。对每个 $F_j$，选择光滑函数 $f_j:N\to\R^+$，满足 $F_j=f_j^{-1}(0)$，且在 $F_j$ 的向外法向量上 $df_j<0$，从而特别有
\SourcePage{108}
$X_{\theta_N}\cdot f_j>0$ 在 $F_j$ 的某个邻域中成立。令 $f=f_1f_2\cdots f_d$。我们断言，在 $\operatorname{int}(N)$ 中充分靠近边界的所有点处，$X_{\theta_N}\cdot f>0$。这一不等式可以改写为
\begin{equation}\label{eq:7.6}
\sum_j\frac{X_{\theta_N}\cdot f_j}{f_j}>0.
\end{equation}
设 $\operatorname{int}(N)$ 中的一列点收敛到某个极限，该极限属于 $F_j$ 当且仅当 $j\in E$，其中 $E\subset\{1,\ldots,d\}$ 为某个非空子集。沿此点列，\eref{7.6}中 $j\in E$ 的各项都趋于 $+\infty$，其余各项有界，故该不等式最终成立。证明了这一断言后，立即可见：
\begin{lemma}\label{st:7.6}
对任意充分小的正数 $\delta$，$M=f^{-1}([\delta,\infty))\subset N$ 配备 $\varpi_N,\theta_N$ 的限制后，是具有接触型边界的正合辛流形。此外，经过足够长的时间，负 Liouville 流将把整个 $N$ 压入 $\operatorname{int}(M)$。
\end{lemma}

\section{经典 Floer 理论}\label{sec:8}
\subsection{Riemann 曲面}\label{sec:8a}
从这里开始，固定一个带角正合辛流形 $M$，在以下整个讨论中把它作为“目标空间”。将 $M$ 上给定的结构记为 $\varpi=\varpi_M$、$\theta=\theta_M$、$I=I_M$。
\begin{definition*}
设 $\widehat S$ 是紧致带边 Riemann 曲面，$\Sigma$ 是有限个边界点所成的集合，并分成两部分 $\Sigma=\Sigma^-\sqcup\Sigma^+$。称 $S=\widehat S\setminus\Sigma$ 为\emph{边界带标点的 Riemann 曲面}，称 $\Sigma^-$、$\Sigma^+$ 中的点分别为它的\emph{入向无穷远点}和\emph{出向无穷远点}；这些点仍作为数据的一部分保留。注意，可以从 $S$ 典范地恢复 $\widehat S$。对于其中最简单的曲面，采用以下记号：
\begin{itemize}
\item $D$ 是 $\CC$ 中的闭单位圆盘；
\item $H$ 是闭上半平面，有一个入向无穷远点；
\item $\overline H$ 是与 $H$ 相同的曲面，但将其无穷远点视为出向的；
\item $Z=\R\times[0,1]$，坐标为 $(s,t)$，入向无穷远点为 $s=-\infty$，出向无穷远点为 $s=+\infty$。
\end{itemize}
边界带标点的 Riemann 曲面的\emph{同构}，是保持入向与出向无穷远点之区分的双全纯映射。
\end{definition*}
\begin{definition*}
$S$ 的一组 \emph{Lagrangian 标记}是以连通分支 $C\subset\partial S$ 为指标的一族正合 Lagrangian 子流形 $\{L_C\}$，其中 $L_C\subset M$。每个 $\zeta\in\Sigma\subset\widehat S$ 都属于两个边界分支 $C_{\zeta,0}$、$C_{\zeta,1}$ 的闭包；这两个分支
\SourcePage{109}
可以相同。因此，给定一组 Lagrangian 标记后，就有与之对应的一对子流形 $(L_{\zeta,0},L_{\zeta,1})$。约定：当 $\zeta\in\Sigma^-$ 时，按边界的自然定向，$C_{\zeta,1}$ 位于 $\zeta$ 之前，$C_{\zeta,0}$ 位于其后；对于 $\zeta\in\Sigma^+$，约定相反。
\end{definition*}

\subsection{Floer 上同调}\label{sec:8b}
固定一个特征为 $2$ 的域 $\K$。对于 $M$ 中的每一对正合 Lagrangian 子流形 $(L_0,L_1)$，都有一个 Floer 上同调群 $HF^{\mathrm{pr}}(L_0,L_1)$，它是有限维 $\K$-向量空间。这里的 $\mathrm{pr}$ 是“初步版本”的缩写，用来与后面将介绍的更精细的理论区分；在后者中，$\operatorname{char}(\K)$ 可以任意，而且 Floer 群带有分次。对每个带有 Lagrangian 标记的边界带标点 Riemann 曲面 $S$，都可以赋予一个同态
\begin{equation}\label{eq:8.1}
\Phi_S:\bigotimes_{\zeta^+\in\Sigma^+}HF^{\mathrm{pr}}(L_{\zeta^+,0},L_{\zeta^+,1})
\longrightarrow\bigotimes_{\zeta^-\in\Sigma^-}HF^{\mathrm{pr}}(L_{\zeta^-,0},L_{\zeta^-,1}).
\end{equation}
这个反变约定可能显得有些奇怪，但要记住，这里讨论的是 Floer\emph{上同调}。这些映射满足开弦拓扑量子场论（TQFT）公理的一种形式，参见例如\bcite{122}；其中最重要的是下面的\emph{复合律}。设 $S_1,S_2$ 是带有 Lagrangian 标记的边界带标点 Riemann 曲面。假设有无穷远点 $\zeta_1^+\in\Sigma_1^+$、$\zeta_2^-\in\Sigma_2^-$，它们配备同一\emph{有序}对子流形。像\sr{fig:8.1}{图~8.1} 那样将两者粘合，得到另一个边界带标点 Riemann 曲面 $S=S_1\#S_2$。由于 TQFT 的另一性质——拓扑不变性——具体怎样完成粘合并不重要。$S$ 带有诱导的 Lagrangian 标记，且
\begin{equation}\label{eq:8.2}
\Phi_S=(\Id\otimes\cdots\otimes\Phi_{S_1}\otimes\cdots\otimes\Id)
\circ(\Id\otimes\cdots\otimes\Phi_{S_2}\otimes\cdots\otimes\Id).
\end{equation}
其规则是：将 $\Phi_{S_2}$ 输出中的 $\zeta_2^-$ 因子，代入 $\Phi_{S_1}$ 输入中的 $\zeta_1^+$ 分量。\znote{101}{反变约定的来源是：Floer 群是路径空间上作用量泛函的 Morse 理论，而带标点的曲面 $S$ 给出的是「推前」——把入向端的闭链沿模空间推到出向端；翻译成上同调的语言，这就成了从出向因子到入向因子的映射，方向随之反转。若把 $L$ 看作对象、$HF(L_0,L_1)$ 看作态射空间，同一件事就是复合的结合律，这也是后文把 $H\F(M)$ 称作范畴、把 $\Phi_S$ 称作复合的原因。[来源：nLab, "Floer homology", https:/\allowbreak /\allowbreak ncatlab.org/\allowbreak nlab/\allowbreak show/\allowbreak Floer+\allowbreak homology（2026-11-15 访问）；并见原书 (8.2)--(8.3)。]}

\subsection{Donaldson–Fukaya 范畴}\label{sec:8c}
一般而言，任意开弦 TQFT 都对应一个 Frobenius 范畴。在 Floer 理论中，由此得到一个 $\K$-线性范畴 $H\F(M)^{\mathrm{pr}}$，称为 Donaldson–Fukaya 范畴的初步版本。对象是正合 Lagrangian 子流形 $L\subset M$，从 $L_0$ 到 $L_1$ 的态射空间是 Floer 群 $HF^{\mathrm{pr}}(L_0,L_1)$。复合
\begin{equation}\label{eq:8.3}
HF^{\mathrm{pr}}(L_1,L_2)\otimes HF^{\mathrm{pr}}(L_0,L_1)\longrightarrow HF^{\mathrm{pr}}(L_0,L_2)
\end{equation}
是\emph{三角形乘积}。用 TQFT 的语言来说，它就是不变量 $\Phi_S$：这里 $S$ 是去掉三个边界点的圆盘，其中一个点入向、两个点出向，并用 $L_0,L_1,L_2$ 标记。恒等态射是 $\Phi_H\in HF^{\mathrm{pr}}(L,L)$，其中在 $\partial H$ 上放置标记 $L$。
\SourcePage{110}
\SourceFigure{figure-8-1}{0.50469176}{8.1}
复合的结合律，以及 $\Phi_H$ 是单位元这一事实，都由 TQFT 的粘合公理推出。$H\F(M)^{\mathrm{pr}}$ 的 Frobenius 性质如下：对于每个 $L$，不变量 $\Phi_{\overline H}$ 给出一个线性映射 $HF^{\mathrm{pr}}(L,L)\to\K$。将其与复合结合，对于任意 $L_0,L_1$，得到非退化配对
\[
HF^{\mathrm{pr}}(L_1,L_0)\otimes_{\K}HF^{\mathrm{pr}}(L_0,L_1)\longrightarrow\K.
\]
更直接地，定义这一配对的曲面，是两个无穷远点都出向的无限条带。接上这样的条带，可以把入向无穷远点换成出向无穷远点。这意味着两类点的区别在一定程度上是人为的。确实，可以不作这种区别，从而使这里的公理体系与 TQFT 的原始表述\bcite{7}一致；但由于后面会说明的技术原因，我们选择保留它。\znote{102}{保留这一区分有两个实际理由：条带状端与非齐次项 $K\in\Omega^1(S,\mathcal H)$ 都要沿 $s$ 方向定义（见 (8.6)--(8.7)），而 Floer 方程 (8.8) 本身就不对称——$\partial_su+J(\partial_tu-X)=0$ 把 $s$ 与 $t$ 区别对待。若抹掉入向与出向的区别，得到的是 TQFT 的原始公理，但那时无法直接写出上链层次的映射 $C\Phi_S$；本章保留区分，正是为了能逐点定义计数：粘合两张曲面必须一进一出，若两端同类便无法直接接上。[来源：原书 §8b 本段与 §8d--8f 的构造，并见原书文献 [7]。]}

还应指出 Donaldson–Fukaya 范畴的两个更具特殊性的性质。首先，对于任意 $L$，都有一个 Frobenius 代数的典范同构，即 Piunikhin–Salamon–Schwarz（PSS）映射\bcite{106}：
\begin{equation}\label{eq:8.4}
H^*(L;\K)\cong HF^{\mathrm{pr}}(L,L).
\end{equation}
理解\eref{8.4}的最佳方式，是将它视为 TQFT 框架的一种自然推广的一部分：在这种推广中，曲面 $S$ 还带有额外的边界标点，对这些点的处理类似于 Gromov–Witten 理论。

其次，每个正合 Lagrangian 子流形的同痕 $\{L_s\}_{0\leq s\leq1}$ 都给出 $HF^{\mathrm{pr}}(L_0,L_1)$ 中的一个特殊元素，具有以下性质：它只依赖于 $\{L_s\}$ 相对于端点的同伦类；对于
\SourcePage{111}
常值同痕，它是恒等态射；同痕作为路径的复合，对应于取相关元素的三角形乘积。直接推论是：这些元素在 $H\F(M)^{\mathrm{pr}}$ 中全都是同构。因此，同痕的 Lagrangian 子流形给出 Donaldson–Fukaya 范畴中同构的对象，进而 Floer 上同调群具有同痕不变性。从形式上说，这一构造以基本 TQFT 结构的另一种扩张为基础；其中允许 Lagrangian 边界条件沿 $\partial S$ 上的点光滑变化。

\subsection{条带状端}\label{sec:8d}
现在开始回顾 Floer 理论的基本基础。正如本章开头列举的参考文献所表明的，这些都是熟悉的材料；但对于后面讨论 Fukaya 范畴，它们是很好的预备练习。

首先，虽然边界带标点的 Riemann 曲面已经足以表述 TQFT，但在实际构造时，采用稍为刚性的结构更好。
\begin{definition*}
将半无限条带记为 $Z^\pm=\R^\pm\times[0,1]$。设 $S$ 是边界带标点的 Riemann 曲面，无穷远点集合为 $\Sigma^\pm$。$S$ 的一组\emph{条带状端}，是对每个 $\zeta\in\Sigma^\pm$ 给定的固有全纯嵌入 $\epsilon_\zeta:Z^\pm\to S$，满足
\[
\epsilon_\zeta^{-1}(\partial S)=\R^\pm\times\{0,1\},\qquad
\lim_{s\to\pm\infty}\epsilon_\zeta(s,\cdot)=\zeta,
\]
并且各 $\epsilon_\zeta$ 的像两两不交。$Z^\pm$ 上的 Riemann 曲面结构可延伸到单点紧化 $\widehat Z^\pm=Z^\pm\cup\{\pm\infty\}$，每个条带状端也都延伸为全纯嵌入 $\widehat\epsilon_\zeta:\widehat Z^\pm\to\widehat S$，将 $\pm\infty$ 映到 $\zeta$。
\end{definition*}
需要这些端有几个原因。其一是，作为理论的分析框架所用的 $W^{1,p}$（$p>2$）范数不具有共形不变性。\znote{103}{选 $p>2$ 的技术原因：二维 Sobolev 嵌入 $W^{1,p}\subset C^0$ 在 $p>2$ 时成立，故 $u$ 连续，边界条件 $u(\partial S)\subset L_C$ 才有逐点意义；$p=2$ 时须改用迹映射陈述（见 §8h）。另一方面 $p=2$ 恰是方程共形不变的情形，此时 $W^{1,2}$ 范数不具共形不变性，模空间在共形变动下不易控制，故实际取 $p>2$，并辅以条带状端上的标准坐标把非齐次项搬过去。[来源：原书 §8d 本段与 §8h 关于 $W^{1,p}$、迹映射的讨论，并见原书文献 [29]。]}同样重要的是，定义 Floer 上同调时出现的非齐次项涉及 $Z=\R\times[0,1]$ 上的标准坐标，因此为了定义运算\eref{8.1}，需要在更一般的 Riemann 曲面的端上赋予类似结构。还有一点是，可以使\sr{fig:8.1}{图~8.1} 的粘合过程更加精确：给定 $\zeta_1^+$、$\zeta_2^-$ 和粘合长度 $l>0$，令
\begin{equation}\label{eq:8.5}
\begin{aligned}
S'_1&=S_1\setminus\epsilon_{\zeta_1^+}((l,\infty)\times[0,1]),\\
S'_2&=S_2\setminus\epsilon_{\zeta_2^-}((-\infty,-l)\times[0,1]),\\
S=S_1\#_l S_2&=(S'_1\sqcup S'_2)/\sim,
\end{aligned}
\end{equation}
其中对 $(s,t)\in[0,l]\times[0,1]$，作识别 $\epsilon_{\zeta_2^-}(s-l,t)\sim\epsilon_{\zeta_1^+}(s,t)$。

\subsection{Floer 数据与扰动数据}\label{sec:8e}
将 $M$ 上所有与 $\varpi$ 相容、且在边界附近与给定的 $I$ 相同的殆复结构 $J$ 所成的空间记为 $\mathcal J$。若 $T$ 是任意流形，则用 $C^\infty(T,\mathcal J)$ 表示由 $T$ 光滑参数化的殆复结构族
\SourcePage{112}
$J=\{J(t)\}_{t\in T}$ 所成的空间，要求它们在 $T\times U$ 上与 $I$ 相同，其中 $U$ 是 $\partial M$ 的某个邻域。类似地，用 $\mathcal H=C_c^\infty(\operatorname{int}(M),\R)$ 表示 $M$ 上在边界附近为零的光滑函数所成的空间；$C^\infty(T,\mathcal H)$ 表示 $T\times M$ 上在某个如上的 $T\times U$ 上为零的函数。
\begin{definition*}
设 $L_0,L_1\subset M$ 是一对正合 Lagrangian 子流形。\emph{Floer 数据}由 $H\in C^\infty([0,1],\mathcal H)$、$J\in C^\infty([0,1],\mathcal J)$ 组成，满足以下性质：若 $X$ 是 $H$ 的时变 Hamilton 向量场，$\phi$ 是它的“流”，则 $\phi^1(L_0)$ 与 $L_1$ 横截相交。
\end{definition*}
设 $S$ 是带有 Lagrangian 标记的边界带标点 Riemann 曲面。假设已经为它选定条带状端，并且对每个无穷远点 $\zeta\in\Sigma$ 所对应的子流形对 $(L_{\zeta,0},L_{\zeta,1})$ 选定 Floer 数据 $(H_\zeta,J_\zeta)$。
\begin{definition*}
$S$ 的\emph{扰动数据}是一对 $(K,J)$，其中 $K\in\Omega^1(S,\mathcal H)$ 满足
\begin{equation}\label{eq:8.6}
K(\xi)|_{L_C}=0\qquad\text{对所有 }\xi\in TC\subset T(\partial S).
\end{equation}
换言之，$K$ 是取值于 $M$ 上函数空间的一形式；如果取一个属于连通分支 $C\subset\partial S$ 的边界点，以及该点处与边界相切的向量 $\xi$，那么 $K(\xi)\in\mathcal H$ 是沿 $L_C\subset M$ 为零的函数。扰动数据的第二部分，就是殆复结构族 $J\in C^\infty(S,\mathcal J)$。此外，$K,J$ 必须与选定的条带状端及 Floer 数据相容，即
\begin{equation}\label{eq:8.7}
\epsilon_\zeta^*K=H_\zeta(t)\,dt,\qquad J(\epsilon_\zeta(s,t))=J_\zeta(t)
\end{equation}
对每个 $\zeta\in\Sigma^\pm$ 和 $(s,t)\in Z^\pm$ 成立。
\end{definition*}
注意，反过来，在 $S=Z$ 上令 $K(s,t)=H(t)\,dt$，便可将 Floer 数据视为平移不变的扰动数据。Floer 数据与扰动数据这两个概念，规定了要加入伪全纯映射方程中的非齐次项的类型。

\subsection{非齐次伪全纯映射}\label{sec:8f}
设有一对正合 Lagrangian 子流形 $L_0,L_1$，以及 Floer 数据 $(H,J)$。像前面一样，以 $X$ 表示 $H$ 的时变向量场。定义 $\mathcal C(L_0,L_1)$ 为满足 $y(0)\in L_0$、$y(1)\in L_1$ 及 $dy/dt=X(t,y(t))$ 的映射 $y:[0,1]\to M$ 所成的集合。它们与 $\phi^1(L_0)\cap L_1$ 中的点一一对应，因此只有有限个。对 $u\in C^\infty(Z,M)$，Floer 方程为
\begin{equation}\label{eq:8.8}
\begin{cases}
\partial_su+J(t,u)(\partial_tu-X(t,u))=0,\\
u(s,0)\in L_0,\quad u(s,1)\in L_1.
\end{cases}
\end{equation}
以 $\M_Z(y_0,y_1)$ 表示\eref{8.8}满足渐近条件 $\lim_{s\to+\infty}u(s,\cdot)=y_1$、$\lim_{s\to-\infty}u(s,\cdot)=y_0$ 的解的集合，其中 $y_0,y_1\in\mathcal C(L_0,L_1)$。规定这种收敛是 $C^\infty$ 意义下的指数速率收敛；由于横截
\SourcePage{113}
相交假设，所有合理的收敛条件都是等价的。通过 $s$ 变量的平移，$\R$ 作用在 $\M_Z(y_0,y_1)$ 上。当 $y_0\ne y_1$ 时，作用是自由的，此时将商空间记为 $\M_Z^*(y_0,y_1)$。在 $y_0=y_1$ 的例外情形，$\M_Z(y_0,y_1)$ 是一个点，唯一的元素是驻定解 $u(s,t)=y_0(t)$；此时约定 $\M_Z^*(y_0,y_1)=\varnothing$。

更一般地，取一个带有 Lagrangian 标记的边界带标点 Riemann 曲面 $S$，为它配备条带状端、每个无穷远点处的 Floer 数据 $(H_\zeta,J_\zeta)$，以及相容的扰动数据 $(K,J)$。$K$ 确定一个向量场值一形式 $Y\in\Omega^1(S,C^\infty(TM))$：对于每个 $\xi\in TS$，$Y(\xi)$ 是 $K(\xi)$ 的 Hamilton 向量场。对 $u\in C^\infty(S,M)$，\emph{非齐次伪全纯映射方程}为
\begin{equation}\label{eq:8.9}
\begin{cases}
Du(z)+J(z,u)\circ Du(z)\circ I_S
=Y(z,u)+J(z,u)\circ Y(z,u)\circ I_S,\\
u(C)\subset L_C\quad\text{对所有 }C\subset\partial S,
\end{cases}
\end{equation}
其中 $I_S$ 是 $S$ 的复结构。可以将它更简洁地写成 $(Du-Y)^{0,1}=0$，或写成 $\overline\partial_J(u)=\nu(z,u(z))$，其中 $\nu=Y^{0,1}$。与\bcite{115}比较可知，我们没有使用最一般的非齐次项 $\nu$。显然，Floer 方程是\eref{8.9}的特例，其中 $S=Z$，数据 $(K,J)$ 平移不变。反过来，\eref{8.7}意味着\eref{8.9}在 $S$ 的条带状端上化为\eref{8.8}。因此，可以考虑满足以下收敛条件的\eref{8.9}的解集 $\M_S(\{y_\zeta\})$：
\begin{equation}\label{eq:8.10}
\lim_{s\to\pm\infty}u(\epsilon_\zeta(s,\cdot))=y_\zeta,
\qquad \zeta\in\Sigma^\pm,\quad y_\zeta\in\mathcal C(L_{\zeta,0},L_{\zeta,1}).
\end{equation}

\subsection{能量}\label{sec:8g}
先考虑 $L_0,L_1$ 及 Floer 数据 $(H,J)$。由于正合性，存在函数 $h_{L_k}\in C^\infty(L_k,\R)$，满足 $dh_{L_k}=\theta|_{L_k}$。路径空间 $\{y\in C^\infty([0,1],M):y(0)\in L_0,\ y(1)\in L_1\}$ 上的 $H$-扰动作用量泛函为
\begin{equation}\label{eq:8.11}
A_H(y)=\left(\int -y^*\theta+H(t,y(t))\,dt\right)+h_{L_1}(y(1))-h_{L_0}(y(0)).
\end{equation}
$A_H$ 的临界点集恰好是 $\mathcal C(L_0,L_1)$，\eref{8.8}则是某个 $L^2$ 度量下的负梯度流方程。\znote{105}{(8.8) 是 $A_H$ 在某个 $L^2$ 度量下的负梯度流方程，这一点值得记住：$A_H$ 的临界点集正是 $\mathcal C(L_0,L_1)$，而 (8.8) 正是从 $y_0$ 流向 $y_1$ 的下降轨道。于是 $E(u)=\int_Z|\partial_su|^2=A_H(y_0)-A_H(y_1)$ 就是 Morse 理论中「能量等于势差」的无穷维版本；它把解的能量换成端点数据，紧性论证才有统一的能量界可用。这也正是 (8.12) 的实质：先验界不是估计出来的，而是由端点决定，这也正是它被称为先验界的理由。[来源：nLab, "Floer homology", https:/\allowbreak /\allowbreak ncatlab.org/\allowbreak nlab/\allowbreak show/\allowbreak Floer+\allowbreak homology（2026-11-15 访问）；并见原书 (8.11)--(8.12)。]}因此，对任意 $u\in\M_Z(y_0,y_1)$，都有先验界
\[
E(u)=\int_Z|\partial u/\partial s|^2=A_H(y_0)-A_H(y_1).
\]
为了将这一点推广到\eref{8.9}，对每个连通分支 $C\subset\partial S$ 选取具有同样性质的函数 $h_{L_C}\in C^\infty(L_C,\R)$。于是，对 $u\in\M_S(\{y_\zeta\})$，有
\begin{equation}\label{eq:8.12}
E(u)\overset{\mathrm{def}}=\int_S\tfrac12|Du-Y|^2
=\sum_{\zeta\in\Sigma^\pm}\mp A_{H_\zeta}(y_\zeta)-\int_S R(z,u(z)).
\end{equation}
在对 $\zeta$ 求和时，每个 $h_{L_C}$ 的选择所含的常数歧义——不连通时为局部常数歧义——都会抵消，因此第一项是有意义的。
\SourcePage{114}
至于第二项，曲率被积形式 $R\in\Omega^2(S,\mathcal H)$ 在 $S$ 的局部坐标 $(s,t)$ 中为
\[
R=\bigl(\partial_sK(\partial_t)-\partial_tK(\partial_s)
+\{K(\partial_s),K(\partial_t)\}\bigr)\,ds\wedge dt.
\]
这里 $\{\cdot,\cdot\}$ 是 Poisson 括号。由于 $R$ 在 $S$ 的端上为零，积分有一个与映射 $u$ 无关的界，因此 $u\in\M_S(\{y_\zeta\})$ 的能量 $E(u)$ 也有这样的界。

\subsection{Cauchy–Riemann 算子}\label{sec:8h}
设 $S$ 是带边 Riemann 曲面，$E\to S$ 是带有辛结构 $\varpi_E$ 和相容复结构 $J_E$ 的向量丛，$F\subset E|_{\partial S}$ 是 Lagrangian 子丛。用 $C^\infty(S,E,F)$ 表示在边界上的限制取值于 $F$ 的截面空间，用 $C^\infty(S,\Omega_S^{0,1}\otimes E)$ 表示 $E$-值 $(0,1)$-形式的空间，不附加任何边界条件。对于较低的光滑性，同样可以定义局部 $W^{1,p}$ 截面的空间 $W_{\mathrm{loc}}^{1,p}(S,E,F)$，其中 $2\leq p<\infty$。当 $p>2$ 时，截面连续，边界条件显然有意义；当 $p=2$ 时，通过限制映射 $W_{\mathrm{loc}}^{1,2}(S,E)\to L_{\mathrm{loc}}^2(\partial S,E)$ 定义边界条件，参见例如\bcite{29}。还可定义局部 $L^p$ 的 $(0,1)$-形式空间 $L_{\mathrm{loc}}^p(S,\Omega_S^{0,1}\otimes E)$。给定 $E$ 上任意联络 $\nabla$，不必是复联络或辛联络，将对应的 Cauchy–Riemann 算子定义为协变微分的 $(0,1)$ 部分：
\begin{equation}\label{eq:8.13}
\overline\partial_\nabla=\nabla^{0,1}:C^\infty(S,E,F)\longrightarrow C^\infty(S,\Omega_S^{0,1}\otimes E).
\end{equation}
它可延伸为 $W_{\mathrm{loc}}^{1,p}(S,E,F)\to L_{\mathrm{loc}}^p(S,\Omega_S^{0,1}\otimes E)$。若 $S$ 紧致，则这一延伸是 Fredholm 算子。

为了在非紧致曲面上建立\eref{8.13}的合理分析理论，需要假设 $S$ 带有条带状端，并且 $(E,F,\nabla)$ 在这些端上的渐近行为受到充分控制。精确表述需要一些预备定义。
\begin{definition*}
对曲面 $S$ 的每个端 $\zeta$，给定向量丛 $E_\zeta\to[0,1]$，配备辛结构 $\varpi_{E_\zeta}$、相容复结构 $J_{E_\zeta}$、\emph{辛联络} $\nabla_\zeta$，以及 $k=0,1$ 时的 Lagrangian 子空间 $\Lambda_{\zeta,k}\subset E_{\zeta,k}$。称 $(E_\zeta,\varpi_{E_\zeta},J_{E_\zeta},\nabla_\zeta,\Lambda_{\zeta,0},\Lambda_{\zeta,1})$ 为 $\zeta$ 处的\emph{极限数据}。通常方便假设这一数据\emph{非退化}，意即沿 $[0,1]$ 的 $\nabla_\zeta$-平行移动把 $\Lambda_{\zeta,0}$ 映到与 $\Lambda_{\zeta,1}$ 横截的子空间。
\end{definition*}
设 $\pi^\pm:Z^\pm\to[0,1]$ 是沿 $t$ 方向的投影。要求向量丛 $E$ 配备指定的同构
\begin{equation}\label{eq:8.14}
\epsilon_\zeta^*E\cong(\pi^\pm)^*E_\zeta,
\end{equation}
并且这些同构在渐近意义下与极限数据相容。也就是说，用\eref{8.14}识别两端后，差 $\varpi_E-\varpi_{E_\zeta}$、$J_E-J_{E_\zeta}$、$\nabla-\nabla_\zeta$ 在 $s\to\pm\infty$ 时以指数速率衰减，例如在 $C^1$ 拓扑中；同时，$F_{\epsilon_\zeta(s,k)}\to\Lambda_{\zeta,k}$ 也在同样意义下成立，这里的收敛在 $(E_\zeta)_k$ 的 Lagrangian 子空间 Grassmann 流形中理解。
\SourcePage{115}
\begin{definition*}
若 $\overline\partial_\nabla$ 在各端满足\eref{8.14}及上述渐近条件，就称其为\emph{容许的}。若所有极限数据还都是非退化的，就称 $\overline\partial_\nabla$ 本身\emph{非退化}。
\end{definition*}
端上的平凡化确定一个沿 $s$ 方向不变的度量，用它来定义\emph{整体} $W^{1,p}$ 截面及 $L^p$ 截面的空间 $W^{1,p}(S,E,F)$、$L^p(S,\Omega_S^{0,1}\otimes E)$。于是\eref{8.13}延伸为连续算子
\begin{equation}\label{eq:8.15}
\overline\partial_\nabla:W^{1,p}(S,E,F)\longrightarrow L^p(S,\Omega_S^{0,1}\otimes E).
\end{equation}
若满足前述非退化条件，这些延伸都是 Fredholm 算子，其核与余核不依赖于 $p$。更一般的做法，是引入在各端带有某个权重 $\epsilon\geq0$ 的空间 $W^{1,p;\epsilon}$、$L^{p;\epsilon}$；这些空间的元素是指数衰减的截面，当 $\epsilon=0$ 时则退化为先前考虑的通常 Sobolev 空间。对于任意容许的 $\overline\partial$ 算子，都存在 $\epsilon_0>0$，使得当 $\epsilon\in(0,\epsilon_0)$ 时，延伸
\begin{equation}\label{eq:8.16}
\overline\partial_\nabla:W^{1,p;\epsilon}(S,E,F)\longrightarrow L^{p;\epsilon}(S,\Omega_S^{0,1}\otimes E)
\end{equation}
是 Fredholm 算子，其核与余核不依赖于权重，比较\bcite{45}，命题 4.1。若 Cauchy–Riemann 算子非退化，同一结论对 $\epsilon\in[0,\epsilon_0)$ 也成立；这意味着在此情形，\eref{8.16}的核与余核分别与\eref{8.15}相同。

有时还将使用带有更高阶导数 $k>1$ 的 Sobolev 空间 $W^{k,p}$，以及加权版本 $W^{k,p;\epsilon}$；上述讨论同样适用于它们。

\subsection{线性化}\label{sec:8i}
\eref{8.9}的解集 $\M_S(\{y_\zeta\})$ 可以置于标准的函数空间框架中。固定 $2<p<\infty$。设 $\mathcal B_S$ 是局部属于 $W^{1,p}$、并在条带状端上以 $W^{1,p}$ 意义收敛于某组 $\{y_\zeta\}$ 的映射 $u:S\to M$ 所成的 Banach 流形。这个流形在点 $u$ 处的切空间正是 $W^{1,p}(S,E,F)$，其中 $E=u^*TM$，而对每个分支 $C\subset\partial S$，有 $F|_C=u^*TL_C$。在 $E$ 上赋予殆复结构 $J_E(z)=J(z,u(z))$，这里 $J$ 是扰动数据的一部分。存在自然的 Banach 向量丛 $\mathcal E_S\to\mathcal B_S$，其纤维是空间 $L^p(S,\Omega_S^{0,1}\otimes u^*TM)$。可以把 $\overline\partial_J-\nu$ 看成这个丛的光滑截面，显然，其零点集就是各模空间 $\M_S(\{y_\zeta\})$ 的并；这也给这些空间赋予了适当的拓扑。模空间中一点处的线性化记为
\begin{equation}\label{eq:8.17}
D_{S,u}:(T\mathcal B_S)_u\longrightarrow(\mathcal E_S)_u,
\end{equation}
它是非退化 Cauchy–Riemann 算子，因而是 Fredholm 算子。
\begin{definition*}
若 $D_{S,u}$ 是满射，就称 $u$ \emph{正则}。若对所有 $u\in\M_S(\{y_\zeta\})$ 都如此，则称模空间本身\emph{正则}。这意味着它是一个适当维数的光滑流形，其维数等于 $D_{S,u}$ 的指标。若所有模空间都正则，
\SourcePage{116}
则称扰动数据 $(K,J)$ 本身\emph{正则}；对 Floer 数据也采用同样的术语。
\end{definition*}
为了将 $D_{S,u}$ 识别为前面所考虑的某个具体算子 $\overline\partial_\nabla$，在 $E$ 上赋予从 $\varpi$ 拉回的辛结构。联络 $\nabla$ 并不唯一，因为算子的定义只涉及它的 $(0,1)$ 分量；一种可能的选择方式如下。设 $\tilde u:S\to S\times M$ 是 $u$ 的图。固定 $Z\in C^\infty(S\times M,T^*S\otimes TM)$，满足 $Z(z,u(z))=du(z)$，且处处有 $(Z-Y)^{0,1}=0$。每个这样的 $Z$ 都对应 $E$ 上的一个联络 $\nabla$。具体地，取 $\widetilde X\in C^\infty(S\times M,TM)$，将其沿 $\tilde u$ 的拉回记为 $X\in C^\infty(S,E)$，令
\begin{equation}\label{eq:8.18}
\nabla X=\tilde u^*[Z,\widetilde X].
\end{equation}
为看出这个定义不依赖于 $\widetilde X$ 的选择，而且确实定义了联络，在 $TM\to S\times M$ 上任取无挠联络 $\widetilde\nabla^0$，将其沿 $\tilde u$ 的拉回记为 $\nabla^0$。借助它，可将\eref{8.18}改写为
\[
\nabla_\xi X=\nabla^0_\xi X-\tilde u^*(\widetilde\nabla^0_X\langle Z,\xi\rangle)
\]
对所有 $\xi\in C^\infty(S,TS)$ 成立；这里 $\langle Z,\xi\rangle\in C^\infty(S\times M,TM)$ 来自自然配对 $TS\otimes(T^*S\otimes TM)\to TM$。这表明 $\nabla$ 等于 $\nabla^0$ 加上一个零阶项，该项依赖于 $Z$ 沿 $u$ 的图的导数。接下来，用 $S$ 上的局部全纯坐标 $z=s+it$ 把计算写得更明确，得到 $\nabla$ 的 $(0,1)$ 部分为
\[
\begin{aligned}
(\nabla X)^{0,1}
={}&\tfrac12\bigl(\nabla^0_{\partial_s}X+J\nabla^0_{\partial_t}X
-\widetilde\nabla^0_X\langle Z,\partial_s\rangle-J\widetilde\nabla^0_X\langle Z,\partial_t\rangle\bigr)\\
={}&\tfrac12\bigl(\nabla^0_{\partial_s}X+J\nabla^0_{\partial_t}X
+\tilde u^*(\widetilde\nabla^0_XJ)\,du(\partial_t)\bigr)\\
&-\tfrac12\tilde u^*\widetilde\nabla^0_X\bigl(Y(\partial_s)+JY(\partial_t)\bigr).
\end{aligned}
\]
这确实是伪全纯映射理论中线性化算子的标准形式，最后一项是为计入\eref{8.9}中的非齐次项而加入的。因此 $D_{S,u}=\overline\partial_\nabla$。为了满足容许性条件，选择 $Z$，使得当 $s+it$ 是某个条带状端 $\epsilon_\zeta:Z^\pm\to S$ 上的标准坐标时，$Z(\partial_s)$ 在 $s\to\pm\infty$ 时以指数速率趋于零，收敛在 $C^2$ 拓扑中理解。此时，$(E,\varpi_E,J_E,F)$ 及 $\nabla$ 的渐近结构由以下极限数据描述：向量丛 $E_\zeta=y_\zeta^*TM$，带有从 $\varpi$ 继承的辛结构，以及殆复结构 $J_{E_\zeta}(t)=J_\zeta(t,y_\zeta(t))$；Lagrangian 子空间 $\Lambda_{\zeta,k}=T(L_{\zeta,k})_{y_\zeta(k)}$；以及沿积分曲线 $y_\zeta$ 线性化 $H_\zeta$ 的“流”所得到的辛联络 $\nabla_\zeta$。

\begin{note}\label{st:8.1}
在某些特殊情形，线性化算子有更直接的描述。假设 $J$ 不依赖于 $z\in S$ 且可积，并且 $Y=0$。那么 $TM$ 及其拉回 $E$ 自然是全纯丛，
\SourcePage{117}
$D_{S,u}$ 就是对应的典范 Dolbeault 算子。另一种情形是 $u\equiv x$ 为常值，而 $K$ 及其一阶导数在 $x$ 处均为零。此时二阶导数 $D^2K$ 可视为 $\Omega^1(S,\operatorname{sym}^2(T M_x)^\vee)$ 中的元素，因而给出平凡丛 $E=S\times T M_x$ 上的辛联络。于是 $D_{S,u}$ 是这个联络相对于复结构 $J_E(z)=J(z,x)$ 的 $(0,1)$ 部分。
\end{note}

\subsection{TQFT 的定义}\label{sec:8j}
对每一对正合 Lagrangian 子流形 $(L_0,L_1)$，选择正则 Floer 数据 $(H_{L_0,L_1},J_{L_0,L_1})$。
\begin{definition*}
定义 \emph{Floer 上链群} $CF^{\mathrm{pr}}(L_0,L_1)$ 为 $\K$-向量空间，对每个 $y\in\mathcal C(L_0,L_1)$ 赋予一个基向量。这个群上的边界算子为
\begin{equation}\label{eq:8.19}
\partial(y_1)=\sum_{y_0}\#\M_Z^*(y_0,y_1)\,y_0,
\end{equation}
其中 $\#$ 表示对孤立点的个数模 $2$ 计数。$HF^{\mathrm{pr}}(L_0,L_1)$ 是 $(CF^{\mathrm{pr}}(L_0,L_1),\partial)$ 的上同调。\znote{109}{这里有两处命名要留意。其一，$CF^{\mathrm{pr}}$ 虽称「上链群」，但 (8.19) 的 $\partial$ 把 $y_1$ 送到 $y_0$，按几何次数是降次的。其二，在 $\operatorname{char}(\K)=2$ 的初步理论里没有分次可谈，同调与上同调只差一次整体平移，命名不影响内容；到了带分次的精细理论（§9），次数由 Maslov 型指标给出，$\partial$ 降低该次数一次，此时必须固定约定，否则符号与次数会一起错位。[来源：原书 §8b、§8j 本段与 §9 关于分次的讨论。]}
\end{definition*}
接下来，对每个带有 Lagrangian 标记 $\{L_C\}$ 及条带状端的边界带标点 Riemann 曲面 $S$，选择与上述 Floer 数据相容的正则扰动数据 $(K_S,J_S)$。定义映射
\begin{equation}\label{eq:8.20}
\begin{aligned}
C\Phi_S:&\ \bigotimes_{\zeta^+\in\Sigma^+}CF^{\mathrm{pr}}(L_{\zeta^+,0},L_{\zeta^+,1})
\longrightarrow\bigotimes_{\zeta^-\in\Sigma^-}CF^{\mathrm{pr}}(L_{\zeta^-,0},L_{\zeta^-,1}),\\
C\Phi_S(\otimes_{\zeta^+}y_{\zeta^+})
={}&\sum_{\{y_{\zeta^-}\}}\#\M_S(\{y_{\zeta^-},y_{\zeta^+}\})\,(\otimes_{\zeta^-}y_{\zeta^-}),
\end{aligned}
\end{equation}
并以 $\Phi_S$ 表示它在上同调上诱导的映射。对这一领域不熟悉的读者应当记住：尽管这里的叙述相当简略，构造的每一步都是一项大工程。仅仅为了定义 $\partial$ 和 $C\Phi_S$，就需要：(i) 横截性结果，说明正则性是 Floer 数据与扰动数据的一般位置条件；(ii) 紧致性结果，保证在正则情形，相关模空间只有有限个孤立点。此后，适当的粘合定理说明：(iii) $\partial^2=0$；(iv) $C\Phi_S$ 是上链映射。进一步还需证明：(v) 所得理论在典范同构意义下不依赖于任何选择；(vi) 它满足复合公理\eref{8.2}。前四个问题在文献中已有充分论述，而且在构造 Fukaya 范畴时还会再次出现，见\secref{9k}、\secref{9l}，因此这里不再讨论。关于最后两个实际上相互交织的步骤，作一些说明是适当的。

首先着手 (v)。先证明，只要端上的 Floer 数据固定，$\Phi_S$ 就不依赖于 $S$ 的共形结构、条带状端的选择及扰动数据。为此，取由 $[0,1]$ 参数化的一族共形结构和扰动数据，在给定的两组之间插值，
\SourcePage{118}
并用所得全部模空间的并来定义一个上链同伦。这可以看作参数化模空间的首次出现，虽然是相当平凡的一例；后面会广泛讨论这种空间。接着转向 (vi)。取配备了条带状端、Lagrangian 边界条件及扰动数据的曲面 $S_1,S_2$，像\eref{8.5}那样沿无穷远点 $\zeta_1^+,\zeta_2^-$ 将它们粘合，形成 $S=S_1\#_lS_2$。假设这两个端的 Floer 数据相同，从而在 $S$ 上自然得到扰动数据。再假设满足适当的正则性条件，且 $l\gg0$，就可以证明复合律在\emph{上链层次}成立。当然，这个等式只对 $S$ 上这组非常特殊的共形结构及扰动数据成立。然而，已经证明，上同调层次的映射 $\Phi_S$ 不依赖于这些选择。对于 Fukaya 范畴这种上链层次的结构，这类论证将不再适用，这意味着我们必须对构造中所用的扰动数据施加精细得多的控制。

有了这些，(v) 中关于 Floer 数据选择无关性的剩余部分就完全是形式上的。考虑无限条带 $Z$，在边界分支 $C_k=\R\times\{k\}$（$k=0,1$）上分别赋予 Lagrangian 标记 $L_k$。在两个端 $\zeta_0=-\infty$、$\zeta_1=+\infty$，为对子流形 $(L_{\zeta_j,0},L_{\zeta_j,1})=(L_0,L_1)$ 配备正则 Floer 数据
\[
(H_{\zeta_j},J_{\zeta_j})=(H^j_{L_0,L_1},J^j_{L_0,L_1}),
\]
并选择相容的正则扰动数据。这给出所得 Floer 上同调群之间的映射
\[
\Phi_Z:HF^{\mathrm{pr}}(L_0,L_1)^1\longrightarrow HF^{\mathrm{pr}}(L_0,L_1)^0,
\]
通常称为\emph{延拓映射}\bcite{118}。上述同伦论证说明它是典范的；应用复合律则说明它是同构。

\begin{note}\label{st:8.2}
在许多情形，可以采用非齐次项为零的 Floer 数据与扰动数据。具体说，只要 $L_0\pitchfork L_1$，就能找到 $H=0$ 的正则 Floer 数据。类似地，考虑带有条带状端及 Lagrangian 边界条件 $\{L_C\}$ 的 Riemann 曲面 $S$。假设每个无穷远点所对应的一对 Lagrangian 子流形都横截相交，而且 $M$ 中没有同时位于所有 $L_C$ 上的点；第二个条件排除了\eref{8.9}的常值解。此时就能找到 $K=0$ 的正则扰动数据。
\end{note}

\subsection{移动边界条件}\label{sec:8k}
\begin{definition*}
设 $S$ 是带有条带状端的 Riemann 曲面。$S$ 的一组\emph{移动边界条件}，是随 $z\in\partial S$ 光滑变化的正合 Lagrangian 子流形族 $L_z\subset M$，并且当 $z\in\partial S\cap\epsilon_\zeta(Z^\pm)$ 时局部不变。因此，每个无穷远点仍有与之对应的一对子流形 $(L_{\zeta,0},L_{\zeta,1})$。
\end{definition*}
在扰动数据
\SourcePage{119}
$(K,J)$ 的定义中，应将\eref{8.6}替换为以下稍复杂的条件：
\begin{equation}\label{eq:8.21}
d(K(\xi)|_{L_z})=\varpi(\cdot,\partial_\xi L_z)\in\Omega^1(L_z)
\qquad\text{对任意 }\xi\in T(\partial S)_z.
\end{equation}
这里 $\partial_\xi L_z$ 是 $C^\infty(TM|_{L_z})$ 中任意一个表示该 Lagrangian 子流形无穷小变形的截面，因此\eref{8.21}右端是这种变形所对应的标准正合一形式。于是可以考虑与前面相同的方程\eref{8.9}。$K$ 的条件使得\eref{8.11}中的边界项可以通过积分处理掉，其余理论无需重大调整便可照常进行，得到推广的映射\eref{8.1}，满足 TQFT 公理的显然变体。\znote{111}{(8.21) 说的是：$L_z$ 随 $z$ 变动时，$K(\xi)$ 在 $L_z$ 上的限制不再是零（那是 (8.6) 的情形），而是一个恰当一形式，其外微分恰为 $\varpi(\cdot,\partial_\xi L_z)$，即该无穷小变形对应的标准正合一形式。这样 (8.11) 中沿边界的积分项可以写成恰当形式而被抵消，能量恒等式与 TQFT 公理才继续成立。紧接着的 (8.22) 又说明：只要各边界分支都是区间，移动边界条件其实可化为固定边界条件，故这一推广增加的是方便，而非实质内容。[来源：原书 (8.21)--(8.22) 及本段论证。]}

实际上，这一推广的幅度比初看起来小得多。为看清这一点，假设 $\partial S$ 的所有连通分支都是开区间。于是，存在由 $z\in S$ 参数化的自同构族 $\phi_z\in\operatorname{Aut}(M,\partial M)$，使得 $\phi_z(L_z)$ 局部不变，并且在条带状端上有 $(\partial/\partial s)\phi_{\epsilon_\zeta(s,t)}=0$。坐标变换
\begin{equation}\label{eq:8.22}
u(z)\longmapsto\phi_z(u(z))
\end{equation}
保持方程\eref{8.9}的类型，同时将移动边界条件化为固定的边界条件。尽管如此，用移动边界条件来思考往往很方便，以下应用说明了这一点。设 $\{L_s\}_{0\leq s\leq1}$ 是正合 Lagrangian 同痕。取 $S=H$，配备移动边界条件 $L_{\psi(s)}$，其中 $\psi\in C^\infty(\R,[0,1])$ 满足 $s\ll0$ 时 $\psi(s)=0$，$s\gg0$ 时 $\psi(s)=1$。由此得到一个特征元素
\begin{equation}\label{eq:8.23}
\Phi_H\in HF^{\mathrm{pr}}(L_0,L_1),
\end{equation}
它实际上就是\secref{8c}所提到的元素。那里陈述的全部性质，都由标准的同伦和粘合论证得到。例如，将\eref{8.23}与三角形乘积复合，对任意 $L$ 得到同构 $HF^{\mathrm{pr}}(L,L_0)\to HF^{\mathrm{pr}}(L,L_1)$ 及 $HF^{\mathrm{pr}}(L_1,L)\to HF^{\mathrm{pr}}(L_0,L)$。事实上，细察粘合论证并应用变换\eref{8.22}，就可将这些同构归结为先前考虑的延拓映射。

\subsection{边界点处的求值}\label{sec:8l}
假设给定一个带有 Lagrangian 标记的边界带标点 Riemann 曲面，并给定额外的有限标点集 $\Pi\subset\partial S$，分为负标点 $\Pi^-$ 和正标点 $\Pi^+$。此时可以引入 TQFT 映射\eref{8.1}的推广：
\begin{equation}\label{eq:8.24}
\begin{aligned}
\Phi_{S,\Pi}:&\ \bigotimes_{\zeta^+\in\Sigma^+}HF^{\mathrm{pr}}(L_{\zeta^+,0},L_{\zeta^+,1})
\otimes\bigotimes_{z^+\in\Pi^+}H^*(L_{C_{z^+}};\K)\\
&\longrightarrow\bigotimes_{\zeta^-\in\Sigma^-}HF^{\mathrm{pr}}(L_{\zeta^-,0},L_{\zeta^-,1})
\otimes\bigotimes_{z^-\in\Pi^-}H^*(L_{C_{z^-}};\K).
\end{aligned}
\end{equation}
\SourcePage{120}
这里 $C_z$ 是 $z$ 所属的边界分支。除了通常将无穷远点粘合外，现在还可以作边界连通和：在相应 Lagrangian 边界条件相同的点 $z_1^+\in\Pi_1^+$、$z_2^-\in\Pi_2^-$ 处，把两个曲面 $S_1,S_2$ 接合起来。对于对应的算子\eref{8.24}，这相当于在相关因子 $H^*(L_{C_z};\K)$ 上施行复合。

最简单的情形是 $S$ 紧致。假设正则性成立，则\eref{8.9}的解所成的模空间 $\M_S$ 是光滑闭流形。求值映射 $\operatorname{ev}_z(u)=u(z)$（$z\in\Pi$）的乘积，在 $\prod_zL_{C_z}$ 中给出一个同调类。利用 Künneth 公式以及针对 $\Pi^-$ 因子的 Poincaré 对偶，可将它解释为多重线性映射\eref{8.24}。然而，如果令非齐次项为零，直接的能量论证就表明，$\M_S$ 只包含常值映射，因而所得不变量并没有多少新意。\znote{112}{这一段把 $H^*(L)$ 也纳入 TQFT：额外的边界标点 $z$ 上可以插入 $L_{C_z}$ 的上同调类，几何上就是把解在 $z$ 处的取值与 $\prod_zL_{C_z}$ 中的闭链求交。$S$ 紧致时，求值映射的乘积给出 $\prod_zL_{C_z}$ 中的同调类，用 Künneth 公式与 Poincaré 对偶即可读出 (8.24)。但若取零非齐次项，能量论证迫使解为常值映射，于是只剩 (8.25) 列出的四种初等映射；真正有用的是带一个出向端与一个边界标点的情形：取 $S=\overline H$ 就得到 PSS 同构 (8.4)。[来源：原书 §8l 本段与 (8.24)--(8.25)，并见原书文献 [120]。]}例如，取 $S=D$，Lagrangian 标记为 $L$，得到
\begin{equation}\label{eq:8.25}
\begin{cases}
\text{单位元 }1\in H^*(L;\K),&|\Pi^-|=1,\ |\Pi^+|=0,\\
\text{积分映射 }H^*(L;\K)\to\K,&|\Pi^-|=0,\ |\Pi^+|=1,\\
\text{恒等映射 }H^*(L;\K)\to H^*(L;\K),&|\Pi^-|=|\Pi^+|=1,\\
\text{杯积},&|\Pi^-|=1,\ |\Pi^+|=2.
\end{cases}
\end{equation}
为简要说明更一般的情形，假设 $\Sigma=\{\zeta\}$ 只含一个出向端，$\Pi=\{z\}$ 只含一个负标点。其思路是：取 Floer 上闭链 $\sum_km_ky_k\in CF^{\mathrm{pr}}(L_{\zeta,0},L_{\zeta,1})$，则对应的求值映射 $\operatorname{ev}_z:\M(y_k)\to L_{C_z}$ 的形式线性组合，在通常的拓扑意义下是 $L_{C_z}$ 上的闭链。其同调类，更准确地说是它的 Poincaré 对偶，就是 Floer 上同调类 $[\sum_km_ky_k]$ 在\eref{8.24}下的像。

一种等价而在技术上可能更简洁的做法，可以追溯到\bcite{120}，即采用 $H^*(L_{C_z};\K)$ 的 Morse 同调模型。这意味着，\eref{8.24}底层上链映射的系数，通过计数交集 $\M(y_k)\cap\operatorname{ev}_z^{-1}(W^u(x))$ 中的点来确定；这里 $W^u(x)$ 是 $-\nabla h_z$ 的不稳定流形，$h_z$ 是在 $L_{C_z}$ 上选取的 Morse 函数。例如，取 $S=\overline H$，Lagrangian 标记为 $L$，便可定义 PSS 映射\eref{8.4}。\secref{8c}中陈述的所有性质，现在都由\eref{8.25}及复合律推出。例如，逆映射恰由带有一个正标点的 $S=H$ 给出。从某种意义上说，这实际上是唯一需要的应用：可以利用粘合及 PSS 同构，把标点 $z$ 替换为满足 $L_{\zeta,0}=L_{\zeta,1}$ 的无穷远点 $\zeta$，从而将\eref{8.24}化为更基本的版本\eref{8.1}。

\SourcePage{121}
\section{Fukaya 范畴：初步版本}\label{sec:9}
\subsection{族}\label{sec:9a}
\begin{definition*}
设 $\widehat S$ 是某个光滑定向带边曲面，$\Sigma=\Sigma^+\sqcup\Sigma^-$ 是有限个边界点所成的集合，$S=\widehat S\setminus\Sigma$。取纤维丛
\[
\pi:\mathcal S\longrightarrow\mathcal R,
\]
其纤维为 $S$，结构群是 $\widehat S$ 上在 $\Sigma$ 处等于恒等映射的保定向微分同胚群。每个纤维 $\mathcal S_r$ 都有自然紧化 $\widehat{\mathcal S}_r$，它们典范地组成光滑固有纤维丛 $\widehat\pi:\widehat{\mathcal S}\to\mathcal R$。此外，$\Sigma^\pm$ 中的每个点都给出截面 $\zeta:\mathcal R\to\partial\widehat{\mathcal S}$，满足 $\mathcal S=\widehat{\mathcal S}\setminus\bigcup\zeta(\mathcal R)$。常将 $\Sigma^\pm$ 中的点与对应的截面等同。现在假设，$\mathcal S$ 的纤维方向切丛上配备了殆复结构 $I_{\mathcal S}$，并能延伸到 $\widehat{\mathcal S}$，从而使每个纤维 $\mathcal S_r$ 成为边界带标点的 Riemann 曲面。在这种情形，称 $\mathcal S$ 为\emph{边界带标点的曲面族}。
\end{definition*}

$\mathcal S\to\mathcal R$ 的一组条带状端，由对每个 $\zeta\in\Sigma^\pm$ 给定的、覆盖 $\mathcal R$ 的固有嵌入 $\epsilon_\zeta:\mathcal R\times Z^\pm\to\mathcal S$ 组成，它们在各纤维上限制为条带状端。这些嵌入自动延伸为光滑、逐纤维全纯的嵌入 $\mathcal R\times\widehat Z^\pm\to\widehat{\mathcal S}$，其在 $\pm\infty\in\widehat Z^\pm$ 处的限制，就是对应截面 $\zeta$。任意族都容许条带状端，根本原因是：对单个边界带标点的 Riemann 曲面，端的选择空间可缩。这与在 Riemann 曲面内部一点附近选取局部坐标的情形不同，后者有本质性的 $S^1$ 自由度，并由此产生引力后裔。

换一个稍有不同的角度，假设已经固定一组端 $\{\epsilon_\zeta\}$。此时通常方便采用使端保持不变的 $\mathcal S\to\mathcal R$ 的局部平凡化，其含义如下。选择一点 $r_0\in\mathcal R$，设 $S=\mathcal S_{r_0}$ 为该点上的纤维，条带状端为 $\epsilon_{S,\zeta}=\epsilon_\zeta(r_0,\cdot):Z^\pm\to S$。对充分小的邻域 $r_0\in U\subset\mathcal R$，存在逐纤维微分同胚 $\Psi:U\times S\to\mathcal S|_U$，它在 $\{r_0\}\times S$ 上为恒等映射，并满足
\begin{equation}\label{eq:9.1}
\epsilon_\zeta(r,s,t)=\Psi(r,\epsilon_{S,\zeta}(s,t)),\qquad r\in U,\ (s,t)\in Z^\pm.
\end{equation}
于是，$\Psi^*I_{\mathcal S}$ 是由 $r\in U$ 参数化的 $S$ 上复结构族，只在紧致子集 $S\setminus\bigcup\epsilon_\zeta(Z^\pm)$ 上变化。

\subsection{Kodaira–Spencer 映射}\label{sec:9b}
设 $S$ 是带有条带状端的 Riemann 曲面。考虑切丛上的标准 $\overline\partial$ 算子，边界条件取在全实子丛 $T(\partial S)$ 中。按\secref{8h}的术语，它是容许的，但不是非退化的，因此在带有充分小权重 $\epsilon>0$ 的加权 Sobolev 空间上成为 Fredholm 算子，见\eref{8.16}。实际上，希望容许稍大一些的定义域
\begin{equation}\label{eq:9.2}
W^{k,p;\epsilon}(S,TS,T(\partial S))^+\supset W^{k,p;\epsilon}(S,TS,T(\partial S)),
\end{equation}
\SourcePage{122}
由在每个条带状端上具有形式 $\epsilon_\zeta^*Z=c_\zeta\partial_s+\widetilde Z_\zeta$ 的向量场 $Z$ 组成，其中 $c_\zeta\in\R$、$\widetilde Z_\zeta\in W^{k,p;\epsilon}$。这只增加有限维，每个端增加一维，对应于常数 $c_\zeta$；因此 $\overline\partial$ 仍可延伸为 Fredholm 算子
\begin{equation}\label{eq:9.3}
\overline\partial^+:W^{k,p;\epsilon}(S,TS,T(\partial S))^+
\longrightarrow W^{k-1,p;\epsilon}(S,\Omega_S^{0,1}\otimes TS).
\end{equation}
将其余核记为 $\mathcal T_S$；它既不依赖于 $(k,p)$，也不依赖于充分小的 $\epsilon>0$。若 $S$ 稳定，即不存在连续的自同构族，则\eref{9.3}是单射，且 $\dim\mathcal T_S=|\Sigma|-3\chi(\widehat S)$。\znote{114}{按\eref{9.3}，$\mathcal T_S$ 是 $\overline\partial^+$ 的余核，可看作 $S$ 的模空间的切空间；当 $S$ 是带 $d+1$ 个标点的圆盘时 $|\Sigma|=d+1$、$\chi(\widehat S)=1$，故 $\dim\mathcal T_S=d-2$，与 $\mathcal R^{d+1}$ 的维数相符，这正是 $d\geq2$ 时模空间非空且无多余参数的维数解释。单射性用到 $S$ 稳定：若存在连续自同构族，$\overline\partial^+$ 会多出核，$\mathcal T_S$ 就不再是切空间。[来源：原书 §(9b) 末尾所引 [72]（带标点曲线的代数变形理论），第 94 页；维数计算并见 §(9c) 对 $\mathcal R^{d+1}$ 的定义。]}

设 $\mathcal S\to\mathcal R$ 是带有条带状端的 Riemann 曲面族。选择一点 $r_0\in\mathcal R$，将其纤维记为 $S=\mathcal S_{r_0}$，并取一个像\eref{9.1}那样使条带状端保持不变的局部平凡化 $\Psi$。对 $S$ 上拉回复结构族 $(\Psi^*I_{\mathcal S})$ 求导，便为 $(T\mathcal R)_{r_0}$ 中的每个切向量给出
\[
C_c^\infty(S,\End^{0,1}(TS))=C_c^\infty(S,\Omega_S^{0,1}\otimes TS)
\]
中的一个元素。向 $\operatorname{coker}(\overline\partial^+)$ 投影，定义 \emph{Kodaira–Spencer 映射}
\begin{equation}\label{eq:9.4}
\kappa_{r_0}:(T\mathcal R)_{r_0}\longrightarrow\mathcal T_S.
\end{equation}
它不依赖于平凡化，也不依赖于条带状端的选择。后一事实从定义中并不立即可见，因为我们偏好采用与条带状端结构紧密相关的函数空间。
\begin{definition*}
若对每个 $r_0\in\mathcal R$，\eref{9.4}都是双射，就称 $\mathcal S\to\mathcal R$ 为\emph{无穷小极小通用族}。
\end{definition*}
为了说明这一名称确有所据，应该证明这样的族具有局部通用性质。这并不困难，但我们实际上用不到它。有兴趣的读者可参看带标点 Riemann 曲面的代数几何变形理论\bcite{72}，第 94 页；两者完全类似。

\subsection{带标点圆盘}\label{sec:9c}
现在将讨论限制于与 Fukaya 范畴相关的一类曲面。
\begin{definition*}
对于 $d\geq0$，一个\emph{带 $(d+1)$ 个标点的圆盘}，是紧化 $\widehat S\cong D$ 为圆盘、并有一个入向和 $d$ 个出向无穷远点的边界带标点 Riemann 曲面 $S$。约定按沿圆盘边界的循环次序给无穷远点编号，使入向点为 $\Sigma^-=\{\zeta_0\}$，出向点为 $\Sigma^+=\{\zeta_1,\ldots,\zeta_d\}$；若 $S$ 有条带状端，则相应编号为 $\epsilon_0,\ldots,\epsilon_d$。$\partial S$ 的连通分支按循环次序记为 $C_0,\ldots,C_d$，其中 $\zeta_0$ 位于 $C_d$ 与 $C_0$ 之间；$d=2$ 的情形见\sr{fig:9.1}{图~9.1}。若 $S$ 带有 Lagrangian 标记，则相应记为 $L_0,\ldots,L_d$。与无穷远点对应的子流形对于是为
\begin{equation}\label{eq:9.5}
(L_{\zeta_k,0},L_{\zeta_k,1})=\begin{cases}(L_0,L_d),&k=0,\\(L_{k-1},L_k),&k>0.\end{cases}
\end{equation}

\end{definition*}
\SourcePage{123}
\SourceFigure{figure-9-1}{0.30946121}{9.1}
在稳定范围 $d\geq2$ 内，带 $(d+1)$ 个标点的圆盘没有非平凡自同构，并存在它们的万有族，记为
\begin{equation}\label{eq:9.6}
\mathcal S^{d+1}\longrightarrow\mathcal R^{d+1}.
\end{equation}
构造很直接。设 $\operatorname{Conf}_{d+1}(\partial D)\subset(\partial D)^{d+1}$ 是圆周上编号与循环次序相容的 $(d+1)$ 元点组所成的构形空间。$\operatorname{Aut}(D)\cong\operatorname{PSL}_2(\R)$ 自由且固有地作用在此空间上，令
\begin{equation}\label{eq:9.7}
\begin{aligned}
\mathcal R^{d+1}&=\operatorname{Conf}_{d+1}(\partial D)/\operatorname{Aut}(D),\\
\widehat{\mathcal S}^{d+1}&=\operatorname{Conf}_{d+1}(\partial D)\times_{\operatorname{Aut}(D)}D.
\end{aligned}
\end{equation}
投影 $\widehat{\mathcal S}^{d+1}\to\mathcal R^{d+1}$ 有截面 $\zeta_k([z_0,\ldots,z_d])=[z_0,\ldots,z_d,z_k]$；按照一般步骤，令 $\mathcal S^{d+1}=\widehat{\mathcal S}^{d+1}\setminus\bigcup_k\zeta_k(\mathcal R^{d+1})$。参数化形式的 Riemann 单值化定理表明，\eref{9.6}确实是万有族：每个带 $(d+1)$ 个标点的圆盘族，都由从其底空间到 $\mathcal R^{d+1}$ 的唯一光滑映射拉回 $\mathcal S^{d+1}$ 而得。类似的无穷小结论也不难证明，它说明\eref{9.6}是无穷小极小通用的。
\begin{definition*}
\emph{条带状端的万有选择}，是对每个 $d\geq2$，在 $\mathcal S^{d+1}\to\mathcal R^{d+1}$ 上选取一组条带状端 $\{\epsilon_k^{d+1}\}_{0\leq k\leq d}$。当然，其万有性意指：任何带 $(d+1)$ 个标点的圆盘族，都通过其分类映射继承条带状端。
\end{definition*}

\clearpage
\subsection{树的术语}\label{sec:9d}
涉及多个带标点圆盘的粘合运算，可以自然地用树来描述。这里汇集有关术语，其中一些将在后面使用。
\begin{definition*}
对于 $d>0$，\emph{$d$ 叶树}是固有嵌入的平面树 $T\subset\R^2$，具有 $d+1$ 条半无限边：指定其中一条称为\emph{根}，其余称为\emph{叶}。半无限边也称为\emph{外边}，其他边称为\emph{内边}。以 $\operatorname{Ve}(T)$ 表示顶点集，以
\SourcePage{124}
$\operatorname{Ed}^{\mathrm{int}}(T)$ 表示内边集。回顾图的\emph{旗标}是由一个顶点与一条相邻边构成的二元组 $f=(v,e)$。用 $\operatorname{Fl}^{\mathrm{int}}(T)$ 表示\emph{内旗标}的集合，即其中的边为内边的旗标。

按从根向叶“向上”流动的方向给 $T$ 定向，见\sr{fig:9.2}{图~9.2}。若 $e$ 的定向背离 $v$，则称旗标 $f=(v,e)$ 为\emph{正}的，否则称为\emph{负}的。每条内边 $e$ 对应两个旗标 $f^\pm(e)$，一正一负。同样，每条外边有唯一的旗标；边为叶时它为正，边为根时它为负。有时方便将与给定顶点 $v$ 相邻的旗标编号为 $f_0(v),\ldots,f_{|v|-1}(v)$，其中 $|v|$ 是顶点的度数。编号总是从唯一的负旗标 $f_0(v)$ 开始，按给定平面嵌入的逆时针方向继续。
\end{definition*}
\SourceFigure{figure-9-2}{0.43615338}{9.2}
\begin{definition*}
若所有顶点 $v$ 都满足 $|v|\geq3$，则称 $d$ 叶树\emph{稳定}；若 $|v|\geq2$，则称其\emph{半稳定}。
\end{definition*}

\subsection{树与粘合}\label{sec:9e}
假设给定一棵 $d$ 叶树 $T$；对每个顶点 $v$，给定带有条带状端的 $|v|$ 标点圆盘 $S_v$；对每条内边 $e$，给定粘合长度 $l_e>0$。从这些数据出发，粘合构造产生一个新的带条带状端的 $(d+1)$ 标点圆盘 $S_l$。

先注意，旗标 $f=(v,e)$ 与 $S_v$ 中的无穷远点之间存在典范双射，使负、正旗标分别对应入向、出向点。具体说，将 $S_v$ 的第 $k$ 个无穷远点对应于旗标 $f_k(v)$。相应调整记号，以 $\epsilon_f$ 表示 $S_v$ 在对应于 $f$ 的无穷远点附近的条带状端。粘合过程的第一步，是从每个内旗标对应的端切掉一段：
\begin{equation}\label{eq:9.8}
S'_v=S_v\setminus\left(\coprod_{\substack{f=(v,e)\\f\in\operatorname{Fl}^{\mathrm{int}}(T)}}\epsilon_f(Z_{l_e}^\pm)\right),
\end{equation}
\SourcePage{125}
其中按旗标为负或为正，分别采用 $Z_{l_e}^-=(-\infty,-l_e)\times[0,1]$ 或 $Z_{l_e}^+=(l_e,\infty)\times[0,1]$。然后，将这些截断端的剩余部分彼此识别：令 $S_l=\coprod_vS'_v/\sim$，其中对每条内边 $e$ 以及 $(s,t)\in[0,l_e]\times[0,1]$，规定
\begin{equation}\label{eq:9.9}
\epsilon_{f^-(e)}(s-l_e,t)\sim\epsilon_{f^+(e)}(s,t).
\end{equation}
$S_l$ 的条带状端继承自 $T$ 的外旗标所对应的端。更精确地，从根开始按循环次序将这些旗标编号为 $f_0,\ldots,f_d$，则 $S_l$ 的第 $k$ 个条带状端为
\begin{equation}\label{eq:9.10}
\epsilon_{l,k}=\epsilon_{f_k}:Z^\pm\longrightarrow\coprod_vS'_v\longrightarrow S_l.
\end{equation}
当 $T$ 只有一个顶点时，粘合过程没有作任何改变，只是返回原曲面；若有两个顶点，它就等于\eref{8.5}中的基本连通和运算；一般情况下，它可看作有限次这种运算的复合，执行次序任意。

\begin{note}\label{st:9.1}
曲面 $S_l$ 继承一个\emph{厚薄分解}。该分解来自粘合过程，并非曲面的内在结构，因此不同于双曲几何所给出的同名分解。薄部由条带状端及长度为 $l_e$ 的有限条带组成：
\[
\begin{aligned}
S_l^{\mathrm{thin}}&=\coprod_{k=0}^d\epsilon_{l,k}(Z^\pm)\ \sqcup\!
\coprod_{e\in\operatorname{Ed}^{\mathrm{int}}(T)}\epsilon_{f^+(e)}([0,l_e]\times[0,1]),\\
S_l^{\mathrm{thick}}&=\coprod_v\left(S_v\setminus\coprod_{k=0}^{|v|-1}\epsilon_{f_k(v)}(Z^\pm)\right).
\end{aligned}
\]

\end{note}
后面将看到，最好采用长度的负指数作为\emph{粘合参数}。更精确地，$\rho_e\in(-1,0)$ 与相应长度的关系为 $l_e=-\log(-\rho_e)/\pi$，因此 $\rho_e\to0$ 对应于 $l_e\to\infty$。粘合后的曲面 $S_\rho=S_l$ 组成一个 $(d+1)$ 标点圆盘族
\[
\mathcal S\longrightarrow\mathcal R=(-1,0)^{\operatorname{Ed}^{\mathrm{int}}(T)},
\]
其条带状端 $\mathcal R\times Z^\pm\to\mathcal S$ 在各纤维上由\eref{9.10}给出。从某种意义上说，也可容许某些粘合参数为零的极限情形。此时的规则是：在粘合过程的两个步骤\eref{9.8}、\eref{9.9}中，都略去退化边，即粘合参数为零的边。结果不再是单个带标点圆盘，而是若干这种圆盘的不交并；它由
\SourcePage{126}
将 $T$ 的所有非退化边压缩后得到的树 $\overline T$ 描述。特别地，$\overline T$ 的每个顶点 $\overline v$ 对应一个连通分支，它是带 $|\overline v|$ 个标点的圆盘。在所有粘合参数都为零的平凡情形，所得正是各 $S_v$ 的不交并。\znote{118}{此处“压缩”（收缩）内边的效果，是把该边两个端点所代表的标点合并成一个：粘合参数非零的边两侧恰被实际粘合，故在树上应把它收缩掉；参数为零的边不作粘合，曲面沿它分裂，但该边仍留在树上，用以记录各连通分支之间的关系。于是每个顶点恰好对应一个连通分支，顶点的度数就是该分支上的标点数。所有参数为零时树不变，曲面回到各 $S_v$ 的不交并，与正文最后一句一致。[来源：原书 §(9e) 的粘合构造，即\eref{9.8}、\eref{9.9}；本注只是该处逻辑的展开。]}将这些纤维加入 $\mathcal S$，得到部分紧化 $\overline{\mathcal S}$；它是光滑带角流形，并带有一个浸没映射 $\overline{\mathcal S}\to\overline{\mathcal R}=(-1,0]^{\operatorname{Ed}^{\mathrm{int}}(T)}$，显然这不是纤维化。

粘合构造的一种直接推广，是对每个顶点 $v$，从带有条带状端的 $|v|$ 标点圆盘族 $\mathcal S_v\to\mathcal R_v$ 出发。粘合得到新的族
\[
\mathcal S\longrightarrow\mathcal R=(-1,0)^{\operatorname{Ed}^{\mathrm{int}}(T)}\times\prod_v\mathcal R_v.
\]
像前面一样，容许粘合参数为零，便得到部分紧化 $\overline{\mathcal S}\to\overline{\mathcal R}$。

\subsection{Deligne–Mumford–Stasheff 紧化：结合多面体}\label{sec:9f}
作为集合，$\mathcal R^{d+1}$ 的 Deligne–Mumford–Stasheff 紧化为
\begin{equation}\label{eq:9.11}
\overline{\mathcal R}^{d+1}=\coprod_T\mathcal R^T.
\end{equation}
这里对所有稳定 $d$ 叶树 $T$ 取并，并令 $\mathcal R^T=\prod_{v\in\operatorname{Ve}(T)}\mathcal R^{|v|}$。为在\eref{9.11}上赋予拓扑与光滑结构，先作某个条带状端的万有选择，再采用粘合构造。具体地，取一棵稳定树 $T$，将万有族 $\mathcal S^{|v|}\to\mathcal R^{|v|}$ 配给每个顶点 $v$。将它们粘合，得到 $(d+1)$ 标点圆盘族
\begin{equation}\label{eq:9.12}
\mathcal S\longrightarrow\mathcal R=(-1,0)^{\operatorname{Ed}^{\mathrm{int}}(T)}\times\mathcal R^T,
\end{equation}
按定义，它由光滑的\emph{粘合映射}
\begin{equation}\label{eq:9.13}
\gamma^T:\mathcal R\longrightarrow\mathcal R^{d+1}
\end{equation}
分类。像前面所说明的，可以容许部分或全部粘合参数为零。结果是退化的粘合曲面，即由压缩后的树 $\overline T$ 控制的若干带标点圆盘的不交并。显然，这确定了\eref{9.11}中 $\overline T$-层的一点，从而定义粘合映射的典范延伸
\begin{equation}\label{eq:9.14}
\overline\gamma^T:\overline{\mathcal R}=(-1,0]^{\operatorname{Ed}^{\mathrm{int}}(T)}\times\mathcal R^T\longrightarrow\overline{\mathcal R}^{d+1}.
\end{equation}
在子集 $\{0\}^{\operatorname{Ed}^{\mathrm{int}}(T)}\times\mathcal R^T$ 上，它代表平凡的、退化程度最大的粘合运算，因此就是\eref{9.11}中 $T$-层上的恒等映射。
\SourcePage{127}
Deligne–Mumford–Stasheff 紧化的拓扑，由所有 $T$ 的映射\eref{9.14}诱导。也就是说，从 $\overline{\mathcal R}^{d+1}$ 到任意空间的映射连续，当且仅当它与所有 $\overline\gamma^T$ 的复合连续。可微结构由类似性质刻画，只需把“连续”替换为“光滑”。这样，$\overline{\mathcal R}^{d+1}$ 成为紧致带角流形，而且不依赖于条带状端的万有选择；此外，\eref{9.11}就是其内部及各连通边界层的分解。\znote{119}{当所有标点都在边界上时，$\overline{\mathcal R}^{d+1}$ 其实就是 Stasheff 结合多面体：它是 $d-2$ 维紧致带角流形，其面按标点的括号化方式编号，$K_3$ 是区间、$K_4$ 是五边形（两个二顶点树构成五边形的边，三个二顶点树构成其顶点，即\secref{9g}的五边形论证）。一族这样的多面体构成消解结合律的拓扑 operad，这正是\eref{9.11}的分层由树给出、且边界层按内边数降维的原因。[来源：nLab "associahedron"，https:/\allowbreak /\allowbreak ncatlab.org/\allowbreak nlab/\allowbreak show/\allowbreak associahedron（2026-11-15 访问）。]}

刚才陈述的性质并不能从构造中立即看出；事实上，甚至 $\overline{\mathcal R}^{d+1}$ 上的拓扑是否 Hausdorff，都可能还不清楚。验证这些性质的最佳方式，是采用 $\overline{\mathcal R}^{d+1}$ 的另一种定义，将它视为复 Deligne–Mumford 空间 $\overline{\mathcal M}_{0,d+1}$ 的实点集 $\overline{\mathcal M}_{0,d+1}(\R)$ 的一部分；这一解释见\bcite{58}、\bcite{60}。按这种方式引入时，$\overline{\mathcal R}^{d+1}$ 自带紧致光滑带角流形的典范结构。
\begin{lemma}\label{st:9.2}
在 $\overline{\mathcal R}^{d+1}$ 上赋予由其嵌入 $\overline{\mathcal M}_{0,d+1}(\R)$ 得到的光滑结构。则延伸后的粘合映射 $\overline\gamma^T$ 光滑，而且在所有粘合参数都为零的子集附近是局部微分同胚。
\end{lemma}
\begin{proof}[证明概要]
先回顾，嵌入 $\mathcal R^{d+1}\subset\overline{\mathcal M}_{0,d+1}(\R)\subset\overline{\mathcal M}_{0,d+1}$ 是通过将带标点圆盘“双倍化”为带有反全纯对合的穿孔 $\CP^1$ 实现的。圆盘上的条带状端也可双倍化，得到 $\CP^1$ 各穿孔附近的局部全纯坐标。将粘合参数复化，就能把\eref{9.14}延伸为映射
\begin{equation}\label{eq:9.15}
(\overline\gamma^T)^{\CC}:\operatorname{int}(D)^{\operatorname{Ed}^{\mathrm{int}}(T)}\times\mathcal R^T\longrightarrow\overline{\mathcal M}_{0,d+1},
\end{equation}
其中 $\operatorname{int}(D)=D\setminus\partial D$ 是开单位圆盘。更精确地，固定 $\mathcal R^T$ 中的一点，复化的粘合构造就在 $\operatorname{int}(D)^{\operatorname{Ed}^{\mathrm{int}}(T)}$ 上产生稳定带标点曲线的全纯族。\eref{9.15}是这个族的分类映射，因此它关于粘合参数全纯，并且在所有粘合参数都非零的开子集上光滑。

为推出它处处光滑，采用 Cauchy 公式的以下初等推论：设 $\overline U=\operatorname{int}(D)^q\times\R^p$、$U=(\operatorname{int}(D)\setminus\{0\})^q\times\R^p$，并对 $r\in\R^p$ 令 $\overline U_r=\operatorname{int}(D)^q\times\{r\}$。假设函数 $g:\overline U\to\CC$ 在 $U$ 上光滑，且在每个 $\overline U_r$ 上的限制都全纯，则 $g$ 光滑。最后，为证明 $\overline\gamma^T$ 是局部微分同胚，需要在所有粘合参数都为零的点处，计算其关于粘合参数的导数。复几何，更确切地说是带节点 Riemann 曲面的变形理论，再次使这一点成为简单的练习。
\end{proof}
有了\stref{引理}{9.2}，很容易看出，最初通过任意条带状端万有选择的粘合映射\eref{9.14}定义的 $\overline{\mathcal R}^{d+1}$ 的拓扑与光滑结构，
\SourcePage{128}
和嵌入 $\overline{\mathcal R}^{d+1}\subset\overline{\mathcal M}_{0,d+1}(\R)$ 给出的结构相同。这就推出了上面陈述的所有性质。

定义 $\mathcal S^T\to\mathcal R^T$ 为万有纤维化 $\mathcal S^{|v|}\to\mathcal R^{|v|}$（$v\in\operatorname{Ve}(T)$）的各拉回的不交并；$\mathcal S^T$ 的纤维是若干 $|v|$ 标点圆盘的不交并。万有族 $\mathcal S^{d+1}\to\mathcal R^{d+1}$ 有一个部分紧化，在集合意义下为
\begin{equation}\label{eq:9.16}
\overline{\mathcal S}^{d+1}=\coprod_T\mathcal S^T\longrightarrow\overline{\mathcal R}^{d+1}.
\end{equation}
通过族的总空间上的粘合映射，可在其中赋予拓扑及光滑结构。按定义，与\eref{9.12}对应的粘合映射 $\gamma^T$ 由逐纤维同构
\[
\Gamma^T:\mathcal S\longrightarrow\mathcal S^{d+1}
\]
覆盖。在解释一般粘合构造时已经指出，$\mathcal S$ 有自然的部分紧化 $\overline{\mathcal S}\to\overline{\mathcal R}$。将它与\eref{9.16}比较，可见 $\Gamma^T$ 延伸为 $\overline\Gamma^T:\overline{\mathcal S}\to\overline{\mathcal S}^{d+1}$。像前面一样，这些映射所诱导的拓扑与光滑结构，和把 $\overline{\mathcal S}^{d+1}$ 嵌入 $\overline{\mathcal M}_{0,d+1}$ 上的稳定曲线万有族得到的结构相同。

\subsection{一致地选择条带状端}\label{sec:9g}
像前面一样，作条带状端的万有选择，并考虑稳定 $d$ 叶树 $T$ 所对应的粘合构造。所得族\eref{9.12}带有两组条带状端，它们先验地可能不同。第一组通过粘合构造，继承自各 $\mathcal S^{|v|}$ 上条带状端的万有选择；第二组则通过分类映射\eref{9.13}，从 $\mathcal S^{d+1}$ 上的万有选择拉回。
\begin{definition*}
若对于每棵 $T$，都存在包含 $\{0\}\times\prod_v\mathcal R^{|v|}$ 的开子集 $\overline U\subset\overline{\mathcal R}$，使上述两组条带状端在 $U=\overline U\cap\mathcal R$ 上相同，就称条带状端的万有选择\emph{一致}。今后，将把这样的 $U$ 或 $\overline U$ 称为\emph{粘合参数充分小的子集}。
\end{definition*}
\begin{lemma}\label{st:9.3}
存在一致的条带状端万有选择。
\end{lemma}
\begin{proof}
证明简单但有些繁琐，因此只说明基本思路。首先值得注意的是，一致性条件具有归纳性质。设 $T$ 是至少有两个顶点的稳定 $d_0$ 叶树；单顶点树给出空的一致性条件。$T$ 的一致性，利用 $d=|v|-1<d_0$ 时的 $\{\epsilon_k^{d+1}\}$，部分地确定了 $\mathcal S^{d_0+1}$ 上的条带状端 $\{\epsilon_k^{d_0+1}\}$。其次，对单棵树 $T$，一致性条件能够满足。细察\stref{引理}{9.2}的证明，可见 $\overline\gamma^T$ 在粘合参数充分小的子集
\SourcePage{129}
$\overline U$ 上的限制是开嵌入，其像为 $\overline{\mathcal R}^{d_0+1}$ 中 $T$-层的一个邻域。特别地，$\gamma^T|_U$ 也是嵌入。因此，假设 $d<d_0$ 时的 $\{\epsilon_k^{d+1}\}$ 已固定，则对于某棵指定的 $T$，总能选取 $\{\epsilon_k^{d_0+1}\}$，使其沿 $\gamma^T|_U$ 的拉回，等于由粘合过程得到的 $\mathcal S|_U$ 上的条带状端。

论证的最后一步，是检查各棵 $T$ 所施加的条件彼此不矛盾。这同样是归纳论证，以 Deligne–Mumford–Stasheff 空间边界层的组合结构为基础。为具体起见，在最简单的非平凡情形 $d_0=4$ 中说明。假设已经选取 $\{\epsilon_k^3\}$ 和 $\{\epsilon_k^4\}$，并对稳定三叶树满足一致性条件。恰有五棵具有两个顶点的稳定四叶树 $T$；顶点度数必为 $3$ 和 $4$。它们给出五个延伸的粘合映射
\[
\overline\gamma^T:(-1,0]\times\mathcal R^3\times\mathcal R^4
=(-1,0]\times\{\text{点}\}\times(-1,1)\hookrightarrow\overline{\mathcal R}^5.
\]
$\overline{\mathcal R}^5$ 是五边形。限制到粘合参数充分小的子集 $\overline U_T$ 后，每个粘合映射都是嵌入，其像覆盖一个边界面的邻域，见\sr{fig:9.3}{图~9.3}。一致性要求：在这些邻域与内部 $\mathcal R^5$ 的交集上，$\{\epsilon_k^5\}$ 由 $\{\epsilon_k^3\}$、$\{\epsilon_k^4\}$ 确定。
\SourceFigure{figure-9-3}{0.40499956}{9.3}
需要证明，这五个条件可以同时满足。任意两个邻域 $\overline\gamma^{T_1}(\overline U_{T_1})$、$\overline\gamma^{T_2}(\overline U_{T_2})$，或者不交，或者在五边形的某个顶点附近重叠；该顶点对应一棵具有三个顶点的稳定四叶树 $T_3$。在第一种情形，一致性条件显然彼此不干扰；在
\SourcePage{130}
第二种情形，利用粘合映射 $\overline\gamma^{T_3}$，以及 $\{\epsilon_k^3\}$ 与 $\{\epsilon_k^4\}$ 之间的一致性，可以验证：当粘合参数充分小时，由 $T_1$、$T_2$ 粘合构造诱导的条带状端相同。因此，相应的 $\{\epsilon_k^5\}$ 的一致性条件能够同时满足。
\end{proof}

\subsection{族的模空间}\label{sec:9h}
固定带角正合辛流形 $M$。
\begin{definition*}
族 $\mathcal S\to\mathcal R$ 的一组 Lagrangian 标记，是各纤维上 Lagrangian 标记的局部常值族。与单个曲面的情形一样，选定标记后，自然为每个 $\zeta\in\Sigma^\pm$ 赋予一对正合 Lagrangian 子流形 $(L_{\zeta,0},L_{\zeta,1})$；它们是在任意纤维中相邻边界分支的标记。假设已经为 $\mathcal S\to\mathcal R$ 配备这些标记及一组条带状端，并为每对 $(L_{\zeta,0},L_{\zeta,1})$ 选定 Floer 数据 $(H_\zeta,J_\zeta)$。于是，这个族的扰动数据是一对 $(K,J)$，其中
\[
K\in\Omega^1_{\mathcal S/\mathcal R}(\mathcal S,\mathcal H),\qquad
J\in C^\infty(\mathcal S,\mathcal J),
\]
在任意纤维上的限制满足\secref{8e}的条件，包括\eref{8.7}。这个记号需要稍作说明：$K$ 是各纤维上一形式的光滑族，取值于 $M$ 上 Hamilton 函数的空间 $\mathcal H$。等价地，若 $T\mathcal S^v\to\mathcal S$ 表示沿纤维的切丛，可以把 $K$ 视为其对偶丛拉回到 $\mathcal S\times M$ 上后的一个截面，并在 $\mathcal S\times\partial M$ 附近及条带状端上附加条件。
\end{definition*}

固定 $\mathcal S\to\mathcal R$ 的扰动数据。对于每组 $\{y_\zeta\}_{\zeta\in\Sigma}$，其中 $y_\zeta\in\mathcal C(L_{\zeta,0},L_{\zeta,1})$，可以考虑模空间
\begin{equation}\label{eq:9.17}
\M_{\mathcal S}(\{y_\zeta\})=
\left\{\begin{array}{l|l}
(r,u)&r\in\mathcal R,\ u\in C^\infty(\mathcal S_r,M),\text{ 且对于扰动数据}\\[-2pt]
&\text{在纤维 }\mathcal S_r\text{ 上的限制，}u\text{ 满足\eref{8.9}、\eref{8.10}}
\end{array}\right\}.
\end{equation}
在某个小开子集 $U\subset\mathcal R$ 上，使用使条带状端保持不变的局部平凡化，可定义一个平凡 Banach 纤维丛 $\mathcal B_{\mathcal S|_U}\to U$，其纤维是\secref{8i}中的 $W^{1,p}$ 映射空间 $\mathcal B_{\mathcal S_r}$。还存在 Banach 向量丛 $\mathcal E_{\mathcal S|_U}\to\mathcal B_{\mathcal S|_U}$，在每个 $r\in U$ 上限制为 $L^p$ 空间所成的丛 $\mathcal E_{\mathcal S_r}$。非齐次 $\overline\partial$ 方程\eref{8.9}给出截面 $\mathcal B_{\mathcal S|_U}\to\mathcal E_{\mathcal S|_U}$，其零点集是所有 $\{y_\zeta\}$ 的 $\M_{\mathcal S|_U}(\{y_\zeta\})$ 的并。在模空间一点 $(r,u)$ 处，该截面的导数称为\emph{扩展线性化算子}：
\begin{equation}\label{eq:9.18}
D_{\mathcal S,r,u}:(T\mathcal B_{\mathcal S|_U})_{(r,u)}
=T_r\mathcal R\times(T\mathcal B_{\mathcal S_r})_u\longrightarrow(\mathcal E_{\mathcal S_r})_u.
\end{equation}
第二个分量正是\eref{8.17}中取 $S=\mathcal S_r$ 的通常线性化算子 $D_{S,u}$；第一个分量则是在保持 $u$ 固定时，\eref{8.9}对 $r$ 变化的导数。注意，扩展算子仍是 Fredholm 算子。对于 $(r,u)$，适当的正则性概念是 $D_{\mathcal S,r,u}$ 为满射，
\SourcePage{131}
而不是要求更强的 $D_{\mathcal S_r,u}$ 为满射；后者是将 $u$ 视为定义域固定为 $\mathcal S_r$ 的非齐次伪全纯映射时的正则性。和通常一样，若所有相关 $(r,u)$ 都满足正则性条件，就称模空间 $\M_{\mathcal S}(\{y_\zeta\})$ 或扰动数据 $(K,J)$ \emph{正则}。\znote{123}{这里要用族版本的满射性，而不是单个纤维上的满射性，因为 $r$ 本身是未知量：\eref{9.18}的定义域里含有因子 $T_r\mathcal R$，若只要求固定 $r$ 时的 $D_{\mathcal S_r,u}$ 满射，就漏掉了族沿底空间方向的形变。换句话说，这是把 $\overline\partial$ 所定义的截面与“底空间方向”作横截，属于横截性理论中“一个映射对一族参数横截”的标准情形；横截性只要求像空间被微分像张满，并不要求两个切空间互补。[来源：nLab "transversal maps"，https:/\allowbreak /\allowbreak ncatlab.org/\allowbreak nlab/\allowbreak show/\allowbreak transversal+\allowbreak maps（2026-11-15 访问）。]}

\begin{note}\label{st:9.4}
这里的表述避开了一个技术困难：只要考虑定义域变化的伪全纯映射，就会出现这个问题。对整个族 $\mathcal S\to\mathcal R$，没有典范定义的 Banach 流形 $\mathcal B_{\mathcal S}$，因为平凡化的变换在局部定义的空间 $\mathcal B_{\mathcal S|_U}$ 之间只诱导 Banach 同胚，而不诱导可微映射。然而，椭圆正则性保证，算子\eref{9.18}以及正则情形下 $\M_{\mathcal S}(\{y_\zeta\})$ 的流形结构，都不依赖于局部平凡化。
\end{note}\znote{123}{过渡映射只是 Banach 同胚而不必是微分同胚，因此无法用它们把各局部空间 $\mathcal B_{\mathcal S|_U}$ 的切丛粘成整个 $\mathcal B_{\mathcal S}$ 的切丛；这就是不存在典范 Banach 流形的原因。但\eref{9.18}仍是有意义的 Banach 丛映射，其 Fredholm 性与满射性在 Banach 同胚下拉回不变；于是零点集在每个局部平凡化下都有光滑结构，而这些结构在重叠处一致。实践中只需固定一个局部平凡化作计算，结论与选择无关。[来源：原书 §(9h)，即\stref{笔记}{9.4}的完整论证。]}

\subsection{一致地选择扰动数据}\label{sec:9i}
现在需要简要回顾粘合构造。假设有一棵 $d$ 叶树 $T$，并对每个顶点 $v\in\operatorname{Ve}(T)$，给定带条带状端和 Lagrangian 标记的 $|v|$ 标点圆盘 $S_v$。回顾旗标 $f=(v,e)$ 与 $S_v$ 的无穷远点之间存在双射。因此，在这里，每个旗标对应一对正合 Lagrangian 子流形。若对每条 $e\in\operatorname{Ed}^{\mathrm{int}}(T)$，两个旗标 $f^\pm(e)=(v^\mp(e),e)$ 对应同一对子流形，就称各 $S_v$ 的\emph{标记相容}。可以给 $\R^2\setminus T$ 的各连通分支赋予正合 Lagrangian 子流形作为标记，以直观表示这种情形，见\sr{fig:9.4}{图~9.4}。由各 $S_v$ 粘合而成的 $(d+1)$ 标点圆盘族 $\mathcal S$，以显然的方式继承 Lagrangian 标记。

进一步假设，各 $S_v$ 配备扰动数据 $(K_v,J_v)$，使每对端 $f^\pm(e)$ 上对应的 Floer 数据相同。于是，通过粘合构造，它们在 $\mathcal S$ 上确定扰动数据 $(K,J)$。这组数据光滑延伸到自然的部分紧化 $\overline{\mathcal S}$；也就是说，存在
\begin{equation}\label{eq:9.19}
\overline K\in\Omega^1_{\overline{\mathcal S}/\overline{\mathcal R}}(\overline{\mathcal S},\mathcal H),\qquad
\overline J\in C^\infty(\overline{\mathcal S},\mathcal J),
\end{equation}
在 $\mathcal S$ 上的限制为 $(K,J)$。由于 $\overline{\mathcal S}\to\overline{\mathcal R}$ 是浸没映射，在部分紧化上，逐纤维一形式的空间仍有意义。在退化程度最大的点 $0\in\overline{\mathcal R}$ 上，纤维为 $\mathcal S_0=\coprod_vS_v$，而 $(\overline K,\overline J)$ 正是给定数据 $(K_v,J_v)$ 的不交并。最后，如果粘合的是族 $\mathcal S_v$ 而不是单个曲面，上述所有结论仍然成立。

固定条带状端的万有选择。此外，对所考虑的辛流形 $M$ 中每一对正合 Lagrangian 子流形 $(L_0,L_1)$，固定 Floer 数据 $(H_{L_0,L_1},J_{L_0,L_1})$。
\begin{definition*}
\emph{扰动数据的万有选择}，是对每个 $d\geq2$ 及每个正合 Lagrangian 子流形的 $(d+1)$ 元组 $(L_0,\ldots,L_d)$，在带有这些 Lagrangian 标记的族 $\mathcal S^{d+1}\to\mathcal R^{d+1}$ 上选择扰动数据
\begin{equation}\label{eq:9.20}
(K_{L_0,\ldots,L_d},J_{L_0,\ldots,L_d}),
\end{equation}
\SourcePage{132}
使其在该族的条带状端上化为子流形对\eref{9.5}的 Floer 数据。
\end{definition*}
\SourceFigure{figure-9-4}{0.70407616}{9.4}
从现在开始，假设条带状端的万有选择是一致的。假设已经作了扰动数据的万有选择。取稳定 $d$ 叶树 $T$，将万有族 $\mathcal S^{|v|}$ 配给其各顶点，并赋予相容的 Lagrangian 标记。于是粘合后的族 $\mathcal S\to\mathcal R$ 带有两组扰动数据，它们先验地可能不同；更精确地，它们在条带状端上总是相同，但在曲面的其余部分可能不同。第一组通过粘合构造从各 $\mathcal S^{|v|}$ 上的扰动数据诱导而来；第二组通过分类映射 $\gamma^T$，从带有正确标记的 $\mathcal S^{d+1}$ 上的扰动数据万有选择拉回。
\begin{definition*}
若在所有上述情形，均满足以下条件，就称扰动数据的万有选择\emph{一致}：
\begin{itemize}
\item 存在粘合参数充分小的子集 $U\subset\mathcal R$，使上述两组扰动数据在曲面 $\mathcal S_r$（$r\in U$）的薄部相同。这里指的是粘合过程中产生的厚薄分解，见
\SourcePage{133}
\stref{笔记}{9.1}。
\item 设 $(K,J)$ 是 $\mathcal S$ 上的第一组扰动数据，即由粘合得到的数据；$(\overline K,\overline J)$ 是其在部分紧化 $\overline{\mathcal S}$ 上如\eref{9.19}的延伸。要求另一组数据，即从 $\mathcal S^{d+1}$ 拉回的数据，也光滑延伸到 $\overline{\mathcal S}$；并且在所有粘合参数都为零的子集 $\{0\}^{\operatorname{Ed}^{\mathrm{int}}(T)}\times\mathcal R^T\subset\overline{\mathcal R}$ 上，这个延伸与 $(\overline K,\overline J)$ 相同。
\end{itemize}

\end{definition*}
为了更好地理解一致性的含义，假设有序列 $r_1,r_2,\ldots\in\mathcal R^{d+1}$，收敛于 $\overline{\mathcal R}^{d+1}\setminus\mathcal R^{d+1}$ 中 $T$-层的一点。当 $k$ 充分大时，可以写成
\begin{equation}\label{eq:9.21}
r_k=\gamma^T(\{\rho_{e,k}\},\{r_{v,k}\}),
\end{equation}
其中 $r_{v,k}\in\mathcal R^{|v|}$，$v\in\operatorname{Ve}(T)$，而 $\rho_{e,k}\in(-1,0)$，$e\in\operatorname{Ed}^{\mathrm{int}}(T)$。当 $k\to\infty$ 时，每个序列 $r_{v,k}$ 收敛，且 $\rho_{e,k}$ 趋于零。更几何地，考虑带标点圆盘
\[
S_k=\mathcal S_{r_k}^{d+1},\qquad S_{v,k}=\mathcal S_{r_{v,k}}^{|v|}.
\]
那么\eref{9.21}说明：\allowbreak 按 $T$ 的边所指定的方式粘合各 $S_{v,k}$，在相邻圆盘间留下有限条带 $[0,-\log(-\rho_{e,k})/\pi]\times[0,1]$；这些条带在 $k\to\infty$ 时越来越长。所得曲面与 $S_k$ 同构。为 $S_k$ 配备某组 Lagrangian 标记 $(L_0,\ldots,L_d)$，为 $S_{v,k}$ 配备相应的相容标记。按定义，扰动数据的万有选择在每个带 Lagrangian 标记的稳定带标点圆盘上给出指定的扰动数据，因此也在这里的 $S_k$ 和 $S_{v,k}$ 上给出数据。

一致性表示：$k$ 充分大时，$S_k$ 上的扰动数据“收敛于” $\coprod_vS_{v,k}$ 上的数据。这个表述是有意义的，因为按定义，$S_k$ 是 $\coprod_vS_{v,k}$ 的某个子集的商。此外，在 $S_k$ 的有限条带以及条带状端上，这种收敛是平凡的：当 $k\gg0$ 时，那里的扰动数据就等于 $\coprod_vS_{v,k}$ 的对应部分上的数据。\znote{125}{一致性的用处是让“取极限”与“计数”可以交换：$p\leq1$ 维模空间的边界点必落在某个退化层上，按\eref{9.21}把该点附近的解展开，得到的正是 $T$ 各顶点上较低阶问题的解；只有两组扰动数据在薄部相同，这个展开才给出\eref{9.23}中所需的系数，$\Ai$ 结合律\eref{1.2}才在链层成立。条件要求对一切 $d$ 同时成立，而不是只对 $p\leq1$ 的情形，原因是同一个族既出现在低维模空间的边界上，也出现在其内部的高维模空间中。[来源：原书 §(9l) 的紧致性论证，尤其\eref{9.29}、\eref{9.30}；另一套处理见 §(10e)。]}
\begin{lemma}\label{st:9.5}
存在一致的扰动数据万有选择。
\end{lemma}
\begin{proof}
证明采用与\stref{引理}{9.3}同类的归纳论证。这里只说明证明中的一点。一般而言，曲面 $S=\mathcal S_r^{d_0+1}$ 可以同时位于不止一个粘合映射 $\gamma^T$ 的像中，而这些映射给出不同的厚薄分解。然而，可以为每棵 $T$ 选择粘合参数充分小的开子集 $U^T\subset(-1,0)^{\operatorname{Ed}^{\mathrm{int}}(T)}\times\mathcal R^T$，使得对每个 $r\in\mathcal R^{d_0+1}$，存在满足 $r\in\gamma^T(U^T)$ 的极大树 $T$。这里的极大性意指，其他具有该性质的树 $\overline T$，都能由压缩 $T$ 的边得到。条带状端选择的一致性意味着：$\gamma^T$ 所诱导的 $S$ 的厚薄分解是极大的；其他满足 $r\in\gamma^{\overline T}(U^{\overline T})$ 的 $\gamma^{\overline T}$ 所给出的分解，由这个分解忘掉若干有限薄条带而得。对于我们的一致性要求，极大
\SourcePage{134}
分解对扰动数据施加最严格的条件。

这对\stref{引理}{9.5}的证明策略有如下启示。假设已经为 $d<d_0$ 选好满足一致性条件的扰动数据\eref{9.20}。那么，对 $\mathcal S^{d_0+1}$ 的构造，应从最复杂、边数最多的树 $T$ 所对应的子集 $\gamma^T(U^T)$ 开始；它们对应于 $\overline{\mathcal R}^{d_0+1}$ 的顶点。随后按紧化空间边界层的维数递增继续构造。
\end{proof}

\begin{note}\label{st:9.6}
乍看起来，尤其将以上讨论与\secref{9g}的对应讨论比较时，似乎施加更强的一致性条件更自然：要求粘合后的族 $\mathcal S\to\mathcal R$ 上的两组扰动数据，在粘合参数充分小时完全相同。为看出它的缺点，取一棵有两个顶点、每个顶点度数均为 $3$ 的树 $T$。为每个顶点配给族 $\mathcal S^3\to\mathcal R^3$；实际上它只包含唯一的三标点圆盘 $S$。采用\eref{8.5}的记号，粘合所得的四标点圆盘为 $S\#_lS$。令所有标记都等于同一个正合 Lagrangian 子流形。如果万有扰动数据满足更强的一致性条件，它就在圆盘 $S\#_lS$ 上赋予两“半”完全相同的扰动数据；更精确地说，是厚部的两个连通分支上的数据相同。在这类扰动数据中，试图使非齐次伪全纯映射模空间正则时，就会出现类似多重覆盖的问题。我们对一致性的定义避开了这个困难。
\end{note}

\subsection{定义}\label{sec:9j}
为定义带角正合辛流形 $M$ 的 Fukaya 范畴初步版本 $\F(M)^{\mathrm{pr}}$，需要选取以下额外数据：
\begin{itemize}
\item 一致的条带状端万有选择 $\{\epsilon_k^{d+1}\}$；
\item 对 $M$ 中每一对正合 Lagrangian 子流形 $(L_0,L_1)$ 选定 Floer 数据 $(H_{L_0,L_1},J_{L_0,L_1})$；
\item 一致的扰动数据万有选择 $(K_{L_0,\ldots,L_d},J_{L_0,\ldots,L_d})$，与先前选择的条带状端及 Floer 数据相容。
\end{itemize}
这些数据还须满足一个正则性条件，现在来说明。对于正合 Lagrangian 子流形的每个 $(d+1)$ 元组 $(L_0,\ldots,L_d)$（$d\geq2$），以及点 $y_0\in\mathcal C(L_0,L_d)$、$y_k\in\mathcal C(L_{k-1},L_k)$（$1\leq k\leq d$），考虑
\begin{equation}\label{eq:9.22}
\M^{d+1}(y_0,\ldots,y_d)=\M_{\mathcal S^{d+1}}(y_0,\ldots,y_d).
\end{equation}
这是族 $\mathcal S^{d+1}\to\mathcal R^{d+1}$ 上方程\eref{8.9}的解所成的模空间\eref{9.17}；该族带有 Lagrangian 标记 $(L_0,\ldots,L_d)$ 和已选定的扰动数据，渐近条件是在条带状端上收敛于 $y_k$。
\begin{definition*}
称\eref{9.22}中的点 $(r\in\mathcal R^{d+1},u:\mathcal S_r^{d+1}\to M)$ 为\emph{非齐次伪全纯多边形}。将记号延伸到 $d=1$，令
\SourcePage{135}
$\M^2(y_0,y_1)=\M_Z^*(y_0,y_1)$，即 Floer 方程的解模去平移后的模空间。
\end{definition*}
所要求的条件是：这些模空间全部正则。对于 $d=1$，已经提到，这对一般位置的 Floer 数据选择成立。对于 $d\geq2$，在一致的扰动数据万有选择这一类别中，同样的结论成立。
\begin{definition*}
$\F(M)^{\mathrm{pr}}$ 的对象是正合 Lagrangian 子流形 $L\subset M$，态射由 Floer 上链空间给出：
\[
\hom_{\F(M)^{\mathrm{pr}}}(L_0,L_1)=CF^{\mathrm{pr}}(L_0,L_1).
\]
一阶映射 $\mu^1=\partial$ 是边界算子，像\eref{8.19}一样，由计数 $\M^2(y_0,y_1)$ 中的孤立点定义。高阶复合类似地定义为
\begin{equation}\label{eq:9.23}
\mu^d(y_d,\ldots,y_1)=\sum_{y_0}\#\M^{d+1}(y_0,\ldots,y_d)\,y_0.
\end{equation}

\end{definition*}
\znote{127}{这个定义只覆盖了“横截”的情形：态射空间 $CF^{\mathrm{pr}}(L_0,L_1)$ 由交点生成，而交点本身依赖于所选的 Floer 数据，也就是依赖于对 $L_0,L_1$ 所作的 Hamilton 扰动。严格说来，$\mu^d$ 不是 $(L_0,\ldots,L_d)$ 的函数，而是后者连同所选数据的函数，§(10a) 用总范畴与协调系统来处理这种依赖。当两个 Lagrangian 子流形不横截相交时，交点不是离散集，计数失去意义；一般情形下 Fukaya 范畴的定义之所以困难，症结正在这里。[来源：nLab "Fukaya category"，https:/\allowbreak /\allowbreak ncatlab.org/\allowbreak nlab/\allowbreak show/\allowbreak Fukaya+\allowbreak category（2026-11-15 访问）。]}
紧致性及粘合论证表明，这一定义良好：每个 $\M^{d+1}(y_0,\ldots,y_d)$ 的零维部分都是有限集；并且满足 $\Ai$ 结合律方程\eref{1.2}。\znote{127}{这里的 $\#$ 指零维模空间的点数按系数域 $\K$ 取的模 $2$ 计数（\stref{笔记}{9.7}始终假定 $\operatorname{char}\K=2$），因此\eref{9.23}的右端是有限和。在分次版本中，$\mu^d$ 的次数为 $2-d$，于是\eref{1.2}中每一项的总次数一致；这与 $\Ai$ 代数通常的约定相同，即 $D_k$ 的次数为 $2-k$（$D_1$ 是微分、$D_2$ 是乘法、$D_3$ 是结合子）。[来源：nLab "A-∞ algebra"，https:/\allowbreak /\allowbreak ncatlab.org/\allowbreak nlab/\allowbreak show/\allowbreak A-infinity-algebra（2026-11-15 访问）。]}因此，$\F(M)^{\mathrm{pr}}$ 是一个不分次的 $\Ai$ 范畴。按定义，$\mu^2$ 在上同调上诱导的映射是三角形乘积\eref{8.3}，所以底层的上同调范畴确实就是先前定义的 $H\F(M)^{\mathrm{pr}}$。这也说明 $\F(M)^{\mathrm{pr}}$ 上同调含幺。\znote{127}{“上同调含幺”指 $H^0(\hom(L,L))$ 中存在元素 $e_L$，使 $\mu^2(e_L,a)=a$、$\mu^2(a,e_L)=a$ 在上同调层成立。它弱于链层的含幺性（存在闭元素 $e_L$ 使 $\mu^2(e_L,-)=\mu^2(-,e_L)=\id$ 在链层成立），因为 $\mu^2$ 只在上同调层才结合，链层上结合律只能到同伦。这里的 $e_L$ 来自 $L$ 与它的小 Hamilton 扰动像的横截交点中最靠近对角线的那一点。[来源：原书 §(1a) 关于上同调范畴与“非含幺”术语的约定，见第一、三处定义。]}

\begin{note}\label{st:9.7}
我们没有明确地定义不分次 $\Ai$ 范畴，但其含义应当清楚。回顾，假设系数域满足 $\operatorname{char}(\K)=2$，所以定义符号时不需要分次。有了这一点，\sr{chapter:I}{第一章}的初等理论只需作表面上的调整便仍然适用，唯一的例外是\stref{推论}{5.8}。读者应当记住，目前的设置只是暂时的简化：最终所有应用都会在分次有良好定义的背景下进行。
\end{note}

\begin{note}\label{st:9.8}
上一段概述的结果为本领域专家所熟知，证明虽然有些长，但并未超出伪全纯曲线理论的标准技术。在本节余下部分，将简要勾勒主要步骤，希望由此传达这些技术的大致面貌。对希望深入细节的读者，以下是一些可能的参考指引。Riemann 曲面族的非齐次项似乎最早用于\bcite{116}。横截性论证的基本参考是\bcite{49}。Floer 的原始论文\bcite{46}、\bcite{45}、\bcite{44}，以及\bcite{103}、\bcite{50}，包含处理 Lagrangian 边界条件所需的一些基本分析。空间 $\mathcal R^{d+1}$ 及其紧化在\bcite{58}中讨论，该文还包含许多与 Fukaya 范畴有关的材料。整套构造中最复杂的部分，
\SourcePage{136}
可能是用来证明 $\Ai$ 结合律方程成立的粘合过程；见\bcite{60}，第 18 节。这个论证推广了 Floer 理论中更熟悉的 $\partial^2=0$ 的证明；其 Lagrangian 版本和 Hamilton 版本分别见\bcite{44}和\bcite{119}。
\end{note}

\subsection{横截性}\label{sec:9k}
固定一致的条带状端万有选择，并为每一对正合 Lagrangian 子流形固定 Floer 数据。假设已经为所有 $d\leq d_0$ 选取扰动数据\eref{9.20}。取一组 Lagrangian 标记 $(L_0,\ldots,L_{d_0})$。希望稍作修改给定扰动数据 $(K,J)=(K_{L_0,\ldots,L_{d_0}},J_{L_0,\ldots,L_{d_0}})$，使对应模空间\eref{9.22}正则。形式上，$(K,J)$ 的无穷小变形由一对
\begin{equation}\label{eq:9.24}
(\delta K,\delta J)\in\Omega^1_{\mathcal S^{d_0+1}/\mathcal R^{d_0+1}}(\mathcal S^{d_0+1},\mathcal H)
\times C^\infty(\mathcal S^{d_0+1},T_J\mathcal J)
\end{equation}
给出，要求 $(\delta K,\delta J)$ 在每个曲面 $\mathcal S_r^{d_0+1}$ 的条带状端上为零；并且对连通分支 $C_k\subset\partial\mathcal S_r^{d_0+1}$ 的每个切向量 $\xi$，有 $\delta K(\xi)|_{L_k}=0$。这样的无穷小变形很容易“取指数”成为实际的变形：
\begin{equation}\label{eq:9.25}
(K',J')=(K+\delta K,\ J\exp(-J\cdot\delta J)).
\end{equation}
为了保持与 $d<d_0$ 时的扰动数据选择的一致性，需要对容许的变形空间作一些限制。对每棵稳定 $d_0$ 叶树 $T$，固定粘合参数充分小的开子集 $U^T\subset(-1,0)^{\operatorname{Ed}^{\mathrm{int}}(T)}\times\mathcal R^T$。只要各 $U^T$ 取得充分小，就可选开子集 $\Omega\subset\mathcal S^{d_0+1}$，使对每个 $r\in\mathcal R^{d_0+1}$，交集 $\Omega_r=\Omega\cap\mathcal S_r^{d_0+1}$ 非空，并具有以下性质：$\Omega_r$ 与 $\mathcal S_r^{d_0+1}$ 的条带状端不交；若对某棵 $T$ 有 $r\in\gamma^T(U^T)$，则 $\Omega_r$ 位于 $\mathcal S_r^{d_0+1}$ 对应分解的厚部。\znote{128}{把 $\Omega_r$ 限制在厚部，是因为薄部与条带状端上的扰动数据已被低阶数据决定：按 §(9i) 的一致性条件，若 $r$ 落在某个 $\gamma^T(U^T)$ 中，$\mathcal S_r^{d_0+1}$ 就带有由 $d<d_0$ 阶数据粘合而成的结构，那里不能自由改动。厚部是圆盘“真正有 $d_0+1$ 个标记”的部分，扰动数据在那里仍可任意选取；\eref{9.27}之后用 $\delta$ 函数逼近的论证正依赖这一点。[来源：原书 §(9i) 的一致性条件与 §(9e) 的厚薄分解，见\eref{9.19}、\eref{9.20}。]}

考虑\eref{9.24}中各对 $(\delta K,\delta J)$ 所成的子空间 $\mathcal T$，要求它们在 $\Omega$ 外为零，并相对于部分紧化 $\overline{\mathcal S}^{d_0+1}$ 渐近衰减。最后一个条件意味着，$(\delta K,\delta J)$ 光滑延伸为 $\overline{\mathcal S}^{d_0+1}$ 上的一对 $(\overline{\delta K},\overline{\delta J})$，它及其任意阶导数都在 $\overline{\mathcal S}^{d_0+1}\setminus\mathcal S^{d_0+1}$ 上为零。对于每个这样的 $(\delta K,\delta J)$，变形后的扰动数据\eref{9.25}仍满足关于粘合映射 $\gamma^T$ 的一致性。具体说，因 $(K',J')=(K,J)$ 在 $\Omega$ 外成立，曲面 $\mathcal S_r^{d_0+1}$ 薄部上的条件仍满足；一致性要求的另一部分涉及相对于 $\overline{\mathcal S}^{d_0+1}$ 的渐近行为，而 $(K',J')$ 与原 $(K,J)$ 的渐近行为也相同。

\begin{proof}[正则性论证]
还需证明，这类变形足以使带有标记 $(L_0,\ldots,L_{d_0})$ 的非齐次伪全纯 $(d_0+1)$ 边形模空间正则。这是 Sard–Smale 定理的标准应用。
\SourcePage{137}
对于原扰动数据 $(K,J)$ 的模空间中任意一点 $(r,u)\in\M^{d_0+1}(y_0,\ldots,y_{d_0})$，都有一个万有线性化算子 $D_{\mathcal T,\mathcal S^{d_0+1},r,u}$。它包含\eref{9.18}，但另有变量，用来表示修改 $(K,J)$ 时 $(du-Y)^{0,1}$ 如何变化。为严格起见，应当为族 $\mathcal S^{d_0+1}$ 选择局部平凡化，见\stref{笔记}{9.4}；为简便起见，这里略去。具体公式为
\begin{equation}\label{eq:9.26}
\begin{aligned}
\overline D_{\mathcal S^{d_0+1},r,u}:&\ \mathcal T\times T_r\mathcal R^{d_0+1}
\times(T\mathcal B_{\mathcal S_r^{d_0+1}})_u\longrightarrow(\mathcal E_{\mathcal S^{d_0+1},r})_u,\\
(\delta K,\delta J,\rho,X)\longmapsto{}&-(\delta Y)^{0,1}
+\tfrac12\delta J\circ(du-Y)\circ I_S+D_{\mathcal S^{d_0+1},r,u}(\rho,X),
\end{aligned}
\end{equation}
其中 $\delta Y$ 是由 $\delta K$ 得到的、取值于 Hamilton 向量场的逐纤维一形式，$I_S$ 是 $S=\mathcal S_r^{d_0+1}$ 上的复结构。主要步骤是证明\eref{9.26}为满射。注意，由于最后一项是 Fredholm 算子，像自动闭合。\eref{9.26}余核中的元素可表示为
\[
\chi\in L^q(S,(\Omega_S^{0,1}\otimes u^*TM)^\vee),\qquad 1/p+1/q=1,
\]
特别地，它必须满足 $D_{S,u}^*\chi=0$。由椭圆正则性，$\chi$ 至少在 $\operatorname{int}(S)$ 中光滑；更仔细考察其性质，实际上可证明它直到边界都光滑，但这里不需要这一点。此外，对所有 $\delta Y$，它满足
\begin{equation}\label{eq:9.27}
\int_S\langle(\delta Y)^{0,1},\chi\rangle=0.
\end{equation}
由构造，存在非空开集 $\Omega_r\subset S$，在其中 $\delta K$ 的选择基本上任意。特别地，对于任意 $z\in\Omega_r\cap\operatorname{int}(S)$ 及任意 $\alpha\in\Omega^{0,1}(S)_z\otimes T M_{u(z)}$，存在一列 $(\delta K,\delta J)$，使相应的 $(\delta Y)^{0,1}$ 弱收敛于 $\delta$ 函数 $\delta_z\otimes\alpha$。将这一列代入\eref{9.27}，可见 $\chi_z=0$。因此，$\chi$ 在 $\Omega_r$ 上为零；它处处为零，则由方程
\[
D_{S,u}^*\chi=0
\]
的解的唯一延拓性质推出。
\end{proof}

\begin{note}\label{st:9.9}
刚才给出的证明概要，至少略去了一个例行却非平凡的技术点。容许的 $(\delta J,\delta K)$ 所成的空间 $\mathcal T$ 是 Fréchet 空间。为了应用 Sard–Smale 定理，需要进一步限制到一个稠密 Banach 子空间，例如采用 Floer 的 $C_\epsilon^\infty$ 拓扑的一种版本，将关于 $\overline{\mathcal S}^{d_0+1}$ 所要求的渐近行为计入其中。这样，\eref{9.26}的定义域在前两个变量上缩减为稠密子空间。然而，因为仍保留第三部分所给出的 Fredholm 算子，满射性不受影响。
\end{note}

\clearpage
\subsection{紧致性}\label{sec:9l}
非齐次伪全纯多边形的模空间有自然的 Gromov 紧化，其结构在一定程度上与 $\mathcal R^{d+1}\subset\overline{\mathcal R}^{d+1}$ 的结构相仿。实际上，后者可以视为
\SourcePage{138}
$M=L_0=\cdots=L_d=\text{一点}$ 的平凡情形。下面先在集合意义下定义这些紧化，不加证明地陈述其主要拓扑性质，再说明这如何推出各映射 $\mu^d$ 构成一个 $\Ai$ 结构。

像\eref{9.22}的定义中那样固定 $L_k,y_k$，再固定一棵半稳定 $d$ 叶树 $T$。假设对每个顶点 $v\in\operatorname{Ve}(T)$ 都给定一个 $|v|$ 标点圆盘 $S_v$。若 $|v|\geq3$，该圆盘确定一点 $r_v\in\mathcal R^{|v|}$，并存在唯一同构 $S_v\cong\mathcal S_{r_v}^{|v|}$。在余下的 $|v|=2$ 情形，有 $S_v\cong Z$，但这一同构只在模去 $s$ 方向平移的意义下确定。此外，要求这些圆盘带有 Lagrangian 标记，使每个旗标 $f=(v,e)$ 都有一对指定子流形 $(L_{f,0},L_{f,1})$。像前面讨论粘合时一样，标记应彼此相容，使粘合后的 $(d+1)$ 标点圆盘上的标记为给定的 $L_k$。最后，对每个 $f$，要求有一点 $\tilde y_f\in\mathcal C(L_{f,0},L_{f,1})$，满足 $\tilde y_{f^+(e)}=\tilde y_{f^-(e)}$，且与根、叶相邻的旗标所对应的点，按自然次序为 $y_0,\ldots,y_d$。
\begin{definition*}
\emph{非齐次伪全纯稳定多边形}，是一组映射 $u_v:S_v\to M$（$v\in\operatorname{Ve}(T)$），满足以下条件：
\begin{itemize}
\item 若 $|v|\geq3$，可将 $u_v$ 看成映射 $\mathcal S_{r_v}^{|v|}\to M$，要求它是对应的非齐次伪全纯映射方程的解，即
\[
(r_v,u_v)\in\M^{|v|}(\tilde y_{f_0(v)},\ldots,\tilde y_{f_{|v|-1}(v)}),
\]
其中 $f_0(v),\ldots,f_{|v|-1}(v)$ 是按自然次序排列的 $v$ 周围的旗标，扰动数据取自我们的万有选择。
\item 若 $|v|=2$，则 $u_v:Z\to M$ 应当代表
\[
u_v\in\M^2(\tilde y_{f_0(v)},\tilde y_{f_1(v)}),
\]
即它是给定 Floer 数据的 Floer 方程的非驻定解。因为该方程具有平移不变性，这一表述不依赖于同构 $S_v\cong Z$ 的选择。
\end{itemize}
\end{definition*}
定义 $\M^T(y_0,\ldots,y_d)$ 为所有这种非齐次伪全纯稳定多边形的等价类的集合。等价关系是各带标点圆盘 $S_v$ 的同构，并要求各映射 $u_v$ 在此同构下对应相等。具体说，若某个 $v$ 满足 $|v|=2$，对相应的 $u_v$ 作平移，并不改变点 $\{u_v\}\in\M^T(y_0,\ldots,y_d)$。在集合意义下，将 Gromov 紧化定义为对所有 $T$ 的不交并：
\begin{equation}\label{eq:9.28}
\overline\M^{d+1}(y_0,\ldots,y_d)=\coprod_T\M^T(y_0,\ldots,y_d).
\end{equation}
与\eref{9.11}不同，半稳定树 $T$ 有无限多棵，因为总能继续插入二价顶点。然而，无论怎样选取 $(y_0,\ldots,y_d)$，\eref{9.28}中实际非空的层
\SourcePage{139}
只有有限多个。为说明这一点，对各分量的能量\eref{8.12}求和，得到
\[
\sum_vE(u_v)=A_{H_{L_0,L_d}}(y_0)-\sum_{k=1}^dA_{H_{L_{k-1},L_k}}(y_k)+(\text{曲率项}).
\]
曲率项只由 $|v|\geq3$ 的顶点贡献；考察扰动数据万有选择的一致性条件，可以推出它在整个空间\eref{9.28}上有一致界。另一方面，存在 $\epsilon>0$，使每个 $|v|=2$ 的分量 $u_v$ 对能量至少贡献 $\epsilon$。具体说，考察 $0\leq i<j\leq d$ 时的所有元素对 $y_0,y_1\in\mathcal C(L_i,L_j)$，在所有为正的作用量差 $A_{H_{L_i,L_j}}(y_0)-A_{H_{L_i,L_j}}(y_1)$ 中取最小值作为 $\epsilon$。这两点合起来，就给出了\eref{9.28}中非空层所对应树 $T$ 的二价顶点数的上界。\znote{139}{$\epsilon>0$ 的存在性用到交点集的有限性：$\mathcal C(L_i,L_j)$ 是紧致 Lagrangian 子流形的横截交点集，故为有限集，于是正的作用量差只有有限多个，其最小值即为 $\epsilon$。这个下界的作用是给二价顶点数一个与树无关的上界：每插入一个二价顶点至少消耗能量 $\epsilon$，而总能量被右端的作用量差与曲率项的一致界控制，从而说明\eref{9.28}中非空层只有有限多个，$\mu^d$ 的计数才是有限和。[来源：原书 §(9l)，能量恒等式见\eref{8.12}。]}

从现在起，假设所有相关模空间都正则。用 $\M^{d+1}(y_0,\ldots,y_d)^p$ 表示 $\M^{d+1}(y_0,\ldots,y_d)$ 的 $p$ 维部分，它是该空间中的开闭子集；类似地定义 $\M^T(y_0,\ldots,y_d)^p$。\eref{9.28}上存在自然的稳定映射拓扑，在这个拓扑下整个空间紧致。并且，每个
\begin{equation}\label{eq:9.29}
\overline\M^{d+1}(y_0,\ldots,y_d)^p
=\coprod_T\M^T(y_0,\ldots,y_d)^{p-|\operatorname{Ed}^{\mathrm{int}}(T)|}
\end{equation}
都是开闭子集，同时是 $p$ 维带角拓扑流形。更精确地，单顶点树对应的部分可与 $\M^{d+1}(y_0,\ldots,y_d)^p$ 等同，它构成内部；其余部分都是余维数为 $|\operatorname{Ed}^{\mathrm{int}}(T)|$ 的边界层。\eref{9.29}的可微结构问题相当微妙：每一层当然都是光滑流形，而且可以定义产生各边界层领状邻域的粘合映射。但是，粘合参数的选择并不那么明确，因此这些映射可能不像 $\overline{\mathcal R}^{d+1}$ 的情形那样很好地拼合。幸好，构造 Fukaya 范畴只用到维数 $p\leq1$ 的模空间，此时这个问题无关紧要。

先考虑 $p=0$，可见只有单顶点树对\eref{9.29}有贡献。该空间的紧致性意味着 $\M^{d+1}(y_0,\ldots,y_d)^0$ 是有限集，从而说明定义\eref{9.23}中的点计数是合理的。再看 $p=1$，边界点来自恰有两个顶点的半稳定树，即\sr{fig:9.5}{图~9.5} 中的 $T=T(m,n)$。设
\begin{equation}\label{eq:9.30}
\begin{aligned}
(r_1,u_1,r_2,u_2)\in{}&\M^{d+2-m}(y_0,y_1,\ldots,y_n,\tilde y,y_{n+m+1},\ldots,y_d)^0\\
&\times\M^{m+1}(\tilde y,y_{n+1},\ldots,y_{n+m})^0
\end{aligned}
\end{equation}
\SourcePage{140}
是 $T(m,n)$-层中的一点；这意味着 $\tilde y\in\mathcal C(L_n,L_{n+m})$。对所有可能的 $\tilde y$ 取并后，这样的 $(r_1,u_1,r_2,u_2)$ 的个数模 $2$，恰好是表达式
\[
\mu^{d+1-m}(y_d,\ldots,y_{n+m+1},\mu^m(y_{n+m},\ldots,y_{n+1}),y_n,\ldots,y_1)
\]
中 $y_0$ 的系数。对 $m,n$ 求和便计入全部边界点，因而必得到偶数；这正是以 $\K$ 为系数的 $\Ai$ 结合律方程\eref{1.2}。
\SourceFigure{figure-9-5}{0.46730719}{9.5}
从更具体的角度考察最后这个 Gromov 紧化情形，是很有启发性的。假设有序列 $(r^i,u^i)\in\M^{d+1}(y_0,\ldots,y_d)^1$，收敛于某个边界点\eref{9.30}。这里包含两种相关但不同的现象：
\begin{itemize}
\item 当 $m=1$ 或 $m=d$ 时，是\emph{不稳定情形}。为简化记号，只讨论前者。$u^i$ 是映射 $\mathcal S_{r^i}^{d+1}\to M$，$u_1$ 是映射 $\mathcal S_{r_1}^{d+1}\to M$，其中 $r^i,r_1\in\mathcal R^{d+1}$。事实上，$r^i\to r_1$，而 $u^i$ 在紧致子集上收敛于 $u_1$。另一分量 $u_2\in\M^2(\tilde y,y_{n+1})$ 是 Floer 方程的解，它以如下方式表现为“无穷远处的气泡”。令
\[
\epsilon^i=\epsilon_{n+1}^{d+1}|_{\{r^i\}\times Z^+}:Z^+\longrightarrow\mathcal S_{r^i}^{d+1}
\]
为 $\mathcal S_{r^i}^{d+1}$ 的第 $n+1$ 个出向条带状端。则存在 $\sigma^i\to+\infty$，使重新参数化的映射 $\tilde u^i(s,t)=u^i(\epsilon^i(s+\sigma^i,t))$ 在紧致子集上收敛于 $u_2$。这个表述有意义，因为由对任意扰动数据施加的条件\eref{8.7}，各 $\tilde u^i$ 都是 Floer 方程的局部解，并且当 $i\to\infty$ 时，它们的定义域穷竭整个 $Z$。
\item 在余下的\emph{稳定情形} $1<m<d$，$r^i$ 收敛于 $\partial\overline{\mathcal R}^{d+1}$ 中的点 $(r_1,r_2)$。按拓扑的定义，这意味着存在收敛
\SourcePage{141}
序列 $r_k^i\to r_k$ 以及粘合参数序列 $\rho^i\to0$，使得对充分大的 $i$，有 $r^i=\gamma^T(\rho^i,r_1^i,r_2^i)$。将对应的带标点圆盘记为
\[
S_1^i=\mathcal S_{r_1^i}^{d+2-m},\qquad S_2^i=\mathcal S_{r_2^i}^{m+1}.
\]
于是 $S^i$ 由 $S_1^i,S_2^i$ 按\eref{8.5}粘合而成，粘合长度为 $l^i=-\log(-\rho^i)/\pi$。把 $u^i$ 限制到 $S^i$ 中由粘合继承的各部分，得到 $S_1,S_2$ 上部分定义的映射序列，它们在紧致子集上分别收敛于 $u_1,u_2$。扰动数据选择的一致性使这一表述有意义，因为它要求 $S^i$ 和 $S_k^i$ 上所用的扰动数据彼此关联。
\end{itemize}

\section{一些基本性质}\label{sec:10}
\subsection{良定义性}\label{sec:10a}
抽象地说，需要考虑以某个集合 $I$ 为指标、彼此等价的范畴所构成的系统。这个概念有以下三个版本，只有中间一个不是显然的。\znote{133}{三个版本的区别在于“复合的相容性”被要求到什么程度：严格系统直接令 $F^{i_2,i_0}$ 等于复合；弱系统只要求任意两个范畴等价，不指定等价的选择，因而没有任何相容性可言；协调系统则给定同构 $T$ 并要求\eref{10.1}交换。后者不显然，是因为两个复合同构之差一般给出一个障碍类（在群作用的情形即 §(10c) 讨论的障碍）；只有当 $T$ 由某个严格自然变换 $S$ 构造出来时（如本段的 $T=L_{K^{i_2}}(R_{H^{i_0}}(S^{i_1}))$），障碍才自动消失。[来源：原书 §(10a)，并见 §(10c) 的\stref{笔记}{10.4}。]}
\begin{definition*}
\begin{itemize}
\item \emph{严格系统}是一族范畴 $A^i$（$i\in I$），连同函子 $F^{i_1,i_0}:A^{i_0}\to A^{i_1}$，满足 $F^{ii}=\Id_{A^i}$ 及 $F^{i_2,i_1}\circ F^{i_1,i_0}=F^{i_2,i_0}$。
\item \emph{协调系统}由范畴 $A^i$、函子 $F^{i_1,i_0}$，以及函子同构 $T^{i_2,i_1,i_0}:F^{i_2,i_1}\circ F^{i_1,i_0}\to F^{i_2,i_0}$ 组成，满足两个条件：首先，$F^{ii}=\Id_{A^i}$；类似地，只要 $i_2=i_1$ 或 $i_1=i_0$，$T^{i_2,i_1,i_0}$ 就是恒等变换。其次，对所有 $(i_3,i_2,i_1,i_0)$，以下图交换：
\SourceEquationDiagram{diagram-10-1}{0.57115323}{10.1}
这里 $L,R$ 分别是作用于自然变换的左复合和右复合；它们是\secref{1e}中讨论的 $\Ai$ 构造的经典对应物。
\item \emph{弱系统}是一族范畴 $A^i$，任意两个都等价，但不指定等价的选择。
\end{itemize}

\end{definition*}
假设有范畴 $A^{\mathrm{tot}}$ 及一族全子范畴 $A^i$，使每个嵌入 $H^i:A^i\to A^{\mathrm{tot}}$ 都是等价。暂时固定某个 $i$，对每个
\SourcePage{142}
$X\in\Ob A^{\mathrm{tot}}$，选择一个与之同构的对象 $K^i(X)\in\Ob A^i$，以及指定的同构 $S^i\in\Hom_{A^{\mathrm{tot}}}(K^i(X),X)$；若 $X$ 已属于 $A^i$，就作显然的选择。很容易将 $K^i$ 赋予函子 $A^{\mathrm{tot}}\to A^i$ 的结构，使 $K^i\circ H^i=\Id_{A^i}$。于是 $S^i$ 成为自然变换 $H^i\circ K^i\to\Id_{A^{\mathrm{tot}}}$，并且 $R_{H^i}(S^i)$、$L_{K^i}(S^i)$ 都是平凡的恒等自然变换。对每个 $i$ 这样做，再令
\[
F^{i_1,i_0}=K^{i_1}\circ H^{i_0},\qquad
T^{i_2,i_1,i_0}=L_{K^{i_2}}\bigl(R_{H^{i_0}}(S^{i_1})\bigr),
\]
就能把各 $A^i$ 组成一个协调范畴系统。

严格系统的概念不作改变就适用于 $\Ai$ 范畴；弱系统也一样，只需将等价换为拟等价。对于协调系统，要求有 $\Ai$ 函子 $\F^{i_1,i_0}:\A^{i_0}\to\A^{i_1}$，满足 $\F^{ii}=\Id_{\A^i}$，并有自然变换的上同调类
\[
[T^{i_2,i_1,i_0}]\in\Hom_{H^0(\operatorname{fun}(\A^{i_0},\A^{i_2}))}
(\F^{i_2,i_1}\circ\F^{i_1,i_0},\F^{i_2,i_0}).
\]
同样，当 $i_2=i_1$ 或 $i_1=i_0$ 时，这个类应为恒等元 $[E_{\F^{i_2,i_0}}]$；而且\eref{10.1}的对应图，采用\secref{1e}所定义的左、右复合函子后，在 $H^0(\operatorname{fun}(\A^{i_0},\A^{i_3}))$ 中交换。

与前面一样，固定一个 $\Ai$ 范畴 $\A^{\mathrm{tot}}$ 及一族与它拟等价的全 $\Ai$ 子范畴 $\A^i$，即可构造协调系统。具体说，以 $\mathcal H^i$ 表示嵌入，先像\stref{定理}{2.9}之前的讨论那样，构造拟等价 $\mathcal K^i:\A^{\mathrm{tot}}\to\A^i$，使 $\mathcal K^i\circ\mathcal H^i=\Id_{\A^i}$。由\stref{引理}{1.7}可知，映射
\begin{equation}\label{eq:10.2}
\begin{aligned}
H(\LG_{\mathcal K^i}):&\ \Hom_{H(\operatorname{fun}(\A^{\mathrm{tot}},\A^{\mathrm{tot}}))}
(\mathcal H^i\circ\mathcal K^i,\Id_{\A^{\mathrm{tot}}})\\
&\longrightarrow\Hom_{H(\operatorname{fun}(\A^{\mathrm{tot}},\A^i))}(\mathcal K^i,\mathcal K^i)
\end{aligned}
\end{equation}
是同构。选择在\eref{10.2}下映到恒等自然变换 $[E_{\mathcal K^i}]$ 的类 $[S^i]$。先考察 $\LG_{\mathcal K^i}\circ\RG_{\mathcal H^i}=\RG_{\mathcal H^i}\circ\LG_{\mathcal K^i}$，再应用一次\stref{引理}{1.7}，可知 $H(\RG_{\mathcal H^i})([S^i])=[E_{\mathcal H^i}]$。其余构造与普通范畴的情形相同。

转向 Floer 理论。回顾，给定带角正合辛流形 $M$，$\F(M)^{\mathrm{pr}}$ 的定义涉及额外选择，见\secref{9j}的汇总。设 $I$ 是所有可能选择的集合，以 $\{\F(M)^{\mathrm{pr},i}\}$ 表示所得 $\Ai$ 范畴族。若转到上同调层次，则 $H\F(M)^{\mathrm{pr},i}$ 构成上述意义下的严格系统。这是 Floer 上同调及其 TQFT 结构的良定义性的结果，\secref{8j}已简要讨论过。

为了处理 $\F(M)^{\mathrm{pr}}$ 本身，方便采用上面描述的形式技巧。定义一个总 Fukaya 范畴 $\F(M)^{\mathrm{pr},\mathrm{tot}}$，其对象是由正合 Lagrangian 子流形 $L$ 与指标 $i\in I$ 组成的二元组 $(L,i)$。态射空间照常取 Floer 复形，对每一对 $((L_0,i_0),(L_1,i_1))$ 分别选择正则 Floer 数据，要求当 $i_0=i_1=i$ 时，数据与将 $(L_0,L_1)$
\SourcePage{143}
视为 $\F(M)^{\mathrm{pr},i}$ 中对象时相同。类似地，通过适当选择条带状端与扰动数据来定义 $d\geq2$ 阶复合，要求当涉及的所有对象都有同一个指标 $i$ 时，选择与 $\F(M)^{\mathrm{pr},i}$ 相同。这稍微修改了标准设置，因为允许 $(d+1)$ 穿孔圆盘上的条带状端依赖于标记 $(L_k,i_k)$。然而，利用\secref{9i}中相容标记族的粘合，很容易适当调整\secref{9g}的一致性条件。结果是，总范畴带有典范的全忠实嵌入
\begin{equation}\label{eq:10.3}
\F(M)^{\mathrm{pr},i}\longrightarrow\F(M)^{\mathrm{pr},\mathrm{tot}}.
\end{equation}
对于任意固定的正合 Lagrangian 子流形 $L$ 及指标 $(i_0,i_1)$，对象 $(L,i_0)$、$(L,i_1)$ 在 $H\F(M)^{\mathrm{pr},\mathrm{tot}}$ 中同构。这是上同调层次的陈述，其证明归结为标准的延拓映射论证。因此，嵌入\eref{10.3}都是拟等价。于是可以利用这一点，把 $\{\F(M)^{\mathrm{pr},i}\}$ 组成 $\Ai$ 范畴的协调系统。严格说，它们是不分次的 $\Ai$ 范畴；不过\stref{笔记}{9.7}适用于这里，也适用于本节后面类似的情形。

此外，只要某个关于 Fukaya 范畴的代数构造具有函子性并保持拟等价，这个结论就同样适用于该构造，例如取导出范畴。\znote{135}{对 $\Ai$ 范畴“取导出范畴”，通常指向 $D(\mathcal A)=H^0(\operatorname{mod}\mathcal A)$，其中 $\operatorname{mod}\mathcal A$ 是 $\mathcal A$ 上右 $\Ai$ 模所成的 $\Ai$ 范畴，$H^0$ 只取零次上同调。欲使它成为三角范畴并使拟等价成为等价，还需再取分裂闭包（Karoubi 包络），即把幂等元的像形式地添进去。Yoneda 嵌入 $X\mapsto\hom_{\mathcal A}(-,X)$ 把 $\mathcal A$ 全忠实地放入 $\operatorname{mod}\mathcal A$，其像在分裂闭包下生成整个导出范畴；分裂闭包的存在性正是“幂等元必分裂”这条万有性质的体现。[来源：nLab "Karoubi envelope"，https://\allowbreak ncatlab.org/\allowbreak nlab/\allowbreak show/\allowbreak Karoubi+\allowbreak envelope 与 nLab "Yoneda embedding"，https://\allowbreak ncatlab.org/\allowbreak nlab/\allowbreak show/\allowbreak Yoneda+\allowbreak embedding（2026-11-15 访问）。]}剩下的问题是：构成协调系统的函子与自然变换本身，在多大程度上是典范的？这留给有兴趣的读者。
\begin{note}\label{st:10.1}
在普通范畴的框架中，有一种反向构造：从范畴 $A^i$ 的协调系统出发，构造一个共同范畴 $A^{\mathrm{tot}}$，以及等价 $A^i\to A^{\mathrm{tot}}$。在 $\Ai$ 背景下，这看来不能成立，原因在于我们的协调系统定义本质上太弱，只停留在上同调函子范畴 $H(\operatorname{fun})$ 的层次。
\end{note}

\subsection{自等价}\label{sec:10b}
这个主题与前一个紧密相关；事实上，可以用群胚统一处理。需要讨论群 $G$ 作用于范畴 $A$ 的含义。
\begin{definition*}
\begin{itemize}
\item \emph{严格 $G$ 作用}是一族函子 $F^g:A\to A$（$g\in G$），满足 $F^e=\Id_A$、$F^{g_2g_1}=F^{g_2}\circ F^{g_1}$。
\item \emph{协调 $G$ 作用}由函子 $F^g$ 及同构 $T^{g_2,g_1}:F^{g_2}\circ F^{g_1}\to F^{g_2g_1}$ 组成，满足 $F^e=\Id_A$；当 $g_2=e$ 或 $g_1=e$ 时，$T^{g_2,g_1}$ 为恒等变换；并且以下图交换：
\SourceEquationDiagram{diagram-10-4}{0.46730719}{10.4}
\SourcePage{144}
\item \emph{弱 $G$ 作用}是一族函子 $F^g$，满足 $F^e$ 同构于 $\Id_A$、$F^{g_2g_1}$ 同构于 $F^{g_2}\circ F^{g_1}$，但不指定这些同构的选择。
\end{itemize}

\end{definition*}
假设范畴 $A^{\mathrm{free}}$ 带有严格 $G$ 作用 $\{\Phi^g\}$，$A$ 是其全子范畴，且嵌入 $H:A\to A^{\mathrm{free}}$ 为等价。像\secref{10a}那样，选择 $K:A^{\mathrm{free}}\to A$ 及函子同构 $S:H\circ K\to\Id_{A^{\mathrm{free}}}$。于是，在 $A$ 上得到诱导的协调 $G$ 作用：
\begin{equation}\label{eq:10.5}
F^g=K\circ\Phi^g\circ H,\qquad
T^{g_2,g_1}=L_{K\circ\Phi^{g_2}}\bigl(R_{\Phi^{g_1}\circ H}(S)\bigr).
\end{equation}
现在应当很清楚，如何将这些定义推广到 $\Ai$ 框架。注意，在此还有比严格作用更基本的另一种可能，即第\sr{sec:6}{6 节}中的朴素群作用；不过它们在当前讨论中不起作用。为说明\eref{10.5}的 $\Ai$ 对应物满足协调关系，论证如下。对应于\eref{10.4}的图的两条复合路径，都具有形式
\[
\LG^1_{\mathcal K\circ\Phi^{g_3}}\bigl(\RG^1_{\Phi^{g_2}\circ\Phi^{g_1}\circ\mathcal H}(\text{某个表达式})\bigr),
\]
括号内的表达式分别为
\begin{equation}\label{eq:10.6}
\begin{aligned}
&\mu^2_{\Q}\bigl(S,\LG^1_{\mathcal H\circ\mathcal K}(\LG^1_{\Phi^{g_2}}(\RG^1_{\Phi^{g_2^{-1}}}(S)))\bigr),\\
&\mu^2_{\Q}\bigl(\LG^1_{\Phi^{g_2}}(\RG^1_{\Phi^{g_2^{-1}}}(S)),
\RG^1_{\Phi^{g_2}\circ\mathcal H\circ\mathcal K\circ\Phi^{g_2^{-1}}}(S)\bigr),
\end{aligned}
\end{equation}
其中 $\Q=\operatorname{fun}(\A^{\mathrm{free}},\A^{\mathrm{free}})$。应用\stref{引理}{1.4}，取
\[
\begin{gathered}
\G_1=\mathcal H\circ\mathcal K,\qquad
\G_2=\Phi^{g_2}\circ\mathcal H\circ\mathcal K\circ\Phi^{g_2^{-1}},\\
T_1=S,\qquad T_2=\LG^1_{\Phi^{g_2}}(\RG^1_{\Phi^{g_2^{-1}}}(S)),
\end{gathered}
\]
可知在 $H(\Q)$ 层次，\eref{10.6}中的两项相同。
\begin{note}\label{st:10.2}
与\stref{笔记}{10.1}所述相应，在普通范畴的框架中存在反向构造，将 $A$ 上的协调群作用变为等价范畴 $A^{\mathrm{strict}}$ 上的严格作用；特殊情形 $G=\Z/2$ 见\secref{14b}。同样，在 $\Ai$ 背景下，这看来不能成立，至少对上面给出的协调性定义如此。
\end{note}
固定一个带角正合辛流形 $M$，令 $G=\operatorname{Aut}(M,\partial M)$ 为其相对于边界的自同构群。每个 $\phi\in G$ 都在整个 Floer TQFT 上诱导典范自同构，特别是在 $H\F(M)^{\mathrm{pr}}$ 上诱导自同构；它们构成严格 $G$ 作用。在上链层次，按如下方式进行。定义一个 Fukaya 型的不分次 $\Ai$ 范畴 $\F(M)^{\mathrm{pr},\mathrm{free}}$，其对象是由正合 Lagrangian 子流形 $L\subset M$ 与 $\gamma\in G$ 组成的二元组 $(L,\gamma)$。为每组有限个对象选择 Floer 数据与扰动数据，借此定义态射空间 $CF^{\mathrm{pr}}((L_0,\gamma_0),(L_1,\gamma_1))$ 及它们之间的复合映射，并要求
\SourcePage{145}
这些选择关于 $G$ 作用 $\phi(L,\gamma)=(\phi(L),\phi\circ\gamma)$ 等变。这并无困难，因为该作用在对象上自由。结果是，$\F(M)^{\mathrm{pr},\mathrm{free}}$ 带有严格 $G$ 作用，由高阶项 $d>1$ 全为零的 $\Ai$ 函子 $\phi_*$ 组成。显然，可将 $\F(M)^{\mathrm{pr}}$ 作为全子范畴嵌入 $\F(M)^{\mathrm{pr},\mathrm{free}}$；因为任意两个 $(L,\gamma_0)$、$(L,\gamma_1)$ 在上同调层次同构，这个嵌入是拟等价。因此，得到 $\F(M)^{\mathrm{pr}}$ 上诱导的协调 $G$ 作用。\znote{137}{“该作用在对象上自由”是这里能选出等变数据的关键：若某个 $\phi\neq\id$ 固定了某个 $L$，则同一组 Floer 数据必须同时满足 $L$ 与 $\phi(L)=L$ 两套边界条件，横截性会因此失效，这正是一般等变横截性问题的困难所在。改用对象 $(L,\gamma)$ 后，$\phi$ 把 $(L,\gamma)$ 送到 $(\phi L,\phi\gamma)$，只要 $\phi\neq\id$ 两者就不同，于是可以在每个轨道上任选数据再用 $G$ 推移。代价是对象集变大：$\F(M)^{\mathrm{pr},\mathrm{free}}$ 的对象数约为 $\F(M)^{\mathrm{pr}}$ 的 $|G|$ 倍。[来源：原书 §(10b) 本段。]}

\subsection{Hamilton 同痕}\label{sec:10c}
下一个主题是\secref{8k}意义下延拓映射的高阶推广。这里只给出简要轮廓，略去技术细节。

在上面引入的 $\Ai$ 范畴 $\F(M)^{\mathrm{pr},\mathrm{free}}$ 中工作。设 $(\phi^s)$（$0\leq s\leq1$）是 $G=\operatorname{Aut}(M,\partial M)$ 中从 $\phi^0=\Id$ 到 $\phi^1=\phi$ 的同痕。按照\eref{8.23}的构造，可为每个 $(L,\gamma)$ 赋予一个上链
\[
T^0\in CF^{\mathrm{pr}}((L,\gamma),\phi(L,\gamma)).
\]
也就是说，取上半平面 $H$，配备由同痕 $\phi^s(L)$ 给出的移动 Lagrangian 边界条件，以及与给定 Floer 数据相容的扰动数据。注意，扰动数据的选择必然破坏 $H$ 的全纯自同构群。因此，将 $H$ 看作内部带有额外标点的曲面可能更自然，例如取 $z=i$，使曲面稳定。

沿着这个思路，更一般地考虑一个 $(d+1)$ 标点圆盘 $S$，它还带有额外的内部标点 $z$、一组条带状端，以及一组对象 $(L_k,\gamma_k)$（$k=0,\ldots,d$）。选择 $0=s_0\leq s_1\leq\cdots\leq s_{d+1}=1$，并为 $S$ 配备移动边界条件，使其在每个分支 $C_k$（$0\leq k\leq d$）上从 $\phi^{s_k}(L_k)$ 变化到 $\phi^{s_{k+1}}(L_k)$。然后在 $S$ 上选择扰动数据，使入向条带状端上的行为由对象对 $((L_0,\gamma_0),\phi(L_d,\gamma_d))$ 的 Floer 数据确定，出向端上的行为则分别由 $(\phi^{s_k}(L_{k-1},\gamma_{k-1}),\phi^{s_k}(L_k,\gamma_k))$ 的 Floer 数据确定。注意，由 $\F(M)^{\mathrm{pr},\mathrm{free}}$ 的定义，后者是将 $((L_{k-1},\gamma_{k-1}),(L_k,\gamma_k))$ 的 Floer 数据经 $\phi^{s_k}$ 诱导而得。

希望所有这些都光滑依赖于复结构及标点选择。因此，不应只考虑单个 $(S,z)$，而应考虑所有带额外一个内部标点的 $(d+1)$ 标点圆盘所成的模空间 $\mathcal R^{d+1,1}$，以及它上面的万有族。这个空间有 Deligne–Mumford 型紧化
\begin{equation}\label{eq:10.7}
\overline{\mathcal R}^{d+1,1}=\coprod_{(T,v_*)}\left(\mathcal R^{|v_*|,1}\times\prod_{v\ne v_*}\mathcal R^{|v|}\right),
\end{equation}
其各层以带有一个指定顶点 $v_*$ 的 $d$ 叶树为指标；$v_*$ 可以具有任意正度数，其他顶点则至少为三价。
\SourcePage{146}
条带状端、各 $\{s_k\}$ 及扰动数据的选择，必须满足与\eref{10.7}相关的适当一致性条件，这里不明确表述。随后，考虑对应的二元组 $(r\in\mathcal R^{d+1,1},u:\mathcal S_r^{d+1,1}\to M)$ 的模空间，其中 $u$ 是非齐次伪全纯映射。对该空间中满足适当渐近行为的点计数，定义映射
\[
\begin{aligned}
T^d:\ &CF^{\mathrm{pr}}((L_{d-1},\gamma_{d-1}),(L_d,\gamma_d))\otimes\cdots\\
&\cdots\otimes CF^{\mathrm{pr}}((L_0,\gamma_0),(L_1,\gamma_1))
\longrightarrow CF^{\mathrm{pr}}((L_0,\gamma_0),\phi(L_d,\gamma_d)).
\end{aligned}
\]
这里再次使用了以下事实：$CF^{\mathrm{pr}}(\phi^{s_k}(L_{k-1},\gamma_{k-1}),\phi^{s_k}(L_k,\gamma_k))$ 与 $CF^{\mathrm{pr}}((L_{k-1},\gamma_{k-1}),(L_k,\gamma_k))$ 可以典范地等同。各 $\{T^d\}$ 满足方程
\begin{equation}\label{eq:10.8}
\begin{aligned}
&\sum_{m,n}\mu^{d-m+1}\bigl(\phi_*^1(a_d),\ldots,\phi_*^1(a_{n+m+1}),
T^m(a_{n+m},\ldots,a_{n+1}),a_n,\ldots,a_1\bigr)\\
&\quad+\sum_{m,n}T^{d-m+1}\bigl(a_d,\ldots,a_{n+m+1},
\mu^m(a_{n+m},\ldots,a_{n+1}),a_n,\ldots,a_1\bigr)=0.
\end{aligned}
\end{equation}
\sr{fig:10.1}{图~10.1} 给出了两类求和项所对应的退化，并着重显示移动及固定的边界条件。由于 $\Ai$ 函子 $\phi_*$ 的高阶项为零，而且目前不需要考虑符号，\eref{10.8}恰好是 $T$ 成为从恒等函子到 $\phi_*$ 的自然变换的条件。像\secref{8k}末尾已经指出的，每个 $T^0$ 都是 $H\F(M)^{\mathrm{pr},\mathrm{free}}$ 中的同构，因此由\stref{引理}{1.6}，$T$ 在 $\Ai$ 函子的上同调范畴中为同构。将这个结果转移到与之拟等价的 $\Ai$ 范畴 $\F(M)^{\mathrm{pr}}$，得到：
\begin{proposition}\label{st:10.3}
设 $\phi\in\operatorname{Aut}(M,\partial M)$ 是与恒等映射同痕的自同构，则由 $\phi$ 诱导的 $\Ai$ 函子在 $H(\operatorname{fun}(\F(M)^{\mathrm{pr}},\allowbreak\F(M)^{\mathrm{pr}}))$ 中同构于恒等函子。
\end{proposition}
换言之，$\operatorname{Aut}(M,\partial M)$ 在 $\F(M)^{\mathrm{pr}}$ 上的协调作用，下降为 $\pi_0(\operatorname{Aut}(M,\partial M))$ 的弱作用。
\begin{note}\label{st:10.4}
考虑群 $\Gamma$ 在范畴 $A$ 上弱作用的一般情形。设 $\HH^0(A,A)^\times$ 是从 $\Id_A$ 到自身的函子同构所成的集合，以复合作为运算，它成为 Abel 群。将这一作用赋予协调结构的障碍，是一个群上同调类，更精确地说，是 $H^2(\Gamma;\HH^0(A,A)^\times)$ 中的元素。现在假设 $\Gamma=\pi_0(G)$、$G=\operatorname{Aut}(M,\partial M)$；为简便起见，取 $A=H\F(M)^{\mathrm{pr}}$ 为 Donaldson–Fukaya 范畴。此时，障碍类由 $H^2(\pi_0(G);\pi_1(G))$ 中一个典范类经同态 $\pi_1(G)\to\HH^0(A,A)^\times$ 诱导而来；这个典范类对任意拓扑群都存在，而同态利用\bcite{123}、\bcite{1}的方法定义。在 $\Ai$ 范畴层次，大概也存在相应陈述，
\SourcePage{147}
\SourceFigure{figure-10-1}{0.55869171}{10.1}
但我们尚未探究细节。无论如何，即使如此也仍然只是整个问题的一部分。与\stref{笔记}{10.1}一样，我们对 $\Ai$ 范畴上协调群作用的定义应当视为初步近似；可以期待，最终的概念将涉及 $G$ 的整个同伦型，而不仅仅是次数不超过 $1$ 的同伦群。
\end{note}

\subsection{接触型边界的情形}\label{sec:10d}
接续\secref{7b}，讨论具有接触型边界的正合辛流形。给定这样的流形 $(M,\varpi,\theta)$，可以将 $\mathcal J$ 的定义改为所有相容、且在边界附近为接触型的殆复结构所成的空间。这意味着，在所得 Floer 数据与扰动数据概念中，殆复结构族 $J(t)$ 或 $J(z)$ 在 $\partial M$ 附近不再保持不变。由\stref{引理}{7.3}，以及非齐次项在边界附近为零这一事实，方程\eref{8.8}、\eref{8.9}的解仍然留在 $M$ 的内部。以这种方式定义 $\F(M)^{\mathrm{pr}}$，会发现它本质上不依赖于任何殆复结构的选择；当然，这里的含义是不同选择组成协调系统。\znote{139}{这里用到接触型边界的凸性：$\partial M$ 附近可取 Liouville 结构，使 $\varpi=d(t\theta)$（$t\in(-\epsilon,0]$），于是 $d\theta$ 限制在 $J$-复直线上为正；对解作关于 $t$ 的最大值原理型论证，就排除了它触及边界。因此才可以在 $\partial M$ 附近自由改动 $J$ 而不改变模空间。\stref{引理}{7.3} 把这一论证表述为“解落在 $\operatorname{int}(M)$ 中”。[来源：原书 §(7b) 与\stref{引理}{7.3}，并见\eref{8.8}、\eref{8.9} 处的讨论。]}
\SourcePage{148}
\begin{proposition}\label{st:10.5}
设 $M$ 是具有接触型边界的正合辛流形，$N$ 是带角正合辛流形，$\iota:M\hookrightarrow\operatorname{int}(N)$ 是嵌入，满足 $\iota^*\varpi_N=\varpi_M$、$\iota^*\theta_N=\theta_M+d(\text{某个函数})$。则 $\F(M)^{\mathrm{pr}}$ 拟同构于 $\F(N)^{\mathrm{pr}}$ 的全 $\Ai$ 子范畴，后者由位于 $\iota(\operatorname{int}(M))$ 内的正合 Lagrangian 子流形组成。
\end{proposition}
\begin{proof}
可以给 $\theta_N$ 加上一个在 $\partial N$ 附近为零的函数的微分，而完全不影响 $\F(N)^{\mathrm{pr}}$。这样就将情形归结为 $\iota^*\theta_N=\theta_M$。此后，将 $\iota$ 简单视为包含映射。在 $N$ 上固定殆复结构 $J$，使其在 $\partial N$ 附近与 $I_N$ 相同，且 $J|_M$ 在 $\partial M$ 附近为接触型。

定义所讨论的 $\F(N)^{\mathrm{pr}}$ 的子范畴时，可以选择所有殆复结构在 $N\setminus\operatorname{int}(M)$ 的一个邻域中等于 $J$，并选择 Hamilton 函数与扰动数据，使它们在这个集合的一个邻域中为零。于是可应用\stref{引理}{7.5}，说明所有非齐次伪全纯多边形都留在 $\operatorname{int}(M)$ 中。在此情形，两个范畴实际上相等。对于一般的数据选择，拟同构结论由前面关于良定义性的讨论得到。
\end{proof}
\begin{corollary}\label{st:10.6}
设 $M,N$ 是具有接触型边界的正合辛流形，$\widehat M,\widehat N$ 是像\eref{7.1}那样接上锥得到的非紧致流形。假设存在正合辛同构 $\widehat\phi:\widehat M\to\widehat N$，则范畴 $\F(M)^{\mathrm{pr}}$、$\F(N)^{\mathrm{pr}}$ 拟等价。\znote{140}{要点是：$\F(M)^{\mathrm{pr}}$ 实际只依赖接触型边界流形的“完成”$\widehat M=M\cup_{\partial M}([0,\infty)\times\partial M)$，而 $\widehat\phi$ 是完成之间的正合辛同构。证明方式是把\stref{命题}{10.5}用两次：$M$ 嵌入 $\widehat M$、$N$ 嵌入 $\widehat N$，而 $\widehat\phi$ 把 $\iota_M(\operatorname{int}M)$ 送到 $\iota_N(\operatorname{int}N)$ 的紧致子集内，于是 $\F(M)^{\mathrm{pr}}$ 与 $\F(N)^{\mathrm{pr}}$ 都与 $\widehat M$ 上的同一个全子范畴拟同构。[来源：原书 §(10d)，完成与\eref{7.1} 的定义见 §(7b)。]}
\end{corollary}
\begin{proof}
考虑完备化的有限部分
\[
\widehat M^\sigma=M\cup_{\partial M}([0,\sigma]\times\partial M),\qquad
\widehat N^\tau=N\cup_{\partial N}([0,\tau]\times\partial N),
\]
选取 $0\ll\sigma\ll\tau$，使 $N\subset\widehat\phi(\widehat M^\sigma)\subset\widehat N^\tau$。首先，$M$ 的 Fukaya 范畴与 $\widehat M^\sigma$ 的拟等价；类似地，$N$ 的与 $\widehat N^\tau$ 的拟等价。这很容易看出：利用负 Liouville 流，两个流形可在辛形式及其原始一形式作共形缩放后彼此识别。若相应地缩放全部 Floer 数据及扰动数据，两者的 Fukaya 范畴实际上相同。

现在应用\stref{命题}{10.5}，它说明 $\F(\widehat M^\sigma)^{\mathrm{pr}}$ 拟同构于 $\F(\widehat N^\tau)^{\mathrm{pr}}$ 的某个全子范畴。该子范畴显然包含所有位于 $N$ 内的正合 Lagrangian 子流形，而在同痕意义下，这已经包含一切，因此得到所需拟等价。
\end{proof}
由\stref{引理}{7.2}，得到：
\begin{corollary}\label{st:10.7}
在具有接触型边界的正合辛流形这一类中，$\F(M)^{\mathrm{pr}}$ 在 $M$ 的变形下不变。
\end{corollary}
\SourcePage{149}
\begin{note}\label{st:10.8}
稍作变化，也可以使用完备化上自然的一类殆复结构 $\widehat J$，直接定义范畴 $\F(\widehat M)^{\mathrm{pr}}$。由\stref{引理}{7.4}可知，$\F(M)^{\mathrm{pr}}$ 与 $\F(\widehat M)^{\mathrm{pr}}$ 拟等价。这种做法的一个优点，是得到在无穷远处为接触型的正合辛自同构所成的群 $\operatorname{Aut}(\widehat M)$ 在 $\F(M)^{\mathrm{pr}}$ 上的协调作用。
\end{note}
\begin{corollary}\label{st:10.9}
设 $N$ 是带角正合辛流形。像\stref{引理}{7.6}那样，在充分小的 $\delta$ 处将角磨圆，得到具有接触型边界的正合辛流形 $M\subset\operatorname{int}(N)$。则 $\F(M)^{\mathrm{pr}}$ 与 $\F(N)^{\mathrm{pr}}$ 拟等价。
\end{corollary}
\begin{proof}
这立即由\stref{命题}{10.5}及每个正合 Lagrangian 子流形都能移入 $M$ 内部这一事实推出。
\end{proof}
\begin{corollary}\label{st:10.10}
在拟等价意义下，带角正合辛流形的 Fukaya 范畴在变形下不变。
\end{corollary}
\begin{proof}
取 $N$ 及一个变形 $(\varpi_N^s,\theta_N^s,I_N^s)$。考察\stref{引理}{7.6}的证明，可见存在 $M\subset\operatorname{int}(N)$，使得对所有 $s$，$(\varpi_N^s,\theta_N^s)$ 在 $M$ 上的限制都给出具有接触型边界的正合辛流形结构。\stref{推论}{10.7}说明 $\F(M)^{\mathrm{pr}}$ 不依赖于 $s$，而\stref{推论}{10.9}使这一结果可以转移到 $\F(N)^{\mathrm{pr}}$。
\end{proof}
\begin{corollary}\label{st:10.11}
设 $\phi:N\to\widetilde N$ 是带角正合辛流形的同构。则存在从 $\F(N)^{\mathrm{pr}}$ 到 $\F(\widetilde N)^{\mathrm{pr}}$ 的拟等价，将每个正合 Lagrangian 子流形 $L$ 映到 $\phi(L)$。
\end{corollary}
在给出证明之前，应先澄清一个可能产生混淆的地方。若同构与殆复结构相容，即边界附近有 $\phi_*(I_N)=I_{\widetilde N}$，则\secref{10a}中的形式论证就足以证明结论。事实上，这些论证给出稍强的陈述，因为所得拟等价在相关函子范畴中自动在同构意义下典范。一般情形的困难，恰恰来自\secref{7a}所定义的同构概念不考虑殆复结构。
\begin{proof}
将角磨圆，得到具有接触型边界的正合辛流形 $M\subset\operatorname{int}(N)$。限制 $\phi|_M$ 显然是正合辛嵌入。取\stref{推论}{10.9}中拟等价的拟逆，再应用\stref{命题}{10.5}，得到 $\Ai$ 函子
\begin{equation}\label{eq:10.9}
\F(N)^{\mathrm{pr}}\longrightarrow\F(M)^{\mathrm{pr}}\longrightarrow\F(\widetilde N)^{\mathrm{pr}},
\end{equation}
\SourcePage{150}
它是上同调全忠实的。将 $\delta$ 取得充分小，可使 $\widetilde N$ 的任意正合 Lagrangian 子流形都能通过负 Liouville 流移入 $\phi(\operatorname{int}(M))$，于是\eref{10.9}就是拟等价。按定义，它把 $L$ 映到 $\phi(\widetilde L)$，其中 $\widetilde L$ 与 $L$ 同痕。但因为 $\phi(L)$、$\phi(\widetilde L)$ 是 $H\F(\widetilde N)^{\mathrm{pr}}$ 中同构的对象，纯代数论证便可修改此函子，使 $L$ 确实映到 $\phi(L)$。
\end{proof}

\subsection{同伦方法}\label{sec:10e}
Floer 理论中还有另一种证明良定义性定理的方法。这种方法相当直接，历史上也早于延拓映射，但在具体实施时，技术上可能颇为棘手。虽然它并非我们当前目的所必需，我们仍介绍这一方法，因为它本身就很重要。

设对于某个带角正合辛流形 $M$，给定 Fukaya 范畴的两个版本 $\F(M)^{\mathrm{pr},s}$，$s=0,1$。假定两者采用相同的条带状端万有选择，因而区别来自所使用的 Floer 数据及扰动数据。进一步假设，对每一对 $(L_0,L_1)$，都存在 Floer 数据的光滑同伦
\begin{equation}\label{eq:10.10}
(H^s_{L_0,L_1},J^s_{L_0,L_1}),\qquad 0\leq s\leq1,
\end{equation}
连接定义 $\F(M)^{\mathrm{pr},0}$ 和 $\F(M)^{\mathrm{pr},1}$ 时使用的数据。注意，第一个假设只是技术性的，第二个假设却相当强，因为它排除了生灭现象。\znote{142}{生灭现象指参数 $s$ 经过某个临界值时，$\phi^{s,1}(L_0)$ 与 $L_1$ 由横截变为不横截（或反之），交点因此成对地产生或湮灭。若允许这种情形，交点集便不能随 $s$ 连续地等同，\eref{10.11}那样的同构也不存在，构造必须改用带生灭的参数族；而\eref{10.10}要求对每个 $s$ 都是 Floer 数据（交点横截），所以这里是较简单的情形。[来源：原书 §(10e) 本段；另一种处理见 §(8k) 的延拓映射。]}为解释这一点，固定 $(L_0,L_1)$，以下略去这些下标。设 $X^s$ 是对应于 $H^s$ 的 Hamilton 向量场，它还依赖另一个时间变量 $t\in[0,1]$。以 $\phi^{s,t}$ 表示沿后一个变量积分得到的 Hamilton 微分同胚族，即 $\phi^{s,0}=\Id_M$，且 $\partial\phi^{s,t}/\partial t=X^{s,t}(\phi^{s,t})$。由于所有\eref{10.10}都必须是 Floer 数据，交集 $\phi^{s,1}(L_0)\cap L_1$ 都是横截的，故交点随 $s$ 光滑变化。换言之，设 $CF^{\mathrm{pr}}(L_0,L_1)^s$ 是对某个 $s\in[0,1]$ 用\eref{10.10}定义的 Floer 上链群，$\mathcal C(L_0,L_1)^s$ 是其标准生成元集合，由满足 $y^s(0)\in L_0$、$y^s(1)\in L_1$ 及 $dy^s/dt=X^{s,t}(y^s(t))$ 的映射 $y^s:[0,1]\to M$ 组成。同样，以 $\mathcal C(L_0,L_1)^{[0,1]}$ 表示光滑映射 $y:[0,1]^2\to M$ 的集合，要求每个限制 $y^s=y(s,\cdot)$ 都属于 $\mathcal C(L_0,L_1)^s$；以 $CF^{\mathrm{pr}}(L_0,L_1)^{[0,1]}$ 表示该集合生成的 $\K$ 向量空间。于是，对任意 $s$，限制映射 $\mathcal C(L_0,L_1)^{[0,1]}\to \mathcal C(L_0,L_1)^s$ 都是双射，从而诱导同构
\begin{equation}\label{eq:10.11}
CF^{\mathrm{pr}}(L_0,L_1)^{[0,1]}\cong CF^{\mathrm{pr}}(L_0,L_1)^s.
\end{equation}
若以为仅凭这些同构便能得到 $\Ai$ 函子，就过于天真了。实际上，必须考虑加入 $s$ 参数的非齐次伪全纯映射的适当模空间。
\SourcePage{151}
先从一阶构造，也就是上链映射的构造开始。这个构造相当清楚，可以作为一般情形的模型。给定 $y_0,y_1\in \mathcal C(L_0,L_1)^{[0,1]}$，记 $\M_Z^*(y_0,y_1)^s=\M_Z^*(y_0^s,y_1^s)$，表示关于数据 $(H^s,J^s)$ 的 Floer 方程的非驻定解模空间，其极限为 $y_0^s,y_1^s$，并照例除去平移作用。同样，记 $\M_Z^*(y_0,y_1)^{[0,1]}$ 为二元组 $(s,u)$ 的空间，其中 $s\in[0,1]$，$u\in\M_Z^*(y_0,y_1)^s$。对于每个 $\nu\geq1$，现在引入相关的模空间 $\mathcal H^\nu(y_0,y_1)^{[0,1]}$。它的点是 $\nu$ 元组 $((s_1,u_1),\ldots,(s_\nu,u_\nu))$，满足 $s_1\leq\cdots\leq s_\nu$，且对某些 $y_0=\widetilde y_0,\widetilde y_1,\ldots,\widetilde y_\nu=y_1$，每个 $u_i\in\M_Z^*(\widetilde y_{i-1},\widetilde y_i)^{s_i}$。可以把 $\mathcal H^\nu(y_0,y_1)$ 的元素抽象地画成由双标点圆盘组成的链，每个圆盘带有一个变量 $s_k$，见\sr{fig:10.2}{图~10.2}。为书写方便，形式上规定：当 $y_0=y_1$ 时，$\mathcal H^0(y_0,y_1)^{[0,1]}$ 是单点；否则为空集。
\SourceFigure{figure-10-2}{0.58569168}{10.2}
可以将这一构造看作平凡族 $[0,1]\times Z\to[0,1]$ 的带参数模问题，只是与\secref{9h}的设置不同，端所对应的 Floer 数据也随参数 $s$ 变化。不过，基本结构仍然相同；特别地，每个 $(s,u)\in\M_Z^*(y_0,y_1)^{[0,1]}$ 都有对应的扩展线性化算子，它像\eref{9.18}那样将额外变量纳入考虑。正则性定义为这个算子的满射性，并蕴含模空间的局部光滑性。如果同伦\eref{10.10}选得足够一般，正则性便处处成立，以下假定如此。为了说明空间 $\mathcal H^\nu(y_0,y_1)^{[0,1]}$ 具有良好性质，还必须证明它们的边界面本身也正则。这些边界面是乘积
\[
\prod_{i=1}^\nu\M_Z^*(\widetilde y_{i-1},\widetilde y_i)^{[0,1]}
\]
中由一个或多个形如 $s_i=s_{i+1}$ 的方程截出的子流形。使用稍复杂一些的标准横截性论证，可证明只要各 $(s_i,u_i)$ 两两不同，就能实现这一点。幸运的是，在目前情形中这个条件自动满足，因为每个 $u_i$ 都使 $H^{s_i}$ 扰动后的作用量泛函严格下降。

能量论证表明，当 $\nu$ 充分大时，$\mathcal H^\nu(y_0,y_1)^{[0,1]}$ 为空。此外，对固定的 $\nu$，带额外 $s$ 参数的 Gromov 紧性说明，$\mathcal H^\nu(y_0,y_1)^{[0,1]}$ 中只有有限多个孤立点。因此可以
\SourcePage{152}
定义映射 $\Gamma:CF^{\mathrm{pr}}(L_0,L_1)^1\to CF^{\mathrm{pr}}(L_0,L_1)^0$，令
\begin{equation}\label{eq:10.12}
\Gamma(y_1^1)=\sum_{y_0}\sum_{\nu\geq0}\#\mathcal H^\nu(y_0,y_1)^{[0,1]}\,y_0^0.
\end{equation}
\znote{144}{这与\eref{9.23}的计数方式相同，只是多了一个参数 $s$：$\nu$ 元组 $(s_1,u_1),\ldots,(s_\nu,u_\nu)$ 相当于把一条链上各段 Floer 条带按 $s$ 排好，$\#$ 仍是零维模空间的模 $2$ 点数。$\nu$ 充分大时空间为空（能量论证），故外层对 $\nu$ 的求和有限；对每个 $\nu$，带 $s$ 参数的 Gromov 紧性又保证孤立点只有有限多个。这两点合起来使\eref{10.12}是良定义的 $\K$-线性映射。[来源：原书 §(10e) 与 §(9j) 的\eref{9.23}。]}
\begin{lemma}\label{st:10.12}
$\Gamma$ 是（不分次）链复形的同构。
\end{lemma}
\begin{proof}
证明思路很直接。记 $\mathcal H^\nu(y_0,y_1)^{[0,1],1}$ 为 $\mathcal H^\nu(y_0,y_1)^{[0,1]}$ 的所有一维连通分支的并。这是带边界的一维流形，边界点出现于以下情形：(i) $s_1=0$；(ii) $s_\nu=1$；或 (iii) 对某个 $i$，$s_{i-1}=s_i$。取 Gromov 紧化 $\overline{\mathcal H}^\nu(y_0,y_1)^{[0,1],1}$ 后，又加入一些边界点，对应于某个 $u_i$ 断裂为一对 Floer 轨线，即对某个 $\widetilde y$，得到 $\M_Z^*(\widetilde y_{i-1},\widetilde y)^{s_i}\times\M_Z^*(\widetilde y,\widetilde y_i)^{s_i}$ 的孤立点。在边界点的代数计数中，这些点恰好抵消上述情形 (iii) 的贡献，剩下等式
\[
\begin{aligned}
&\sum_{\nu\geq0}\sum_{\widetilde y}\#\M_Z^*(y_0,\widetilde y)^0\,\#\mathcal H^\nu(\widetilde y,y_1)^{[0,1]}\\
&\qquad+\sum_{\nu\geq0}\sum_{\widetilde y}\#\mathcal H^\nu(y_0,\widetilde y)^{[0,1]}\,\#\M_Z^*(\widetilde y,y_1)^1=0,
\end{aligned}
\]
这恰好说明 $\Gamma$ 是链映射。要证明它是同构，论证大致如下。如果模空间 $\M_Z^*(y_0,y_1)^{[0,1]}$ 没有孤立点，则由定义 $\Gamma=\Id$。再看次简单的情形，即孤立点恰好只出现在参数 $s$ 的一个取值处。这时 $\Gamma=\Id+(\text{下三角部分})$，其中第二项关于 $A_{H^s}$ 作用量的值所给出的滤过严格下降。最后，利用\eref{10.12}与同伦复合的良好相容性，将一般情形归结为已经讨论的情形。\znote{144}{“$\Gamma=\Id+$下三角”之所以给出同构，是因为按 $A_{H^s}$ 的作用量值把生成元排成有限降链后，$\Gamma$ 在该滤过下的矩阵对角元为 $1$，非对角元只出现在作用量严格更小处；这样的矩阵可逆，其逆按作用量差逐级（Neumann 级数）给出，每一级只有有限多项。这正是证明要从“无孤立点故 $\Gamma=\Id$”逐步放宽到一般情形的理由。[来源：原书 §(10e)；作用量滤过与\eref{8.12} 的能量恒等式相容。]}沿此思路的完整论证，见下文\stref{笔记}{10.14}所列参考文献。
\end{proof}
对于伪全纯 $(d+1)$ 边形（$d\geq2$），相应构造在技术上稍复杂一些。对每个正合 Lagrangian 子流形的 $(d+1)$ 元组 $(L_0,\ldots,L_d)$，选择
\begin{equation}\label{eq:10.13}
\sigma_{L_0,\ldots,L_d,k},\quad1\leq k\leq d,\qquad
(K^s_{L_0,\ldots,L_d},J^s_{L_0,\ldots,L_d}),\quad0\leq s\leq1,
\end{equation}
满足下述性质；讨论这些性质时，略去下标 $L_0,\ldots,L_d$。每个 $\sigma_k$ 是光滑函数 $\mathcal R^{d+1}\times[0,1]\to[0,1]$，当 $s=0,1$ 时满足 $\sigma_k(r,s)=s$，称为\emph{延迟函数}。$(K^s,J^s)$ 是万有族 $\mathcal S^{d+1}$ 上扰动数据的同伦，在端点 $s=0,1$ 应还原为 $\F(M)^{\mathrm{pr},s}$ 使用的数据，并且在由延迟函数确定的 $s$ 变量变换之后，与 Floer 数据\eref{10.10}相容。具体地，在每个条带状端 $\eps_k^{d+1}:\mathcal R^{d+1}\times Z^\pm\to\mathcal S^{d+1}$ 的像上，要求
\begin{equation}\label{eq:10.14}
(\eps_k^{d+1})^*(K^s,J^s)=
\begin{cases}
(H^s_{L_0,L_d}(t)\,dt,J^s_{L_0,L_d}(t)),&k=0,\\
(H^{\widetilde s}_{L_{k-1},L_k}(t)\,dt,J^{\widetilde s}_{L_{k-1},L_k}(t)),
&k>0,\quad\widetilde s=\sigma_k(r,s).
\end{cases}
\end{equation}
\SourcePage{153}
这里的记号有些冲突：$t$ 是半无穷条带 $Z^\pm$ 的第二个坐标，然而 $s$ 并不是第一个坐标（该坐标在下面的论证中不起作用），而是控制同伦的变量。入向端与出向端的 Floer 数据\eref{10.14}一般对应于不同的取值 $\widetilde s\ne s$，这一事实影响曲面的粘合方式。现在说明怎样将\secref{9i}最初引入的扰动数据一致性选择概念调整到这里。

设 $T$ 为稳定的 $d$ 叶树。对每个顶点 $v\in\operatorname{Ve}(T)$，固定一个 $|v|$ 标点圆盘 $S_v$，以 $r_v\in\mathcal R^{|v|}$ 表示对应的模参数。为这些圆盘配备彼此相容的 Lagrangian 标记。再固定 $s\in[0,1]$。随后可用下面的递归规则为每个顶点赋予参数值 $s_v\in[0,1]$：首先，若 $v$ 是离根最近的顶点，令 $s_v=s$。其次，设 $e$ 为连接顶点 $v^-(e)$ 与 $v^+(e)$ 的任意内边；回忆\secref{9d}，边的定向指向远离根的方向。假设对应的旗标 $f^+(e)=(v^-(e),e)$，按典范次序是邻接 $v^-(e)$ 的第 $k_e$ 个正旗标，则令
\begin{equation}\label{eq:10.15}
s_{v^+(e)}=\sigma_{k_e}(r_{v^-(e)},s_{v^-(e)}),
\end{equation}
其中 $\sigma_{k_e}$ 是以 $S_{v^-(e)}$ 的 $|v^-(e)|$ 个 Lagrangian 标记为指标的延迟函数之一。现在以粘合参数 $\{\rho_e\}$ 将各 $S_v$ 粘合起来，以 $r\in\mathcal R^{d+1}$ 表示所得 $(d+1)$ 标点圆盘 $S$ 的分类模点。回忆，$S$ 的出向无穷远点可与 $T$ 的叶子识别。若叶子对应的旗标为 $f_{k_1}(v_1),\ldots,f_{k_d}(v_d)$，就在这些无穷远点上分别赋予数值 $s^1=\sigma_{k_1}(r_{v_1},s_{v_1}),\ldots,s^d=\sigma_{k_d}(r_{v_d},s_{v_d})\in[0,1]$。一致性条件的第一部分只涉及延迟函数：当粘合参数充分小时，这些数值应分别等于 $S$ 自然对应的数值 $\sigma_1(r,s),\ldots,\sigma_d(r,s)$。等式两边都略去了 Lagrangian 标记，但选择都是自然的：由于已为所有顶点选定相容标记，每个 $S_{v_k}$ 都带有 Lagrangian 标记，而 $S$ 通过粘合继承这些标记。接着，在每个 $S_v$ 上配备时间 $s_v$ 处的扰动数据\eref{10.13}。条件\eref{10.14}与\eref{10.15}共同保证这些数据在对应条带状端上具有相同的行为，因而可以粘合，得到 $S$ 上的扰动数据。另一方面，也可以通过分类同构 $S\cong\mathcal S_r^{d+1}$ 拉回\eref{10.13}，为同一个圆盘配备扰动数据。一致性条件要求这两种数据满足与\secref{9i}相同的关系：当粘合参数充分小时，在 $S$ 的薄部完全相等，在厚部渐近一致；后一个条件用万有族的部分紧化来定义。

选定满足这种一致性条件的\eref{10.13}后，就能引入对应的模空间。给定 $(L_0,\ldots,L_d)$ 及 $y_0\in \mathcal C(L_0,L_d)^{[0,1]}$、$y_k\in \mathcal C(L_{k-1},L_k)^{[0,1]}$，令 $\M^{d+1}(y_0,\ldots,y_d)^{[0,1]}$ 为
\SourcePage{154}
三元组 $(r,s,u)$ 的空间，其中 $r\in\mathcal R^{d+1}$、$s\in[0,1]$，而 $u:\mathcal S_r^{d+1}\to M$ 满足采用扰动数据\eref{10.13}的非齐次伪全纯映射方程；其在入向条带状端上的极限为 $y_0^s$，在各出向端上的极限为 $y_k^{\widetilde s}$，其中 $\widetilde s=\sigma_{L_0,\ldots,L_d,k}(r,s)$。与前面考虑的 $\mathcal H^\nu(y_0,y_1)^{[0,1]}$ 类似，还存在另一族密切相关的模空间 $\mathcal H^T(y_0,\ldots,y_d)^{[0,1]}$。其定义部分地类似于\secref{9l}讨论的紧化，我们借用那里的某些记号。固定半稳定 $d$ 叶树 $T$。$\mathcal H^T(y_0,\ldots,y_d)^{[0,1]}$ 的一个点由下列数据组成：
\begin{itemize}
\item 对每个满足 $|v|\geq3$ 的顶点，给定 $r_v\in\mathcal R^{|v|}$，从而给定一个 $|v|$ 标点圆盘 $S_v$。对其余顶点，令 $S_v=Z$。
\end{itemize}
可以用唯一的方式给各 $S_v$ 配备相容的 Lagrangian 标记，使粘合后的 $(d+1)$ 标点圆盘上的诱导标记等于给定的 $L_k$。完成这一选择后，下一项数据是：
\begin{itemize}
\item 对每个 $v$，给定数值 $s_v\in[0,1]$，满足\eref{10.15}的不等式版本。也就是说，对于每条内边 $e$，
\begin{equation}\label{eq:10.16}
s_{v^+(e)}\geq
\begin{cases}
\sigma_{k_e}(r_{v^-(e)},s_{v^-(e)}),&|v^-(e)|\geq3,\\
s_{v^-(e)},&|v^-(e)|=2.
\end{cases}
\end{equation}
\end{itemize}
接着，回忆在这里每个旗标 $f$ 都对应一对 Lagrangian 子流形 $(L_{f,0},L_{f,1})$。
\begin{itemize}
\item 对每个 $f$，指定 $\widetilde y_f\in \mathcal C(L_{f,0},L_{f,1})^{[0,1]}$，满足下列条件。若 $f^-$ 是对应于 $T$ 的根的唯一旗标，则 $\widetilde y_{f^-}=y_0$。同样，对应于叶子的 $d$ 个旗标上的 $\widetilde y_f$ 分别为 $y_1,\ldots,y_d$。若 $f^\pm(e)$ 是一条内边的两个旗标，则 $\widetilde y_{f^-(e)}=\widetilde y_{f^+(e)}$。
\item 最后，对每个 $v$，给定映射 $u_v:S_v\to M$；当 $|v|=2$ 时，它是\eref{8.8}的非驻定解，当 $|v|\geq3$ 时，它是\eref{8.9}的解。相应的 Floer 数据和扰动数据分别由\eref{10.10}和\eref{10.13}在 $s=s_v$ 处给出；条带状端上的极限应为 $v$ 周围各旗标 $f$ 对应的 $\widetilde y_f$，在上述相应参数处取值。等价地，
\[
\begin{cases}
(s_v,u_v)\in\M_Z^*(\widetilde y_{f_0(v)},\widetilde y_{f_1(v)})^{[0,1]},&|v|=2,\\
(r_v,s_v,u_v)\in\M^{|v|}(\widetilde y_{f_0(v)},\ldots,\widetilde y_{f_{|v|-1}(v)})^{[0,1]},&|v|\geq3.
\end{cases}
\]
\end{itemize}
\sr{fig:10.3}{图~10.3}示意各曲面 $S_v$ 及其 Lagrangian 标记，并画出各 $s_v$ 之间相应的不等式。照例，若两组映射 $\{u_v:S_v\to M\}$ 只相差各不稳定分支 $S_v\cong Z$ 上的平移，则视为相同。
\SourcePage{155}
\SourceFigure{figure-10-3}{0.64592238}{10.3}
假定\eref{10.13}中的选择足够一般，就可以仿照\eref{10.12}定义多线性映射
\begin{equation}\label{eq:10.17}
\begin{aligned}
&\G^d:CF^{\mathrm{pr}}(L_{d-1},L_d)^1\otimes\cdots\otimes CF^{\mathrm{pr}}(L_0,L_1)^1
\longrightarrow CF^{\mathrm{pr}}(L_0,L_d)^0,\\
&\G^d(y_d^1,\ldots,y_1^1)=\sum_{y_0}\sum_T\#\mathcal H^T(y_0,\ldots,y_d)^{[0,1]}\,y_0^0.
\end{aligned}
\end{equation}
\begin{proposition}\label{st:10.13}
\eref{10.12}中的映射 $\G^1=\Gamma$，连同\eref{10.17}中的高阶项 $\G^d$，构成（不分次）$\Ai$ 函子
\[
\G:\F(M)^{\mathrm{pr},1}\longrightarrow\F(M)^{\mathrm{pr},0}.
\]
\end{proposition}
注意，由\stref{引理}{10.12}，线性部分 $\G^1$ 都是同构；若通过\eref{10.11}识别 $s=0,1$ 时的 Floer 上链空间 $CF^{\mathrm{pr}}(L_0,L_1)^s$，则 $\G$ 是\secref{1c}意义下的形式
\SourcePage{156}
微分同胚。特别地，$\F(M)^{\mathrm{pr},0}$ 与 $\F(M)^{\mathrm{pr},1}$ 是拟同构的 $\Ai$ 范畴。
\begin{proof}[命题~\ref*{st:10.13}的证明概要]
证明策略基本上与\stref{引理}{10.12}相同。一维模空间 $\mathcal H^T(y_0,\ldots,y_d)^{[0,1],1}$ 的边界点出现于以下情形：(i) 离根最近的顶点 $v$ 满足 $s_v=0$；(ii) 离叶子最近的某个顶点 $v$ 满足 $s_v=1$；或 (iii) 条件\eref{10.16}中的某个不等式成为等式。在前两种情形中，像\sr{fig:10.4}{图~10.4}那样将树分成若干部分，便得出相关代数贡献恰好是 $\Ai$ 函子方程\eref{1.6}中的两项。情形 (iii) 的贡献仍与 Gromov 紧化引入的额外点相抵消。
\SourceFigure{figure-10-4}{0.65215314}{10.4}
当然，这个形式上的纲要并不构成完整证明。不过，其中的大部分分析并不比定义 Fukaya 范畴时使用的分析复杂多少。因此，我们只讨论横截性问题的一个方面，即促使我们引入延迟函数的那个方面。为更清楚地说明这一点，暂时只考虑平凡延迟函数 $\sigma_k(r,s)=s$。

首先，可以一般地选择\eref{10.13}中的扰动数据，使所有空间 $\M^{d+1}(y_0,\ldots,y_d)^{[0,1]}$ 都正则。剩下的问题，是保证 $\mathcal H^T(y_0,\ldots,y_d)^{[0,1]}$ 的边界面仍然正则；这些边界面是乘积
\begin{equation}\label{eq:10.18}
\prod_{v\in\operatorname{Ve}(T)}\M^{|v|}(\widetilde y_{f_0(v)},\ldots,\widetilde y_{f_{|v|-1}(v)})^{[0,1]}
\end{equation}
中由等式
\begin{equation}\label{eq:10.19}
s_{v^+(e)}=s_{v^-(e)}
\end{equation}
截出的子集。只需要考虑维数不超过 $1$ 的模空间，使问题稍微容易一些。例如，在零维情形中，
\SourcePage{157}
所需陈述是：若 $\{u_v:S_v\to M\}\in\mathcal H^T(y_0,\ldots,y_d)^{[0,1]}$ 满足每个 $(r_v,s_v,u_v)$ 都是\eref{10.18}中相应因子的孤立点，则对任何边 $e$，\eref{10.19}都不成立。标准横截性论证说明，除非对应于相关顶点的两个映射相同，即
\begin{equation}\label{eq:10.20}
S_{v^-(e)}\cong S_{v^+(e)},\qquad u_{v^-(e)}=u_{v^+(e)},
\end{equation}
否则可以通过一般地选择扰动数据实现这一点。不幸的是，与前面对映射 $\Gamma$ 的讨论中相应的步骤不同，这个剩余情形根本无法排除，因为存在正则的伪全纯 $(d+1)$ 边形，其入向端的极限与一个或多个出向端的极限相同。现在允许非平凡的延迟函数；此时不再考虑\eref{10.19}，而是考虑\eref{10.16}中某些不等式成为等式的情形。一般地会有如下性质：\emph{若 $(r,s,u)$ 是某个模空间 $\M^{d+1}(y_0,\ldots,y_d)^{[0,1]}$（$d\geq2$）的孤立点，则延迟函数对所有 $k=1,\ldots,d$ 都满足 $\sigma_k(r,s)\ne s$。}这个条件显然排除了\eref{10.20}这样的病态行为，从而可以完成横截性论证。\znote{149}{病态情形\eref{10.20}的实质是：两个相邻顶点的参数相等时，它们的解可以拼成同一个圆盘，于是所谓“边界层”并不真的落在边界上，横截性论证落空。延迟函数 $\sigma_k$ 的作用正是把上、下圆盘的参数故意错开，使 $\sigma_k(r,s)\ne s$ 恒成立；这样\eref{10.19}的等号条件在任何孤立点上都不出现，横截性论证才得以完成。注意 $\sigma_k$ 只改变族沿 $s$ 方向的参数化，不改变任何模空间本身，故不影响已证得的结构。[来源：原书 §(10e) 关于\eref{10.13}、\eref{10.16}、\eref{10.20} 的讨论。]}对于一维模空间，平行的讨论只略复杂一些，这里略去。
\end{proof}
\begin{note}\label{st:10.14}
Floer 同调的同伦方法已经出现在 Floer 最初的论文\bcite{44}中。向 Fukaya 范畴的推广遵循 Fukaya–Oh–Ohta–Ono\bcite{60}的定理 15.19，不过我们以更明确的横截性论证替代了虚扰动理论。顺带指出，在非正合情形的 Floer 同调中，横截性问题更为严重，目前只找到间接的解决办法；见\bcite{94}，或参见\bcite{78}中作为模型的 Morse–Novikov 理论。
\end{note}

\clearpage
\section{指标与行列式线}\label{sec:11}\label{section:11}
\subsection{行列式线}\label{sec:11a}
本节的讨论主要围绕行列式线展开。行列式线是对应于 Fredholm 算子的一维实向量空间。实际上，我们关心的是这条线的两种可能定向所成的集合，而非线本身。特别地，每当我们说两个行列式线（或稍后引入的类似对象）之间的同构是典范的，总是指在相差一个正实数倍的意义下典范。作了这一措辞上的说明之后，设 $D$ 为两个实 Banach 空间之间的 Fredholm 算子。$D$ 的行列式线定义为其核与余核的最高次外幂的张量积，确切地说，
\begin{definition*}
\[
\det(D)=\lambda^{\mathrm{top}}(\operatorname{coker}(D)^\vee)\otimes\lambda^{\mathrm{top}}(\ker(D)).
\]
\end{definition*}
\SourcePage{158}
由于只关心定向，对余核取对偶是多余的；但我们仍愿意保留这种写法，使公式与复行列式线的公式相对应。\znote{150}{行列式线是 Fredholm 指标理论的标准工具。对 Fredholm 算子 $D:E\to F$，取 $\lambda(\ker D)\otimes\lambda(\operatorname{coker}D)^\vee$（$\lambda$ 表示最高次外幂）为它的行列式线；当算子连续形变时，这些一维空间拼成底空间上的线丛，即行列式线丛。nLab 把这个构造从有限维的行列式 $\det T:\wedge^nV\to\wedge^nW$ 出发加以推广，指出向量丛的最高次外幂丛正是行列式线丛，并由此讨论 Quillen 行列式线丛以及 Grassmann 流形上的行列式丛，可用作理解本条定义的背景。[来源：nLab "determinant line bundle"，https:/\allowbreak /\allowbreak ncatlab.org/\allowbreak nlab/\allowbreak show/\allowbreak determinant+\allowbreak line+\allowbreak bundle（2026-11-15 访问）]}
存在典范同构
\begin{equation}\label{eq:11.1}
\det(D_1)\otimes\det(D_2)\cong\det(D_1\oplus D_2),
\end{equation}
其定义是将所有向量视为奇向量，并使用 Koszul 符号规则。具体地，若 $\{v_{k,\allowbreak i}\}$ 是 $\ker(D_k)$ 的基，$\{w_{k,\allowbreak i}^\vee\}$ 是 $\operatorname{coker}(D_k)^\vee$ 的基，则同构为
\begin{equation}\label{eq:11.2}
\begin{aligned}
&\left(\bigwedge_iw_{1,i}^\vee\otimes\bigwedge_iv_{1,i}\right)
\otimes\left(\bigwedge_iw_{2,i}^\vee\otimes\bigwedge_iv_{2,i}\right)\\
&\quad\longmapsto(-1)^{\dim(\operatorname{coker}D_2)\cdot\operatorname{index}(D_1)}
\bigwedge_iw_{2,i}^\vee\wedge\bigwedge_iw_{1,i}^\vee
\otimes\bigwedge_iv_{1,i}\wedge\bigwedge_iv_{2,i}.
\end{aligned}
\end{equation}
它严格满足结合律，并在分次意义下满足交换律。后者指的是，两个同构的复合
\[
\det(D_1)\otimes\det(D_2)\cong\det(D_1\oplus D_2)
=\det(D_2\oplus D_1)\cong\det(D_2)\otimes\det(D_1)
\]
与直接交换因子的映射相差符号 $(-1)^{\operatorname{index}(D_1)\operatorname{index}(D_2)}$。

暂且考虑有限维向量空间之间的线性映射 $D:H\to K$。可以用唯一的方式选择识别 $t_D:\det(D)\to\det(0)=\lambda^{\mathrm{top}}(K^\vee)\otimes\lambda^{\mathrm{top}}(H)$，使之与 $D$ 在左、右两侧复合同构相容，并满足以下三条额外规则：若 $D=0$，则 $t_D$ 为恒等映射；若 $D=\id:H\to H$，其中 $H$ 为一维向量空间，则 $t_D:\R=\det(\id)\to H^\vee\otimes H$ 将 $1$ 映为 $v^\vee\otimes v$，这里 $\langle v^\vee,\allowbreak v\rangle>0$；最后，这些同构与直和相容，也就是说下图交换：
\SourceDiagram{diagram-11a}{0.46523027}
其中 $0_k$ 是与 $D_k$ 具有相同定义域和目标空间的零映射。另一种做法，是对任意 $D:H\to K$ 直接定义 $t_D$：取 $H$ 的基 $e_1,\allowbreak \ldots,\allowbreak e_m$ 与 $K$ 的基 $f_1,\allowbreak \ldots,\allowbreak f_n$，满足 $D(e_1)=f_1,\allowbreak \ldots,\allowbreak D(e_k)=f_k$、$D(e_{k+1})=0,\allowbreak \ldots,\allowbreak D(e_m)=0$，再取它们的对偶基 $e_i^\vee,\allowbreak f_i^\vee$。则
\begin{equation}\label{eq:11.3}
\begin{aligned}
t_D:\ &(f_n^\vee\wedge\cdots\wedge f_{k+1}^\vee)\otimes(e_{k+1}\wedge\cdots\wedge e_m)\in\det(D)\\
&\longmapsto(f_n^\vee\wedge\cdots\wedge f_1^\vee)\otimes(e_1\wedge\cdots\wedge e_m)\in\det(0).
\end{aligned}
\end{equation}
有了这些准备，回到一般情形。设 $\mathfrak F(H,\allowbreak K)$ 是实 Banach 空间 $H,\allowbreak K$ 之间所有 Fredholm 算子所成的空间。\emph{行列式线丛} $\underline{\det}\to\mathfrak F(H,\allowbreak K)$ 的纤维为 $\det(D)$，其拓扑结构可由下述性质刻画。首先，若 $H,\allowbreak K$ 有限维，同构 $t_D$ 应给出 $\underline{\det}$ 的一个平凡化。其次，上面定义的映射
\SourcePage{159}
$\det(D_1)\otimes\det(D_2)\cong\det(D_1\oplus D_2)$，是在 $\mathfrak F(H_1,\allowbreak K_1)\times\mathfrak F(H_2,\allowbreak K_2)$ 上的行列式线丛同构。最后，若限制到满射算子的子空间，则等式 $\det(D)=\lambda^{\mathrm{top}}(\ker(D))$ 给出行列式线丛与核向量丛的最高次外幂之间的同构；这是常用的通过稳定化定义行列式线丛的方法的一个小变体。不必固定两个 Banach 空间，也可以考虑两个 Banach 向量丛以及它们之间的 Fredholm 映射，从而得到底空间上的行列式线丛。
\begin{note}\label{st:11.1}
这里引入的符号约定采用\bcite{87}的约定；亦见\bcite{77}附录 D。在\bcite{36}第一卷第 62、362 页中，它被称为“观点 II”。提一下另一种可能的约定，也许有助于理解：将\eref{11.2}右侧换成（即去掉符号）
\[
\bigwedge_iw_{1,i}^\vee\wedge\bigwedge_iw_{2,i}^\vee\otimes\bigwedge_iv_{1,i}\wedge\bigwedge_iv_{2,i}
\]
并将\eref{11.3}中 $f_i^\vee$ 的次序反转。这样得到另一个线丛 $\underline{\det}'$，它的纤维与 $\underline{\det}$ 相同；两个线丛同构，却并非以平凡的方式同构，恒等映射 $\underline{\det}\to\underline{\det}'$ 并不连续。
\end{note}

\subsection{路径空间上的线丛与谱流}\label{sec:11b}
设 $N$ 为闭流形，$\Sigma\subset N$ 为超曲面，其上有一个实线丛 $\sigma\to\Sigma$，并选定同构
\begin{equation}\label{eq:11.4}
(\sigma^{\otimes2})^\vee\cong\nu_\Sigma.
\end{equation}
这里 $\nu_\Sigma$ 是法丛，因此\eref{11.4}赋予 $\Sigma$ 一个余定向。由这些数据，可以在光滑路径 $[0;1]\to N$ 所成的空间 $\mathcal P$ 上得到局部常值函数 $I:\mathcal P\to\Z$ 和实线丛 $\delta\to\mathcal P$；这些路径的两个端点都位于 $\Sigma$ 之外。用上同调的语言说，就是将 Thom 同构与求值映射 $\mathcal P\times[0;1]\to N$ 结合起来：若
\[
\begin{aligned}
S:H^*(\Sigma)&\cong H^{*+1}(N,N\setminus\Sigma)\\
&\longrightarrow H^{*+1}(\mathcal P\times[0;1],\mathcal P\times\{0;1\})\cong H^*(\mathcal P),
\end{aligned}
\]
则 $I=S(1)$，同样有 $w_1(\delta)=S(w_1(\sigma))$。事实上，$I$ 就是与 $[\Sigma]$ 的相交数，因而在几何上可以解释为交点的计数。下面沿着同一思路，说明如何显式构造 $\delta$ 的一个代表。

取一条与 $\Sigma$ 横截相交的路径 $\alpha\in\mathcal P$。在有限多个穿越时刻 $s\in\alpha^{-1}(\Sigma)\subset(0;1)$ 处，都有非零法向量 $\alpha'(s)$，所以由\eref{11.4}得到 $\sigma_{\alpha(s)}$ 上的非退化二次型 $q_\alpha(s)$。它的符号就是与 $\Sigma$ 的局部相交数，因此 $I(\alpha)=\sum_s\operatorname{sign}(q_\alpha(s))$。定义
\begin{equation}\label{eq:11.5}
\delta_\alpha=\bigotimes_s\sigma_{\alpha(s)}^{\operatorname{sign}(q_\alpha(s))},
\end{equation}
其中 $\sigma^{-1}$ 是逆线丛，即对偶丛；张量积约定按穿越时刻 $s$ 的递增次序排列。要使 $\delta_\alpha$ 成为一个线丛的纤维，
\SourcePage{160}
就要对一般位置的同伦 $\{\alpha_r\}$、$0\leq r\leq1$ 定义平行移动映射 $\delta_{\alpha_0}\to\delta_{\alpha_1}$。其中非平凡的一步，是规定当交集 $\alpha_r\cap\Sigma$ 发生一般位置的分岔，即交点的产生或消失时，映射应如何变化。假设在 $(\tilde r,\allowbreak \tilde s)$ 处产生一对交点。在该点，$(d/ds)^2\alpha_r$ 给出一个定义良好的非零法向量，从而在 $\tilde\sigma=\sigma_{\alpha_{\tilde r}(\tilde s)}$ 上得到非零二次型 $\tilde q$。取充分小的 $\epsilon>0$，令 $r_\pm=\tilde r\pm\epsilon$，并考虑刚产生的两个交点，其时刻为 $(r_+,\allowbreak s_\pm)$，其中 $s_-<s_+$。它们靠近分岔点；若用局部平凡化将这两点的 $\sigma$ 纤维与 $\tilde\sigma$ 识别，则 $q_{\alpha_{r_+}}(s_\pm)$ 的符号为 $\pm\operatorname{sign}(\tilde q)$。其他交点的符号继承自 $\alpha_{r_-}$，因此有选定的同构
\begin{equation}\label{eq:11.6}
\begin{aligned}
\delta_{\alpha_{r_+}}\cong{}&
\bigotimes_{s<\tilde s}\sigma_{\alpha_{r_-}(s)}^{\operatorname{sign}(q_{\alpha_{r_-}}(s))}
\otimes\tilde\sigma^{-\operatorname{sign}(\tilde q)}\otimes\tilde\sigma^{\operatorname{sign}(\tilde q)}\\
&\otimes\bigotimes_{s>\tilde s}\sigma_{\alpha_{r_-}(s)}^{\operatorname{sign}(q_{\alpha_{r_-}}(s))}.
\end{aligned}
\end{equation}
取 $\tilde\sigma$ 中的非零向量 $\tilde v$ 及其在 $\tilde\sigma^{-1}$ 中的对偶向量 $\tilde v^{-1}$；依照 $\tilde q$ 的符号，在\eref{11.6}中插入 $\tilde v^{-1}\otimes\tilde v$ 或 $\tilde v\otimes\tilde v^{-1}$，便定义了平行移动 $\delta_{\alpha_{r_-}}\to\delta_{\alpha_{r_+}}$。对交点消失的过程，使用逆向规则。要证明这确实定义了路径空间上的线丛，还应验证关于“同伦的同伦”的一致性条件；这很直接，故省略。

典型例子如下。令 $N=\operatorname{sym}^2(\R^n)^\vee$ 为 $\R^n$ 上二次型所成的空间，按余秩分层为 $N=\Sigma^0\cup\Sigma^1\cup\cdots$，其中 $\Sigma^0$ 是非退化二次型组成的开稠密子集。各二次型的零空间构成向量丛 $\sigma^d\to\Sigma^d$，并有自然同构 $(\operatorname{sym}^2\sigma^d)^\vee\cong\nu_{\Sigma^d}$。将上述构造用于 $\Sigma=\overline{\Sigma^1}=\Sigma^1\cup\Sigma^2\cup\cdots$ 和 $\sigma=\sigma^1$，便在端点位于 $\Sigma^0$ 中的路径空间 $\mathcal P$ 上得到函数 $I$ 与线丛 $\delta$。这里稍有困难，因为 $\Sigma$ 在余维二处存在奇性，而 $\sigma$ 只在其光滑部分上定义。处理方法是先对始终位于 $\Sigma^0\cup\Sigma^1$ 中的路径 $\alpha$ 定义 $\delta$，再通过涉及下一层 $\Sigma^2$ 的链环的一项计算，证明它能延拓到整个 $\mathcal P$。细节省略。

与行列式线丛的联系来自谱流构造。\znote{152}{谱流是自伴 Fredholm 算子连续族的整值不变量：沿参数变化时，其特征值穿过零的次数按方向带符号计数，从负变正贡献 $+1$，从正变负贡献 $-1$；对闭路径，它就是绕行一圈后净穿过零的特征值个数。nLab 指出，这一概念最早出现在引入 $\eta$ 不变量的 Atiyah–Patodi–Singer 工作中（1975），其后由 Phillips、Booß-Bavnbek–Lesch–Phillips 等人对无界自伴算子作了系统处理。正文中的 $I(\alpha)$ 正是这一思想在算子族 $d_\alpha$ 上的实现。[来源：nLab "spectral flow"，https:/\allowbreak /\allowbreak ncatlab.org/\allowbreak nlab/\allowbreak show/\allowbreak spectral+\allowbreak flow（2026-11-15 访问）]}
固定 $\R^n$ 上的内积，以及函数 $\psi\in C^\infty(\R,\allowbreak [0;1])$，满足：$s\ll0$ 时 $\psi(s)=0$，$s\gg0$ 时 $\psi(s)=1$，且在 $\psi(s)\in(0;1)$ 的所有点都有 $\psi'(s)>0$。每个 $\alpha\in\mathcal P$ 都给出 Fredholm 算子
\[
d_\alpha=d/ds+\alpha(\psi(s))\in\mathfrak F(W^{1,2}(\R,\R^n),L^2(\R,\R^n)).
\]
\begin{lemma}\label{st:11.2}
$I(\alpha)=\operatorname{index}(d_\alpha)$；此外，$\delta$ 与算子族 $\{d_\alpha\}$ 的行列式线丛典范同构。
\end{lemma}
\begin{proof}
取与各层 $\Sigma^d$ 横截的 $\alpha\in\mathcal P$；特别地，对所有 $d\geq2$ 都有 $\alpha^{-1}(\Sigma^d)=\varnothing$。以 $\nu_{\alpha(s)}$ 表示对称矩阵 $\alpha(s)$ 的负特征子空间。首先扩张 $d_\alpha$ 的定义域，加入函数
\SourcePage{161}
$e^{-s\lambda}v$ 的线性张成空间，其中 $\lambda$ 为 $\alpha(1)$ 的任意负特征值，$v$ 为对应特征向量。所得算子
\[
W^{1,2}(\R,\R^n)\oplus\nu_{\alpha(1)}\longrightarrow L^2(\R,\R^n)
\]
为满射，因为可从 $-\infty$ 积分到 $s$ 求得原像；其核可通过在靠近 $-\infty$ 的某点求值而与 $\nu_{\alpha(0)}$ 识别。于是
\begin{equation}\label{eq:11.7}
\begin{aligned}
\dim\nu_{\alpha(0)}&=\operatorname{index}(d_\alpha)+\dim\nu_{\alpha(1)},\\
\lambda^{\mathrm{top}}(\nu_{\alpha(0)})&\cong\det(d_\alpha)\otimes\lambda^{\mathrm{top}}(\nu_{\alpha(1)}),
\end{aligned}
\end{equation}
其中同构是典范的。在 $\alpha$ 穿越 $\Sigma^1$ 的有限多个时刻 $s_j$（$1\leq j\leq m$），都有一维零空间 $\sigma_{\alpha(s_j)}$，带有非退化二次型 $q_\alpha(s_j)$，它正是 $\alpha'(s_j)$ 的限制。容易看出，对充分小的 $\epsilon>0$ 有同构
\[
\begin{cases}
\sigma_{\alpha(s_j)}\oplus\nu_{\alpha(s_j+\epsilon)}\cong\nu_{\alpha(s_j-\epsilon)},&q_\alpha(s_j)>0,\\
\nu_{\alpha(s_j+\epsilon)}\cong\sigma_{\alpha(s_j)}\oplus\nu_{\alpha(s_j-\epsilon)},&q_\alpha(s_j)<0.
\end{cases}
\]
在所有其他时刻，$\nu_{\alpha(s)}$ 随 $s$ 光滑变化；沿路径追踪这些变化便得到
\[
\begin{aligned}
\dim\nu_{\alpha(0)}&=\sum_j\operatorname{sign}(q_\alpha(s_j))+\dim\nu_{\alpha(1)},\\
\lambda^{\mathrm{top}}(\nu_{\alpha(0)})&\cong\bigotimes_j\sigma_{\alpha(s_j)}^{\operatorname{sign}(q_\alpha(s_j))}
\otimes\lambda^{\mathrm{top}}(\nu_{\alpha(1)}).
\end{aligned}
\]
结合\eref{11.7}，即得所需等式与同构。全部构造只是线性代数，因此容易验证此同构与\eref{11.6}中的平行移动相容。\znote{153}{谱流只依赖算子族两端，在族的形变下保持不变，这正是正文断言 $I(\alpha)$ 仅由 $\alpha$ 的同伦类决定的原因。Robbin 与 Salamon 证明，把算子族与相应的 Lagrangian 子空间族对应起来之后，谱流恰好等于该路径的 Maslov 指标；他们在《The spectral flow and the Maslov index》（Bull. London Math. Soc. 27，1995）中系统论述了这一对应，同作者早前的《The Maslov index for paths》（Topology 32，1993）则给出了路径 Maslov 指标的严格定义。两文均见于 nLab 的 Maslov 指标条目所附文献。[来源：nLab "Maslov index"，https:/\allowbreak /\allowbreak ncatlab.org/\allowbreak nlab/\allowbreak show/\allowbreak Maslov+\allowbreak index（2026-11-15 访问）]}
\end{proof}

下面把上述论证改写成更适合推广到无穷维情形的形式。首先，$\delta_\alpha$ 与 $d_\alpha$ 的定义显然可推广到由任意闭区间 $C$ 参数化的路径 $\alpha$。在满足 $\alpha(s)\notin\Sigma$ 的点 $s\in\operatorname{int}(C)$ 处切开，将 $C$ 分成 $C_1\cup C_2$。依定义，$I(\alpha)=I(\alpha|C_1)+I(\alpha|C_2)$，并有典范同构 $\delta_\alpha\cong\delta_{\alpha|C_1}\otimes\delta_{\alpha|C_2}$。另一方面，由\eref{11.7}可知
\begin{equation}\label{eq:11.8}
\begin{aligned}
\operatorname{index}(d_\alpha)&=\operatorname{index}(d_{\alpha|C_1})+\operatorname{index}(d_{\alpha|C_2}),\\
\det(d_\alpha)&\cong\det(d_{\alpha|C_1})\otimes\det(d_{\alpha|C_2}).
\end{aligned}
\end{equation}
作一个不改变交点 $\alpha(s_j)$ 及该处相交符号的小变形后，可以假定对每个 $j$，都存在邻域
\SourcePage{162}
$[s_j-\epsilon;s_j+\epsilon]$，使其中的 $\alpha(s)$ 可同时对角化。将 $[0;1]$ 分成如下区间：
\[
\begin{aligned}
C_0&=[0;s_1-\epsilon],&C_1&=[s_1-\epsilon;s_1+\epsilon],\\
C_2&=[s_1+\epsilon;s_2-\epsilon],&\ldots&&\\
C_{2m}&=[s_m+\epsilon;1].
\end{aligned}
\]
由于可以同时对角化，算子 $d_{\alpha|C_{2j-1}}$ 可化为标量算子并作显式分析。结果是，每个这样的算子或者单射，或者满射，其一维核或余核（由相交符号决定）可直接与 $\sigma_{\alpha(s_j)}$ 识别。对于其他区间，有将 $d_{\alpha|C_{2j}}$ 变形成可逆算子的一个选定方法：只需在 $N\setminus\Sigma$ 内，将 $\alpha|C_{2j}$ 收缩为取值 $\alpha(s)$ 的常值路径，其中 $s\in C_{2j}$。最后运用\eref{11.8}将各部分重新组合，便完成论证。

\subsection{\texorpdfstring{$\bar\partial$}{∂̄} 算子的行列式线}\label{sec:11c}
需要回顾若干熟知的指标理论结果，它们可视为\eref{11.8}的二维对应物。设 $S_k$（$k=1,\allowbreak 2$）为带条带状端的两个 Riemann 曲面，其上有在\secref{8h}意义下非退化的 Cauchy–Riemann 算子 $\bar\partial_{\nabla_k}$。假设 $S_1$ 有出向端 $\zeta_1^+$，$S_2$ 有入向端 $\zeta_2^-$，且两端上的极限数据一致。于是可粘合向量丛及其他数据，当 $l\gg0$ 时，在曲面 $S=S_1\#_lS_2$ 上得到非退化 Cauchy–Riemann 算子 $\bar\partial_\nabla=\bar\partial_{\nabla_1}\#_l\bar\partial_{\nabla_2}$。该算子在相差小扰动的意义下唯一；特别地，其行列式线 $\det(\bar\partial_\nabla)$ 与构造中的选择无关。粘合定理断言
\begin{equation}\label{eq:11.9}
\begin{aligned}
\operatorname{index}(\bar\partial_\nabla)&=\operatorname{index}(\bar\partial_{\nabla_1})+\operatorname{index}(\bar\partial_{\nabla_2}),\\
\det(\bar\partial_\nabla)&\cong\det(\bar\partial_{\nabla_1})\otimes\det(\bar\partial_{\nabla_2}),
\end{aligned}
\end{equation}
其中第二式为典范同构，且对于两次粘合满足结合律。证明\eref{11.9}的标准方法是先假定两个组成算子均为满射。当 $l\gg0$ 时，可将 $\ker(\bar\partial_{\nabla_k})$ 中的元素拼接为方程 $\bar\partial_\nabla=0$ 的近似解，再经正交投影得到同构
\begin{equation}\label{eq:11.10}
\ker(\bar\partial_\nabla)\cong\ker(\bar\partial_{\nabla_1})\oplus\ker(\bar\partial_{\nabla_2}).
\end{equation}
对极限 $l\to\infty$ 运用紧性论证可知 $\operatorname{coker}(\bar\partial_\nabla)=0$，由此得出所需结论。满射假设可通过有限维稳定化论证去除。
\SourcePage{163}
偶尔还会用到另一种粘合方法，它在 Gromov–Witten 理论中更为常见。按粘合是在内部点还是边界点进行，有两个略有不同的版本；这里同时讨论。仍从两个带条带状端的曲面 $S_k$ 出发，其上有作用于向量丛 $E_k$ 的算子 $\bar\partial_{\nabla_k}$，并有 Lagrangian 子丛 $F_k$。这次不要求端上的渐近同构，而是假定有点 $z_k\in S_k$（均为内部点，或均为边界点），以及保持辛结构和复结构的纤维识别 $E_{1,\allowbreak z_1}\cong E_{2,\allowbreak z_2}$；边界情形还要求它将两个 Lagrangian 子空间彼此映到一起。取连通和，便得到带算子 $\bar\partial_\nabla$ 的 Riemann 曲面 $S$：当 $z_k$ 为内部点时取通常的连通和，否则取边界连通和。此时\eref{11.10}的对应映射只对在 $z_k$ 处取值相同的核元素对 $(X_1,\allowbreak X_2)$ 定义，所得\eref{11.9}的版本如下。对边界点，
\begin{equation}\label{eq:11.11}
\begin{aligned}
\operatorname{index}(\bar\partial_\nabla)&=\operatorname{index}(\bar\partial_{\nabla_1})+\operatorname{index}(\bar\partial_{\nabla_2})-n,\\
\det(\bar\partial_\nabla)\otimes\lambda^{\mathrm{top}}(F_{1,z_1})&\cong\det(\bar\partial_{\nabla_1})\otimes\det(\bar\partial_{\nabla_2});
\end{aligned}
\end{equation}
对内部点，
\begin{equation}\label{eq:11.12}
\begin{aligned}
\operatorname{index}(\bar\partial_\nabla)&=\operatorname{index}(\bar\partial_{\nabla_1})+\operatorname{index}(\bar\partial_{\nabla_2})-2n,\\
\det(\bar\partial_\nabla)\otimes\lambda^{\mathrm{top}}(E_{1,z_1})&\cong\det(\bar\partial_{\nabla_1})\otimes\det(\bar\partial_{\nabla_2}).
\end{aligned}
\end{equation}
利用辛形式赋予 $E_{1,\allowbreak z_1}$ 的定向，可将\eref{11.12}中的 $\lambda^{\mathrm{top}}(E_{1,\allowbreak z_1})$ 因子去掉。还要注意，内部点情形中，$S$ 与 $\bar\partial_\nabla$ 的构造依赖另一个角度参数，它记录切空间 $T(S_k)_{z_k}$ 的识别。不过，\eref{11.12}中的同构对该参数连续，因此改变角度会给出一个由 $S^1$ 参数化的算子族，其行列式线丛沿该族的单值变换平凡。
\begin{note}\label{st:11.3}
像\eref{11.9}这样的粘合公式，几乎出现于所有 Floer 同调型理论的构造中，例如见\bcite{119}、\bcite{48}及\bcite{37}第 3.3 节。此外，关于一般行列式线丛的切拼性质，还有大量文献；其中许多使用 Atiyah–Patodi–Singer 型边界条件，而非较简单但限制较强的管状端框架。入门介绍可见\bcite{122}第二讲，一般定理见\bcite{105}、\bcite{121}，Riemann 曲面情形见\bcite{77}附录 D。
\end{note}

\subsection{线丛上的\texorpdfstring{$\bar\partial$}{∂̄} 算子}\label{sec:11d}
在闭 Riemann 曲面上，负次数线丛没有非平凡全纯截面。一个标准论证将这一事实推广到带边界的 Riemann 曲面，以及形如
\SourcePage{164}
$\bar\partial+\{\text{零阶项}\}$ 的算子。具体而言，设 $S$ 为边界非空的紧连通 Riemann 曲面，$E\to S$ 为 Hermite 线丛，$F\subset E|\partial S$ 为实子丛，$\bar\partial_\nabla$ 为相应的 $\bar\partial$ 算子。选取 $E$ 的平凡化后，可将 $F$ 视为有限个映射 $\rho_1,\allowbreak \ldots,\allowbreak \rho_l:S^1\to\RP^1$，每个边界分支对应一个。用 $\mu(\rho_j)$ 表示它们的次数。假设 $\bar\partial_\nabla X=0$ 有一个非零解。其零点孤立，且每个零点的阶 $\nu(z)$ 有限；简单的绕数论证给出
\begin{equation}\label{eq:11.13}
2\sum_{z\in\operatorname{int}(S)}\nu(z)+\sum_{z\in\partial S}\nu(z)=\mu(\rho_1)+\cdots+\mu(\rho_l).
\end{equation}
特别地，
\begin{lemma}\label{st:11.4}
若 $\mu(\rho_1)+\cdots+\mu(\rho_l)<0$，则 $\bar\partial_\nabla$ 为单射。
\end{lemma}
对于带条带状端的曲面 $S$ 及其上的非退化 $\bar\partial_\nabla$ 算子，也有对应版本。为简单起见，假定在每个端 $\epsilon_\zeta:Z^\pm\to S$ 上，$(E,\allowbreak \nabla)$ 的拉回平凡，边界子丛局部常值。必要时将平凡化 $\epsilon_\zeta^*E\cong Z^\pm\times\CC$ 旋转一个常量角度后，可写成
\begin{equation}\label{eq:11.14}
F_{\epsilon_\zeta(s,0)}=\R,\qquad F_{\epsilon_\zeta(s,1)}=e^{i\sigma_\zeta}\R,
\end{equation}
其中 $\sigma_\zeta\in(-\pi;0)$。现在仍可为紧化 $\hat S$ 的每个边界分支得到一个映射 $\rho_j$：沿着 $F$ 走，并在每个无穷远点附近，用经过 $e^{it\sigma_\zeta}\R$、$t\in[0;1]$ 的同伦连接\eref{11.14}中的两条直线。
\begin{lemma}\label{st:11.5}
若 $\mu(\rho_1)+\cdots+\mu(\rho_l)-|\Sigma^-|<0$，则 $\bar\partial_\nabla$ 为单射。
\end{lemma}
\begin{proof}
考虑 $\bar\partial_\nabla X=0$ 的非零解。在每个条带状端上，该算子就是具有边界条件\eref{11.14}的标准算子 $\partial_s+i\partial_t$，因此 $X_\zeta(s,\allowbreak t)=X(\epsilon_\zeta(s,\allowbreak t))$ 可作 Fourier 展开
\[
X_\zeta(s,t)=\begin{cases}
\sum_{m>0}c_{\zeta,m}e^{(\sigma_\zeta+\pi m)z},&\zeta\in\Sigma^-,\\
\sum_{m\geq0}c_{\zeta,m}e^{(\sigma_\zeta-\pi m)z},&\zeta\in\Sigma^+,
\end{cases}
\]
其中系数 $c_{\zeta,\allowbreak m}\in\R$ 在 $m\to\infty$ 时快速衰减。令 $\nu(\zeta)$ 为满足 $c_{\zeta,\allowbreak m}\ne0$ 的最小 $m$。它确定了当 $|s|$ 足够大时，非零函数 $X_\zeta|\{s\}\times[0;1]$ 的拓扑行为；由此得到\eref{11.13}的修正项，给出
\[
2\sum_{z\in\operatorname{int}(S)}\nu(z)+\sum_{z\in\partial S}\nu(z)+\sum_{\zeta\in\Sigma}\nu(\zeta)=\mu(\rho_1)+\cdots+\mu(\rho_l).
\]
因为 $\zeta\in\Sigma^-$ 时 $\nu(\zeta)>0$，故结论成立。
\end{proof}
\SourcePage{165}
\begin{note}\label{st:11.6}
$\bar\partial_\nabla X=0$ 的非平凡解只有孤立零点、且每个零点的重数有限，这一事实是 Aronszajn 唯一延拓定理的推论：内部点见\bcite{49}，边界点见\bcite{91}。同一定理，或更直接的 Taylor 级数计算，还说明局部重数均为正。最后应当指出，即使不假设各条带状端上的数据平凡，\stref{引理}{11.5}仍成立。证明原理相同，但此时无法显式写出 $X_\zeta$，所以 $\nu(\zeta)$ 的定义需要更精细地分析 $X_\zeta$ 的渐近行为，见\bcite{114}。
\end{note}

\subsection{指标理论与 Lagrangian Grassmann 流形}\label{sec:11e}
设 $V=(V,\allowbreak \omega_V)$ 是带辛结构的 $2n$ 维向量空间，$\operatorname{Gr}(V)\cong U_n/O_n$ 是其无定向线性 Lagrangian 子空间的 Grassmann 流形。映射 $B\to\operatorname{Gr}(V)$ 等价于平凡辛向量丛 $B\times V\to B$ 的一个 Lagrangian 子丛，因此可以将 $\operatorname{Gr}(V)$ 的上同调类视为相对示性类。我们关心的是 Maslov 类 $\mu\in H^1(\operatorname{Gr}(V))$，以及由所有无定向子空间的通常 Grassmann 流形诱导的第二 Stiefel–Whitney 类 $w_2\in H^2(\operatorname{Gr}(V);\Z/2)$；第一 Stiefel–Whitney 类是 $\mu$ 的模 $2$ 约化，因此不必另行考虑。与 $\mu$ 配对给出同构 $\pi_1(\operatorname{Gr}(V))\to\Z$。当 $n\geq3$ 时，$w_2$ 同样给出同构 $\pi_2(\operatorname{Gr}(V))\to\Z/2$。低维例外是：$n=1$ 时 $U_1/O_1=S^1$，所以 $\pi_2(\operatorname{Gr}(V))=0$；$n=2$ 时 $U_2/O_2=S^1\times_{\Z/2}S^2$，所以 $\pi_2(\operatorname{Gr}(V))\cong\Z$。最后这个群可由二重覆盖 $U_2/SO_2$ 上的 Euler 类检测，而该类模 $2$ 约化为第二 Stiefel–Whitney 类，所以 $w_2:\pi_2(U_2/O_2)\to\Z/2$ 仍为满射。\znote{158}{Maslov 类是 Lagrangian Grassmann 流形的基本上同调类。nLab 指出 $H^1(\mathrm{LGrass},\Z)\cong\Z$，其生成元称为万有 Maslov 类 $u\in H^1(\mathrm{LGrass},\Z)$；对 Lagrangian 子流形 $Y$，把切丛的分类映射 $i:Y\to\mathrm{LGrass}$ 拉回，就得到 $i^*u\in H^1(Y,\Z)$，即 $Y$ 的 Maslov 类。这一上同调解释出自 Arnol'd 1967 年的工作，它把 Maslov 原来的渐近分析不变量提升为万有示性类。正文中 $\mu\in H^1(\operatorname{Gr}(V))$ 正是该类的具体化，因此与 $\mu$ 配对给出同构 $\pi_1(\operatorname{Gr}(V))\cong\Z$。对一条环路而言，配对得到的整数就是这条环路的 Maslov 指标：它把环路对应的一族 Lagrangian 子空间与一个固定参考子空间作交，再按穿越方向带符号求和；所得整数模 $2$ 约化即第一 Stiefel–Whitney 类，故正文说后者不必另行考虑。[来源：nLab "Maslov index"，https:/\allowbreak /\allowbreak ncatlab.org/\allowbreak nlab/\allowbreak show/\allowbreak Maslov+\allowbreak index（2026-11-15 访问）]}

还需要考察带基点环路空间 $\Omega\operatorname{Gr}(V)$ 与自由环路空间 $\mathcal L\operatorname{Gr}(V)$。它们各有以 $\Z$ 编号的连通分支 $\Omega_k\operatorname{Gr}(V)$ 或 $\mathcal L_k\operatorname{Gr}(V)$，由 Maslov 指标 $\mu(\rho)=\langle\mu,\allowbreak [\rho]\rangle=k$ 区分。当 $n\geq3$ 时，容易看出 $\pi_1(\Omega_k\operatorname{Gr}(V))\cong\Z/2$，并且
\[
\pi_1(\mathcal L_k\operatorname{Gr}(V))\cong\pi_2(\operatorname{Gr}(V))\times\pi_1(\operatorname{Gr}(V))\cong\Z/2\times\Z.
\]
特别地，
\begin{equation}\label{eq:11.15}
H^1(\mathcal L_k\operatorname{Gr}(V);\Z/2)\cong\Z/2\oplus\Z/2,
\end{equation}
其显式生成元为 $T(w_2)$ 与 $U(\mu)$ 的模 $2$ 约化。这里 $(T,\allowbreak U)$ 表示拉回映射
\[
H^*(\operatorname{Gr}(V))\longrightarrow H^*(S^1\times\mathcal L\operatorname{Gr}(V))
\cong H^{*-1}(\mathcal L\operatorname{Gr}(V))\oplus H^*(\mathcal L\operatorname{Gr}(V))
\]
的两个分量。

为与指标理论联系，作如下构造。固定 $V$ 上与 $\omega_V$ 相容的复结构 $I_V$、一个只有一个边界分支的紧 Riemann 曲面 $S$，以及保持定向的识别 $\partial S\cong S^1$。对每个 $\rho\in\mathcal L\operatorname{Gr}(V)$，令 $D_{S,\allowbreak \rho}$ 为作用于平凡复向量丛 $E=S\times V$ 的截面的 $\bar\partial$ 算子，边界条件由 Lagrangian 子丛 $F_\rho\subset E|\partial S$ 给出，$(F_\rho)_{\exp(2\pi is)}=\rho(s)$。考虑随 $\rho$ 变化的算子族 $D_{S,\allowbreak \rho}$，
\SourcePage{166}
便得到行列式线丛 $\underline{\det}_S\to\mathcal L\operatorname{Gr}(V)$。
\begin{lemma}\label{st:11.7}
$\operatorname{index}(D_{S,\allowbreak \rho})=n\chi(S)+\mu(\rho)$。此外，行列式线丛的同构类由下式给出：
\begin{equation}\label{eq:11.16}
w_1(\underline{\det}_S)=T(w_2)+(T(\mu)-1)U(\mu).
\end{equation}
\end{lemma}
\begin{proof}
第一条是熟知的带边界曲面的 Riemann–Roch 定理。第二条同样可视为“Real”算子族的一般指标定理\bcite{9}的特例。另一种论证不那么直接，却更为显式，它依赖以下几个要点。

\noindent(i) \emph{沿边界捏合分离。}假设 $S$ 的亏格大于零。可将它视为圆盘 $D$ 与同亏格闭 Riemann 曲面的连通和。闭曲面上的 $\bar\partial$ 算子显然是 $\CC$ 线性的，故其行列式线具有典范平凡化。于是由\eref{11.12}得到同构 $\underline{\det}_S\cong\underline{\det}_D$。

\noindent(ii) \emph{直和公式。}令 $\Delta:\operatorname{Gr}(V)\times\operatorname{Gr}(W)\to\operatorname{Gr}(V\oplus W)$ 为直和运算 $(\Lambda_1,\allowbreak \Lambda_2)\mapsto\Lambda_1\oplus\Lambda_2$，令
\[
\mathcal L_{k,l}\Delta:\mathcal L_k\operatorname{Gr}(V)\times\mathcal L_l\operatorname{Gr}(W)\longrightarrow\mathcal L_{k+l}\operatorname{Gr}(V\oplus W)
\]
为自由环路空间上的诱导映射。由\eref{11.1}，
\[
(\mathcal L_{k,l}\Delta)^*\underline{\det}_{S,V\oplus W}\cong\underline{\det}_{S,V}\boxtimes\underline{\det}_{S,W}.
\]
这里在记号中加入了所涉及的向量空间，以明确彼此关系。另一方面，运用 Whitney 和公式，
\[
\begin{aligned}
(\mathcal L_{k,l}\Delta)^*T(w_2)
&=T(w_2\times1+1\times w_2+\mu\times\mu)\\
&=T(w_2)\times1+1\times T(w_2)+U(\mu)\times T(\mu)+T(\mu)\times U(\mu)\\
&=(T(w_2)+lU(\mu))\times1+1\times(T(w_2)+kU(\mu)),\\
(\mathcal L_{k,l}\Delta)^*U(\mu)
&=U(\Delta^*\mu)=U(\mu)\times1+1\times U(\mu).
\end{aligned}
\]
因此
\[
\begin{aligned}
&(\mathcal L_{k,l}\Delta)^*(T(w_2)+(k+l-1)U(\mu))\\
&\quad=(T(w_2)+(k-1)U(\mu))\times1+1\times(T(w_2)+(l-1)U(\mu)).
\end{aligned}
\]
结合\eref{11.15}中生成元的表达式，上述计算还表明，只要 $V,\allowbreak W$ 中至少一个空间的维数不小于 $6$，$\mathcal L_{k,\allowbreak l}\Delta$ 在 $H^1(-;\Z/2)$ 上诱导的映射就是单射；事实上维数不小于 $4$ 已经足够，不过这里不强调这一点。

\noindent(iii) \emph{$n=1$ 时的计算。}假设 $n=1$，此时 Maslov 数 $\mu(\rho)$ 就是\secref{11d}考虑的绕数。例如，若 $S=D$ 且 $\mu(\rho)=-1$，则由\stref{引理}{11.4}，$D_{S,\allowbreak \rho}$ 为单射，再由指标
\SourcePage{167}
公式可知它可逆。再考虑 $S=D$、$\mu(\rho)=0$ 的情形。\eref{11.13}说明任意非零 $X\in\ker(D_{S,\allowbreak \rho})$ 都是处处非零的函数。这意味着任意两个非零核元素成比例，因为它们的商是 $S$ 上边界取实值的全纯函数，必为常数。因此，在一个边界点求值诱导典范同构 $\ker(D_{S,\allowbreak \rho})\cong\rho(0)$。另一方面，指标公式说明算子为满射，故也有 $\det(D_{S,\allowbreak \rho})\cong\rho(0)$。

\noindent(iv) \emph{一个高维情形。}考虑 $n\geq3$ 时的 $\mathcal L_0\operatorname{Gr}(V)$，并假设 $S=D$。\eref{11.15}的第二个直和因子来自常值环路所成的子空间 $\operatorname{Gr}(V)\subset\mathcal L_0\operatorname{Gr}(V)$。若 $\rho$ 为常值，$(E,\allowbreak F_\rho)$ 显然分解为一维直和分量，(iii) 的论证于是给出 $\det(D_{S,\allowbreak \rho})\cong\lambda^{\mathrm{top}}(\rho(0))$，这说明\eref{11.16}两边限制在 $\operatorname{Gr}(V)$ 上时一致。另一方面，可以构造与 $w_2$ 配对非零的映射 $S^2\to\operatorname{Gr}(V)$，将它视为环路 $S^1\to\Omega\operatorname{Gr}(V)$ 时，所得 $S^1$ 上的行列式线丛非平凡。\bcite{60}第 6 章或\bcite{34}详细进行了这一构造，这里不再重复。这两个观察合起来表明，在 $\mathcal L_0\operatorname{Gr}(V)$ 上有 $w_1(\underline{\det}_S)=T(w_2)+U(\mu)$。

有了这些准备，\eref{11.16}的证明如下。(i) 表明只需考虑 $S=D$。由 (iii)，已知等式在 $\mathcal L_0\operatorname{Gr}(\CC)$ 和 $\mathcal L_{-1}\operatorname{Gr}(\CC)$ 上成立。由 (ii)，若等式在 $\mathcal L_k\operatorname{Gr}(V\oplus\CC)$ 上成立，则它在 $\mathcal L_k\operatorname{Gr}(V)$ 和 $\mathcal L_{k+1}\operatorname{Gr}(V)$ 上也成立；此外，只要 $\dim(V)\geq6$，还可以反向推论。从 (iv) 出发并运用这些观察，便容易得出一般情形。
\end{proof}
\begin{note}\label{st:11.8}
将 Hilbert 空间族 $W^{1,\allowbreak 2}(D,\allowbreak E,\allowbreak F_\rho)$ 平凡化之后——由 Kuiper 定理，这总是可以做到——算子 $D_{D,\allowbreak \rho}$ 在固定 $\rho(0)$ 时定义映射 $\Omega(U_n/O_n)\to\mathfrak F(H,\allowbreak H)$。如\bcite{34}所述，应将它视为 Bott 周期性同伦等价
\[
\mathfrak F(H,H)\simeq\Z\times BO_\infty\simeq\Omega(U_\infty/O_\infty)
\]
的一个近似，其中允许指标相差一个常数平移。
\end{note}

\clearpage
\subsection{Arnol’d 分层}\label{sec:11f}
乘积 $\operatorname{Gr}(V)^2=\operatorname{Gr}(V)\times\operatorname{Gr}(V)$ 按相对位置有一个 Schubert 型分层。每层 $\Sigma^d=\{(\Lambda_0,\allowbreak \Lambda_1):\dim(\Lambda_0\cap\Lambda_1)=d\}$ 上，都有由这些交空间构成的秩 $d$ 向量丛 $\sigma^d$，并带典范同构
\begin{equation}\label{eq:11.17}
(\operatorname{sym}^2\sigma^d)^\vee\cong\nu_{\Sigma^d}.
\end{equation}
可以对 $\operatorname{Gr}(V)^2$ 运用\secref{11b}的路径空间构造，取 $\Sigma=\overline{\Sigma^1}$、$\sigma=\sigma^1$。由此得到局部常值整数函数 $I_Z$，称为\emph{路径的 Maslov 指标}，以及实线丛 $\delta_Z$；它们定义在路径空间 $\mathcal P^2\operatorname{Gr}(V)$ 上，其元素是满足 $s=0,\allowbreak 1$ 时 $\lambda_0(s)\pitchfork\lambda_1(s)$ 的路径 $(\lambda_0,\allowbreak \lambda_1):[0;1]\to\operatorname{Gr}(V)^2$。如同先前
\SourcePage{168}
二次型的情形——事实上该情形可作为局部模型——这里也有涉及 $\Sigma^2$ 的一个容易克服的小困难。

值得把上述构造说得更明确一些。取路径 $(\lambda_0,\allowbreak \lambda_1)$ 和 $s\in[0;1]$。对 $k=0,\allowbreak 1$ 及充分小的 $|r-s|$，选择线性映射族 $\phi_{k,\allowbreak r,\allowbreak s}:\lambda_k(s)\to\lambda_k(r)$，使 $\phi_{k,\allowbreak s,\allowbreak s}=\Id$。\emph{穿越型}是 $\lambda_0(s)\cap\lambda_1(s)$ 上如下定义的二次型：
\[
q_{\lambda_0,\lambda_1}(s):v\longmapsto-\left(\frac d{dr}\right)_{r=s}\omega_V(\phi_{0,r,s}(v),\phi_{1,r,s}(v)).
\]
它与 $\phi_{k,\allowbreak r,\allowbreak s}$ 的选择无关。用\eref{11.17}来表述，若 $(\lambda_0(s),\allowbreak \lambda_1(s))\in\Sigma^d$、$d>0$，则穿越型就是对应于 $(\lambda_0'(s),\allowbreak \lambda_1'(s))$ 法向分量的二次型。对一般位置的路径，空间 $\lambda_0(s)\cap\lambda_1(s)$ 至多一维，而一维交空间上的穿越型非零。\znote{160}{穿越型是 Robbin–Salamon 定义路径 Maslov 指标的出发点。对两族 Lagrangian 子空间 $\lambda_0(s),\lambda_1(s)$，在交点 $s$ 处取两条曲线的速度向量，沿交空间得到一个二次型，其符号记录了交空间在 $s$ 附近如何随参数变化；把路径按穿越点切成小段后，每一段的贡献由该处穿越型的符号给出，Maslov 指标就等于这些符号之和，这与正文用一般机制\eref{11.5}从二次型分层得到的 $I_Z$ 一致。该构造见 Robbin–Salamon《The Maslov index for paths》（Topology 32，1993），nLab 的 Maslov 指标条目列出了该文。[来源：nLab "Maslov index"，https:/\allowbreak /\allowbreak ncatlab.org/\allowbreak nlab/\allowbreak show/\allowbreak Maslov+\allowbreak index（2026-11-15 访问）]}再应用一般机制\eref{11.5}，便得到 $I_Z(\lambda_0,\allowbreak \lambda_1)$ 和 $(\delta_Z)_{\lambda_0,\allowbreak \lambda_1}$。

与\secref{11b}平行，下面说明这些不变量的指标理论含义。选择该处所用的函数 $\psi$。对每个 $(\lambda_0,\allowbreak \lambda_1)\in\mathcal P^2\operatorname{Gr}(V)$，赋予算子 $D_{Z,\allowbreak \lambda_0,\allowbreak \lambda_1}$：它是无穷条带上平凡向量丛 $E=Z\times V$ 的 $\bar\partial$ 算子，边界条件由 Lagrangian 子丛 $F_{\lambda_0,\allowbreak \lambda_1}\subset E|\partial Z$ 给出，其中 $(F_{\lambda_0,\allowbreak \lambda_1})_{(s,\allowbreak k)}=\lambda_k(\psi(s))$。由于端点处的横截条件，它是非退化 Cauchy–Riemann 算子，因而为 Fredholm 算子。于是有
\begin{lemma}\label{st:11.9}
$I_Z(\lambda_0,\allowbreak \lambda_1)=\operatorname{index}(D_{Z,\allowbreak \lambda_0,\allowbreak \lambda_1})$；且 $(\delta_Z)_{\lambda_0,\allowbreak \lambda_1}$ 与 $\det(D_{Z,\allowbreak \lambda_0,\allowbreak \lambda_1})$ 典范同构。
\end{lemma}
\begin{proof}
证明策略与\stref{引理}{11.2}证明的第二种表述相同，故不重复整个论证。这里直接使用粘合公式\eref{11.9}，代替导出\eref{11.8}的初等常微分方程计算。剩下的任务是找出足够多能够显式写出指标与行列式线的例子。一种方法如下；还可采用其他方法，其中有些在技术上稍简单。假设 $V=\R^{2n}=\CC^n$，带标准辛结构和标准复结构 $I_V=i$；对所有 $s$，$\lambda_1(s)=\R^n\times\{0\}$；$\lambda_0(s)$ 是图像 $\{x_2=\alpha(s)x_1\}$，其中 $\alpha(s)$ 为一族对称 $n\times n$ 矩阵。等价地，$\lambda_0(s)=\exp(iA(s))\R^n$，其中
\[
A(s)=\begin{pmatrix}\alpha(s)&0\\0&0\end{pmatrix}.
\]
若对 $(s,\allowbreak t)\in Z$ 应用辛线性映射族 $\Phi(s,\allowbreak t)=\exp((1-t)iA(\psi(s)))$ 改变 $E$ 的平凡化，算子便变为
\begin{equation}\label{eq:11.18}
\begin{aligned}
D_{Z,\lambda_0,\lambda_1}X={}&(\partial_s-J(s,t)iA(\psi(s)))X+J(s,t)\partial_tX\\
&+\Phi(s,t)^{-1}(\partial_s\Phi(s,t))X,
\end{aligned}
\end{equation}
其中 $J(s,\allowbreak t)=\Phi(s,\allowbreak t)^{-1}i\Phi(s,\allowbreak t)$ 是 $\R^{2n}$ 上一族相容复结构。可以将 $J(s,\allowbreak t)$ 连续变形回常值族 $i$；这
\SourcePage{169}
给出\eref{11.18}的一条 Fredholm 算子同伦，故指标与行列式线（在典范同构意义下）不变。同样，\eref{11.18}的最后一项紧支且为零阶项，因而是紧算子，也可通过变形去掉。最终得到
\[
\widetilde D_ZX=(\partial_s+A(\psi(s)))X+i\partial_tX,
\]
以及常值边界条件 $X(s,\allowbreak 0),\allowbreak X(s,\allowbreak 1)\in\R^n\times\{0\}$。椭圆正则性论证说明，每个核元素都可延拓为 $\R^2$ 上的光滑函数，满足 $X(s,\allowbreak t+2)=X(s,\allowbreak t)$、$X(s,\allowbreak -t)=\overline{X}(s,\allowbreak t)$；余核也有类似性质。将定义域与值域限制到这些子空间后，可沿 $t$ 方向作 Fourier 变换，把 $\widetilde D_Z$ 分解为如下单变量微分算子的直和：
\begin{equation}\label{eq:11.19}
\begin{gathered}
d/ds+\alpha(\psi(s)),\qquad
d/ds+\begin{pmatrix}\alpha(\psi(s))&\pi l\cdot\Id\\\pi l\cdot\Id&0\end{pmatrix},\\
l=1,2,\ldots.
\end{gathered}
\end{equation}
\secref{11b}已经研究过这类算子的指标理论。对\eref{11.19}而言，唯一非平凡的贡献来自第一个直和分量，更确切地说，来自使 $\alpha(s)$ 的零空间非平凡的那些 $s$。简单计算表明，这些贡献恰是 $(\lambda_0,\allowbreak \lambda_1)$ 的穿越型。
\end{proof}
\begin{note}\label{st:11.10}
Lagrangian Grassmann 流形的分层由\bcite{2}引入。\stref{引理}{11.9}的第一条是 Floer 的指标定理\bcite{45}。总体而言，本节叙述紧随\bcite{112}、\bcite{113}。上面最后的论证可以视为从 Floer 理论约化到 Morse 理论\bcite{47}的线性版本。
\end{note}

\subsection{一个变体}\label{sec:11g}
现在略微修改前述构造，将条带 $Z$ 换成上半平面 $H$。以 $\mathcal P^-\operatorname{Gr}(V)$ 表示满足以下两个性质的路径 $\lambda:[0;1]\to\operatorname{Gr}(V)$ 所成的空间：
\begin{itemize}
\item $\lambda(0)$ 与 $\lambda(1)$ 横截相交。
\item 路径对 $(\lambda,\allowbreak \lambda(1))$ 在 $s=1$ 处的穿越型负定，其中第二个分量表示取值为 $\lambda(1)$ 的常值路径。
\end{itemize}
该空间上有自然的局部常值函数 $I_H$ 和线丛 $\delta_H$。为归结到先前的框架，注意上述第二个条件保证存在 $s'<1$，使对所有 $s\in[s';1)$ 都有 $\lambda(s)\pitchfork\lambda(1)$；特别地，将 $(\lambda,\allowbreak \lambda(1))$ 限制到 $[0;s']$，便得到 $\mathcal P^2\operatorname{Gr}(V)$ 的一个元素。可以选取光滑依赖于 $\lambda$ 的 $s'$，于是限制给出映射 $\mathcal P^-\operatorname{Gr}(V)\to\mathcal P^2\operatorname{Gr}(V)$。定义 $I_H,\allowbreak \delta_H$ 为 $I_Z,\allowbreak \delta_Z$ 沿此映射的拉回。更明确地，
\SourcePage{170}
用 $q_\lambda$ 简记 $q_{\lambda,\allowbreak \lambda(1)}$，则对一般位置的路径 $\lambda\in\mathcal P^-\operatorname{Gr}(V)$ 有
\begin{equation}\label{eq:11.20}
\begin{aligned}
I_H(\lambda)&=\sum_{s<1}\operatorname{sign}(q_\lambda(s)),\\
(\delta_H)_\lambda&=\bigotimes_{s<1}(\lambda(s)\cap\lambda(1))^{\operatorname{sign}(q_\lambda(s))}.
\end{aligned}
\end{equation}
在指标理论方面，仍选择前述函数 $\psi$，考虑 $E=H\times V$ 上的 $\bar\partial$ 算子，边界条件为 $F_\lambda\subset\partial H\times V=\R\times V$，$(F_\lambda)_s=\lambda(\psi(s))$。由于 $\lambda$ 的端点横截，它是非退化算子。将其记为 $D_{H,\allowbreak \lambda}$，则有：
\begin{lemma}\label{st:11.11}
$I_H(\lambda)=\operatorname{index}(D_{H,\allowbreak \lambda})$；且 $(\delta_H)_\lambda$ 与行列式线 $\det(D_{H,\allowbreak \lambda})$ 典范同构。
\end{lemma}
\begin{proof}
先看一个简单例子，即形如
\begin{equation}\label{eq:11.21}
\lambda(s)=\exp(\alpha sI_V)\lambda(0)
\end{equation}
的路径，其中 $\lambda(0)$ 为 Lagrangian 子空间，$\alpha\in(-\pi;0)$。此时没有计入\eref{11.20}的交点，所以引理断言 $D_{H,\allowbreak \lambda}$ 的指标为零，并有选定同构 $\det(D_{H,\allowbreak \lambda})\cong\R$。首先，指标与 $\alpha$ 的选择无关，故可假设 $\alpha=-\pi/2$。在此情形，可将两个 $D_{H,\allowbreak \lambda}$ 粘合成圆盘上的算子 $D_{D,\allowbreak \rho}$，其相应 Lagrangian 子空间环路满足 $\mu(\rho)=-n$；于是\eref{11.9}与\stref{引理}{11.7}给出 $2\operatorname{index}(D_{H,\allowbreak \lambda})=\operatorname{index}(D_{D,\allowbreak \rho})=0$。其次，要说明 $D_{H,\allowbreak \lambda}$ 为单射。取直和分量后，可归结到最小秩 $n=1$，此时适用\stref{引理}{11.5}；用该引理的记号，我们的路径 $\lambda$ 可补成满足 $\mu(\rho)=0$ 的环路 $\rho$，而曲面 $H$ 满足 $|\Sigma^-|=1$。显然，这一情形特别简单，也很容易找到其他同样可行的方法。

对一般情形，$\mathcal P^-\operatorname{Gr}(V)$ 中任意路径都可在不改变 $I_H$ 或 $\delta_H$ 的情况下，变形成对某个 $s'<1$ 在 $[s';1]$ 上满足\eref{11.21}的路径。因此只需考虑已有这一性质的路径。将 $(\lambda,\allowbreak \lambda(1))$ 限制在 $[0;s']$ 上，得到 $\mathcal P^2\operatorname{Gr}(V)$ 的元素，因而得到 $Z$ 上的 $\bar\partial$ 算子；其指标与行列式线由\stref{引理}{11.9}给出，因此与\eref{11.20}一致。另一方面，对应于 $\lambda|[s';1]$ 的 $H$ 上算子属于刚才讨论的情形。将 $H$ 的端与 $Z$ 的出向端连接，便可粘合这两部分，所得算子在相差无关紧要的变形的意义下就是 $D_{H,\allowbreak \lambda}$。由\eref{11.9}即得结论。
\end{proof}

\subsection{抽象膜结构}\label{sec:11h}
\SourcePage{171}
这一对类 $(\mu,\allowbreak w_2)$ 由一个到 Eilenberg–MacLane 空间乘积的映射表示：
\begin{equation}\label{eq:11.22}
\operatorname{Gr}(V)\longrightarrow K(\Z,1)\times K(\Z/2,2).
\end{equation}
拉回两个万有纤维化的乘积，得到纤维化 $\operatorname{Gr}(V)^\#\to\operatorname{Gr}(V)$，其纤维为 $\Z\times\RP^\infty$。当 $n\geq3$ 时，$\operatorname{Gr}(V)^\#$ 可与 Postnikov 分解意义下 $\operatorname{Gr}(V)$ 的 $2$ 连通“覆盖”识别。\znote{163}{抽象膜的思想是把 Lagrangian 子空间连同其两条“障碍类”一起提升。$K(\Z,1)$ 分类整系数的一维上同调类，在这里就是 Maslov 类 $\mu$；$K(\Z/2,2)$ 分类二阶的 $\Z/2$ 类，在这里是第二 Stiefel–Whitney 类 $w_2$。提升存在当且仅当这两条障碍同时消失，其纤维因而是 $\Z\times\RP^\infty$：$\Z$ 记录 Maslov 指标的选择，$\RP^\infty=K(\Z/2,1)$ 记录 $w_2$ 的抵消方式。$n\geq3$ 时 $\operatorname{Gr}(V)^\#$ 是 $2$ 连通的，故提升在低维上唯一。nLab 分别给出了这两类障碍的上同调解释。[来源：nLab "Maslov index"，https:/\allowbreak /\allowbreak ncatlab.org/\allowbreak nlab/\allowbreak show/\allowbreak Maslov+\allowbreak index（2026-11-15 访问）；nLab "twisted spin structure"，https:/\allowbreak /\allowbreak ncatlab.org/\allowbreak nlab/\allowbreak show/\allowbreak twisted+\allowbreak spin+\allowbreak structure（2026-11-15 访问）]}当 $n\leq2$ 时，它仍连通，并分别满足
\begin{equation}\label{eq:11.23}
\begin{cases}
\pi_1(\operatorname{Gr}(V)^\#)\cong\Z/2,\quad\pi_2(\operatorname{Gr}(V)^\#)=0,&n=1,\\
\pi_1(\operatorname{Gr}(V)^\#)=0,\quad\pi_2(\operatorname{Gr}(V)^\#)\cong\Z,&n=2.
\end{cases}
\end{equation}
\begin{definition*}
称 $\operatorname{Gr}(V)^\#$ 的点为\emph{抽象线性 Lagrangian 膜}，记为 $\Lambda^\#$，其底层 Lagrangian 子空间记为 $\Lambda$。
\end{definition*}
\begin{lemma}\label{st:11.12}
取环路 $\rho^\#\in\mathcal L\operatorname{Gr}(V)^\#$，令 $D_{D,\allowbreak \rho}$ 为圆盘 $D$ 上对应于底层环路 $\rho$ 的 $\bar\partial$ 算子。则 $\operatorname{index}(D_{D,\allowbreak \rho})=n$；此外，$\rho^\#$ 确定一个同构
\[
\det(D_{D,\rho})\cong\lambda^{\mathrm{top}}(\rho(0)).
\]
\end{lemma}
\begin{proof}
能够提升到 $\operatorname{Gr}(V)^\#$ 的环路必满足 $\mu(\rho)=0$，所以指标结论是\stref{引理}{11.7}的特例。映射 $\mathcal L\operatorname{Gr}(V)^\#\to\mathcal L_0\operatorname{Gr}(V)$ 的纤维同伦等价于 $\Z\times\mathcal L\RP^\infty$。忘掉第一个因子，并只考虑第二个因子的 $\pi_0$，便得到一个诱导的二重覆盖，其上可关联实线丛 $\delta_D\to\mathcal L_0\operatorname{Gr}(V)$。由构造，它由 $T(w_2)\in H^1(\mathcal L_0\operatorname{Gr}(V);\Z/2)$ 分类，因此由\stref{引理}{11.7}，它同构于纤维为
\begin{equation}\label{eq:11.24}
\det(D_{D,\rho})^\vee\otimes\lambda^{\mathrm{top}}(\rho(0))
\end{equation}
的自然线丛。为固定两者间的一个同构，只需考察 $\mathcal L_0\operatorname{Gr}(V)$ 的任意一个连通子集，例如常值环路所成的子集。回顾\stref{引理}{11.7}的证明：当 $\rho$ 为常值时，$D_{D,\allowbreak \rho}$ 为满射，核同构于 $\rho(0)$，故 $\det(D_{D,\allowbreak \rho})^\vee\otimes\lambda^{\mathrm{top}}(\rho(0))$ 有自然选定的定向。另一方面，$\mathcal L\operatorname{Gr}(V)^\#$ 在 $\rho$ 上的纤维可与 $\Z\times\mathcal L\RP^\infty$ 识别，该识别在相差同伦及 $\Z$ 因子中的平移的意义下唯一；所以选择 $\pi_0(\mathcal L\RP^\infty)\cong\pi_1(\RP^\infty)$ 的平凡元素，便赋予 $(\delta_D)_\rho$ 一个定向。结合两者，得到所需的线丛同构。最后，依定义，$\delta_D$ 拉回到 $\mathcal L\operatorname{Gr}(V)^\#$ 后典范平凡，从而给出\eref{11.24}拉回后的平凡化。
\end{proof}
\SourcePage{172}
对于底层子空间 $(\Lambda_0,\allowbreak \Lambda_1)$ 横截相交的任意一对抽象线性 Lagrangian 膜 $(\Lambda_0^\#,\allowbreak \Lambda_1^\#)$，可以如下赋予一个整数——\emph{绝对指标} $i(\Lambda_0^\#,\allowbreak \Lambda_1^\#)$——以及一个一维实向量空间——\emph{定向空间} $o(\Lambda_0^\#,\allowbreak \Lambda_1^\#)$。选择 $\operatorname{Gr}(V)^\#$ 中连接两点的路径 $\lambda^\#$，使其投影 $\lambda$ 属于 $\mathcal P^-\operatorname{Gr}(V)$，并令
\begin{equation}\label{eq:11.25}
i(\Lambda_0^\#,\Lambda_1^\#)=I_H(\lambda),\qquad o(\Lambda_0^\#,\Lambda_1^\#)=(\delta_H)_\lambda.
\end{equation}
当然，必须证明它们与 $\lambda^\#$ 的选择无关；对 $o$ 而言，是在典范同构意义下无关。以 $\Omega^-(\operatorname{Gr}(V),\allowbreak \Lambda_0,\allowbreak \Lambda_1)$ 表示 $\mathcal P^-\operatorname{Gr}(V)$ 中从 $\Lambda_0$ 到 $\Lambda_1$ 的路径所成的子空间，以 $\Omega^-(\operatorname{Gr}(V)^\#,\allowbreak \Lambda_0^\#,\allowbreak \Lambda_1^\#)$ 表示对应的 $\operatorname{Gr}(V)^\#$ 中的路径空间。它们分别同伦等价于通常的环路空间 $\Omega\operatorname{Gr}(V)$ 与 $\Omega\operatorname{Gr}(V)^\#$。当 $n\geq3$ 时，已知 $\Omega\operatorname{Gr}(V)^\#$ 单连通，立即得到所需唯一性。当 $n=2$ 时，$w_2$ 在 $\pi_2(\operatorname{Gr}(V))$ 上仍非平凡，这意味着映射 $\Z=\pi_2(\operatorname{Gr}(V)^\#)\to\pi_2(\operatorname{Gr}(V))=\Z$ 为乘以 $2$。因此，线丛 $\delta_H$ 拉回到 $\Omega^-(\operatorname{Gr}(V)^\#,\allowbreak \Lambda_0^\#,\allowbreak \Lambda_1^\#)$ 后平凡，同样得到结论。最后，当 $n=1$ 时，由\eref{11.23}有两个不同的同伦类；二者给出相同的 $i$，而对于定向空间，我们直接规定：改变 $\lambda^\#$ 的同伦类，会给出 $(\delta_H)_\lambda$ 的一个反定向自同构。一种等价而或许更系统的做法，是通过稳定化处理这些低维情形。\znote{164}{定向空间 $o(\Lambda_0^\#,\Lambda_1^\#)$ 是一个一维实向量空间，按\stref{引理}{11.11}它典范同构于路径上某个 $\bar\partial$ 算子的行列式线。给一维实空间选定一个生成元就是给它定向，所以“定向空间”这个名字恰是行列式线加上定向数据的合称：一旦为 $o$ 选定生成元（等价地，为相应的行列式线选定定向），两个膜之间的定向数据就确定了。Floer 理论中的定向问题正是把诸如此类的线丛的定向相容地选出来；nLab 关于行列式线丛的条目讨论了这类线的整体性质。[来源：nLab "determinant line bundle"，https:/\allowbreak /\allowbreak ncatlab.org/\allowbreak nlab/\allowbreak show/\allowbreak determinant+\allowbreak line+\allowbreak bundle（2026-11-15 访问）]}

设 $S$ 为带条带状端的 Riemann 曲面，配有平凡辛向量丛 $E=S\times V$，以及由映射 $\nu:\partial S\to\operatorname{Gr}(V)$ 给出的 Lagrangian 子丛 $F_\nu\subset E|\partial S$；$\nu$ 在条带状端上局部常值。后一条件意味着，对每个无穷远点 $\zeta$，有一对 Lagrangian 子空间 $(\Lambda_{\zeta,\allowbreak 0},\allowbreak \Lambda_{\zeta,\allowbreak 1})$，使 $\nu(\epsilon_\zeta(s,\allowbreak k))=\Lambda_{\zeta,\allowbreak k}$，其中 $(s,\allowbreak k)\in\R^\pm\times\{0;1\}$。假设每一对均满足 $\Lambda_{\zeta,\allowbreak 0}\pitchfork\Lambda_{\zeta,\allowbreak 1}$，从而相应的 $\bar\partial$ 算子 $D_{S,\allowbreak \nu}$ 非退化。此外，假设给定 $\nu$ 到 $\operatorname{Gr}(V)^\#$ 的提升 $\nu^\#$，它在各端上也局部常值，因而为这些 Lagrangian 子空间提供选定原像 $\Lambda_{\zeta,\allowbreak k}^\#$。
\begin{proposition}\label{st:11.13}
以 $\hat C_1,\allowbreak \ldots,\allowbreak \hat C_k$ 表示紧化 $\hat S$ 的边界分支，令 $e_j=|\Sigma^-\cap\hat C_j|$ 为每个分支上的入向无穷远点数。对每个 $j$ 固定一个非无穷远点 $z_j\in\hat C_j$。则
\begin{equation}\label{eq:11.26}
\begin{aligned}
\operatorname{index}(D_{S,\nu})&=n(\chi(S)-|\Sigma^-|)+\sum_{\zeta\in\Sigma^\pm}\mp i(\Lambda_{\zeta,0}^\#,\Lambda_{\zeta,1}^\#),\\
\det(D_{S,\nu})&\cong\bigotimes_j\left(\lambda^{\mathrm{top}}(\nu_{z_j})^{\otimes(1-e_j)}
\otimes\bigotimes_{\zeta\in\Sigma^\pm\cap\hat C_j}o(\Lambda_{\zeta,0}^\#,\Lambda_{\zeta,1}^\#)^{\mp1}\right).
\end{aligned}
\end{equation}
\end{proposition}
\SourcePage{173}
\begin{proof}
已证明两个特殊情形。$S=D$ 时就是\stref{引理}{11.12}。另一个情形 $S=H$ 中，先对 $\nu:\partial H=\R\to\operatorname{Gr}(V)$ 作一个紧支的小变形，使之成为 $\nu(s)=\lambda(\psi(s))$，其中 $\lambda\in\Omega^-(\operatorname{Gr}(V),\allowbreak \Lambda_{\zeta,\allowbreak 0},\allowbreak \Lambda_{\zeta,\allowbreak 1})$，并对提升 $\nu^\#,\allowbreak \lambda^\#$ 作相应处理。这诱导 Fredholm 算子的变形，因此不改变\eref{11.26}左侧；也显然保持右侧，因为右侧与路径的任何选择无关。然后运用\stref{引理}{11.11}及 $i,\allowbreak o$ 的定义。还需要一个相关观察。取带出向无穷远点的上半平面，分别记为 $\bar H$ 与 $\bar\zeta$。赋予它反向路径 $\bar\nu^\#(s)=\nu^\#(-s)$，相应的 Lagrangian 子空间对仍与之前相同，$\Lambda_{\bar\zeta,\allowbreak k}=\Lambda_{\zeta,\allowbreak k}$。于是可以粘合 $\bar H\#H\cong D$，且与全部附加结构相容。应用\eref{11.9}得到
\begin{equation}\label{eq:11.27}
\begin{aligned}
n&=\operatorname{index}(D_{\bar H,\bar\nu})+i(\Lambda_{\bar\zeta,0}^\#,\Lambda_{\bar\zeta,1}^\#),\\
\lambda^{\mathrm{top}}(\bar\nu_z)&\cong\det(D_{\bar H,\bar\nu})\otimes o(\Lambda_{\bar\zeta,0}^\#,\Lambda_{\bar\zeta,1}^\#),
\end{aligned}
\end{equation}
其中 $z$ 为 $\partial\bar H$ 的任意点。这证明\eref{11.26}对 $\bar H$ 也成立。

一般策略如下。可以将紧化 $\hat S$ 视为 $k$ 个圆盘 $\hat D_1,\allowbreak \ldots,\allowbreak \hat D_k$ 与同亏格 $g$ 的闭曲面的连通和。以 $D_1,\allowbreak \ldots,\allowbreak D_k$ 表示这些圆盘的有限部分，即带边界标记点的部分。如\stref{引理}{11.7}证明的步骤 (i)，沿平行于 $\partial\hat D_j$ 的曲线捏合分离，得到
\[
\begin{aligned}
\operatorname{index}(D_{S,\nu})&=n(2-2g)+\sum_j(\operatorname{index}(D_{D_j,\nu_j})-2n),\\
\det(D_{S,\nu})&\cong\bigotimes_j\det(D_{D_j,\nu_j}),
\end{aligned}
\]
其中 $\nu_j$ 是 $\nu$ 在属于 $\hat C_j$ 的各分支上的限制。这就把问题归结为 $\hat S$ 为圆盘的情形，以下均作此假设。对每个出向无穷远点 $\zeta^+\in\Sigma^+$，选择 $\Omega^-(\operatorname{Gr}(V)^\#,\allowbreak \Lambda_{\zeta^+,\allowbreak 0}^\#,\allowbreak \Lambda_{\zeta^+,\allowbreak 1}^\#)$ 中的一条路径。取带相应边界数据的 $H$ 的副本，粘合到 $S$ 上，填补各出向端。若只有一个入向无穷远点，粘合后仍得到 $H$ 上的算子；由已证明的情形可见
\[
\begin{aligned}
i(\Lambda_{\zeta^-,0}^\#,\Lambda_{\zeta^-,1}^\#)&=\operatorname{index}(D_{S,\nu})+\sum_{\zeta^+}i(\Lambda_{\zeta^+,0}^\#,\Lambda_{\zeta^+,1}^\#),\\
o(\Lambda_{\zeta^-,0}^\#,\Lambda_{\zeta^-,1}^\#)&\cong\det(D_{S,\nu})\otimes\bigotimes_{\zeta^+}o(\Lambda_{\zeta^+,0}^\#,\Lambda_{\zeta^+,1}^\#).
\end{aligned}
\]
若没有入向无穷远点，或入向无穷远点至少有两个，则采取类似但稍复杂的办法，额外粘合 $\bar H$ 的副本，将
\SourcePage{174}
$S$ 完全紧化，再利用\eref{11.27}得出
\begin{equation}\label{eq:11.28}
\begin{aligned}
n={}&\sum_{\zeta^-}(n-i(\Lambda_{\zeta^-,0}^\#,\Lambda_{\zeta^-,1}^\#))
+\operatorname{index}(D_{S,\nu})+\sum_{\zeta^+}i(\Lambda_{\zeta^+,0}^\#,\Lambda_{\zeta^+,1}^\#),\\
\lambda^{\mathrm{top}}(\nu_{z_0})\cong{}&\bigotimes_{\zeta^-}\left(\lambda^{\mathrm{top}}(\nu_{z_j})\otimes o(\Lambda_{\zeta^-,0}^\#,\Lambda_{\zeta^-,1}^\#)^{-1}\right)\\
&\otimes\det(D_{S,\nu})\otimes\bigotimes_{\zeta^+}o(\Lambda_{\zeta^+,0}^\#,\Lambda_{\zeta^+,1}^\#).
\end{aligned}
\end{equation}
这里 $z_j$ 是 $S$ 的边界点。沿紧化曲面上的 Lagrangian 子空间环路作平行移动，可将各 $\lambda^{\mathrm{top}}(\nu_{z_j})$ 相互识别。由于该环路容许膜结构，它的 Maslov 指标为零，因而可定向，所以这里没有歧义。\znote{165}{这里用到了膜结构的全体数据：$\Lambda^\#$ 自带 Maslov 指标与 $\Z/2$ 障碍，故环路一旦提升到 $\operatorname{Gr}(V)^\#$ 中，它的 Maslov 指标与第二 Stiefel–Whitney 类便同时为零。前者使 $\lambda^{\mathrm{top}}$ 沿环路的平行移动没有单值性，后者使定向可以一致地搬运；两者合起来，才允许把各个 $\lambda^{\mathrm{top}}(\nu_{z_j})$ 相互等同。换言之，$\Z$ 的相位障碍与 $\Z/2$ 的符号障碍在此抵消。nLab 关于扭曲 spin 结构的条目说明了 $w_2$ 与 twist 如何合成新的障碍类。[来源：nLab "twisted spin structure"，https:/\allowbreak /\allowbreak ncatlab.org/\allowbreak nlab/\allowbreak show/\allowbreak twisted+\allowbreak spin+\allowbreak structure（2026-11-15 访问）]}这样就将\eref{11.28}化为与\eref{11.26}一致的形式。
\end{proof}

目前的表述有一个严重缺点：$\operatorname{Gr}(V)^\#$ 只在同伦等价意义下定义，使得抽象线性 Lagrangian 膜的概念带有一定的含混性。因此，下面用一个等价但更具几何意味的框架来代替，即分次与 Pin 结构。

\subsection{Pin 结构}\label{sec:11i}
$\operatorname{Pin}_n$ 是 $O_n$ 关于 $\Z/2$ 的中心扩张。这样的扩张实际有两个，它们拓扑相同，却是不同的 Lie 群。对我们的目的，其中一种稍方便，因此规定 $\operatorname{Pin}_n$ 为由 $w_2\in H^2(BO_n;\Z/2)$ 在拓扑群扩张意义下分类的群。\znote{167}{两种约定源于对二次型取 $q_\pm=\pm\sum_i(x^i)^2$：$\operatorname{Pin}_+(n)$ 满足 $v^2=+\|v\|^2$，$\operatorname{Pin}_-(n)$ 满足 $v^2=-\|v\|^2$，两者都是 $O_n$ 的二重覆盖，因而是同胚的拓扑空间，却有不同的群结构。正文取 $v^2=\|v\|^2e$，故这里的 $\operatorname{Pin}_n$ 即 $\operatorname{Pin}_{0,n}$，也记作 $\operatorname{Pin}_n^+$。这一约定还决定了反射的提升是二阶还是四阶元素：$\operatorname{Pin}_+$ 中反射的原像满足 $a^2=e$，而另一版本 $\operatorname{Pin}_{n,0}=\operatorname{Pin}_n^-$ 中它满足 $a^4=e$，且其分类类为 $w_2+w_1^2$，正如下一则笔记所述。两个版本给出的 Pin 结构障碍因此相差 $w_1^2$，这正是选用哪一版本会改变障碍计算的原因。[来源：nLab "pin group"，https:/\allowbreak /\allowbreak ncatlab.org/\allowbreak nlab/\allowbreak show/\allowbreak Pin+\allowbreak group（2026-11-15 访问）]}简要回顾其具体定义。将 $\R^n$ 嵌入 Clifford 代数 $\operatorname{Cl}(\R^n)$，其中向量满足 $v^2=\|v\|^2e$，$e\in\operatorname{Cl}(\R^n)$ 为单位元。$\operatorname{Pin}_n\subset\operatorname{Cl}(\R^n)^\times$ 是由 $S^{n-1}\subset\R^n$ 生成的乘法子群。\znote{167}{nLab 的 pin 群条目给出的定义与此一致：$\operatorname{Pin}(V,q)$ 是 Clifford 代数乘法群 $\operatorname{GL}_1(\operatorname{Cl}(V,q))$ 的子群，由满足 $q(v_i)=1$ 的向量之积 $v_1\cdots v_n$ 生成；其中偶数个因子之积再构成 $\operatorname{Spin}(V,q)$。作为抽象群，它由这些单位向量生成，关系来自 Clifford 关系 $v^2=q(v)$；投影 $\operatorname{Pin}(V,q)\to O(V)$ 由扭曲伴随作用给出，核恰为 $\{\pm1\}$。正文中的 $v$ 满足 $v^2=\|v\|^2e$，即 $q(v)=1$，故正是这里的生成元。这些关系在 $n$ 固定后是有限的，故 $\operatorname{Pin}(V,q)$ 是有限维 Lie 群，维数与 $O(V)$ 相同。标准参考是 Lawson–Michelsohn《Spin geometry》第 I 章第 2 节，nLab 该条目也列出了此书。[来源：nLab "pin group"，https:/\allowbreak /\allowbreak ncatlab.org/\allowbreak nlab/\allowbreak show/\allowbreak Pin+\allowbreak group（2026-11-15 访问）]}同态 $\operatorname{Pin}_n\to O_n$ 是扭曲伴随作用在 $\R^n\subset\operatorname{Cl}(\R^n)$ 上的限制：元素 $v\in S^{n-1}\subset\operatorname{Pin}_n$ 按 $w\mapsto-vwv$ 作用，即关于垂直于 $v$ 的超平面反射。当然，可以用任意内积空间 $F$ 代替 $\R^n$，定义具有同样性质的群 $\operatorname{Pin}(F)$。考察反射的复合可得：
\begin{lemma}\label{st:11.14}
若 $A\in O(F)$ 为对合，则 $A$ 在 $\operatorname{Pin}(F)$ 中的原像与子空间 $F^{\mathrm{anti}}=\ker(\Id+A)$ 的定向典范对应。任一这样的原像的平方为
\begin{equation}\label{eq:11.29}
(-1)^{\frac12\operatorname{rank}(F^{\mathrm{anti}})(\operatorname{rank}(F^{\mathrm{anti}})-1)}e.
\end{equation}
特别地，反射的原像是二阶元素。
\end{lemma}
\begin{note}\label{st:11.15}
Pin 的一般参考文献有\bcite{93}，其中我们的 $\operatorname{Pin}_n$ 被称为 $\operatorname{Pin}_{0,\allowbreak n}$；以及\bcite{86}，其中称它为 $\operatorname{Pin}_n^+$。另一个版本记为 $\operatorname{Pin}_{n,\allowbreak 0}$ 或 $\operatorname{Pin}_n^-$，具有以下性质：它关联于类 $w_2+w_1^2$；它位于
\SourcePage{175}
满足 $v^2=-\|v\|^2e$ 的、更为熟知的 Clifford 代数内；最后，在该群中反射的原像是四阶元素。
\end{note}
\begin{definition*}
如通常一样，实内积向量丛 $F^n\to B$ 上的 \emph{Pin 结构}，是一个 $\operatorname{Pin}_n$ 主丛 $P^\#\to B$，连同同构 $P^\#\times_{\operatorname{Pin}_n}\R^n\cong F$。
\end{definition*}
同态 $\operatorname{Pin}_n\to O_n$ 诱导二重覆盖 $P^\#\to P$，其中 $P$ 为 $F$ 的标准正交标架丛。由最初对 $\operatorname{Pin}_n$ 的刻画，这种结构存在的障碍是 $w_2(F)$。若该类消失，则 Pin 结构的同构类所成的集合是 $H^1(B;\Z/2)$ 上的仿射空间。\znote{168}{Pin 结构的标准图景是正交标架丛分类映射的提升：标架丛给出映射 $B\to BO_n$，Pin 结构即它穿过 $\operatorname{BPin}_n\to BO_n$ 的提升，二重覆盖的障碍正是 $w_2$。若两个提升存在，它们的差别给出一个到 $K(\Z/2,1)$ 的映射，即 $H^1(B;\Z/2)$ 中的一个类；反过来，任给这样的类都可用它扭出新的提升，所以全部 Pin 结构构成 $H^1(B;\Z/2)$ 的仿射空间。nLab 关于扭曲 spin 结构的条目以同样方式刻画了 spin 结构为 Cech 圈穿过 $\mathrm{BSpin}\to BSO$ 的提升。[来源：nLab "twisted spin structure"，https:/\allowbreak /\allowbreak ncatlab.org/\allowbreak nlab/\allowbreak show/\allowbreak twisted+\allowbreak spin+\allowbreak structure（2026-11-15 访问）]}具体地，给定 Pin 结构 $P^\#$ 和实线丛 $\beta\to B$，可以取纤维积 $P^\#\times_{\Z/2}S(\beta)$ 形成新的 Pin 结构；这里 $\Z/2$ 以 $\pm1$ 作用于与 $\beta$ 关联的二重覆盖 $S(\beta)$，并通过中心子群 $\ker(\operatorname{Pin}_n\to O_n)=\{\pm e\}$ 作用于 $P^\#$。为简便，将新结构记为 $P^\#\otimes\beta$。

任意向量丛 $F$ 都可关联一个群丛 $\operatorname{Pin}(F)$。下面澄清一个可能混淆的地方。给定 $F$ 上的 Pin 结构 $P^\#$，考虑群丛 $\operatorname{Aut}(P^\#)\to B$，其在 $x$ 上的纤维由 $\operatorname{Pin}_n$ 主齐性空间 $P_x^\#$ 的全部自同构组成。可将其视为伴随丛
\[
\operatorname{Aut}(P^\#)=P^\#\times_{\operatorname{Pin}_n}\operatorname{Pin}_n
=P\times_{O_n}\operatorname{Pin}_n\hookrightarrow P\times_{O_n}\operatorname{Cl}(\R^n).
\]
这里 $\operatorname{Pin}_n$ 通过共轭作用于自身；此作用下降为 $O_n$ 作用，并嵌入 $\operatorname{Cl}(\R^n)$ 上的非扭曲伴随作用。例如，关于超平面 $v^\perp$ 的反射按 $w\mapsto vwv$ 作用，其在 $\R^n\subset\operatorname{Cl}(\R^n)$ 上的限制等于该反射乘以 $-1$。所以 $\operatorname{Aut}(P^\#)$ 同构的并非 $\operatorname{Pin}(F)$，而是 $\operatorname{Pin}(F\otimes\lambda^{\mathrm{top}}(F))$，并构成如下交换图：
\SourceEquationDiagram{diagram-11-30}{0.51507637}{11.30}
更一般地，将 Pin 结构 $P^\#$ 分别与两个线丛 $\beta_0,\allowbreak \beta_1$ 作张量积，所得主丛之间的 $\operatorname{Pin}_n$ 等变映射所成的丛为
\begin{equation}\label{eq:11.31}
\operatorname{Iso}(P^\#\otimes\beta_0,P^\#\otimes\beta_1)
\cong\operatorname{Pin}(F\otimes\lambda^{\mathrm{top}}(F))\times_{\Z/2}S(\beta_0^\vee\otimes\beta_1).
\end{equation}
\begin{note}\label{st:11.16}
还有关于参考类 $w\in H^2(B;\Z/2)$ 扭曲的 Pin 结构的一般理论，其存在障碍为 $w_2(F)+w$。一个 gerbe 式定义如下：选择分别表示 $w$ 与丛 $F$ 的 Čech 上链 $\check w\in\check C^2(B;\Z/2)$ 与 $\check F\in\check C^1(B;O_n)$。扭曲 Pin 结构
\SourcePage{176}
是 $\check C^1(B;\operatorname{Pin}_n)$ 中投影到 $\check F$ 的元素，而它未能成为上闭链的偏差由 $\check w$ 编码；因为 $\Z/2\subset\operatorname{Pin}_n$ 是中心子群，所以这一定义有意义。\znote{169}{扭曲 spin 结构的一般定义是：预先指定障碍类 $\alpha\in H^2(B;\Z/2)$，要求标架丛的 Čech 圈 $g:B\to BO_n$ 满足 $[w_2(g)]=\alpha$；所有这样的圈连同其间的同伦构成一个 groupoid，它由 $(\infty,1)$-拉回刻画。障碍本身是一个 $\mathrm{B}\Z/2$ 主 $2$-丛，即束 gerbe（Murray–Singer 称之为 spin gerbe）。正文这里的 gerbe 式定义正是该说法的 Čech 版本，其中 $\check w$ 就是这个 $2$-丛的过渡数据。具体地说，$\check w$ 记录的是沿三重交叠把局部提升粘起来时出现的符号偏差：它不必是上闭链，其 coboundary 正是障碍类 $\alpha$，这也解释了为何在扭曲的图景里只能要求提升“相差 $\check w$”而不是严格相容。[来源：nLab "twisted spin structure"，https:/\allowbreak /\allowbreak ncatlab.org/\allowbreak nlab/\allowbreak show/\allowbreak twisted+\allowbreak spin+\allowbreak structure（2026-11-15 访问）]}

另一种描述较不一般，但借助先前讨论的材料也许更易理解。假设有光滑超曲面 $\Delta\subset B$，其上有实线丛 $\beta$，而 $w$ 由 $w_1(\beta)\in H^1(\Delta;\Z/2)$ 经推前得到。为简化叙述，再假设 $\Delta$ 有余定向。沿 $\Delta$ 切开 $B$，得到有两个边界部分 $\partial_+\widetilde B,\allowbreak \partial_-\widetilde B\cong\Delta$ 的流形 $\widetilde B$。此时，$F$ 上的扭曲 Pin 结构，就是拉回丛 $\widetilde F\to\widetilde B$ 上的 Pin 结构 $\widetilde P^\#$，连同同构 $\widetilde P^\#|\partial_+\widetilde B\cong\widetilde P^\#|\partial_-\widetilde B\otimes\beta$，它覆盖显然的识别 $\widetilde F|\partial_-\widetilde B=F|\Delta=\widetilde F|\partial_+\widetilde B$。
\end{note}

Pin 与 $\bar\partial$ 算子指标理论的关联，源自它和 $w_2$ 的关系。为看一个基本例子，暂时回到引理 \sr{st:11.7}{11.7}、\sr{st:11.12}{11.12} 的情形。考虑环路 $\rho\in\mathcal L_k\operatorname{Gr}(V)$、相应向量丛 $F_\rho\to S^1$，以及圆盘上的算子 $D_{D,\allowbreak \rho}$。令 $\widetilde F_\rho$ 为 $F_\rho$ 沿端点识别映射 $[0;1]\to S^1$ 的拉回。希望赋予它
\begin{equation}\label{eq:11.32}
\text{Pin 结构 }\widetilde P_\rho^\#,\quad\text{以及同构 }\widetilde P_1^\#\cong\widetilde P_0^\#\otimes\lambda^{\mathrm{top}}(\rho(1))^{\otimes k}.
\end{equation}
当 $k$ 为偶数时，这等价于 $F_\rho$ 上的一个 Pin 结构。当 $k$ 为奇数时，它是\stref{笔记}{11.16}中扭曲 Pin 结构定义的特例，其中 $F=F_\rho$，超曲面为点 $\Delta=\{0\}\subset S^1$，$\beta\to\{0\}$ 就是实直线 $\lambda^{\mathrm{top}}(\rho(1))$。这显然可推广到环路族，此时\eref{11.32}的存在障碍来自类
\begin{equation}\label{eq:11.33}
T(w_2)+kU(\mu)\in H^1(\mathcal L_k\operatorname{Gr}(V);\Z/2).
\end{equation}
\begin{lemma}\label{st:11.17}
$F_\rho$ 上一个形如\eref{11.32}的扭曲 Pin 结构的选择，确定同构 $\det(D_{D,\allowbreak \rho})\cong\lambda^{\mathrm{top}}(\rho(0))$。
\end{lemma}
\begin{proof}
对每个 $\rho$，扭曲 Pin 结构有两个可能的同构类；它们构成 $\mathcal L_k\operatorname{Gr}(V)$ 的一个二重覆盖，由\eref{11.33}分类。由\stref{引理}{11.7}，它同构于纤维为 $\det(D_{D,\allowbreak \rho})^\vee\otimes\lambda^{\mathrm{top}}(\rho(0))$ 的线丛所对应的覆盖。因此，要固定两个丛间的同构，只需考察它们在 $\mathcal L_k\operatorname{Gr}(V)$ 的某个非空连通子集上的限制。

这与\stref{引理}{11.12}的策略十分相似，事实上 $k=0$ 的情形只是该证明的改写。考虑常值环路 $\rho$，此时 $F_\rho$ 平凡，因而其上有选定的平凡 Pin 结构。规则是：若 Pin 结构平凡，就选择由 $\ker(D_{D,\allowbreak \rho})\cong\rho(0)$ 诱导的同构 $\det(D_{D,\allowbreak \rho})\cong
\SourcePage{177}
\lambda^{\mathrm{top}}(\rho(0))$；若 Pin 结构非平凡，就反转其符号。

再假设 $k=-1$。取正交复分解 $V=V_+\oplus V_-$，其中 $V_-$ 的复维数为一，考虑形如 $\rho(s)=\Lambda_+\oplus\exp(-\pi sI_V)\Lambda_-$ 的环路，其中 $\Lambda=\Lambda_+\oplus\Lambda_-$ 为相应分解的 Lagrangian 子空间。$F_\rho$ 上有一个显然的联络，其单值变换为 $A=\Id_{\Lambda_+}\oplus(-\Id_{\Lambda_-})$。一个扭曲 Pin 结构\eref{11.32}的选择，等价于将 $A$ 提升为 Pin 结构同构 $\widetilde P_0^\#\otimes\lambda^{\mathrm{top}}(\rho(0))\to\widetilde P_1^\#$。运用\eref{11.31}可知，它对应于 $A\otimes\Id$ 在 $\operatorname{Pin}(\Lambda\otimes\lambda^{\mathrm{top}}(\Lambda))\times_{\Z/2}S(\lambda^{\mathrm{top}}(\Lambda))$ 中的一个原像。由\stref{引理}{11.14}，两处 $\lambda^{\mathrm{top}}(\Lambda)$ 抵消，所以这归结为 $\Lambda_-$ 的一个定向。另一方面，将算子分解为 $V_\pm$ 两部分，并运用\stref{引理}{11.7}证明的 (iii)，可见 $D_{D,\allowbreak \rho}$ 为满射，其核同构地映到 $\Lambda_+$。因此，$\Lambda_-$ 的定向典范确定同构 $\det(D_{D,\allowbreak \rho})\cong\lambda^{\mathrm{top}}(\Lambda)$，正如所需。它对分解 $V=V_+\oplus V_-$ 和 $\Lambda=\Lambda_+\oplus\Lambda_-$ 的选择连续。

对一般的 $k$，考虑环路 $\rho(s)=\Lambda_+\oplus\exp(k\pi sI_V)\Lambda_-$。论证中 Pin 结构的部分只依赖 $k$ 的奇偶性，因此与之前相同。在行列式线丛方面，可用粘合公式\eref{11.11}得到典范同构 $\det(D_{D,\allowbreak \rho})\cong\lambda^{\mathrm{top}}(\Lambda_+)\otimes(\Lambda_-)^{\otimes(1+k)}$。
\end{proof}

\subsection{重新考察膜}\label{sec:11j}
固定 $(V,\allowbreak I_V)$ 上的\emph{二次复体积形式}，即复直线之间的同构
\[
\eta_V^2:\lambda_{\CC}^{\mathrm{top}}(V)^{\otimes2}\longrightarrow\CC.
\]
等价地，可将它视为在相差符号的意义下确定的复体积形式 $\pm\eta_V:\lambda_{\CC}^{\mathrm{top}}(V)\to\CC$。相关的\emph{平方相位映射}——在特殊 Lagrangian 几何中广为人知——为
\begin{equation}\label{eq:11.34}
\alpha_V:\operatorname{Gr}(V)\to S^1,\qquad
\alpha_V(\Lambda)=\frac{\eta_V(v_1\wedge\cdots\wedge v_n)^2}{|\eta_V(v_1\wedge\cdots\wedge v_n)|^2},
\end{equation}
其中 $v_1,\allowbreak \ldots,\allowbreak v_n$ 是 $\Lambda$ 的任意一组基。该定义良好，因为非零元素 $v_1\wedge\cdots\wedge v_n\in\lambda_{\CC}^{\mathrm{top}}(V)$ 在相差 $\R^\times$ 倍的意义下与基无关，而平方去除了符号歧义。\znote{170}{平方相位映射是特殊 Lagrangian 几何的基本对象。nLab 关于 Lagrangian Grassmann 流形的条目指出，不稳定地 $H^1(\operatorname{LGrass}(\CC^n),\Z)\cong[U(n)/O(n),U(1)]$，由 $\det^2:U(n)\to U(1)$ 给出——因为 $O(n)$ 的行列式为 $\pm1$，$\det^2$ 过商良好。$\eta_V(v_1\wedge\cdots\wedge v_n)^2$ 正是 $\det^2$ 的具体实现，所以\eref{11.34}的 $\alpha_V$ 表示 Maslov 类 $\mu$；这也解释了为何取平方恰好消去基的符号歧义。从分次的角度看，$\alpha_V$ 的相位正是膜结构中实值相位 $\alpha^\#$ 的底：$\exp(2\pi i\alpha^\#)=\alpha_V(\Lambda)$ 把 $\alpha_V$ 拉回到 $\R$ 上成为多值函数，其单值化所需的整值修正由 Maslov 指标控制，所以 $\alpha_V$ 同时含有 $\Z/2$ 的符号信息与 $\Z$ 的 Maslov 信息。[来源：nLab "Lagrangian Grassmannian"，https:/\allowbreak /\allowbreak ncatlab.org/\allowbreak nlab/\allowbreak show/\allowbreak Lagrangian+\allowbreak Grassmannian（2026-11-15 访问）]}
\begin{definition*}
\emph{线性 Lagrangian 膜}是三元组 $(\Lambda,\allowbreak \alpha^\#,\allowbreak P^\#)$，由 Lagrangian 子空间 $\Lambda$、满足 $\exp(2\pi i\alpha^\#)=\alpha_V(\Lambda)$ 的实数 $\alpha^\#$，以及 $\operatorname{Pin}_n$ 主齐性空间 $P^\#$ 连同同构 $P^\#\times_{\operatorname{Pin}_n}\R^n\cong\Lambda$ 组成。称 $\alpha^\#$ 为\emph{实值相位}或\emph{分次}。
\end{definition*}
还需要考虑这种膜的族。具体地，设 $E\to B$ 为满足 $2c_1(E)=0$ 的辛向量丛。该条件保证，一旦固定相容复结构 $I_E$，复线丛 $\lambda_{\CC}^{\mathrm{top}}(E)^{\otimes2}$ 就
\SourcePage{178}
平凡，从而 $E$ 容许逐纤维的二次复体积形式。固定一个这样的形式 $\eta_E^2$，以 $\alpha_E$ 表示 Lagrangian Grassmann 流形丛 $\operatorname{Gr}(E)\to B$ 上相应的平方相位映射。
\begin{definition*}
Lagrangian 子丛 $F\subset E$ 的\emph{膜结构}，由光滑函数 $\alpha^\#:B\to\R$ 和 $F$ 上的 Pin 结构 $P^\#$ 组成，其中 $\exp(2\pi i\alpha^\#(b))=\alpha_E(F_b)$。
\end{definition*}
显然，膜结构存在的障碍是由平方相位 $b\mapsto\alpha_E(F_b)$ 表示的类 $\mu_F\in H^1(B)$，以及 $w_2(F)\in H^2(B;\Z/2)$。若这些障碍消失，则 $F$ 上膜结构的同构类构成 $H^0(B)\times H^1(B;\Z/2)$ 上的仿射空间。对应运算是：给分次加上一个局部常值的 $\Z$ 值函数，以及将 Pin 结构与实线丛作张量积。回到初始数据，应指出 $E$ 的子丛上的膜结构概念依赖于 $\eta_E^2$。更确切地说，只取决于它作为 $\lambda_{\CC}^{\mathrm{top}}(E)^{\otimes2}$ 平凡化的同伦类：变形 $\eta_E^2$ 时，平方相位函数连续变化，并可提升为实值相位的同伦。因此，实际的选择自由度由 $H^1(B)$ 给出。
\begin{note}\label{st:11.18}
令 $\xi_E\subset\lambda_{\CC}^{\mathrm{top}}(E)$ 为使 $\pm\eta_E$ 取实值的一维实子空间。$\xi_E$ 的非平凡性是为 $\eta_E^2$ 找到全局平方根的障碍。对任意 $b\in B$ 及 $F_b$ 的任意基 $\{v_k\}$，有 $v_1\wedge\cdots\wedge v_n\in\alpha_E(F_b)^{1/2}\xi_E$。假定给定膜结构，将它乘以 $\exp(-\pi i\alpha^\#)$，便得到同构 $\lambda^{\mathrm{top}}(F)\cong\xi_E$。因此，分次确定 $F$ 相对于 $\xi_E$ 的一个定向。给分次加上 $\sigma\in\Z$，使该定向改变 $(-1)^\sigma$ 倍。
\end{note}

建立一般框架之后，回到较受限的情形：平凡向量丛 $E=B\times V$，且 $I_E=I_V$、$\eta_E^2=\eta_V^2$ 均为常值。在此背景下，现在的膜概念与先前定义的抽象膜有何关系？容易看出，映射\eref{11.34}表示 Maslov 类 $\mu_V$。此外，短正合列 $\Z/2\to\operatorname{Pin}_n\to O_n$ 诱导纤维化 $B\operatorname{Pin}_n\to BO_n$，纤维为 $B\Z/2=\RP^\infty$；它可与万有纤维化沿表示 $w_2$ 的映射 $BO_n\to K(\Z/2,\allowbreak 2)$ 的拉回识别。这只是 $w_2$ 与 Pin 结构关系的同伦理论表述。因此，可将 $\operatorname{Gr}(V)^\#$ 取为两个拉回 $\alpha_V^*(\R\to S^1)$ 与 $i_V^*(B\operatorname{Pin}_n\to BO_n)$ 的纤维积，其中 $i_V$ 是将 $\operatorname{Gr}(V)$ 嵌入通常的 $n$ 维子空间 Grassmann 流形，再嵌入 $BO_n$ 得到的映射。于是，对给定 $F\subset B\times V$，相应映射 $B\to\operatorname{Gr}(V)$ 到 $\operatorname{Gr}(V)^\#$ 的提升的同伦类，与二元组 $(\alpha^\#,\allowbreak P^\#)$ 的同构类之间存在双射。这在同伦层面回答了我们的问题，虽大体令人满意，但并不完全。不过，与其寻找形式上更完善的陈述，不如简要回顾主要应用；这会具体说明，膜结构的新定义在实际作用上等价于旧定义。
\SourcePage{179}
取一对线性 Lagrangian 膜 $\Lambda_k^\#=(\Lambda_k,\allowbreak \alpha_k^\#,\allowbreak P_k^\#)$（$k=0,\allowbreak 1$），使 $(\Lambda_0,\allowbreak \Lambda_1)$ 横截。如\secref{11h}，可赋予它们绝对指标和定向空间。具体地，取路径 $\lambda\in\Omega^-(\operatorname{Gr}(V),\allowbreak \Lambda_0,\allowbreak \Lambda_1)$，在相应丛 $F_\lambda\subset[0;1]\times V$ 上取膜结构 $(\alpha_\lambda^\#,\allowbreak P_\lambda^\#)$，使分次满足 $\alpha_\lambda^\#(k)=\alpha_k^\#$，并取同构 $f_k:(P_\lambda^\#)_k\cong P_k^\#$，两处均为 $k=0,\allowbreak 1$。其余完全如\eref{11.25}。值得细想其中一些细节，以澄清膜结构在哪里起作用。假设给定路径 $\lambda$。存在唯一函数 $\alpha_\lambda^\#:[0;1]\to\R$，满足 $\alpha_\lambda^\#(0)=\alpha_0^\#$ 和 $\exp(2\pi i\alpha_\lambda^\#(s))=\alpha_V(\lambda(s))$。按我们的规定，还须有 $\alpha_\lambda^\#(1)=\alpha_1^\#$；这对 $\lambda$ 相对端点的同伦类施加限制，实际上选出一个优先同伦类。相反，为 $F_\lambda$ 赋予 Pin 结构没有障碍；但考虑路径的单参数族时，在 $[0;1]^2$ 上的相应向量丛上相对边界赋予 Pin 结构，就会对同伦类施加非平凡限制，保证定向空间不会出现非平凡单值变换。
\begin{note}\label{st:11.19}
实际上，只取决于 $(P_\lambda^\#,\allowbreak f_0,\allowbreak f_1)$ 的同构类，所以有效地只有两种可能。要从一种变到另一种，只需将 $f_0$ 或 $f_1$ 乘以非平凡元素 $-e\in\ker(\operatorname{Pin}_n\to O_n)$；这将诱导定向空间的反定向自同构。
\end{note}
\begin{example}\label{st:11.20}
考察最低维情形 $n=1$。考虑带标准生成元 $v=1$ 的 $\R$。依定义，$\operatorname{Pin}_1=\{\pm e,\allowbreak \pm v\}\subset\operatorname{Cl}(\R)=\R[v]/(v^2-1)$，所以投影 $\operatorname{Pin}_1\to O_1$ 分裂；此后固定像为 $\{e,\allowbreak v\}$ 的分裂。因此，实线丛 $F\to B$ 上的 $\operatorname{Pin}_1$ 结构等价于在底空间 $B$ 上选择另一个实线丛。特别地，线性 Lagrangian 膜由实直线 $\Lambda\subset V$、实值相位 $\alpha^\#$ 和另一个抽象实直线 $\beta$ 组成。对于横截的两个这样的膜，
\begin{equation}\label{eq:11.35}
\begin{aligned}
i(\Lambda_0^\#,\Lambda_1^\#)&=\lfloor\alpha_1^\#-\alpha_0^\#\rfloor+1,\\
o(\Lambda_0^\#,\Lambda_1^\#)&\cong\Hom(\beta_0,\beta_1)\otimes\Lambda_1^{\otimes i(\Lambda_0^\#,\Lambda_1^\#)}.
\end{aligned}
\end{equation}
这里 $\lfloor\cdot\rfloor$ 为向下取整。只需选择合适路径，并展开\eref{11.25}，即可得到这些公式。
\end{example}

\subsection{平移运算}\label{sec:11k}
对任意 $\sigma\in\Z$，线性 Lagrangian 膜有一个 $\sigma$ 重平移运算，将 $\Lambda^\#=(\Lambda,\allowbreak \alpha^\#,\allowbreak P^\#)$ 映为
\begin{equation}\label{eq:11.36}
S^\sigma\Lambda^\#=(\Lambda,\alpha^\#-\sigma,P^\#\otimes\lambda^{\mathrm{top}}(\Lambda)^{\otimes\sigma}).
\end{equation}
$\sigma=1$ 时简记为 $S\Lambda^\#$。\znote{172}{平移 $S^\sigma$ 把实值相位减去 $\sigma$，同时把 Pin 结构张量上 $\lambda^{\mathrm{top}}(\Lambda)$ 的 $\sigma$ 次幂；两者必须同步进行，否则 $\exp(2\pi i\alpha^\#)=\alpha_V(\Lambda)$ 不再成立。这个运算说明分次并非一个绝对数字：它只在相差一个整数的意义下由 $\Lambda$ 决定，而 $S^\sigma$ 恰把这个选择整体平移。绝对指标 $i$ 正是把这种“相差整数”的分次固定成 $\Z$ 分级的工具，\eref{11.37}的第一条说明平移 $\sigma$ 使指标减少 $\sigma$。nLab 的 Maslov 指标条目指出该指标取整值，并对路径的同伦不变。[来源：nLab "Maslov index"，https:/\allowbreak /\allowbreak ncatlab.org/\allowbreak nlab/\allowbreak show/\allowbreak Maslov+\allowbreak index（2026-11-15 访问）]}其关键性质是：
\SourcePage{180}
\begin{lemma}\label{st:11.21}
对满足 $\Lambda_0\pitchfork\Lambda_1$ 的任意线性 Lagrangian 膜 $\Lambda_0^\#,\allowbreak \Lambda_1^\#$，
\begin{equation}\label{eq:11.37}
\begin{aligned}
i(\Lambda_0^\#,S^\sigma\Lambda_1^\#)&=i(\Lambda_0^\#,\Lambda_1^\#)-\sigma,\\
o(\Lambda_0^\#,S^\sigma\Lambda_1^\#)&\cong o(\Lambda_0^\#,\Lambda_1^\#).
\end{aligned}
\end{equation}
\end{lemma}
\begin{proof}
指标结论直接可得，故只考虑定向空间。为简单起见，假设 $\sigma=1$。在 Lagrangian Grassmann 流形中，取从 $\Lambda_0$ 到 $\Lambda_1$ 的路径 $\lambda$，以及从 $\Lambda_1$ 回到自身的环路 $\rho$，将其串接记为 $\lambda\cdot\rho$。相应地，在上半平面和闭圆盘上有算子 $D_{H,\allowbreak \lambda}$、$D_{H,\allowbreak \lambda\cdot\rho}$ 和 $D_{D,\allowbreak \rho}$，关系如下：作边界连通和，将 $s\in\partial H$、$s\gg0$ 与 $1\in\partial D$ 识别，则按此方式粘合 $D_{H,\allowbreak \lambda}$ 与 $D_{D,\allowbreak \rho}$ 所得算子可与 $D_{H,\allowbreak \lambda\cdot\rho}$ 识别。

现在假设 $\lambda$ 使函数 $\alpha_V(\lambda(t))$ 有一个实值提升，在端点取值 $\alpha_0^\#,\allowbreak \alpha_1^\#$。再假设 $\mu(\rho)=-1$，于是将该实值提升延拓到 $\lambda\cdot\rho$ 后，新端点的值为 $\alpha_1^\#-1$。接着考虑相应丛 $F_\lambda\to[0;1]$ 和 $\widetilde F_\rho\to[0;1]$；后者的记号与\secref{11i}一致。可以为它们赋予 Pin 结构 $P_\lambda^\#,\allowbreak \widetilde P_\rho^\#$ 和同构
\begin{equation}\label{eq:11.38}
\begin{cases}(P_\lambda^\#)_0\cong P_0^\#,\\(P_\lambda^\#)_1\cong P_1^\#,
\end{cases}\qquad
\begin{cases}(\widetilde P_\rho^\#)_0\cong P_1^\#,\\(\widetilde P_\rho^\#)_1\cong P_1^\#\otimes\lambda^{\mathrm{top}}(\Lambda_1).
\end{cases}
\end{equation}
利用\eref{11.38}的第二、第三个同构将两者结合，得到 $F_{\lambda\cdot\rho}$ 上的 Pin 结构，且在端点分别有到 $P_0^\#$ 和 $P_1^\#\otimes\lambda^{\mathrm{top}}(\Lambda_1)$ 的典范同构。

由\stref{引理}{11.11}及定向空间的定义，$F_\lambda$ 和 $F_{\lambda\cdot\rho}$ 上 Pin 结构的这些选择确定同构 $\det(D_{H,\allowbreak \lambda})\cong o(\Lambda_0^\#,\allowbreak \Lambda_1^\#)$ 与 $\det(D_{H,\allowbreak \lambda\cdot\rho})\cong o(\Lambda_0^\#,\allowbreak S\Lambda_1^\#)$。另一方面，$\widetilde P_\rho^\#$ 是\eref{11.32}意义下的扭曲 Pin 结构，由\stref{引理}{11.17}给出同构 $\det(D_{D,\allowbreak \rho})\cong\lambda^{\mathrm{top}}(\Lambda_1)$。结合\eref{11.11}，得到
\[
\begin{aligned}
o(\Lambda_0^\#,S\Lambda_1^\#)\otimes\lambda^{\mathrm{top}}(\Lambda_1)
&\cong\det(D_{H,\lambda\cdot\rho})\otimes\lambda^{\mathrm{top}}(\Lambda_1)\\
&\cong\det(D_{H,\lambda})\otimes\det(D_{D,\rho})\\
&\cong o(\Lambda_0^\#,\Lambda_1^\#)\otimes\lambda^{\mathrm{top}}(\Lambda_1).
\end{aligned}
\]
\end{proof}
\begin{note}\label{st:11.22}
也可不引用\stref{引理}{11.17}，而直接照用其部分证明。具体地，取某个分解 $\Lambda_1=\Lambda_{1,\allowbreak +}\oplus\Lambda_{1,\allowbreak -}$，其中 $\dim(\Lambda_{1,\allowbreak -})=1$，并令 $\rho(s)=\Lambda_{1,\allowbreak +}\oplus\exp(-\pi sI_V)\Lambda_{1,\allowbreak -}$。由\eref{11.20}可见，$(\delta_H)_{\lambda\cdot\rho}$ 比 $(\delta_H)_\lambda$ 多出一个张量因子 $\Lambda_{1,\allowbreak -}$。但扭曲 Pin 结构的选择恰对应于同一空间的定向，因而可识别 $(\delta_H)_{\lambda\cdot\rho}\cong(\delta_H)_\lambda$。此论证避免了指标理论，较为初等；不过，对于理解\eref{11.37}与\eref{11.26}中同构的相互作用等目的，原证明更方便。
\end{note}

\SourcePage{181}
\subsection{定向算子}\label{sec:11l}
将\stref{命题}{11.13}适配到膜的新定义很直接。具体地，仍取前述 $(S,\allowbreak \nu)$，并为 Lagrangian 子丛 $F_\nu\subset\partial S\times V$ 赋予膜结构。对每个条带状端 $\zeta$，便得到一对线性 Lagrangian 膜 $(\Lambda_{\zeta,\allowbreak 0}^\#,\allowbreak \Lambda_{\zeta,\allowbreak 1}^\#)$，其绝对指标和定向空间按\eref{11.26}确定 $D_{S,\allowbreak \nu}$ 的指标与行列式线。本节最后的任务是将这一结论推广到稍大的一类算子，以便用于非齐次伪全纯映射的线性化算子。

回顾\secref{8h}中定义非退化 $\bar\partial$ 算子的数据：$S$ 是带条带状端的 Riemann 曲面，$(E,\allowbreak \omega_E,\allowbreak J_E)\to S$ 是 Hermite 向量丛，$\nabla$ 是 $E$ 上的联络，$F\subset E|\partial S$ 是 Lagrangian 子丛。它们在条带状端上满足适当的渐近收敛条件。为方便应用，在每个端添加区间 $\epsilon_\zeta(\{\pm\infty\}\times[0;1])$ 来紧化，得到带角曲面 $\bar S$。以 $\overline{\partial S}$ 表示 $\partial S$ 在 $\bar S$ 中的闭包，即 $\partial S$ 与各点 $\epsilon_\zeta(\pm\infty,\allowbreak k)$（$k=0,\allowbreak 1$）的并。由渐近行为，上述结构分别在 $\bar S$ 上、最后一个则在 $\overline{\partial S}$ 上，有典范延拓 $(\bar E,\allowbreak \omega_{\bar E},\allowbreak J_{\bar E},\allowbreak \bar\nabla,\allowbreak \bar F)$；限制到无穷远区间时，这些延拓恰是最初所说的极限数据 $(E_\zeta,\allowbreak \omega_{E_\zeta},\allowbreak J_{E_\zeta},\allowbreak \nabla_\zeta,\allowbreak \Lambda_{\zeta,\allowbreak k})$。现在再添加以下数据：
\begin{itemize}
\item 另一个与 $\omega_{\bar E}$ 相容的复结构 $I_{\bar E}$，连同逐纤维二次复体积形式 $\eta_{\bar E}^2$ 及相应平方相位函数 $\alpha_{\bar E}$；
\item 子丛 $\bar F\subset\bar E|\overline{\partial S}$ 上的膜结构，由分次 $\bar\alpha^\#$ 和 Pin 结构 $\bar P^\#$ 组成。
\end{itemize}
取一个条带状端 $\zeta$。考察 $k=0,\allowbreak 1$ 对应的点 $\epsilon_\zeta(\pm\infty,\allowbreak k)$，上述数据使每个 $\Lambda_{\zeta,\allowbreak k}$ 成为线性 Lagrangian 膜 $\Lambda_{\zeta,\allowbreak k}^\#$。起初它们位于不同的辛向量空间 $E_{\zeta,\allowbreak k}$ 中，但可用辛联络 $\nabla_\zeta$ 的平行移动，结合相应纤维 $\bar E_z$、$z\in\epsilon_\zeta(\{\pm\infty\}\times[0;1])$ 上的平方相位映射 $\alpha_{\bar E}$，将 $\Lambda_{\zeta,\allowbreak 0}^\#$ 移到 $E_{\zeta,\allowbreak 1}$ 中，从而定义绝对指标和定向空间。我们断言：考虑到这一点后，算子 $\bar\partial_\nabla$ 满足\eref{11.26}的对应公式。

证明相当直接。选择辛平凡化 $\bar E\cong\bar S\times V$，使之在无穷远区间上与联络 $\nabla_\zeta$ 相容。接着，要将 $I_{\bar E}$ 变形成常值复结构 $I_V$，并相应地将 $\eta_{\bar E}^2$ 变形成 $\eta_V^2$。后一部分在 $H^1(S)$ 中有一个显然的障碍，但只要适当选取最初平凡化的同伦类，该障碍就消失。二次复体积形式的变形诱导平方相位映射的同伦，分次也应相应地随之移动。结果是，现在由映射 $\bar\rho:\overline{\partial S}\to\operatorname{Gr}(V)$ 给出的 Lagrangian 子丛 $\bar F$，作为
\SourcePage{182}
$S\times(V,\allowbreak \omega_V,\allowbreak I_V)$ 的子丛带有膜结构。最后，将 $J_{\bar E}$ 变形成常值复结构 $I_V$，将 $\bar\nabla$ 变形成平凡联络。这诱导 Fredholm 算子空间内从 $\bar\partial_\nabla$ 到标准 Dolbeault 算子 $\bar\partial$ 的同伦，将问题归结到先前考虑的特殊情形。$S=H$ 时，得到\stref{引理}{11.11}的一个小推广：
\begin{equation}\label{eq:11.39}
\operatorname{index}(\bar\partial_\nabla)=i(\Lambda_0^\#,\Lambda_1^\#),\qquad
\det(\bar\partial_\nabla)\cong o(\Lambda_0^\#,\Lambda_1^\#).
\end{equation}
\begin{definition*}
这一情形足够重要，值得单独命名：称 $\bar\partial_\nabla$ 为线性 Lagrangian 膜对 $(\Lambda_0^\#,\allowbreak \Lambda_1^\#)$ 的\emph{定向算子}。\znote{174}{定向算子是带 Lagrangian 边界条件的非退化 $\bar\partial$ 算子的一般名称：它的指标等于两个边界膜的绝对指标，它的行列式线与定向空间同构。nLab 关于 Maslov 指标的条目列举了这类指标定理的经典来源（Robbin–Salamon 的两篇文章），其中路径情形正是 $S=H$ 的特例。名称中的“定向”指\eref{11.39}第二条：算子的行列式线被膜数据定向化以后，Floer 理论中相交点所携带的符号才能一致地确定。[来源：nLab "Maslov index"，https:/\allowbreak /\allowbreak ncatlab.org/\allowbreak nlab/\allowbreak show/\allowbreak Maslov+\allowbreak index（2026-11-15 访问）]}
\end{definition*}

\section{Fukaya 范畴（完整版本）}\label{sec:12}\label{section:12}
\subsection{Lagrangian 膜}\label{sec:12a}
设 $(M,\allowbreak \omega_M)$ 为满足 $2c_1(M)=0$ 的辛流形，配有相容殆复结构 $I_M$。选择二次复体积形式 $\eta_M^2$，即双典范丛 $K_M^2=\lambda_{\CC}^{\mathrm{top}}(TM)^{\otimes-2}$ 的一个处处非零截面；按假设，该丛平凡。令 $\alpha_M:\operatorname{Gr}(TM)\to S^1$ 为相应的平方相位映射，见\eref{11.34}。
\begin{definition*}
\emph{Lagrangian 膜}是三元组 $L^\#=(L,\allowbreak \alpha^\#,\allowbreak P^\#)$，由 Lagrangian 子流形 $L\subset M$ 及子丛 $TL\subset TM|L$ 上的膜结构 $(\alpha^\#,\allowbreak P^\#)$ 组成。具体而言，就是满足 $\exp(2\pi i\alpha^\#(x))=\alpha_M(TL_x)$ 的函数 $\alpha^\#:L\to\R$，即 $L$ 的分次，以及 $TL$ 的 Pin 结构 $P^\#$。
\end{definition*}
这种结构的基本理论是\secref{11j}一般讨论的特例。为与辛拓扑中更常用的概念联系，应指出：作为分次存在障碍的 Maslov 类 $\mu_L\in H^1(L)$，经边界映射 $H^1(L)\to H^2(M,\allowbreak L)$ 映为相对第一 Chern 类的两倍 $2c_1(M,\allowbreak L)$。另一方面，$\mu_L$ 的模二约化为 $w_1(L)-w_1(\xi_M|L)$，其中 $\xi_M\subset\lambda_{\CC}^{\mathrm{top}}(TM)$ 是使 $\pm\eta_M$ 取实值的实线丛。这与\stref{笔记}{11.18}密切相关：所有 Lagrangian 膜 $L^\#$ 都相对于 $\xi_M|L$ 典范定向。特别地，若 $\xi_M$ 平凡，等价地 $\eta_M^2$ 有全局平方根 $\eta_M$，则 Lagrangian 膜具有通常的定向；更确切地说，平方根的选择固定了这些定向。用 $L^\#\otimes\beta$ 表示以实线丛 $\beta\to L$ 改变 Pin 结构 $P^\#$ 的运算，以 $S^\sigma L^\#$ 表示逐点由\eref{11.36}给出的 $\sigma$ 重平移。\znote{174}{膜的定义把两件独立的几何数据拼在一起：$\alpha^\#$ 是平方相位映射的提升，它存在的充要条件是 Maslov 类 $\mu_L$ 为零；$P^\#$ 则是 $TL$ 上的 Pin 结构，用来处理 $TL$ 一般不可定向的情形。当 $\xi_M$ 平凡、从而每个 $L$ 可定向时，Pin 结构退化为通常的定向（至多差一个整体符号），这也是下文常直接用定向代替 Pin 数据的理由。分清这两项各管什么，才不会把“分次”与“定向”混为一谈。[来源：Seidel, Fukaya categories and Picard–Lefschetz theory, §11j 与例题 11.20。]}
\begin{note}\label{st:12.1}
Lagrangian 膜概念实质上只依赖 $\eta_M^2$ 的同伦类，有由 $H^1(M)$ 给出的实际选择自由度。此外，定义还有两个变体。首先，可固定类
\SourcePage{183}
$d\in H^1(M;\Z/2)$，用实线丛 $\delta$ 表示它，并将 $TL$ 上的 Pin 结构换成 $TL\oplus\delta|L$ 上的 Pin 结构。其次，给定 $w\in H^2(M;\Z/2)$，可考虑\stref{笔记}{11.16}意义下关于 $w|L$ 的扭曲 Pin 结构；也可将这两个想法结合。这里不继续展开，有兴趣的读者可参见\bcite{60}。不过，仍想概述如何将\secref{11h}的构造整体化，为这些可能性建立统一框架。

给定逐纤维分别限制为 Maslov 类与第二 Stiefel–Whitney 类的上同调类 $\mu_M\in H^1(\operatorname{Gr}(TM))$ 与 $w_M\in H^2(\operatorname{Gr}(TM);\Z/2)$，可如\eref{11.22}构造纤维化 $\operatorname{Gr}(TM)^\#\to\operatorname{Gr}(TM)$，其纤维为 $\Z\times\RP^\infty$。于是，抽象 Lagrangian 膜可定义为 Lagrangian 子流形 $L\subset M$，连同典范截面 $L\to\operatorname{Gr}(TM)|L$ 到 $\operatorname{Gr}(TM)^\#$ 的提升。为与此前术语联系，注意平方相位映射在二次体积形式 $\eta_M^2$ 的同伦类与上同调类 $\mu_M$ 之间建立双射。同样，最一般的类 $w_M$ 具有形式 $w_2(F)+w_1(F)\cdot\pi^*d+\pi^*w$，其中 $F\to\operatorname{Gr}(TM)$ 是万有丛，$\pi:\operatorname{Gr}(TM)\to M$ 是投影，$(d,\allowbreak w)$ 如前。\znote{175}{纤维 $\Z\times\RP^\infty$ 记录的正是膜结构的两部分：$\Z$ 分量给出 $\alpha^\#$ 的整值提升（$\alpha^\#$ 可整体加整数而不改变 $\alpha_M$），而 $\RP^\infty=K(\Z/2,1)$ 分类 Pin 结构，因为实线丛的同构类由 $H^1(\cdot;\Z/2)$ 刻画。所谓抽象膜，就是把 $L$ 的 Gauss 映射连同这两项额外拓扑数据一起提升为 $\operatorname{Gr}(TM)^\#$ 的截面；具体膜则对应于这些截面中满足相容条件的那些。注意 $\Z$ 分量并非多余：平移 $S^k$ 正是在这个因子上作加法，所以膜与它的平移在抽象膜空间里分居不同连通分支。[来源：Seidel, Fukaya categories and Picard–Lefschetz theory, §11h，尤见构造 (11.22)。]}
\end{note}

\subsection{模空间的定向}\label{sec:12b}
从现在起，回到通常的框架：$(M,\allowbreak \omega_M,\allowbreak \theta_M,\allowbreak I_M)$ 为带角正合辛流形，仍带有 $\eta_M^2$，且只考虑正合 Lagrangian 膜，即底层 Lagrangian 子流形正合的膜。

设 $(L_0^\#,\allowbreak L_1^\#)$ 为一对正合 Lagrangian 膜，$L_0,\allowbreak L_1$ 横截相交。在每个交点 $y\in L_0\cap L_1$ 的切空间 $(TM_y,\allowbreak \omega_{M,\allowbreak y},\allowbreak I_{M,\allowbreak y},\allowbreak \eta_{M,\allowbreak y}^2)$ 中，有显然的一对线性 Lagrangian 膜
\begin{equation}\label{eq:12.1}
\Lambda_{k,y}^\#=(\Lambda_{k,y}=(TL_k)_y,\ \alpha_{k,y}^\#=\alpha_k^\#(y),\ P_{k,y}^\#=(P_k^\#)_y).
\end{equation}
取\eref{11.25}定义的指标与定向空间，记为 $i(y)=i(\Lambda_{0,\allowbreak y}^\#,\allowbreak \Lambda_{1,\allowbreak y}^\#)$、$o(y)=o(\Lambda_{0,\allowbreak y}^\#,\allowbreak \Lambda_{1,\allowbreak y}^\#)$。更一般地，允许任意 $(L_0,\allowbreak L_1)$，选择 Floer 数据 $(H,\allowbreak J)$，并考虑 $y\in\mathcal C(L_0,\allowbreak L_1)$。此时\eref{12.1}中的 $\Lambda_{k,\allowbreak y}^\#$ 起初位于不同向量空间 $TM_{y(k)}$ 中；不过，与\secref{11l}的讨论平行，可利用 $H$ 的 Hamilton “流”在 $y$ 附近的线性化，结合 $\operatorname{Gr}(TM_{y(t)})$ 上的平方相位函数，将 $\Lambda_{0,\allowbreak y}^\#$ 移到 $TM_{y(1)}$ 中。然后如前定义 $i(y)$ 与 $o(y)$。由\eref{11.39}，也可将它们视为任意选取的、对应于 $(\Lambda_{0,\allowbreak y}^\#,\allowbreak \Lambda_{1,\allowbreak y}^\#)$ 的定向算子的指标与行列式线，将这种算子记为 $D_{H,\allowbreak y}$。

假设 $S$ 为配有条带状端及膜标记 $L_0^\#,\allowbreak \ldots,\allowbreak L_d^\#$ 的 $(d+1)$ 点圆盘；膜标记就是连同膜结构的 Lagrangian 标记。为 $(L_0,\allowbreak L_d)$ 及 $(L_{k-1},\allowbreak L_k)$（$k=1,\allowbreak \ldots,\allowbreak d$）选择 Floer 数据，并在 $S$ 上选择相容
\SourcePage{184}
扰动数据。在正则点 $u\in\M_S(y_0,\allowbreak \ldots,\allowbreak y_d)$ 处，有
\begin{equation}\label{eq:12.2}
\begin{aligned}
\dim_u\M_S(y_0,\ldots,y_d)&=i(y_0)-i(y_1)-\cdots-i(y_d),\\
(\lambda^{\mathrm{top}}T\M_S(y_0,\ldots,y_d))_u
&\cong\lambda^{\mathrm{top}}\ker(D_{S,u})\\
&\cong o(y_0)\otimes o(y_1)^\vee\otimes\cdots\otimes o(y_d)^\vee.
\end{aligned}
\end{equation}
第二式的意思是：模空间相对于各端的定向空间典范定向。\eref{12.2}是\stref{命题}{11.13}的推论，不过，直接重述相关粘合论证也许更有启发。为 $k=1,\allowbreak \ldots,\allowbreak d$ 选择定向算子 $D_{H,\allowbreak y_k}$。依定义，它们是 $H$ 上的 $\bar\partial$ 算子，大致说来，其边界条件子丛带有膜结构。线性化算子 $D_{S,\allowbreak u}$ 带类似数据。因此，依次将 $D_{S,\allowbreak u}$ 与 $D_{H,\allowbreak y_d},\allowbreak \ldots,\allowbreak D_{H,\allowbreak y_1}$ 粘合，并与给定膜结构相容，便得到 $(\Lambda_{0,\allowbreak y_0}^\#,\allowbreak \Lambda_{1,\allowbreak y_0}^\#)$ 的一个定向算子，记为 $D_{H,\allowbreak y_0}$。粘合公式与\eref{11.39}立即给出
\begin{equation}\label{eq:12.3}
\begin{aligned}
i(y_0)&=\operatorname{index}(D_{S,u})+i(y_1)+\cdots+i(y_d),\\
o(y_0)&\cong\det(D_{S,u})\otimes o(y_d)\otimes\cdots\otimes o(y_1).
\end{aligned}
\end{equation}
正则情形下，$D_{S,\allowbreak u}$ 满射，且 $\ker(D_{S,\allowbreak u})\cong T\M_S(y_0,\allowbreak \ldots,\allowbreak y_d)_u$，故得\eref{12.2}。
\begin{note}\label{st:12.2}
更精确地，这里的约定是对偶带有典范配对 $o(y_k)^\vee\otimes o(y_k)\to\R$。特别地，从\eref{12.2}中的同构出发，在右侧按递减次序与所有 $o(y_k)$ 作张量积，就还原出\eref{12.3}。
\end{note}

需要知道，这些定向在粘合非齐次伪全纯映射这一基本运算下如何变化。取膜标记为 $(L_0^\#,\allowbreak \ldots,\allowbreak L_n^\#,\allowbreak L_{n+m}^\#,\allowbreak \ldots,\allowbreak L_d^\#)$ 的 $(d+2-m)$ 点圆盘 $S_1$，以及标记为 $(L_n^\#,\allowbreak \ldots,\allowbreak L_{n+m}^\#)$ 的 $(m+1)$ 点圆盘 $S_2$。以某个粘合长度 $l$，将 $S_1$ 的第 $n+1$ 个出向端与 $S_2$ 的入向端粘合，得到继承标记 $(L_0^\#,\allowbreak \ldots,\allowbreak L_d^\#)$ 的 $(d+1)$ 点圆盘 $S$。为所有相关的 Lagrangian 子流形对选择 Floer 数据，在 $S_1,\allowbreak S_2$ 上选择相容扰动数据，并诱导 $S$ 上的扰动数据。假设对某个 $\tilde y$，$u_1\in\M_{S_1}(y_0,\allowbreak \ldots,\allowbreak y_n,\allowbreak \tilde y,\allowbreak y_{n+m+1},\allowbreak \ldots,\allowbreak y_d)$ 与 $u_2\in\M_{S_2}(\tilde y,\allowbreak y_{n+1},\allowbreak \ldots,\allowbreak y_{n+m})$ 是各自模空间的正则点。只要 $l$ 足够大，先用单位分解直接拼接 $u_k$，再应用隐函数定理，粘合过程就产生正则点 $u\in\M_S(y_0,\allowbreak \ldots,\allowbreak y_d)$。这一过程延拓到模空间的邻近点，实际上给出 $(u_1,\allowbreak u_2)$ 与 $u$ 的邻域之间的局部微分同胚 $g$。即使从这段粗略描述中，也应能理解其微分 $Dg$ 近似于线性化算子的\eref{11.9}
\SourcePage{185}
所给映射。为得出所需结论，考虑下图：
\SourceEquationDiagram{diagram-12-4}{0.65423006}{12.4}
水平同构是\eref{12.2}，右侧向上箭头是 $o(\tilde y)$ 与其对偶之间的典范配对。则\eref{12.4}在相差 Koszul 约定规定的符号的意义下交换，该符号为 $(-1)^\Delta$，其中
\[
\Delta=\left(\sum_{k>n+m}i(y_k)\right)\operatorname{index}(D_{S_2,u_2}).
\]
此符号的来源可追溯到行列式线的定义，见\secref{11a}，尤其是\eref{11.2}。\znote{177}{这个 Koszul 符号是记账方式而非笔误：行列式线是分次线，两个奇次因子交换次序会带来 $-1$。粘合公式 (11.9) 把 $D_{S,u}$ 与各 $D_{H,y_k}$ 依次拼接，每交换一次记一个符号，累积起来即 $(-1)^\Delta$，其中 $\Delta$ 由外侧交点的指标之和与 $D_{S_2,u_2}$ 的指标相乘得到。把这一条弄懂，12g 中 $\ddagger$、$\diamond$、$\dagger$ 三个符号就有了统一的来源，不必逐条死记。读者不妨按 (12.22) 自己推一遍 $\ddagger=m(d-n)+m+n$，看每一项来自哪一次交换。[来源：Seidel, Fukaya categories and Picard–Lefschetz theory, §11a–11b，尤见 (11.2) 与命题 11.13。]}

上述讨论可以沿两个方向推广。为叙述方便，分别处理这两个方向，将二者结合这一直接的工作留给读者。

\subsection{更一般的 Riemann 曲面}\label{sec:12c}
第一个想法是扩大 Riemann 曲面 $S$ 的范围。先看紧化 $\hat S$ 的每个边界圆周都恰有一个入向无穷远点的情形。对适当模空间 $\M_S(\{y_\zeta\})$ 的正则点 $u$，\stref{命题}{11.13}表明
\[
\begin{aligned}
\dim_u\M_S(\{y_\zeta\})&=n(\chi(S)-|\Sigma^-|)+\sum_{\zeta\in\Sigma^\pm}\mp i(y_\zeta),\\
(\lambda^{\mathrm{top}}T\M_S(\{y_\zeta\}))_u&\cong\bigotimes_{\zeta\in\Sigma^\pm}o(y_\zeta)^{\mp1},
\end{aligned}
\]
而\eref{12.4}的推广没有困难。只要每个边界圆周包含奇数个入向无穷远点，同样成立；不过，粘合时追踪符号行为更棘手，因为\stref{命题}{11.13}的相关部分多了一个步骤。

余下情形稍不直接。对 $\hat S$ 的每个包含偶数个入向无穷远点的边界圆周，选择属于该圆周的 $\partial S$ 的一个分支，再在该分支上选一个标记点。将这些选择记为 $(C_j,\allowbreak z_j)$，再运用\stref{命题}{11.13}得到
\begin{equation}\label{eq:12.5}
(\lambda^{\mathrm{top}}T\M_S(\{y_\zeta\}))_u
\cong\bigotimes_j(\lambda^{\mathrm{top}}TL_{C_j})_{u(z_j)}\otimes\bigotimes_{\zeta\in\Sigma^\pm}o(y_\zeta)^{\mp1}.
\end{equation}
消除额外因子的显然办法是确保 $L_{C_j}$ 有定向，例如假定 $\xi_M$ 平凡。另一种技巧是将额外因子吸收到定向空间的改变中。为使记号简单，
\SourcePage{186}
只说明一个特殊情形：$S$ 为具有 $e$ 个入向和 $d$ 个出向无穷远点的圆盘，其中 $e$ 为偶数，$d>0$，排列如\sr{fig:12.1}{图~12.1}。将涉及的 Lagrangian 膜记为 $L_k^{\pm,\allowbreak \#}$，伪全纯映射 $u$ 的极限记为 $y_k^\pm$。将标记点 $z$ 向 $\zeta_d^+$ 移动，使 $(TL_d^+)_{u(z)}$ 渐近等于 $\Lambda_{y_d^+,\allowbreak 1}$，因而可识别
\begin{equation}\label{eq:12.6}
\begin{aligned}
&(\lambda^{\mathrm{top}}TL_d^+)_{u(z)}\otimes o(\Lambda_{y_d^+,0}^\#,\Lambda_{y_d^+,1}^\#)\\
&\qquad\cong o(\Lambda_{y_d^+,0}^\#,\Lambda_{y_d^+,1}^\#\otimes\lambda^{\mathrm{top}}\Lambda_{y_d^+,1}).
\end{aligned}
\end{equation}
据此，\eref{12.5}变为
\begin{equation}\label{eq:12.7}
\begin{aligned}
&\lambda^{\mathrm{top}}T\M_S(y_1^-,\ldots,y_e^-,y_1^+,\ldots,y_d^+)\\
&\quad\cong o(\Lambda_{y_e^-,0}^\#,\Lambda_{y_e^-,1}^\#)\otimes\cdots\otimes o(\Lambda_{y_1^-,0}^\#,\Lambda_{y_1^-,1}^\#)\\
&\qquad\otimes o(\Lambda_{y_1^+,0}^\#,\Lambda_{y_1^+,1}^\#)^\vee\otimes\cdots
\otimes o(\Lambda_{y_d^+,0}^\#,\Lambda_{y_d^+,1}^\#\otimes\lambda^{\mathrm{top}}\Lambda_{y_d^+,1})^\vee.
\end{aligned}
\end{equation}
\SourceFigure{figure-12-1}{0.30323044}{12.1}

\subsection{带标记点圆盘的族}\label{sec:12d}
沿另一个方向，取膜标记为 $(L_0^\#,\allowbreak \ldots,\allowbreak L_d^\#)$ 的 $(d+1)$ 点圆盘族 $\mathcal S\to\RG$。选择扰动数据，考虑相应的一个参数化模空间 $\M_{\mathcal S}(y_0,\allowbreak \ldots,\allowbreak y_d)$。正则点 $(r,\allowbreak u)$ 处的切空间现在由\eref{9.18}中的扩展算子 $D_{\mathcal S,\allowbreak r,\allowbreak u}$ 的核给出。将该算子的第一个有限维分量变形成零，得到
\begin{equation}\label{eq:12.8}
\begin{aligned}
\dim_{(r,u)}\M_{\mathcal S}(y_0,\ldots,y_d)&=\dim_r\RG+i(y_0)-i(y_1)-\cdots-i(y_d),\\
(\lambda^{\mathrm{top}}T\M_{\mathcal S}(y_0,\ldots,y_d))_{(r,u)}
&\cong(\lambda^{\mathrm{top}}T\RG)_r\otimes o(y_0)\otimes o(y_1)^\vee\otimes\cdots\otimes o(y_d)^\vee.
\end{aligned}
\end{equation}
假设有两个族 $\mathcal S_k\to\RG_k$（$k=1,\allowbreak 2$），如前分别标为 $(L_0^\#,\allowbreak \ldots,\allowbreak L_n^\#,\allowbreak L_{n+m}^\#,\allowbreak \ldots,\allowbreak L_d^\#)$ 与 $(L_n^\#,\allowbreak \ldots,\allowbreak L_{n+m}^\#)$，并配有适当扰动数据。小节 \sr{sec:9e}{(9e)} 与
\SourcePage{187}
\sr{sec:9i}{(9i)} 所述的族粘合过程，在两顶点树这一特殊情形下，给出诱导族 $\mathcal S\to\RG$，其中 $\RG=(-1;0)\times\RG_1\times\RG_2$。对应的非齐次伪全纯映射粘合，取正则点 $(r_1,\allowbreak u_1)\in\M_{\mathcal S_1}(y_0,\allowbreak \ldots,\allowbreak y_n,\allowbreak \tilde y,\allowbreak y_{n+m+1},\allowbreak \ldots,\allowbreak y_d)$ 与 $(r_2,\allowbreak u_2)\in\M_{\mathcal S_2}(\tilde y,\allowbreak y_{n+1},\allowbreak \ldots,\allowbreak y_{n+m})$，以及足够接近零的粘合参数 $\rho\in(-1;0)$，产生元素 $(r,\allowbreak u)\in\M_{\mathcal S}(y_0,\allowbreak \ldots,\allowbreak y_d)$。此外，它延拓为 $(-1;0)\times\M_{\mathcal S_1}\times\M_{\mathcal S_2}$ 与 $\M_{\mathcal S}$ 之间的局部微分同胚 $g$。注意，由于其中用到了隐函数定理，$r$ 通常不等于 $(\rho,\allowbreak r_1,\allowbreak r_2)$。不过两点很接近；更精确地，$\rho\to0$ 时两者渐近接近。特别地，可识别
\begin{equation}\label{eq:12.9}
(\lambda^{\mathrm{top}}T\RG)_r\cong\R\otimes(\lambda^{\mathrm{top}}T\RG_1)_{r_1}\otimes(\lambda^{\mathrm{top}}T\RG_2)_{r_2}.
\end{equation}
\eref{12.4}的适当推广是下图：
\SourceEquationDiagram{diagram-12-10}{0.64176854}{12.10}
右侧向上箭头结合了配对 $o(\tilde y)^\vee\otimes o(\tilde y)\cong\R$ 与\eref{12.9}。该图在相差 $(-1)^\diamond$ 的意义下交换，其中
\[
\diamond=\dim(\RG_2)\operatorname{index}(D_{S_1,u_1})+
\left(\sum_{k>n+m}i(y_k)\right)\operatorname{index}(D_{S_2,u_2}).
\]
这里的指标是曲面 $S_k=(\mathcal S_k)_{r_k}$ 上通常线性化算子的指标，并非涉及参数空间的扩展算子的指标。

\subsection{膜的 Floer 上同调}\label{sec:12e}
先引入一个初等记号。固定系数域 $\K$。对任意一维实向量空间 $\delta$，以 $|\delta|_{\K}$ 表示由 $\delta$ 的两个定向生成、并施加两个生成元之和为零这一关系的一维 $\K$ 向量空间。显然，$\R$ 线性同构 $c:\delta_1\to\delta_2$ 诱导 $\K$ 线性同构 $|c|_{\K}:|\delta_1|_{\K}\to|\delta_2|_{\K}$。分别称其为 $\delta$ 与 $c$ 的 $\K$ \emph{规范化}。规范化也适用于族，将实线丛变成纤维为 $\K$ 的局部系数系统，对它们之间的同构亦然。可将它理解为通过同态 $\R^\times\to\{\pm1\}\subset\K^\times$、$t\mapsto t/|t|$ 取关联丛。

升级后的 Floer 上同调为 $M$ 中任意一对正合 Lagrangian 膜 $(L_0^\#,\allowbreak L_1^\#)$ 关联有限维分次 $\K$ 向量空间 $HF^*(L_0^\#,\allowbreak L_1^\#)$；取 $\operatorname{char}(\K)=2$ 并将所有
\SourcePage{188}
分次部分直和，就还原出 $HF^{\mathrm{pr}}(L_0,\allowbreak L_1)$。下面回顾\sr{sec:8}{第~8~节}讨论的若干基本性质，并说明有膜结构时如何适配。
\begin{itemize}
\item \emph{三角形乘积。}存在分次、结合且含幺的乘积
\begin{equation}\label{eq:12.11}
HF^*(L_1^\#,L_2^\#)\otimes HF^*(L_0^\#,L_1^\#)\longrightarrow HF^*(L_0^\#,L_2^\#).
\end{equation}
\item \emph{改变膜结构。}若对所有涉及的膜施加同一个运算，Floer 上同调群不变，\eref{12.11}也不变。这里“运算”可以是给分次加上同一 $\Z$ 中常数，也可以是将各 Pin 结构与某个共同线丛 $\beta\to M$ 的限制作张量积。另有不那么显然的典范同构，它也说明“平移”这一名称合理：
\begin{equation}\label{eq:12.12}
HF^*(L_0^\#,S^\sigma L_1^\#)\cong HF^{*+\sigma}(L_0^\#,L_1^\#).
\end{equation}
它们也与乘积相容，即下图交换：
\SourceEquationDiagram{diagram-12-13}{0.55453786}{12.13}
\item \emph{PSS 同构。}存在典范同构 $HF^*(L^\#,\allowbreak L^\#)\cong H^*(L;\K)$，在 $L_0^\#=L_1^\#=L_2^\#=L^\#$ 时将\eref{12.11}化为通常的杯积。更一般地，若用任意两个线丛 $\beta_0,\allowbreak \beta_1$ 扭曲 $L^\#$，结果是 $L$ 的、取值于相应局部系数系统的上同调：
\begin{equation}\label{eq:12.14}
HF^*(L^\#\otimes\beta_0,L^\#\otimes\beta_1)\cong H^*(L;|\beta_0^\vee\otimes\beta_1|_{\K}).
\end{equation}
\item \emph{对偶性。}对任意 $L^\#$，Poincaré 对偶与 PSS 同构给出映射
\begin{equation}\label{eq:12.15}
\begin{aligned}
\int_{L^\#}:HF^n(L^\#,L^\#\otimes\lambda^{\mathrm{top}}(TL))
&\cong H^n(L;|\lambda^{\mathrm{top}}TL|_{\K})\\
&\cong H_0(L;\K)\longrightarrow\K.
\end{aligned}
\end{equation}
由此构造的次数为 $-n$ 的配对
\[
\begin{gathered}
HF^*(L_1^\#,L_0^\#\otimes\lambda^{\mathrm{top}}(TL_0))\otimes HF^*(L_0^\#,L_1^\#)\longrightarrow\K,\\
\langle a_2,a_1\rangle=\int_{L_0^\#}a_2a_1
\end{gathered}
\]
非退化且分次对称。后一个性质需要略作解释。利用选定同构 $\lambda^{\mathrm{top}}(TL_k)\cong\xi_M|L_k$，可将 $a_2$ 视为 $HF^*(L_1^\#,\allowbreak L_0^\#\otimes\xi_M|L_0)$ 的元素，将 $a_1$ 视为
\SourcePage{189}
$HF^*(L_0^\#\otimes\xi_M|L_0,\allowbreak L_1^\#\otimes\lambda^{\mathrm{top}}(TL_1))$ 的元素。于是 $a_1a_2\in HF^*(L_1^\#,\allowbreak L_1^\#\otimes\lambda^{\mathrm{top}}(TL_1))$，因此下式有意义：
\[
\langle a_1,a_2\rangle=(-1)^{|a_1|\cdot|a_2|}\langle a_2,a_1\rangle.
\]
\item \emph{同痕不变性。}从 $L_0^\#$ 到 $L_1^\#$ 的正合 Lagrangian 膜同痕给出 $HF^0(L_0^\#,\allowbreak L_1^\#)$ 的一个选定元素。这些元素的性质仍如\secref{8k}；特别地，可用来证明 $HF^*$ 具有同痕不变性。
\end{itemize}
\begin{definition*}
以正合 Lagrangian 膜为对象、Floer 上同调群为态射，并用\eref{12.11}给出复合，定义 \emph{Donaldson–Fukaya 范畴} $H\F(M)$。
\end{definition*}
它是 $\K$ 线性的 $\Z$ 分次范畴；取 $\operatorname{char}(\K)=2$ 并忘掉分次，便退化为 $H\F(M)^{\mathrm{pr}}$。由于 Poincaré 对偶配对具有扭曲形式，一般情况下它不是 Frobenius 范畴，除非 $\xi_M$ 平凡。\znote{180}{配对右端落在 $\lambda^{\mathrm{top}}(TL_0)$ 的规范化里，而不是平凡地落在 $\K$：Frobenius 范畴要求存在与复合相容的非退化配对 $\operatorname{hom}(A,B)\cong\operatorname{hom}(B,A)^\vee$，这里多出的扭转使该同构带上一个线丛因子，故只有 $\xi_M$ 平凡、扭转可平凡化时才成立。第 12j 节把这一缺陷代数化：Donaldson–Fukaya 范畴满足的是弱 Calabi–Yau 性质，对偶由 Serre 函子而非恒等函子实现，次数为 $-\dim M$。换言之，Frobenius 性质并未完全丢失，而是被 $\lambda^{\mathrm{top}}$ 的扭转吸收掉了。[来源：nLab “Serre functor”，https:/\allowbreak /\allowbreak ncatlab.org/\allowbreak nlab/\allowbreak show/\allowbreak Serre+\allowbreak functor（2026-11-15 访问）；并见原书 §12j。]}

\subsection{确定分次与符号}\label{sec:12f}
有了全部准备，$HF^*(L_0^\#,\allowbreak L_1^\#)$ 的定义相当直接。选择正则 Floer 数据，对每个 $k\in\Z$ 定义分次 Floer 上链群
\begin{equation}\label{eq:12.16}
CF^k(L_0^\#,L_1^\#)=\bigoplus_{i(y)=k}|o(y)|_{\K}.
\end{equation}
设 $u\in\M_Z(y_0,\allowbreak y_1)$ 为 Floer 方程的非驻定解。有显然的短正合列
\[
0\longrightarrow\R\longrightarrow T\M_Z(y_0,y_1)_u\longrightarrow T\M_Z^*(y_0,y_1)_{[u]}\longrightarrow0,
\]
其中 $\R$ 项对应平移方向；更精确地，约定其正生成元为 $\partial u/\partial s$。由此可识别 $\M_Z$ 与 $\M_Z^*$ 的切空间的最高次外幂，再由\eref{12.2}，得到 $\M_Z^*(y_0,\allowbreak y_1)$ 相对于 $o(y_0)\otimes o(y_1)^\vee$ 的选定定向。特别地，若 $[u]$ 是 $\M_Z^*(y_0,\allowbreak y_1)$ 的孤立点，就得到同构
\begin{equation}\label{eq:12.17}
c_u:o(y_1)\longrightarrow o(y_0).
\end{equation}
将边界算子 $\partial$ 的 $(y_0,\allowbreak y_1)$ 系数定义为所有这种 $c_u$ 的 $\K$ 规范化之和：
\begin{equation}\label{eq:12.18}
\partial_{y_0,y_1}=\sum_u|c_u|_{\K}:|o(y_1)|_{\K}\longrightarrow|o(y_0)|_{\K}.
\end{equation}
这是一种较精巧的、不依赖基的说法，意思是每个 $u$ 依照 $c_u$ 的符号贡献 $\pm1$。再看\eref{12.2}即可知，$\partial$ 的次数为所需的 $1$。\znote{181}{$\partial$ 的每个矩阵元都是数出来的：一个孤立 Floer 轨线给出同构 $c_u:o(y_1)\to o(y_0)$，规范化后贡献 $\pm1$。这样写的好处是符号自动一致：模空间的定向由 $o(y_0)\otimes o(y_1)^\vee$ 给出，而 $\partial^2=0$ 归结为紧化边界上两个端点的定向相反，下文用粘合映射的微分 $Dg$ 逐项验证。同一套“计数 + 行列式线”模式在 12g 定义 $\mu^d$、13c 定义 $\tilde\mu^d$ 时反复出现，值得在此处就熟悉。这种先数点、再取规范化和的写法，使 $\mu^d$ 的符号完全由模空间定向决定，而不依赖任何临时约定。[来源：Seidel, Fukaya categories and Picard–Lefschetz theory, §12f 与 §9j。]}

为说明\eref{12.18}合理，取一对解 $u_1\in\M_Z(y_0,\allowbreak \tilde y)$、$u_2\in\M_Z(\tilde y,\allowbreak y_1)$。将 $u_2$ 向右平移足够远后，
\SourcePage{190}
可粘合二者，得到 $u\in\M_Z(y_0,\allowbreak y_1)$。用 Gromov 紧化\eref{9.28}来说，$[u]$ 是一维模空间 $\M_Z^*(y_0,\allowbreak y_1)$ 的点，靠近紧化的边界点 $([u_1],\allowbreak [u_2])$。由粘合映射 $g$ 的结构，可见其微分 $Dg_{(u_1,\allowbreak u_2)}$ 具有以下性质：$(\partial u_1/\partial s,\allowbreak \partial u_2/\partial s)$ 映为平移切向量 $\partial u/\partial s$，而 $(-\partial u_1/\partial s,\allowbreak \partial u_2/\partial s)$ 映为远离边界点 $(u_1,\allowbreak u_2)$ 的向量。沿相反方向 $(\partial u_1/\partial s,\allowbreak -\partial u_2/\partial s)$ 移动 $(u_1,\allowbreak u_2)$，对应于将两个部分拉开，也就是令粘合长度趋于无穷。结合\eref{12.4}可知，通过同构\eref{12.2}由 $c_{u_1}\circ c_{u_2}:o(y_1)\to o(y_0)$ 得到的 $T\M_Z^*(y_0,\allowbreak y_1)_{[u]}$ 的定向，始终指向远离 Gromov 边界的方向。因此，同一连通分支的两个边界点贡献相反的映射 $|c_{u_1}\circ c_{u_2}|_{\K}$，从而 $\partial^2=0$。如通常一样，定义 $HF^*(L_0^\#,\allowbreak L_1^\#)=H(CF^*(L_0^\#,\allowbreak L_1^\#),\allowbreak \partial)$。

现取膜标记为 $(L_0^\#,\allowbreak \ldots,\allowbreak L_d^\#)$ 的 $(d+1)$ 点圆盘 $S$，选择正则 Floer 数据与扰动数据。在每个孤立点 $u\in\M_S(y_0,\allowbreak \ldots,\allowbreak y_d)$ 处，相对定向\eref{12.2}给出同构 $c_u:o(y_d)\otimes\cdots\otimes o(y_1)\to o(y_0)$。如\eref{12.18}，用它们定义映射
\[
C\Phi_S:CF^*(L_{d-1}^\#,L_d^\#)\otimes\cdots\otimes CF^*(L_0^\#,L_1^\#)\longrightarrow CF^*(L_0^\#,L_d^\#).
\]
为证明 $C\Phi_S$ 为链映射，考察 Gromov 紧化的两类边界点：
\begin{itemize}
\item 第一类为 $(u_1,\allowbreak [u_2])\in\M_S(y_0,\allowbreak \ldots,\allowbreak y_n,\allowbreak \tilde y,\allowbreak y_{n+2},\allowbreak \ldots,\allowbreak y_d)\times\M_Z^*(\tilde y,\allowbreak y_{n+1})$。平移向量 $(0,\allowbreak \partial u_2/\partial s)$ 在 $Dg$ 下的像是 $\M_S(y_0,\allowbreak \ldots,\allowbreak y_d)$ 的一个远离边界的切向量。因此，由 $c_{u_1}\circ c_{u_2}$ 通过\eref{12.2}得到的局部定向，与标准向外定向相差 $(-1)$ 乘以\eref{12.4}中的 Koszul 符号，其中 $\operatorname{index}(D_{S_2,\allowbreak u_2})=1$。
\item 第二类为 $([u_1],\allowbreak u_2)\in\M_Z^*(y_0,\allowbreak \tilde y)\times\M_S(\tilde y,\allowbreak y_1,\allowbreak \ldots,\allowbreak y_d)$。$(\partial u_1/\partial s,\allowbreak 0)$ 的像指向边界，且 Koszul 符号平凡。
\end{itemize}
合起来得到
\begin{equation}\label{eq:12.19}
\sum_n(-1)^*C\Phi_S(y_d,\ldots,\partial(y_{n+1}),y_n,\ldots,y_1)
=\partial(C\Phi_S(y_d,\ldots,y_1)),
\end{equation}
其中 $*=i(y_{n+2})+\cdots+i(y_d)$。沿同样思路作类似但更简单的论证，可知复合律\eref{8.2}成立，且不带任何符号。如\secref{8b}，利用 $C\Phi_S$ 诱导的上同调映射 $\Phi_S$，定义 $HF^0(L^\#,\allowbreak L^\#)$ 中的恒等态射（$d=0$）、三角形乘积（$d=2$）及延拓映射（$d=1$）。

也值得提一下，如何利用 $S$ 上扰动数据的族 $(K^s,\allowbreak J^s)$、$s\in[0;1]$ 构造链同伦。考虑
\SourcePage{191}
相应的参数化模空间，用其中的孤立点定义映射
\[
h:CF^*(L_{d-1}^\#,L_d^\#)\otimes\cdots\otimes CF^*(L_0^\#,L_1^\#)\longrightarrow CF^*(L_0^\#,L_d^\#)[-1].
\]
模空间的局部定向来自\eref{12.8}，粘合方面使用\eref{12.10}，取 $\RG=[0;1]$。与先前论证相比，唯一差别是第二类情形中 Koszul 约定引入另一个符号 $(-1)^{\dim(\RG)}=-1$。结果是
\begin{equation}\label{eq:12.20}
\begin{aligned}
C\Phi_S^1-C\Phi_S^0={}&\sum_n(-1)^*h(y_d,\ldots,\partial(y_{n+1}),y_n,\ldots,y_1)\\
&+\partial(h(y_d,\ldots,y_1)),
\end{aligned}
\end{equation}
其中 $C\Phi_S^k$ 是由 $(K^k,\allowbreak J^k)$ 导出的链映射，$*$ 如\eref{12.19}。这已足以建立基本的良定义性质。

沿\secref{12c}的讨论，只要 $\hat S$ 的每个边界圆周都有奇数个入向无穷远点，就可毫无困难地将理论推广到其他曲面 $S$，由此得到适当意义下的拓扑量子场论。若要允许不具备这一性质的曲面，则须对膜施加适当的定向假设，或使用\eref{12.6}的技巧。例如，取 $S=\bar H$，其端为出向端，膜标记为 $L^\#$。\eref{12.7}说明，每个孤立元素 $u\in\M_S(y)$ 都给出同构
\[
c_u:o(\Lambda_{y,0}^\#,\Lambda_{y,1}^\#\otimes\lambda^{\mathrm{top}}(\Lambda_{y,1}))\longrightarrow\K.
\]
将它们规范化后合并，便得到次数为 $-n$ 的链映射 $CF^*(L^\#,\allowbreak L^\#\otimes\lambda^{\mathrm{top}}(TL))\to\K$。诱导的上同调映射恰为\eref{12.15}。

\subsection{Fukaya 范畴}\label{sec:12g}
回顾一下也许有益：Fukaya 范畴中出现 $\Ai$ 结合关系\eref{1.2}，很大程度上来自 Deligne–Mumford–Stasheff 紧化 $\overline\RG^{d+1}$ 的递归结构。更确切地，每棵有两个顶点的稳定 $d$ 叶树——按\sr{fig:9.5}{图~9.5}的记号为 $T(m,\allowbreak n)$——都标记一个余维一边界面
\begin{equation}\label{eq:12.21}
\overline\RG^{T(m,n)}\cong\overline\RG^{d+2-m}\times\overline\RG^{m+1}\subset\partial\overline\RG^{d+1}.
\end{equation}
这些面对应于\eref{1.2}中 $1<m<d$ 的各项。\znote{183}{$\Ai$ 关系的几何来源就是模空间边界的分层：$\overline\RG^{d+1}$ 的每个余维一边界面都对应一种把 $d+1$ 个点分成两组的括号化方式，边界算子在该面上的取值正是相邻括号化之差，求和为零即 (1.2)。这也是为什么 Fukaya 范畴的 $\mu^d$ 必须满足无穷多条关系：紧化本身就携带全部高阶数据。$\Ai$ 代数的标准定义（张量余代数上次数 $2-k$ 的余导子 $D$ 满足 $D^2=0$）正是这套分层结构的代数化身。[来源：nLab “A-infinity algebra”，https:/\allowbreak /\allowbreak ncatlab.org/\allowbreak nlab/\allowbreak show/\allowbreak A-infinity+\allowbreak algebra（2026-11-15 访问）。]}

由\eref{12.8}可见，定义 $\Ai$ 结构需要选择各空间 $\RG^{d+1}$ 的定向。选定后，比较\eref{12.21}中的定向，会对每个 $(m,\allowbreak n)$ 给出一个符号。目标是安排这些符号，使它们与\eref{12.10}产生的 Koszul 符号，以及定义 $\mu^d$ 时引入的其他辅助符号一起，最终给出正确版本的\eref{1.2}。从\eref{9.7}将 $\RG^{d+1}$ 定义为 $\operatorname{Conf}_{d+1}(\partial D)$ 关于 $\operatorname{Aut}(D)$ 的商开始。取 $\operatorname{Conf}_{d+1}(\partial D)$ 中由 $z_0,\allowbreak z_1,\allowbreak z_2$ 固定的构形 $(z_0,\allowbreak \ldots,\allowbreak z_d)$ 组成的子空间。$(z_3,\allowbreak \ldots,\allowbreak z_d)\in(\partial D)^{d-2}$ 的自然边界定向赋予该子空间一个定向。另一方面，它是 $\operatorname{Aut}(D)$ 作用的全局切片，因此可经投影与 $\RG^{d+1}$ 识别。将诱导定向赋予 $\RG^{d+1}$。一个简单但繁琐的
\SourcePage{192}
计算表明，\eref{12.21}中 $\overline\RG^{d+2-m}\times\overline\RG^{m+1}$ 的乘积定向，与 $\partial\overline\RG^{d+1}$ 的边界定向相差 $(-1)^\ddagger$，其中
\begin{equation}\label{eq:12.22}
\ddagger=m(d-n)+m+n.
\end{equation}
这里省略该计算。

假设已如\secref{9j}，完成建立初步版本 $\F(M)^{\mathrm{pr}}$ 所需的全部选择。真正 Fukaya 范畴 $\F(M)$ 定义的头几步显而易见：对象为 $M$ 中的正合 Lagrangian 膜，态射空间为
\[
\operatorname{hom}(L_0^\#,L_1^\#)=CF^*(L_0^\#,L_1^\#),
\]
一阶运算为加上符号的 Floer 边界映射，
\[
\mu^1(a)=(-1)^{|a|}\partial(a).
\]
注意，这使用的是对应于底层 Lagrangian 子流形对 $(L_0,\allowbreak L_1)$ 的 Floer 数据。原则上，也可使 Floer 数据及扰动数据依赖于膜结构。目前这只会徒增复杂性，但后面讨论 Fukaya 范畴的某些变体时将采用这种做法，见\secref{14e}。接着，取膜 $(L_0^\#,\allowbreak \ldots,\allowbreak L_d^\#)$，其中 $d\geq2$，固定 $y_0\in\mathcal C(L_0,\allowbreak L_d)$、$y_k\in\mathcal C(L_{k-1},\allowbreak L_k)$。运用\eref{12.8}及 $\RG^{d+1}$ 的给定定向，可知模空间 $\M^{d+1}(y_0,\allowbreak \ldots,\allowbreak y_d)$ 的维数为 $i(y_0)-i(y_1)-\cdots-i(y_d)+2-d$，且相对于 $o(y_0)\otimes o(y_1)^\vee\otimes\cdots\otimes o(y_d)^\vee$ 定向。现在假设 $(r,\allowbreak u)$ 属于零维模空间。则 $i(y_0)=i(y_1)+\cdots+i(y_d)+2-d$，这正是运算 $\mu^d$ 所需的分次；此外，得到选定同构
\begin{equation}\label{eq:12.23}
c_{r,u}:o(y_d)\otimes\cdots\otimes o(y_1)\longrightarrow o(y_0).
\end{equation}
如\eref{12.18}，将 $\mu^d$ 的 $(y_0,\allowbreak \ldots,\allowbreak y_d)$ 系数定义为相应规范化 $|c_{r,\allowbreak u}|_{\K}$ 之和，但额外乘以符号 $(-1)^\dagger$，其中
\begin{equation}\label{eq:12.24}
\dagger=i(y_1)+2i(y_2)+\cdots+di(y_d).
\end{equation}
立即指出，由于\eref{1.3}的约定，底层上同调范畴确实与此前定义的 $H\F(M)$ 一致。\znote{184}{$\mu^1=(-1)^{|a|}\partial$ 与 $(-1)^\dagger$ 都是纯约定：前者来自第 1 节的符号约定 (1.3)，后者把几何上得到的定向符号调整成 (1.2) 的标准形式。改动这些约定只会把整个 $\Ai$ 结构逐项乘以 $\pm1$，得到同构的 $\Ai$ 结构；因此符号选择不影响范畴，只影响书写。Remark 12.4 把这一点说得更明白：$\RG^{d+1}$ 的整体定向与 $(-1)^\dagger$ 是构造中仅有的两处人为选择，几何本身并不挑出优先的约定。[来源：Seidel, Fukaya categories and Picard–Lefschetz theory, §12g，尤见 (12.24) 与 Remark 12.4。]}
\begin{proposition}\label{st:12.3}
各 $\mu^d$ 构成一个 $\Ai$ 结构。
\end{proposition}
\begin{proof}
如上一小节较简单的例子，证明基于粘合论证，从\eref{9.30}中一对孤立点 $(r_k,\allowbreak u_k)$、$k=1,\allowbreak 2$ 出发。只考虑稳定情形 $1<m<d$，有意略去不稳定情形；纳入后者所需的额外论证紧随\eref{12.19}与\eref{12.20}的证明，留给读者作为练习。
\SourcePage{193}
取同构 $c_{r_k,\allowbreak u_k}$，将第二个的输出插入第一个的第 $n+1$ 个输入以复合。从定向空间看，就是将 $o(\tilde y)$ 与其对偶配对。为简洁，将结果记为 $\tilde c$。由上述一般讨论，$\tilde c$ 诱导一维模空间 $\M^{d+1}(y_0,\allowbreak \ldots,\allowbreak y_d)$ 的一个定向。用足够接近零的参数 $\rho$，将 $(r_k,\allowbreak u_k)$ 粘合成其中的点 $(r,\allowbreak u)$。用 Gromov 紧化\eref{9.28}来说，$(r,\allowbreak u)$ 靠近边界点 $(r_1,\allowbreak u_1,\allowbreak r_2,\allowbreak u_2)$，因而存在指向该边界的典范局部定向。两个定向相差符号 $(-1)^{\S}$。确定该符号之后，由构造可知，对所有点对 $(r_k,\allowbreak u_k)$ 求规范化 $(-1)^{\S}|\tilde c|_{\K}$ 的和为零。要由此得到 $(r_1,\allowbreak u_1,\allowbreak r_2,\allowbreak u_2)$ 对
\[
\mu^{d+1-m}(a_d,\ldots,\mu^m(a_{n+m},\ldots,a_{n+1}),a_n,\ldots,a_1)
\]
的 $(y_0,\allowbreak \ldots,\allowbreak y_d)$ 系数的贡献，还须加入\eref{12.24}中的相应符号：
\[
\begin{aligned}
\dagger_1={}&i(y_1)+2i(y_2)+\cdots+(n+1)i(\tilde y)\\
&+(n+2)i(y_{n+m+1})+\cdots+(d+1-m)i(y_d)\\
={}&i(y_1)+2i(y_2)+\cdots+ni(y_n)\\
&+(n+1)(i(y_{n+1})+\cdots+i(y_{n+m})+2-m)\\
&+(n+2)i(y_{n+m+1})+\cdots+(d+1-m)i(y_d),\\
\dagger_2={}&i(y_{n+1})+2i(y_{n+2})+\cdots+mi(y_{n+m}).
\end{aligned}
\]
剩下要证明
\begin{equation}\label{eq:12.25}
\S+\dagger_1+\dagger_2\equiv\sgnmark_n+(2i(y_1)+\cdots+(d+1)i(y_d))\pmod2,
\end{equation}
右侧第一项给出\eref{1.2}的标准符号，第二项则与 $(m,\allowbreak n)$ 无关。结合前面的论证，便建立结合性方程。

为进行所需计算，考虑\eref{12.10}，取 $\RG_1=\RG^{d+2-m}$、$\RG_2=\RG^{m+1}$，但通过对应于 $T(m,\allowbreak n)$ 的映射\eref{9.13}，将 $\RG=(-1;0)\times\RG_1\times\RG_2$ 换成 $\RG^{d+1}$。相应地，令 $\M_{\mathcal S_1}=\M^{d+2-m}$、$\M_{\mathcal S_2}=\M^{m+1}$，并将 $\M_{\mathcal S}$ 换成 $\M^{d+1}$。图的每个角上都有一维实向量空间，按如下方式定向：
\begin{itemize}
\item 右上角取 $\tilde c$ 与 $\RG^{d+1}$ 的定向，再用水平箭头将之转移为 $\lambda^{\mathrm{top}}T\M^{d+1}$ 的元素。按先前讨论，这正是 $\tilde c$ 诱导的模空间定向。
\item 左下角的向量空间与 $\R$ 典范识别，用水平箭头将其转移为右下角的一个平凡化。由 $c_{r_k,\allowbreak u_k}$ 的定义，这等价于将 $\RG^{d+2-m}$ 和 $\RG^{m+1}$ 的定向与 $c_{r_1,\allowbreak u_1}\otimes c_{r_2,\allowbreak u_2}$ 结合。
\end{itemize}
\SourcePage{194}
接着考察这些定向是否与图中的竖直箭头相容。
\begin{itemize}
\item 取左下角模空间的平凡定向，沿作为粘合映射线性化的向上箭头映射，就得到 $\M^{d+1}$ 指向 Gromov 边界的定向。依定义，它与上述定向相差 $(-1)^{\S}$。
\item 右侧一列两个定向的唯一差别，来自相关族的底空间之间的粘合映射\eref{9.13}，所以等于 $(-1)^\ddagger$。
\end{itemize}
最后，整张图在相差已知 Koszul 符号 $(-1)^\diamond$ 的意义下交换，此时
\[
\diamond=m(d-m-1)+m\sum_{k>n+m}i(y_k).
\]
故 $\S\equiv\diamond+\ddagger\pmod2$。直接计算便得到\eref{12.25}。
\end{proof}
\begin{note}\label{st:12.4}
构造中有两项人为选择：$\RG^{d+1}$ 的定向，更确切地，是边界层应给出定向差\eref{12.22}的要求，它在相差无关紧要的整体符号改变的意义下确定这些定向；以及\eref{12.24}。调整这些项，可以得到目前其他任一种 $\Ai$ 结构的符号约定；几何本身并不挑出优先约定。\stref{命题}{12.3}的另一证明见\bcite{60}第 22 节。
\end{note}

\subsection{平移运算与\texorpdfstring{$\Ai$}{A∞} 结构}\label{sec:12h}
现在证明，逐点由\eref{11.36}定义的正合 Lagrangian 膜平移运算，确实对应正确的代数概念。更明确地，它表示相应的 $\Ai$ 模，该模在\secref{3d}中作为\eref{3.2}的特例定义。这一事实在上同调层次的推论，已在 Floer 上同调的讨论中不加证明地陈述，见\eref{12.12}与\eref{12.13}。

考虑一对正合 Lagrangian 膜 $L_0^\#,\allowbreak L_1^\#$，以及后者的平移 $SL_1^\#$。取 $y\in\mathcal C(L_0,\allowbreak L_1)$。一个有用的记号区分，是用 $\tilde y$ 表示与膜对 $(L_0^\#,\allowbreak SL_1^\#)$ 关联的同一条路径 $y$。这意味着除了 $i(y),\allowbreak o(y)$，还有指标与定向空间 $i(\tilde y),\allowbreak o(\tilde y)$。\stref{引理}{11.21}说明 $i(\tilde y)=i(y)-1$，且有典范同构
\begin{equation}\label{eq:12.26}
\sigma_y:o(y)\longrightarrow o(\tilde y).
\end{equation}
回顾该引理的证明，以指标理论表述，\eref{12.26}来自以下事实：定向算子 $D_{H,\allowbreak \tilde y}$ 是 $D_{H,\allowbreak y}$ 与圆盘上指标为 $n-1$ 的算子 $D_{D,\allowbreak \rho}$ 按此顺序的边界连通和。事实上，最好将 $D_{D,\allowbreak \rho}$ 的定义域限制为在 $\partial D$ 的某个固定点消失的函数空间；这使指标降为 $-1$，并同时去掉粘合公式\eref{11.11}中的额外因子 $\lambda^{\mathrm{top}}(F)$。
\SourcePage{195}
基本断言是：若 $u\in\M_Z^*(y_0,\allowbreak y_1)$ 为孤立 Floer 轨线，$\tilde u$ 为视作 $\M_Z^*(\tilde y_0,\allowbreak \tilde y_1)$ 元素的同一映射，则相应同构\eref{12.17}与\eref{12.26}交换，即
\begin{equation}\label{eq:12.27}
\sigma_{y_0}\circ c_u=c_{\tilde u}\circ\sigma_{y_1}.
\end{equation}
两侧都来自将 $D_{H,\allowbreak \tilde y_0}$ 视为由 $D_{Z,\allowbreak u}$（$u$ 的线性化算子）、$D_{H,\allowbreak y_1}$ 与 $D_{D,\allowbreak \rho}$ 粘合而成。两次出现的次序相同，故\eref{12.27}由\eref{11.9}与\eref{11.11}的结合性得到。使用\eref{12.26}的 $\K$ 规范化，再乘以 $(-1)^{i(y)}$，作识别
\begin{equation}\label{eq:12.28}
CF^*(L_0^\#,SL_1^\#)\cong CF^{*+1}(L_0^\#,L_1^\#).
\end{equation}
依定义，它与 $\partial$ 反交换，因而与 $\mu^1$ 交换。\znote{187}{这一小节把“平移”从几何操作提升为代数结构：$S$ 作用在膜上，(12.28) 给出链层面的同构，而 $(-1)^{i(y)}$ 与 $\star=i(y_1)+\cdots+i(y_{d-1})$ 正是使 $\mu^d$ 与之相容所需的修正。用第 3d 节的语言说，$S$ 给每个对象配一个平移对象，(12.28) 及其与 $\mu^d$ 的相容性正是 $\Ai$ 范畴中平移对象应满足的公理；12e 的 (12.12) 只是它在上同调层的影子。若把 $S$ 看成 $\Ai$ 范畴的自等价，则 (12.28) 就是它到一个平移函子的自然同构在链层面的实现。[来源：Seidel, Fukaya categories and Picard–Lefschetz theory, §3d（$\Ai$ 模的定义，(3.2)）与 §12h。]}

再考虑 $(d+1)$ 元膜组 $(L_0^\#,\allowbreak \ldots,\allowbreak L_d^\#)$，其中 $d\geq2$，以及孤立点 $u\in\M^{d+1}(y_0,\allowbreak \ldots,\allowbreak y_d)$。要将 $L_d^\#$ 平移为 $SL_d^\#$，并如前引入记号 $\tilde y_0,\allowbreak \tilde y_d,\allowbreak \tilde u$。对应于\eref{12.27}的陈述是：定向空间之间的多重线性映射 $c_u,\allowbreak c_{\tilde u}$ 满足
\[
\sigma_{y_0}\circ c_u=(-1)^\star c_{\tilde u}\circ(\sigma_{y_d}\otimes\id\otimes\cdots\otimes\id),
\]
其中 $\star=i(y_1)+\cdots+i(y_{d-1})$。论证与前面相同，只是比较两侧时，在复合次序中必须把 $D_{D,\allowbreak \rho}$ 移过 $D_{H,\allowbreak y_1},\allowbreak \ldots,\allowbreak D_{H,\allowbreak y_{d-1}}$，从而产生 Koszul 型符号。因此 $\mu^d$ 与\eref{12.28}交换，这正是一般 $\Ai$ 范畴中平移对象的定义。为看清这一点，注意除 $\star$ 之外，相关符号量还有：由\eref{12.28}对 $y_d,\allowbreak y_0$ 的定义得到的 $i(y_d)$ 与 $i(y_0)=i(y_1)+\cdots+i(y_d)+2-d$，以及来自\eref{12.24}的 $i(y_1)+\cdots+di(y_d)$ 和 $i(y_1)+\cdots+di(\tilde y_d)$。最终所有符号抵消，得到所需结论。

\clearpage
\subsection{自同构}\label{sec:12i}
\sr{sec:10}{第~10~节}关于 $\F(M)^{\mathrm{pr}}$ 的全部观察都可推广到 $\F(M)$。当然，有些陈述须略作改变，以容纳膜结构，但通常很直接。作为实际演示，考察一个较有吸引力的例子：辛自同构群在 Fukaya 范畴上的作用，最初在\secref{10b}讨论。

先说一些一般事实。取任意辛流形 $(M,\allowbreak \omega)$，配有殆复结构 $I_M$ 和二次复体积形式 $\eta_M^2$。
\begin{definition*}
$M$ 的\emph{分次自同构}是二元组 $(\phi,\allowbreak \gamma^\#)$，由辛自同构 $\phi$ 和映射 $\gamma^\#:\operatorname{Gr}(TM)\to\R$ 组成，对所有 $\Lambda\in\operatorname{Gr}(TM)$ 满足
\[
\exp(2\pi i\gamma^\#(\Lambda))=\frac{\alpha_M(D\phi(\Lambda))}{\alpha_M(\Lambda)}.
\]
\end{definition*}
\SourcePage{196}
这种自同构的复合定义为 $\phi_2^\#\circ\phi_1^\#=(\phi_2\circ\phi_1,\allowbreak \gamma_2^\#\circ D\phi_1+\gamma_1^\#)$。分次自同构自然作用于 Lagrangian 膜的集合：若 $(L,\allowbreak \alpha^\#,\allowbreak P^\#)$ 为膜，则像 $\phi(L)$ 继承分次 $x\mapsto\gamma^\#(TL_{\phi^{-1}(x)})+\alpha^\#(\phi^{-1}(x))$，以及推前 Pin 结构 $(\phi^{-1})^*P^\#$。这还不是最一般的情形。
\begin{definition*}
\emph{膜自同构}是三元组 $\phi^\#=(\phi,\allowbreak \gamma^\#,\allowbreak \beta)$，其中 $(\phi,\allowbreak \gamma^\#)$ 如前，$\beta\to M$ 为实线丛。
\end{definition*}
这种自同构在膜上的作用还包括 Pin 结构的扭曲：
\begin{equation}\label{eq:12.29}
P^\#\longmapsto(\phi^{-1})^*(P^\#\otimes\beta|L).
\end{equation}
一个相关例子是平移运算 $S$；由\stref{笔记}{11.18}，它可写为 $(\Id_M,\allowbreak \gamma^\#=-1,\allowbreak \beta=\xi_M)$ 的作用。

以下假设 $M$ 为带角正合辛流形，考虑由底层 $\phi$ 属于 $\operatorname{Aut}(M,\allowbreak \partial M)$ 的 $(\phi,\allowbreak \gamma^\#)$ 组成的群 $\operatorname{Aut}^{\mathrm{gr}}(M,\allowbreak \partial M)$。赋予它从 $\operatorname{Aut}(M,\allowbreak \partial M)\times C^\infty(\operatorname{Gr}(TM),\allowbreak \R)$ 诱导的拓扑。$\operatorname{Aut}^{\mathrm{gr}}(M,\allowbreak \partial M)$ 在 $H\F(M)$ 上有典范严格作用，并下降为 $\pi_0(\operatorname{Aut}^{\mathrm{gr}}(M,\allowbreak \partial M))$ 的弱作用。对 Fukaya 范畴本身，使用小节 \sr{sec:10b}{(10b)}、\sr{sec:10c}{(10c)} 的同一构造，得到 $\operatorname{Aut}^{\mathrm{gr}}(M,\allowbreak \partial M)$ 在 $\F(M)$ 上的协调作用，同样下降为 $\pi_0(\operatorname{Aut}^{\mathrm{gr}}(M,\allowbreak \partial M))$ 的弱作用。现在加入一般膜自同构，但不区分 $M$ 上同构的线丛。这定义 $\operatorname{Aut}^{\mathrm{gr}}(M,\allowbreak \partial M)$ 关于 $H^1(M;\Z/2)$ 的扩张，记为 $\operatorname{Aut}^\#(M,\allowbreak \partial M)$。该群在 Fukaya 范畴、甚至 Donaldson–Fukaya 范畴上的自然作用，一般都只有弱作用，因为\eref{12.29}要求在同构类 $[\beta]$ 中选取代表 $\beta$。这一现象在实际应用中并不十分重要，不过后面讨论对合时，细心的读者仍会看到它的痕迹，见\sr{sec:14}{第~14~节}。\znote{188}{“弱作用”指群作用只在同构意义下成立：公理 $g(hx)\cong(gh)x$、$ex\cong x$ 只要求自然同构而非等式。膜自同构多出的线丛 $\beta$ 使 (12.29) 必须先选代表，故作用只能弱；把线丛自身的对称群 $H^0(M;\Z/2)$ 一并记入，$\operatorname{Aut}^\#(M,\partial M)$ 就成为弱 $2$ 群——群公理被范畴化，元素之间的同构本身也成为数据。也正因如此，只有 $\pi_0$ 才是通常意义上的群。笔记 12.5 的计算说明，平凡线丛的自同构 $-\Id_\beta$ 给出的正是自然变换 $-E$。[来源：Seidel, Fukaya categories and Picard–Lefschetz theory, §10b–10c 与笔记 12.5。]}
\begin{note}\label{st:12.5}
为说明代表 $\beta$ 的选择确有影响，考虑三个正合 Lagrangian 膜 $L_k^\#=(L_k,\allowbreak \alpha_k^\#,\allowbreak P_k^\#)$（$k=0,\allowbreak 1,\allowbreak 2$），其中后两个同构：具有相同的 Lagrangian 子流形及分次，Pin 结构同构但不一定严格相等。Pin 同构 $g:P_1^\#\to P_2^\#$ 的选择确定 Floer 上同调群同构
\begin{equation}\label{eq:12.30}
HF^*(L_0^\#,L_1^\#)\cong HF^*(L_0^\#,L_2^\#).
\end{equation}
为简单起见，假定 $L_0$ 与 $L_1=L_2$ 横截相交，这样相关 Floer 数据中的 Hamilton 项可取为零。固定交点 $y$，令 $\Lambda_k^\#=\Lambda_{k,\allowbreak y}^\#$（$k=0,\allowbreak 1,\allowbreak 2$）为由 $L_k^\#$ 在该点得到的线性膜。

选择处于正确同伦类中的路径 $\lambda\in\Omega^-(\operatorname{Gr}(TM_y),\allowbreak \Lambda_0,\allowbreak \Lambda_1)$，以及相应向量丛 $F_\lambda\to[0;1]$ 上的 Pin 结构 $P_\lambda^\#$。可将
\SourcePage{197}
$(\delta_H)_\lambda$ 同时与定向空间 $o(\Lambda_0^\#,\allowbreak \Lambda_1^\#)$、$o(\Lambda_0^\#,\allowbreak \Lambda_2^\#)$ 识别，见\stref{引理}{11.11}。但这种识别依赖辅助数据的选择，即同构 $f_0:(P_\lambda^\#)_0\to P_0^\#$ 和 $f_k:(P_\lambda^\#)_1\to P_k^\#$、$k=1,\allowbreak 2$。为得到\eref{12.30}底层的链映射，应取 $f_2=g_y\circ f_1$，并考虑诱导的定向空间映射。反转 $g$ 的符号会改变 $f_2$ 的符号，因而按\stref{笔记}{11.19}的讨论，改变识别 $(\delta_H)_\lambda\cong o(\Lambda_0^\#,\allowbreak \Lambda_2^\#)$，最终将\eref{12.30}乘以 $-1$。
\end{note}

改善这一情形的一种办法，是把 $\operatorname{Aut}^\#(M,\allowbreak \partial M)$ 看作更精细的对象——“弱 $2$ 群”，将线丛的内部对称群 $H^0(M;\Z/2)$ 也纳入考虑。这与一般范畴论十分吻合，因为范畴论本就认为“一个范畴的自同构群”是这类对象。用这一语言，\stref{笔记}{12.5}中的计算说明：平凡线丛 $\beta=M\times\R$ 的自同构 $-\Id_\beta$ 给出恒等函子到自身的自然变换 $-E$。

\subsection{补充说明}\label{sec:12j}
Fukaya 范畴还有许多其他性质，这里无法讨论。不过，其中两项值得作最基本的介绍，至少因为它们在上同调层次的对应物已在 Floer 理论的讨论中出现。

第一项是 PSS 同构 $HF^*(L^\#,\allowbreak L^\#)\cong H^*(L;\K)$。相应的链层次陈述说：若 $L^\#$ 是单个正合 Lagrangian 膜，则只考察这个对象得到的 $\Ai$ 代数 $\B=\operatorname{hom}_{\F(M)}(L^\#,\allowbreak L^\#)$，拟同构于用 $L$ 的经典拓扑所得到的代数，例如奇异上链的微分分次代数，或等价地，使用 Morse 同调及 Fukaya 的梯度树运算\bcite{54}得到的代数。该陈述的早期形式见\bcite{58}，其中通过构造模空间的同构，基本上旨在得到相关代数结构的相等。另一种办法是在\bcite{33}及其前期工作的方式下，混合使用伪全纯多边形与梯度树；它的分析技术较不棘手，但只能证明拟同构。还有允许扭曲系数的小变体，其中包括\eref{12.14}的链层次版本。

第二项是 Donaldson–Fukaya 范畴的 Frobenius 性质的对应物，所以假设 $\xi_M$ 平凡。对任意 $\Ai$ 范畴 $\A$，有典范 $\Ai$ 函子 $\mathcal D:\A\to\Q=\modcat(\A)$，称为\emph{抽象 Serre 函子}。对象 $Y$ 在 $\mathcal D$ 下的像为对偶 $\Ai$ 模 $\Y^\vee$，满足 $\Y^\vee(X)=\operatorname{hom}_{\A}(Y,\allowbreak X)^\vee$，结构映射也有类似公式。\znote{189}{Serre 函子的定义性性质是自然同构 $\operatorname{hom}_{\A}(A,B)\cong\operatorname{hom}_{\A}(B,SA)^\vee$；在凝聚层的导出范畴中它由 $(\cdot\otimes\omega_X)[n]$ 实现，因此“$n$ 维弱 Calabi–Yau”就是把 Serre 函子识别为平移 $[n]$，即 $\mathcal D\cong\mathcal I[n]$。这里 $\mathcal D$ 把对象 $Y$ 送到对偶模 $Y^\vee$，正是把 12e 中带 $\lambda^{\mathrm{top}}(TL)$ 扭转的 Poincaré 对偶配对搬到 $\Ai$ 层面，所以弱 CY 性质可读作“Frobenius 性质差一个 $[n]$”。把 $\A$ 换成 $\F(M)$ 即得 Donaldson–Fukaya 范畴，而 $\mathcal D$ 与 Yoneda 嵌入 $\mathcal I$ 的差别正是那个扭转在模层面的体现，也是 12j 全节的核心。[来源：nLab “Serre functor”，https:/\allowbreak /\allowbreak ncatlab.org/\allowbreak nlab/\allowbreak show/\allowbreak Serre+\allowbreak functor（2026-11-15 访问）；并见原书 §12j。]}
\begin{definition*}
若在 $H(\operatorname{fun}(\A,\allowbreak \Q))$ 中，存在从 Yoneda 嵌入 $\mathcal I$ 到 $\mathcal D$ 的次数为 $-n$ 的
\SourcePage{198}
同构，则称 $\A$ 为 $n$ 维\emph{弱 Calabi–Yau 范畴}。
\end{definition*}
自然变换 $\mathcal I\to\mathcal D$ 自然来自对偶 Hochschild 上同调 $HH(\A,\allowbreak \A^\vee)$ 中的类。在我们的情形 $\A=\F(M)$ 中，相关类通过考察只有出向无穷远点、并额外带一个内部标记点的带标记点圆盘来定义，方式如\secref{10c}。随后由\stref{引理}{1.6}及 Floer 上同调中的 Poincaré 对偶，得到相应自然变换为同构。

\section{曲面上的多边形}\label{sec:13}\label{section:13}
\subsection{自动正则性}\label{sec:13a}
设 $S$ 为带条带状端的连通 Riemann 曲面。考虑 $S$ 上所有在某个紧子集之外与给定 $j_0=I_S$ 一致的复结构 $j$ 所成的空间 $\widetilde{\RG}$。乘积 $\widetilde{\mathcal S}=\widetilde{\RG}\times S\to\widetilde{\RG}$ 的纤维上有自然的复结构，因此成为带条带状端的 Riemann 曲面族。此外，该族上有群 $\widetilde{\mathcal G}$ 的自然作用，该群由 $S$ 的如下微分同胚组成：在每个条带状端的无穷远附近，它们限制为 $s$ 方向的平移。

为 $S$，因而也为 $\widetilde{\mathcal S}$，赋予 Lagrangian 标记，使每个无穷远点 $\zeta$ 都满足 $L_{\zeta,\allowbreak 0}\pitchfork L_{\zeta,\allowbreak 1}$。选择平凡扰动数据 $K=0$、$J=I_M$，即殆复结构的常值族。取伪全纯映射 $u\in\M_S(\{y_\zeta\})$，将 $(j_0,\allowbreak u)$ 视为相应参数化模空间 $\M_{\widetilde{\mathcal S}}(\{y_\zeta\})$ 的元素。$\widetilde{\mathcal G}$ 的无穷小作用定义从 Lie 代数 $L\widetilde{\mathcal G}$ 到 Zariski 切空间 $T^{\mathrm{Zar}}\M_{\widetilde{\mathcal S}}(\{y_\zeta\})_{(j_0,\allowbreak u)}$ 的线性映射 $d^0$；使用“Zariski 切空间”这一说法，是因为尚未作正则性假设。另一方面，依定义该 Zariski 切空间是扩展线性化算子 $d^1=D_{\widetilde{\mathcal S},\allowbreak j_0,\allowbreak u}$ 的核。因此可将两个算子视为如下复形的边界映射：
\begin{equation}\label{eq:13.1}
0\longrightarrow L\widetilde{\mathcal G}\xrightarrow{d^0}T\widetilde{\RG}_{j_0}\times(T\mathcal B_S)_u
\xrightarrow{d^1}(\mathcal E_S)_u\longrightarrow0.
\end{equation}
到目前为止，对无穷维空间 $\widetilde{\RG}$、$\widetilde{\mathcal G}$ 的处理相当非正式；现在至少在切空间层面，希望通过选择合适的 Hilbert 空间完备化使其精确。按先前定义，$\widetilde{\mathcal G}$ 的 Lie 代数由 $S$ 上与边界相切的光滑向量场 $Z$ 组成，它们在每个条带状端的无穷远附近满足 $\epsilon_\zeta^*Z=c_\zeta\partial_s$，其中 $c_\zeta\in\R$ 为常数。现在改取 $(k,\allowbreak p)=(2,\allowbreak 2)$、权重为充分小 $\epsilon>0$ 的加权 Sobolev 完备化，即\eref{9.2}记号中的 $W^{2,\allowbreak 2;\epsilon}(S,\allowbreak TS,\allowbreak T(\partial S))^+$。同样，朴素地说，$\widetilde{\RG}$ 的切空间是与复结构反交换的 $TS$ 的紧支自同态空间。将这些自同态解释为复张量积线丛 $\Omega_S^{0,\allowbreak 1}\otimes TS$ 的光滑截面，并完备化为
\SourcePage{199}
$W^{1,\allowbreak 2;\epsilon}(\Omega_S^{0,\allowbreak 1}\otimes TS)$。最后，对作为 $d^1$ 第二个分量的线性化算子 $D_{S,\allowbreak u}$，使用相同的 Sobolev 拓扑。\eref{13.1}的完备化于是为
\begin{equation}\label{eq:13.2}
\begin{gathered}
W^{2,2;\epsilon}(S,TS,T(\partial S))^+\\
\big\downarrow d^0\\
W^{1,2;\epsilon}(S,\Omega_S^{0,1}\otimes TS)\oplus W^{1,2;\epsilon}(S,E,F)\\
\big\downarrow d^1\\
L^{2;\epsilon}(S,\Omega_S^{0,1}\otimes E),
\end{gathered}
\end{equation}
其中 $E=u^*TM$，$F\subset E|\partial S$ 由相关 Lagrangian 子流形的切空间给出，且
\[
d^0(Z)=(2I_S\circ\bar\partial Z,Du\circ Z),\qquad
d^1(Y,X)=\tfrac12I_M\circ Du\circ Y+D_{S,u}(X).
\]
由于 Sobolev 空间的选择是为当前问题临时作出的，首先应验证这些算子确实有界，不过这很直接。事实上，还能得到更多结论：$TS$ 上边界条件为 $T(\partial S)$ 的标准 Cauchy–Riemann 算子 $\bar\partial$ 容许，故在加权 Sobolev 空间上成为 Fredholm 算子，见\secref{9b}。$D_{S,\allowbreak u}$ 本身是非退化 $\bar\partial_\nabla$ 算子，因此也为 Fredholm 算子。其余分量为零阶项，因而紧。结论是\eref{13.2}为 Fredholm 复形。

选择曲面族 $\mathcal S\to\RG$，这次底空间有限维，使它成为\secref{9b}意义下 $S=\mathcal S_{r_0}$ 的无穷小极小通用变形。令 $\mathcal G$ 为 $S$ 的全纯自同构群。
\begin{lemma}\label{st:13.1}
\eref{13.2}的上同调群具有如下解释：$H^0,\allowbreak H^1$ 分别为映射 $L\mathcal G\to\ker(D_{\mathcal S,\allowbreak r_0,\allowbreak u})$ 的核与余核；$H^2=\operatorname{coker}(D_{\mathcal S,\allowbreak r_0,\allowbreak u})$。
\end{lemma}
\begin{proof}
非正式地，想法是 $\RG$ 为 $\widetilde{\mathcal G}$ 在 $\widetilde{\RG}$ 上作用的局部切片。为使论证严格，选择在条带状端上恒定的 $\mathcal S$ 的局部平凡化。这给出 $\RG$ 到 $\widetilde{\RG}$ 的局部嵌入，因而给出
\SourcePage{200}
切空间之间的映射。于是有各列正合的交换图
\SourceEquationDiagram{diagram-13-3}{0.36969191}{13.3}
依定义，$T\RG_{r_0}$ 同构地投影到\eref{9.3}的余核，因此\eref{13.3}的底行是满足 $H^0=L\mathcal G$、$H^1=0$ 的复形。上同调长正合列给出所需结论。
\end{proof}
\begin{lemma}\label{st:13.2}
若 $M$ 为二维且 $u$ 非常值，则对任意无穷小极小通用族 $\mathcal S\to\RG$，$D_{\mathcal S,\allowbreak r_0,\allowbreak u}$ 均为满射。
\end{lemma}
\begin{proof}
$u$ 为 Riemann 曲面之间的非常值全纯映射，所以除有限多个点外，$Du:TS_z\to TM_{u(z)}$ 都是同构。因此 $d^1$ 的第一个分量具有稠密像。由于第二个分量为 Fredholm 算子，便知 $d^1$ 满射。再运用\stref{引理}{13.1}最后一条即可。
\end{proof}
称\stref{引理}{13.2}为“自动正则性”结论，因为它对任意 $I_M$ 成立，且不需要扰动伪全纯映射方程。\znote{192}{自动正则性的关键在 $\dim M=2$：此时 $S$ 与 $M$ 的切空间都是复一维，$u$ 非常值全纯就使 $Du$ 除有限点外为同构，于是 $d^1$ 的第一个分量已提供足够多的方向，第二个 Fredholm 分量再补足余下部分，满射性无需额外假设。若 $M$ 维数更高，$Du$ 只是单射，$d^1$ 的像一般并不稠密，必须靠扰动才能保证正则。这也解释了为什么组合 Floer 理论只在曲面情形如此干净。[来源：Seidel, Fukaya categories and Picard–Lefschetz theory, §13a 与引理 13.2。]}
\begin{note}\label{st:13.3}
偏好代数几何的读者也许会觉得下面的层论解释有帮助。为简单起见，转到类似的情形：全纯映射 $u:S\to M$，其中 $S$ 是闭 Riemann 曲面，$M$ 是复流形。此时\eref{13.2}变为
\[
\begin{gathered}
W^{2,2}(S,TS)\longrightarrow
W^{1,2}(S,\Omega_S^{0,1}\otimes TS)\oplus W^{1,2}(S,u^*TM)\\
\longrightarrow L^2(S,\Omega_S^{0,1}\otimes u^*TM).
\end{gathered}
\]
它实际上计算凝聚层的两项复形 $\{Du:TS\to u^*TM\}$ 的超上同调。若 $u$ 非常值，则 $Du$ 作为层映射为单射，所以实际上计算的是商层 $u^*TM/TS$ 的上同调，再平移一阶。最后，若 $M$ 是曲面，则该商为挠层，因此一次上同调消失。
\end{note}

\subsection{组合 Floer 上同调}\label{sec:13b}
设 $M$ 为二维带角正合辛流形。$M$ 中连通正合 Lagrangian 子流形是
\SourcePage{201}
满足 $\int_L\theta=0$ 的简单闭曲线 $S^1\cong L\subset\operatorname{int}(M)$，其中 $\theta\in\Omega^1(M)$ 为辛形式 $\omega$ 的指定原始形式。若曲线有多个连通分支，则每个分支上的积分均须为零。注意，由 Stokes 定理，这排除了零同调的曲线。固定一族相互处于一般位置的正合 Lagrangian 子流形 $\symbfit{L}=(L_1,\allowbreak \ldots,\allowbreak L_m)$，即任意两条横截相交，且无三重交点。定义\sr{sec:9}{第~9~节}初步版本的、不分次也不带符号的 Fukaya 范畴 $\F(M)^{\mathrm{pr}}$ 时，可以将 Floer 数据
\begin{equation}\label{eq:13.4}
(H_{L_{i_0},L_{i_1}},J_{L_{i_0},L_{i_1}}),\qquad i_0<i_1,
\end{equation}
取为平凡数据 $H=0,\allowbreak J=I_M$，扰动数据
\begin{equation}\label{eq:13.5}
(K_{L_{i_0},\ldots,L_{i_d}},J_{L_{i_0},\ldots,L_{i_d}}),\qquad i_0<\cdots<i_d,
\end{equation}
也可如此选择。验证这一做法可行涉及两方面。第一方面由\stref{引理}{13.2}解决，它保证对应于\eref{13.4}和\eref{13.5}的 Floer 轨线模空间及全纯多边形模空间分别正则。这些模空间没有常值映射，因为边界上的 Lagrangian 标记 $L_{i_k}$ 各不相同且处于一般位置。第二方面是\eref{13.5}与万有扰动数据选择的一致性条件相容，见\secref{9i}。这是以下初等观察的推论：取稳定 $d$ 叶树 $T$，并对每个顶点取带 Lagrangian 标记和条带状端的 $|v|$ 点圆盘 $S_v$。假设标记相容，使各 $S_v$ 可粘合为带 Lagrangian 标记的 $(d+1)$ 点圆盘 $S$。若 $S$ 的标记为 $(L_{i_0},\allowbreak \ldots,\allowbreak L_{i_d})$，且 $i_0<\cdots<i_d$，则每个组成部分 $S_v$ 的标记均为其有序子集，因而具有同样形式。

接着，假定\eref{13.4}与\eref{13.5}均取平凡数据。断言：\secref{5n}意义下的有向 $\Ai$ 子范畴 $\F(\symbfit{L})^{\mathrm{pr},\allowbreak \to}\subset\F(M)^{\mathrm{pr}}$ 与复结构 $I_M$ 无关。换言之，固定 $i_0<\cdots<i_d$ 及交点 $y_0\in L_{i_0}\cap L_{i_d}$、$y_k\in L_{i_{k-1}}\cap L_{i_k}$（$1\leq k\leq d$），则 $\M^{d+1}(y_0,\allowbreak \ldots,\allowbreak y_d)$ 中孤立点的个数在 $I_M$ 的同伦下不变。这由模空间的直接配边论证证明，也可视为\secref{10e}讨论的平凡情形：由\stref{引理}{13.2}，参数化模空间 $\M^{d+1}(y_0,\allowbreak \ldots,\allowbreak y_d)^{[0;1]}$ 光滑纤维化于 $[0;1]$ 上，故不含孤立点，因此\stref{命题}{10.13}的 $\Ai$ 函子为恒等函子。

既已建立 $\F(\symbfit{L})^{\mathrm{pr},\allowbreak \to}$ 原则上具有拓扑性质，下一步是得到显式公式。这方面已有充分研究，见本章开头的参考文献，因此这里只作简述。令 $\Delta^{d+1}\subset\R^2$ 为有定向的凸 $(d+1)$ 边形；$d=1$ 时为二边形，或半月形。其角点和边的编号与 $(d+1)$ 点圆盘相同。对前述点 $y_0,\allowbreak \ldots,\allowbreak y_d$，
\SourcePage{202}
考虑保持正定向的光滑浸入 $\Delta^{d+1}\to M$ 所成的集合 $\operatorname{Imm}^{d+1}(y_0,\allowbreak \ldots,\allowbreak y_d)$，要求第 $k$ 个角点映到 $y_k$，第 $k$ 条边映到 $L_{i_k}$，见\sr{fig:13.1}{图~13.1}。坚持要求映射直到边界都是浸入，
\SourceFigure{figure-13-1}{0.40084572}{13.1}
即每个角点附近的像是顶角小于 $\pi$ 的楔形区域。固定各角点的 $\Delta^{d+1}$ 微分同胚群作用于 $\operatorname{Imm}^{d+1}(y_0,\allowbreak \ldots,\allowbreak y_d)$，其商可与 $\M^{d+1}(y_0,\allowbreak \ldots,\allowbreak y_d)$ 的零维部分识别。证明如下。为方便，假定选择 $I_M$，使每个两曲线交点附近都有局部全纯坐标，在其中两曲线成为 $\CC$ 中的实轴和虚轴。这没有问题，因为已知 $\F(\symbfit{L})^{\mathrm{pr},\allowbreak \to}$ 与 $I_M$ 无关；也可处理任意 $I_M$，代价是多用一些分析，如\bcite{114}所示。指标论证\bcite{43}证明，每个孤立 $u\in\M^{d+1}(y_0,\allowbreak \ldots,\allowbreak y_d)$ 都是浸入，且其像在无穷远点附近呈楔形。再稍加考察 $u$ 在条带状端上的收敛速率，便知可将其紧化为浸入 $\tilde u:\Delta^{d+1}\to M$。反方向从浸入出发构造全纯映射，依赖以下事实：在 $\Delta^{d+1}$ 上给定任意复结构后，去掉角点，剩下的就是 $(d+1)$ 个边界穿孔的圆盘；这很容易归结为闭圆盘的单值化定理。\znote{194}{这一步把分析问题化成组合问题：$\Delta^{d+1}$ 上的任意复结构在去掉角点后给出边界穿孔的圆盘，而单值化定理保证它与标准圆盘共形等价，边界点按顺序对应。于是多边形的复结构不带来额外模参数，零维模空间的点与浸入多边形及其角点重参数化一一对应，$\mu^d$ 因而可以只靠数多边形来计算。13c 能用纯组合数据定义 $\tilde\mu^d$，依据正在于此。[来源：Seidel, Fukaya categories and Picard–Lefschetz theory, §13b（原文所引单值化定理）。]}
\begin{note}\label{st:13.4}
自然会问，由一般位置的一族 $\symbfit{L}$ 组成的 $\F(M)^{\mathrm{pr}}$ 的满 $\Ai$ 子范畴，是否也有组合描述。看来很可能有，但答案应比这里给出的复杂得多，
\SourcePage{203}
因为交点处的常值全纯多边形也有贡献，而此时\stref{引理}{13.2}失效：其模空间的维数过大，虽光滑，却不正则。\znote{195}{常值映射是唯一的例外：$u\equiv y$ 时 $Du=0$，引理 13.2 证明中“$Du$ 除有限点外为同构”完全失效，$D_{S,u}$ 的满射性没有保障。几何上这些点对应退化多边形：所有角点都落在同一交点附近，多边形缩成一点。它们在一般位置下仍对 $\mu^d$ 有贡献，所以 $\F(M)^{\mathrm{pr}}$ 的满 $\Ai$ 子范畴不能简单取组合模型，这正是 13b 只处理有向子范畴、而不直接采用组合定义的原因。[来源：Seidel, Fukaya categories and Picard–Lefschetz theory, §13b，Remark 13.4 与引理 13.2。]}
\end{note}

\subsection{膜结构}\label{sec:13c}
现在假设 $M$ 带有二次复体积形式；在这一维数中，它就是处处非零的 $C^\infty$ 二次微分 $\eta_M^2:TM\otimes_{\CC} TM\to\CC$。相应平方相位映射可视为实射影化切丛的平凡化 $\alpha_M:\RP(TM)=\operatorname{Gr}(TM)\to S^1$。一种好的直观图像来自实子丛 $\xi_M=\alpha_M^{-1}(1)$，它给出曲面上的无定向叶状结构。用这一语言，$M$ 中的 Lagrangian 膜就是嵌入闭曲线 $L$，配有
\begin{itemize}
\item 拉回丛 $\pi^*(TM|L)$ 的实子丛 $H$，其中 $\pi:[0;1]\times L\to L$ 是投影，满足 $H|\{0\}\times L=\xi_M|L$、$H|\{1\}\times L=TL$；
\item 实线丛 $\beta\to L$。
\end{itemize}
$H$ 是 $TM|L$ 的子丛同伦；从几何看，就是同时旋转 $L$ 的所有切空间，使之与 $\xi_M$ 平行的一种方式。这种同伦存在的障碍是 $L$ 相对于 $\xi_M$ 的整体旋转数；要求该数为零，再次排除了零同调的简单闭曲线。为与一般分次概念联系，注意：给定这种 $H$ 后，存在唯一函数 $a:[0;1]\times L\to\R$，满足 $a(0,\allowbreak x)=0$ 与 $\exp(2\pi ia(r,\allowbreak x))=\alpha_M(H_{r,\allowbreak x})$，于是 $\alpha^\#(x)=a(1,\allowbreak x)$ 是 $L$ 的一个分次。至于 $\beta$，由\stref{例题}{11.20}已知，它等价于 Pin 结构的一个选择。\znote{195}{曲面上的膜结构可以完全用初等数据描述：$H$ 是从 $\xi_M|L$ 到 $TL$ 的子丛同伦，旋转角的累计给出分次（见 (13.6) 前的构造）；$\beta$ 则代替 Pin 结构，因为在实一维情形，$TL$ 上的 Pin 结构恰好等价于一个实线丛（例题 11.20）。这解释了 13c 为何能完全避开 12a 的 $\operatorname{Gr}(TM)^\#$ 语言：二维时 $\alpha_M:\RP(TM)\to S^1$ 已是射影化切丛的平凡化，$\xi_M$ 就是它的水平叶状结构。由 $H$ 与 $\beta$ 还能看出，膜的自同构群恰是保持这两项数据的复形，这与 12i 中的 $\operatorname{Aut}^\#$ 相呼应。[来源：Seidel, Fukaya categories and Picard–Lefschetz theory, §13c 与例题 11.20。]}

令 $\symbfit{L}^\#$ 为处于一般位置的一族正合 Lagrangian 膜 $L_k^\#=(L_k,\allowbreak H_k,\allowbreak \beta_k)$，$1\leq k\leq m$。现在为先前构造加入分次与符号，得到组合定义的有向 $\Ai$ 范畴。首先，为每个交点 $y\in L_{i_0}\cap L_{i_1}$（$i_0<i_1$）赋予绝对指标 $i(y)$ 和定向空间 $o(y)$。将同伦 $H_{i_0}(1-\cdot,\allowbreak y)$ 与 $H_{i_1}(\cdot,\allowbreak y)$ 组合为 $\RP(TM_y)$ 中从 $(TL_{i_0})_y$ 到 $(TL_{i_1})_y$ 的路径。令 $\measuredangle(y)\in\R\setminus\pi\Z$ 为该路径旋转的总角度，并设
\begin{equation}\label{eq:13.6}
\begin{aligned}
i(y)&=\lfloor\measuredangle(y)/\pi\rfloor+1,\\
o(y)&=\Hom(\beta_{i_0,y},\beta_{i_1,y})\otimes(TL_{i_1})_y^{\otimes i(y)}.
\end{aligned}
\end{equation}
然后如\eref{12.16}定义态射空间 $CF^*(L_{i_0}^\#,\allowbreak L_{i_1}^\#)$。接着，取角点位于 $y_0,\allowbreak \ldots,\allowbreak y_d$ 的浸入多边形 $\tilde u:\Delta^{d+1}\to M$。用\eref{13.6}作简单计算可见，这种多边形的存在蕴含等式 $i(y_0)=i(y_1)+\cdots+i(y_d)+2-d$。用 $\Delta^{d+1}$ 在各角点的边界定向，赋予 $(TL_{i_d})_{y_0}$ 和 $(TL_{i_k})_{y_k}$、$k=1,\allowbreak \ldots,\allowbreak d$ 定向，如\sr{fig:13.1}{图~13.1}中的箭头所示；注意 $y_0$ 处与 $y_1,\allowbreak \ldots,\allowbreak y_d$ 处的约定不同。沿 $\Delta^{d+1}$ 各边的平行移动定义同构 $\beta_{i_k,\allowbreak y_k}\to\beta_{i_k,\allowbreak y_{k+1}}$ 及 $\beta_{i_d,\allowbreak y_d}\to\beta_{i_d,\allowbreak y_0}$。由\eref{13.6}可见，这两个初等事实
\SourcePage{204}
结合给出同构 $o(y_d)\otimes\cdots\otimes o(y_1)\to o(y_0)$，记为 $\tilde c_{\tilde u}$。将全部规范化 $|\tilde c_{\tilde u}|_{\K}$ 相加，定义
\[
\tilde\mu^d:CF^*(L_{i_{d-1}}^\#,L_{i_d}^\#)\otimes\cdots\otimes CF^*(L_{i_0}^\#,L_{i_1}^\#)
\longrightarrow CF^{*+2-d}(L_{i_0}^\#,L_{i_d}^\#).
\]
一个基于具有一个凹角的多边形分解的初等论证表明，这些映射满足 $\Ai$ 结合方程\eref{1.2}。\znote{196}{这里的论证是“带一个凹角的多边形分解”：沿一条弦把多边形切成两块，正好对应把一个余维一边界面拆成两个低阶模空间，凹角保证切开后各部分仍是凸多边形。于是 $d$ 阶关系归结为 $d-1$ 阶与更低阶关系的组合，与 12g 中 Deligne–Mumford 紧化的论证完全平行，只是把“模空间边界”换成了“多边形剖分”。这类只数多边形的范畴是纯组合的，文献中正是用它在带二次微分的曲面上研究稳定性条件。[来源：arXiv:1409.8611，Haiden–Katzarkov–Kontsevich, Flat surfaces and stability structures，摘要（2026-11-15 访问）；并见原书 §13c。]}将由此定义的有向 $\Ai$ 范畴记为 $\widetilde{\F}(\symbfit{L}^\#)^\to$。
\begin{example}\label{st:13.5}
考虑 $M=S^1\times[-1;1]$；\sr{fig:13.2}{图~13.2}把它画成两个竖直边相粘合的正方形，并画出其上的曲线 $L_0,\allowbreak L_1$。使用对应于水平直线叶状结构 $\xi_M$ 的常值复体积形式。选择从 $\xi_M|L_0$ 到 $TL_0$ 的平凡同伦。另一条曲线中，$\xi_M|L_1$ 与 $TL_1$ 的夹角处处小于 $\pi/2$，所以二者之间有唯一的“短”同伦，采用这一同伦。最后取平凡线丛 $\beta_0,\allowbreak \beta_1$。则 $i(y_0)=0$，
\SourceFigure{figure-13-2}{0.35099962}{13.2}
$i(y_1)=1$，$o(y_0)=\R$，且此识别是典范的，因为指标为偶数；另有 $o(y_1)=(TL_1)_{y_1}$。\sr{fig:13.1}{图~13.1}的约定说明，计算二边形 $u$ 对边界算子的贡献时，应使用\sr{fig:13.2}{图~13.2}所示的定向识别 $o(y_1)\cong\R$。另一个二边形 $v$ 要求相反识别，所以二者贡献抵消，得到
\[
HF^0(L_0^\#,L_1^\#)=\K,\qquad HF^1(L_0^\#,L_1^\#)=|(TL_1)_{y_1}|_{\K}.
\]
因为 $L_0,\allowbreak L_1$ 作为正合 Lagrangian 膜同痕，一般理论说明存在典范同构 $HF^*(L_0^\#,\allowbreak L_1^\#)\cong H^*(L_1;\K)$。确实，$H^0(L_1;\K)\cong\K$ 是典范的，而 $H^1(L_1;\K)$ 的自然生成元恰是 $L_1$ 的定向，因此 $H^1(L_1;\K)\cong|(TL_1)_{y_1}|_{\K}$。鼓励读者试验改变膜结构的效果，并将结果与\eref{12.12}、\eref{12.14}比较。
\end{example}

下面将上述组合构造与真正 Fukaya 范畴的有向子范畴 $\F(\symbfit{L}^\#)^\to\subset\F(M)$ 联系。对象显然
\SourcePage{205}
相同，态射空间也相同，因为\eref{13.6}只是\eref{11.35}的改写；如先前所见，零维模空间的点 $(r,\allowbreak u)\in\M^{d+1}(y_0,\allowbreak \ldots,\allowbreak y_d)$ 与浸入 $(d+1)$ 边形 $\tilde u$ 的重参数化等价类之间存在典范双射。唯一差别是：按\secref{12g}的一般方法构造 $\Ai$ 结构映射 $\mu^d$ 时，使用的是通过行列式线粘合论证在分析上定义的同构 $c_{r,\allowbreak u}$，见\eref{12.23}。更确切地，$\mu^d$ 由全部贡献 $|(-1)^\dagger c_{r,\allowbreak u}|_{\K}$ 相加得到，其中 $\dagger$ 如\eref{12.24}。据此，以 $\gamma^d(\tilde u)\in\{\pm1\}$ 表示 $\tilde c_{\tilde u}$ 与 $(-1)^\dagger c_{r,\allowbreak u}$ 之间的符号差。

断言 $\gamma^d(\tilde u)=\gamma^d(k_d,\allowbreak \ldots,\allowbreak k_1)$ 实际只由相关交点的绝对指标 $k_j=i(y_j)$ 决定。首先，符号差只依赖浸入多边形附近几何数据的局部行为；其次，若将 $M$ 换成某个覆盖，取 $L_{i_k}$ 的原像并任取多边形的一个提升，符号差不变。由此可将符号的确定归结到多边形嵌入而非仅浸入的情形。此时可沿各边平凡化线丛 $\beta_{i_k}$，说明 $\gamma^d(\tilde u)$ 对所有线丛选择都相同。剩下的局部拓扑数据只有绝对指标，故断言成立。\znote{197}{$\gamma^d$ 只依赖指标，理由有三层：其一，符号差是局部量，只取决于多边形附近的几何数据；其二，覆盖不改变符号差，故可把浸入提升为嵌入；其三，嵌入后可沿各边平凡化诸 $\beta_{i_k}$，线丛的选取不再影响符号。剩下的只有各角点的绝对指标 $k_j=i(y_j)$。这与 12g 中 $\dagger$、$\ddagger$、$\diamond$ 只依赖指标的现象同源：定向空间带来的符号最终都由 Maslov 指标决定。[来源：Seidel, Fukaya categories and Picard–Lefschetz theory, §13c；并与 §12g 的符号计算对照。]}

还有一项先验信息：$\mu^d$ 与 $\tilde\mu^d$ 均满足 $\Ai$ 结合方程。这对 $\gamma^d$ 施加相当强的限制。例如，考察\sr{fig:13.3}{图~13.3}(i) 所示的分解，可见
\[
\begin{aligned}
&\gamma^d(k_d,\ldots,k_{n+1},k_n+1,\ldots,k_1)\gamma^1(k_n)\\
&\qquad=\gamma^d(k_d,\ldots,k_{n+1}+1,k_n,\ldots,k_1)\gamma^1(k_{n+1}),\\
&\gamma^d(k_d+1,k_{d-1},\ldots,k_1)\gamma^1(k_d)\\
&\qquad=\gamma^d(k_d,\ldots,k_1)\gamma^1(k_1+\cdots+k_d+2-d).
\end{aligned}
\]
\SourceFigure{figure-13-3}{0.4880764}{13.3}
作初等变换，得到
\[
\frac{\gamma^d(k_d,\ldots,k_1)}{\gamma^d(0,\ldots,0)}
=\frac{\Gamma(k_1)\Gamma(k_2)\cdots\Gamma(k_d)\Gamma(k_1+\cdots+k_d+2-d)}{\Gamma(2-d)},
\]
\SourcePage{206}
其中 $k\geq0$ 时 $\Gamma(k)=\gamma^1(0)\gamma^1(1)\cdots\gamma^1(k-1)$，$k<0$ 时 $\Gamma(k)=\gamma^1(k)\gamma^1(k+1)\cdots\gamma^1(-1)$。对\sr{fig:13.3}{图~13.3}(ii) 作类似论证，得到
\[
\begin{aligned}
&\gamma^d(k_d,\ldots,k_{n+1},k_n,\ldots,k_1)\gamma^1(k_n)\\
&\qquad=\gamma^{d-1}(k_d,\ldots,k_n+k_{n+1},\ldots,k_1)\gamma^2(k_{n+1},k_n),
\end{aligned}
\]
{\looseness=-1
从而 $\gamma^d(0,\allowbreak \ldots,\allowbreak 0)=\gamma^2(0,\allowbreak 0)^{d-1}\Gamma(2-d)$。因此，$\F(\symbfit{L}^\#)^\to$ 与 $\widetilde{\F}(\symbfit{L}^\#)^\to$ 之间有一个简单的线性同构，它将每个 Floer 上链群 $CF^k$ 乘以 $\Gamma(k)\gamma^2(0,\allowbreak 0)$。这已足够满足我们的目的，因此不尝试显式确定剩下的符号。\par}

\clearpage
\section{辛对合}\label{sec:14}\label{section:14}
\subsection{等变膜}\label{sec:14a}
设 $(M,\allowbreak \omega_M)$ 为辛流形，配有殆复结构 $I_M$、二次复体积形式 $\eta_M^2$，以及保持 $(\omega_M,\allowbreak I_M,\allowbreak \eta_M^2)$ 的对合 $\iota_M$。\znote{199}{此处 $\iota_M$ 称为辛对合：它是二阶辛同胚，且保持 $I_M$，故在每点切空间上给出复线性对合，其不动点集 $M^{\mathrm{fix}}$ 因而是复子流形，法向切空间带复结构；当法向复维数为一（最简单的非平凡情形）时 $M^{\mathrm{fix}}$ 的实余维为二，这就是后文假设 14.3 的几何来源；保持 $\eta_M^2$ 则保证不动点集上仍留有可用的二次复体积形式，后文讨论分次时正要用到它。若法向复维数更高，下述整套论证需要更多技术条件，故假设 14.3 只处理余维二的情形。[来源：原书 §(14c) 假设 14.3，印刷页 203。]}希望将 $\iota_M$ 对 Lagrangian 子流形的自然作用与 Pin 结构的改变结合。这在纯形式层面造成一些困难，下面予以处理。
\begin{definition*}
\emph{预等变膜}由以下数据组成：其主体，即 Lagrangian 膜 $L_+^\#=(L_+,\allowbreak \alpha_+^\#,\allowbreak P_+^\#)$；其影子，即分次 Lagrangian 子流形 $(L_-=\iota_M(L_+),\allowbreak \alpha_-^\#=\alpha_+^\#\circ\iota_M)$ 上的 Pin 结构 $P_-^\#$；以及一对同构
\[
f_+:P_+^\#\longrightarrow\iota_M^*(P_-^\#\otimes\lambda^{\mathrm{top}}(TL_-)),\qquad
f_-:P_-^\#\longrightarrow\iota_M^*(P_+^\#\otimes\lambda^{\mathrm{top}}(TL_+)),
\]
使复合
\begin{equation}\label{eq:14.1}
\begin{aligned}
P_+^\#&\xrightarrow{f_+}\iota_M^*(P_-^\#\otimes\lambda^{\mathrm{top}}(TL_-))\\
&\xrightarrow{\iota_M^*(f_-)\otimes\Id}
P_+^\#\otimes\lambda^{\mathrm{top}}(TL_+)\otimes\iota_M^*\lambda^{\mathrm{top}}(TL_-)
\cong P_+^\#
\end{aligned}
\end{equation}
等于 Pin 结构的非平凡自同构 $-e\in\ker(\operatorname{Pin}_n\to O_n)$；最后一步使用显然的同构 $\iota_M^*(TL_-)\cong TL_+$。\znote{199}{$\operatorname{Pin}_n\to O_n$ 是二重覆叠，核同构于 $\Z/2$，其非平凡元素记作 $-e$（$e$ 为恒等元），它实现标架丛二重覆叠的换层。按二次型 $\pm\sum x_i^2$ 的选取，Pin 群有两种约定 $\operatorname{Pin}_\pm(n)$，二者相差一个符号；因此“平方等于 $e$ 还是 $-e$”必须逐处指明，定义中特意写 $-e$ 而不写“平凡”，正是为了固定这一约定。当 $n$ 较小时两种约定还可能给出同构的群，但一般情形下并不相同，因此对照其他文献时必须先核对它采取哪一种，否则符号会整体差一个。[来源：nLab "pin group"，https:/\allowbreak /\allowbreak ncatlab.org/\allowbreak nlab/\allowbreak show/\allowbreak pin+\allowbreak group（2026-11-15 访问）]}
\end{definition*}
为使措辞简洁，有时称之为“预等变膜 $L_+^\#$”，但总默认辅助数据 $P_-^\#$ 与 $f_-$ 已作某个具体选择；也可用 $f_+$，因为两个同构中任一个都唯一确定另一个。
\begin{note}\label{st:14.1}
每个 Lagrangian 膜 $L_+^\#$ 都可以显然地成为预等变膜：令 $P_-^\#=\iota_M^*(P_+^\#\otimes\lambda^{\mathrm{top}}(TL_+))$，取 $f_-$ 为恒等，并相应调整 $f_+$。
\end{note}
预等变膜上有自然对合 $\epsilon$，交换正、负两部分；简记为 $\epsilon(L_+^\#)=L_-^\#$。这里使用了另一顺序的复合 $(\iota_M^*(f_+)\otimes\Id)\circ f_-$ 也等于 $-e$ 这一事实。
\begin{definition*}
\emph{等变膜}是此对合的不动点。具体地，它由不变 Lagrangian 子流形 $L$、$\iota_M(L)=L$，一个 $\iota_M$ 不变分次 $\alpha^\#$，以及 Pin 结构 $P^\#$ 连同同构 $f:P^\#\to\iota_M^*(P^\#\otimes\lambda^{\mathrm{top}}(TL))$ 组成，使其“平方”
\SourcePage{207}
\begin{equation}\label{eq:14.2}
\begin{aligned}
P^\#&\xrightarrow{f}\iota_M^*(P^\#\otimes\lambda^{\mathrm{top}}(TL))\\
&\xrightarrow{\iota_M^*(f)\otimes\Id}P^\#\otimes\lambda^{\mathrm{top}}(TL)\otimes\iota_M^*\lambda^{\mathrm{top}}(TL)\cong P^\#
\end{aligned}
\end{equation}
等于 $-e$。
\end{definition*}
给定 $\iota_M$ 不变 Lagrangian 子流形 $L$，考察它容许等变膜结构的条件。假设 $H^1(L)$ 中的 Maslov 类消失，使 $L$ 容许分次 $\alpha^\#$。若 $L$ 连通，则差 $\alpha^\#\circ\iota_M-\alpha^\#$ 是 $\iota_M$ 反不变的常值函数，因此自动为零。另一方面，若 $L$ 有两个被 $\iota_M$ 交换的连通分支，可调整其中一个分支上的 $\alpha^\#$ 使该差消失。一般情形由这两种情形得到；结论是，得到不变分次没有额外障碍。

Pin 部分更有趣。最直接地，可以逐步处理：先为 $L$ 赋予 Pin 结构，再试图找到同构 $f:P^\#\to\iota_M^*(P^\#\otimes\lambda^{\mathrm{top}}(TL))$，最后确定它的平方。障碍分别位于 $H^2(L;\Z/2)$、$H^1(L;\Z/2)$ 和 $H^0(L;\Z/2)$；但后两项是次级障碍，依赖先前选择，并不取遍所有值。也可将这些方面结合成一个更典范的单一障碍，即等变上同调类
\begin{equation}\label{eq:14.3}
w_2^{\mathrm{eq}}(L)+w_1^{\mathrm{eq}}(L)u+u^2\in H_{\Z/2}^2(L;\Z/2).
\end{equation}
这里 $u$ 为 $H^*(B\Z/2;\Z/2)=\Z/2[u]$ 的生成元，$w_k^{\mathrm{eq}}(L)$ 为 $TL$ 的等变 Stiefel–Whitney 类，依赖其等变向量丛结构。若 $\iota_M|L$ 没有不动点，有一个稳定化技巧，可将问题归结到非等变 Pin 结构的标准问题，从而说明\eref{14.3}。具体地，令 $\pi:L\to\bar L$ 为到商 $\bar L=L/(\Z/2)$ 的覆盖，$\beta\to\bar L$ 为关联实线丛。\znote{199}{$\Z/2$ 在 $L$ 上自由作用，故商 $\bar L=L/(\Z/2)$ 是光滑流形、$\pi$ 是二重覆叠；$\beta$ 是符号表示所关联的实线丛，其拉回 $\pi^*\beta$ 平凡。于是 $TL$ 的等变结构与 $T\bar L\oplus\beta$ 的结构相当，再添三个 $\beta$ 使作用自由，就把等变 Pin 结构化为通常的 Pin 结构；$\R^3$ 的维数三正是引理 11.14 中 $\operatorname{rank}(F^{\mathrm{anti}})=3$ 情形的来源。反过来看，正因为 $\pi^*\beta$ 平凡、且 $\Z/2$ 通过符号表示作用在 $\beta$ 上，$L$ 上的 Pin 结构与 $\bar L$ 上的 Pin 结构加 $\beta$ 的扭曲才只是同一份数据的两种写法。[来源：原书 §(14a)，印刷页 199–200，并见引理 11.14。]}假设在 $T\bar L\oplus\beta\oplus\beta\oplus\beta$ 上有 Pin 结构 $\bar P^\#$，即
\begin{equation}\label{eq:14.4}
\begin{aligned}
w_2(T\bar L\oplus\beta\oplus\beta\oplus\beta)
&=w_2(T\bar L)+w_1(T\bar L)w_1(\beta)+w_1(\beta)^2\\
&\in H^2(\bar L;\Z/2)
\end{aligned}
\end{equation}
必须为零。取拉回 $\pi^*\bar P^\#$，它是 $TL\oplus\R^3$ 上的 Pin 结构，令 $P^\#$ 为其在 $TL$ 上的限制。该拉回有显然的自同构，覆盖 $D(\iota_M|L)\oplus(-\Id_{\R^3})$ 在标架丛上的作用。希望将其与覆盖 $\Id_{TL}\oplus(-\Id_{\R^3})$ 的辅助自同构 $g$ 复合，再将复合限制到 $P^\#$。由\eref{11.30}与\stref{引理}{11.14}，$g$ 的选择对应于丛 $\R^3\otimes\lambda^{\mathrm{top}}(TL\oplus\R^3)$ 的定向，也就是 $TL$ 本身的定向。将这一因子计入后，例如选择局部定向，再拼接所得 $P^\#$ 的局部自同构，就得到 $P^\#$ 的自然自同构 $f:P^\#\to\iota_M^*(P^\#\otimes\lambda^{\mathrm{top}}(TL))$。依定义，$\pi^*\bar P^\#$ 的原始自同构为对合；由
\SourcePage{208}
\stref{引理}{11.14}中 $\operatorname{rank}(F^{\mathrm{anti}})=3$ 的情形，有 $g^2=-e$；二者交换，因而\eref{14.2}具有所需性质。论证也可反向进行，说明等变膜结构的存在蕴含\eref{14.4}消失，而在 $H_{\Z/2}^*(L)\cong H^*(\bar L)$ 下，它恰对应于\eref{14.3}。

对一般的、非自由作用的情形，可以通过 Borel 构造将任意 $\Z/2$ 作用转为自由作用，立即说明\eref{14.3}消失是 $L$ 上等变膜结构存在的必要条件。\znote{200}{Borel 构造把任意 $\Z/2$ 作用换成自由作用：取 $L\times_{\Z/2}E\Z/2$（对角作用），它是同伦商，其普通上同调即等变上同调 $H^*_{\Z/2}(L)$；当 $\Z/2$ 自由作用时它同伦等价于商空间 $L/(\Z/2)$。因此“等变上同调类消失”可以在这一自由作用空间上按通常的示性类理论验证，这便是正文用它论证必要性的道理。同伦商的构造不依赖于 $E\Z/2$ 的具体模型，故所得障碍类只依赖作用本身，可作为不变量使用。[来源：nLab "Borel construction"，https:/\allowbreak /\allowbreak ncatlab.org/\allowbreak nlab/\allowbreak show/\allowbreak Borel+\allowbreak construction（2026-11-15 访问）]}也可用这种方法证明充分性，但需要进一步讨论 Borel 构造上的 Pin 结构，这里省略。

为补充前述讨论，更直接地考察不变 Lagrangian 子流形 $L$ 的不动点集 $L^{\mathrm{fix}}$ 上的情形。将 $TL|L^{\mathrm{fix}}=F^{\mathrm{inv}}\oplus F^{\mathrm{anti}}$ 分成 $D\iota_M$ 不变与反不变部分。假设 $L$ 上有 Pin 结构 $P^\#$ 及同构 $f:P^\#\to\iota_M^*(P^\#\otimes\lambda^{\mathrm{top}}(TL))$。由\eref{11.30}，$f$ 在 $L^{\mathrm{fix}}$ 上的限制由
\[
\operatorname{Pin}(TL\otimes\lambda^{\mathrm{top}}(TL))\times_{\Z/2}S(\lambda^{\mathrm{top}}(TL))
\]
的一个截面给出，它覆盖 $D(\iota_M|L)\otimes\Id\in O(TL\otimes\lambda^{\mathrm{top}}(TL))$。再用\stref{引理}{11.14}可知，这种截面对应于 $(F^{\mathrm{anti}}\otimes\lambda^{\mathrm{top}}(TL))\oplus\lambda^{\mathrm{top}}(TL)$ 的定向，或等价地，对应于下列线丛的平凡化：
\begin{equation}\label{eq:14.5}
\lambda^{\mathrm{top}}(F^{\mathrm{anti}})^{\otimes\operatorname{rank}(F^{\mathrm{anti}})}
\otimes\lambda^{\mathrm{top}}(F^{\mathrm{inv}})^{\otimes(\operatorname{rank}(F^{\mathrm{anti}})+1)}\longrightarrow L^{\mathrm{fix}}.
\end{equation}
此外，若 $\frac12\operatorname{rank}(F^{\mathrm{anti}})(\operatorname{rank}(F^{\mathrm{anti}})+1)$ 为偶数，$f$ 的平方为 $e$；若为奇数，则为 $-e$。这里的平方按\eref{14.2}理解，所以相比\eref{11.29}多出符号 $(-1)^{\operatorname{rank}(F^{\mathrm{anti}})}$；它来自 $\det(D(\iota_M|L))$，是配对 $\iota_M^*(\lambda^{\mathrm{top}}(TL))\otimes\lambda^{\mathrm{top}}(TL)\to\R$ 出现的结果。顺便指出，这一论证的大部分也可从抽象障碍理论\eref{14.3}出发进行，只需注意，令 $r=\operatorname{rank}(F^{\mathrm{anti}})$、$F=TL|L^{\mathrm{fix}}$，有
\[
\begin{aligned}
&(w_2^{\mathrm{eq}}(L)+w_1^{\mathrm{eq}}(L)u+u^2)|L^{\mathrm{fix}}\\
&\quad=w_2(F)+(r\,w_1(F^{\mathrm{anti}})+(r+1)w_1(F^{\mathrm{inv}}))u
+\left(\frac{r(r+1)}2+1\right)u^2\\
&\quad\in H^*(L^{\mathrm{fix}};\Z/2)[u].
\end{aligned}
\]
最终，只有一个特殊情形与应用有关，将其记录下来以备后用：
\begin{lemma}\label{st:14.2}
设 $L$ 为 $M$ 中 $\iota_M$ 不变的 Lagrangian 子流形。假设 $L$ 连通，其不动点集 $L^{\mathrm{fix}}$ 也连通、非空且余维为一。再假设 $L$ 上有同构于 $\iota_M^*(P^\#\otimes\lambda^{\mathrm{top}}(TL))$ 的 Pin 结构 $P^\#$。则同构 $f:P^\#\to\iota_M^*(P^\#\otimes\lambda^{\mathrm{top}}(TL))$ 的选择与 $L^{\mathrm{fix}}\subset L$ 的余定向双射对应。此外，任一这样的同构按\eref{14.2}取平方，均为 $-e$。
\end{lemma}
\SourcePage{209}
显然，$M$ 的 $\Z/2$ 等变分次辛自同构作用于预等变膜集合，并与对合 $\epsilon$ 严格交换，特别地保持等变膜。等变线丛 $\beta\to M$ 也可作用：给定这种线丛及预等变膜，将 $P_\pm^\#$ 换成 $P_\pm^\#\otimes\beta|L_\pm$，利用 $\beta$ 的等变结构给出的同构 $\beta|L_\pm\cong\iota_M^*(\beta|L_\mp)$，将 $f_\pm$ 变成新的 Pin 结构对之间的同构。例如，取 $\beta=\R\times M$，带非平凡 $\Z/2$ 作用，对应于 $\Z/2$ 的非平凡实特征标；与它作张量积，等价于将两个 $f_\pm$ 均乘以 $-e$。回到一般 $\beta$，将此作用施加于等变膜，并考察不动点集 $L^{\mathrm{fix}}\subset L$，其效果是使\eref{14.5}的平凡化改变一个符号，该符号为 $\Z/2$ 在 $\beta|L^{\mathrm{fix}}$ 上的作用符号。

\subsection{等变 Donaldson–Fukaya 范畴}\label{sec:14b}
值得看看上述定义如何融入小节 \sr{sec:10b}{(10b)} 与 \sr{sec:12i}{(12i)} 中较抽象的讨论。

假设给定带协调 $\Z/2$ 作用的范畴 $A$。该作用由自等价 $E:A\to A$ 和函子同构 $T:E^2\to\Id_A$ 组成，满足唯一的协调条件 $R_E(T)=L_E(T)$，更具体地，
\begin{equation}\label{eq:14.6}
T_{E(X)}=E(T_X)\in\Hom_A(E^3(X),E(X)),\qquad\text{对所有 }X.
\end{equation}
于是可以如下定义新范畴 $A^{\mathrm{strict}}$。对象形如 $(X_\pm,\allowbreak F_+,\allowbreak F_-)$，其中 $X_\pm$ 为 $A$ 的对象，$F_\pm:X_\pm\to E(X_\mp)$ 为同构，使 $E(F_-)\circ F_+:X_+\to E^2(X_+)$ 等于 $T_{X_+}^{-1}$。由\eref{14.6}，从 $X_-$ 出发的相应性质自动成立。$A^{\mathrm{strict}}$ 中从 $(X_{0,\allowbreak \pm},\allowbreak F_{0,\allowbreak \pm})$ 到 $(X_{1,\allowbreak \pm},\allowbreak F_{1,\allowbreak \pm})$ 的态射，就是 $A$ 中从 $X_{0,\allowbreak +}$ 到 $X_{1,\allowbreak +}$ 的态射；立即可见 $A^{\mathrm{strict}}$ 等价于 $A$。不过，现在有了严格 $\Z/2$ 作用：函子 $E^{\mathrm{strict}}$ 交换对象的正、负两部分，并将
\[
\begin{aligned}
\Phi\in\Hom_A(X_{0,+},X_{1,+})\longmapsto{}
F_{1,-}^{-1}\circ E(\Phi)\circ F_{0,-}\in\Hom_A(X_{0,-},X_{1,-}).
\end{aligned}
\]
特别地，$A^{\mathrm{strict}}$ 中在 $E^{\mathrm{strict}}$ 下不变的对象，由 $A$ 的单个对象 $X$ 及同构 $F:X\to E(X)$ 给出，其中 $E(F)\circ F=T_X^{-1}$。若 $(X_0,\allowbreak F_0)$ 与 $(X_1,\allowbreak F_1)$ 为两个这样的对象，则态射群 $\Hom_{A^{\mathrm{strict}}}(X_0,\allowbreak X_1)=\Hom_A(X_0,\allowbreak X_1)$ 带自然对合 $E^{\mathrm{strict}}(\Phi)=F_1^{-1}\circ E(\Phi)\circ F_0$。\znote{201}{把协调作用换成严格作用的手法在范畴论中很常见：让对象额外携带“实现同构”的数据（这里是 $F_\pm$），态射则忘掉这些数据。于是新模式与原范畴等价，而群作用成为严格函子作用，不再需要额外的协调同构；代价是每个对象的严格化不唯一（相差一个自动同构），但对同调范畴无影响。第 (10b)、(12i) 小节用过同样的构造。[来源：原书 §(10b)、§(12i)。]}

转到几何情形：$M$ 为带角正合辛流形，带二次复体积形式 $\eta_M^2$。假设有保持全部给定结构的对合 $\iota_M$；这里包括 $\omega_M,\allowbreak I_M$ 及一形式 $\theta_M$，以后谈到正合辛流形的对合时，均以此为默认假设。添加
\SourcePage{210}
平凡分次，即零分次，以及线丛 $\beta=\xi_M$，将它变成膜自同构 $\iota_M^\#$。由一般理论，这定义 Donaldson–Fukaya 范畴的唯一自等价，作为我们的 $E$。此前考虑的是边界附近等于恒等的自同构；不过，除了正合性外，只需较弱的条件，即自同构在边界附近保持 $I_M$，而这里显然满足。用\stref{笔记}{11.18}的同构 $\xi_M|L\cong\lambda^{\mathrm{top}}(TL)$，可将 $E$ 对对象的作用写为
\begin{equation}\label{eq:14.7}
(L,\alpha^\#,P^\#)\longmapsto
(\iota_M(L),\alpha^\#\circ\iota_M,\iota_M^*(P^\#\otimes\lambda^{\mathrm{top}}(TL))).
\end{equation}
因为 $\iota_M^\#\circ\iota_M^\#=(\Id_M,\allowbreak 0,\allowbreak \xi_M\otimes\iota_M^*\xi_M)$ 平凡，$E^2=E\circ E$ 同构于恒等函子。更确切地，$\xi_M\otimes\iota_M^*\xi_M$ 的平凡化给出同构 $T:E^2\to\Id$。通过同构 $\lambda^{\mathrm{top}}(D\iota_M):\xi_M\to\iota_M^*\xi_M$ 有一个显然的平凡化，但我们要将它乘以 $-1$，并取相应的 $T$。唯一非平凡的协调条件\eref{14.6}归结为简单的事实：这一平凡化本身 $\Z/2$ 不变，其中作用是沿 $\iota_M$ 拉回并交换两个因子。

现在显然，预等变膜的定义正是上述抽象代数构造的一个实例。特别地，有 $\K$ 线性分次范畴 $H\F(M)^{\mathrm{strict}}$，对象为预等变膜，态射空间为 $HF^*(L_{0,\allowbreak +}^\#,\allowbreak L_{1,\allowbreak +}^\#)$。它等价于 $H\F(M)$，其唯一优势是带严格的 $\Z/2$ 作用。对对象，该作用就是之前描述的对合 $\epsilon$；对态射，规定
\begin{equation}\label{eq:14.8}
\begin{aligned}
\epsilon:HF^*(L_{0,+}^\#,L_{1,+}^\#)
&\xrightarrow{\cong}HF^*(L_{0,+}^\#\otimes\lambda^{\mathrm{top}}(TL_{0,+}),L_{1,+}^\#\otimes\lambda^{\mathrm{top}}(TL_{1,+}))\\
&\xrightarrow{\cong}HF^*(\iota_M(L_{0,+}^\#\otimes\lambda^{\mathrm{top}}(TL_{0,+})),
\iota_M(L_{1,+}^\#\otimes\lambda^{\mathrm{top}}(TL_{1,+})))\\
&\xrightarrow{\cong}HF^*(L_{0,-}^\#,L_{1,-}^\#).
\end{aligned}
\end{equation}
第一个同构使用 $\lambda^{\mathrm{top}}(TL_{k,\allowbreak +})$ 与同一个线丛 $\xi_M$ 的限制典范同构这一事实，比较\eref{14.7}；第二个来自分次辛自同构下的函子性；第三个来自与 $f_{0,\allowbreak -},\allowbreak f_{1,\allowbreak -}$ 的复合。

关于 $\Z/2$ 在 $H\F(M)^{\mathrm{strict}}$ 上的作用，还可陈述一些性质，基本上是\secref{12e}诸性质的等变版本。证明直接，故省略。
\begin{itemize}
\item \emph{改变符号。}令 $\beta\to M$ 为带非平凡 $\Z/2$ 作用的平凡线丛。显然，对所有预等变膜 $L_{k,\allowbreak +}^\#$，均有 $HF^*(L_{0,\allowbreak +}^\#,\allowbreak L_{1,\allowbreak +}^\#\otimes\beta|L_{1,\allowbreak +})=HF^*(L_{0,\allowbreak +}^\#,\allowbreak L_{1,\allowbreak +}^\#)$，但这些同构与\eref{14.8}反交换。特别地，若 $L_0^\#,\allowbreak L_1^\#$ 是真正的等变膜，则将其中一个与 $\beta$ 作张量积，会反转 $\Z/2$ 在 $HF^*(L_0^\#,\allowbreak L_1^\#)$ 上作用的符号，因而
\SourcePage{211}
在底域满足 $\operatorname{char}(\K)\ne2$ 时，交换不变部分与反不变部分。\znote{202}{把膜与 $\beta$ 张量积，相当于把所有 $f_\pm$ 同时乘以 $-e$。由于 $f_\pm$ 正是联系主部 $L_+$ 与影子 $L_-$ 的数据，这一操作在链层次给出同构，但它与 $\epsilon$ 的作用反交换，故在上同调上把 $\Z/2$ 作用变为负作用；特征不等于二时，正、负本征空间随之互换。这也说明 $\beta$ 的 $\Z/2$ 作用是本质的，不能通过重新选择平凡化消去。从表示论看，这是 $\Z/2$ 的一维实表示只有两个、而张量积给出群同态这一最简单事实的直接体现。[来源：原书 §(14a)，印刷页 201。]}
\item \emph{PSS 同构。}若 $L^\#$ 为等变膜，则 PSS 同构 $HF^*(L^\#,\allowbreak L^\#)\cong H^*(L;\K)$ 与两侧的 $\Z/2$ 作用相容。
\item \emph{同痕不变性。}取正合预等变 Lagrangian 膜的同痕，即族 $\{(L_{s,\allowbreak +}^\#,\allowbreak P_{s,\allowbreak -}^\#,\allowbreak f_{s,\allowbreak +},\allowbreak f_{s,\allowbreak -})\}$，$0\leq s\leq1$。如通常一样，$\{L_{s,\allowbreak +}^\#\}$ 确定 $HF^0(L_{0,\allowbreak +}^\#,\allowbreak L_{1,\allowbreak +}^\#)$ 的一个元素，其影子确定 $HF^0(L_{0,\allowbreak -}^\#,\allowbreak L_{1,\allowbreak -}^\#)$ 的一个元素。这两个元素实际上在 $\Z/2$ 作用 $\epsilon$ 下互为像。特别地，若两个正合等变膜在此类中同痕，则它们之间在 $H\F(M)^{\mathrm{strict}}$ 中存在不变同构。
\end{itemize}

\subsection{等变横截性}\label{sec:14c}
接下来要解决的问题，是通过上链层次的朴素群作用实现 $\Z/2$ 对称性。我们选择施加一系列几何约束，以简化这一问题的解决，而不追求最大的普遍性。

\StatementNumber{3}\begin{assumption}\label{st:14.3}
对于给定的带角正合辛流形 $M$ 及其对合 $\iota_M$，要求不动点集 $M^{\mathrm{fix}}$ 的实余维为二，且其法丛 $\nu$ 平凡。实际上，还要固定该丛的一个复平凡化 $\nu\to\CC$，并加上下述部分可积性条件：这一平凡化应是某个映射 $\pi:U\to\CC$ 的导数，其中 $U\subset M$ 是不动点集的某个邻域；映射满足 $\pi\circ\iota_M=-\pi$，且相对于 $I_M$ 计算的 $\bar\partial\pi$ 沿 $M^{\mathrm{fix}}$ 连同一阶导数一起消失。
\end{assumption}

假设给定 $I_M$ 及平凡化 $D\pi|M^{\mathrm{fix}}$。若 $N(I_M)$ 是 Nijenhuis 张量，则只要存在具有上述性质的 $\pi$，便有 $D\pi\circ N(I_M)$ 在 $M^{\mathrm{fix}}$ 上消失；这解释了“部分可积性”这一名称。\znote{203}{“部分可积性”指 $I_M$ 只在法向上可积。映射 $\pi$ 在 $M^{\mathrm{fix}}$ 上只被要求到一阶全纯（$\bar\partial\pi$ 及其一阶导数消失），因为 $M^{\mathrm{fix}}$ 是实子流形，而紧流形上整体全纯的 $\pi$ 只能是常值。存在这样的 $\pi$ 迫使法向部分的 Nijenhuis 张量消失，即 $D\pi\circ N(I_M)=0$；这正是复结构在 $M^{\mathrm{fix}}$ 的复法向方向上可积的表现。若 $I_M$ 沿法向也不可积，则不存在这样的 $\pi$；所以假设 14.3 实质上是对 $I_M$ 的一个限制，而非纯技术性的方便条件。[来源：原书 §(14c) 假设 14.3 及其后的说明，印刷页 203–204。]}

\StatementNumber{4}\begin{assumption}\label{st:14.4}
在所有 $\iota_M$ 不变的正合 Lagrangian 子流形 $L$ 中，只考虑完全不与不动点集相交，或从恒定方向趋近不动点集的那些。后一个条件是指，存在某个 $\sigma\in\R$，称为 $L$ 的\emph{法向角}，使对所有 $x\in L^{\mathrm{fix}}$ 都有
\[
D\pi(TL_x)=e^{i\sigma}\R\subset\CC.
\]
$\sigma$ 的可能取值有 $\Z$ 个，彼此相差 $\pi$ 的整数倍，但假设已经为每个有关的 Lagrangian 子流形选定其中一个。
\end{assumption}

继续讨论之前，需要引入一些术语。令 $H$ 为 $M$ 上的 $\iota_M$ 不变函数。在每一点 $x\in M^{\mathrm{fix}}$，法向的 Hessian $D^2H$
\SourcePage{212}
是 $\nu_x$ 上良定义的二次型，因为一阶导数必须消失。若对 $Z\in\nu_x$ 有
\begin{equation}\label{eq:14.9}
(D^2H)(Z)=h\,|Z|^2,
\end{equation}
则称 $H$ 以速度 $h\in\R$ \emph{无穷小地旋转} $\nu_x$。这里的 $|\cdot|$ 是将 $\nu$ 与 $TM|M^{\mathrm{fix}}$ 的反不变部分认同所得的 Hermite 度量。等价地，取 Hamilton 向量场 $X_H$，经 $D\pi$ 投影得到函数 $D\pi(X_H):U\to\CC$；则\eref{14.9}是说，在 $x$ 附近的一阶展开中有 $D\pi(X_H)=ih\pi$。

\clearpage
\StatementNumber{5}\begin{assumption}\label{st:14.5}
取满足\stref{假设}{14.4}的一对不变正合 Lagrangian 子流形 $(L_0,L_1)$，两者均与不动点集相交，法向角分别为 $\sigma_0,\sigma_1$。考虑 $\Z/2$ 不变的 Floer 数据 $(H,J)$ 时，施加下述附加要求。首先，与对 $I_M$ 的假设一样，要求相对于任意 $J_t$ 计算的导数 $\bar\partial\pi$ 沿 $M^{\mathrm{fix}}$ 连同一阶导数一起消失。其次，令 $\phi$ 为 $H$ 的“流”，则应存在满足
\begin{equation}\label{eq:14.10}
\sigma_1-\sigma_0-h\in(-\pi,0)
\end{equation}
的 $h\in\R$，以及开子集 $V\subset[0,1]\times M^{\mathrm{fix}}$，使其对所有 $t\in[0,1]$ 均包含 $\{t\}\times\phi^t(L_0^{\mathrm{fix}})$，并使 $H(t,\cdot)$ 对每个 $(t,x)\in V$ 都以速度 $h$ 无穷小地旋转 $\nu_x$。
\end{assumption}

这里只要求在足够大的开子集上控制法向 Hessian，因为按照 Floer 数据的定义，$H(t,\cdot)$ 必须在 $\partial M$ 附近为零，故不可能在整个 $M^{\mathrm{fix}}$ 上满足\eref{14.9}。关键在于，对 $x\in L_0^{\mathrm{fix}}$，映射 $D\phi_x^t:\nu_x\to\nu_{\phi^t(x)}$ 是乘以 $e^{iht}$ 的旋转。特别地，$\phi^1(L_0)$ 从方向 $e^{i(\sigma_0+h)}\R$ 趋近不动点集。仍由 Floer 数据的定义，该方向必须与 $L_1$ 横截；单是这一点就已蕴含 $\sigma_1-\sigma_0-h\in\R\setminus\pi\Z$，当然，$\phi^1(L_0)\cap L_1\cap M^{\mathrm{fix}}=\varnothing$ 的情形除外。附加要求\eref{14.10}选定 $\R\setminus\pi\Z$ 的一个“室”，即一个连通分支。这一条件看起来可能相当空泛，因为总能改变 $\sigma_0$ 或 $\sigma_1$ 来得到所需结果；但同时考虑多个 Floer 数据时，它就有了实际意义。

具有这些性质的函数仍十分丰富。特别地，可以任意指定 $H$ 在不动点集上，以及在该集合任意小邻域之外的值，但仍须在边界附近为零，且当然要 $\iota_M$ 不变。对殆复结构也一样，理由如下。必要时先缩小 $M^{\mathrm{fix}}$ 的邻域 $U$，使 $\pi|U$ 成为具有辛纤维的浸没。切丛 $TM|U$ 分解成 $TM^v=\ker(D\pi)$ 及其辛正交补 $TM^h$。后者由 $D\pi$ 平凡化，且在 $M^{\mathrm{fix}}$ 上的限制恰为 $\nu$。由于 $I_M$ 与平凡化 $D\pi|M^{\mathrm{fix}}$ 相容，沿不动点集，$I_M$ 的 $(hv)$ 分量为零，$(hh)$ 分量等于 $i$。对 $\bar\partial\pi$ 的附加要求
\SourcePage{213}
是说，这两个条件还应延伸到一阶导数；实际上，对 $(hh)$ 分量，这一点已由 $\iota_M$ 对称性推出。寻找具有同样性质的其他殆复结构时，显然可以自由指定其沿 $M^{\mathrm{fix}}$ 的 $(vv)$ 分量。

\StatementNumber{6}\begin{lemma}\label{st:14.6}
令 $L_0,L_1$ 为不变正合 Lagrangian 子流形，$(H,J)$ 为\stref{假设}{14.5}中的不变 Floer 数据。取 $u\in\M_Z(y_0,y_1)^*$，其像完全包含在 $M^{\mathrm{fix}}$ 中。则 $u$ 在 $M$ 中正则，当且仅当它在 $M^{\mathrm{fix}}$ 中正则。
\end{lemma}
\begin{proof}
线性化算子是 $\Z/2$ 等变的，因此分解为
\[
D_{Z,u}=D^-_{Z,u}\oplus D^+_{Z,u}.
\]
不变部分容易辨认。具体地，将 $u$ 视为 $M^{\mathrm{fix}}$ 中关于 Lagrangian 子流形 $(L_0^{\mathrm{fix}},L_1^{\mathrm{fix}})$ 及 Floer 数据 $(H|M^{\mathrm{fix}},J|M^{\mathrm{fix}})$ 的 Floer 方程解，并记为 $u^{\mathrm{fix}}$，便有
\begin{equation}\label{eq:14.11}
D_{Z,u^{\mathrm{fix}}}=D^+_{Z,u}.
\end{equation}
对于反不变部分，可用 $D\pi$ 将有关向量丛认同为 $E=Z\times\CC\to Z$，边界上的子丛 $F\subset E|\partial Z$ 则成为 $F_{s,0}=e^{i\sigma_0}\R$、$F_{s,1}=e^{i\sigma_1}\R$。沿 $M^{\mathrm{fix}}$，精确到一阶有 $D\pi\circ J_t=i\circ D\pi$。此外，当 $|s|\gg0$ 时，$x=u(s,t)$ 接近 $\phi^t(L_0)$，因而 $(t,x)\in V$；这意味着 $H(t,\cdot)$ 的 Hamilton 向量场 $X(t,\cdot)$ 在 $x$ 处精确到一阶满足 $D\pi(X(t,x))=ih\pi(x)$。因此，在条带状端上，
\begin{equation}\label{eq:14.12}
D^-_{Z,u}=\partial_s+i\partial_t+h.
\end{equation}
可以变形该算子，使\eref{14.12}处处成立，从而得到 $\operatorname{index}(D^-_{Z,u})=0$，例如见\stref{引理}{11.9}。另一方面，\stref{引理}{11.5}在 $\mu=0$、$|\Sigma^-|=1$ 时断言 $D^-_{Z,u}$ 单射，故它必可逆。结合\eref{14.11}便得到结论。
\end{proof}

有了这一结果，按照常规即可证明：\stref{假设}{14.5}所定义的这一相当特殊类别中的 Floer 数据一般是正则的。给定其中一个数据，先只考察不动点集 $M^{\mathrm{fix}}$，作微小扰动，使其在该处正则。由\stref{引理}{14.6}，这保证 $M$ 中所有不变 Floer 轨线正则。其余轨迹可用标准横截性论证的一种变体处理，见\bcite{85}命题 5.13。\znote{205}{把线性化算子按 $\Z/2$ 分解，是处理等变横截性的标准手法：等变算子与对合交换，故核与余核都带 $\Z/2$ 分解，整体满射当且仅当两部分各自满射。不变部分恰是 $M^{\mathrm{fix}}$ 中相应问题的算子\eref{14.11}；反不变部分因指标为零且单射而自动可逆。于是 $M$ 中的正则性完全由 $M^{\mathrm{fix}}$ 中的正则性决定，这正是引理 14.6 的实质，也说明为何只须在不动点集上作扰动。[来源：原书 §(14c) 引理 14.6 及其证明，印刷页 205。]}

令 $(L_0,\ldots,L_d)$，$d\geq2$，为满足\stref{假设}{14.4}的 $\iota_M$ 不变正合 Lagrangian 子流形，每一个都与不动点集相交，法向角分别为 $\sigma_0,\ldots,\sigma_d$。假设已按\stref{假设}{14.5}为子流形对 $(L_0,L_d)$ 及
\SourcePage{214}
$(L_{k-1},L_k)$，$1\leq k\leq d$，选定 Floer 数据。考虑带有 Lagrangian 标记 $(L_0,\ldots,L_d)$ 的 $(d+1)$ 标点圆盘万有族 $\mathcal S^{d+1}\to\RG^{d+1}$，并为其配备 $\iota_M$ 不变的扰动数据 $(K,J)$，与之前选定的 Floer 数据相容。

\StatementNumber{7}\begin{lemma}\label{st:14.7}
令 $(r,u)\in\M^{d+1}(y_0,\ldots,y_d)=\M_{\mathcal S^{d+1}}(y_0,\ldots,y_d)$ 为非齐次伪全纯多边形，且 $u$ 的像完全包含在 $M^{\mathrm{fix}}$ 中。则 $u$ 在 $M$ 中正则，当且仅当它在 $M^{\mathrm{fix}}$ 中正则。
\end{lemma}
\begin{proof}
如前，须证明算子 $D_{\mathcal S^{d+1},r,u}$ 的反不变部分可逆。第一步仍将这一部分写为 $D^-_{S,u}:W^{1,p}(S,E,F)\to L^p(S,E)$，其中 $S=\mathcal S^{d+1}_r$、$E=S\times\CC$，且 $F|C_k=e^{i\sigma_k}\R$。在每个条带状端上，该算子的行为由 Floer 数据规定，即如\eref{14.12}有 $\epsilon_{S,k}^*D^-_{S,u}=\partial_s+i\partial_t+h_k$，其中 $h_k$ 满足相应的条件\eref{14.10}。随后结合给出指标为零的\stref{命题}{11.13}与\stref{引理}{11.5}即可。
\end{proof}

与\stref{引理}{14.6}之后的讨论一样，这说明一般选取 $\iota_M$ 不变扰动数据 $(K,J)$，就能使所有模空间 $\M_{\mathcal S^{d+1}}$ 正则。事实上，这次论证稍简单，因为 $(d+1)$ 标点圆盘没有自同构。

\StatementNumber{8}\begin{note}\label{st:14.8}
如\stref{笔记}{11.6}所述，可以证明\stref{引理}{11.5}的一个版本，对有关算子在条带状端上的行为无需如此精细的控制。这又允许削弱这里施加的某些条件，将其换成纯拓扑的对应条件，只要它们恰好足以保证反不变线性化算子的指标为零即可。沿另一个方向，可以考虑不动点集具有更高余维的 $\Z/2$ 作用；但那时必须保留对反不变算子的显式控制，事实上还须将这种控制扩展到整个 Riemann 曲面，因为次数论证已不再适用。
\end{note}

\subsection{分次与符号}\label{sec:14d}
仍处于\stref{假设}{14.3}的情形，但另要求 $\dim(M^{\mathrm{fix}})=\dim(M)-2>0$，以免为平凡情形另作例外规定。假设 $M$ 带有二次复体积形式 $\eta_M^2$，不必 $\iota_M$ 不变。由于法丛平凡，它在不动点集上诱导同类形式；用局部平方根表示，就是
\[
\eta_M|M^{\mathrm{fix}}=\pm\pi^*(dz)\wedge\eta_{M^{\mathrm{fix}}}.
\]
令 $L$ 为满足\stref{假设}{14.4}且与不动点集相交的不变正合 Lagrangian 子流形。$L$ 上任意膜结构 $(\alpha^\#,P^\#)$ 都在不动点集
\SourcePage{215}
$L^{\mathrm{fix}}\subset M^{\mathrm{fix}}$ 上诱导膜结构。其分次为 $\alpha^{\mathrm{fix},\#}=\alpha^\#|L^{\mathrm{fix}}-\sigma/\pi$，其中 $\sigma$ 为法向角。\znote{207}{分次的这一改变量来自法向方向。按本书的约定，膜结构的分次记录的是从标架丛到 $\lambda^{\mathrm{top}}$ 平凡化的相位；而由 $L^{\mathrm{fix}}$ 的标架补上法向量 $e^{i\sigma}$ 才得到 $L$ 的标架，这一操作使 Maslov 势移动旋转角 $\sigma$（以弧度计），除以 $\pi$ 后即为 $\sigma/\pi$。因此只有当 $\sigma$ 是 $\pi$ 的整数倍时，$L^{\mathrm{fix}}$ 才原样继承 $L$ 的分次；一般情况下必须作此平移，引理 14.9 中的指标与定向空间才与 $L$ 上的相应对象一致。[来源：原书 §(14d)，印刷页 207。]}对于 Pin 结构，取与切丛关联的标架丛 $P\to L$ 和 $P^{\mathrm{fix}}\to L^{\mathrm{fix}}$，将 $L^{\mathrm{fix}}$ 的标架 $(X_1,\ldots,X_{n-1})$ 映为 $(e^{i\sigma},X_1,\ldots,X_{n-1})$，从而把后一个丛嵌入前一个；第一项是法向量。沿该嵌入拉回二重覆盖 $P^\#\to P$，并赋予继承的 $\operatorname{Pin}_{n-1}\subset\operatorname{Pin}_n$ 作用，就得到所需的 Pin 结构 $P^{\mathrm{fix},\#}$。

现在令 $(L_0^\#,L_1^\#)$ 为两个这样的带膜结构 Lagrangian 子流形，$(H,J)$ 为\stref{假设}{14.5}中的 Floer 数据。如前所述，将其限制到不动点集，便定义 $(L_0^{\mathrm{fix}},L_1^{\mathrm{fix}})$ 的 Floer 数据。取某个完全位于 $M^{\mathrm{fix}}$ 中的 $y\in\mathcal C(L_0,L_1)$，将同一对象视为 $\mathcal C(L_0^{\mathrm{fix}},L_1^{\mathrm{fix}})$ 的元素时记为 $y^{\mathrm{fix}}$。不动点集带有上述诱导膜结构时，有

\StatementNumber{9}\begin{lemma}\label{st:14.9}
$i(y^{\mathrm{fix}})=i(y)$ 且 $o(y^{\mathrm{fix}})\cong o(y)$。
\end{lemma}
\begin{proof}
按照\eref{11.39}，把 $i(y)$ 和 $o(y)$ 视为一对线性 Lagrangian 膜 $\Lambda_{k,y}^\#$ 的定向算子的指标与行列式线；这些线性膜来自 $L_k^\#$ 在 $y$ 处的切空间。我们断言，可以把定向算子选成 $\Z/2$ 等变的，使不变部分为 $y^{\mathrm{fix}}$ 的定向算子，反不变部分则可逆。这与\stref{引理}{14.6}和\stref{引理}{14.7}的证明类似。为简化记号，只在 $H=0$ 时写出细节。此时 $y$ 是 $L_0$ 与 $L_1$ 的不变交点，可如下构造定向算子 $D_{H,y}$。取向量丛 $E=H\times TM_y$，赋予非常值的复结构 $J_E$，其在无穷远处由 Floer 数据确定；也就是说，若 $\epsilon_\zeta$ 是 $H$ 唯一的条带状端，则 $s\ll0$ 时有 $J_E(\epsilon_\zeta(s,t))=J_t$。在 $E$ 上赋予平凡联络 $\nabla$。子丛 $F=F_\lambda\subset E|\partial H$ 由 $\operatorname{Gr}(TM_y)$ 中连接 $\lambda(0)$ 与 $\lambda(1)$ 的路径 $\lambda$ 给出。最后，要使其成为真正的定向算子，还须为 $\lambda$ 选择膜结构 $(\alpha_\lambda^\#,P_\lambda^\#)$；这当然限制了路径的同伦类。

在该同伦类内，可以选择 $J_E$ 在 $D\iota_M|TM_y$ 下不变，并使路径 $\lambda$ 由不变 Lagrangian 子空间组成，于是 $D_{H,y}=\bar\partial_\nabla$ 为 $\Z/2$ 等变。此外，$\lambda$ 的反不变部分，即 $\nu_y\cong\CC$ 中的一族直线 $\lambda^-(s)$，可取为 $\lambda^-(s)=e^{i\sigma(s)}\R$，其中函数 $\sigma:[0,1]\to\R$ 从 $\sigma(0)=\sigma_0$ 走到 $\sigma(1)=\sigma_1$。这样，不变部分按照上面处理 Lagrangian 子流形的同一程序继承诱导膜结构，使 $D^+_{H,y}$ 成为 $y^{\mathrm{fix}}$ 的定向算子。另一方面，条件\eref{14.10}在此就是 $\sigma_1-\sigma_0\in(-\pi,0)$，它蕴含 $\operatorname{index}(D^-_{H,y})=0$。再由\stref{引理}{11.5}，$D^-_{H,y}$ 可逆。
\end{proof}

\SourcePage{216}
\StatementNumber{10}\begin{note}\label{st:14.10}
也许值得用另一种方式重述此论证：避开指标理论，改用\eref{11.25}中指标和定向空间的初等定义。具体地，选择 $\lambda$，使其除上述条件外还属于 $\mathcal P^-\operatorname{Gr}(V)$。于是指标与行列式线只依赖于交 $\lambda(s)\cap\lambda(1)$。然而，可以使函数 $\sigma$ 严格递减；此时 $\sigma_1-\sigma_0\in(-\pi,0)$ 意味着对所有 $s<1$，均有 $\lambda^-(s)\cap\lambda^-(1)=0$。因此，这些交全部集中在不变部分。
\end{note}

假设 Floer 数据正则。取 $u\in\M_Z^*(y_0,y_1)$，它位于 $M^{\mathrm{fix}}$ 内，且为模空间的孤立点。如\stref{引理}{14.6}的证明，将同一映射视为 $\M_Z^*(y_0^{\mathrm{fix}},y_1^{\mathrm{fix}})$ 的元素时，记作 $u^{\mathrm{fix}}$。接下来断言：在\stref{引理}{14.9}的同构下，$u$ 与 $u^{\mathrm{fix}}$ 对 Floer 微分的贡献具有相同符号。用\secref{12f}的记号，就是下图交换：
\SourceEquationDiagram{diagram-14-13}{0.22015361}{14.13}
为看出这一点，须回顾\eref{12.2}：$c_u$ 是把线性化算子 $D_{Z,u}$ 与定向算子 $D_{H,y_1}$ 粘合，再将所得算子变形为 $D_{H,y_0}$ 得到的。在此，这些算子全是 $\Z/2$ 等变的；其不变部分可与 $M^{\mathrm{fix}}$ 中对应的算子认同，而后者给出 $c_{u^{\mathrm{fix}}}$；其反不变部分可逆，因而对行列式线的贡献平凡。此外，从 $D_{Z,u}\#D_{H,y_1}$ 到 $D_{H,y_0}$ 的同伦经过的都是同样形式的算子，故反不变部分始终可逆。同一论证也适用于\stref{引理}{14.7}情形中不变孤立伪全纯多边形的贡献 $c_{r,u}$。

\subsection{等变 Fukaya 范畴}\label{sec:14e}
现在把讨论的两条线索结合起来。令 $M$ 为带角正合辛流形，配备满足\stref{假设}{14.3}的对合 $\iota_M$，并带有不变的二次复体积形式 $\eta_M^2$。

考虑 $M$ 上预等变正合 Lagrangian 膜，并附加如下条件：若其等变，则须满足\stref{假设}{14.4}；除非它们与不动点集不交，否则这还包含一个法向角的选择。为任意一对这样的膜选择 Floer 数据，要求其关于 $\Z/2$ 作用 $\epsilon$ 自然，即
\begin{equation}\label{eq:14.14}
(H_{L_{0,-}^\#,L_{1,-}^\#},J_{L_{0,-}^\#,L_{1,-}^\#})
=(\iota_M^*H_{L_{0,+}^\#,L_{1,+}^\#},\iota_M^*J_{L_{0,+}^\#,L_{1,+}^\#}).
\end{equation}
\SourcePage{217}
若所涉及的两个膜均等变，则意味着 Floer 数据须 $\iota_M$ 不变；在这种情况下，还按\stref{假设}{14.5}限制选择。也要求所有 Floer 数据正则。对等变膜对，\stref{引理}{14.6}之后已说明能够做到；在所有其他情况下也没有困难，因为\eref{14.14}联系的是两个不同的 Floer 数据，只要选择其中一个正则，另一个显然也有同样性质。Floer 数据选择中的对称性意味着存在典范同构
\begin{equation}\label{eq:14.15}
\epsilon:CF^*(L_{0,+}^\#,L_{1,+}^\#)\longrightarrow CF^*(L_{0,-}^\#,L_{1,-}^\#),
\end{equation}
它们与微分严格交换，且平方为恒等。随后，为对象的所有 $(d+1)$ 元组选择扰动数据，满足通常的一致性要求及\eref{14.14}的对应条件。由\stref{引理}{14.7}，可在这一类别内达到正则性，所得复合映射 $\mu^d$ 与\eref{14.15}严格交换。由此得到 $\Ai$ 范畴 $\F(M)^{\mathrm{naive}}$，称为\emph{等变 Fukaya 范畴}；其对象是满足上述条件的预等变膜，态射空间为 $CF^*(L_{0,+}^\#,L_{1,+}^\#)$。此 $\Ai$ 范畴带有由\eref{14.15}给出的朴素 $\Z/2$ 作用。为与之前的构造比较，转到上同调层次；显然存在与严格 $\Z/2$ 作用相容的全忠实嵌入 $H\F(M)^{\mathrm{naive}}\to H\F(M)^{\mathrm{strict}}$。事实上，这是一个等价：不满足\stref{假设}{14.4}的不变 Lagrangian 子流形，不能实现为 $H\F(M)^{\mathrm{naive}}$ 中的等变膜，但可作为预等变膜纳入，而在上同调层次这些对象同构。换言之，这样的子流形作为 $H\F(M)^{\mathrm{naive}}$ 的对象存在，但无法成为 $\Z/2$ 作用的不动点。这样做的好处是，$H\F(M)^{\mathrm{naive}}$ 确实是某个带朴素 $\Z/2$ 作用的 $\Ai$ 范畴的上同调范畴；$H\F(M)^{\mathrm{strict}}$ 则只是上同调层次的构造。\znote{209}{两者的差别只在链层次可见：$\F(M)^{\mathrm{naive}}$ 上的 $\Z/2$ 作用由链映射\eref{14.15}给出，是严格作用；而 $H\F(M)^{\mathrm{strict}}$ 只记住上同调上的作用，无法再提升到链层次。二者的上同调范畴等价，故就“分次范畴”而言无差别；但凡需要在上链层次作等变构造（取商、扭曲复形、等变 Floer 理论等），都必须使用朴素版本，否则会丢失高阶同伦信息。这正是不满足\stref{假设}{14.4}的不变 Lagrangian 子流形在 $H\F(M)^{\mathrm{naive}}$ 中不成为不动点的原因：它们只是同构于不动点，而同构并非链层次的相等。[来源：原书 §(14e)，印刷页 209。]}

为更详细地描述\eref{14.15}，取 $y_+\in\mathcal C(L_{0,+}^\#,L_{1,+}^\#)$，并考虑相关的一对线性 Lagrangian 膜 $\Lambda_{k,+}^\#$，即 $L_{k,+}^\#$ 在 $y_+(k)$ 处的切空间。这里把膜结构加入记号，因为 Hamilton 函数 $H$，因而集合 $\mathcal C$，现在依赖于这些结构。对 $\mathcal C(L_{0,-}^\#,L_{1,-}^\#)$ 中相应的元素 $y_-=\iota_M(y_+)$，由对象上 $\Z/2$ 作用的定义，有关的线性 Lagrangian 膜 $\Lambda_{k,-}^\#$ 是 $L_{k,-}^\#$ 的切空间。作为等变膜结构的一部分，有同构
\begin{equation}\label{eq:14.16}
(f_{k,-})_{y_-(k)}:\Lambda_{k,-}^\#\longrightarrow
\Lambda_{k,+}^\#\otimes\lambda^{\mathrm{top}}(TL_{k,+})_{y_+(k)},
\end{equation}
覆盖 $D\iota_M:(TL_{k,-})_{y_-(k)}\to(TL_{k,+})_{y_+(k)}$ 的作用。由\stref{笔记}{11.18}，有选定的同构 $\lambda^{\mathrm{top}}(TL_{k,+})_{y_+(k)}\cong\xi_{y_+(k)}$，还可沿 $y$ 平行移动，认同 $\xi_{y_+(0)}\cong\xi_{y_+(1)}$。将所涉及的两个线性 Lagrangian 膜同时与同一实直线作张量积，不改变绝对指标，也不改变定向空间，后者的“不改变”是在典范同构意义下而言。因此，\eref{14.16}
\SourcePage{218}
诱导同构 $o(y_-)\to o(y_+)$。将这些同构的逆作 $\K$ 规范化，就组成\eref{14.15}。这构成 $\Z/2$ 作用，归结于对\eref{14.2}的条件，以及下述事实：若用非平凡元素 $-e\in\ker(\operatorname{Pin}_n\to O_n)$ 同时作用于两个线性 Lagrangian 膜 $\Lambda_{k,+}^\#$，则在 $o(y)$ 上诱导的自同构是恒等。

一个也许值得特别关注的情形，是两个膜 $L_k^\#$ 都等变，且 $y\in\mathcal C(L_0^\#,L_1^\#)$ 在 $\iota_M$ 的作用下不变。此时，上述构造给出 $o(y)$ 的自同构，其符号由局部几何数据决定。与\stref{引理}{14.9}的证明一样，作简化假设 $H=0$，故 $y$ 为不变交点。回顾\stref{引理}{14.2}：同构 $f_k$ 为不动部分 $L_k^{\mathrm{fix}}$ 的法丛确定选定的平凡化。经 $D\pi$ 投影，这些平凡化给出直线 $e^{i\sigma_k}\R$ 的定向。若该定向与生成元 $e^{i\sigma_k}$ 所给定向一致，置 $w_k=+1$，否则置 $w_k=-1$。于是，

\StatementNumber{11}\begin{lemma}\label{st:14.11}
若 $w_0w_1=+1$，则 $|o(y)|_{\K}$ 属于 $CF^*(L_0^\#,L_1^\#)$ 的 $\epsilon$ 不变部分；否则属于反不变部分。
\end{lemma}
\begin{proof}
按照\stref{笔记}{14.10}的策略，选取从 $\lambda(0)=(TL_0)_y$ 到 $\lambda(1)=(TL_1)_y$ 的路径 $\lambda\in\mathcal P^-\operatorname{Gr}(V)$，完全由 $\Z/2$ 不变 Lagrangian 子空间组成。记 $F_\lambda$ 为 $[0,1]$ 上的关联向量丛。但这次要稍稍改变反不变直和项的同伦类：要求 $\lambda^-(s)=e^{i\sigma(s)}\R$，其中 $\sigma(0)=\sigma_0$、$\sigma(1)=\sigma_1+k\pi$，$k\in\Z$ 满足 $(-1)^k=w_0w_1$。为 $F_\lambda$ 选择膜结构，由实值相位 $\alpha_\lambda^\#$ 和 Pin 结构 $P_\lambda^\#$ 组成，在端点上限制为给定的线性 Lagrangian 膜。为实现\eref{14.15}，还要有同构
\[
f_\lambda:P_\lambda^\#\longrightarrow P_\lambda^\#\otimes\lambda^{\mathrm{top}}(F_\lambda),
\]
覆盖 $D\iota_M$ 在 $F_\lambda$ 上的作用。\stref{引理}{14.2}表明，这样的同构与 $\lambda^-$ 的定向一一对应。显然可调整 $f_\lambda$，使其在 $s=0$ 时与\eref{14.16}一致；追踪诱导的定向，并利用关于 $\lambda^-$ 同伦类的假设，便知在 $s=1$ 时也一致。因此，可在\eref{11.25}中考察 $D\iota_M$ 对\eref{11.20}的作用，读出 $o(y)$ 上诱导的自同构。在此，对每个 $s<1$ 时的非零交 $\lambda^-(s)\cap\lambda^-(1)$，都会产生一个 $(-1)$ 因子。这种点的个数的奇偶性由 $w_0w_1=(-1)^k$ 确定，即与 $k$ 相同，因而得到所述结论。等价地，也可如\stref{引理}{14.9}的证明，将此论证转回指标理论的语言。
\end{proof}

\chapter{Picard–Lefschetz 理论}\label{chapter:III}
\SourcePage{219}
本书所综合的第三项内容，是 Lefschetz 纤维化的辛几何。这是一个极其丰富的课题，我们的论述远不足以充分展现它。这里采用的定义基本与\bcite{131}相同。事实上，许多讨论都是对那篇论文的扩充或评述；这也意味着本书的处理有些偏于一面。有兴趣获得更全面认识的读者，可以从\bcite{40}、\bcite{64}、\bcite{11}、\bcite{17}入手。\znote{211}{Picard–Lefschetz 理论把全纯函数奇点的研究与其定义域复流形的拓扑联系起来，可视为复几何中与 Morse 理论平行的一套方法。其出发点是 Picard–Lefschetz 公式：临界点附近局部形如 $z_1^2+\cdots+z_n^2$，绕过临界值的环路诱导的单值变换，在消失圈上表现为一次 Dehn 扭转（相差一个符号）。由此，纤维化的单值变换可以完全用消失圈的相交数描述，这正是本章把 Lefschetz 纤维化问题化为组合数据的基础。[来源：nLab "Picard-Lefschetz theory"，https:/\allowbreak /\allowbreak ncatlab.org/\allowbreak nlab/\allowbreak show/\allowbreak Picard-Lefschetz+\allowbreak theory（2026-11-15 访问）。]}

前两节回顾一些初等方面。与\bcite{131}相比，主要增加的是 Donaldson 提出的匹配圈概念，另见\bcite{16}。粗略地说，若某个 Lefschetz 纤维化的一般纤维本身又是第二个 Lefschetz 纤维化的总空间，就会出现匹配圈。我们把这种结构称为双纤维化；在最低的非平凡维数中，它是\bcite{12}、\bcite{13}所考察的辛分歧覆盖的一种较简单的版本。\znote{211}{“一般纤维本身又是总空间”意味着有链 $E\to S\to W$，其中 $E\to S$ 与 $S\to W$ 都是带奇点的纤维化。此时固定 $w\in W$，$E_w\to S_w$ 又给出一个纤维化，而 $S_w$ 是曲面、$E_w$ 比它高两维。匹配圈就是在这个结构里构造的 Lagrangian 球面：在 $S_w$ 中取一条连接两个临界值的路径，把两端纤维中的消失圈沿路径平行移动并粘合。路径的选取与粘合方式不同，所得球面可能不同，这正是 Donaldson 的构造所利用的自由度。[来源：原书 §(15e) 双纤维化的定义，印刷页 218；匹配圈见 \bcite{16}。]}

随后，考虑以 Lefschetz 纤维化的全纯截面为基础的 Floer TQFT 扩展，并借此导出一个基本比较定理，联系第\sr{sec:5}{5 节}的代数理论与第\sr{sec:16}{16 节}的辛几何。这个定理在对象层次将扭转函子 $T$ 的作用与辛 Dehn 扭转 $\tau$ 的作用认同；事实上，它是\bcite{131}中长正合列的另一种表述。我们不重述该正合列的完整证明，但会较详细地解释，如何把结果变成这里所需的形式。

给定以圆盘为底空间的 Lefschetz 纤维化 $\pi:E\to D$，可以取沿某个纤维 $E_z$ 分歧的二重覆盖；代数几何中称之为“底扩张”。实际上，改变 $z$ 就得到一族这样的覆盖，这是双纤维化的一个特别简单的例子。二重覆盖显然有 $\Z/2$ 作用，可对它应用第\sr{sec:6}{6 节}和第\sr{sec:14}{14 节}的等变技术。我们将用这一想法定义 $\Ai$ 范畴 $\F(\pi)$，称为 \emph{Lefschetz 纤维化的 Fukaya 范畴}，它是纤维 $E_z$ 与总空间 $E$ 的 Fukaya 范畴之间的中介。诚然，借二重覆盖技巧给出的定义并非最自然，但它使大量已有工具得以重新使用。

接下来考察经典代数几何意义下的 Lefschetz 束，并证明引言中陈述的结果。最后一节讨论来自代数几何中一个稍有不同领域的一系列例子，即 $(A_m)$ 型简单奇点的 Milnor 纤维。它们首先具有历史意义，因为这是最早有人为导出 Fukaya 范畴提出显式代数模型的情形之一\bcite{85}；后来发现，那些模型只在实维数 $\geq4$ 时正确。到这里，讨论相较全书其余部分不再那么自足，但希望能借此展示具体计算中使用的几何与代数技巧。关于同一套一般工具的更复杂
\SourcePage{220}
应用，请读者参看文献，例如\bcite{14}、\bcite{15}、\bcite{130}、\bcite{59}。

\section{初步概念}\label{sec:15}
\subsection{带奇点的纤维化}\label{sec:15a}
我们将考察映射 $\pi:E\to B$，其中 $E$、$B$ 为带角正合辛流形，$B$ 连通，$\dim(E)\geq\dim(B)$，且 $\pi$ 为 $(I_E,I_B)$ 全纯映射。另作三个假设，它们共同保证映射沿边界具有合理的行为。
\begin{itemize}
\item \emph{与 $\partial B$ 横截。}令 $C\subset\partial B$ 为任意边界层。对每个满足 $y=\pi(x)\in C$ 的 $x\in E$，有 $D\pi(TE_x)+TC_y=TB_y$。
\end{itemize}
这显然蕴含 $\pi^{-1}(C)$ 为 $E$ 的边界层，余维与 $C$ 相同。所有这样的层之并为 $\partial^vE=\pi^{-1}(\partial B)$，称为 $E$ 的\emph{垂直边界}。所有不包含在 $\partial^vE$ 中的边界面之并，称为 $E$ 的\emph{水平边界}，记为 $\partial^hE$。
\begin{itemize}
\item \emph{沿 $\partial^hE$ 正则。}若 $F$ 是 $E$ 的一个不包含在 $\partial^vE$ 中的边界面，则 $\pi|F:F\to B$ 为光滑纤维化。
\end{itemize}
这蕴含每个纤维 $E_y$ 在其边界 $\partial E_y=E_y\cap\partial^hE$ 附近光滑。此前集中讨论微分拓扑，但还须处理辛方面的性质。在任意点 $x\in E$，切空间分解为
\begin{equation}\label{eq:15.1}
TE_x=TE_x^h\oplus TE_x^v,
\end{equation}
其中垂直部分 $TE_x^v=\ker(D\pi_x)$，水平部分 $TE_x^h$ 为其关于 $\omega_E$ 的正交补；等价地，也可取关于 $\omega_E$ 与 $I_E$ 所确定度量的正交补。我们在\secref{14c}中曾短暂遇到这样的分解，但在那里它只起技术作用。最后一项要求是
\begin{itemize}
\item \emph{$\partial^hE$ 的水平性。}若 $x$ 位于某个不包含在 $\partial^vE$ 中的边界面 $F$ 上，则 $TE_x^h\subset TF_x$。
\end{itemize}
满足以上三个条件的映射 $\pi:E\to B$，称为\emph{（正合辛）带奇点的纤维化}。

考察几个简单推论。记 $\operatorname{Crit}(\pi)\subset E$ 为临界点集，$\operatorname{Critv}(\pi)=\pi(\operatorname{Crit}(\pi))\subset B$ 为临界值集。如上所述，沿 $\partial^hE$ 正则意味着 $\operatorname{Crit}(\pi)\subset E\setminus\partial^hE$。当 $\dim(B)=2$ 时，与 $\partial B$ 横截还蕴含 $\operatorname{Crit}(\pi)\cap\partial^vE=\varnothing$；但底空间维数更高时就不再如此。

\SourcePage{221}
\StatementNumber{1}\begin{lemma}\label{st:15.1}
若 $y\notin\operatorname{Critv}(\pi)$，则纤维 $E_y$，配备限制得到的 $\omega_{E_y}=\omega_E|E_y$、$\theta_{E_y}=\theta_E|E_y$、$I_{E_y}=I_E|E_y$，是带角正合辛流形。
\end{lemma}
\begin{proof}
唯一需要非平凡验证的，是负 Liouville 向量场沿 $\partial E_y=E_y\cap\partial^hE$ 指向内部。已知 $E$ 的负 Liouville 向量场沿每个边界面 $F$ 严格向内。现假设 $F$ 不属于 $\partial^vE$。由于 $TE^h$ 平行于 $F$，该向量场的 $TE^v$ 分量也沿 $F$ 严格向内。而这一分量恰是 $E_y$ 的负 Liouville 向量场，从而得到结论。
\end{proof}

分解\eref{15.1}在 $B\setminus\operatorname{Critv}(\pi)$ 上定义一个联络。由于 $TE^h$ 平行于 $\partial^hE$，该联络具有良定义的平行移动映射；即对每条路径 $\beta:[0,1]\to B\setminus\operatorname{Critv}(\pi)$，得到一个选定的微分同胚
\begin{equation}\label{eq:15.2}
h_\beta:E_{\beta(0)}\longrightarrow E_{\beta(1)},
\end{equation}
它光滑依赖于路径的选择。标准计算表明，$h_\beta$ 实际上是带角正合辛流形的同构。固定 $\pi$ 的一个正则值 $*$，记相应纤维为 $M=E_*$，令 $B'\subset B\setminus\operatorname{Critv}(\pi)$ 为 $*$ 的一个小邻域。利用适当的一族平行移动映射，可得 $B'$ 上的局部平凡化 $E'=B'\times M\to E$。拉回 $\theta_E,\omega_E$ 得到的微分形式 $\theta_{E'},\omega_{E'}$ 具有形式
\begin{equation}\label{eq:15.3}
\theta_{E'}=\theta_M+\kappa+d(\text{某个函数}),\qquad
\omega_{E'}=\omega_M+d\kappa,
\end{equation}
其中一形式 $\kappa$ 在 $(TE')^v$ 方向上消失。可将其视为底空间上取值于函数的一形式
\[
\kappa\in\Omega^1(B',C^\infty(M,\R)).
\]
事实上，从函数转为 Hamilton 向量场之后，$\kappa$ 恰是 $(TE')^h$ 的联络一形式。\znote{213}{由\eref{15.3}，扭曲完全由 $\kappa$ 记录：$\omega_{E'}$ 与 $\omega_M$ 之差恰为 $d\kappa$；而 $\kappa$ 在垂直方向上消失，故 $d\kappa$ 在每条纤维上的限制为零，它只记录不同纤维之间的比较。$\kappa$ 闭（局部为恰当形式）时，平凡化就是正合辛同构；一般情形下 $d\kappa$ 度量水平分布与辛结构的失配。这也解释了为何下文把“沿纤维恒定的变形”实现为改变联络：只须指定新的 $\kappa$，而不改变每条纤维的辛形式与原始形式。[来源：原书 §(15a) 与 §(15b)，印刷页 213–215。]}为看出这一点，取 $B'$ 上向量场 $Y$，令 $Y^h=(Y,X)$ 为其到 $E'$ 的唯一水平提升。在 Cartan 公式 $i_{Y^h}(d\kappa)=\mathcal L_{Y^h}\kappa-d(i_{Y^h}\kappa)$ 中，Lie 导数仍是纤维方向上消失的一形式。因此，$Y^h$ 的定义性质可写为
\[
i_X\omega_M-d(i_Y\kappa)=0\qquad\text{在每个纤维 }\{y\}\times M\text{ 上}.
\]
所以，$X$ 在每个纤维上都是函数 $-\kappa(Y)$ 的 Hamilton 向量场，证明了所声称的 $\kappa$ 的性质。上述考察表明，在 $B$ 的正则部分上，我们得到的是带选定联络的 $\operatorname{Aut}(M)$ 纤维丛。注意，$B$ 上的辛结构在这里不起作用。

\StatementNumber{2}\begin{note}\label{st:15.2}
文献\bcite{131}对带奇点的纤维化的定义，把我们的第三个条件，即 $\partial^hE$ 的水平性，换成下述更强的版本：
\SourcePage{222}
\begin{itemize}
\item \emph{$\partial^hE$ 附近的平凡性。}令 $M=E_*$ 为某个纤维。则存在 $B\times\partial M$ 的开邻域 $V'\subset B\times M$、$\partial^hE$ 的开邻域 $V\subset E$，以及逐纤维的微分同胚 $V'\to V$，在 $V'\cap(\{*\}\times M)$ 上为恒等，使 $\theta_E$ 的拉回为 $\theta_B+\theta_M$，$I_E$ 的拉回为乘积殆复结构 $I_B\times I_M$。
\end{itemize}
对 $\theta_E$ 的条件意味着，联络 $TE^h$ 在 $\partial^hE$ 附近平凡，并且在底空间的正则部分上给出结构群到 $\operatorname{Aut}(M,\partial M)$ 的典范约化。这一性质使纤维化稍易处理，但也有自然构造不满足它。
\end{note}

\subsection{同构与变形}\label{sec:15b}
若两个带奇点的纤维化的总空间及底空间分别存在带角正合辛流形同构，并与纤维化映射相容，则称这两个纤维化同构。同样，显然可以定义这类纤维化的一般变形概念，即所涉及全部几何数据的单参数族。这里不展开一般定义，而集中于两类较特殊的变形；稍后陈述若干技术结果时会用到它们。第一类的主要特征是保持正则纤维的同构类，第二类则完全不改变纤维的结构。

固定流形 $E,B$ 及其间映射 $\pi$。在它们上取正合辛结构的光滑族 $(\omega_E^s,\theta_E^s,I_E^s)$、$(\omega_B^s,\theta_B^s,I_B^s)$，$s\in[0,1]$，使 $\pi$ 对每个 $s$ 都成为带奇点的纤维化。若存在 $E$ 上光滑函数 $R^s$，使由
\begin{equation}\label{eq:15.4}
\omega_E^s(\cdot,Y^s)=\frac{\partial}{\partial s}\theta_E^s-dR^s
\end{equation}
定义的向量场 $Y^s$ 平行于 $\partial^hE$，则称该变形\emph{沿 $\partial^hE$ 行为良好}。等价地，取乘积 $\widetilde\pi=\Id\times\pi:\widetilde E=[0,1]\times E\to\widetilde B=[0,1]\times B$，并在总空间上赋予 $\omega_{\widetilde E}=d(\theta_E^s+R^s\,ds)$。这显然在 $\ker(D\widetilde\pi)$ 上为辛，故得到如\eref{15.1}的 $T\widetilde E$ 分解；\eref{15.4}是说，该分解的水平部分 $T\widetilde E^h$ 平行于 $\partial^h\widetilde E=[0,1]\times\partial^hE$。这允许对 $[0,1]\times(B\setminus\operatorname{Critv}(\pi))$ 中的路径定义平行移动映射。特别地，沿 $[0,s]\times\{y\}$ 移动，得到正合辛同构
\begin{equation}\label{eq:15.5}
k_y^s:(E_y,\omega_{E_y}^0,\theta_{E_y}^0)\longrightarrow
(E_y,\omega_{E_y}^s,\theta_{E_y}^s).
\end{equation}
更具体地，它们是通过积分 $Y^s|E_y$ 关于 $\omega_E^s$ 的垂直分量定义的；若在 Moser 论证中，用 $\partial_s\theta_{E_y}^s-dR^s|E_y$ 作为 $\partial_s\omega_{E_y}^s$ 的原始形式，则使用的正是这个向量场。$k_y^s$ 光滑依赖于 $(s,y)$，并在 Hamilton 同痕意义下，与沿 $B\setminus\operatorname{Critv}(\pi)$ 中路径的平行移动交换。换言之，$E$ 的正则部分的纤维丛结构，在沿 $\partial^hE$ 行为良好的变形下不变。

\SourcePage{223}
另一种限制更强的概念，是\emph{沿纤维恒定的变形}。如前，这由单参数族 $(\omega_E^s,\theta_E^s,I_E^s)$ 和 $(\omega_B^s,\theta_B^s,I_B^s)$ 给出，但现在要求一形式 $\partial_s\theta_E^s-dR^s$ 在 $TE^v$ 上消失。显然，这意味着\eref{15.5}中的映射是恒等。此外，还要求在 $\partial^hE$ 的某个邻域内，$I_E^s|TE^v$ 与 $s$ 无关。

取带奇点的纤维化 $\pi:E\to S$，其中 $S$ 为带边界曲面。令 $\kappa$ 为 $E$ 上一形式，在 $TE^v$ 方向上消失，且在 $\partial E\cup\operatorname{Crit}(\pi)$ 的某个邻域中为零。对 $s\in[0,1]$ 及某个常数 $c\gg0$，考虑
\begin{equation}\label{eq:15.6}
\begin{aligned}
\theta_S^s&=(cs+1)\theta_S,&\omega_S^s&=d\theta_S^s,& I_S^s&=I_S,\\
\theta_E^s&=\theta_E+s\kappa+cs\pi^*\theta_S,&\omega_E^s&=d\theta_E^s.
\end{aligned}
\end{equation}
由于 $\omega_E^s|TE^v=\omega_E|TE^v$ 非退化，得到依赖于 $s$ 的分解 $TE_x=TE_x^{h,s}\oplus TE_x^v$，它是\eref{15.1}的变形。取 $c$ 充分大，可保证 $\omega_E^s|TE_x^{h,s}$ 关于底空间诱导的定向为正，从而 $\omega_E^s$ 为辛。取 $I_E^s$ 为唯一的 $\omega_E^s$ 相容殆复结构，在 $TE^v$ 方向等于 $I_E$，在 $TE^{h,s}$ 上则与 $I_S$ 的拉回一致。由于 $\partial E$ 附近有 $TE^{h,s}=TE^h$，$I_E^s$ 在该处与 $I_E$ 一致，故边界弱 $I_E^s$ 凸。其次，$\theta_E^s,\omega_E^s$ 关联的负 Liouville 向量场沿 $\partial^hE$ 严格向内，因为该性质只依赖于直和项 $TE^v$；见\stref{引理}{15.1}的证明。现在考虑边界圆周 $C\subset\partial S$ 及相应的 $F=\pi^{-1}(C)\subset\partial^vE$。对 $x\in F$，$TE_x^h$ 包含一维子空间 $TF_x\cap TE_x^h$；将其投影到 $S$ 并使用该处的边界定向，就得到典范定向。负 Liouville 向量场沿 $F$ 指向内部，等价于
\begin{equation}\label{eq:15.7}
\theta_E|(TF_x\cap TE_x^h)>0.
\end{equation}
由于 $\theta_S$ 本身使 $S$ 成为带角正合辛流形，有 $\pi^*\theta_S|(TF_x\cap TE_x^h)>0$。结合 $\partial^vE$ 附近的等式 $\theta_E^s=\theta_E+cs\pi^*\theta_S$，这意味着\eref{15.7}的对应条件在整个变形中持续成立。现在已经验证，$(E,\omega_E^s,\theta_E^s,I_E^s)$ 对每个 $s$ 都是带角正合辛流形，故\eref{15.6}是沿纤维恒定的变形。从纤维丛观点看，其效果是通过 $\kappa$ 改变辛联络。为稍后使用，记下一个具体应用：

\StatementNumber{3}\begin{lemma}\label{st:15.3}
令 $\pi:E\to S$ 为带奇点的纤维化，$S$ 为带边界曲面。令 $\beta:[0,1]\to\operatorname{int}(S)\setminus\operatorname{Critv}(\pi)$ 为嵌入路径，$(\phi^s)$ 为 $\operatorname{Aut}(E_{\beta(0)},\partial E_{\beta(0)})$ 中从 $\phi^0=\id$ 开始的同痕。则存在沿纤维恒定的变形，在 $s=0$ 时为给定纤维化结构，使沿 $\beta$ 所得的平行移动映射族满足
\[
h_\beta^s=h_\beta\circ\phi^s.
\]
\SourcePage{224}
此外，若开子集 $U\subset S$ 包含 $\beta$ 的某一部分，则可安排变形，使 $\partial_s(\theta_E^s-\pi^*\theta_S^s)$ 在 $\pi^{-1}(U)$ 之外消失。
\end{lemma}

\subsection{相对二次体积形式}\label{sec:15c}
令 $\pi:E\to B$ 为带奇点的纤维化。相对二次体积形式丛为
\[
\mathcal K_{E/B}^2=\pi^*\lambda_{\CC}^{\mathrm{top}}(TB)^{\otimes2}\otimes
\lambda_{\CC}^{\mathrm{top}}(TE)^{\otimes-2}.
\]
假设有该丛的处处非零截面 $\eta_{E/B}^2$；最简单的取得方式，是用总空间和底空间上的二次体积形式之商 $\eta_E^2/\eta_B^2$，但一般情形不必如此。若 $b$ 是 $\pi$ 的正则值，则短正合列 $0\to T(E_b)\to TE|E_b\to\pi^*TB_b\to0$ 诱导典范同构
\[
\mathcal K_{E/B}^2|E_b\cong\mathcal K_{E_b}^2.
\]
因此，$\eta_{E/B}^2$ 在光滑纤维上确定二次复体积形式 $\eta_{E_b}^2$。类似地，得到相对平方相位映射 $\alpha_{E/B}:\operatorname{Gr}(TE^v)\to S^1$，在临界点之外定义，其在每个光滑纤维上的限制与相应的 $\alpha_{E_b}$ 一致。在此情况下，平行移动映射\eref{15.2}取得自然的分次，即函数 $\gamma_{h_\beta}^\#:\operatorname{Gr}(E_{\beta(0)})\to\R$，使
\begin{equation}\label{eq:15.8}
\exp(2\pi i\,\gamma_{h_\beta}^\#(\Lambda))
=\frac{\alpha_{E_{\beta(1)}}(Dh_\beta(\Lambda))}{\alpha_{E_{\beta(0)}}(\Lambda)}.
\end{equation}
从 $\Id_{E_{\beta(0)}}$ 的零函数出发，对所有 $s$ 连续延伸到 $\beta|[0,s]$，便得到这些函数。在 $\beta$ 为闭路径的特殊情形，结果是 $h_\beta$ 成为\secref{12i}所定义的分次辛自同构。\znote{216}{相对二次体积形式的作用是为纤维化提供“沿纤维的典范数据”。因为纤维的切空间由正合列 $0\to T(E_b)\to TE|E_b\to\pi^*TB_b\to0$ 给出，$\mathcal K_{E/B}^2$ 限制到每条光滑纤维上就是 $\mathcal K_{E_b}^2$；所以一个相对截面同时为所有光滑纤维指定二次复体积形式，不必逐条纤维选择。相位 $\alpha_{E/B}$ 是 $\mathcal K_{E/B}$ 的平方根在 Lagrangian Grassmann 流形上的相位，它使平行移动能够带上分次：\eref{15.8}把端点相位之比的对数取为 $\gamma^\#_{h_\beta}$，沿路径连续延拓即得，而闭路径情形下所得正是分次辛自同构。若没有相对数据，就只能逐条纤维比较，无法保证与平行移动相容。[来源：原书 §(15c)，印刷页 216。]}

\subsection{Lefschetz 纤维化}\label{sec:15d}
令 $S$ 为具有非空边界的连通紧 Riemann 曲面。选择一形式 $\theta_S$，满足 $\theta_S|\partial S>0$ 及 $\omega_S=d\theta_S>0$，就使其成为带角正合辛流形。\emph{（正合）Lefschetz 纤维化}是带奇点的纤维化 $\pi:E^{2n+2}\to S$，其临界点一般，也称为非退化，且局部可积。这意味着 $I_E$ 在 $\operatorname{Crit}(\pi)$ 的某个邻域中可积，且 $D\pi$ 作为复线性映射丛 $\Hom_{\CC}(TE,\pi^*TS)$ 的截面与零截面横截。第二个条件等价于每个临界点处的复 Hessian $D^2\pi$ 非退化。此外，还假设每个纤维中至多有一个临界点，故投影 $\operatorname{Crit}(\pi)\to\operatorname{Critv}(\pi)$ 为双射；最后这一假设只是为方便起见，很容易去除，代价是后面的记号稍显繁琐。

非退化性蕴含临界点孤立，故 $\operatorname{Crit}(\pi)$ 是 $\operatorname{int}(E)$ 的有限子集，类似地，$\operatorname{Critv}(\pi)$ 是 $\operatorname{int}(S)$ 的有限子集。在每个
\SourcePage{225}
临界点及其临界值附近，存在局部全纯坐标，使 $\pi$ 成为标准二次映射
\begin{equation}\label{eq:15.9}
Q(x)=x_1^2+\cdots+x_{n+1}^2.
\end{equation}
一般而言，$\omega_E$ 在这些坐标中并不标准。不过，可以找到沿 $\partial^hE$ 行为良好的纤维化变形，事实上它局限在临界点附近，使在变形的另一端，Kähler 形式在给定的全纯 Morse 坐标中成为标准形式。这使下述基本模型更为重要：

\StatementNumber{4}\begin{example}\label{st:15.4}
令 $Q:\CC^{n+1}\to\CC$ 如\eref{15.9}，并令 $k:\CC^{n+1}\to\R^+$ 为函数 $k(x)=(|x|^4-|Q(x)|^2)/4$。对固定的 $r,s>0$，定义
\begin{equation}\label{eq:15.10}
E=\{x\in\CC^{n+1}:|Q(x)|\leq r,\ k(x)\leq s\},
\end{equation}
赋予 $\CC^{n+1}$ 上标准辛形式的限制、它的标准原始形式 $(i/4)(z\,d\bar z-\bar z\,dz)$、给定复结构 $I_E=i$，以及限制 $Q$ 所得的映射 $\pi:E\to S=rD$。边界面为
\[
\partial^vE=\{x\in E:|\pi(x)|=r\},\qquad
\partial^hE=\{x\in E:k(x)=s\}.
\]
截断函数 $k$ 的选取使 $TE^h$ 平行于 $\partial^hE$。为验证这一点，注意 $TE_x^h$ 在 $\CC$ 上由 $(\nabla Q)_x=2\bar x$ 生成，再核对 $dk_x(\bar x)=0$、$dk_x(i\bar x)=0$ 即可。每个非奇异纤维 $E_z$，$z\ne0$，都与余切向量在标准度量下长度 $\leq\sqrt{s}$ 的子集 $B_{\sqrt{s}}^*S^n\subset T^*S^n$ 辛同构。明确地，
\[
B_{\sqrt{s}}^*S^n=\{(u,v)\in\R^{n+1}\times S^n:
\langle u,v\rangle=0,\ |u|\leq\sqrt{s}\},
\]
带辛形式 $du\wedge dv$；当 $z>0$ 时，同构 $E_z\to B_{\sqrt{s}}^*S^n$ 由下式给出：
\[
\phi_z(x)=(-\operatorname{im}(x)|\operatorname{re}(x)|,
\operatorname{re}(x)|\operatorname{re}(x)|^{-1}).
\]
遗憾的是，尽管负 Liouville 向量场，即负径向向量场，确实沿 $\partial E$ 指向内部，$E$ 却并非弱 $I_E$ 凸；纤维是弱凸的，总空间则不是。因此，这还不完全是这里定义的正合 Lefschetz 纤维化的例子，虽然从纯辛的观点看，它具有全部所需性质。当然，可以改变 $I_E$ 来改善这种情况，但实际上没有必要，因为最终 $E$ 只用作局部模型。
\end{example}

\StatementNumber{5}\begin{note}\label{st:15.5}
辛 Lefschetz 纤维化的其他定义，例如见\bcite{39}或\bcite{64}。通常假设底空间和纤维均闭，这在技术上略为简单。局部标准形\eref{15.9}的存在是全纯 Morse 引理\bcite{8}的推论；更准确地说，对底空间上任意选取的全纯坐标，都可在总空间上找到坐标，使 $\pi=Q$。使这些坐标中的辛结构标准化的变形，构造于\bcite{131}引理 1.6。\stref{例题}{15.4}的几何在\bcite{131}第 1.2 节中有详尽讨论。
\end{note}

\SourcePage{226}
\subsection{双纤维化}\label{sec:15e}
现在考察形如下图的结构：
\SourceEquationDiagram{diagram-15-11}{0.2388459}{15.11}
这里 $\Psi$ 为 Lefschetz 纤维化，总空间维数为 $2n+4\geq6$。中间空间 $\mathcal S$ 为带角正合辛四维流形，$\varpi$ 与 $\psi$ 均为带奇点的纤维化。假设 $\psi$ 没有临界点，且其纤维为带边界的连通曲面；特别地，$\mathcal S$ 只有余维 $\leq2$ 的角。对 $\varpi$，只要求它在 $\partial^h\mathcal S$ 上没有临界值，此处这部分边界是关于 $\psi$ 定义的。\znote{218}{“双”纤维化指两个纤维化叠在一起：$\varpi:\mathcal E\to\mathcal S$ 与 $\psi:\mathcal S\to\mathcal W$ 都以曲面为底，复合 $\Psi=\psi\circ\varpi$ 是 Lefschetz 纤维化。固定 $w$ 后，$\mathcal E_w\to\mathcal S_w$ 本身又是以曲面为底的纤维化；当 $w$ 取 $\Psi$ 的临界值时，$\mathcal E_w$ 上恰有一个临界点，这一纤维化随之退化。这样，关于 $2n+4$ 维总空间的问题便可化归为曲面 $\mathcal S_w$ 上的曲线配置：消失圈与匹配路径都生活在曲面里，而 $\Psi$ 的单值变换由 $\operatorname{Critv}(\varpi)$ 在曲面间的运动描述。[来源：原书 §(15e)，印刷页 218。]}取点 $w\in\mathcal W$，将 $\varpi$ 限制到该点上方的纤维，记所得映射为
\begin{equation}\label{eq:15.12}
\mathcal E_w\xrightarrow{\varpi_w}\mathcal S_w.
\end{equation}
除非 $w$ 是 $\Psi$ 的临界值，否则它就是带奇点的纤维化；其余必要性质均由关于 $\varpi,\psi,\Psi$ 的假设推出。因此，可将 $\mathcal E\to\mathcal W$ 视为一个 Lefschetz 纤维化，其光滑纤维又在曲面 $\mathcal S_w$ 上纤维化。若还满足下述附加条件，就称\eref{15.11}为\emph{（正合）Lefschetz 双纤维化}：
\begin{itemize}
\item $\varpi$ 的临界点一般，即 $D\varpi$ 在 $\Hom_{\CC}(T\mathcal E,\varpi^*T\mathcal S)$ 中，与零截面以及复秩为一的线性映射所成的层均横截。若 $x$ 是 $\varpi$ 的临界点但为 $\Psi$ 的正则点，则 $I_{\mathcal E}$ 在纤维 $\mathcal E_w$，$w=\Psi(x)$，上的限制在 $x$ 附近局部可积。此外，对每个 $w\in\partial\mathcal W$，\eref{15.12}均为 Lefschetz 纤维化。
\item 取 $w\in\operatorname{Critv}(\Psi)$，令 $x$ 为 $\mathcal E_w$ 中唯一临界点。要求 $D^2\Psi_x$ 在 $\ker(D\varpi_x)\subset T\mathcal E_x$ 上的限制非退化。殆复结构 $I_{\mathcal S}$ 在 $\varpi(x)$ 的某个邻域中可积。此外，$\mathcal E_w\setminus\{x\}$ 中 $\varpi_w$ 的所有临界点均须非退化，其 $\varpi$ 值两两不同，且均不同于 $\varpi(x)$。
\end{itemize}
下面逐步解释这些条件的含义。一般性与维数计算表明，$D\varpi$ 从不为零，其复秩为一的点所成的集合 $\operatorname{Crit}(\varpi)$ 是光滑二维子流形。事实上，

\StatementNumber{6}\begin{lemma}\label{st:15.6}
$\operatorname{Crit}(\varpi)$ 是光滑的 $I_{\mathcal E}$ 全纯曲线。
\end{lemma}
\begin{proof}
对 $x\in\operatorname{Crit}(\varpi)$ 及 $X\in T\mathcal E_x$，考虑部分二阶导数
\begin{equation}\label{eq:15.13}
\bar D^2\varpi_{x,X}:\ker(D\varpi_x)\longrightarrow\operatorname{coker}(D\varpi_x).
\end{equation}
$\operatorname{Crit}(\varpi)$ 的切空间由使它为零的 $X$ 组成。局部坐标中的简单计算表明，$\bar D^2\varpi_{x,I_{\mathcal E}X}=I_{\mathcal S}\bar D^2\varpi_{x,X}$。
\end{proof}

\SourcePage{227}
\StatementNumber{7}\begin{lemma}\label{st:15.7}
令 $x\in\operatorname{Crit}(\varpi)\setminus\operatorname{Crit}(\Psi)$，且 $\Psi(x)=w$。则 $x$ 是 $\varpi_w$ 的临界点；此外，交 $\operatorname{Crit}(\varpi)\cap\mathcal E_w$ 在 $x$ 处横截，当且仅当 $x$ 作为 $\varpi_w$ 的临界点非退化。
\end{lemma}
\begin{proof}
由于 $x$ 不是 $\Psi$ 的临界点，$D\varpi_x$ 的像与 $\mathcal S_w$ 的切空间横截。这蕴含 $\ker(D\varpi_x)=T(\mathcal E_w)_x$，且包含映射诱导的 $T(\mathcal S_w)_{\varpi(x)}\to\operatorname{coker}(D\varpi_x)$ 为同构。从而有 $D(\varpi_w)_x=0$，并且对 $X\in T(\mathcal E_w)_x$，可将\eref{15.13}与通常的二阶导数 $D^2(\varpi_w)_{x,X}$ 认同。若交 $\operatorname{Crit}(\varpi)\cap\mathcal E_w$ 横截，则对任何非零 $X\in T(\mathcal E_w)_x$，\eref{15.13}必非零，因而意味着该临界点非退化。同一论证也给出逆命题。
\end{proof}

\StatementNumber{8}\begin{lemma}\label{st:15.8}
在 $\Psi$ 的每个临界点 $x$ 处，映射 $\varpi|\operatorname{Crit}(\varpi):\operatorname{Crit}(\varpi)\to\mathcal S$ 为局部嵌入，而 $\Psi|\operatorname{Crit}(\varpi):\operatorname{Crit}(\varpi)\to\mathcal W$ 具有一个二阶分歧点。
\end{lemma}
\begin{proof}
由于 $D\varpi_x$ 秩为一，而 $D\Psi_x$ 为零，$D\psi_{\varpi(x)}$ 将 $\operatorname{coker}(D\varpi_x)$ 同构映到 $T\mathcal W_w$，其中 $w=\Psi(x)$。在此认同下，\eref{15.13}成为 $\Psi$ 在 $x$ 处的通常二阶导数；这意味着 $\operatorname{Crit}(\varpi)$ 的切空间在 $\CC$ 上由一个非零 $X$ 生成，满足对所有 $Y\in\ker(D\varpi_x)$ 有 $D^2\Psi_x(X,Y)=0$。已假设 $D^2\Psi_x$ 在 $\ker(D\varpi_x)$ 上非退化，故 $D\varpi_x(X)\ne0$，证明了第一项断言。非退化性还蕴含 $D^2\Psi_x(X,X)\ne0$，说明 $\Psi|\operatorname{Crit}(\varpi)$ 在 $x$ 处的二阶导数非零。
\end{proof}

\StatementNumber{9}\begin{lemma}\label{st:15.9}
在 $\Psi$ 的每个临界点及其在 $\mathcal S,\mathcal W$ 中的值附近，存在局部全纯坐标，使
\[
\varpi(x_1,\ldots,x_{n+2})=(x_1^2+x_2^2+\cdots+x_{n+2}^2,x_1),\qquad
\psi(y_1,y_2)=y_1.
\]
\end{lemma}
\begin{proof}
这是带参数全纯 Morse 引理的推论，该引理的证明与不带参数的情形相同。先在 $\mathcal S$ 与 $\mathcal W$ 上，以临界点的 $\varpi$ 值和 $\Psi$ 值为中心引入局部全纯坐标，使 $\psi$ 具有所需形式 $\psi(y_1,y_2)=y_1$；按假设，$I_{\mathcal S}$ 局部可积，故可做到。在这些坐标中局部写成 $\varpi=(\varpi_1,\varpi_2):\mathcal E\to\CC^2$，会发现 $D\varpi_2$ 在临界点非零，故水平集 $\{\varpi_2=0\}$ 光滑，且 $\varpi_1=\Psi$ 在此水平集上的限制在 $x$ 处具有非退化临界点。因此，附近的水平集 $\{\varpi_2=y_2\}$ 也光滑，$\varpi_1$ 限制到每一个水平集均有唯一非退化临界点，并全纯依赖于 $y_2$。于是可在每个水平集上找到全纯依赖于 $y_2$ 的局部坐标 $x_2,\ldots,x_{n+2}$，使 $\varpi_1$ 的限制具有通常的二次标准形。加上 $x_1=\varpi_2$，就得到 $\mathcal E$ 上一组坐标，使 $\varpi(x_1,\ldots,x_{n+2})=(f(x_1)+x_2^2+\cdots+x_{n+2}^2,x_1)$，
\SourcePage{228}
其中 $f$ 为某个全纯函数。由假设，$x=0$ 是 $\varpi_1=\Psi$ 的非退化临界点，故对某个 $c\ne0$，有 $f(t)=ct^2+O(t^3)$。最后，将 $x_1$ 换成 $f^{1/2}(x_1)$，并将 $y_2$ 换成 $f^{1/2}(y_2)$ 即可。
\end{proof}

现在讨论这些技术结果对于临界曲线及其像的几何意味着什么：
\SourceDiagram{diagram-15-critical}{0.43615338}
由于假设 $\operatorname{Critv}(\varpi)\cap\partial^h\mathcal S=\varnothing$，$\operatorname{Crit}(\varpi)$ 的边界等于它与 $\Psi^{-1}(\partial\mathcal W)$ 的交；又因 $\omega_{\mathcal E}$ 正合，$\operatorname{Crit}(\varpi)$ 的每个连通分支都必须有边界。这意味着没有分支能包含在 $\Psi$ 的某个纤维中，因此 $\Psi|\operatorname{Crit}(\varpi)$ 是有限分歧覆盖。由\stref{引理}{15.7}，对 $w\in\partial\mathcal W$ 要求 $\varpi_w$ 是 Lefschetz 纤维化，意味着在 $\partial\operatorname{Crit}(\varpi)$ 的某个邻域中，映射 $\varpi|\operatorname{Crit}(\varpi)$ 为嵌入，$\Psi|\operatorname{Crit}(\varpi)$ 为局部同构。应用一些基本伪全纯曲线理论，可见 $\varpi|\operatorname{Crit}(\varpi)$ 几乎处处单射，故其像 $\operatorname{Critv}(\varpi)$ 是没有多重覆盖分支的 $I_{\mathcal S}$ 全纯曲线。此外，曲线在边界 $\partial\operatorname{Critv}(\varpi)=\operatorname{Critv}(\varpi)\cap\partial^v\mathcal S$ 附近光滑，且 $\psi|\operatorname{Critv}(\varpi)$ 为局部同构。结合\stref{引理}{15.7}与\stref{引理}{15.8}可知，对 $w\in\operatorname{Critv}(\Psi)$，$\operatorname{Critv}(\varpi)$ 在 $\mathcal S_w$ 的某个邻域中也光滑，且每个这样的 $\mathcal S_w$ 包含 $\psi|\operatorname{Critv}(\varpi)$ 的唯一分歧点，该点还是二阶的。\sr{fig:15.1}{图~15.1}概括了这一情形。
\SourceFigure{figure-15-1}{0.53376866}{15.1}

定义\emph{伪临界值集} $\operatorname{Fakev}(\Psi,\varpi)$：它由那些 $w\in\mathcal W\setminus\operatorname{Critv}(\Psi)$ 组成，使 $\operatorname{Critv}(\varpi)$ 在 $\mathcal S_w$ 上具有奇点，或在 $\mathcal S_w$ 附近光滑但与其不横截相交。等价地，这些正是 $\mathcal E_w$ 光滑而\eref{15.12}不是 Lefschetz 纤维化的 $w$。由前面的讨论，$\operatorname{Fakev}(\Psi,\varpi)$ 是 $\operatorname{int}(\mathcal W)$ 的有限子集；特别地，$\mathcal W\setminus(\operatorname{Critv}(\Psi)\cup\operatorname{Fakev}(\Psi,\varpi))$ 仍连通。\znote{220}{伪临界值集是这一理论中的“例外集”：在 $w\notin\operatorname{Critv}(\Psi)$ 处 $\mathcal E_w$ 光滑，但 $\mathcal E_w\to\mathcal S_w$ 可能并非 Lefschetz 纤维化，因为 $\varpi_w$ 的临界点可能退化，或者 $\operatorname{Critv}(\varpi)$ 与 $\mathcal S_w$ 相切。前面关于临界曲线的讨论说明这样的 $w$ 只有有限多个：$\operatorname{Critv}(\varpi)$ 是无多重覆盖分支的全纯曲线，只在 $\Psi$ 的临界值处出现二阶分歧点。故从 $\mathcal W$ 中挖去有限多个点后，消失圈、匹配路径与单值变换的讨论全部可用；后文对一般 $w$ 陈述结论正是依据这一点。[来源：原书 §(15e)，印刷页 220。]}

\clearpage
\section{消失圈与匹配圈}\label{sec:16}
\subsection{标架}\label{sec:16a}
令 $(M^{2n},\omega)$ 为任意辛流形，$V\subset M$ 为 Lagrangian 球面。$V$ 的\emph{标架}是一个微分同胚 $v:S\to V$，其中 $S$ 为某个 $(n+1)$ 维内积空间的单位球面；$S$ 依赖于 $V$。注意，这与例如 Pontryagin–Thom 构造中带标架流形的概念不同。类似地，若有两个带标架 Lagrangian 球面 $v_k:S_k\to V_k$，$k=0,1$，则它们之间的\emph{带标架同痕} $I=(i,j_0,j_1)$
\SourcePage{229}
由一族 Lagrangian 嵌入 $i^s:S\to M$，$0\leq s\leq1$，以及两个等距映射 $j_k:S\to S_k$ 组成，使 $v_k\circ j_k=i^k$；这里 $S$ 仍是某个内积空间的单位球面。

\StatementNumber{1}\begin{note}\label{st:16.1}
假设 $V_0=V_1=V$ 且 $S_0=S_1=S$ 相同，但有两个不同映射 $v_0,v_1$。显然，若 $v_k$ 在微分同胚空间 $\operatorname{Diff}(S,V)$ 中同痕，则两个带标架 Lagrangian 球面同痕。另一方面，若 $v_k$ 相差 $S$ 的某个反定向等距映射，则令所有 $s$ 的 $i^s=v_0$，并取 $j_0=\Id$、$j_1=v_1^{-1}\circ v_0$，就得到带标架 Lagrangian 球面的同痕。因此，只要关心的仅是同痕类，而我们使用这一概念时基本都是如此，标架信息就归结为 $\pi_0(\operatorname{Diff}(S,V)/O(S))$ 中的一个类。已知 $n<3$ 时 $\operatorname{Diff}(S^n)/O_{n+1}$ 可缩，故这些维数中的标架并不重要。但对较大的 $n$，奇异球面的存在意味着 $\pi_0(\operatorname{Diff}(S^n)/O_{n+1})$ 往往非平凡。\znote{221}{可缩只在 $n<3$ 成立，不要推广成所有维数或写成 $n\leq3$。可缩意味着任意两个标架都同痕，标架信息于是完全消失；$n\geq3$ 时 $\pi_0$ 不再平凡，奇异球面（与 $S^n$ 同胚但不微分同胚的球面）正是其非平凡性的来源。因此后文凡出现“带标架”之处，在 $n\geq3$ 都不能随手丢弃标架，只有在 $n<3$ 才可以。[来源：原书 Remark 16.1，印刷页 221。]}
\end{note}

\subsection{消失圈}\label{sec:16b}
令 $\pi:E^{2n+2}\to S$ 为 Lefschetz 纤维化，仍照例假设正合。\emph{消失路径}是终止于一个临界值的嵌入路径 $\gamma:[0,1]\to S$；更确切地，要求 $\gamma^{-1}(\operatorname{Critv}(\pi))=\{1\}$。对每条这样的路径，可赋予其 \emph{Lefschetz 顶针} $\Delta_\gamma$，即 $E$ 中唯一的嵌入 Lagrangian $(n+1)$ 维球体，满足 $\pi(\Delta_\gamma)=\gamma([0,1])$、$\pi(\partial\Delta_\gamma)=\{\gamma(0)\}$。边界 $V_\gamma=\partial\Delta_\gamma$ 是 $E_{\gamma(0)}$ 中的 Lagrangian 球面，称为 $\gamma$ 的\emph{消失圈}。由于它是 $E$ 中一个 Lagrangian 圆盘的边界，任何消失圈都自动正合；当然，这仅在 $n=1$ 时非平凡。\znote{221}{三点提醒：一、顶针 $\Delta_\gamma$ 是 $(n+1)$ 维球体（实维数），不是球面；其边界 $V_\gamma=\partial\Delta_\gamma$ 才是 $n$ 维球面，二者维数差一。二、$\pi|\Delta_\gamma$ 是到 $\gamma([0,1])$ 的同胚，故 $\pi(\partial\Delta_\gamma)=\{\gamma(0)\}$ 说明整块边界落在起始纤维 $E_{\gamma(0)}$ 内，这正是“消失”一词的几何内容：绕临界值一周，圈便收缩成一点。三、“正合”不是额外假设：由 Stokes 公式，$\theta|_{V_\gamma}=df$，其中 $f$ 是 $\Delta_\gamma$ 上的辛面积函数；$n>1$ 时 $V_\gamma$ 单连通，故并无实质内容。[来源：原书 §(16b)，印刷页 221–222。]}继续之前，记下一个初等事实：

\SourcePage{230}
\StatementNumber{2}\begin{lemma}\label{st:16.2}
对任意 Lagrangian 子流形 $F\subset E$，点 $x\in F$ 是 $\pi$ 的临界点，当且仅当它是 $\pi|F$ 的临界点。
\end{lemma}

对于具有上述性质的圆盘 $\Delta_\gamma$，函数 $p=\gamma^{-1}\circ\pi|\Delta_\gamma:\Delta_\gamma\to[0,1]$ 必在某个内点达到极大值；由\stref{引理}{16.2}，该点必是 $E_{\gamma(1)}$ 中唯一的奇点 $x_1$。类似论证表明，$p$ 在该点的 Hessian 非退化，故 $p$ 是只有一个临界点的 Morse 函数。在 $\Delta_\gamma$ 上赋予一个度量，使其在临界点附近的局部 Morse 坐标中标准化，便可用 $\nabla p$ 的流给出微分同胚
\begin{equation}\label{eq:16.1}
B=B(T_{x_1}\Delta_\gamma)\longrightarrow\Delta_\gamma,
\end{equation}
从 $x_1$ 处切空间的单位球映到整个 Lefschetz 顶针。特别地，限制到边界便得到标架 $v_\gamma:S=\partial B\to V_\gamma$，其在同痕意义下不依赖于度量的选择。将 $v_\gamma$ 视为消失圈结构的一部分，因而消失圈总被视为带标架正合 Lagrangian 球面。

我们从 Lefschetz 顶针的定义性质推出了一些结论，但其存在唯一性仍有待验证。现在来做这件事，先从唯一性开始。以下事实很有用：

\clearpage
\StatementNumber{3}\begin{lemma}\label{st:16.3}
令 $\beta$ 为底空间 $S$ 正则部分中的嵌入路径，$F\subset E$ 为在 $\beta$ 上纤维化的子流形，且每个 $F_{\beta(s)}\subset E_{\beta(s)}$ 都为 Lagrangian。则整个 $F$ 为 Lagrangian，当且仅当沿 $\beta$ 任意一段的平行移动把相应的 $F_{\beta(s)}$ 互相映到彼此。
\end{lemma}

对于 Lefschetz 顶针，这意味着沿 $\gamma$ 的平行移动把水平集 $\Delta_\gamma\cap E_{\gamma(t)}=p^{-1}(t)$，$0\leq t<1$，互相映到彼此。当 $t\to1$ 时，这些水平集收缩到 $\{x_1\}$。因此，$\Delta_\gamma$ 必包含在下述集合中：
\[
\Delta'_\gamma=\{x\in E_{\gamma(t_0)},\ 0\leq t_0<1:
\lim_{t_1\to1}h_{\gamma|[t_0,t_1]}(x)=x_1\}\cup\{x_1\}.
\]
$\Delta'_\gamma$ 的结构未必一眼可见，因为平行移动向量场在 $x_1$ 有极点。为避开这一困难，选择光滑函数 $h:S\to\R$，使 $h(\gamma(t))=0$，且对任意 $t$，$(\gamma'(t),\nabla h_{\gamma(t)})$ 为正定向切向量基。令 $X_H$ 为 $H=h\circ\pi$ 的 Hamilton 向量场。容易看出，它在 $\pi^{-1}(\gamma([0,1)))$ 上与平行移动向量场指向相同方向，同时在 $x_1$ 具有驻点。故 $\Delta'_\gamma$ 包含在
\begin{equation}\label{eq:16.2}
\Delta''_\gamma=\{x\in E_{\gamma(t)},\ 0\leq t\leq1:
\lim_{s\to\infty}\phi_H^s(x)=x_1\}.
\end{equation}
\SourcePage{231}
由于驻点 $x_1$ 双曲，可用稳定流形定理证明，$\Delta''_\gamma$ 是嵌入的 $(n+1)$ 维球体，边界为球面 $\Delta''_\gamma\cap E_{\gamma(0)}$。按假设 $\Delta_\gamma$ 也具有同样性质，所以包含关系 $\Delta_\gamma\subset\Delta'_\gamma\subset\Delta''_\gamma$ 必为等式，证明了 Lefschetz 顶针的唯一性。对于存在性，只须注意\eref{16.2}为 Lagrangian，因为辛流 $\phi_H$ 在其上收缩，且它具有其余全部所需性质。\znote{223}{这是“夹逼”式论证，两处跳步值得点明。其一，$h$ 的选取使 $X_H$ 与平行移动同向，才有 $\Delta'_\gamma\subset\Delta''_\gamma$。其二，$\Delta''_\gamma$ 是嵌入球体、且其边界落在 $E_{\gamma(0)}$ 内，靠的是 $x_1$ 双曲（Hessian 非退化）时稳定流形定理的标准结论，不能只从 $\Delta'_\gamma$ 的定义读出；$x_1$ 处平行移动向量场有极点，正是这一步必须换用 Hamilton 流的原因。存在性反过来取 $\Delta''_\gamma$ 为候选：$\phi_H$ 在其上收缩，故 $\omega$ 沿流的切向消失，从而得到 Lagrangian。[来源：原书 §(16b)，印刷页 222。]}

\StatementNumber{4}\begin{note}\label{st:16.4}
作为上述讨论的副产物，任意一段 $\gamma|[t_0,1]$ 的消失圈与整个路径的消失圈通过平行移动联系：
\begin{equation}\label{eq:16.3}
V_\gamma=h_{\gamma|[0,t_0]}^{-1}(V_{\gamma|[t_0,1]}).
\end{equation}
利用这一观察，可对不必嵌入、但仍满足 $\gamma^{-1}(\operatorname{Critv}(\pi))=\{1\}$ 及 $\gamma'(1)\ne0$ 的路径 $\gamma$ 定义 Lefschetz 顶针和消失圈。具体地，$t_0$ 充分接近 $1$ 时，$\gamma|[t_0,1]$ 为嵌入路径，可用平行移动延伸 $\Delta_{\gamma|[t_0,1]}$。这样定义的 Lefschetz 顶针是浸入而不必嵌入的 Lagrangian 球体；但其边界恰是\eref{16.3}，具有与嵌入情形相同的性质。
\end{note}

\StatementNumber{5}\begin{example}\label{st:16.5}
取\stref{例题}{15.4}定义的简单模型 $\pi:E\to S$。它尚不完全是 Lefschetz 纤维化，但缺少的条件，即全纯凸性，与眼前的问题无关。唯一临界值为 $0$，对任何消失路径 $\gamma$，均可显式确定 Lefschetz 顶针：
\begin{equation}\label{eq:16.4}
\Delta_\gamma=\bigcup_{0\leq t\leq1}\sqrt{\gamma(t)}S^n;
\end{equation}
这里 $\sqrt zS^n=\{x\in\CC^{n+1}:x=\pm\sqrt z\,y\text{，其中某个 }y\in S^n\subset\R^{n+1}\}$。要看出\eref{16.4}成立，可用\stref{例题}{15.4}中的函数 $k$，它在平行移动下不变；再注意 $k^{-1}(0)$ 恰是对所有 $z\in S$ 取子集 $\sqrt zS^n$ 所成的并。
\end{example}

\subsection{Dehn 扭转}\label{sec:16c}
令 $M$ 为带角正合辛流形。对每个带标架正合 Lagrangian 球面 $V\subset M$，可赋予其 \emph{Dehn 扭转}，亦称 Picard–Lefschetz 变换，$\tau_V\in\operatorname{Aut}(M,\partial M)$；它在同痕意义下唯一。更准确地说，构造 $\tau_V$ 时作出的附加选择所成的空间可缩。定义时利用 Lagrangian 管状邻域定理，将标架 $v:S\to V$ 延伸为辛嵌入 $\psi:B_s^*S\to\operatorname{int}(M)$，其中 $B_s^*S$ 为 $S$ 上长度 $\leq s$ 的余切向量所成子集，$s>0$。现在取函数 $b:\R\to\R$，满足
\begin{equation}\label{eq:16.5}
\begin{cases}
b(r)=0,&r\geq s,\\
b(r)-b(-r)=-r,&\text{处处成立}.
\end{cases}
\end{equation}
\SourcePage{232}
在 $\psi(B_s^*S)\setminus V$ 上定义 $H(x)=b(|\psi^{-1}(x)|)$，并在 $M\setminus V$ 的其余部分延伸为零。则 $\tau_V$ 是该 Hamilton 函数的时间 $2\pi$ 流，并以唯一可能的方式延伸到 $V$ 本身。为获得更具体的图景，回顾 $T^*S\setminus S$ 上长度函数 $|\xi|$ 的 Hamilton 流可认同为归一化测地流：将任意非零切向量沿其出发的测地线以单位速度移动，而不考虑该向量的长度。因此，$b(|\xi|)$ 的流以速度 $b'(|\xi|)$ 移动切向量。当长度趋于零时，导数趋于 $-1/2$，故极限在零截面上给出对径映射。要证明零截面上的延伸光滑，而不只是连续，可写成 $b(r)=\beta(r^2)-r/2$，其中 $\beta$ 光滑；由于 $b(r)+r/2$ 为偶函数，这是可能的。$\beta(|\xi|^2)$ 的流在整个余切丛上光滑，而 $-|\xi|^2/2$ 的时间 $2\pi$ 流成为对径映射。此外，两 Hamilton 函数均为 $|\xi|$ 的函数，因此 Poisson 交换。由此可将 $b(|\xi|)$ 的时间 $2\pi$ 映射写成两个辛自同构之复合，每一个都显然在零截面附近光滑。\znote{224}{原书此处印作 $|\xi|^2/2$，但它与同一句的分解 $b(r)=\beta(r^2)-r/2$ 以及“导数趋于 $-1/2$”都不相容：$|\xi|^2/2$ 的流是测地流的反向重参数化，单位球面上测地线的周期为 $2\pi$，故时间 $2\pi$ 后回到原处，得到恒等而非对径映射。给出对径映射的应是 $|\xi|/2$：速度为 $-1/2$，时间 $2\pi$ 走弧长 $-\pi$。因此这一句应按 $|\xi|/2$ 理解，$|\xi|^2/2$ 疑为原书笔误。此外，真正在零截面上光滑的只是 $\beta(|\xi|^2)$ 这一支；$-|\xi|/2$ 在零截面上不可微，但由偶性拆分可知其时间 $2\pi$ 流光滑延拓，这正是本段要说明的。[来源：原书 §(16c)，印刷页 224；周期与导数均为初等核对。]}更仔细考察这一构造，可见 $\tau_V$ 实际上正合；这只有在最低维数，即 $M$ 为曲面时才非平凡，此时也可直接观察 Dehn 扭转的标准图形来证明。由定义，类 $[\tau_V]\in\pi_0(\operatorname{Aut}(M,\partial M))$ 在 $V$ 作为带标架正合 Lagrangian 球面的同痕下不变。特别地，它不依赖于 $V$ 的定向。

\StatementNumber{6}\begin{note}\label{st:16.6}
$[\tau_V]$ 是否依赖于标架，是一个颇为微妙的问题。对它诱导的 Fukaya 范畴自等价，已知标架在 $\Ai$ 函子同构意义下不影响结果。这里不证明这一断言，但一个仅涉及对象上作用的较弱版本，可由\stref{推论}{17.17}推出。
\end{note}

取正合 Lefschetz 纤维化 $\pi:E\to S$ 及消失路径 $\gamma$。令 $V_\gamma$ 为相应消失圈，位于 $M=E_{\gamma(0)}$。在 $S\setminus\operatorname{Critv}(\pi)$ 中取绕 $\gamma$ 去而复返的回路 $\lambda$，它逆时针绕过 $\gamma(1)$，见\sr{fig:16.1}{图~16.1}。则沿 $\lambda$ 的单值变换实质上就是沿 $V_\gamma$ 的 Dehn 扭转；更准确地，
\begin{equation}\label{eq:16.6}
h_\lambda\simeq\tau_{V_\gamma},
\end{equation}
其中 $\simeq$ 表示 $\operatorname{Aut}(M)$ 中的同痕。若按\stref{笔记}{15.2}假设纤维化在 $\partial^hE$ 附近平凡，则可稍加强这一断言：同痕可在群 $\operatorname{Aut}(M,\partial M)$ 内进行。这就是\emph{辛 Picard–Lefschetz 定理}。这里不给出证明，但也许值得概述基本策略。首先注意，按\secref{15b}所述作沿 $\partial^hE$ 行为良好的 Lefschetz 纤维化结构变形，不影响结论的有效性，因为它只在同痕意义下改变等式两边。如\secref{15d}所述，可以找到这种变形，使
\SourcePage{233}
$E$ 上辛结构在临界点附近的局部全纯 Morse 坐标中标准化。在这些局部坐标中，以\stref{例题}{16.5}为模型，可显式写出一条短路径趋向临界点时的消失圈。类似地，通过积分有关向量场，可确定绕临界点充分小圆周的单值变换，结果正是沿消失圈的 Dehn 扭转。随后，用一个简单的平行移动论证，就能对原来路径 $\gamma,\lambda$ 推出相应结论。\znote{225}{证明策略可压缩为四步：先用 §(15b) 的沿水平边界良好的变形把辛形式在临界点附近标准化，两边各自只改同痕类，故合法；在标准模型里借例题 16.5 写下显式的短消失路径及其消失圈；对绕临界值的小圆周直接积分向量场，验证单值变换就是沿该消失圈的 Dehn 扭转；最后用平行移动把局部结论搬回任意 $\gamma,\lambda$。要紧的是 $\lambda$ 须在 $\gamma$ 上往返（doubling），这样它恰好绕临界值一周；若绕 $k$ 周则给 $\tau_{V_\gamma}^k$，而不再是同痕意义下的 $\tau_{V_\gamma}$ 本身。[来源：原书 §(16c) 与 §(16d)，印刷页 224–226。]}
\SourceFigure{figure-16-1}{0.1495383}{16.1}

\StatementNumber{7}\begin{note}\label{st:16.7}
显然，用于证明\eref{16.6}的变形并不完全唯一，但可能选择所成的空间可缩。这说明，在两个辛自同构之间，实际上存在唯一的保持端点不动的同痕同伦类。
\end{note}

\clearpage
\subsection{特选基}\label{sec:16d}
现在集中考察底空间为圆盘 $S=D$ 的正合 Lefschetz 纤维化。选择基点 $*\in S\setminus\operatorname{Critv}(\pi)$，并将其视为位于边界上或靠近边界。意思是，若 $*\notin\partial S$，还要选择嵌入的\emph{基路径} $\chi:[0,1]\to S\setminus\operatorname{Critv}(\pi)$，从 $\chi(0)=*$ 通向某点 $\chi(1)\in\partial S$。

消失路径的\emph{特选基}是从 $*$ 出发的消失路径的有序族 $\symbfit\gamma=(\gamma_1,\ldots,\gamma_m)$，$m=|\operatorname{Critv}(\pi)|$，满足下述条件：
\begin{itemize}
\item 不同的 $\gamma_k$ 只在 $*$ 相交，出发方向 $\R^+\gamma_k'(0)\subset TS_*$ 两两不同。
\item 若 $*\in\partial S$，可将 $\R^+\gamma_k'(0)$ 看成上半平面 $H$ 中从原点出发的半直线，故有自然的顺时针顺序。要求它与 $\gamma_k$ 在 $\symbfit\gamma$ 中的顺序一致。另一种情形 $*\notin\partial S$ 下，
\SourcePage{234}
要求每个 $\gamma_k$ 与 $\chi$ 只在 $*$ 相交，且 $\R^+\gamma_k'(0)\ne\R^+\chi'(0)$。此时，$\R^+\gamma_k'(0)$ 是割缝平面 $\CC\setminus\R^+$ 中从原点出发的半直线，其中去掉的方向为 $\chi'(0)$ 的方向。这给出一个顺序，仍施加同样的条件。
\end{itemize}
所有可能的基的同痕类所成集合，带有辫群 $\operatorname{Br}_m$ 通过所谓 \emph{Hurwitz 移动}给出的作用：第 $k$ 个标准生成元 $\beta_k$，$1\leq k<m$，对应于第 $k$ 个初等 Hurwitz 移动 $\symbfit\gamma\mapsto\symbfit\delta=\mathscr L_k\symbfit\gamma$。该移动保持几乎所有路径不变，即 $j\ne k,k+1$ 时 $\delta_j=\gamma_j$，只按\sr{fig:16.2}{图~16.2}变换余下两条路径。容易验证辫关系成立。
\SourceFigure{figure-16-2}{0.44653798}{16.2}

还有一个相关但稍有不同的观点。令 $\mathcal D\subset\operatorname{Diff}^+(S)$ 为这样的保定向微分同胚组成的子群：在 $*$ 上为恒等，适用时在基路径 $\chi$ 上也为恒等，并将子集 $\operatorname{Critv}(\pi)$ 映到自身。显然，$\pi_0(\mathcal D)$ 作用于特选基的同痕类集合；利用 $\operatorname{Diff}(S,\partial S)$ 的可缩性，很容易看出此作用单传递。同类论证还表明，$\pi_0(\mathcal D)$ 同构于 $\operatorname{Br}_m$；然而，此同构并非典范，因此须稍加小心。先回顾：对任意避开 $*$ 及 $\chi$，并满足 $\mu^{-1}(\operatorname{Critv}(\pi))=\{-1,1\}$ 的嵌入路径 $\mu:[-1,1]\to S$，可赋予其\emph{半扭转} $t_\mu\in\mathcal D$，在 $\mu([-1,1])$ 的小邻域之外为恒等，其局部行为由\sr{fig:16.3}{图~16.3}刻画。显然，其中涉及一些选择，但所得类 $[t_\mu]\in\pi_0(\mathcal D)$ 唯一，并在相应路径类别内 $\mu$ 的同痕下不变。给定特选基 $\symbfit\gamma$，可将 $\gamma_k$ 的反向路径与 $\gamma_{k+1}$ 接合，并磨平拐角，定义路径 $\mu_k$，$1\leq k<m$。将标准生成元 $\beta_k$ 映到沿 $\mu_k$ 的半扭转，就给出同构 $\rho_{\symbfit\gamma}:\operatorname{Br}_m\to\pi_0(\mathcal D)$。关于这一同构，初等移动可简写为
\[
\mathscr L_k\symbfit\gamma=t_{\mu_k}^{-1}(\symbfit\gamma)
=\rho_{\symbfit\gamma}(\beta_k^{-1})(\symbfit\gamma).
\]
\SourcePage{235}
\SourceFigure{figure-16-3}{0.28038431}{16.3}
另一方面，这种移动也影响 $\rho_{\symbfit\gamma}$，通过共轭将其变为
\[
\rho_{\mathscr L_k\symbfit\gamma}(\beta)
=t_{\mu_k}^{-1}\circ\rho_{\symbfit\gamma}(\beta)\circ t_{\mu_k}
=\rho_{\symbfit\gamma}(\beta_k^{-1}\beta\beta_k).
\]
结合这两个观察可见，对所有 $k_1,\ldots,k_r\in\{1,\ldots,m-1\}$ 和 $\sigma_1,\ldots,\sigma_r\in\{\pm1\}$，有
\[
\mathscr L_{k_r}^{\sigma_r}\cdots\mathscr L_{k_1}^{\sigma_1}\symbfit\gamma
=\rho_{\symbfit\gamma}(\beta_{k_1}^{-\sigma_1}\cdots\beta_{k_r}^{-\sigma_r})(\symbfit\gamma).
\]
用文字说，以一系列初等移动定义的元素 $\beta\in\operatorname{Br}_m$ 对特选基 $\symbfit\gamma$ 的作用，等于微分同胚 $\rho_{\symbfit\gamma}(\beta^{-1})$ 将 $\symbfit\gamma$ 映到的像。这特别表明，任意两个基都由 Hurwitz 移动联系。\znote{227}{“特选”二字的含义在此：基点固定后，消失路径的排列顺序不是任意的，须与出发方向在 $*$ 处的顺时针顺序一致；所有基的同痕类因此构成一个辫群作用下的单一轨道，$\operatorname{Br}_m$ 经 Hurwitz 移动传递地作用其上。一般位置式的条件（$\gamma_k$ 只在 $*$ 相交、方向互异、避开 $\chi'(0)$）保证移动可逐个进行而不出现退化。注意 $\rho_{\symbfit\gamma}$ 本身依赖基，换基时要共轭，这正是上面共轭公式的来源，也是“辫单值”这一名称的由来：辫群元素记录的是临界值绕基点转动的相对方式。[来源：原书 §(16d)，印刷页 226–227。]}

现在引入消失圈的相应概念。在固定的带角正合辛流形内，考虑由 $m$ 个带标架正合 Lagrangian 球面组成的有序集族 $\symbfit V=(V_1,\ldots,V_m)$。对所涉及的 Lagrangian 球面施加带标架正合同痕所得的集族，不加区分。在此背景下，初等 Hurwitz 移动是某个 $1\leq k<m$ 所给出的变换 $\symbfit V\mapsto\symbfit W=\mathscr L_k\symbfit V$：对所有 $j\ne k,k+1$ 置 $W_j=V_j$，并置
\begin{equation}\label{eq:16.7}
W_k=\tau_{V_k}(V_{k+1}),\qquad W_{k+1}=V_k.
\end{equation}
如通常一样，一般 Hurwitz 移动是这些变换及其逆的任意复合。与先前话题的联系应该很清楚：任意消失路径特选基 $\symbfit\gamma$ 都确定一个有序集族 $\symbfit V=\symbfit V^{\symbfit\gamma}$，由消失圈 $V_k=V_{\gamma_k}$ 组成；称为相应的\emph{消失圈特选基}。\stref{笔记}{16.4}与\eref{16.6}表明，消失路径的初等 Hurwitz 移动 $\mathscr L_k$ 在同痕意义下诱导特选基的相应移动：
\[
\symbfit V^{\mathscr L_k\symbfit\gamma}\simeq\mathscr L_k\symbfit V^{\symbfit\gamma}.
\]
\SourcePage{236}
将集族 $\symbfit V$ 的 \emph{Hurwitz 等价}定义为由 Hurwitz 移动及同痕生成的关系。用这些术语，所得结论就是：$\symbfit V^{\symbfit\gamma}$ 的 Hurwitz 等价类是 Lefschetz 纤维化的不变量；严格地说，是纤维化连同基点选择，以及必要时的基路径选择，的不变量。在适当意义下，它在\secref{15b}定义的沿水平边界行为良好的变形下也不变。

\subsection{构造 Lefschetz 纤维化}\label{sec:16e}
定义集族 $\symbfit V$ 上的 Hurwitz 移动时，并未预先要求它们来自消失圈。结果表明，这并无区别：任意这样的集族都可实现为某个纤维化的特选基。现在将此说得更准确些。

\clearpage
\StatementNumber{8}\begin{lemma}\label{st:16.8}
令 $V$ 为带角正合辛流形 $M$ 中带标架的正合 Lagrangian 球面。取圆盘 $S=D$，及从某点 $\gamma(0)\in\partial S$ 到 $\gamma(1)=0$ 的嵌入路径。则存在正合 Lefschetz 纤维化 $\pi:E\to S$，满足 $\operatorname{Critv}(\pi)=\{0\}$，并带有认同 $E_{\gamma(0)}=M$，使 $\gamma$ 的消失圈在此认同下对应于 $V$，且与标架相容。此外，该纤维化在 $\partial^hE$ 附近平凡。
\end{lemma}

构造很简单。从\stref{例题}{15.4}的模型出发，其纤维为球面余切丛的紧部分 $B_s^*S^n$，其中 $s$ 充分小。如 Dehn 扭转的定义，再用 Lagrangian 管状邻域定理找到延伸标架 $v:S^n\to V$ 的嵌入 $\psi:B_s^*S^n\to M$。借此把局部模型与平凡纤维化 $D\times(M\setminus V)$ 粘合。纤维上的辛形式一致，但在总空间上，还须用适当的截断函数把它们拼接。所需的只是，在相应 Lefschetz 顶针\eref{16.4}的邻域中保留局部模型的二形式；随后，所有消失路径都会给出所需的 Lagrangian 球面。

一旦能够构造具有单个指定消失圈的 Lefschetz 纤维化，就很容易将它们与非奇异部分粘合，得到一般结果：

\StatementNumber{9}\begin{lemma}\label{st:16.9}
令 $\symbfit V=(V_1,\ldots,V_m)$ 为带角正合辛流形 $M$ 中带标架正合 Lagrangian 球面的集族。在圆盘 $S=D$ 上选择基点 $*$，必要时还选择基路径 $\chi$，并选择消失路径特选基 $(\gamma_1,\ldots,\gamma_m)$。则存在正合 Lefschetz 纤维化 $\pi:E\to S=D$，其临界值恰是端点 $\gamma_1(1),\ldots,\gamma_m(1)$；它带有认同 $E_*=M$，使带标架消失圈 $V_{\gamma_k}$ 在此认同下对应于 $V_k$。此外，该纤维化在 $\partial^hE$ 附近平凡。
\end{lemma}

\stref{引理}{16.8}有一个唯一性方面的对应结果。具体地，假设在 $D$ 上有一个 Lefschetz 纤维化，只有一个临界点，位于 $0$ 上方的纤维中，
\SourcePage{237}
且在 $\partial^hE$ 附近平凡。最后这个要求是技术性的，重要性很有限；它源于 $\partial E$ 应弱 $I_E$ 凸这一条件，而该条件在辛联络变形下表现得不太方便。则在沿 $\partial^hE$ 行为良好的变形意义下，该纤维化完全由 $1$ 上方的纤维及其中的消失圈确定；消失圈视为同痕意义下的带标架 Lagrangian 球面。在\stref{引理}{16.9}更复杂的情形中也有相应结论。

\subsection{分次问题}\label{sec:16f}
令 $M^{2n}$ 为带角正合辛流形，带二次复体积形式 $\eta_M^2$。若 $V\subset M$ 为可赋予分次的带标架正合 Lagrangian 球面，并不需要具体选择某个分次，则 Dehn 扭转 $\tau_V$ 自然地成为\secref{12i}意义下的分次辛自同构；实值分次函数 $\gamma_{\tau_V}^\#$ 由下述条件确定：在 $V$ 的某个邻域之外恒为零，而 $\tau_V$ 在该处为恒等。证明并不困难：上同调检查表明，维数 $n>1$ 时不存在障碍，余下的情形则是初等的。\znote{229}{“分次函数” $\gamma^\#_{\tau_V}$ 应与\secref{12i}的定义对照读：$\tau_V$ 在 $V$ 之外为恒等，故 $\gamma^\#_{\tau_V}$ 支撑在 $V$ 附近，它不是 $\tau_V$ 的“像”，而是记录 $\tau_V$ 把分次 Lagrangian 的相位移动了多少的实值函数。$n>1$ 时 $V$ 的余维至少为 $2$，绕过 $V$ 的回路不会被迫穿过 $V$，相位提升的障碍自动消失；$n=1$ 时 $V$ 是余维 $1$ 的圈，$M\setminus V$ 被真正切开，两侧的相位提升可以相差整数，故须单独核对。$\eta_{E/S}^2=\eta_E^2/\eta_S^2$ 只是“相对”二字的含义：底空间的贡献由 $\eta_S^2$ 除掉，留在纤维上的是相对量。[来源：原书 §(16f)，印刷页 229。]}

另一方面，考虑配备相对二次复体积形式 $\eta_{E/S}^2$ 的 Lefschetz 纤维化。它在光滑纤维上诱导体积形式，关于这些形式，任意消失圈均可赋予分次，但没有自然选定的分次。可如前论证：除最低维情形外，其余情况都平凡，而最低维情形可直接解决。不过，还有一个更优美的论证，与证明消失圈正合性的论证类似：令 $\gamma$ 为消失路径。在 $\pi^{-1}(\gamma)\subset E$ 的某个邻域中，写成 $\eta_{E/S}^2=\eta_E^2/\eta_S^2$；事实上，由假设 $TS$ 平凡，可在整个空间中这样写，但此处不需要。由于 $\Delta_\gamma$ 可缩，与 $\eta_E^2$ 关联的平方相位函数
\[
\alpha_{\Delta_\gamma}:\Delta_\gamma\longrightarrow S^1
\]
当然可以提升到 $\R$。该函数在边界上的限制与 $V_\gamma$ 的平方相位函数只差一个依赖于 $\gamma'(0)$ 的常数，故得到所需事实。最后，考察上面概述的\eref{16.6}证明可见：若为单值变换赋予\eref{15.8}定义的分次，并按刚才描述的方式为 Dehn 扭转赋予分次，则那里构造的同痕提升到分次辛自同构群 $\operatorname{Aut}^{\mathrm{gr}}(M)$。

最后，若给定 $M$、$\eta_M^2$，以及一族均可赋予分次的带标架正合 Lagrangian 球面，则按\stref{引理}{16.9}构造的相应 Lefschetz 纤维化可赋予相对二次复体积形式，且相对于 $E_*=M$ 在同伦意义下唯一。同样，无须具体选择分次，只要它存在即可。

\StatementNumber{10}\begin{note}\label{st:16.10}
高维 Dehn 扭转为辛的事实，由 Arnol'd\bcite{3}发现。相应地，消失圈为 Lagrangian 这一事实
\SourcePage{238}
由 Donaldson 向作者指出，尽管此前也很可能已有人知道。至于 Picard–Lefschetz 定理，它在拓扑层次，即将单值变换视为微分同胚的层次，是经典结果。众多可用参考文献中的两项是\bcite{92}、\bcite{5}。辛版本见\bcite{131}命题 1.15。消失圈的特选基也是经典概念，在奇点理论考察的局部情形中尤其如此\bcite{5}、\bcite{4}。在辛层次上可从这种基重建总空间的观察，则较为近期，见本章开头引用的论文。Dehn 扭转的分次在\bcite{125}中有更详细讨论。
\end{note}

\subsection{匹配圈}\label{sec:16g}
令 $\pi:E\to S$ 为 Lefschetz 纤维化，仍照例假设正合。取嵌入路径 $\mu:[-1,1]\to\operatorname{int}(S)$，满足 $\mu^{-1}(\operatorname{Critv}(\pi))=\{-1,1\}$。可将它分为具有相同起点的一对消失路径 $\gamma_\pm(t)=\mu(\pm t)$，$t\in[0,1]$，因而得到一对消失圈 $V_{\gamma_\pm}\subset M=E_{\mu(0)}$。先看一个特别简单的情形，即两者相等。此时，
\begin{equation}\label{eq:16.8}
\Sigma_\mu=\Delta_{\gamma_-}\cup\Delta_{\gamma_+}
\end{equation}
是总空间 $E$ 的光滑 Lagrangian 子流形。由 Lefschetz 顶针的定义，对所有 $-1<t<1$，沿 $\mu$ 的平行移动把交 $\Sigma_\mu\cap E_{\mu(t)}$ 互相映到彼此，这在重叠处 $\Delta_{\gamma_-}\cap\Delta_{\gamma_+}=V_{\gamma_\pm}$ 附近给出局部坐标 $(-1,1)\times V_{\gamma_\pm}$。由于它是两个球沿边界粘合的结果，$\Sigma_\mu$ 必为同伦球面。若假设 $V_{\gamma_\pm}$ 的标架同痕，则它在微分意义下是标准球面。事实上，给定两个标架之间的同痕，可把\eref{16.1}的微分同胚合在一起，得到 $\Sigma_\mu$ 的标架。将\eref{16.8}称为\emph{朴素匹配圈构造}。

要求两个消失圈严格相等似乎很不现实，但可以稍微放宽假设，代价是失去对所得球面在 $E$ 中位置的控制。具体地，假设 $V_{\gamma_\pm}$ 在 $M$ 中作为带标架正合 Lagrangian 球面同痕。此时称 $\mu$ 为\emph{匹配路径}。两个消失圈之间的任意带标架正合 Lagrangian 同痕 $I=(i,j_0,j_1)$，也相应称为\emph{匹配同痕}。\znote{230}{匹配圈与消失圈的区别在于“控制权”：消失圈由纤维化完全决定，而匹配圈要先把两个消失圈在中心纤维里同痕到一起，再沿 $\mu$ 把两块顶针接起来，因此只能确定到 $E$ 中的同痕类，位置不受控制——这正是本节开头所说“失去对所得球面在 $E$ 中位置的控制”。$I=(i,j_0,j_1)$ 中的 $i^s$ 是把 $V_{\gamma_-}$ 同痕到 $V_{\gamma_+}$ 的一族 Lagrangian 嵌入，$j_0,j_1$ 则把两端的标架对齐；所以 $I$ 既是位置的粘合数据，也是标架的粘合数据，舍去它 $\Sigma_\mu$ 就不再是带标架球面。[来源：原书 §(16g)，印刷页 230。]}初等辛几何表明，$\operatorname{Aut}(M,\partial M)$ 中存在同痕 $(\phi^s)$，使 $\phi^0=\Id$、$\phi^s\circ i^0=i^s$。由\stref{引理}{15.3}，存在给定正合 Lefschetz 纤维化的变形，沿纤维恒定并集中在 $E_{\mu(-1/4)}$ 的某个邻域，使沿 $\mu|[-1/2,0]$ 的平行移动映射相应变形为
\[
h_{\mu|[-1/2,0]}^s=\phi^s\circ h_{\mu|[-1/2,0]};
\]
右边的平行移动是初始结构，即 $s=0$ 时的平行移动。对消失圈的影响为 $V_{\gamma_-}^s=\phi^s(V_{\gamma_-})$，而 $V_{\gamma_+}$ 保持不变。在变形的另一端 $s=1$，得到一个两个消失圈相等的正合 Lefschetz 纤维化结构，因而可用朴素匹配圈构造在总空间中定义
\SourcePage{239}
Lagrangian 球面。由于在 $i^0,i^1$ 对有关向量空间给出的认同下，消失圈的两个标架也相同，朴素匹配圈便带有标架。对反向变形应用\stref{引理}{7.1}，得到从 $(E,\omega_E^1)$ 到 $(E,\omega_E=\omega_E^0)$ 的共形辛嵌入 $\rho$，并用它把朴素匹配圈映回原来的辛形式。在带标架正合 Lagrangian 球面的同痕意义下，结果不依赖于 $(\phi^s)$、Lefschetz 纤维化变形以及该嵌入 $\rho$ 的选择。称其为与 $\mu,I$ 关联的\emph{匹配圈}，记为 $\Sigma_{\mu,I}\subset E$。为简化记号，在有明显的同痕选择，或不关心采用哪一种选择时，仍写作 $\Sigma_\mu$。

\StatementNumber{11}\begin{note}\label{st:16.11}
同痕 $I$ 在构造中的作用不易估量。保持端点不动地变形 $I$，不影响匹配圈的同痕类。在最低维情形 $n=1$，曲面 $M$ 中正合嵌入圆周所成的空间具有平凡的 $\pi_1$，故 $\Sigma_{\mu,I}$ 与 $I$ 无关。更高维中的相应问题似乎很难回答。不过，后面结合\stref{命题}{18.21}与\stref{命题}{18.13}将看到，$\Sigma_{\mu,I}$ 作为 Donaldson–Fukaya 范畴对象的同构类与 $I$ 无关。
\end{note}

一般而言，无法保证 $E$ 中给定的 Lagrangian 球面能表示为匹配圈；不过，Lefschetz 束背景下此类渐近结论可参看\bcite{16}。此外，不同同痕类的匹配路径可以给出同痕的匹配圈。这不是偶发现象，而是任何足够复杂的 Lefschetz 纤维化中都可以预期出现的情况。下例出自 Auroux：

\StatementNumber{12}\begin{example}\label{st:16.12}
为简单起见，考虑最低维情形 $\dim(E)=4$。假设底空间 $S$ 中有两条嵌入路径 $\mu,\nu$，各自连接一对临界值，四个临界值两两不同；两条路径恰相交一次，显然是在各自的某个内点，并且横截相交。假设 $\mu$ 为匹配路径，而 $\nu$ 两半对应的消失圈横截相交且只有一个交点。令 $t_\nu$ 为沿 $\nu$ 的半扭转。由于沿相交一次的圆周的 Dehn 扭转满足辫关系，$\mu'=t_\nu^3(\mu)$ 仍为匹配路径，且相应的 Lagrangian 二维球面 $\Sigma_{\mu'}$ 与 $\Sigma_\mu$ 同痕。一种直接但稍不直观的做法是：取所涉及全部 Lefschetz 顶针的邻域；从胞腔附着的观点考察它；认出它与 $T^*S^2$ 的某个紧部分辛同构；再应用\bcite{74}，证明 Lagrangian 二维球面在同痕意义下的唯一性。这一构造还有一个略微不那么有趣的变体：假设 $\nu$ 两半对应的消失圈不交，并取 $\mu'=t_\nu^2(\mu)$。
\end{example}

\SourcePage{240}
\subsection{路径的接合}\label{sec:16h}
令 $\mu_0,\mu_1$ 为两条匹配路径，只在一个共同端点相交。假设路径的定向满足 $\mu_0(1)=\mu_1(-1)$；这只是为方便起见，定向在下面的构造中不起作用。考虑将 $\mu_k$ 拼起来，再移到共同端点右侧所得到的路径 $\mu=\mu_1\#\mu_0$，见\sr{fig:16.4}{图~16.4}。这是一种不对称的运算，依赖于如何排列这两条路径。显然，$\mu$ 仍为匹配路径。我们要解释的结果是：相应的匹配圈 $\Sigma_\mu$ 由 $\Sigma_{\mu_0}$ 沿 $\Sigma_{\mu_1}$ 作一次 Dehn 扭转得到。
\SourceFigure{figure-16-4}{0.25753818}{16.4}
为准确陈述这一结论，须把标架和同痕考虑在内。每个 $\mu_k$ 分为一对消失路径 $\gamma_k^\pm$，因而有四个相应的消失圈。方便起见，假设 $\mu|[-1,-1/2]$ 在显然的重新参数化下与 $\mu_0|[-1,0]$ 一致，类似地，$\mu|[1/2,1]$ 与 $\mu_1|[0,1]$ 一致。接合路径 $(\mu|[-1/2,0])^{-1}\cdot\gamma_0^+$ 与 $(\mu|[0,1/2])\cdot\gamma_1^-$ 保持端点不动地同伦；由此在 $E_{\mu(0)}$ 中得到带标架同痕
\begin{equation}\label{eq:16.9}
h_{\mu|[-1/2,0]}(V_{\gamma_0^+})\simeq
h_{\mu|[0,1/2]}^{-1}(V_{\gamma_1^-}).
\end{equation}
选择定义匹配圈 $\Sigma_{\mu_k,I_k}$ 所需的同痕 $I_k$。结合\eref{16.9}，便得到 $\mu$ 两半 $\gamma^\pm$ 的消失圈之间的同痕 $I$：
\[
V_{\gamma^-}=h_{\mu|[-1/2,0]}(V_{\gamma_0^-})
\overset I\simeq h_{\mu|[0,1/2]}^{-1}(V_{\gamma_1^+})=V_{\gamma^+},
\]
并用它定义另一个匹配圈 $\Sigma_{\mu,I}$。

\StatementNumber{13}\begin{lemma}\label{st:16.13}
在 $E$ 中存在带标架 Lagrangian 球面的同痕
\begin{equation}\label{eq:16.10}
\Sigma_{\mu,I}\simeq\tau_{\Sigma_{\mu_1,I_1}}(\Sigma_{\mu_0,I_0}).
\end{equation}
\end{lemma}
\begin{proof}
这是通过变形到更简单特殊情形来论证的典型例子。若结论对某个正合 Lefschetz 纤维化成立，则对其任意沿 $\partial^hE$ 行为良好的变形也成立；因为一旦为变形后的 Lefschetz 纤维化结构找到同痕，就可以
\SourcePage{241}
结合\stref{引理}{7.1}那样的一族共形辛嵌入，得到原结构的同痕。由\stref{引理}{15.3}及证明 Picard–Lefschetz 公式时的技巧，允许作下述假设：$\mu_0$ 两半对应的消失圈相等，$\mu_1$ 也如此；同痕 $I_0,I_1$ 恒定；且在 $E_{\mu_0(1)}$ 中唯一临界点附近，存在全纯坐标，使 $\pi$ 成为标准二次型\eref{15.9}，$\omega_E$ 成为 $\CC^{n+1}$ 上标准 Kähler 形式。仍不影响结论，还可变形所涉及的路径，使在 $\mu_0(1)$ 附近的局部坐标中，$\mu_0,\mu_1$ 分别为负、正实半轴，且 $\mu$ 在该点的小邻域之外与 $\mu_0\cup\mu_1$ 一致。于是，在 $E_{\mu_0(1)}$ 中临界点的小邻域之外，$\Sigma_{\mu,I}=\Sigma_{\mu_0,I_0}\cup\Sigma_{\mu_1,I_1}$；而在该点附近，由\stref{例题}{16.5}得到如下局部图景：
\begin{equation}\label{eq:16.11}
\Sigma_{\mu_0,I_0}=i\R^{n+1},\qquad
\Sigma_{\mu_1,I_1}=\R^{n+1},\qquad
\Sigma_{\mu,I}=\bigcup_t\sqrt{\mu(t)}S^n,
\end{equation}
其中 $\sqrt\mu$ 是 $\CC$ 中位于闭的右下象限内的光滑嵌入路径，满足 $t<-\epsilon$ 时 $\sqrt{\mu(t)}=it$，$t>\epsilon$ 时 $\sqrt{\mu(t)}=t$，$\epsilon$ 为某个小的正数。\znote{233}{局部模型\eref{16.11}是理解整个构造的关键，值得记牢：$\Sigma_{\mu_0}$ 在标准坐标里是虚方向的 $\R^{n+1}$，$\Sigma_{\mu_1}$ 是实方向的 $\R^{n+1}$，接合后的 $\Sigma_\mu=\bigcup_t\sqrt{\mu(t)}S^n$ 在 $t<-\epsilon$ 时落在 $i\R^{n+1}$ 中、$t>\epsilon$ 时落在 $\R^{n+1}$ 中，中间由半径随 $t$ 变化的球面插值。于是 $\Sigma_\mu$ 相当于把 $\Sigma_{\mu_0}$ 与 $\Sigma_{\mu_1}$ 各挖去一个开球，再沿两个半径 $\epsilon$ 的 $n$ 维球面用柄粘起来——这正是连通和，也是它与 $\Sigma_{\mu_1}\#\Sigma_{\mu_0}$ 同痕的直观原因。[来源：原书 §(16h)，印刷页 233。]}\eref{16.11}中第三个子流形，就是 Lagrangian 连通和或 Lagrangian 手术构造中出现的 \emph{Lagrangian 柄}。同样用 $\#$ 表示该连通和，并采用\bcite{125}小节 2e 的约定，有
\begin{equation}\label{eq:16.12}
\Sigma_{\mu,I}=\Sigma_{\mu_1,I_1}\#\Sigma_{\mu_0,I_0}.
\end{equation}
而\eref{16.10}与\eref{16.12}的右边在 Lagrangian 同痕意义下相同\bcite{124}，附录。同痕 $I_0,I_1,I$ 均恒定，所以容易验证这与标架相容。
\end{proof}

\StatementNumber{14}\begin{note}\label{st:16.14}
既然使用了记号 $\mu=\mu_1\#\mu_0$，忽略同伦后，\eref{16.10}可更直观地写为 $\Sigma_{\mu_1\#\mu_0}=\Sigma_{\mu_1}\#\Sigma_{\mu_0}$。还有另一种重要解释。令 $t_{\mu_1}$ 为沿 $\mu_1$ 的半扭转。由于 $\mu$ 与 $t_{\mu_1}(\mu_0)$ 同痕，我们证明的是
\[
\Sigma_{t_{\mu_1}(\mu_0)}\simeq\tau_{\Sigma_{\mu_1}}(\Sigma_{\mu_0}).
\]
在这种形式下，结论可推广到任意相交的匹配路径 $\mu_0,\mu_1$。但证明与上述特殊情形完全不同，此处也确无使用它的需要，故不作介绍；不过，\stref{引理}{18.3}可视为另一个特殊情形。
\end{note}

\subsection{将消失圈视为匹配圈}\label{sec:16i}
取正合 Lefschetz 双纤维化，记号如\eref{15.11}。第一个观察，是关于
\SourcePage{242}
临界值曲线 $\operatorname{Critv}(\varpi)\subset\mathcal S$ 及其到 $\mathcal W$ 的投影的拓扑的一个简单推论。对每个 $w\in\mathcal W\setminus(\operatorname{Critv}(\Psi)\cup\operatorname{Fakev}(\Psi,\varpi))$，考虑曲面 $\mathcal S_w$ 及其中的有限点构形 $\operatorname{Critv}(\varpi)\cap\mathcal S_w=\operatorname{Critv}(\varpi_w)$。由于交横截，该构形随 $w$ 光滑变化。沿任意路径 $\beta:[0,1]\to\mathcal W\setminus(\operatorname{Critv}(\Psi)\cup\operatorname{Fakev}(\Psi,\varpi))$ 追踪它的移动，就可定义\emph{辫平行移动映射}，即保定向微分同胚
\[
k_\beta:(\mathcal S_{\beta(0)},\operatorname{Critv}(\varpi_{\beta(0)}))
\longrightarrow(\mathcal S_{\beta(1)},\operatorname{Critv}(\varpi_{\beta(1)})),
\]
在同痕意义下唯一。可使它为辛，但由于维数低，这样做不会得到更多信息。现假设 $\gamma:[0,1]\to\mathcal W$ 为 $\Psi$ 的消失路径，且避开 $\operatorname{Fakev}(\Psi,\varpi)$。\stref{引理}{15.8}表明，当 $t\to1$ 时，$\operatorname{Critv}(\varpi)\cap\mathcal S_{\gamma(t)}$ 中的两个点，在覆盖 $\psi|\operatorname{Critv}(\varpi):\operatorname{Critv}(\varpi)\to\mathcal W$ 的一个分歧点处汇合。当 $s$ 接近 $1$ 时，取 $\mathcal S_{\gamma(s)}$ 中连接这两个点的短路径，它在不考虑定向的同痕意义下唯一。其在 $k_{\gamma|[0,s]}$ 下的原像，是嵌入路径 $\mu:[-1,1]\to\mathcal S_{\gamma(0)}$，满足 $\mu^{-1}(\operatorname{Critv}(\varpi_{\gamma(0)}))=\{-1,1\}$。称之为\emph{与 $\gamma$ 关联的匹配路径}，理由如下：

\StatementNumber{15}\begin{lemma}\label{st:16.15}
$\mu$ 是 $\varpi_{\gamma(0)}$ 的匹配路径，且适当选择同痕 $I$ 后，相应匹配圈 $\Sigma_{\mu,I}$ 在 $\mathcal E_{\gamma(0)}$ 中作为带标架 Lagrangian 球面，与消失圈 $V_\gamma$ 同痕。
\end{lemma}
\begin{proof}
只须证明：$s$ 接近 $1$ 时，在 $\mathcal E_{\gamma(s)}$ 中，连接 $\operatorname{Critv}(\varpi_{\gamma(s)})$ 中两个相近点的短路径是匹配路径，且所得匹配圈为 $\gamma|[s,1]$ 的消失圈。随后将一切从 $\mathcal E_{\gamma(s)}$ 移回 $\mathcal E_{\gamma(0)}$，就得到原结论。

考虑 $\mathcal E_{\gamma(1)}$ 中 $\Psi$ 的唯一临界点。在该点及其像附近引入具有\stref{引理}{15.9}所述性质的局部全纯坐标。如\eref{16.6}或\stref{引理}{16.13}的证明，可局部修改 $\mathcal E$ 上辛形式，使其在这些坐标中标准化。也可变形路径 $\gamma$，使在 $\Psi(x)\in\mathcal W$ 附近的局部坐标中，$t$ 接近 $1$ 时有 $\gamma(t)=1-t$。考虑 $s$ 接近 $1$ 时的短消失路径 $\gamma|[s,1]$。如\stref{例题}{16.5}，相应消失圈为
\[
V_{\gamma|[s,1]}=\sqrt{1-s}\,S^n\subset\CC^{n+1},
\]
另一方面，同一个球面也是 $\mathcal S_{\gamma(s)}$ 中连接 $\varpi_{\gamma(s)}$ 的两个临界值 $(y_1=1-s,y_2=\pm\sqrt{1-s})$ 的直线段所对应的匹配圈。验证这一点毫无困难，因为在此可用朴素程序\eref{16.8}定义匹配圈。最后，如\stref{引理}{16.13}的证明，应用\stref{引理}{7.1}回到原来的辛形式。当然，对该形式，匹配圈与消失圈不再严格相等，但论证给出两者之间的带标架同痕。\znote{234}{这条引理回答了第 16 节标题中“与”字的含义：消失圈与匹配圈在 Lefschetz 纤维化内部本是同一类对象，只是造法不同——消失圈沿 $\Psi$ 的一条消失路径生成，匹配圈则在曲面纤维 $\mathcal S_{\gamma(s)}$ 中取两个临界值之间的短路径、再经辫平行移动 $k_{\gamma|[0,s]}$ 拉回。注意上面的公式 $V_{\gamma|[s,1]}=\sqrt{1-s}\,S^n\subset\CC^{n+1}$ 中半径随 $s\to1$ 趋于零，正是“收缩到临界点”的显式写法；而同一条短直线段（$y_1=1-s$ 到 $y_2=\pm\sqrt{1-s}$）给出的匹配圈恰是同一个球面。[来源：原书 §(16i) 与引理 16.15，印刷页 234。]}
\end{proof}

\SourcePage{243}
\StatementNumber{16}\begin{note}\label{st:16.16}
$V_\gamma$ 的同痕类只依赖于 $\gamma$ 在全体消失路径空间中的同痕类；特别地，即使把 $\gamma$ 移过 $\operatorname{Fakev}(\Psi,\varpi)$ 的一个点，也不改变它。但辫平行移动通常会改变，因而关联的匹配路径 $\mu$ 也会改变。这并不矛盾，因为不同的匹配路径可以给出同一个匹配圈。事实上，\stref{例题}{16.12}中提到的这一现象，正是以这种方式从最简单的一般伪临界值产生的；在这些值处，曲线 $\operatorname{Critv}(\varpi)$ 具有尖点或结点。
\end{note}

\StatementNumber{17}\begin{note}\label{st:16.17}
如已提到，匹配圈由 Donaldson 引入，尚未发表。Lagrangian 连通和，尤其二维情形，研究于\bcite{109}。近来，它在镜像对称中也有了一定重要性，见\bcite{140}、\bcite{139}；后者包含公式\eref{16.12}的一个版本。熟悉辫单值变换技术，例如 Auroux 的分歧覆盖\bcite{11}与 Donaldson 束之间关系的读者，会觉得\secref{16i}的内容显而易见。
\end{note}

\section{伪全纯截面}\label{sec:17}
\subsection{截面计数}\label{sec:17a}

推广 Floer TQFT 标准框架的一种方式，是将目标空间换成 Lefschetz 纤维化，并以伪全纯截面代替映射。先在最简单的情形，即底曲面紧致时解释这一推广，再扩展框架以允许条带状端。这样做的目的，是将代数扭转构造 $T$ 与几何 Dehn 扭转 $\tau$ 联系起来，见\stref{推论}{17.17}。我们会特别仔细地以相当典范的形式建立这种认同，使之易于推广到 $\Z/2$ 等变背景。

令 $\pi:E^{2n+2}\to S$ 为 Lefschetz 纤维化，仍照例正合。\emph{Lagrangian 边界条件}是一个包含在 $\partial^vE$ 中且与 $\partial^hE$ 不交的子流形 $F^{n+1}\subset E$，并要求存在 $\alpha_F\in\Omega^1(\partial S)$ 及 $h_F\in C^\infty(F,\R)$，使
\begin{equation}\label{eq:17.1}
\theta_E|F=(\pi|F)^*(\alpha_F+\theta_S|\partial S)+dh_F.
\end{equation}
稍加解释。$F$ 显然在 $E$ 中为 Lagrangian。由于 $\pi$ 在 $\partial^vE$ 中没有临界点，$\pi|F:F\to\partial S$ 必为纤维丛，见\stref{引理}{16.2}。每个 $F_z\subset E_z$，$z\in\partial S$，实际上都是正合 Lagrangian 子流形，且这些子流形由沿边界的平行移动互相映到彼此，见\stref{引理}{16.3}。逆命题也在很大程度上成立：假设 $F_z\subset E_z$，$z\in\partial S$，是一族由边界平行移动互相映到彼此的正合 Lagrangian 子流形，令 $F$ 为它们的并。则 $\theta_E|F$ 闭，且可找到函数 $h_F$，使 $\theta_E|F-dh_F$ 限制到每个 $F_z$ 时消失。若
\SourcePage{244}
再假设 $F_z$ 连通，则 $\theta_E|F-dh_F$ 必是 $\partial S$ 上某个一形式的拉回，故 $F$ 为 Lagrangian 边界条件。\znote{236}{“正合 Lagrangian 边界条件”的准确含义是这个分解，而不只是“$F$ 为 Lagrangian”。它由三件事合作给出：$\pi$ 在 $\partial^vE$ 上无临界点，故 $\pi|F$ 是纤维丛（引理 16.2）；每条纤维 $F_z$ 正合、且由平行移动彼此相联（引理 16.3）；再加上 $F_z$ 连通，使拼出的 $\theta_E|F-dh_F$ 只能来自 $\partial S$ 上的一形式。少了连通性，各纤维上的 $h_F$ 之间可差局部常数，$F$ 就未必满足\eref{17.1}。[来源：原书 §(17a)，印刷页 236。]}

固定这样的 $F$。\emph{相对扰动数据}是一对 $(K,J)$，其中：$K$ 为 $E$ 上一形式，在 $TE^v$ 中的向量上为零，在 $\partial^hE\cup\operatorname{Crit}(\pi)$ 的某个邻域中完全消失，且满足 $K|F=0\in\Omega^1(F)$；$J$ 为 $E$ 上的 $\omega_E$ 相容殆复结构，使 $\pi$ 伪全纯，并在 $\partial^hE\cup\operatorname{Crit}(\pi)$ 的某个邻域中与 $I_E$ 一致。给定这些数据，边界位于 $F$ 中的\emph{非齐次伪全纯截面}是映射 $u:S\to E$，满足
\begin{equation}\label{eq:17.2}
\begin{cases}
\pi(u(z))=z,\\
Du(z)+J(u)\circ Du(z)\circ I_S=Y(u)+J(u)\circ Y(u)\circ I_S,\\
u(\partial S)\subset F.
\end{cases}
\end{equation}
这里 $Y$ 是丛 $\Hom(\pi^*TS,TE^v)$ 的截面，按下述方式由 $K$ 确定：给定任意 $\xi\in TS_z$，$z=\pi(x)$，在 $E_z$ 的正则部分上将其提升为切向量场 $X$；函数 $K(X)$ 与提升的选择无关，定义 $Y(\xi)_x$ 为该函数在 $E_z$ 上关联的 Hamilton 向量场于 $x$ 处的值。由于要求 $K$ 在临界点附近为零，$Y$ 也如此，故可平凡地延伸过这些点。

$Y$ 在临界点附近消失，因而可连续延伸过这些点；在 $\partial^hE$ 附近也消失。

\StatementNumber{1}\begin{note}\label{st:17.1}
相对扰动数据有一个简单的几何解释。考虑修改后的二形式
\begin{equation}\label{eq:17.3}
\omega_E^{\mathrm{new}}=\omega_E-dK,
\end{equation}
它仍闭且逐纤维为辛；其关联联络为
\begin{equation}\label{eq:17.4}
TE_x^{h,\mathrm{new}}=\{X=X^h+X^v\in TE_x=TE_x^h\oplus TE_x^v:
X^v=Y(D\pi(X))\}.
\end{equation}
关于殆复结构，注意按假设 $J$ 保持分解 $TE^h\oplus TE^v$，且在第一个直和项上等于 $I_E$。故这里的选择实质上是 $TE^v$ 上殆复结构的选择。存在唯一的 $J^{\mathrm{new}}$，使 $\omega_E^{\mathrm{new}}(\cdot,J^{\mathrm{new}}\cdot)$ 对称、$\pi$ 为 $J^{\mathrm{new}}$ 全纯，且在 $TE^v$ 上有 $J^{\mathrm{new}}=J$。用这些术语，\eref{17.2}就是截面在通常意义下为 $J^{\mathrm{new}}$ 全纯的方程。联络观点也有助于理解解 $u$ 的一些基本性质。例如，能量的先验估计为
\[
E(u)=\int_S\tfrac12|(Du)^v-Y(u)|^2
=\int_{\partial S}(\alpha_F+\theta_S)-\int_Su^*R,
\]
\SourcePage{245}
其中 $(\cdot)^v$ 是关于 $TE=TE^h\oplus TE^v$ 的垂直分量，而曲率项在局部平凡化\eref{15.3}中由下式给出：
\begin{equation}\label{eq:17.5}
\begin{aligned}
R(\partial_s,\partial_t)=\bigl(&\partial_s(K-\kappa)(\partial_t)-\partial_t(K-\kappa)(\partial_s)\\
&+\{(K-\kappa)(\partial_s),(K-\kappa)(\partial_t)\}\bigr)\,ds\wedge dt.
\end{aligned}
\end{equation}
等价地，使用修改后的联络\eref{17.4}，可写为 $E(u)=\int_S\tfrac12|(Du)^{v,\mathrm{new}}|^2$；此时\eref{17.5}确实是该联络的曲率。\znote{237}{$R$ 的表达式值得逐项读：前两项是通常的 $d(K-\kappa)$，第三项 $\{(K-\kappa)(\partial_s),(K-\kappa)(\partial_t)\}$ 来自 $K-\kappa$ 取值为垂直向量场、取值之间不可交换的二次修正；漏掉它，$R$ 就不闭、能量恒等式也不成立。能量恒等式的意义是：非齐次项 $Y$ 只贡献垂直方向的能量，而右端把它换成边界项减去曲率积分。于是 $R$ 有界时能量受边界项控制，这正是下文紧性论证的出发点；也解释了为何要求 $K$ 在 $\partial^hE\cup\operatorname{Crit}(\pi)$ 附近为零——否则 $R$ 在那里没有一致界。[来源：原书 §(17a)，印刷页 237。]}
\end{note}

\StatementNumber{2}\begin{note}\label{st:17.2}
在此背景下，也有移动边界条件的对应概念。这种边界条件只是任意一族正合 Lagrangian 子流形 $F_z\subset E_z$，$z\in\partial S$。适当的相对扰动数据概念是一对 $(K,J)$，其中不再要求 $K$ 在 $F$ 上的限制为零，而要求：对任意 $X\in TF|F_z$，$z\in\partial S$，若 $D\pi(X)$ 平行于边界，则
\begin{equation}\label{eq:17.6}
d(i_XK)|F_z=\omega_E(\cdot,X)|F_z.
\end{equation}
除此之外，非齐次伪全纯截面的定义不变。其含义应该很清楚：$F_z$ 未必在 $E$ 上给定辛联络的平行移动下保持，但对于满足\eref{17.6}的某个 $K$，转到修改后的形式\eref{17.3}，就恰好恢复了这种相容性。
\end{note}

记 $\M_{E/S}$ 为\eref{17.2}所有解的模空间。在其中任意点 $u$ 处，该方程的线性化由向量丛 $u^*(TE^v)$ 上的 $\bar\partial$ 算子 $D_{E/S,u}$ 给出，边界值在 $\partial S$ 上由 $u^*(TF^v)$ 给定。如通常一样，若 $D_{E/S,u}$ 对所有 $u$ 都满射，则称 $(K,J)$ 正则；此时模空间为闭光滑流形。一般选择 $(K,J)$ 就具有这一性质。取任意这样的选择，可通过对孤立点作模二计数，定义一个简单模型不变量：
\begin{equation}\label{eq:17.7}
\Phi_{E/S}=\#\M_{E/S}\in\Z/2.
\end{equation}

\clearpage
\StatementNumber{3}\begin{example}\label{st:17.3}
令 $\pi:E\to S$ 为\stref{例题}{15.4}和\stref{例题}{16.5}中的局部模型。赋予它 Lagrangian 边界条件 $F$，取为 $|z|=r$ 时标准消失圈 $F_z=\sqrt zS^n$ 之并。在 $F$ 上有 $x_j=(\pi(x)/r)\bar x_j$，故
\[
\theta_E|F=\frac i4\sum_j(x_j\,d\bar x_j-\bar x_j\,dx_j)
=\pi^*\left(\frac{i}{8r}(z\,d\bar z-\bar z\,dz)\right).
\]
选取 $\alpha_F$，使 $\alpha_F+\theta_S|\partial S$ 对应于右边的形式，并取 $h_F=0$，就表明 $F$ 是 Lagrangian 边界条件。当然，这一点此前已经知道，因为 $\sqrt zS^n$ 在任意方向的平行移动下都保持。
\SourcePage{246}
取标准复结构 $J=I_E$ 和平凡非齐次项 $K=0$。所得模空间 $\M_{E/S}$ 可显式确定：它由映射
\begin{equation}\label{eq:17.8}
u_a:S\longrightarrow E\subset\CC^{n+1},\qquad
u_a(z)=r^{-1/2}az+r^{1/2}\bar a
\end{equation}
组成，其中 $a\in\CC^{n+1}$ 满足 $a_1^2+\cdots+a_{n+1}^2=0$ 及 $\|a\|^2=1/2$，见\bcite{131}引理 2.16。此外，这些全都是模空间的正则点。严格地说，还须稍加小心，即\eref{15.10}中的 $s$ 应充分大，使所有这些截面的像确实包含在 $E$ 中。这与 $\partial^hE$ 并非全纯凸有关，也正是这一点使 $E$ 不能成为完全有效的正合 Lefschetz 纤维化例子。暂且忽略这一缺陷，就会发现不变量\eref{17.7}因维数理由而平凡，因为模空间没有孤立点。
\end{example}

这一计算是一个更一般消失结果的模型。令 $\pi:E\to S$ 为底空间 $S$ 是圆盘、且只有一个临界点的 Lefschetz 纤维化。选择 $\partial S$ 上一点 $z$，及从 $z$ 到临界值的消失路径 $\gamma$。记 $M=E_z$、$V=V_\gamma$。由\eref{16.6}，沿边界回路 $\lambda$ 的单值变换 $h_\lambda$ 与 $\tau_V$ 同痕，且同痕在保持端点不动的同伦意义下唯一。这意味着，可以如下为纤维化赋予移动 Lagrangian 边界条件 $F$：从 $F_z=V$ 出发，通过平行移动绕边界几乎一周，再用同痕 $h_\lambda(V)\simeq\tau_V(V)=V$ 将结果接回 $F_z$。称之为\emph{标准边界条件}。

\StatementNumber{4}\begin{lemma}\label{st:17.4}
$\Phi_{E/S}=0$。
\end{lemma}
\begin{proof}
这里只概述证明；细节见\bcite{131}第 2.3 节。为简化记号，假设 $S=D$ 为单位圆盘，且 $\operatorname{Critv}(\pi)=\{0\}$。作不影响不变量的局部变形后，可假设临界点附近存在全纯 Morse 坐标，辛形式在其中标准化。考虑将给定纤维化限制到更小圆盘 $S^r=rD$ 所得的 Lefschetz 纤维化族 $\pi^r:E^r\to S^r$，$r\in(0,1]$。从 $F$ 出发，可为它们赋予光滑边界条件族 $F^r$，在 $r$ 小时等于\stref{例题}{17.3}局部模型中的边界条件。为 $(E^r,F^r)$ 取一族扰动数据，在 $r$ 小时平凡。此时，紧致性论证\bcite{131}引理 2.15 表明，在极限 $r\to0$ 中，$\M_{E^r/S^r}$ 的所有元素都包含在临界点任意小的指定邻域内，因而回到\stref{例题}{17.3}的情形。最后，用常规配边论证将结果带回原来的 $(E,F)$。
\end{proof}

\SourcePage{247}
作为补充，再提及一个较初等但普遍性较弱的证明，它只依赖于模空间的维数。先需要关于 $(E,F)$ 拓扑的一个基本事实：

\StatementNumber{5}\begin{lemma}\label{st:17.5}
若相对第一 Chern 类 $2c_1(M,V)\in H^2(M,V)$ 消失，则相应的 $2c_1(E,F)\in H^2(E,F)$ 等于 $H^2(S,\partial S)$ 标准生成元的拉回的 $(n+1)$ 倍。
\end{lemma}
\begin{proof}
由\secref{15c}可知，$E$ 可赋予二次复体积形式，故 $2c_1(E)=0$，这蕴含 $2c_1(E,F)$ 来自 $H^1(F)$ 的元素。对 $n>1$，这意味着它是 $H^2(S,\partial S)$ 某个元素在 $\pi^*$ 下的像。对 $n=1$，注意 $2c_1(E,F)$ 到 $H^2(\Delta,V)$ 的限制必因 $\Delta$ 的 Lagrangian 性质而消失，也能得出同样结论；这里 $\Delta$ 为 Lefschetz 顶针。有了这一认识，论证便归结为局部模型中的计算，这里省略。
\end{proof}

模空间 $\M_{E/S}$ 在点 $u$ 处的预期维数由算子 $D_{E/S,u}$ 的指标给出。应用相对 Riemann–Roch 公式，即稍作改写的\stref{引理}{11.7}，及\stref{引理}{17.5}，得到
\[
\operatorname{index}(D_{E/S,u})
=n+\int_S2c_1(u^*(TE^v),u^*(TF^v))=2n-1,
\]
这当然与局部模型中看到的一致。特别地，无论选择何种正则相对扰动数据，$\M_{E/S}$ 都不会有孤立点。这与前一个更一般的配边论证所得结果有所不同；前者只保证这些孤立点的贡献相互抵消。\znote{239}{指标公式给出的是模空间的“形式维数” $2n-1$，与例题 17.3 局部模型中看到的维数一致；要点是 $2n-1>0$，故正则情形下根本没有孤立点，$\Phi_{E/S}=0$ 便因维数理由自动成立。两种论证的差别正在此处：配边论证只说明孤立点的贡献成对抵消，维数论证则说明孤立点压根不存在。真正的困难始终是 $\M_{E/S}$ 的紧性——紧性失效时，形式维数与实际维数可以不符，孤立点仍可能出现，这才是引理 17.4 必须单独证明的原因。[来源：原书 §(17a)，印刷页 239。]}

\subsection{带条带状端的纤维化}\label{sec:17b}
令 $M$ 为带角正合辛流形，$\widehat\pi:\widehat E\to\widehat S$ 为 Lefschetz 纤维化。假设底空间上有有限边界标点集 $\Sigma=\Sigma^-\sqcup\Sigma^+$，并在每个标点处将纤维 $\widehat E_\zeta$ 与 $M$ 认同。还假设纤维化在标点附近局部平凡。这意味着，每个 $\zeta$ 有邻域 $\widehat U=\widehat U_\zeta\subset\widehat S$，使
\[
\begin{aligned}
(&\widehat\pi|\widehat U:\widehat E|\widehat U\to\widehat U,
\omega_{\widehat E},\theta_{\widehat E},I_{\widehat E})\\
&=(\widehat U\times M\to\widehat U,
\omega_{\widehat S}+\omega_M,\theta_{\widehat S}+\theta_M,I_{\widehat S}\times I_M).
\end{aligned}
\]
照常取 $S=\widehat S\setminus\Sigma$；对每个 $\zeta$ 选择条带状端 $\epsilon_\zeta$，其像位于 $U_\zeta=\widehat U_\zeta\setminus\{\zeta\}$ 内。于是 $\widehat\pi$ 的限制 $\pi:E\to S$ 在各条带状端上带有典范平凡化 $Z^\pm\times M\to E$。称这样的对象为\emph{带条带状端的 Lefschetz 纤维化}。

\StatementNumber{6}\begin{note}\label{st:17.6}
许多情况下，标点附近的局部平凡性并不预先成立，但可如下通过变形补救。令 $\widehat\pi:\widehat E\to\widehat S$
\SourcePage{248}
为在 $\partial^h\widehat E$ 附近平凡的 Lefschetz 纤维化。选择 $\zeta\in\partial\widehat S$，置 $M=\widehat E_\zeta$。通过改变该纤维附近的辛结构及殆复结构，可在 $\zeta$ 附近达到局部平凡性；更准确地，用\secref{15b}的术语，这是沿纤维恒定的变形。与\stref{引理}{15.3}类似，在局部平凡化 $\widehat U\times M\to\widehat E|\widehat U$ 中有
\begin{equation}\label{eq:17.9}
\theta_{\widehat E}=\theta_M+\theta_{\widehat S}+\kappa+dR,
\end{equation}
其中 $\kappa$ 与函数 $R$ 均在 $\widehat U\times\partial M$ 附近消失。希望用 $\psi(\kappa+dR)$ 替换\eref{17.9}中的这两项，$\psi$ 是 $\widehat U$ 上在 $\zeta$ 附近为零的适当截断函数。修改后的一形式的外导数在总空间上未必仍为辛，但加上底空间某个适当项的拉回，就容易补救。类似地，在局部平凡化中，殆复结构 $I_{\widehat E}$ 在 $\widehat U\times\partial M$ 附近等于 $I_{\widehat S}\times I_M$；容易修改它，使其在整个纤维 $\widehat E_\zeta$ 附近都具有相同的乘积形式。
\end{note}

须将 Lagrangian 边界条件与相对扰动数据的概念扩展到条带状端背景。对前者，取 $F\subset E$，使其到底空间 $S$ 的投影为固有映射，并满足\eref{17.1}。则族 $F_z$ 在每个条带状端上自动局部恒定，即存在一对 Lagrangian 子流形 $(L_{\zeta,0},L_{\zeta,1})$，使对所有 $s$ 有 $F_{\epsilon_\zeta(s,k)}=L_{\zeta,k}$。此外，在条带状端上必有 $\alpha_F=0$。现假设已为每个 $(L_{\zeta,0},L_{\zeta,1})$ 固定 Floer 数据 $(H_\zeta,J_\zeta)$。则相对扰动数据仍为如前的一对 $(K,J)$，但另要求在每个条带状端上有 $K(s,t,x)=H_\zeta(t,x)\,dt$、$J(s,t,x)=i\times J_\zeta(t,x)$。取\eref{17.2}的解 $u$，利用每个端上的平凡化，写成 $(u\circ\epsilon_\zeta)(s,t)=(s,t,u_\zeta(s,t))$，其中 $u_\zeta:Z^\pm\to M$。于是\eref{17.2}归结为 $u_\zeta$ 的 Floer 方程，故可合理要求其渐近收敛于某个 $y_\zeta\in\mathcal C(L_{\zeta,0},L_{\zeta,1})$。\znote{240}{条带状端背景下的三个技术要点值得串起来看。一、“投影固有”不可省：$F_z$ 须有界地覆盖到 $S$ 的端点，端上的极限 $L_{\zeta,k}$ 才存在。二、端上 $F$ 自动局部恒定，故每条端只记录一对 $(L_{\zeta,0},L_{\zeta,1})$，$\alpha_F$ 也随之在端上为零；这一步同样依赖投影固有性。三、条件 $K=H_\zeta(t,x)\,dt$、$J=i\times J_\zeta$ 是让\eref{17.2}在端上恰好退化为 Floer 方程的自然选择，于是每条解在端上有极限 $y_\zeta$，计数才落在 $CF^{\mathrm{pr}}$ 上。[来源：原书 §(17b)，印刷页 240–241。]}所得模空间记为 $\M_{E/S}(\{y_\zeta\})$。假设满足适当正则性条件，就可如\eref{8.20}，对 $\M_{E/S}(\{y_\zeta\})$ 中孤立点作模二计数，定义链同态
\begin{equation}\label{eq:17.10}
C\Phi_{E/S}:\bigotimes_{\zeta^+\in\Sigma^+}CF^{\mathrm{pr}}(L_{\zeta^+,0},L_{\zeta^+,1})
\longrightarrow\bigotimes_{\zeta^-\in\Sigma^-}CF^{\mathrm{pr}}(L_{\zeta^-,0},L_{\zeta^-,1}).
\end{equation}
在 Floer 上同调上诱导的映射 $\Phi_{E/S}$ 与扰动数据的选择无关，并满足 TQFT 公理的直接推广。特别地，有一个复合律：将两个 Lefschetz 纤维化 $\pi_k:E_k\to S_k$，通过把 $S_1$ 的一个出向端与 $S_2$ 的一个入向端认同来粘合。当然，$S_k$ 上辛结构通常不一致，但可如\stref{笔记}{17.6}，通过局部拼接处理。

\SourcePage{249}
这一框架既包含前面考虑的 $S$ 紧致的情形，也包含第\sr{sec:8}{8 节}的理论；后者对应于只允许平凡纤维化 $E=S\times M$。借用这些情形中的想法，还可沿若干方向进一步扩展：
\begin{itemize}
\item \emph{移动边界条件。}如紧致情形中所讨论，但另须要求 $F_z$ 在每个条带状端上关于 $z$ 局部恒定。
\item \emph{带参数的模空间。}在基本构造中已经出现 Lefschetz 纤维化族及相应模空间，用它们为不同相对扰动数据选择所给出的映射\eref{17.10}构造链同伦。一般结构是：考察依赖于参数空间中 $r\in R$ 的族 $E^r\to S^r$，带 Lagrangian 边界条件 $F^r$ 和扰动数据 $(K^r,J^r)$，而条带状端结构关于 $r$ 恒定。为简单起见，假设 $R$ 紧致。相应的 Floer 上链群之间的映射，满足形如 $\partial\circ(\text{映射})+(\text{映射})\circ\partial+(\text{边界项})=0$ 的等式，其中边界项来自将族限制到 $\partial R$。特别地，$R=[0,1]$ 时，得到对应于两个端点的两映射之间的链同伦。后面会看到几个双参数例子；进一步说，可以沿第\sr{sec:9}{9 节}的思路建立一般理论，但对所得代数结构的意义，目前只理解了一部分。
\item \emph{求值映射。}除条带状端外，还可在 $\partial S$ 上固定有限点集 $\Pi=\Pi^-\sqcup\Pi^+$，并考虑这些点处的求值映射 $\operatorname{ev}_{z^\pm}(u)=u(z^\pm)$。由此得到的不变量，如\eref{8.24}，在张量积中具有额外因子 $H^*(F_{z^\pm};\K)$。例如，在\stref{例题}{17.3}的局部模型中，任意点处的求值映射可认同为从切球面丛到底空间的投影 $S(TS^n)\to S^n$。从根本上说，引入标点并不使理论更一般，因为处理技术细节并使用 PSS 同构后，可用条带状端替代这些标点；但它在计算中往往更方便。对此情形还有复合律的一种变体：沿标点取底曲面的边界连通和。
\end{itemize}

\subsection{一个消失定理}\label{sec:17c}
最后提到的复合律，会在后续讨论中多次使用，因此值得说得更准确些。为简单起见，令 $\pi_k:E_k\to S_k$，$k=1,2$，为两个 Lefschetz 纤维化，第一个带条带状端，第二个具有紧底空间，各自带有可能移动的 Lagrangian 边界条件 $F_k$。此外，还要有边界标点 $z_1\in\partial S_1$、$z_2\in\partial S_2$，使纤维 $(E_k)_{z_k}$
\SourcePage{250}
都可与共同的 $M$ 认同，且 $(F_k)_{z_k}$ 成为同一子流形 $L\subset M$。还假设两个纤维化在标点附近局部平凡，边界条件局部恒定。选择相对扰动数据 $(K_k,J_k)$，在 $(E_k)_{z_k}$ 附近平凡，即 $K_k=0$、$J_k=I_{E_k}$。现在在标点处作 $S_1,S_2$ 的边界连通和。可从两种方式理解：一种是先去掉 $z_k$，在所得穿孔曲面上引入条带状端，再对某个粘合长度 $l$，如\eref{8.5}进行端的标准粘合；另一种是在 $z_k$ 附近取以半圆盘 $D_+=D\cap H$ 为模型的局部坐标，再对某个 $\rho\in(-1,0)$ 使用形如 $z\mapsto\rho z^{-1}$ 的过渡映射。两种图景完全等价，具有通常的关系 $l=-\log(-\rho)/\pi$。\znote{242}{原书此处印作 $l=-\log(\rho)/\pi$，但同句已声明 $\rho\in(-1,0)$，负实数在实对数下无定义，故此式只能按 $-\log|\rho|/\pi$ 读，译文径写 $-\log(-\rho)/\pi$。因子 $\pi$ 而非 $2\pi$ 来自模型：标点附近取的是半圆盘 $D_+=D\cap H$，其角为 $\pi$，而 $z\mapsto\rho z^{-1}$ 是穿过该角的过渡映射；绕标点一周的“角度”与 $|\rho|$ 的对数成正比，比例即半圆的 $\pi$。相应地，$|\rho|$ 越小，粘合长度 $l$ 越大，这与 $\rho\to0$ 对应长粘合的说法一致。[来源：原书 §(17c)，印刷页 242；对数定义域为初等核对。]}提及它们，是因为两者各有直观性，并暗示两个略有不同但同样可行的分析框架：第一种得到通常 Floer 理论的 Morse–Bott 版本，第二种接近相对 Gromov–Witten 理论中全纯圆盘的粘合程序。不论哪一种，所得空间均带有 Lefschetz 纤维化 $\pi:E\to S$，以及诱导的边界条件 $F$ 和相对扰动数据 $(K,J)$；基本粘合定理\bcite{131}等式 (2.28) 说明，若满足适当正则性条件，且 $l$ 大，或 $\rho$ 接近 $0$，则
\begin{equation}\label{eq:17.11}
\M_{E/S}(\{y_\zeta\})^0\cong
\coprod_{p+q=n}\M_{E_1/S_1}(\{y_\zeta\})^p\times_L\M_{E_2/S_2}^q.
\end{equation}
这里上标表示维数：右边的维数如此安排，使关于求值映射 $\operatorname{ev}_{z_1}:\M_{E_1/S_1}(\{y_\zeta\})\to(F_1)_{z_1}\cong L$ 及 $\operatorname{ev}_{z_2}:\M_{E_2/S_2}\to(F_2)_{z_2}\cong L$ 所取的纤维积为零维。上述正则性假设包括 $\operatorname{ev}_{z_1},\operatorname{ev}_{z_2}$ 相互处于一般位置。现在将其转成相应不变量的结论。对 $E_2$，考虑由 $\operatorname{ev}_{z_2}$ 表示的同调类，记为 $\Phi_{E_2/S_2,z_2}\in H_*(L;\K)$。对 $E$，则直接取\eref{17.10}。

\StatementNumber{7}\begin{lemma}\label{st:17.7}
若 $\Phi_{E_2/S_2,z_2}$ 消失，则 $\Phi_{E/S}$ 也消失。
\end{lemma}
\begin{proof}
证明直接：取以 $\operatorname{ev}_{z_2}$ 为边界的（伪）循环 $b:B\to L$。则 $\M_{E_1/S_1}(\{y_\zeta\})\times_L B$ 的零维部分给出 Floer 上链群之间的映射，它是 $\Phi_{E/S}$ 的零同伦。
\end{proof}

\clearpage
\StatementNumber{8}\begin{note}\label{st:17.8}
我们对该引理的所有应用中，紧的部分都如\stref{引理}{17.4}。此时可加上假设 $2c_1(M,L)=0$，并得到稍强的结果。若 $n>1$，空间 $\M_{E_2/S_2}$ 的维数为 $2n-1>n$，所以纤维积\eref{17.11}为空。在余下的 $n=1$ 情形，唯一非平凡的项为 $p=0,q=1$。由\stref{引理}{17.4}，已知求值映射
\[
\operatorname{ev}_{z_2}:\M_{E_2/S_2}^1\longrightarrow L
\]
\SourcePage{251}
具有偶数次。对每个孤立点 $u_1\in\M_{E_1/S_1}(\{y_\zeta\})^0$，由于横截性假设，$\operatorname{ev}_{z_1}(u_1)$ 是 $\operatorname{ev}_{z_2}$ 的正则值，故满足 $\operatorname{ev}_{z_2}(u_2)=\operatorname{ev}_{z_1}(u_1)$ 的映射 $u_2$ 的数目为偶数。这说明纤维积仍包含偶数个点。结论是，对这一特定的扰动数据选择，上链层次的不变量 $C\Phi_{E/S}$ 本身就是零。
\end{note}

\subsection{基本态射}\label{sec:17d}
令 $M$ 为带角正合辛流形，$V$ 为其中带标架的正合 Lagrangian 球面。为方便利用\stref{笔记}{17.8}中上链层次的消失结果，假设 $2c_1(M,V)$ 消失。置 $L_0=V$，另取两个正合 Lagrangian 子流形 $L_1,L_2$，要求 $L_2$ 与 $\tau_V(L_1)$ 同痕；更准确地，还要指定一个同痕选择。为子流形对 $(L_0,L_1)$、$(L_0,L_2)$、$(L_1,L_2)$ 固定 Floer 数据，并固定定义乘积 $\mu^2:CF^{\mathrm{pr}}(L_1,L_2)\otimes CF^{\mathrm{pr}}(L_0,L_1)\to CF^{\mathrm{pr}}(L_0,L_2)$ 所需的扰动数据。为便于与代数理论比较，用 $\Y_0,\Y_1,\Y_2$ 表示 $\F^{\mathrm{pr}}(M)$ 中由 $L_0,L_1,L_2$ 给出的对象。

固定圆盘上的 Lefschetz 纤维化 $\pi:E\to S$，纤维为 $M$，只有一个临界点，消失圈为 $V$。更准确地，要求该纤维化带有标点 $*\in\partial S$，并将该点上方纤维 $E_*$ 与 $M$ 认同；还带有从 $*$ 出发的消失路径 $\gamma$，使 $V_\gamma=V$。为简化技术，进一步假设纤维化在 $*$ 附近局部平凡。$E$ 将是后续全部构造的基本组成部分，这些构造基本上是往它上面粘合各种平凡部分。

\StatementNumber{9}\begin{note}\label{st:17.9}
回顾，对任意给定的 $M,V$，都存在这样的纤维化，而且可以“手工”构造；概要见\secref{16e}，完整论述见\bcite{131}引理 1.10。仔细关注构造细节，还能获得更多控制：辛联络在临界点的小邻域之外平坦；任意消失路径的消失圈均等于 $V$；任意绕临界值一周的回路的单值变换，都是一个 Dehn 扭转 $\tau_V$，而不只是与之同痕。不过，这些性质对于我们的构造都并非必要。
\end{note}

作为第一个也是最初等的应用，定义 Floer 上闭链
\begin{equation}\label{eq:17.12}
\boxed{\begin{gathered}
c\in CF^{\mathrm{pr}}(L_1,L_2),\\ \mu^1(c)=0.
\end{gathered}}
\end{equation}
\SourcePage{252}
为此，取给定的 $E\to S$ 并去掉 $*$。所得 Lefschetz 纤维化 $E^c\to S^c$ 有一个无穷远点，将其视为入向点，记为 $\zeta^{c,-}$，并选择条带状端 $\epsilon_{\zeta^{c,-}}$。存在逐纤维同痕意义下典范的移动 Lagrangian 边界条件 $F^c$，定义如下：在 $\epsilon_{\zeta^{c,-}}(s,0)$，$s\ll0$，上方的纤维中从 $L_1$ 出发，通过平行移动沿 $\partial S^c$ 几乎绕行一周。结果与 $\tau_V(L_1)$ 同痕，因而也与 $L_2$ 同痕；用该同痕将其连续延伸，使在 $\epsilon_{\zeta^{c,-}}(s,1)$，$s\ll0$，上方的纤维中等于 $L_2$。

\sr{fig:17.1}{图~17.1}示意这一构造。图形记号的含义是：去掉奇异纤维，再沿虚线切开底空间，得到平凡辛纤维化。随后将单值变换看成切口两侧之间的认同。在此类图中，通常不区分同痕的对象，所以直接将单值变换写为 $\tau_V$。
\SourceFigure{figure-17-1}{0.37176883}{17.1}
于是，\eref{17.12}就是相应不变量 $C\Phi_{E^c/S^c}$，取某个与已有 Floer 数据选择相容的正则相对扰动数据。

第二项数据是 $\mu^2(c,\cdot)$ 与零映射之间的链同伦，即
\begin{equation}\label{eq:17.13}
\boxed{\begin{gathered}
k:CF^{\mathrm{pr}}(L_0,L_1)\longrightarrow CF^{\mathrm{pr}}(L_0,L_2),\\
\mu^1(k(\cdot))+k(\mu^1(\cdot))+\mu^2(c,\cdot)=0.
\end{gathered}}
\end{equation}
如所预期，这是由依赖于参数 $r\in[0,1]$ 的带条带状端 Lefschetz 纤维化族构造的。在一端 $r=0$，取三标点圆盘上的平凡纤维化，带 Lagrangian 边界条件 $(L_0,L_1,L_2)$。选取适当扰动数据，其关联不变量正好给回 $\mu^2$。现取两个出向端中标记为 $(L_1,L_2)$ 的那个，与 $E^c\to S^c$ 的入向端粘合。粘合长度应充分大，使零维伪全纯截面模空间的粘合公式在上链层次成立。所得纤维化记为 $E^{k,0}\to S^{k,0}$，配备移动边界条件 $F^{k,0}$
\SourcePage{253}
和相对扰动数据，其相应不变量恰为 $\mu^2(c,\cdot)$。\sr{fig:17.2}{图~17.2}给出两个图景，只在我们设想中切开底空间的方式上有所不同。在另一端 $r=1$，按如下方式进行。
\SourceFigure{figure-17-2}{0.48184563}{17.2}

取 $Z$ 上的平凡纤维化。赋予移动 Lagrangian 边界条件，在 $\R\times\{0\}$ 上等于 $L_0$，在 $\R\times\{1\}$ 上实现从 $L_2$ 到 $\tau_V(L_1)$ 的同痕。这意味着，两端关联的 Lagrangian 子流形对分别是 $(L_0,L_2)$ 与 $(L_0,\tau_V(L_1))$。对于后一对，取 $(L_0,L_1)$ 的 Floer 数据在 $\tau_V$ 下的像，同时记住 $\tau_V(L_0)=\tau_V(V)=V=L_0$。还要标记边界点 $z^+\in\R\times\{0\}$。现取原来的 $\pi:E\to S$，赋予标准边界条件 $F$，使其在 $*$ 上方纤维中等于 $L_0$，且在该点附近局部恒定。最后，沿 $z^+$ 与 $z^-=*$ 取边界连通和，将两个纤维化粘合，结果记为 $E^{k,1}\to S^{k,1}$，见\sr{fig:17.3}{图~17.3}。这是\stref{引理}{17.7}讨论的过程的一个实例，事实上\stref{笔记}{17.8}也适用；因此，选择适当
\SourcePage{254}
小的粘合参数，可保证 $E^{k,1}$ 的所有零维伪全纯截面模空间，在 $n>1$ 时为空，在 $n=1$ 时包含偶数个点。
\SourceFigure{figure-17-3}{0.3489227}{17.3}

直接比较\sr{fig:17.2}{图~17.2}和\sr{fig:17.3}{图~17.3}可见，确有单参数纤维化族 $E^{k,r}\to S^{k,r}$ 及边界条件 $F^{k,r}$，连接我们在 $r\in\{0,1\}$ 时构造的对象。这些纤维化的无穷远结构，即条带状端上的结构，对所有 $r$ 相同。选择一族相对扰动数据，在 $r=0,1$ 时与之前构造的相同，并考虑所得带参数伪全纯截面模空间。计数其中孤立点，就定义同态 $k:CF^{\mathrm{pr}}(L_0,L_1)\cong CF^{\mathrm{pr}}(L_0,\tau_V(L_1))\to CF^{\mathrm{pr}}(L_0,L_2)$，它是 $C\Phi_{E^{k,0}/S^{k,0}}$ 与 $C\Phi_{E^{k,1}/S^{k,1}}$ 之间的链同伦，故恰具有所需性质\eref{17.13}。

构造 $c$ 时必须作一些选择，因为边界条件及扰动数据都不完全典范。即使指定了它们，定义族 $E^{k,r}\to S^{k,r}$ 时仍有进一步选择，又会影响 $k$。但两种情况下，参数空间都道路连通，故能在任意两种可能之间光滑移动。具体地，若 $(c,k)$ 与 $(\tilde c,\tilde k)$ 用两种不同选择定义，则可考虑多一个参数的模空间，从而找到
\begin{equation}\label{eq:17.14}
\boxed{\begin{gathered}
\gamma\in CF^{\mathrm{pr}}(L_1,L_2),\\
\kappa:CF^{\mathrm{pr}}(L_0,L_1)\longrightarrow CF^{\mathrm{pr}}(L_0,L_2),\\
\mu^1(\gamma)+c+\tilde c=0,\\
\mu^1(\kappa(\cdot))+\kappa(\mu^1(\cdot))+\mu^2(\gamma,\cdot)+k+\tilde k=0.
\end{gathered}}
\end{equation}
这只是此前考察的直接扩展，因此不
\SourcePage{255}
解释细节。下面改用第\sr{sec:5}{5 节}的代数理论来解释所得结果。

\StatementNumber{10}\begin{note}\label{st:17.10}
注意，Floer 理论框架尚未引入分次，而且我们在 $\operatorname{char}(\K)=2$ 下工作，符号是平凡的；上面的公式清楚地体现了二者的缺席。不过，\stref{笔记}{9.7}在此也适用：首先，为不分次 $\Ai$ 范畴改写代数理论相当容易；其次，当前背景只是暂时的，后面会加以修改，纳入适当的分次与符号。因此，我们可以直接使用第\sr{sec:5}{5 节}的结果。
\end{note}

具体地，\stref{引理}{5.3}(i) 结合\eref{17.12}与\eref{17.13}表明，$(c,k)$ 在 $\Ai$ 模的上同调范畴中给出态射 $T_{\Y_0}(\Y_1)\to\Y_2$。类似地，\stref{引理}{5.3}(ii) 与\eref{17.14}说明该态射典范；即它与用来定义 $c,k$ 的 Lefschetz 纤维化和模空间的定义中的辅助选择无关。\znote{247}{这一段是本节的落点：$c$ 是上闭链、$k$ 把 $\mu^2(c,\cdot)$ 同伦到零，两者合起来正是“$\mathcal T_{Y_0}(Y_1)\to Y_2$ 是 $\Ai$ 模态射”在链级的完整内容，用第 5 节的语言即给出同态 $(\Id,c,k,\ldots)$。之所以要求 $L_2$ 与 $\tau_V(L_1)$ 同痕，是因为 $c$ 的边界条件沿底空间绕一圈后正好把 $L_1$ 送到 $\tau_V(L_1)$；若 $L_2$ 不与之同痕，$c$ 便没有落点。这一构造是 §(17g)–(17j) 反复使用的部件，也是定理 17.20 的出发点；该定理断言：在拟等价意义下，$\Tw\F^\to(V^\#)$ 是 $V$ 的 Hurwitz 等价类的不变量，其证明的代数核心正是“几何 Dehn 扭转 $\tau_V$ 对应扭转函子 $T_{Y_0}$”。[来源：原书 §(17d) 与 §(17g)，印刷页 244–247；定理编号见原书目录。]}

\clearpage
\subsection{一个辅助链复形}\label{sec:17e}
另取正合 Lagrangian 子流形 $Q$。在此前选择之外，为子流形对 $(Q,L_k)$，$k=0,1,2$，固定 Floer 数据，并固定定义三角形乘积
\[
\begin{aligned}
\mu^2 &:CF^{\mathrm{pr}}(L_0,L_1)\otimes CF^{\mathrm{pr}}(Q,L_0)\longrightarrow CF^{\mathrm{pr}}(Q,L_1),\\
\mu^2 &:CF^{\mathrm{pr}}(L_1,L_2)\otimes CF^{\mathrm{pr}}(Q,L_1)\longrightarrow CF^{\mathrm{pr}}(Q,L_2)
\end{aligned}
\]
所需的扰动数据，以及与这些数据一致的四标点圆盘万有族上的数据，后者给出
\[
\mu^3:CF^{\mathrm{pr}}(L_1,L_2)\otimes CF^{\mathrm{pr}}(L_0,L_1)\otimes CF^{\mathrm{pr}}(Q,L_0)
\longrightarrow CF^{\mathrm{pr}}(Q,L_2).
\]
略加改变上述想法，即可定义映射
\begin{equation}\label{eq:17.15}
\boxed{\begin{gathered}
d:CF^{\mathrm{pr}}(Q,L_1)\longrightarrow CF^{\mathrm{pr}}(Q,L_2),\\
l:CF^{\mathrm{pr}}(L_0,L_1)\otimes CF^{\mathrm{pr}}(Q,L_0)\longrightarrow CF^{\mathrm{pr}}(Q,L_2),\\
\mu^1(d(\cdot))+d(\mu^1(\cdot))=0,\\
d(\mu^2(\cdot,\cdot))+\mu^1(l(\cdot,\cdot))
+l(\mu^1(\cdot),\cdot)+l(\cdot,\mu^1(\cdot))=0.
\end{gathered}}
\end{equation}
为以更紧凑的形式编码这些关系，定义
\begin{equation}\label{eq:17.16}
\begin{gathered}
D=CF^{\mathrm{pr}}(L_0,L_1)\otimes CF^{\mathrm{pr}}(Q,L_0)
\oplus CF^{\mathrm{pr}}(Q,L_1)\oplus CF^{\mathrm{pr}}(Q,L_2),\\
\partial=\begin{pmatrix}
\mu^1\otimes\Id+\Id\otimes\mu^1&0&0\\
\mu^2&\mu^1&0\\ l&d&\mu^1
\end{pmatrix}\in\End(D).
\end{gathered}
\end{equation}
于是\eref{17.15}等价于 $\partial^2=0$。映射 $d,l$ 实际上就是构造 Floer 上同调长正合列所用的映射，见\bcite{131}第 3.3、3.4 节，\znote{247}{\eref{17.16} 是长正合列的“链级包装”：$\partial$ 的第一行是 $\mu^1$ 在张量积上的 Leibniz 作用，第二行是 $\mu^2$，第三行是 $l$ 与 $d$；$\partial^2=0$ 与\eref{17.15}的两式逐一对应。这种写法的好处是把“$\mu^1$ 加更高阶修正”压成单个微分，从而可以整体做线性代数：下文与代数一侧的锥比较时，正需要把 $(d,l)$ 放进矩阵的第三行，而不是逐条处理关系式。[来源：原书 §(17e)，印刷页 247–248。]}
\SourcePage{256}
而\eref{17.16}在其证明中起关键作用，见\bcite{131}引理 2.32。不过，由于这里与\bcite{131}之间有记号及少许技术差异，仍值得简要回顾定义。

对于 $d$，从 $\pi:E\to S$ 出发，不仅去掉 $\zeta^{d,-}=*$，还去掉另一个边界点 $\zeta^{d,+}$；例如从 $*$ 沿正方向稍移一点得到的点，这样做的好处是纤维化在该点附近仍局部平凡。结果为 Lefschetz 纤维化 $E^d\to S^d\cong Z$，有一个入向条带状端和一个出向条带状端。相应边界条件 $F^d$ 如下定义：在“短”的边界分支上恒定等于 $Q$；在“长”的分支上，$\epsilon_{\zeta^{d,+}}(s,1)$，$s\geq0$，上方纤维中为 $L_1$，$\epsilon_{\zeta^{d,-}}(s,1)$，$s\leq0$，上方则为 $L_2$，利用平行移动及同痕 $\tau_V(L_1)\simeq L_2$ 在两者间插值。\sr{fig:17.4}{图~17.4}给出该纤维化的两个图景，只在底空间的画法上不同。赋予相对扰动数据，并定义 $d=C\Phi_{E^d/S^d}$ 为所得不变量。
\SourceFigure{figure-17-4}{0.52546097}{17.4}

$l$ 的构造与 $k$ 的构造十分相似。取三标点圆盘上的平凡纤维化，带边界条件 $(Q,L_0,L_1)$，将其入向端与 $E^d\to S^d$ 的出向端粘合，得到纤维化 $E^{l,1}\to S^{l,1}$，带边界条件 $F^{l,1}$ 及相对扰动数据。另一方面，可取三标点圆盘上的另一个平凡纤维化，这次前两个边界分支上的边界条件分别恒定为 $Q,L_0$，第三个分支上则实现从 $\tau_V(L_1)$ 到 $L_2$ 的同痕。在 $L_0$ 分支上标记一个边界点，再在 $*$ 处与带标准边界条件 $F$ 的纤维化 $E\to S$ 作边界连通和，得到另一个纤维化 $E^{l,0}\to S^{l,0}$；如前，\stref{笔记}{17.8}适用。存在明显的族 $E^{l,r}\to S^{l,r}$ 在两者之间插值，取 $l$ 为所得带参数模空间给出的不变量，见\sr{fig:17.5}{图~17.5}。
\SourceFigure{figure-17-5}{0.51092253}{17.5}

\subsection{补充说明}\label{sec:17f}
如前可证明，定义 $(d,l)$ 时所涉及的选择实质上无关紧要。更准确地，若 $(d,l)$ 与 $(\tilde d,\tilde l)$
\SourcePage{257}
由不同选择定义，则可考虑在两种构造版本之间插值的附加模空间，从而得到
\[
\boxed{\begin{gathered}
\delta:CF^{\mathrm{pr}}(Q,L_1)\longrightarrow CF^{\mathrm{pr}}(Q,L_2),\\
\lambda:CF^{\mathrm{pr}}(L_0,L_1)\otimes CF^{\mathrm{pr}}(Q,L_0)\longrightarrow CF^{\mathrm{pr}}(Q,L_2),\\
\mu^1(\delta(\cdot))+\delta(\mu^1(\cdot))+d+\tilde d=0,\\
\delta(\mu^2(\cdot,\cdot))+\mu^1(\lambda(\cdot,\cdot))
+\lambda(\mu^1(\cdot),\cdot)\\
{}+\lambda(\cdot,\mu^1(\cdot))+l+\tilde l=0.
\end{gathered}}
\]
其代数含义是，在链复形\eref{17.16}之间得到同构
\begin{equation}\label{eq:17.17}
\begin{pmatrix}
\Id\otimes\Id&0&0\\0&\Id&0\\\lambda&\delta&\Id
\end{pmatrix}:D\longrightarrow\widetilde D.
\end{equation}
更一般地，可考察 $(d,l)$ 如何依赖于 $L_1$。假设有与 $L_1$ 同痕的另一个子流形 $\widetilde L_1$，以及相应版本 $(\tilde d,\tilde l)$。取
\SourcePage{258}
延拓映射
\[
\begin{aligned}
a &:CF^{\mathrm{pr}}(Q,L_1)\longrightarrow CF^{\mathrm{pr}}(Q,\widetilde L_1),\\
b &:CF^{\mathrm{pr}}(L_0,L_1)\longrightarrow CF^{\mathrm{pr}}(L_0,\widetilde L_1);
\end{aligned}
\]
它们在链同伦意义下与三角形乘积相容，即存在 $h:CF^{\mathrm{pr}}(L_0,L_1)\otimes CF^{\mathrm{pr}}(Q,L_0)\to CF^{\mathrm{pr}}(Q,\widetilde L_1)$，使
\[
\begin{aligned}
\mu^1(h(\cdot,\cdot))+h(\mu^1(\cdot),\cdot)+h(\cdot,\mu^1(\cdot))\\
{}+\mu^2(b(\cdot),\cdot)+a(\mu^2(\cdot,\cdot))=0.
\end{aligned}
\]
从一般观点看，延拓映射来自平凡纤维化 $M\times Z\to Z$，其中边界条件在 $\R\times\{0\}$ 上恒定等于 $Q$ 或 $L_0$，在 $\R\times\{1\}$ 上则实现从 $\widetilde L_1$ 到 $L_1$ 的同痕；\secref{8j}中已指出这一点。可取三穿孔圆盘上的平凡纤维化，带边界条件 $(Q,L_0,L_1)$，并将其入向端与定义 $a$ 的纤维化的出向端粘合。另一方面，可对三穿孔圆盘取边界条件 $(Q,L_0,\widetilde L_1)$，将第二个出向端与定义 $b$ 的纤维化的入向端粘合。如此构造的两个纤维化可相互变形，所得族给出 $h$。

同样的想法可如下联系 $d$ 与 $\tilde d$：将定义 $a$ 的纤维化的入向端，与 $E^{\tilde d}\to S^{\tilde d}$ 的出向端粘合，再将结果变形回原来的 $E^d\to S^d$。结果是单参数纤维化族 $E^{\alpha,r}\to S^{\alpha,r}$，如\sr{fig:17.6}{图~17.6}；其相应不变量具有形式
\[
\boxed{\begin{gathered}
\alpha:CF^{\mathrm{pr}}(Q,L_1)\longrightarrow CF^{\mathrm{pr}}(Q,L_2),\\
\mu^1(\alpha(\cdot))+\alpha(\mu^1(\cdot))+\tilde d(a(\cdot))+d(\cdot)=0.
\end{gathered}}
\]
\SourceFigure{figure-17-6}{0.62515317}{17.6}

\SourcePage{259}
同类而稍复杂的论证也适用于 $l,\tilde l$，结果具有形式
\[
\boxed{\begin{gathered}
\beta:CF^{\mathrm{pr}}(L_0,L_1)\otimes CF^{\mathrm{pr}}(Q,L_0)\longrightarrow CF^{\mathrm{pr}}(Q,L_2),\\
\mu^1(\beta(\cdot,\cdot))+\beta(\mu^1(\cdot),\cdot)+\beta(\cdot,\mu^1(\cdot))\\
{}+\tilde l(b(\cdot),\cdot)+l(\cdot,\cdot)+\alpha(\mu^2(\cdot,\cdot))
+\tilde d(h(\cdot,\cdot))=0.
\end{gathered}}
\]
合在一起，便得到链映射
\begin{equation}\label{eq:17.18}
\begin{pmatrix}
b\otimes\Id&0&0\\h&a&0\\\beta&\alpha&\Id
\end{pmatrix}:D\longrightarrow\widetilde D.
\end{equation}
\eref{17.18}中的对角项是延拓映射，因而为拟同构，所以整个映射也为拟同构。最后，类似论证表明，$D$ 在 $L_2$ 的同痕下于拟同构意义内保持不变；为避免重复，省略这一部分。

\clearpage
\subsection{联系两种构造}\label{sec:17g}
给定 $(c,k)$，可考虑 $\tilde d=\mu^2(c,\cdot)$ 及 $\tilde l=\mu^2(k(\cdot),\cdot)+\mu^3(c,\cdot,\cdot)$。它们满足\eref{17.15}，并可与 $(d,l)$ 联系。这可视为此前关于\eref{17.15}的解的唯一性断言的一个极限情形，但由于对应用很重要，我们采用独立记号，给出完整细致的论述。具体地，将构造
\begin{equation}\label{eq:17.19}
\boxed{\begin{gathered}
\xi:CF^{\mathrm{pr}}(Q,L_1)\longrightarrow CF^{\mathrm{pr}}(Q,L_2),\\
\psi:CF^{\mathrm{pr}}(L_0,L_1)\otimes CF^{\mathrm{pr}}(Q,L_0)\longrightarrow CF^{\mathrm{pr}}(Q,L_2),\\
\mu^1(\xi(\cdot))+\xi(\mu^1(\cdot))+d+\mu^2(c,\cdot)=0,\\
\mu^1(\psi(\cdot,\cdot))+\psi(\mu^1(\cdot),\cdot)+\psi(\cdot,\mu^1(\cdot))\\
{}+\xi(\mu^2(\cdot,\cdot))+l(\cdot,\cdot)
+\mu^2(k(\cdot),\cdot)+\mu^3(c,\cdot,\cdot)=0.
\end{gathered}}
\end{equation}
对于 $\xi$，先令 $E^{\xi,0}\to S^{\xi,0}$ 就是 $E^d\to S^d$，带之前采用的边界条件及相对扰动数据。在变形的另一端，取三标点圆盘上的平凡纤维化，三个边界分支上的边界条件为 $Q,L_1,L_2$，将其第二个出向端与 $E^c\to S^c$ 的入向端粘合。在两者间插值的族显而易见，见\sr{fig:17.7}{图~17.7}。
\SourceFigure{figure-17-7}{0.53169174}{17.7}

开始定义 $\psi$ 时，先写出五个单参数 Lefschetz 纤维化族，配备边界条件及扰动数据，
\SourcePage{260}
使它们合起来构成闭链：第一族的终点是第二族的起点，依此类推；图示见\sr{fig:17.8}{图~17.8}。第一族是之前定义的 $E^{l,r}\to S^{l,r}$。为了得到第二族，记为 $E^{\xi(\mu^2),r}\to S^{\xi(\mu^2),r}$，取 $E^{\xi,r}\to S^{\xi,r}$，将唯一的出向端与三标点圆盘上的平凡纤维化的入向端粘合；后者带边界条件 $Q,L_0,L_1$，以及定义 $\mu^2$ 乘积的扰动数据。因此，只要粘合长度充分大，计数该纤维化族的伪全纯截面，恰好给出 $\xi(\mu^2(\cdot,\cdot))$。回顾，$E^{l,1}\to S^{l,1}$ 是取 $E^d\to S^d$ 再粘入三标点圆盘上的平凡纤维化得到的；另一方面，由定义有 $(E^{\xi,0}\to S^{\xi,0})=(E^d\to S^d)$。这意味着，若适当选取粘合参数，则
\begin{equation}\label{eq:17.20}
(E^{\xi(\mu^2),0}\to S^{\xi(\mu^2),0})=(E^{l,1}\to S^{l,1}),
\end{equation}
其中等式也包括边界条件与扰动数据。在另一端 $r=1$，得到纤维化 $E^{\xi(\mu^2),1}\to S^{\xi(\mu^2),1}$，给出表达式 $\mu^2(c,\mu^2(\cdot,\cdot))$。这立即暗示下一步：在 $\RG^4$ 中取一个充分大的闭区间 $[0,1]$，考虑其上的四标点圆盘万有族，边界分支标记为 $(Q,L_0,L_1,L_2)$。\znote{252}{$\RG^4$ 是四个边界标点的圆盘模空间：按 (9.7)，$\RG^{d+1}=\operatorname{Conf}_{d+1}(\partial D)/\operatorname{Aut}(D)$，实维数为 $d-2$，故 $\RG^4$ 是一维区间，其 Deligne–Mumford 紧化把两端各缩成一点。所谓“充分大的闭区间”即伸到这两端附近，使某个标点被挤出的两种退化（粘合参数趋于零）也落在族内，于是这个单参数族的两端恰好与相邻的族拼接。[来源：Seidel, Fukaya categories and Picard–Lefschetz theory, §9 中 (9.7) 与 §12g。]}将其变成带适当边界条件与扰动数据的平凡 Lefschetz 纤维化族，再将每个圆盘的第三个出向端与 $E^c\to S^c$ 的入向端粘合。将结果记为 $E^{\mu^3(c\otimes\Id),r}\to S^{\mu^3(c\otimes\Id),r}$，因为计数该族的伪全纯截面就得到 $\mu^3(c,\cdot,\cdot)$。适当选择粘合参数，显然可使
\SourcePage{261}
\SourceFigure{figure-17-8}{0.60438396}{17.8}
\[
(E^{\mu^3(c\otimes\Id),0}\to S^{\mu^3(c\otimes\Id),0})
=(E^{\xi(\mu^2),1}\to S^{\xi(\mu^2),1}),
\]
等式的意义与\eref{17.20}相同。在另一端 $r=1$，得到的 Lefschetz 纤维化的截面定义同态 $\mu^2(\mu^2(c,\cdot),\cdot)$。对于第四族 $E^{\mu^2(k\otimes\Id),r}\to S^{\mu^2(k\otimes\Id),r}$，取边界分支标记为 $(Q,L_0,L_2)$ 的三标点圆盘，将其变成平凡 Lefschetz 纤维化，随后
\SourcePage{262}
将它的第二个出向端与纤维化 $E^{k,r}\to S^{k,r}$ 的入向端粘合。由 $k$ 的定义，
\[
(E^{\mu^2(k\otimes\Id),0}\to S^{\mu^2(k\otimes\Id),0})
=(E^{\mu^3(c\otimes\Id),1}\to S^{\mu^3(c\otimes\Id),1}).
\]
另一方面，在 $r=1$ 时，可将 $E^{\mu^2(k\otimes\Id),1}$ 视为如下构造：从三标点圆盘上的平凡纤维化出发，两个边界分支上的边界条件恒定等于 $Q,L_0$，余下分支上则给出从 $\tau_V(L_1)$ 到 $L_2$ 的同痕。在 $L_0$ 分支上标记一个边界点，如 $k$ 的定义那样，与带标准边界条件的 $E$ 作边界连通和。注意，$E^{l,0}\to S^{l,0}$ 的定义实质上相同。细节，例如标点位置，可能不同，但显然可用第五个也就是最后一个族在两者之间插值，记为 $E^{0,r}\to S^{0,r}$；其性质是，每个 $E^{0,r}$ 都由与 $E$ 作边界连通和得到，且粘合参数\emph{始终保持很小}。由定义，这使我们的链闭合：
\[
\begin{aligned}
(E^{0,0}\to S^{0,0})&=(E^{\mu^2(k\otimes\Id),1}\to S^{\mu^2(k\otimes\Id),1}),\\
(E^{0,1}\to S^{0,1})&=(E^{l,0}\to S^{l,0}).
\end{aligned}
\]
\stref{笔记}{17.8}的消失结果以适当的带参数版本适用于 $E^{0,r}\to S^{0,r}$，因此相应零维伪全纯截面模空间的计数平凡。最后，整个闭链形成以 $S^1$ 为参数的 Lefschetz 纤维化族，显然可用 $D^2$ 上的族填充。用所得带参数模空间定义 $\psi$，则\eref{17.19}由构造成立。

\StatementNumber{11}\begin{note}\label{st:17.11}
尝试将这一程序纳入更一般、更标准化的框架颇有启发。令 $\RG^{3,1}$ 为带一个额外内部标点的三标点圆盘模空间。等价地，也更适合眼前目的，可固定内部点及第一个边界点，从而消除 $\operatorname{Aut}(D)$ 对称性，并将 $\RG^{3,1}$ 看成另外两个边界点的构形空间。取圆盘 $D$ 为标准 Lefschetz 纤维化 $E\to S$ 的底空间，内部标点为唯一临界值，第一个边界点为 $*$。对任意构形 $(z_1,z_2)\in\RG^{3,1}$，可将 $E$ 限制到 $S\setminus\{*,z_1,z_2\}$，使该限制成为带条带状端的 Lefschetz 纤维化。为这一纤维化族赋予由 $Q,L_0,L_1$ 及从 $\tau_V(L_1)$ 到 $L_2$ 的同痕组成的 Lagrangian 边界条件，并考虑由所得带参数模空间给出的不变量。

为理解它所满足的关系，需要考察 Deligne–Mumford 型紧化 $\overline{\RG}^{3,1}$，它是一个六边形。与 $\overline{\RG}^5$ 一样，可将其实现为复 Deligne–Mumford 空间 $\overline{\M}_{0,5}$ 关于某个反全纯对合的实点集，
\SourcePage{263}
该对合将复共轭与两个标点的交换结合。对应于各层的退化示于\sr{fig:17.9}{图~17.9}。\znote{255}{六条边对应 $\overline{\RG}^{3,1}$ 的两类余维一退化：内部标点（即临界值）撞上某个边界标点，三种；两个边界标点相撞，也是三种；合计正是六，与图 17.9 周围六幅退化图一一对应。这也解释了何以只有边 (i)、(ii) 直接对应图 17.8 中的族：其余四条边给出的族只到“并的、保持端点不动的变形”为止，其贡献并入同一条链映射恒等式，而不成为独立项。[来源：Seidel, Fukaya categories and Picard–Lefschetz theory, §17g 与图 17.9。]}对于边 (i)、(ii)，所得单参数族是\sr{fig:17.8}{图~17.8}中相应族在粘合参数趋于零时的极限，给出\eref{17.19}中的项 $\mu^3(c\otimes\Id)$ 与 $\mu^2(k\otimes\Id)$。图的其余部分稍复杂：给出 $\xi(\mu^2)$ 项的族对应于边 (vi) 的一部分，$E^{0,r}\to S^{0,r}$ 也类似地覆盖边 (iii) 的一部分；这两条边的其余部分，以及 (iv)、(v)，在\sr{fig:17.8}{图~17.8}中没有直接对应物，但它们的并可保持端点不动地变形为给出 $l$ 的单参数族。因此，从图形上说，\sr{fig:17.8}{图~17.8}的五个族可看成\sr{fig:17.9}{图~17.9}中的虚线边。实际意义是：若直接使用 $\overline{\RG}^{3,1}$，则得到的是\eref{17.15}与\eref{17.19}稍复杂的版本，含有更多项。
\SourceFigure{figure-17-9}{0.60853781}{17.9}
\end{note}

\StatementNumber{12}\begin{note}\label{st:17.12}
在后续应用中，由于分次的存在，条件 $2c_1(M,V)=0$ 总会成立，见\secref{12a}的讨论。不过，必要时也可不用这一条件，而加入
\SourcePage{264}
涉及适当填充链的修正项，如\bcite{131}引理 3.12 或上述\stref{引理}{17.7}。
\end{note}

\subsection{无上同调性}\label{sec:17h}
将\bcite{131}引理 2.32 证明中的代数论证与该文第 3 章中的几何论证结合起来，即可证明\eref{17.16}所定义的复形 $D$ 无上同调；鉴于\stref{引理}{5.3}(iii)，这一点对我们很重要。\znote{256}{“无上同调”与“$t$ 是同构”实为一事，桥梁是映射锥：引理 5.3 的证明指出 $t^1$ 的映射锥是复形 $\widetilde D$，其第三个分量把 (5.8) 中的 $d(c_1)+l(c_0,b)$ 换成 $\mu^2_{\A}(c,c_1)+\mu^2_{\A}(k(c_0),b)+\mu^3_{\A}(c,c_0,b)$，而 $(\xi,\psi)$ 正是补出这一替换的链同构。17g 节的 $\tilde d=\mu^2(c,\cdot)$、$\tilde l=\mu^2(k(\cdot),\cdot)+\mu^3(c,\cdot,\cdot)$ 就是这份锥数据，故\eref{17.16}无上同调即 $t$ 为拟同构。注意 $D$ 与 $\widetilde D$ 只有第三个分量不同，其余两项完全一样。[来源：Seidel, Fukaya categories and Picard–Lefschetz theory, §5c 引理 5.3(iii) 及其证明。]}严格说来，\bcite{131}只在若干附加技术条件下证明了无上同调性，即：
\begin{enumerate}[label=(\roman*)]
\item $Q,L_0,L_2$ 处于一般位置。对子流形对 $(Q,L_0)$、$(L_0,L_2)$ 和 $(Q,L_2)$，Floer 数据的 Hamilton 项应为零。此外，对作用量泛函\eref{8.11}在交点处的取值还有一个条件：$x_0\in\C(Q,L_2)$ 时的各数 $A(x_0)$，以及 $(x_1,x_2)\in\C(Q,L_0)\times\C(L_0,L_2)$ 时的各数 $A(x_1)+A(x_2)$，合在一起必须两两不同。
\item 必须能将标架 $L_0=V\cong S^n$ 延拓为余切球丛的辛嵌入 $B_s^*S^n\to\operatorname{int}(M)$，其中某个 $s>0$，并使 $\theta_M$ 的拉回为 $T^*S^n$ 上的标准一形式。此外，在这个管状邻域内，$Q$ 与 $L_2$ 都应为若干余切纤维的并。
\item $L_1=\tau_V^{-1}(L_2)$，其中 Dehn 扭转使用先前选定的管状邻域定义。决定扭转速度的函数 $b$，见\eref{16.5}，必须满足某些条件；这些条件依赖于 (i) 中作用量值之间的最小差，以及点 $x_1,x_2$ 关于 $V$ 上球面度量的最小距离。
\item 对子流形对 $(L_0,L_1)$ 和 $(Q,L_1)$，Floer 数据的 Hamilton 项应为零；类似地，三元组 $(Q,L_0,L_1)$ 的扰动数据应满足 $K=0$。
\item 最后，定义 $(d,l)$ 时使用的辅助几何数据，即 Lefschetz 纤维化、Lagrangian 边界条件及相对扰动数据，必须按一种精确受控的方式选取。
\end{enumerate}
首先观察到，只要允许作下列改变，就能满足所有这些条件：对各 $L_k$ 作 Lagrangian 同痕；将 $\theta_M$ 换为 $\theta_M+dh$，其中函数 $h$ 支撑于 $\operatorname{int}(M)$；以及改变用来定义 Fukaya 范畴的 Floer 数据和扰动数据。这些改变都不影响上同调层次上对象的同构类，无论是在 Fukaya 范畴中，还是在其 $\Ai$ 模范畴中。因此，应用\stref{引理}{5.3}(iii)，已经可以得出 $T_{Y_0}(\Y_1)\cong\Y_2$。对大多数用途而言，这已经足够，但仍不完全令人满意。因此，我们将利用\secref{17f}的结果作更细致的分析。首先，(v) 显然并非必要，因为正如\eref{17.17}所示，$D$ 的同构类与这些选择无关。
\SourcePage{265}
类似地，\eref{17.18}说明它与 $L_1$ 无关，因而可以去掉 (iii)；由于 (ii) 中的管状邻域只用于定义 $L_1=\tau_V^{-1}(L_2)$，这个条件也变得多余。此外，换成与之同痕的 $\widetilde L_1$，也就使得与该子流形相关的 Floer 数据可以完全自由选择，而此前这些数据受 (iv) 限制。对 $L_1$ 所用的论证同样表明，可以在 $L_2$ 的同痕类中移动它，而不影响 $D$ 的拟同构类。最后，正如\stref{引理}{5.3}(iii)中已经提到的，$(d,l)$ 通过\eref{17.19}与 $(c,k)$ 联系这一事实，意味着\eref{17.16}对于某个 $Q$ 无上同调，就对所有与之同痕的 $\widetilde Q$ 也无上同调。合起来，这就使 (i) 以及 (iv) 的剩余部分都不再必要。因此，我们已经证明：\emph{对于任意 $(L_0,L_1,L_2,Q)$，无论定义 $(d,l)$ 时如何作出选择，\eref{17.16}中的复形 $D$ 都无上同调。}用\stref{引理}{5.3}的语言来说，这意味着由 $(c,k)$ 定义的态射 $T_{Y_0}(\Y_1)\to\Y_2$ 始终为同构。

\StatementNumber{13}\begin{note}\label{st:17.13}
作为这一讨论的补充，我们粗略说明怎样构造满足\stref{引理}{5.4}各条件的附加代数数据 $(g,p,q)$。在上同调层次上，$[g]$ 是延拓同构
\[
HF^{\mathrm{pr}}(L_0,L_2)\longrightarrow HF^{\mathrm{pr}}(L_0,\tau_V(L_1))\cong HF^{\mathrm{pr}}(L_0,L_1)
\]
的 Poincaré 对偶。可以用一个技巧简化任务：适当选择 $(L_2,L_0)$ 的 Floer 数据，使其上链群自然地对偶于 $(L_0,L_1)$ 的上链群。若 $L_2$ 就等于 $\tau_V(L_1)$，实现方法是显然的：取 $(L_0,L_1)$ 的 Floer 数据 $(H,J)$，将其换为 $(-H(-t),J(-t))$，再对其施以 $\tau_V$。一般而言，$L_2$ 只与 $\tau_V(L_1)$ 同痕，但可以将这一同痕纳入上述 Floer 数据的修正版中。现在，取 $g$ 为恒等映射 $CF^{\mathrm{pr}}(L_2,L_0)^\vee\to CF^{\mathrm{pr}}(L_0,L_1)$ 的对偶。由此，其余各项及其所需性质可以写成
\begin{equation}\label{eq:17.21}
\boxed{\begin{aligned}
p &:CF^{\mathrm{pr}}(L_2,L_0)\longrightarrow CF^{\mathrm{pr}}(L_1,L_0),\\
q &:CF^{\mathrm{pr}}(L_2,L_0)\otimes CF^{\mathrm{pr}}(L_0,L_1)\longrightarrow CF^{\mathrm{pr}}(L_0,L_0),\\
0&=\mu^1(p(\cdot))+p(\mu^1(\cdot))+\mu^2(\cdot,c),\\
0&=\mu^1(q(\cdot,\cdot))+q(\mu^1(\cdot),\cdot)+q(\cdot,\mu^1(\cdot))\\
 &\quad+\mu^3(\cdot,c,\cdot)+\mu^2(\cdot,k(\cdot))+\mu^2(p(\cdot),\cdot)\\
 &\quad+e\langle\cdot,\cdot\rangle.
\end{aligned}}
\end{equation}
在最后一行中，$\langle\cdot,\cdot\rangle$ 为对偶配对，而 $e\in CF^{\mathrm{pr}}(L_0,L_0)$ 是上同调单位的一个代表。从这些等式的第一个关系式，就能清楚看出如何调整前面 $k$ 的定义来构造 $p$。对于 $q$，考虑\sr{fig:17.9}{图~17.9}中那样的、以 $\RG^{3,1}$ 为底空间的双参数族，不过此时令 $Q=L_0$。边 (i)、(ii) 和 (vi) 的贡献给出\eref{17.21}倒数第二行中的各项，而 (iii) 与 (v) 的贡献消失。于是，
\SourcePage{266}
最后剩下且最有意思的是 (iv)。有关几何设置由两个 Lefschetz 纤维化组成，并将它们在边界标点上方的纤维作一等同，见\sr{fig:17.10}{图~17.10}。左边那一块记为 $E^e\to S^e$，其标点位置可变，因此得到一个带附加参数的求值映射。假设 $(L_0,L_0)$ 的 Floer 数据满足：Hamilton 函数为 $L_0=S^n$ 上的标准 Morse 函数，从而 $CF^{\mathrm{pr}}(L_0,L_0)$ 有两个生成元 $e=y_0$ 与 $y_n$。那么，模空间 $\M_{E^e/S^e}(y_0)$ 和 $\M_{E^e/S^e}(y_n)$ 的维数分别为 $n-1$ 与 $2n-1$。取其中第一个空间，考虑带参数的求值映射；这里用平行移动将 $\partial S^e$ 任意点上的 Lagrangian 边界条件与 $L_0$ 等同：
\[
\operatorname{ev}:\R\times\M_{E^e/S^e}(y_0)\longrightarrow L_0.
\]
可以证明，$\operatorname{ev}$ 在 $L_0\setminus\{y_0,\tau_V(y_0)\}$ 上的度数为奇数；缺少的这两点对应于 $z$ 趋于无穷时的极限。证明方法是先归约到闭圆盘的情形，再归约到标准局部模型；参见\bcite{131}笔记 2.17(iii)。一旦得到这一点，就容易看出 (iv) 的贡献确实是 $e\langle\cdot,\cdot\rangle$。
\SourceFigure{figure-17-10}{0.458999506}{17.10}
\end{note}

\subsection{相对膜结构}\label{sec:17i}
设 $\pi:E\to S$ 为 Lefschetz 纤维化；我们先讨论 $S$ 紧的情形，稍后再处理带条带状端的版本。假设 $E$ 带有相对二次复体积形式 $\eta^2_{E/S}$，以 $\alpha_{E/S}$ 表示相对平方相位映射。再假设 $F$ 为 $E$ 上的 Lagrangian 边界条件。
\begin{definition*}
$F$ 的\emph{相对分次}是映射 $x\mapsto\alpha_{E/S}(TF_x^v)$ 的一个提升 $\alpha^\#:F\to\R$。$F$ 上的\emph{相对膜结构}是二元组 $(\alpha^\#,P^\#)$，由相对分次与相对 Pin 结构组成；后者就是向量丛
\SourcePage{267}
$TF^v\to F$ 上的一个 Pin 结构 $P^\#$。
\end{definition*}
这类结构存在的障碍可作如下理解：在每个边界圆周 $C$ 上选一点 $z$。每个 $F_z\subset E_z$ 都必须能关于 $\eta^2_{E_z}$ 赋予分次；此外，若 $h_C$ 为沿该边界的单值变换，并如\secref{15c}中那样带有自然分次，则其像 $h_C(F_z)$ 必须与 $F_z$ 作为分次 Lagrangian 子流形相等。类似地，每个 $F_z$ 必须有 Pin 结构 $P_z^\#$，使 $(h_C|F_z)^*P_z^\#$ 同构于 $P_z^\#$；选定一个这样的同构，就确定了 $F|C$ 的相对 Pin 结构。

\StatementNumber{14}\begin{example}\label{st:17.14}
取\stref{例题}{17.3}中的局部模型，为简单起见令 $r=1$。标准相对二次复体积形式是 $\eta_E=dx_1\wedge\cdots\wedge dx_{n+1}$ 与 $\eta_S=dz$ 的商的平方。对于每个 $x\in F_z$，有
\[
\alpha_{E/S}(TF_x^v)=\frac{\eta_E(v_1,\ldots,v_{n+1})^2}{\eta_S(iz\partial_z)^2},
\]
其中 $\{v_1,\ldots,v_{n+1}\}$ 为整个 $F$ 的切空间的一组标准正交基。具体地，可选 $v_1,\ldots,v_n$ 为 $\sqrt zS^n\subset\sqrt z\R^{n+1}$ 的切向量，并取 $v_{n+1}=ix\in i\sqrt z\R^{n+1}$，从而 $\alpha_{E/S}(TF_x^v)=z^{n-1}$。因此，$n=1$ 时 $F$ 存在相对分次，而 $n>1$ 时不存在。\znote{259}{$\alpha_{E/S}$ 取值于 $S^1$，相对分次是它的实值提升，绕临界值一圈的卷绕数就是障碍。此处 $\alpha_{E/S}(TF_x^v)=z^{n-1}$ 绕 $z=0$ 卷绕 $n-1$ 次，故 $n>1$ 时无整体提升；$n=1$ 时 $\alpha_{E/S}\equiv1$，常值提升即为分次。又因 $\alpha$ 是平方相位，其卷绕数是相对第一 Chern 数的两倍，这正是\stref{笔记}{17.12}中 $2c_1(M,V)=0$ 的来源。[来源：Seidel, Fukaya categories and Picard–Lefschetz theory, §17i。]}这反映了如下事实：任意 Dehn 扭转 $\tau_V$ 赋予其自然分次后，都会将中心球面 $V$ 的分次平移 $1-n$，见\bcite{125}引理 5.7。相对 Pin 结构则始终存在；事实上，$n=1$ 时有四种不同选择，$n>1$ 时有两种。
\end{example}

分次的存在保证：对任意边界值落在 $F$ 中的截面 $u$，丛对 $(u^*(TE^v),u^*(TF^v))$ 的相对第一 Chern 类为零。因此，假定正则性成立，模空间的维数就是 $n\chi(S)$。与前面一样，在 $\partial S$ 的每个分支上选一点，记为 $z_j$。给定相对 Pin 结构后，就得到 $\M_{E/S}$ 相对于各 $z_j$ 处求值映射之积的一个优选定向。这是\stref{命题}{11.13}的推论，只须将抽象膜结构的语言转换成当前的语言。于是，这些带优选相对定向的求值映射，通过 Poincaré 对偶表示一个上同调类
\begin{equation}\label{eq:17.22}
\Phi_{E/S,z_1,\ldots,z_k}\in H^{n(k-\chi(S))}\Bigl(\prod_jF_{z_j}\Bigr).
\end{equation}
取任意实线丛 $\beta\to\partial S$。可以通过 $P^\#\mapsto P^\#\otimes(\pi|F)^*\beta$ 改变相对 Pin 结构；这对\eref{17.22}的影响是将不变量乘以 $(-1)^i$，其中 $i$ 是 $\beta$ 非平凡的分支数。最后，这一整套讨论都可毫无困难地推广到移动边界条件。

对我们而言，\eref{17.22}最重要的特殊情形是 $S$ 为圆盘、纤维维数为 $2n=2$。后一个假设意味着，线性化算子 $D_{E/S,u}$ 是线丛上的 $\bar\partial$ 算子，因此\secref{11d}中的论证
\SourcePage{268}
适用。特别地，由\eref{11.13}可知，对每个 $u\in\M_{E/S}$，$\ker(D_{E/S,u})$ 都是一维的，并由 $u^*(TE^v)$ 的一个处处非零截面 $X$ 生成。在 $X$ 给出的平凡化下，$D_{E/S,u}$ 就是平凡线丛上具有常值边界条件的标准 $\bar\partial$ 算子。回顾\stref{引理}{11.12}的证明可见，若 $u^*(TF^v)$ 上的 Pin 结构平凡，则相对定向 $\ker(D_{E/S,u})\cong(TF_z)_{u(z)}$ 由映射 $X\mapsto X(z)$ 给出；否则，定向与之相反。这意味着，可按如下方式计算局部常值函数 $\Phi_{E/S,z}\in H^0(F_z)$：选一点 $x\in F_z$，然后计数满足 $u(z)=x$ 的截面 $u\in\M_{E/S}$。若 $u^*(TF^v)$ 上的 Pin 结构平凡，则赋予重数 $+1$；否则赋予 $-1$。

\StatementNumber{15}\begin{example}\label{st:17.15}
回到\stref{例题}{17.3}，假设 $n=1$。边界条件可显式参数化为
\[
S^1\times S^1\longrightarrow F,\qquad
(t_1,t_2)\longmapsto\bigl(e^{it_1/2}\cos(t_1/2+t_2),e^{it_1/2}\sin(t_1/2+t_2)\bigr).
\]
截面模空间是两个圆周 $\M_{E/S,\pm}$ 的不交并，它们由\eref{17.8}中取 $a=\tfrac12(e^{it},\pm ie^{it})$ 得到。在这些坐标下，沿边界的求值映射成为
\[
\begin{aligned}
\partial S\times\M_{E/S,+}&\longrightarrow F,&(s,t)&\longmapsto(t_1,t_2)=(s,t-s),\\
\partial S\times\M_{E/S,-}&\longrightarrow F,&(s,t)&\longmapsto(t_1,t_2)=(s,t).
\end{aligned}
\]
若固定 $F$ 中一点，比如 $(t_1,t_2)=(0,0)$，则每族中都恰有一个截面 $u_\pm$ 经过该点，即 $t=0$ 所对应的截面。注意，在同调中，圆周 $u_\pm(\partial S)\subset F$ 相差一份纤维 $F_z$。根据前面的一般性观察，可得如下结论：若选择的相对 Pin 结构在 $F_z$ 上\emph{非平凡}，则 $u_+$ 与 $u_-$ 的贡献相互抵消，不变量为零。否则，$\Phi_{E/S,z}=\pm2$；两个符号都会出现，对应的两个 Pin 结构相差 $\partial S$ 上非平凡线丛的拉回。\znote{260}{$u_+(\partial S)$ 与 $u_-(\partial S)$ 在 $F$ 中只相差一份纤维类 $F_z$，故相对 Pin 结构限制到 $u_\pm^*(TF^v)$ 上时恰相差一次符号；结合前面\stref{引理}{11.12}的观察（Pin 结构平凡给 $+1$、非平凡给 $-1$），两条截面的贡献相消。只有当 Pin 结构在 $F_z$ 上平凡时二者同号，总数才得 $\pm2$。[来源：Seidel, Fukaya categories and Picard–Lefschetz theory, §17i 与\stref{引理}{11.12}。]}
\end{example}

固定带角正合辛流形 $M$ 及其上的二次复体积形式 $\eta_M^2$。取一个带条带状端的 Lefschetz 纤维化，使每个端上的平凡化都以 $Z^\pm\times M$ 为模型。为它赋予一个可能移动的 Lagrangian 边界条件 $F$，以 $(L_{\zeta,0},L_{\zeta,1})$ 表示与每个端 $\zeta$ 相关的 Lagrangian 子流形对。现在假设，它们带有膜结构 $(L_{\zeta,0}^\#,L_{\zeta,1}^\#)$，且 $E$ 带有相对二次复体积形式 $\eta^2_{E/S}$，在每个端上都与 $\eta_M^2$ 一致。
\begin{definition*}
此时，$F$ 上的\emph{相容相对膜结构}由下列数据组成：相对分次 $\alpha^\#$，在每个端上等于各 $L_{\zeta,k}$ 的分次；相对 Pin 结构 $P^\#$；以及该结构与各 $L_{\zeta,k}$ 在端上所带的 Pin 结构之间的指定同构。
\end{definition*}
有了这种结构，便可对行列式线丛作与\secref{12c}类似的分析。当每个边界圆周恰含一个入向无穷远点时，结论最容易陈述；我们所有的应用都属于这种情形。此时，对
\SourcePage{269}
$\M_{E/S}(\{y_\zeta\})$ 中孤立点作带符号计数，就得到\eref{17.10}的精细版本，即上链映射
\[
C\Phi_{E/S}:\bigotimes_{\zeta_+\in\Sigma^+}CF^*(L_{\zeta_+,0}^\#,L_{\zeta_+,1}^\#)
\longrightarrow\bigotimes_{\zeta_-\in\Sigma^-}CF^*(L_{\zeta_-,0}^\#,L_{\zeta_-,1}^\#)[n(\chi(S)-|\Sigma^-|)],
\]
其中 $CF^*$ 是以任意域 $\K$ 为系数的 $\Z$ 分次 Floer 上链复形。还可以允许更一般的曲面，但此时要么必须为某些 $F_z$ 选定定向，见\eref{12.5}；要么必须像\eref{12.7}中那样，改变与端相关的某些膜结构。

\subsection{代数扭转与几何扭转}\label{sec:17j}
假设有一个以圆盘为底空间、纤维为 $M$、恰有一个临界点的 Lefschetz 纤维化。若纤维维数 $2n>2$，这样的纤维化总存在从 $M$ 的二次复体积形式继承而来的相对二次复体积形式；若 $2n=2$，则当且仅当消失圈 $V$ 能赋予分次时才存在。考虑到这一点，\secref{17d}至\secref{17h}中的全部论证都可照常进行，所得符号也都与\stref{引理}{5.3}中预期的符号一致；这一验证直接但略显繁琐，这里略去。如\stref{例题}{17.14}所示，当 $2n>2$ 时，标准边界条件不能赋予相对分次，但这不成问题，因为在这些维数中，\stref{笔记}{17.8}的消失论证不依赖任何抵消。在剩下的情形中，\stref{例题}{17.14}适用，因此若要在 $\operatorname{char}(\K)\ne2$ 的情形下工作，就必须在消失圈 $V=L_0$ 上选择非平凡 Pin 结构。\secref{17h}中的主要结果，在补上正确的分次和符号后，如下所述。

\StatementNumber{16}\begin{theorem}\label{st:17.16}
设 $V=L_0$ 为 $M$ 中带标架的正合 Lagrangian 球面，并带有膜结构 $L_0^\#$；当 $n=1$ 且 $\operatorname{char}(\K)\ne2$ 时，该结构中的 Pin 结构应非平凡。设 $L_1,L_2$ 为另外两个正合 Lagrangian 子流形，固定一个同痕 $L_2\simeq\tau_V(L_1)$，并假设它们也带有与该同痕相容的膜结构 $L_k^\#$，其中为 $\tau_V$ 赋予自然分次。以 $Y_k$ 表示 $L_k^\#$ 给出的 $\F(M)$ 的对象，以 $\Y_k$ 表示其 Yoneda $\Ai$ 模。那么，通过计数纤维化 $E^c\to S^c$ 与 $E^{k,r}\to S^{k,r}$ 的截面，可得到
\[
c\in CF^0(L_1^\#,L_2^\#),\qquad
k:CF^*(L_0^\#,L_1^\#)\longrightarrow CF^*(L_0^\#,L_2^\#)[-1],
\]
使按\stref{引理}{5.3}定义的相应 $[t]:T_{Y_0}(\Y_1)\to\Y_2$ 在 $H^0(\modcat(\F(M)))$ 中为同构。
\end{theorem}

为得到更具体的陈述，现在像\eref{5.6}中那样，以扭曲复形代替 $\Ai$ 模。再假设 $L_2=\tau_V(L_1)$，并略去一些细节，\stref{定理}{17.16}就成为下列“$T=\tau$”陈述：
\StatementNumber{17}\begin{corollary}\label{st:17.17}
在导出范畴 $D\F(M)$ 中，有 $T_{L_0^\#}(L_1^\#)\cong\tau_{L_0}(L_1^\#)$。
\end{corollary}
\SourcePage{270}
$T_{L_0^\#}(L_1^\#)$ 属于正合三角\eref{5.4}，由正合三角的一般性质，可将其换成 $\tau_{L_0}(L_1^\#)$。注意，三角中的一个态射为\secref{5b}所定义的纯代数求值态射 $[f]$。由\stref{引理}{5.3}，第二个态射为 $[c]$，因此有：
\clearpage
\enlargethispage{\baselineskip}
\StatementNumber{18}\begin{corollary}\label{st:17.18}
在 $D\F(M)$ 中存在正合三角
\SourceEquationDiagram{diagram-17-23}{0.413307247}{17.23}
\end{corollary}

\StatementNumber{19}\begin{note}\label{st:17.19}
\eref{17.23}中的第三个映射是下列空间中的元素：
\[
\begin{aligned}
&\Hom_{D\F(M)}\bigl(\tau_{L_0}(L_1^\#),HF^*(L_0^\#,L_1^\#)\otimes L_0^\#[1]\bigr)\\
&\quad\cong\bigl(HF^*(L_0^\#,L_1^\#)\otimes HF^*(\tau_{L_0}(L_1^\#),L_0^\#)\bigr)^1\\
&\quad\cong\bigl(HF^*(L_0^\#,L_1^\#)\otimes HF^*(L_1^\#,L_0^\#)\bigr)^n.
\end{aligned}
\]
最后一个同构使用了\stref{例题}{17.14}中提到的等式 $\tau_{L_0}(L_0^\#)=L_0^\#[1-n]$。\znote{262}{两个指数是同一笔平移账：目标中的 $[1]$ 使整个 $\Hom$ 落在 $\bigl(HF^*\otimes HF^*\bigr)^1$；又由 $\tau_{L_0}(L_0^\#)=L_0^\#[1-n]$，把 $\tau_{L_0}$ 同时作用到两个变元上得 $HF^*(\tau_{L_0}(L_1^\#),L_0^\#)\cong HF^*(L_1^\#,L_0^\#)[n-1]$，次数整体升 $n-1$，指数便由 $1$ 变成 $n$。这类次数账在正合三角里最容易错：只要把 $\tau_{L_0}(L_1^\#)$ 换成 $L_1^\#$ 而忘记 $L_0^\#$ 同时平移，指数就差 $n-1$。[来源：Seidel, Fukaya categories and Picard–Lefschetz theory, §17j 与例 17.14。]}自然会预期，这就是由一对 Poincaré 对偶基给出的元素；等价地，可以将其视为\eref{5.20}中所用的那种对偶求值映射。原则上，将\stref{笔记}{5.5}与\stref{笔记}{17.13}结合起来，似乎就能证明情况确实如此。不过，我们尚未核查全部细节；尤其是符号问题颇为棘手，会以非平凡的方式影响结果，参见\stref{笔记}{3.12}。
\end{note}

最后，给出一个更有分量的应用。设 $V^\#=(V_1^\#,\ldots,V_m^\#)$ 为 $M$ 中一组带标架的正合 Lagrangian 球面，每个都带有膜结构；当 $n=1$ 且 $\operatorname{char}(\K)\ne2$ 时，要求所有 Pin 结构非平凡。设 $W^\#=\LG_kV^\#$ 为对它作一次初等 Hurwitz 移动\eref{16.7}得到的集合，并赋予自然诱导的膜结构；其中的分次使用 $\tau_{V_k}$ 的自然分次。设 $\F^\to(V^\#)$ 与 $\F^\to(W^\#)$ 为 $\F(M)$ 中相应的有向 $\Ai$ 子范畴。\stref{推论}{17.17}表明，$W^\#$ 的各对象在 $D^b(\F(M))$ 中，与对 $V^\#$ 施以变异\eref{5.24}所得集合的各对象同构。由\stref{引理}{5.21}，换成同构对象不会改变有向 $\Ai$ 范畴的拟同构类。因此，\stref{引理}{5.23}适用，并说明 $\Tw\F^\to(V^\#)$ 与 $\Tw\F^\to(W^\#)$ 拟等价。另一个更初等得多的事实是：平移一个或多个 $V_j$ 的分次，也不会改变 $\Tw\F^\to(V^\#)$。将结论归纳如下：
\SourcePage{271}
\StatementNumber{20}\begin{theorem}\label{st:17.20}
在拟等价意义下，$\Tw\F^\to(V^\#)$ 是 $V$ 的 Hurwitz 等价类的不变量。
\end{theorem}
将其应用于消失圈的特选基，立即得到下述结论：
\StatementNumber{21}\begin{corollary}\label{st:17.21}
设 $\pi:E\to S$ 为底空间 $S$ 是圆盘的 Lefschetz 纤维化，并带有相对二次复体积形式 $\eta^2_{E/S}$。选取基点 $*$，为 $M=E_*$ 赋予相应的形式 $\eta_M^2$。设 $\symbfit{\gamma}$ 为从 $*$ 出发的消失路径特选基，以 $V_{\symbfit{\gamma}}$ 表示 $M$ 中相应的消失圈集合。为每个消失圈赋予分次和 Pin 结构；若 $n=1$ 且 $\operatorname{char}(\K)\ne2$，则选择非平凡 Pin 结构。将所得集合记为 $V_{\symbfit{\gamma}}^\#$，并设 $\F^\to(V_{\symbfit{\gamma}}^\#)$ 为 $\F(M)$ 中相应的有向 $\Ai$ 子范畴。那么，在拟等价意义下，扭曲复形范畴 $\Tw\F^\to(V_{\symbfit{\gamma}}^\#)$ 与 $\symbfit{\gamma}$ 及分次的选择无关，因此是该 Lefschetz 纤维化的不变量。故导出范畴 $D\F^\to(V_{\symbfit{\gamma}}^\#)$ 也具有同样的性质，只须将拟等价换成三角范畴的等价。
\end{corollary}
\StatementNumber{22}\begin{note}\label{st:17.22}
\stref{推论}{17.17}最初由 Kontsevich 猜想提出。显然，人们会期待一个更强的版本成立，即 $\tau$ 诱导的自同构，应在 $H^0(\operatorname{fun}(\F(M),\F(M)))$ 中典范地同构于 $T$。这里不证明这一点；不过，预期它将是 Wehrheim–Woodward 正在进行的 Lagrangian 对应研究的一个推论。
\end{note}

\clearpage
\subsection{等变情形}\label{sec:17k}
先作一些简单的几何准备。设 $M$ 为带 $\Z/2$ 作用 $\iota_M$ 的正合辛流形。
\begin{definition*}
$M$ 中的\emph{不变带标架正合 Lagrangian 球面}，是一个 $\iota_M$-不变的正合 Lagrangian 子流形 $V$，加上等变微分同胚 $v:S\to V$；这里 $S$ 是某个带线性 $\Z/2$ 作用的内积空间中的单位球面。
\end{definition*}
由于 Lagrangian 管状邻域定理可以与有限群作用相容，而 Dehn 扭转的局部模型关于 $O_{n+1}$ 等变，故沿这样的球面的 Dehn 扭转 $\tau_V$ 自然地成为一个等变正合辛自同构。给定带逐纤维对合的正合 Lefschetz 纤维化，其相应平行移动映射等变，消失圈为等变带标架球面，并且 Picard–Lefschetz 公式\eref{16.6}的适当等变版本成立。反过来，给定 $M$ 中一个或多个等变带标架球面，可以使\stref{引理}{16.8}和\stref{引理}{16.9}中的 Lefschetz 纤维化带有逐纤维对合。

现在假设处于\secref{14e}的情形，从而可定义等变 Fukaya 范畴。将其记为 $\A=\F(M)^{\mathrm{naive}}$，将其上的 $\Z/2$ 作用记为 $\eps_{\A}$。假设有一个等变带标架正合 Lagrangian
\SourcePage{272}
球面 $V=L_0$，并可将其赋予等变膜结构 $L_0^\#=(L_0,\alpha_0^\#,P_0^\#,f_0)$。还要求有另外两个等变膜 $L_k^\#=(L_k,\alpha_k^\#,P_k^\#,f_k)$，$k=1,2$，以及一个通过不变正合 Lagrangian 子流形的同痕 $\tau_V(L_1)\simeq L_2$，并可使之与等变膜结构相容。此时，可以考虑由 $L_k^\#$ 给出的 $\eps_{\A}$-不动对象 $Y_k\in\Ob\A$，以及相应的等变 Yoneda 模 $\Y_k\in\Ob\modcat^{\mathrm{eq}}(\A)$。正如\secref{6c}所指出的，扭曲 $\Ai$ 模 $T_{Y_0}(\Y_1)$ 仍然等变。\stref{定理}{17.16}的等变版本说明，$(c,k)$ 给出等变模态射
\begin{equation}\label{eq:17.24}
[t^{\mathrm{eq}}]\in\Hom_{H^0(\modcat^{\mathrm{eq}}(\A))}(T_{Y_0}(\Y_1),\Y_2),
\end{equation}
并且它在该范畴中为同构。虽然这里没有新的内容，稍作解释仍有帮助。首先，$t^{\mathrm{eq}}$ 由与\stref{引理}{5.3}(i)相同的公式定义。再回顾，若忘掉 $\Z/2$ 作用，$\A$ 就与 $\F(M)$ 拟等价。更精确地说，将每个 Lagrangian 膜赋予预等变结构，就可将 $\F(M)$ 等同于 $\A$ 的一个全 $\Ai$ 子范畴，该子范畴与整个 $\A$ 拟等价；赋予这种结构总是可能的，见\stref{笔记}{14.1}。由此得到的函子
\begin{equation}\label{eq:17.25}
\modcat^{\mathrm{eq}}(\A)\longrightarrow\modcat(\A)\longrightarrow\modcat(\F(M))
\end{equation}
在上同调意义下全忠实；第一部分忘掉等变性，在链层次上全忠实，而第二部分是向 $\F(M)\subset\A$ 的限制，由\stref{引理}{2.8}为拟等价。此外，可以适当安排嵌入 $\F(M)\subset\F(M)^{\mathrm{naive}}$，使 Lagrangian 膜 $L_k^\#$ 映到对应的等变膜。此时，$t^{\mathrm{eq}}$ 在\eref{17.25}下的像，就是前面考虑过的模同态 $t$，因此如上所述，$[t^{\mathrm{eq}}]$ 是同构。

\StatementNumber{23}\begin{lemma}\label{st:17.23}
$[t^{\mathrm{eq}}]$ 属于\eref{17.24}中 $\Hom$ 群的 $\Z/2$-不变部分。
\end{lemma}
\znote{264}{“naive”的代价正在这里：$\eps_{\A}$ 只在链同伦意义下作用，$c$ 与 $\eps_{\A}(c)$ 并不相等，仅相差一个上边界。证明正是具体给出 $(\gamma,\kappa)$，使 $\mu^1_{\A}(\gamma)=c-\eps_{\A}(c)$，并让 $k$ 与 $\eps_{\A}(k)$ 相差一个同伦，于是 $t^{\mathrm{eq}}$ 与其 $\Z/2$-像同伦。故“不变”只能在上同调层次断言，不能写成链层次的等式。[来源：Seidel, Fukaya categories and Picard–Lefschetz theory, §14e 与 §17k。]}
\begin{proof}
完全有理由认为这显然成立：此前已证明 $[t]$ 是典范的，因此它应在几何设置固有的任何对称下保持不变，等变版本也应如此。我们仍要查看细节，确保其中没有意外。可以使 Lefschetz 纤维化 $E^c\to S^c$ 与 $E^{k,r}\to S^{k,r}$，以及它们的边界条件，都带有与 $M$ 上作用相容的逐纤维 $\Z/2$ 作用。所有涉及的 Lagrangian 子流形都是等变膜，按 $\F(M)^{\mathrm{naive}}$ 的构造，它们带有 $\Z/2$-不变 Floer 数据。剩下的选择是相对扰动数据。此处，我们希望避开等变横截性问题，改作如下论证。先更一般地假设 $L_1^\#,L_2^\#$ 为预等变膜。首先考虑带有逐纤维
\SourcePage{273}
对合的 $E^c\to S^c$。以 $F$ 表示用来定义 $c$ 的 Lagrangian 边界条件，以 $F^\#$ 表示其相对膜结构，以 $(K,J)$ 表示相应的相对扰动数据。可将逐纤维对合同时作用于 $(K,J)$ 与 $F^\#$，再将后者的 Pin 结构与线丛 $\lambda^{\mathrm{top}}(TF^v)$ 作张量积。在这个改变后的情形中计数非齐次伪全纯截面，所得为
\[
CF^0\bigl(\iota_M(L_1^\#\otimes\lambda^{\mathrm{top}}(TL_1)),\iota_M(L_2^\#\otimes\lambda^{\mathrm{top}}(TL_2))\bigr)
\]
中的一个元素，按定义它就是 $\eps_{\A}(c)$。现在，若 $L_1^\#$ 与 $L_2^\#$ 均为等变膜，同构 $f_k$ 就给出自然等同
\[
CF^0\bigl(\iota_M(L_1^\#\otimes\lambda^{\mathrm{top}}(TL_1)),\iota_M(L_2^\#\otimes\lambda^{\mathrm{top}}(TL_2))\bigr)=CF^0(L_1^\#,L_2^\#),
\]
它来自 Pin 结构的同构 $P_k^\#\to\iota_M^*(P_k^\#\otimes\lambda^{\mathrm{top}}(TL_k))$。在当前情形中，这些同构与同痕 $\tau_V(L_1)\simeq L_2$ 相容，也就是说，它们给出 Lagrangian 边界条件 $F$ 的相对 Pin 结构之间的一个同构。由此可清楚看出，$\eps_{\A}(c)$ 与 $c$ 的区别仅在于相对扰动数据的选择。$k$ 的情形完全类似。因此，可以应用\eref{17.14}中的构造，插入适当符号，得到 $(\gamma,\kappa)$，满足
\[
\begin{aligned}
&\mu^1_{\A}(\gamma)+c-\eps_{\A}(c)=0,\\
&(-1)^{|c_0|-1}\bigl(\mu^1_{\A}(\kappa(c_0))-\kappa(\mu^1_{\A}(c_0))+\mu^2_{\A}(\gamma,c_0)\bigr)\\
&\qquad-\eps_{\A}(\kappa(\eps_{\A}(c_0)))+\kappa(c_0)=0
\qquad\text{对所有 }c_0\in CF^*(L_0^\#,L_1^\#).
\end{aligned}
\]
有了这些数据，\stref{引理}{5.3}(ii)的论证就给出 $t^{\mathrm{eq}}$ 与它在 $\Z/2$ 作用下的像之间的同伦，从而证明，在上同调层次上，$[t^{\mathrm{eq}}]$ 确实不变。
\end{proof}

\section{Lefschetz 纤维化的 Fukaya 范畴}\label{sec:18}\label{section:18}
\subsection{二重分歧覆盖}\label{sec:18a}
固定一个正合 Lefschetz 纤维化 $\pi:E^{2n+2}\to S$，其中 $S=D$ 为 $\CC$ 中带标准复结构的单位圆盘。选一个很小的 $\delta>0$，以 $S^{\mathrm{shrink}}\subset S$ 表示半径 $1-4\delta$ 的圆盘，以 $E^{\mathrm{shrink}}\subset E$ 表示它在总空间中的原像。假设 $\pi$ 的全部临界值都在 $S^{\mathrm{shrink}}$ 的内部；还假设，$E$ 上的负 Liouville 流经过某个正时间后，将整个 $E$ 推入 $E^{\mathrm{shrink}}$。当 $\delta$ 充分小时，这两个条件必定满足。取 $*=1-3\delta\in\operatorname{int}(S)$ 为基点，$M=E_*$ 为该点上的纤维。需要选择基路径 $\chi$ 时，始终使用从 $*$ 到 $1\in\partial S$ 的直线段。\znote{266}{四个半径的层级容易混淆，先列清楚：$S$ 是半径 $1$ 的圆盘；$\mathcal W=\{w:|w|\le1-2\delta\}$ 是 §18b 中族参数的取值圆盘；基点 $*=1-3\delta$ 落在 $\mathcal W$ 内但在 $S^{\mathrm{shrink}}$ 之外；而 $S^{\mathrm{shrink}}$ 与 $E^{\mathrm{shrink}}$ 的门槛是 $1-4\delta$。于是含 $*$ 的圆周 $|w|=1-3\delta$ 整条落在 $\mathcal W$ 内，引理 18.1 中的单值变换才是族 $\Psi$ 的合法环路；同时 $E^{\mathrm{shrink}}$ 上的一切估计都留出了 $\delta$ 的余量。后文还会出现第五个半径 $1-\delta$，用于定义 $\widetilde U$，不要与前四个混同。[来源：原书 §18a、§18b（印刷页 266、269）]}

设 $\widetilde S$ 为在 $*$ 处分歧的 $S$ 的二重覆盖；类似地，设 $\widetilde E$ 为沿 $M=E_*$ 分歧的 $E$ 的二重覆盖。当然，$\widetilde S$ 仍与 $D$ 同构，但我们
\SourcePage{274}
不作这一等同，而写成
\begin{equation}\label{eq:18.1}
\begin{aligned}
\widetilde S&=\{y\in\CC:y^2+*\in S\},\\
\widetilde E&=\{(y,x)\in\widetilde S\times E:y^2+*=\pi(x)\}.
\end{aligned}
\end{equation}
$\widetilde E$ 带有诱导的殆复结构，因为它是乘积 $\widetilde S\times E$ 的殆复子流形；这一结构使 $\partial\widetilde E$ 弱凸。关于该结构，有下列伪全纯映射：
\[
\begin{aligned}
\iota_{\widetilde S}(y)&=-y,&\iota_{\widetilde E}(y,x)&=(-y,x)&&\text{覆盖对合；}\\
\beta_{\widetilde S}(y)&=y^2+*:\widetilde S\to S,&\beta_{\widetilde E}(y,x)&=x:\widetilde E\to E&&\text{商映射；}\\
\widetilde\pi(y,x)&=y:\widetilde E\to\widetilde S&&&&\text{$\pi$ 的提升。}
\end{aligned}
\]
$\widetilde\pi$ 的临界点是 $\pi$ 的各临界点的原像；由于局部结构不变，它们仍然非退化。

为完成构造，还要引入适当的正合辛结构。回顾，定义纤维化的数据中，底空间 $S$ 带有某个 $\theta_S\in\Omega^1(S)$ 及 $\omega_S=d\theta_S\in\Omega^2(S)$。从\eref{18.1}可见，$\widetilde S$ 包含以原点为中心、半径 $\sqrt{2\delta}$ 的圆盘。设 $\kappa$ 为支撑在更小的半径 $\sqrt\delta$ 的圆盘内的一形式，关于 $\iota_{\widetilde S}$ 不变，且 $d\kappa$ 在原点处为正。我们将 $\kappa$ 看作很小的形式，因此保留在必要时将最初选取的形式乘以某个很小正常数的余地。令
\begin{equation}\label{eq:18.2}
\theta_{\widetilde S}=\beta_{\widetilde S}^*\theta_S+\kappa,\qquad
\omega_{\widetilde S}=d\theta_{\widetilde S}=\beta_{\widetilde S}^*\omega_S+d\kappa.
\end{equation}
由于 $d\kappa$ 很小，且在原点处为正，$\omega_{\widetilde S}$ 处处为正；又因为 $\kappa$ 支撑于内部，$\theta_{\widetilde S}$ 沿边界为正。这些条件保证\eref{18.2}使 $\widetilde S$ 成为带边界正合辛流形。与此相应，为 $\widetilde E$ 赋予
\begin{equation}\label{eq:18.3}
\theta_{\widetilde E}=\beta_{\widetilde E}^*\theta_E+\widetilde\pi^*\kappa,\qquad
\omega_{\widetilde E}=d\theta_{\widetilde E}=\beta_{\widetilde E}^*\omega_E+\widetilde\pi^*d\kappa.
\end{equation}
第一项 $\beta_{\widetilde E}^*\omega_E$ 只在分歧纤维 $\{y=0\}$ 的法向方向上不满足辛条件。因此，由与前面相同的论证，加入充分小的 $\widetilde\pi^*d\kappa$ 项，就使它处处为辛形式。注意，由于\eref{18.3}的第二项来自底空间，$\widetilde\pi$ 正则部分上的相应辛联络，是 $\pi$ 的辛联络的拉回。这说明 $\partial^h\widetilde E$ 是水平的；同时，使用\stref{引理}{15.1}中的相同论证，可知负 Liouville 流沿这一部分边界指向内部。在 $\partial^v\widetilde E$ 附近，负 Liouville 向量场是
\SourcePage{275}
$E$ 上相应向量场的拉回，因此也指向内部。这样，就验证了 $\widetilde\pi:\widetilde E\to\widetilde S$ 仍为正合 Lefschetz 纤维化。

在以下讨论中，$\widetilde E$ 中两类正合 Lagrangian 子流形将起突出作用。分别将它们记为 (U) 型与 (B) 型，对应于无分歧与有分歧：
\begin{definition*}[(U) 型]
覆盖 $\beta_{\widetilde S}$ 在 $S^{\mathrm{shrink}}$ 上的限制无分歧，因而平凡，也就是说，总空间由底空间的两个副本组成。$\beta_{\widetilde E}$ 在 $E^{\mathrm{shrink}}$ 上的限制也如此。记
\begin{equation}\label{eq:18.4}
\begin{aligned}
\widetilde S^{\mathrm{shrink}}&=\beta_{\widetilde S}^{-1}(S^{\mathrm{shrink}})=\widetilde S^{\mathrm{shrink}}_+\sqcup\widetilde S^{\mathrm{shrink}}_-,\\
\widetilde E^{\mathrm{shrink}}&=\beta_{\widetilde E}^{-1}(E^{\mathrm{shrink}})=\widetilde E^{\mathrm{shrink}}_+\sqcup\widetilde E^{\mathrm{shrink}}_-.
\end{aligned}
\end{equation}
更具体地说，满足 $|y^2+*|\le1-4\delta$ 的点 $y$ 必有 $\operatorname{im}(y)\ne0$，该虚部的符号区分了\eref{18.4}的两个分支。这些 $y$ 还满足 $|y|\ge\sqrt\delta$，因此在各二重分歧覆盖的限制版本上，\eref{18.2}、\eref{18.3}中的 $\kappa$ 项为零。\znote{267}{这两步都是作者的「显然」，补上推导：若 $y$ 为实数，则 $y^2+*\ge1-3\delta>1-4\delta$，与 $|y^2+*|\le1-4\delta$ 矛盾，故 $\operatorname{im}(y)\ne0$，\eref{18.4} 按虚部符号分成两支正由此而来；又 $|y^2|\ge|*|-|y^2+*|\ge(1-3\delta)-(1-4\delta)=\delta$，即 $|y|\ge\sqrt\delta$。因 $\kappa$ 支撑在 $\sqrt\delta$ 圆盘内，它在 $\widetilde S^{\mathrm{shrink}}$ 上恒为零，故\eref{18.2}、\eref{18.3}在该区域就是纯粹的拉回——这正是 (U) 型可写成两个不相交副本、且每支仍是正合 Lagrangian 的原因。[来源：原书 §18a（印刷页 267）]}一个显然的推论是：若 $L$ 为 $E$ 中包含于 $E^{\mathrm{shrink}}$ 的正合 Lagrangian 子流形，则它在 $\widetilde E$ 中的原像由 $L$ 的两个副本组成，且它们也都是正合 Lagrangian 子流形。类似于前面，记
\[
\widetilde L=\beta_{\widetilde E}^{-1}(L)=\widetilde L_+\sqcup\widetilde L_-.
\]
\end{definition*}
\begin{definition*}[(B) 型]
设 $\gamma:[0,1]\to S$ 为满足 $\gamma(0)=*$ 的消失路径，$\Delta_\gamma\subset E$ 为其 Lefschetz 顶针。那么
\begin{equation}\label{eq:18.5}
\widetilde\Delta_\gamma=\beta_{\widetilde E}^{-1}(\Delta_\gamma)
\end{equation}
为 $\widetilde E$ 中的光滑 Lagrangian 球面。可将其表示为匹配圈，从而看出这一点，具体如下。在相差定向反转的意义下，存在唯一的光滑路径 $\widetilde\gamma:[-1,1]\to\widetilde S$，满足 $\beta_{\widetilde S}(\widetilde\gamma(t))=\gamma(t^2)$。$\widetilde\gamma$ 的端点是临界值 $\gamma(1)$ 的原像，因此其两个半段 $\widetilde\gamma_\pm(t)=\widetilde\gamma(\pm t):[0,1]\to\widetilde S$ 均为 $\widetilde\pi$ 的消失路径。如上所述，辛联络是 $\pi$ 的辛联络的拉回，由此可知，$\widetilde\gamma_\pm$ 的每个 Lefschetz 顶针都投影到 $\gamma$ 的 Lefschetz 顶针。这意味着，$\widetilde\gamma_\pm$ 的消失圈重合，等于 $M$ 内 $\gamma$ 的消失圈；而\eref{18.5}就是 $\widetilde\gamma$ 的朴素意义下的匹配圈，见\secref{16g}开头。
\end{definition*}
可以用另一个正则值 $w\in S\setminus\operatorname{Critv}(\pi)$ 代替 $*=1-3\delta$ 作为分歧点。证明二重分歧覆盖为 Lefschetz 纤维化的技术细节要求 $|w|\le1-2\delta$。构造 (U) 型 Lagrangian 子流形时，须假设 $L$ 落在
\SourcePage{276}
$|\pi(x)|\le|w|-\delta$ 的区域内。此外，$\widetilde L_+$ 与 $\widetilde L_-$ 的区分依赖于 $\sqrt w$ 的选择。当然，这意味着，若令 $w$ 沿例如半径 $1-3\delta$ 的圆周移动，这两个分支就会互换。下面将对此作严格说明，见\stref{引理}{18.1}。

\clearpage
\subsection{二重覆盖族}\label{sec:18b}
现在构造一个 Lefschetz 双纤维化，将 $\widetilde\pi:\widetilde E\to\widetilde S$ 及其他分歧点所对应的类似纤维化包含在内，也包括退化的，即奇异的情形。令 $\mathcal W\subset\CC$ 为半径 $1-2\delta$ 的圆盘。考虑沿对角线分歧的 $\mathcal W\times S$ 的二重覆盖，以及沿 $\pi$ 的图像分歧的 $\mathcal W\times E$ 的二重覆盖；更精确地说，使用对角线或图像中相应坐标落在 $\mathcal W\subset S$ 内的部分：
\[
\begin{aligned}
\Sfamily&=\{(w,y)\in\mathcal W\times\CC:y^2+w\in S\},\\
\Efamily&=\{(w,y,x)\in\Sfamily\times E:y^2+w=\pi(x)\}.
\end{aligned}
\]
$\Sfamily$ 带有自然复结构，而 $\Efamily$ 带有殆复结构；二者的边界均弱凸。有下列伪全纯映射：
\[
\begin{aligned}
\iota_{\Sfamily}(w,y)&=(w,-y):\Sfamily\to\Sfamily,\\
\iota_{\Efamily}(w,y,x)&=(w,-y,x):\Efamily\to\Efamily
&&\text{覆盖对合；}\\
\beta_{\Sfamily}(w,y)&=y^2+w:\Sfamily\to S,\\
\beta_{\Efamily}(w,y,x)&=x:\Efamily\to E
&&\text{忘掉 $w$ 的商映射；}\\
\varpi(w,y,x)&=(w,y):\Efamily\to\Sfamily
&&\text{对每个 $w$，为 $\pi$ 的提升；}\\
\psi(w,y)&=w:\Sfamily\to\mathcal W,\\
\Psi(w,y,x)&=w:\Efamily\to\mathcal W
&&\text{向参数 $w$ 的投影。}
\end{aligned}
\]
注意 $\Psi=\psi\circ\varpi$。$\varpi$ 的奇点为
\begin{equation}\label{eq:18.6}
\begin{aligned}
\operatorname{Crit}(\varpi)&=\coprod_{x\in\operatorname{Crit}(\pi)}\{y^2+w=\pi(x)\},\\
\operatorname{Critv}(\varpi)&=\coprod_{z\in\operatorname{Critv}(\pi)}\{y^2+w=z\}.
\end{aligned}
\end{equation}
在第一个集合中，$\Psi$ 的临界点恰为形如 $(w=\pi(x),0,x)$ 的那些点，故 $\operatorname{Critv}(\Psi)=\operatorname{Critv}(\pi)$。\znote{268}{原书此处印作 $\operatorname{Critv}(\Psi)=\operatorname{Crit}(\pi)$，两侧所在的空间并不相同：临界值落在 $\mathcal W\subset\CC$ 中，而临界点是 $E$ 中的点，故应视为笔误。与之相符的正确关系是 $\operatorname{Critv}(\Psi)=\operatorname{Critv}(\pi)$，译文即按此印出，可与\eref{18.6}第二式（$\operatorname{Critv}(\varpi)=\coprod_{z\in\operatorname{Critv}(\pi)}\{y^2+w=z\}$）对照。按《译校说明》的体例，原书笔误不回改，仅在此注明。[来源：原书印刷页 268 与\eref{18.6}；原书自身前后不一致，属笔误]}若在满足 $\pi(x)=z$ 的点 $x\in\operatorname{Crit}(\pi)$ 附近选择局部全纯 Morse 坐标，则 $\Efamily$ 的定义方程成为
\begin{equation}\label{eq:18.7}
w=z-y^2+x_1^2+\cdots+x_{n+1}^2;
\end{equation}
消去变量 $w$ 后，有 $\varpi(x,y)=(z-y^2+x_1^2+\cdots+x_{n+1}^2,y)$。由此容易看出，$\varpi$ 与 $\Psi$ 的临界点都是一般的。此外，将 $\Psi$ 在临界点 $(z,0,x)$ 处的 Hessian 限制到
\SourcePage{277}
$\ker(D\varpi)$，即与 $y=0$ 相切的子空间上，仍然非退化。这些计算说明，就殆复几何而言，
\SourceEquationDiagram{diagram-18-8}{0.232615134}{18.8}
满足 Lefschetz 双纤维化的全部性质；事实上，由于没有伪临界点，这是双纤维化的一个特别简单的例子。尽管整个讨论十分初等，仍值得用更直观的方式重述要点。对每个 $w$，限制
\begin{equation}\label{eq:18.9}
\varpi_w:\Efamily_w=\Psi^{-1}(w)\longrightarrow\Sfamily_w=\psi^{-1}(w)
\end{equation}
是沿二重分歧覆盖 $\beta_{\Sfamily}|\Sfamily_w:\Sfamily_w\to S$ 对 $\pi$ 作的拉回；特别地，$w=*$ 时，就回到前面定义的 $\widetilde\pi:\widetilde E\to\widetilde S$。只要 $w\notin\operatorname{Critv}(\pi)$，$\Efamily_w$ 就光滑，且\eref{18.9}具有一般临界点；它们成对出现，对应于 $\pi$ 的各临界点。更精确地说，固定 $w$ 后，每个 $x\in\operatorname{Crit}(\pi)$ 对应于两点 $(w,\pm\sqrt{\pi(x)-w},x)$。当 $w$ 趋于临界值 $z=\pi(x)$ 时，两个平方根汇合，其汇合处就是 $\Efamily_z$ 的唯一奇点。

取使 $\mathcal W$ 成为带边界正合辛流形的 $\theta_{\mathcal W},\omega_{\mathcal W}$，并假设它们充分大。考虑如下的一形式 $\theta_{\Sfamily}$：将 $\theta_{\mathcal W}$ 沿 $\psi$ 的拉回、$\kappa$ 沿第二投影的拉回，以及 $\theta_S$ 沿 $\beta_{\Sfamily}$ 的拉回相加。令 $\omega_{\Sfamily}=d\theta_{\Sfamily}$。已知 $\omega_{\Sfamily}|\Sfamily_*=\omega_{\widetilde S}$ 为辛形式，在其他任意纤维上的限制也如此。因此，由于加入了底空间上充分大的辛形式 $\omega_{\mathcal W}$ 的拉回，$\omega_{\Sfamily}$ 本身也是辛形式。用\stref{引理}{15.3}中的相同论证，特别参见\eref{15.7}，可知负 Liouville 向量场沿位于 $\partial\mathcal W$ 上方的边界面 $\partial^v\Sfamily$ 指向内部。为处理另一边界面，考虑部分平凡化
\begin{equation}\label{eq:18.10}
\mathcal W\times\widetilde S\supset\mathcal W\times\widetilde U\longrightarrow\Sfamily,\qquad
(w,y)\longmapsto(w,\sqrt{y^2+*-w}),
\end{equation}
它定义在子集 $\widetilde U=\{|y^2+*|\ge1-\delta\}\subset\widetilde S$ 上；选择在 $w=*$ 时给出恒等映射的唯一平方根。\eref{18.10}的像是逐纤维边界 $\partial^h\Sfamily$ 的邻域。由于在 $\mathcal W\times\widetilde U$ 上有 $|\sqrt{y^2+*-w}|\ge\sqrt\delta$，该部分平凡化与给定辛结构及其原始一形式相容，具体来说，$\theta_{\Sfamily}$ 沿\eref{18.10}的拉回等于 $\theta_{\mathcal W}+(\theta_{\widetilde S}|\widetilde U)$。由此容易推出，$\Sfamily$ 是带角正合辛流形，且 $\psi$ 的水平切空间平行于 $\partial^h\Sfamily$。

相应地，令 $\theta_{\Efamily}=\varpi^*(\theta_{\mathcal W}+\kappa)+\beta_{\Efamily}^*\theta_E$。与前面相同的论证说明，若 $\omega_{\mathcal W}$ 充分大，则 $\omega_{\Efamily}=d\theta_{\Efamily}$ 除了可能在
\SourcePage{278}
$\Psi$ 的临界点附近以外都是辛形式。现在，如\eref{18.7}中那样，在局部坐标 $(x,y)$ 下考察一个临界点，则 $\omega_{\Efamily}$ 等于 $y$ 坐标中的 $d\kappa$、$x$ 坐标中的 $\omega_E$，再加上 $\omega_{\mathcal W}$ 沿 $\Psi(y,x)=z-y^2+x_1^2+\cdots+x_{n+1}^2$ 的拉回；这显然是 Kähler 形式。接下来考察 $\Efamily$ 的各个边界面。(i) 对于 $|w|=1-2\delta$ 所定义的面，即 $\Psi^{-1}(\partial\mathcal W)$，前面处理 $\partial^v\Sfamily$ 的论证适用，说明负 Liouville 向量场指向内部。(ii) 在 $|y^2+w|=1$ 所定义的面附近，即 $x\in\partial^vE$ 处，有类似于\eref{18.10}的部分平凡化。因此，Liouville 向量场的行为归约为 $\mathcal W\times\widetilde E$ 中的行为。注意，该面是 $\Efamily$ 关于 $\Psi$ 的水平边界的一部分；这一平凡化还说明，$\Psi$ 的辛联络在该面附近平凡。(iii) 剩余的面由满足 $x\in\partial^hE$ 的点 $(w,y,x)$ 组成。在这样的点处，考虑
\begin{equation}\label{eq:18.11}
D\beta_{\Efamily}:T\Efamily=T\Efamily^h\oplus T\Efamily^v\longrightarrow TE=TE^h\oplus TE^v.
\end{equation}
第二项 $T\Efamily^v$ 为沿 $\varpi$ 的纤维的切空间，因此同构地映到 $TE^v$。第一项 $T\Efamily^h$ 是它关于 $\omega_{\Efamily}$ 的正交补。它实际上也就是关于 $\beta_{\Efamily}^*\omega_E$ 的正交补，因为这两个形式之差从 $\varpi$ 的底空间拉回。因此，$D\beta_{\Efamily}$ 保持分解\eref{18.11}。由于映射 $\beta_{\Efamily}$ 与 $\partial^hE$ 横截，当前边界面上的切空间是 $T(\partial^hE)$ 在\eref{18.11}下的原像，因而包含整个 $T\Efamily^h$。这说明，与 $\varpi$ 相关的辛联络的水平切空间，在 $(w,y,x)$ 处平行于边界；显然，$\Psi$ 也因此具有相同性质。此外，$T\Efamily$ 中的负 Liouville 向量与 $TE$ 中相应向量场的差别，只是一个 $T\Efamily^h$ 分量。由于该分量与边界相切，负 Liouville 向量必指向内部。这样，已经证明 $\Efamily$ 是带角正合辛流形，$\varpi,\Psi$ 是带奇点的正合纤维化；这就验证了\eref{18.8}满足 Lefschetz 双纤维化定义中关于辛结构的要求。为便于后面引用，记录关于 $\Psi$ 的几何的一个简单观察：

\StatementNumber{1}\begin{lemma}\label{st:18.1}
设 $\lambda$ 为沿 $\mathcal W$ 中半径 $|*|=1-3\delta$ 的圆周正向绕行一周的闭路径。考虑 $\Efamily_*=\widetilde E$ 上相应的单值变换自同构。在子集 $\widetilde E^{\mathrm{shrink}}$ 上，该映射与对合 $\iota_{\widetilde E}$ 一致。
\end{lemma}
\begin{proof}
在任意满足 $|w|=1-3\delta$ 且 $|\pi(x)|\le1-4\delta$ 的点 $(w,x,y)$，有 $|y|=|\sqrt{\pi(x)-w}|\ge\sqrt\delta$，因此辛形式中 $d\kappa$ 的贡献为零。于是，$\Psi$ 的辛联络的限制，等于沿投影到 $\mathcal W\times E$ 的平凡联络的拉回。因此，沿
\SourcePage{279}
$\lambda(t)=(1-3\delta)e^{2\pi it}$ 的平行移动为
\begin{equation}\label{eq:18.12}
\begin{aligned}
h_{\lambda|[0,t]}:(*,y,x)&\longmapsto\bigl(\lambda(t),\sqrt{y^2+\lambda(0)-\lambda(t)},x\bigr)\\
&=\left(\lambda(t),y\sqrt{\frac{\pi(x)-\lambda(t)}{\pi(x)-\lambda(0)}},x\right),
\end{aligned}
\end{equation}
其中选择平方根使\eref{18.12}在 $t=0$ 时为恒等映射。当 $t$ 从 $0$ 变到 $1$ 时，$\lambda(t)-\pi(x)$ 沿原点绕行一周，因此平方根改变符号。\znote{271}{这里补两句。其一，$\lambda$ 是半径 $1-3\delta$ 的整圆，而讨论中的点满足 $|\pi(x)|\le1-4\delta$，故 $\lambda(t)-\pi(x)$ 恒不为零，绕行一周是良定义的；其二，平方根变号正是 $\widetilde S\to S$ 绕分歧点 $*$ 一周交换两叶，即对合 $\iota_{\widetilde E}$：$\widetilde E$ 上正支与负支的区分只能在单连通区域内作出。§18a 末尾关于「令 $w$ 沿半径 $1-3\delta$ 的圆周移动会使 $\widetilde L_\pm$ 互换」的提醒，是同一事实的另一面。[来源：原书 §18a、§18b（印刷页 268、271）]}
\end{proof}

\subsection{二重覆盖中的消失圈与匹配圈}\label{sec:18c}
取 $\pi$ 的消失路径 $\gamma:[0,1]\to S$，从 $\gamma(0)=*$ 出发，并始终落在较小的圆盘 $\mathcal W\subset S$ 内。由于 $\Psi:\Efamily\to\mathcal W$ 与 $\pi$ 有相同的临界值，也可将 $\gamma$ 视为该 Lefschetz 纤维化的消失路径。因此，有相应的消失圈 $\widetilde V_\gamma$，它是 $\widetilde E=\Efamily_*$ 中的 Lagrangian 球面。另一方面，已在\secref{18a}中看到，$\gamma$ 的原像经适当参数化成为 $\widetilde\gamma:[-1,1]\to\widetilde S$，是 $\widetilde\pi$ 的匹配路径，其相应匹配圈为 $\widetilde\Delta_\gamma$。\znote{272}{这里用了 §16g 的匹配圈构造：匹配路径 $\widetilde\gamma$ 的两个半段各有一个 Lefschetz 顶针，二者的公共边界恰是消失圈（此处即 $M$ 中 $\gamma$ 的消失圈），沿该圈把两个顶针粘合，所得球面就是匹配圈。于是同一族消失路径有两种读法：把 $\gamma$ 看成 $\Psi$ 的消失路径，得到 $\widetilde E$ 中的球面 $\widetilde V_\gamma$；把 $\widetilde\gamma$ 看成 $\widetilde\pi$ 的匹配路径，得到 $\widetilde\Delta_\gamma$。引理 18.2 断言二者只差带标架同痕，这正是全节把 (B) 型膜同时写成顶针与匹配圈的合法性所在。[来源：原书 §16g、§18c（印刷页 272）]}

\StatementNumber{2}\begin{lemma}\label{st:18.2}
作为 $\widetilde E$ 中的带标架 Lagrangian 球面，$\widetilde\Delta_\gamma$ 与 $\widetilde V_\gamma$ 同痕。
\end{lemma}
\begin{proof}
其中大部分内容可以从双纤维化的一般理论，具体而言从\stref{引理}{16.15}推出。考察 $\operatorname{Critv}(\varpi)\subset\Sfamily$ 的结构，见\eref{18.6}，就能看出，与 $\gamma$ 相关的匹配路径与 $\widetilde\gamma$ 同痕。这一论证唯一的问题是，\stref{引理}{16.15}没有刻画其中出现的同痕 $I$。当然，在当前情形中，这个同痕有些多余，因为匹配圈可以用朴素方式构造；不过，仍值得核查各个要素确实能够正确结合。与往常一样，先在唯一临界点 $x\in E_z$，$z=\gamma(1)$，附近作局部变形，使 Kähler 形式 $\omega_E$ 在局部全纯 Morse 坐标下成为标准形式。还可假设，当 $t$ 接近 $1$ 时，$\gamma(t)=z+(1-t)$。这两种改变都允许，因为它们不影响所考虑的两个带标架 Lagrangian 球面的同痕类。取 $s$ 充分接近 $1$，考虑短消失路径 $\gamma|[s,1]$，及其在 $\Sfamily_{\gamma(s)}\to S$ 下的原像；该原像是 $\varpi_{\gamma(s)}$ 的匹配路径。由假设，这一原像实际上就是连接 $z-\gamma(s)$ 的两个平方根的直线段。

取\eref{18.7}中的局部坐标 $(y,x)$，使 $\omega_{\Efamily}=d\kappa+\omega_E+\Psi^*\omega_{\mathcal W}$，其中第二项现在是 $x$ 坐标中的标准形式。这并不完全是 $\CC^{n+2}$ 上的标准形式，但由 $y\in i\R$、$x\in\R^{n+1}$ 定义的子空间，仍是 $\omega_{\Efamily}$-Lagrangian 的。根据前面对 Lefschetz 顶针的刻画，容易看出，该子空间的适当一部分必为 $\gamma|[s,1]$ 的 Lefschetz 顶针。因此，消失圈为
\begin{equation}\label{eq:18.13}
\widetilde V_{\gamma|[s,1]}=\{y\in i\R,\ x\in\R^{n+1}:|y|^2+|x|^2=\gamma(s)-z\},
\end{equation}
\SourcePage{280}
并带有显然的标架。到目前为止，$\gamma|[s,1]$ 都被视为 $\Psi$ 的消失路径。若暂时再将其看作 $\pi$ 的消失路径，相应的 Lefschetz 顶针当然是 $\{x\in\R^{n+1}:|x|^2\le\gamma(s)-z\}$；因此，这个顶针在 $\beta_{\Efamily}|\Efamily_{\gamma(s)}$ 下的原像，按定义就是 $\widetilde\Delta_{\gamma|[s,1]}$，与\eref{18.13}相同。再像\stref{引理}{16.15}中那样，使用适当的平行移动映射，将这一事实传回整条路径 $\gamma$。
\end{proof}

为便于后面引用，再考察另外两种几何情形，其中仍涉及上面定义的各类 Lagrangian 子流形。假设给定 $\pi$ 的三条消失路径 $\gamma_k,\gamma_{k+1},\delta_k$，如\sr{fig:16.2}{图~16.2}所示；当然，此时不需要整个特选基，不过为方便起见仍沿用该记号。以 $\widetilde\Delta_k,\widetilde\Delta_{k+1},\widetilde\Pi_k$ 表示 $\widetilde E$ 中相应的 (B) 型 Lagrangian 子流形。
\StatementNumber{3}\begin{lemma}\label{st:18.3}
作为 $\widetilde E$ 中的带标架 Lagrangian 球面，$\widetilde\Pi_k$ 与 $\tau_{\widetilde\Delta_k}(\widetilde\Delta_{k+1})$ 同痕。
\end{lemma}
\begin{proof}
不失一般性，可假设这些路径包含于 $\mathcal W$。由\stref{引理}{18.2}，将 $\gamma_k$ 看成 $\Psi$ 的消失路径后，$\widetilde\Delta_k$ 就是其消失圈。因此，相应的 Dehn 扭转是 $\Psi$ 沿路径 $\lambda$ 的单值变换，其中 $\lambda$ 沿该消失路径绕去再绕回，即
\[
\tau_{\widetilde\Delta_k}(\widetilde\Delta_{k+1})\simeq h_\lambda(\widetilde\Delta_{k+1}).
\]
按定义，$\delta_k$ 由取 $\gamma_{k+1}$、再将其起点沿 $\lambda$ 移动得到。这意味着存在 $\pi$ 的消失路径的光滑族 $\gamma^s$，满足 $\gamma^s(0)=\lambda(s)$、$\gamma^s(1)=\gamma_{k+1}(1)$，从 $\gamma^0=\gamma_{k+1}$ 开始，到 $\gamma^1=\delta_k$ 结束。它们给出带标架 Lagrangian 球面的光滑族 $\widetilde\Delta^s\subset\Efamily_{\lambda(s)}$，在 $\widetilde\Delta^0=\widetilde\Delta_{k+1}$ 与 $\widetilde\Delta^1=\widetilde\Pi_k$ 之间插值。考察这些球面沿 $\lambda|[s,1]$ 平行移动后的像，即得到 $h_\lambda(\widetilde\Delta_{k+1})$ 与 $\widetilde\Pi_k$ 之间的同痕。
\end{proof}

最后，设 $\mu$ 为 $\pi$ 的匹配路径，$\Sigma_\mu\subset E$ 为其匹配圈。由假设，负 Liouville 流将每个 Lagrangian 子流形移入 $E^{\mathrm{shrink}}$；又因为 $\Sigma_\mu$ 只在相差同痕的意义下定义，故不妨假设它已经位于 $E^{\mathrm{shrink}}$ 内。这样，其原像 $\widetilde\Sigma_\mu=\widetilde\Sigma_{\mu,+}\cup\widetilde\Sigma_{\mu,-}$ 就是 (U) 型 Lagrangian 子流形，由 $\widetilde E$ 中两个不交的带标架 Lagrangian 球面组成。选择从 $*$ 出发、分别终止于 $\mu$ 两端的消失路径 $\gamma_k,\gamma_{k+1}$，如\sr{fig:18.1}{图~18.1}所示。关键要求是，这三条路径围成的三角形内部不含 $\pi$ 的临界值，从而复合路径 $\gamma'_{k+1}=\mu\#\gamma_k$ 是与 $\gamma_{k+1}$ 同痕的消失路径；其形成方式与\secref{16h}相同。设 $\widetilde\Delta_k,\widetilde\Delta_{k+1}$ 为与 $\gamma_k,\gamma_{k+1}$ 相关的 (B) 型子流形。
\StatementNumber{4}\begin{lemma}\label{st:18.4}
作为带标架 Lagrangian 球面，$\widetilde\Delta_{k+1}$ 与 $\tau_{\widetilde\Sigma_{\mu,-}}\tau_{\widetilde\Sigma_{\mu,+}}(\widetilde\Delta_k)$ 同痕。
\end{lemma}
\SourcePage{281}
\SourceFigure{figure-18-1}{0.469384110}{18.1}
\begin{proof}
以 $\widetilde\mu=\widetilde\mu_+\cup\widetilde\mu_-$ 及 $\widetilde\gamma_k,\widetilde\gamma_{k+1},\widetilde\gamma'_{k+1}$ 表示相应路径在 $\widetilde S$ 中的原像。它们都是 $\widetilde\pi$ 的匹配路径，而且
\[
\widetilde\gamma_{k+1}\simeq\widetilde\gamma'_{k+1}=\widetilde\mu_-\#(\widetilde\mu_+\#\widetilde\gamma_k).
\]
两次应用\stref{引理}{16.13}，即可得到所需结论，唯一可能剩下的是标架问题；但这也容易解决：对原 Lefschetz 纤维化 $E$ 的辛联络作变形，使 $\mu$ 的匹配圈可以用朴素方式构造，即同痕 $I$ 为常值。同样，对于匹配路径 $\widetilde\gamma_k,\widetilde\gamma_{k+1}$，同痕也为常值，这与证明\stref{引理}{16.13}时使用的 Lagrangian 连通和论证相容。
\end{proof}

\StatementNumber{5}\begin{note}\label{st:18.5}
以上三个引理都有将对合 $\iota_{\widetilde E}$ 考虑在内的加强版本：
\begin{itemize}
\item \stref{引理}{18.2}中的同痕自然等变，因为在局部模型中成立的是等式，而用于归约到该模型情形的变形与 $\iota_{\Efamily}$ 相容。
\item 对于\stref{引理}{18.3}，只须使用 Picard–Lefschetz 公式的等变版本。这里指出一个有用的附加细节：将该同痕与 $\tau_{\widetilde\Delta_k}|\widetilde\Delta_{k+1}$ 复合，得到等变微分同胚
\begin{equation}\label{eq:18.14}
\widetilde\Delta_{k+1}\longrightarrow\widetilde\Pi_k,
\end{equation}
因此，可以比较两个球面内不动点集的余定向。这些余定向自然对应于 $\widetilde S$ 中路径 $\widetilde\delta_k,\widetilde\gamma_{k+1}$ 的定向。结论是，两个余定向与\eref{18.14}相容，当且仅当它们在唯一公共点 $y=0$ 处的相交数 $\widetilde\delta_k\cdot\widetilde\gamma_{k+1}$ 为 $-1$。为看出这一点，取\stref{引理}{18.3}证明中的环路 $\lambda$，比较相应的 $\widetilde E$ 上的单值变换映射与
\SourcePage{282}
$\widetilde S$ 上的辫单值变换映射。一般而言，辫单值变换是保持 $\operatorname{Critv}(\widetilde\pi)$ 的 $\widetilde S$ 的微分同胚；在当前情形中，它就是沿 $\widetilde\gamma_k$ 的半扭转。若为 $\widetilde\gamma_{k+1}$ 选定定向，并为其像 $\widetilde\delta_k=t_{\widetilde\gamma_k}(\widetilde\gamma_{k+1})$ 赋予诱导定向，则局部相交数确实为 $-1$，见\sr{fig:18.2}{图~18.2}。
\SourceFigure{figure-18-2}{0.365538068}{18.2}
\item 对这一目的而言，我们对\stref{引理}{18.4}的证明不甚理想，因为每次只作一个 Dehn 扭转破坏了对称性。不过，适当改变论证便可轻易补救：先在 $E$ 中相对于 $M$ 构造 Lagrangian 球体的同痕 $\Delta_{k+1}\simeq\tau_{\Sigma_\mu}(\Delta_k)$，再将其提升到二重分歧覆盖。
\end{itemize}
\end{note}

\subsection{分次}\label{sec:18d}
本节余下部分始终假设 $E$ 带有二次复体积形式 $\eta_E^2$。除以底空间 $S$ 上的形式 $dz^2$，得到相对形式 $\eta_{E/S}^2=\eta_E^2/dz^2$。$\eta_E^2$ 到 $\widetilde E$ 的拉回在分歧集 $M=E_*$ 之外非零。为了解该分歧集上的情形，写 $\eta_E^2=\eta_{E/S}^2\wedge dz^2$；由于该纤维非奇异，这在 $M$ 附近有意义。拉回后出现因子 $4y^2dy^2$，因此可在 $\widetilde E$ 上定义二次复体积形式及相应的相对形式，令
\begin{equation}\label{eq:18.15}
\eta_{\widetilde E}^2=\frac{\beta_{\widetilde E}^*\eta_E^2}{4y^2},\qquad
\eta_{\widetilde E/\widetilde S}^2=\eta_{\widetilde E}^2/dy^2.
\end{equation}
注意，这样安排后，非奇异纤维 $\widetilde E_y$ 上诱导的形式，与 $E_z$ 上的形式一致，其中 $z=\beta_{\widetilde S}(y)$。\znote{275}{因子 $4y^2$ 就是换元的 Jacobi 因子：由 $z=y^2+*$ 得 $dz=2y\,dy$，故 $dz^2=4y^2dy^2$；除以它之后 $\eta^2_{\widetilde E}$ 在底坐标 $z$ 下就是 $\beta_{\widetilde E}^*\eta^2_{E/S}\wedge dy^2$。若不作这一步，拉回的 $\beta_{\widetilde E}^*\eta^2_E$ 在分歧集 $M=\{y=0\}$ 上会带上二阶零点，这正是本节开头说「拉回在 $M$ 之外非零」的原因。又注意 $\eta_{\widetilde E}=\beta_{\widetilde E}^*\eta_E/2y$ 的符号依赖分支的选取，笔记 18.8 关于两个分支定向相反的讨论就建立在这一选择上。[来源：原书 §18d（印刷页 275）]}同样的论证可以应用于 $\Efamily$，将其看成沿 $\{w=\pi(x)\}$ 分歧的 $\mathcal W\times E$ 的二重覆盖。其结果
\SourcePage{283}
是，Lefschetz 纤维化 $\Psi:\Efamily\to\mathcal W$ 带有形式
\[
\eta_{\Efamily}^2=\frac{\beta_{\Efamily}^*\eta_E^2\wedge dw^2}{4y^2},\qquad
\eta_{\Efamily/\mathcal W}^2=\eta_{\Efamily}^2/dw^2,
\]
满足：在正则纤维 $\Efamily_w$ 上诱导的形式等于 $(\beta_{\Efamily}|\Efamily_w)^*\eta_E^2/4y^2$；特别地，$w=*$ 时就得到\eref{18.15}。回顾，在这种情形下，沿底空间 $\mathcal W$ 正则部分中任意环路的单值变换映射，自然成为分次辛自同构。

\StatementNumber{6}\begin{lemma}\label{st:18.6}
取\stref{引理}{18.1}中的单值变换，设 $\gamma^\#$ 为其自然分次。那么，对所有 $(y,x)\in\widetilde E^{\mathrm{shrink}}$，都有 $\gamma^\#(y,x)=-1$。
\end{lemma}
\begin{proof}
由 $\eta_{\Efamily/\mathcal W}^2$ 的定义及\eref{18.12}，在每个满足假设的点 $(y,x)$ 处，有
\begin{equation}\label{eq:18.16}
h_{\lambda|[0,t]}^*(\eta_{\Efamily/\mathcal W}^2|\Efamily_{\lambda(t)})
=\eta_{\widetilde E}^2\cdot\frac{\pi(x)-\lambda(0)}{\pi(x)-\lambda(t)}.
\end{equation}
当 $t$ 从 $0$ 变到 $1$ 时，\eref{18.16}中的分式绕原点顺时针转一周，因此\eref{15.8}中的所有相位函数也如此，所需结论随之成立。
\end{proof}

还须重新考察\sr{fig:18.1}{图~18.1}。以 $V_k,\Delta_k,\widetilde\Delta_k$ 表示与 $\gamma_k$ 相关的消失圈、Lefschetz 顶针及 (B) 型 Lagrangian 子流形；以 $V'_{k+1},\Delta'_{k+1},\widetilde\Delta'_{k+1}$ 表示从 $\gamma'_{k+1}$ 得到的相同数据。$V_k$ 与 $\gamma_{k+1}$ 的消失圈同痕，因而也与 $V'_{k+1}$ 同痕；更精确地说，用来构造 $\mu$ 的匹配圈的同痕 $I$，诱导了它们之间的同痕。由于 $\gamma_k$ 与 $\gamma'_{k+1}$ 在基点附近相同，这又诱导了 $\Delta_k\cap U$ 与 $\Delta'_{k+1}\cap U$ 之间的逐纤维同痕，其中 $U\subset E$ 为 $E_*$ 的一个适当小邻域。提升到二重覆盖后，在 $U$ 的原像 $\widetilde U$ 内，得到 $\widetilde\Delta_k\cap\widetilde U$ 与 $\widetilde\Delta'_{k+1}\cap\widetilde U$ 之间的同痕。选择与该局部同痕相容的分次，再将后一分次沿变形传递给 $\widetilde\Delta_{k+1}$。
\StatementNumber{7}\begin{lemma}\label{st:18.7}
为 $\widetilde\Delta_k$ 与 $\widetilde\Delta_{k+1}$ 如此选择分次，并为涉及的 Dehn 扭转赋予自然分次后，\stref{引理}{18.4}中的同痕成为分次 Lagrangian 子流形的同痕。
\end{lemma}
\begin{proof}
通常的变形技巧将问题归约到如下情形：匹配圈 $\Sigma_\mu$ 可以用朴素方式构造，且 $E$ 上的辛形式在 $\gamma_k(1)$ 上方的临界点附近，在局部全纯 Morse 坐标下为标准形式。加入第二个要求，是为了利用如下事实：在\stref{例题}{16.5}的局部模型中，沿底空间任意路径 $\beta$ 的平行移动，都会将消失圈 $\sqrt{\beta(0)}S^n$ 映到 $\sqrt{\beta(1)}S^n$。
\SourcePage{284}
仔细选取 $\gamma'_{k+1}$，使其在临界值 $\gamma_k(1)$ 的一个小邻域之外，与 $\gamma_k$ 或 $\mu$ 一致。此时，沿各平行移动映射追踪，便可看出 $V_k=V'_{k+1}$。因此，仍取上述 $\widetilde U$，有 $\widetilde\Delta_k\cap\widetilde U=\widetilde\Delta'_{k+1}\cap\widetilde U$。按定义，这两个子流形的分次相同。由于 Dehn 扭转集中在 $\widetilde\Sigma_{\mu,\pm}$ 附近，因而在 $\widetilde U$ 上平凡，故
\begin{equation}\label{eq:18.17}
\tau_{\widetilde\Sigma_{\mu,-}}\tau_{\widetilde\Sigma_{\mu,+}}(\widetilde\Delta_k)\cap\widetilde U
=\widetilde\Delta'_{k+1}\cap\widetilde U,
\end{equation}
而由 Dehn 扭转的典范分次诱导的分次仍然一致。现在，\stref{引理}{16.13}给出的\eref{18.17}中两个 Lagrangian 子流形之间的同痕，仍局限于 $\widetilde\Sigma_{\mu,\pm}$ 附近，因此分次相等的性质保持不变。
\end{proof}

\subsection{一些预等变膜与等变膜}\label{sec:18e}
现在重新考察\secref{18a}中引入的 $\widetilde E$ 中两类 Lagrangian 子流形，目的是将它们变成等变 Fukaya 范畴 $\F(\widetilde E)^{\mathrm{naive}}$ 的对象。回顾，为定义该范畴，必须对 $\Z/2$ 作用施加一些技术条件，见\stref{假设}{14.3}。在当前情形中，这些条件显然满足：$\iota_{\widetilde E}$ 的不动点集为 $\{y=0\}\cong M$，其法丛有自然平凡化，由 $D\widetilde\pi|M:T\widetilde E|M\to T\widetilde S_0=\CC$ 给出。

\begin{description}[leftmargin=2em]
\item[(U) 型。]\mbox{}\par 设 $L^\#=(L,\alpha^\#,P^\#)$ 为 $E$ 中的正合 Lagrangian 膜，包含于 $E^{\mathrm{shrink}}$ 的内部。设 $\widetilde L$ 为 $L$ 在 $\widetilde E$ 中的原像。由\eref{18.15}，$L$ 与 $\widetilde L$ 的平方相位映射满足
\[
\alpha_{\widetilde E}(T\widetilde L_{(y,x)})=\frac{|y|^2}{y^2}\alpha_E(TL_x).
\]
回顾，$\widetilde L=\widetilde L_+\cup\widetilde L_-$，其中连通分支 $\widetilde L_\pm$ 位于子集 $\pm\operatorname{im}(y)>0$ 内。由于 $\operatorname{im}(y)$ 始终非零，可以取辐角函数 $\arg:\CC\setminus\R^+\to(0,2\pi)$，令
\begin{equation}\label{eq:18.18}
\widetilde\alpha^\#(y,x)=\alpha^\#(x)-\frac1{2\pi}\arg(y^2),
\end{equation}
从而定义 $\widetilde L$ 的分次。限制后，得到 $\widetilde L_\pm$ 的分次 $\widetilde\alpha^\#_\pm$。至于 Pin 结构，我们约定为 $\widetilde L_+$ 赋予给定 $P^\#$ 的拉回，为 $\widetilde L_-$ 赋予 $P^\#\otimes\lambda^{\mathrm{top}}(TL)$ 的拉回。将它们记为 $\widetilde P_\pm^\#\to\widetilde L_\pm$，将其并记为 $\widetilde P^\#\to\widetilde L$。

据此，容易将 $\widetilde L_+$ 赋予\secref{14a}意义下的预等变膜结构。其主体为 $\widetilde L_+^\#$，影子为 $\widetilde L_-^\#$，同构 $\widetilde f_-:\widetilde P_-^\#\to\iota_{\widetilde E}^*(\widetilde P_+^\#\otimes\lambda^{\mathrm{top}}(T\widetilde L_+))$ 为恒等同构，并调整 $\widetilde f_+$ 使\eref{14.1}成立。交换 $+$ 与 $-$ 后，相同程序给出 $\widetilde L_-$ 的预等变膜结构；取二者的不交并，$\widetilde L$ 就成为等变膜。
\end{description}
\SourcePage{285}
已经可以得出一些初步结论。设 $L_k^\#=(L_k,\alpha_k^\#,P_k^\#)$，$k=0,1$，为 $E$ 中上述类型的两个膜。为该子流形对选择 Floer 数据 $(H,J)$ 时，可以要求在 $E^{\mathrm{shrink}}$ 之外 $H$ 为零、$J$ 与 $I_E$ 一致。这不妨碍模空间 $\M_Z(y_0,y_1)$ 的正则性，因为在 $L_0\cup L_1$ 的邻域内仍可自由选择。对于原像对 $(\widetilde L_{0,+},\widetilde L_{1,+})$，取提升 $(H,J)$ 得到的 Floer 数据 $(\widetilde H,\widetilde J)$；由于所取 Floer 数据在分歧集附近平凡，这没有问题。每个 $\widetilde y_+\in\C(\widetilde L_{0,+},\widetilde L_{1,+})$ 始终落在 $\widetilde E^{\mathrm{shrink}}_+$ 内，并投影到某个 $y\in\C(L_0,L_1)$，从而建立了两个集合之间的双射。该对应保持绝对指标，并在典范同构意义下保持定向空间：
\begin{equation}\label{eq:18.19}
i(\widetilde y_+)=i(y),\qquad o(\widetilde y_+)\cong o(y).
\end{equation}
这里有一个小问题，可作如下处理：虽然 $\eta_{\widetilde E}^2|\widetilde E^{\mathrm{shrink}}_+$ 不等于 $\eta_E^2$ 的拉回，但可选一个将 $\CC\setminus\R^+$ 变形收缩到点 $-1$ 的映射，并将其作用于\eref{18.15}中的 $|y|^2/y^2$ 项，从而将一个形式变形为另一个，允许相差一个整体符号。这个二次复体积形式的同伦诱导分次的同伦；事实上，它将 $\widetilde L_k$ 的分次 $\widetilde\alpha_{k,+}^\#$ 变形为 $\alpha_k^\#-1/2$ 的拉回。Pin 结构当然本来就一致，因此\eref{18.19}容易成立。

对每个 $u\in\M_Z(y_0,y_1)$，映射 $v=\pi\circ u:Z\to S$ 在 $\Omega=v^{-1}(S\setminus S^{\mathrm{shrink}})$ 上全纯，故可对 $\overline\Omega$ 上的 $|v|$ 应用最大值原理。由于 $\Omega$ 相对紧且不与 $\partial Z$ 相交，结论是 $\Omega=\varnothing$，因此 $u$ 仍落在 $E^{\mathrm{shrink}}$ 内。这意味着，投影在这样的 $u$ 与元素 $\widetilde u_+\in\M_Z(\widetilde y_{0,+},\widetilde y_{1,+})$ 之间建立双射。该双射还将相应线性化算子等同，故 $(\widetilde H,\widetilde J)$ 仍然正则。最后，这一等同与用来定义 Floer 微分 $\partial$ 的贡献\eref{12.17}的粘合过程相容，从而定义了 Floer 上链复形的同构。

对原像的 $-$ 部分，相同论证也适用；唯一区别是 Pin 结构被 $\lambda^{\mathrm{top}}(TL_k)\cong\xi_E|L_k\cong\xi_{\widetilde E}|\widetilde L_{k,-}$ 扭曲。但这一扭曲同时作用于涉及的两个 Lagrangian 子流形，故 Floer 复形不受影响。总之，存在典范同构
\begin{equation}\label{eq:18.20}
\begin{aligned}
CF^*(\widetilde L_0^\#,\widetilde L_1^\#)
&\cong CF^*(\widetilde L_{0,+}^\#,\widetilde L_{1,+}^\#)\oplus CF^*(\widetilde L_{0,-}^\#,\widetilde L_{1,-}^\#)\\
&\cong CF^*(L_0^\#,L_1^\#)\otimes_{\K}\K[+,-],
\end{aligned}
\end{equation}
其中 $\K[+,-]$ 指具有两个生成元 $+$ 与 $-$ 的二维向量空间。\znote{277}{两个生成元 $+$ 与 $-$ 都在 $0$ 次：它们不是次数，也不是 §12 的 $[m]$ 平移记号，更不是群环。之所以能写成张量积，是因为\eref{18.20}的两个直和项次数分布完全相同（$\widetilde\alpha_{\pm}^\#$ 由同一个 $\widetilde\alpha^\#$ 限制而来）。与它配套的是下一句：$\widetilde L^\#$ 的等变性给出交换这两个生成元的对合，取不变部分便回到 $CF^*(L_0^\#,L_1^\#)$；命题 18.13 的证明正是靠这一点，把 $\widetilde{\A}^{\mathrm{inv}}$ 与 $E$ 中由较小膜构成的全子范畴等同起来。[来源：原书 §18e、命题 18.13（印刷页 277、282）]}由于 $\widetilde L_k^\#$ 的等变性，Floer 复形带有典范对合；关于\eref{18.20}，它恰为交换生成元 $+$ 与 $-$ 的映射。更一般地，给定任意同类 Lagrangian 膜的 $(d+1)$ 元组，可以选择具有与上述 Floer 数据相同性质的扰动数据。这样，扰动全纯 $(d+1)$ 边形也被限制在 $E^{\mathrm{shrink}}$ 内，意味着\eref{18.20}与 $\Ai$ 结构相容。
\SourcePage{286}
\StatementNumber{8}\begin{note}\label{st:18.8}
假设所有二次复体积形式都存在连续平方根，使各膜自然带有定向，见\stref{笔记}{11.18}。更明确地说，某点处切空间 $T\widetilde L_{(y,x)}$ 的基 $(v_1,\ldots,v_{n+1})$ 为正定向，当且仅当
\[
\exp(\pi i\widetilde\alpha^\#(y,x))=
\frac{\eta_{\widetilde E}(v_1\wedge\cdots\wedge v_{n+1})}{|\eta_{\widetilde E}(v_1\wedge\cdots\wedge v_{n+1})|}.
\]
从 $(y,x)$ 换到 $(-y,x)$ 时，$\eta_{\widetilde E}=\beta_{\widetilde E}^*\eta_E/2y$ 的符号改变，而按定义，$\widetilde\alpha^\#$ 保持不变。因此，两个分支 $\widetilde L_+$ 与 $\widetilde L_-$ 带有相反定向。
\end{note}

\begin{description}[leftmargin=2em]
\item[(B) 型。]\mbox{}\par 设 $\gamma:[0,1]\to D$ 为从 $\gamma(0)=*$ 出发的消失路径。要求 $\gamma'(0)\in\CC\setminus\R^+$，且 $\gamma|(0,1]$ 与基路径 $\chi$ 不交；从现在起，每当考虑 $\pi$ 的消失路径时，都作这些固定假设。设 $\widetilde\Delta_\gamma$ 为相应 Lefschetz 顶针 $\Delta_\gamma$ 在 $\widetilde E$ 中的原像。为其赋予任意分次 $\widetilde\alpha^\#$，以及同构意义下唯一的 Pin 结构 $\widetilde P^\#$。

为使该膜成为等变膜，还须选择一个同构 $f:\widetilde P^\#\to\iota_{\widetilde E}^*(\widetilde P^\#\otimes\lambda^{\mathrm{top}}(T\widetilde\Delta_\gamma))$。由\stref{引理}{14.2}，这类同构典范地对应于不动点集的余定向；在此情形中，该不动点集为 $V_\gamma\subset\widetilde\Delta_\gamma$。由于 $\widetilde\Delta_\gamma$ 纤维化在 $\gamma$ 的原像 $\widetilde\gamma$ 上，有
\begin{equation}\label{eq:18.21}
D\widetilde\pi(T\widetilde\Delta_\gamma|V_\gamma)=\R\widetilde\gamma'(0)=\R\sqrt{\gamma'(0)},
\end{equation}
因此，余定向自然对应于平方根的选择。我们的约定是，始终取\emph{虚部为正}的平方根。\eref{18.21}也表明，$\widetilde\Delta_\gamma$ 满足\stref{假设}{14.4}。最后，定义 $\F(\widetilde E)^{\mathrm{naive}}$ 要求为每个 $\widetilde\Delta_\gamma$ 选择实值法向角，这里取为 $\arg(\gamma'(0))/2\in(0,\pi)$。

如\secref{14d}所指出的，$\widetilde\Delta_\gamma$ 上的等变膜结构诱导其不动点集 $V_\gamma$ 上的膜结构。具体地，分次为
\[
\alpha^\#_{V_\gamma}(x)=\alpha^\#_{\widetilde\Delta_\gamma}(0,x)-\frac1{2\pi}\arg(\gamma'(0)),
\]
而 Pin 结构是在同构意义下唯一能作为 $\Delta_\gamma$ 上某个 Pin 结构的边界的结构。这意味着，当 $n=1$ 时，消失圈带有非平凡 Pin 结构，而不是平凡 Pin 结构；理由是，对任意紧连通曲面 $S$，$H^2(S,\partial S;\Z/2)\cong\Z/2$ 中的相对第二 Stiefel–Whitney 类等于模 $2$ Euler 示性数。对 $S=D$，该类非零，因此不可能存在沿边界平凡的 Pin 结构。
\end{description}
\SourcePage{287}
\StatementNumber{9}\begin{lemma}\label{st:18.9}
(i) 对任意 (B) 型等变膜 $\widetilde\Delta^\#=\widetilde\Delta_\gamma^\#$，有 $HF^*(\widetilde\Delta^\#,\widetilde\Delta^\#)\cong H^*(S^{n+1};\K)$。$\Z/2$ 作用在 $HF^0$ 上为 $\Id$，在 $HF^{n+1}$ 上为 $-\Id$。

(ii) 取消失路径对 $\gamma_k$，$k=0,1$，并设 $\widetilde\Delta_k^\#$ 为相应的 (B) 型等变膜。那么，Poincaré 对偶映射
\[
HF^*(\widetilde\Delta_1^\#,\widetilde\Delta_0^\#)\cong HF^{n+1-*}(\widetilde\Delta_0^\#,\widetilde\Delta_1^\#)^\vee
\]
与 $\Z/2$ 作用反相容，即相差一个整体 $-1$ 符号而等变。
\end{lemma}
\begin{proof}
这些纯粹是上同调层次上的陈述，由\secref{14b}中列出的等变 Floer 上同调的一般性质推出。事实上，$\iota_{\widetilde E}|\widetilde\Delta$ 是 $\widetilde\Delta$ 的反定向自同构，而 PSS 同构与 $\Z/2$ 作用相容，故 (i) 成立。对于 (ii)，回顾 Poincaré 对偶同构来自三角形乘积
\[
\begin{aligned}
&HF^*(\widetilde\Delta_1^\#,\widetilde\Delta_0^\#)\otimes HF^{n+1-*}(\widetilde\Delta_0^\#,\widetilde\Delta_1^\#)\\
&\qquad\longrightarrow HF^{n+1}(\widetilde\Delta_0^\#,\widetilde\Delta_0^\#)\cong H^{n+1}(\widetilde\Delta_0;\K)\cong\K,
\end{aligned}
\]
因此，更典范的写法应为
\[
HF^*(\widetilde\Delta_1^\#,\widetilde\Delta_0^\#)\cong H^{n+1}(\widetilde\Delta_0;\K)\otimes HF^{n+1-*}(\widetilde\Delta_0^\#,\widetilde\Delta_1^\#)^\vee.
\]
在这一形式下，它仍与 $\Z/2$ 作用相容；由 (i)，去掉 $H^{n+1}(\widetilde\Delta_0;\K)$ 因子便引入一个符号。\znote{279}{两条要点。其一，(i) 中两种作用的差别就是「$\Id$ 与 $-\Id$ 不同」的全部内容：若 $\operatorname{char}\K=2$，引理退化为「作用平凡」，此后一切用 $\Z/2$-不变性区分次数、进而判定态射群无上同调（命题 18.14）的论证都会失效。其二，(ii) 的反相容有具体来源：Poincaré 对偶把 $HF^*$ 与 $HF^{n+1-*}$ 配对到 $HF^{n+1}(\widetilde\Delta_0^\#,\widetilde\Delta_0^\#)\cong\K$，而 $\Z/2$ 在这个基类上按 $-\Id$ 作用，故对偶映射整体带一个符号。二者与\eref{18.23}合用，便得 $j>k$ 时 $HF^*(\widetilde\Delta_j^\#,\widetilde\Delta_k^\#)$ 上的作用为 $-\Id$。[来源：原书 §18e、§18f（印刷页 279、282）]}
\end{proof}

如同 (U) 型膜的情形，特殊选择 Floer 数据和扰动数据，有助于使链层次上的 $\Z/2$ 作用更清楚。取两条消失路径 $\gamma_k$，相应的 (B) 型膜为 $\widetilde\Delta_k^\#$，消失圈为 $V_k^\#$。假设
\begin{equation}\label{eq:18.22}
\begin{cases}
\arg(\gamma'_0(0))>\arg(\gamma'_1(0)),\\
\gamma_0([0,1])\cap\gamma_1([0,1])=\{*\},\\
V_0\pitchfork V_1.
\end{cases}
\end{equation}
为 $(\widetilde\Delta_0,\widetilde\Delta_1)$ 选择 Floer 数据 $(H,J)$ 时，要施加比\stref{假设}{14.5}强得多的要求：对每个 $J_t$，$\widetilde\pi$ 都为 $J_t$-全纯，且 $H=0$。第二个条件与\stref{假设}{14.5}相容，因为由\eref{18.22}，法向角满足 $\sigma_1-\sigma_0\in(-\pi,0)$，故\eref{14.10}中可以取 $h=0$。任意交点 $\widetilde y=(y,x)\in\C(\widetilde\Delta_0,\widetilde\Delta_1)$ 在 $\widetilde\pi$ 下必投影到 $\widetilde\gamma_0$ 与 $\widetilde\gamma_1$ 的唯一交点。因此 $y=0$，意味着 $\widetilde y$ 位于 $\widetilde E^{\mathrm{fix}}=M$ 内。按定义，每个 $\widetilde\Delta_k$ 上的等变 Pin 结构，对应于 $\widetilde\Delta_k^{\mathrm{fix}}$ 的法丛平凡化，该平凡化投影到虚部为正的平方根 $\sqrt{\gamma'_k(0)}$。该平方根为 $\exp(i\sigma_k)=\exp(i\arg(\gamma'_k(0))/2)$ 的正实数倍。因此，\stref{引理}{14.11}中的符号 $w_k$ 都为 $+1$，意味着 $|o(\widetilde y)|_{\K}$ 上的 $\Z/2$ 作用平凡。这就说明，整个上链群 $CF^*(\widetilde\Delta_0^\#,\widetilde\Delta_1^\#)$ 都是 $\Z/2$-不变的。

类似地，任意 $\widetilde u\in\M_Z(\widetilde y_0,\widetilde y_1)$ 都投影到全纯映射 $\widetilde v=\widetilde\pi\circ\widetilde u:Z\to\widetilde S$，满足 $\widetilde v(\R\times\{k\})\subset\widetilde\gamma_k([-1,1])$、$\lim_{s\to\pm\infty}\widetilde v(s,\cdot)=0$。任何这样的 $\widetilde v$ 的能量都
\SourcePage{288}
必为零，因此它是常值映射，意味着 $\widetilde u$ 也完全落在 $M$ 中。假设 $(H,J)$ 的选择使 $M$ 上的 Floer 模空间正则，则\stref{引理}{14.6}蕴含 $\widetilde E$ 中的正则性。此外，\stref{引理}{14.9}说明，无论在 $M$ 还是在 $\widetilde E$ 中工作，绝对指标与定向空间都相同；\eref{14.13}对每个孤立的 $\widetilde u$ 贡献于 Floer 微分时所带的符号，也作出了相同结论。总之，存在复形同构
\begin{equation}\label{eq:18.23}
CF^*(\widetilde\Delta_0^\#,\widetilde\Delta_1^\#)\cong CF^*(V_0^\#,V_1^\#),
\end{equation}
其中右端消失圈上的膜结构是从 $\widetilde\Delta_k^\#$ 诱导的。此外，左端的典范 $\Z/2$ 作用平凡。更一般地，假设有消失路径 $\gamma_0,\ldots,\gamma_d$，$d\ge2$，满足
\begin{equation}\label{eq:18.24}
\begin{cases}
\arg(\gamma'_k(0))>\arg(\gamma'_{k+1}(0))\quad\text{对所有 }k,\\
\gamma_j([0,1])\cap\gamma_k([0,1])=\{*\}\quad\text{若 }j\ne k,\\
\text{各 $V_k$ 处于一般位置。}
\end{cases}
\end{equation}
那么，选取满足 $K=0$ 且 $\widetilde\pi$ 为 $J$-全纯的扰动数据 $(K,J)$ 后，就发现所有全纯 $(d+1)$ 边形都落在 $M$ 中。假设正则性成立，见\stref{引理}{14.7}，它们对 $M$ 与 $\widetilde E$ 上的 $\Ai$ 结构也给出相同贡献。

\StatementNumber{10}\begin{note}\label{st:18.10}
由\stref{引理}{18.9}(ii)，\eref{18.23}意味着，对偶群 $HF^*(\widetilde\Delta_1^\#,\widetilde\Delta_0^\#)$ 上的典范 $\Z/2$ 作用为 $-\Id$。原则上，可取非齐次项为零的 Floer 数据，在上链层次上证明这一点。为使其满足\eref{14.10}，其中必有 $h=0$，就须将法向角 $\sigma_1$ 换成 $\sigma_1+\pi$；这意味着\stref{引理}{14.11}中 $w_1=-1$，使相关定向空间上的 $\Z/2$ 作用非平凡。遗憾的是，法向角的改变使这一选择与此前对 $(\widetilde\Delta_0,\widetilde\Delta_1)$ 使用的 Floer 数据不相容。具体地，尝试同时使用这两者，会遇到伪全纯多边形的等变横截性方面的潜在问题，因此这里放弃这种做法。这意味着，在我们的 $\F(\widetilde E)^{\mathrm{naive}}$ 版本中，链层次上的等式\eref{18.23}成立，而 $(\widetilde\Delta_1,\widetilde\Delta_0)$ 的对应陈述只在上同调层次上成立。以下\stref{命题}{18.14}决定了这一选择。
\end{note}

\StatementNumber{11}\begin{note}\label{st:18.11}
在混合情形中，即一个子流形 $\widetilde L=\beta_{\widetilde E}^{-1}(L)$ 为 (U) 型，而另一个 $\widetilde\Delta=\beta_{\widetilde E}^{-1}(\Delta)$ 为 (B) 型时，类似于\eref{18.20}的结果成立。首先，虽然 $\Delta$ 有边界，仍可容易地定义 $(L,\Delta)$ 的 Floer 上同调，方法如下。取在 $E^{\mathrm{shrink}}$ 之外平凡的 Floer 数据 $(H,J)$，使任意 $u\in\M_Z(y_0,y_1)$ 对应的 $v=\pi\circ u$ 在 $v^{-1}(S^{\mathrm{shrink}})$ 之外全纯。此时，$v$ 的像不包含 $*$；否则，由开映射性，它也必须包含附近不属于 $\gamma([0,1])$ 的点 $z$，这意味着相交数 $u\cdot[E_z]$ 为
\SourcePage{289}
正。由于 $S\setminus(S^{\mathrm{shrink}}\cup\gamma([0,1]))$ 连通，同样的正性也会对 $z\in\partial S$ 成立，产生矛盾。结论是，由于 $u$ 不与 $\partial\Delta$ 相交，在 Floer 理论的考察中可以忽略该边界的存在。据此，可以利用同痕不变性，将问题归约到如下情形：$t\in[0,\delta]$ 时 $\gamma(t)=*-t$，而 $\gamma$ 的其余部分落在 $S^{\mathrm{shrink}}$ 内。此时，可如前对 $|v|$ 应用最大值原理，得到
\[
HF^*(\widetilde L^\#,\widetilde\Delta^\#)\cong HF^*(L^\#,\Delta^\#)\otimes_{\K}\K[+,-].
\]
在该等式右端，$\Delta$ 带有唯一的平凡 Pin 结构，分次则从 $\widetilde\Delta$ 的分次下降得到，即反过来应用\eref{18.18}的过程。
\end{note}

\subsection{两个全嵌入}\label{sec:18f}
此时，引入一个技术条件很方便。若在消失路径特选基的每个同伦类中，都有一个基使其相应消失圈处于一般位置，就称 $\pi:E\to S$ 具有\emph{一般平行移动}。给定任意 Lefschetz 纤维化，都能找到 $(\theta_E,\omega_E,I_E)$ 的一个小扰动，使其在 $\partial E$ 及 $M=E_*$ 附近恒定，且扰动后的结构具有一般平行移动。这很容易从\stref{引理}{15.3}及特选基的同伦类只有可数多个这一事实推出。这种扰动不改变 $M$ 或 $E$ 的 Fukaya 范畴；后一情形指在拟同构意义下不变，并使用 Moser 引理。因此，从现在起假设给定的 Lefschetz 纤维化已具有一般平行移动。

从现在起要求 $\operatorname{char}(\K)\ne2$。这使我们能应用第\sr{sec:6}{6 节}的结果；更重要的是，它使\stref{引理}{18.9}中的 $\Id$ 与 $-\Id$ 不同。\znote{282}{这个假设不是解析上的方便，而是构造方式的副产品：全节把 $\widetilde E$ 当作 $E$ 的二重分歧覆盖来用，而对 $\Z/2$ 取不变部分（如引理 18.9(i) 中对 $HF^*(\widetilde\Delta^\#,\widetilde\Delta^\#)$ 取不变部分）必须能区分 $\pm1$ 两个本征值；特征 $2$ 时二者合并，信息全失。§18j 末尾明说这是为整套论证付出的「唯一代价」（印刷页 288）。因此引用定理 18.24、推论 18.25、18.26 时，不能把结论直接搬到 $\operatorname{char}\K=2$ 的情形；§6 的一般代数结果同样以此为前提。[来源：原书 §18f、§18j（印刷页 282、288）]}

取 $\widetilde{\A}\subset\F(\widetilde E)^{\mathrm{naive}}$ 为全 $\Ai$ 子范畴，其对象为预等变膜 $\widetilde L_\pm^\#$ 及等变膜 $\widetilde L^\#,\widetilde\Delta_\gamma^\#$。虽然下面大体上是在重复前面的讨论，我们仍想明确说明构造 $\widetilde{\A}$ 时使用哪些 Floer 数据与扰动数据。对于任意一对 (U) 型等变膜或预等变膜，应选择使\eref{18.20}成立的 Floer 数据。特别地，这蕴含每个这样的膜都是第\sr{sec:6a}{6a 小节}意义下 $H^0(\widetilde{\A})$ 的自由对象。类似地，对于 (U) 型膜的任意 $(d+1)$ 元组，$d\ge2$，取同类扰动数据，使该同构与 $\Ai$ 结构映射相容。其次，对于任意一对 (B) 型等变膜，Floer 数据应为满足\stref{假设}{14.5}的不变数据。不过，在这些膜来自满足\eref{18.22}的路径 $\gamma_0,\gamma_1$ 的特殊情形中，要采用导出同构\eref{18.23}的更受限制的选择。每当消失路径集 $(\gamma_0,\ldots,\gamma_d)$ 满足\eref{18.24}时，对扰动数据的选择也作同样要求。这与整体选择的一致性条件相容，因为这类 $(d+1)$ 元组对于取有序子集封闭。

\SourcePage{290}
\StatementNumber{12}\begin{definition}\label{st:18.12}
\emph{Lefschetz 纤维化的 Fukaya 范畴}记为 $\F(\pi)$，定义为不变部分 $\widetilde{\A}^{\mathrm{inv}}$。
\end{definition}

\StatementNumber{13}\begin{proposition}\label{st:18.13}
存在在上同调上全忠实的 $\Ai$ 函子
\[
\F(E)\longrightarrow\F(\pi),
\]
将每个位于 $E^{\mathrm{shrink}}$ 中的正合 Lagrangian 膜 $L^\#$ 映到其原像 $\widetilde L^\#$。
\end{proposition}
\begin{proof}
由假设，$E$ 中每个正合 Lagrangian 膜都同痕于一个位于 $E^{\mathrm{shrink}}$ 中的膜。因此，由这类膜构成的全 $\Ai$ 子范畴到 $\F(E)$ 的包含是拟等价，由\stref{定理}{2.9}可将其反演。事实上，考察该定理前的论证，就能找到一个逆拟等价，使它在该子范畴自身上的限制为恒等。另一方面，由\eref{18.20}，同一子范畴可等同于 $\widetilde{\A}^{\mathrm{inv}}$ 中由对象 $\widetilde L^\#$ 构成的全 $\Ai$ 子范畴。
\end{proof}

现令 $\symbfit{\gamma}=(\gamma_1,\ldots,\gamma_m)$ 为消失路径的特选基。对每个 $k=1,\ldots,m$，选择一个分次，使相关的 (B) 型 Lagrangian 子流形 $\widetilde\Delta_k$ 成为等变膜，并给相应消失圈 $V_k$ 赋予诱导的膜结构。分别以 $\widetilde{\symbfit{\Delta}}^{\#}$ 和 $\symbfit V^\#$ 表示由此得到的 $\F(\pi)$ 与 $\F(M)$ 的对象集。它们带有相应的有向 $\Ai$ 子范畴 $\F^\to(\widetilde{\symbfit{\Delta}}^{\#})$ 与 $\F^\to(\symbfit V^\#)$。

\StatementNumber{14}\begin{proposition}\label{st:18.14}
$\widetilde{\symbfit{\Delta}}^{\#}$ 是 $\F(\pi)$ 中的例外集。此外，$\F^\to(\widetilde{\symbfit{\Delta}}^{\#})$ 与 $\F^\to(\symbfit V^\#)$ 拟同构，拟同构由保持对象次序的严格含幺 $\Ai$ 函子给出。因此，存在在上同调上全忠实的 $\Ai$ 函子
\[
\F^\to(\symbfit V^\#)\longrightarrow\F(\pi),
\]
将每个 $V_k^\#$ 映到 $\widetilde\Delta_k^\#$。
\end{proposition}
\begin{proof}
先假设消失圈处于一般位置。由\stref{引理}{18.9}(i)，有
\[
\Hom_{H\F(\pi)}(\widetilde\Delta_k^\#,\widetilde\Delta_k^\#)
=HF^*(\widetilde\Delta_k^\#,\widetilde\Delta_k^\#)^{\Z/2}
\cong H^*(S^{n+1};\K)^{\Z/2}=\K,
\]
故每个 $\widetilde\Delta_k^\#$ 都是例外对象。此外，如\stref{笔记}{18.10}中所述，同一引理的 (ii) 及\eref{18.23}说明，$j>k$ 时，$HF^*(\widetilde\Delta_j^\#,\widetilde\Delta_k^\#)$ 上的 $\Z/2$ 作用为 $-\Id$。因此 $\F(\pi)$ 中的相应态射群无上同调，说明这些对象确实构成例外集。\stref{引理}{5.20}于是给出从 $\F^\to(\widetilde{\symbfit{\Delta}}^{\#})$ 到 $\F(\pi)$ 的在上同调上全忠实函子。另一方面，由\eref{18.23}，该有向子范畴与 $\F^\to(\symbfit V^\#)$ 相同。

\SourcePage{291}
在一般情形中，使用 Lefschetz 纤维化具有一般平行移动这一假设。它说明，对任意 $\symbfit{\gamma}$，都有一个与之同伦的 $\symbfit{\gamma}'$，使上述论证适用。每个消失圈 $V'_k$ 都同痕于 $V_k$，故它们成为 $H^0\F(M)$ 的同构对象。由\stref{引理}{5.21}，所得两个有向 $\Ai$ 子范畴拟同构。同样，每个 $\widetilde\Delta'_k$ 都通过 $\Z/2$-不变 Lagrangian 子流形同痕于 $\widetilde\Delta_k$，因而 $H^0\F(\pi)$ 中的相应对象同构。这就得到相应有向子范畴之间，以及全 $\Ai$ 子范畴之间的拟同构。
\end{proof}

\subsection{(U) 型膜与整体单值变换}\label{sec:18g}
与前面一样，令 $\symbfit{\gamma}$ 为消失路径的特选基。给每个 $\gamma_k$ 配上 $\widetilde E$ 中的 (B) 型等变膜 $\widetilde\Delta_k^\#$，其中分次任意选取。此外，设 $L^\#$ 为包含在 $E^{\mathrm{shrink}}$ 中的 $E$ 中正合 Lagrangian 膜，$\widetilde L^\#=\widetilde L_+^\#\cup\widetilde L_-^\#$ 为相应的 (U) 型等变膜。

\StatementNumber{15}\begin{lemma}\label{st:18.15}
在导出范畴 $D(\widetilde{\A})=H^0(\Tw\widetilde{\A})$ 中，有
\begin{equation}\label{eq:18.25}
T_{\widetilde\Delta_1^\#}\cdots T_{\widetilde\Delta_m^\#}(\widetilde L_+^\#)\cong S(\widetilde L_-^\#).
\end{equation}
\end{lemma}
\begin{proof}
若先忽略膜结构和平移 $S$，可作如下简化解释。由\stref{推论}{17.17}，\eref{18.25}左端的对象同构于 Lagrangian 子流形 $\tau_{\widetilde\Delta_1}\cdots\tau_{\widetilde\Delta_m}(\widetilde L_+)$ 所给出的对象。\stref{引理}{18.2}说明，在带标架同痕意义下，$\widetilde\Delta_k$ 是将 $\gamma_k$ 看成 $\Psi:\Efamily\to\mathcal W$ 的消失路径时所对应的消失圈，因此它们的 Dehn 扭转是 $\Psi$ 的单值变换。\znote{284}{这里藏着一个需要点破的等式：为什么若干 Dehn 扭转的复合会同痕于沿圆周的单值变换？因为按 Picard--Lefschetz 公式（§16d），沿 $\lambda$ 的单值变换同痕于以各消失圈为核的 Dehn 扭转、按临界值沿路径出现的先后排成的乘积；而引理 18.2 保证 $\widetilde\Delta_k$ 确实是 $\Psi$ 的消失圈，故这些扭转就是 $\Psi$ 自身的扭转。次序（先 $\widetilde\Delta_1$ 后 $\widetilde\Delta_2$……）与沿 $\lambda$ 遇到临界值的先后一致，不能交换；换成逆序得到的是逆单值变换，这也解释了\eref{18.25}左端乘积为何按 $1,\ldots,m$ 排列。[来源：原书 §16d、§18g（印刷页 284）]}于是，全部 $\tau_{\widetilde\Delta_k}$ 的复合在 $\operatorname{Aut}(\widetilde E)$ 中同痕于 $\Psi$ 沿 $\mathcal W$ 中半径 $1-3\delta$ 的圆的单值变换。由\stref{引理}{18.1}，该单值变换在子集 $\widetilde E^{\mathrm{shrink}}$ 上的限制等于 $\iota_{\widetilde E}$；按定义，该子集包含 $\widetilde L_\pm$。因此，它将 $\widetilde L_+$ 映到 $\widetilde L_-$。

为得到完整陈述，需作更仔细的论证。$T_{\widetilde\Delta_k^\#}$ 在对象上的作用对应于带典范分次的 $\tau_{\widetilde\Delta_k}$ 的作用；由 Picard--Lefschetz 公式的分次版本，这也是自然赋予 $\Psi$ 的相应单值变换的分次。由\stref{引理}{18.6}，$\widetilde L_+^\#=(\widetilde L_+,\widetilde\alpha_+^\#,P_+^\#)$ 在此单值变换下的像是
\[
\bigl(\iota_{\widetilde E}(\widetilde L_+),\,
\iota_{\widetilde E}^*\widetilde\alpha_+^\#-1,\,
\iota_{\widetilde E}^*P_+^\#\bigr).
\]
它与 $\widetilde L_-^\#$ 相差分次中的常数 $-1$，以及 Pin 结构上的 $\otimes\lambda^{\mathrm{top}}(T\widetilde L_-)$。这一差异恰好就是\secref{12a}中定义的 Lagrangian 膜平移运算 $S$ 的作用；在\secref{12h}中，我们已验证它确实拟表示抽象的代数平移。
\end{proof}

\StatementNumber{16}\begin{note}\label{st:18.16}
若将 $S$ 看成使 Lagrangian 子流形的定向反向的运算，\eref{18.25}中出现 $S$ 似乎令人困惑。事实上，
\SourcePage{292}
相反的定向已经包含在 $\widetilde L_\pm^\#$ 的定义中，见\stref{笔记}{18.8}；平移仅仅恢复一个共同的选择。
\end{note}

\StatementNumber{17}\begin{proposition}\label{st:18.17}
$\widetilde L^\#$ 位于 $\Tw\F(\pi)$ 中由 $\widetilde\Delta_1^\#,\ldots,\widetilde\Delta_m^\#$ 生成的三角子范畴中。
\end{proposition}
\begin{proof}
这是\stref{引理}{6.3}的直接应用。$\widetilde L^\#$ 是等变膜，因而是 $\widetilde{\A}$ 中在对合 $\eps_{\widetilde{\A}}$ 下不变的对象。按定义，它同构于 $\widetilde L_+^\#$ 与 $\widetilde L_-^\#=\eps_{\widetilde{\A}}(\widetilde L_+^\#)$ 的直和。此外，$HF^*(\widetilde L_+^\#,\widetilde L_-^\#)=0$，因为底层 Lagrangian 子流形不相交。各 $\widetilde\Delta_k^\#$ 同样是 $\eps_{\widetilde{\A}}$-不变对象，并且由\eref{18.23}，对所有 $k<l$，$HF^*(\widetilde\Delta_k^\#,\widetilde\Delta_l^\#)$ 都是 $\Z/2$-不变的。\stref{引理}{18.15}恰好给出所需的关系\eref{6.2}，而按定义，$\widetilde{\A}^{\mathrm{inv}}=\F(\pi)$。
\end{proof}

\clearpage
\subsection{匹配圈}\label{sec:18h}
现在更仔细地考察上述情形的一个重要特例，即 $L=\Sigma_{\mu,I}$ 是 $S$ 中某条匹配路径 $\mu$ 所对应的匹配圈，而 $\mu$ 相对于各 $\gamma_k$ 的位置如\sr{fig:18.1}{图~18.1} 所示。给 $\widetilde\Delta_k,\widetilde\Delta_{k+1}$ 赋予等变膜结构，使 $V_k,V_{k+1}$ 上诱导的膜结构彼此同痕；这只是对\stref{引理}{18.7}中固定分次方式的重述。类似地，选择 $L$ 的一个分次，并给 $\widetilde L$ 赋予相应的等变膜结构。

\StatementNumber{18}\begin{lemma}\label{st:18.18}
$HF^*(\widetilde\Delta_k^\#,\widetilde\Delta_{k+1}^\#)\cong H^*(S^n;\K)$，其上的 $\Z/2$ 作用平凡。
\end{lemma}
\begin{proof}
这只是\eref{18.23}的特例：$HF^*(\widetilde\Delta_k^\#,\widetilde\Delta_{k+1}^\#)\cong HF^*(V_k^\#,V_{k+1}^\#)$，且 $\Z/2$ 作用平凡；后一群同构于 $H^*(V_k;\K)$，因为两个膜彼此同痕。
\end{proof}

\StatementNumber{19}\begin{lemma}\label{st:18.19}
群 $HF^*(\widetilde L_\pm^\#,\widetilde\Delta_k^\#)$ 都是一维的，且集中于同一个次数。
\end{lemma}
\begin{proof}
仔细考察匹配圈的构造，就应当能看清这一点。利用\stref{引理}{15.3}，可找到一个沿各纤维恒定的变形，从给定的纤维化结构出发，最终到达一个可用常值同痕 $I$ 按朴素方式定义匹配圈的 Lefschetz 纤维化。可要求该变形集中在 $\mu$ 本身的原像附近，因此它提升为 $\widetilde E$ 的同类变形。在变形后的结构中，所考察的 Lagrangian 子流形的对应物，记为 $\widetilde L_\pm^*$ 与 $\widetilde\Delta_k^*$，横截相交于唯一一点；该点是 $E_{\gamma_k(1)}$ 中临界点在 $\widetilde E$ 中的一个原像。由\stref{引理}{7.1}，变形后的二重分歧覆盖带有到原来 $\widetilde E$ 的共形
\SourcePage{293}
辛嵌入；在同痕意义下，它将 $\widetilde L_\pm^*$ 映到 $\widetilde L_\pm$，将 $\widetilde\Delta_k^*$ 映到 $\widetilde\Delta_k$。第一种情形显然成立，因为这类同痕本来就是匹配圈构造的一部分；第二种情形则可由考察整个变形过程中连续产生的消失圈族看出。因此，可对所考察的 Lagrangian 子流形作同痕，使它们横截相交于唯一一点，这当然说明 Floer 上同调是一维的。为说明非平凡部分位于同一次数，只需使用 $\Z/2$ 作用 $\eps$；按定义，它交换这两个 Floer 群。
\end{proof}

为简便，假设 $L$ 的分次已选成使群 $HF^*(\widetilde L_\pm^\#,\widetilde\Delta_k^\#)$ 集中于次数 $1$。

\StatementNumber{20}\begin{lemma}\label{st:18.20}
在 $H(\widetilde{\A})$ 中，有如下形式的正合三角：
\SourceEquationDiagram{diagram-18-26}{0.238845897}{18.26}
\end{lemma}
\begin{proof}
由\stref{引理}{18.7}及\stref{推论}{17.17}，在 $D(\widetilde{\A})$ 中有
\begin{equation}\label{eq:18.27}
\widetilde\Delta_{k+1}^\#\cong T_{\widetilde L_-^\#}T_{\widetilde L_+^\#}(\widetilde\Delta_k^\#).
\end{equation}
此外，$\widetilde L_\pm$ 不相交，故两者之间没有态射，从而
\[
T_{\widetilde L_-^\#}T_{\widetilde L_+^\#}(\widetilde\Delta_k^\#)
\cong T_{\widetilde L_-^\#,\widetilde L_+^\#}(\widetilde\Delta_k^\#),
\]
其中右端的定义见\stref{笔记}{5.1}、\stref{笔记}{5.2}。于是，\eref{18.27}导出 $H(\Tw\widetilde{\A})$ 中如下形式的正合三角：
\SourceDiagram{diagram-18-h-triangle}{0.616845490}
由\stref{引理}{18.19}，左上顶点是 $S^{-1}(\widetilde L_-^\#)\oplus S^{-1}(\widetilde L_+^\#)\cong S^{-1}(\widetilde L^\#)$。现在旋转这个三角即可，见\stref{推论}{3.17}。\znote{285}{$S^{-1}$ 是 §12a 的平移函子（分次减 $1$，并张量 $\lambda^{\mathrm{top}}$ 的逆），它出现是因为引理 18.19 把两个群都定在次数 $1$，而三角的左顶点必须写成对象的直和。图中两条边上的 $[1]$ 同样是分次平移，不是同调次数加一：它记录态射 $\widetilde L^\#\to\widetilde\Delta_k^\#$ 的次数。另注意第三项取 $\widetilde\Delta_{k+1}^\#$ 而非其平移，旋转三角（推论 3.17）之后，方向才与引理所述形式一致；这类「平移指数是否要跟着改」的记账，是本节最易出错之处。[来源：原书 §12a、引理 18.19、推论 3.17（印刷页 285）]}
\end{proof}

\StatementNumber{21}\begin{proposition}\label{st:18.21}
在 $D\F(\pi)$ 中，$\widetilde L^\#$ 同构于 $\Cone(\widetilde\Delta_k^\#\to\widetilde\Delta_{k+1}^\#)$，其中锥可在任意非平凡的零次态射上构造。
\end{proposition}
\begin{proof}
\SourcePage{294}
由\eref{18.26}，在 $D(\widetilde{\A})$ 中得到同构 $\widetilde L^\#\cong\Cone(\widetilde\Delta_k^\#\to\widetilde\Delta_{k+1}^\#)$。在同构意义下，将任意态射乘以 $\K^*$ 中的某个元素不会改变其锥。在这里，态射群是一维的，因此只需证明所用的态射非零。若它为零，两个直和项 $\widetilde L_\pm^\#$ 就必分别同构于 $\widetilde\Delta_k^\#,\widetilde\Delta_{k+1}^\#$，但这立即与前面的 Floer 上同调计算矛盾，例如\stref{引理}{18.19}。最后，由于 $\widetilde L^\#$ 是自由的，可应用\stref{引理}{6.2}，得到 $D(\widetilde{\A}^{\mathrm{inv}})=D\F(\pi)$ 中的同构。
\end{proof}

\StatementNumber{22}\begin{note}\label{st:18.22}
另一种方法是将此情形严格视为\stref{命题}{18.17}的特例。为简便，设 $k=1$。此时需考察适当的局部模型，证明乘积
\[
HF^*(\widetilde\Delta_2^\#,\widetilde L_+^\#)\otimes HF^n(\widetilde\Delta_1^\#,\widetilde\Delta_2^\#)
\longrightarrow HF^{*+n}(\widetilde\Delta_1^\#,\widetilde L_+^\#)
\]
是同构；所涉及的各空间实际都是一维的。由此，再结合 $j>2$ 时 $HF^*(\widetilde\Delta_j^\#,\widetilde L_+^\#)=0$，可知\stref{引理}{6.3}证明中的对象 $U$，在这一特殊情形中归约为 $\Cone(\widetilde\Delta_1^\#\to\widetilde\Delta_2^\#)$。
\end{note}

\subsection{Hurwitz 移动}\label{sec:18i}
设 $\gamma_k,\gamma_{k+1},\delta_k$ 为\sr{fig:16.2}{图~16.2} 所示的消失路径，$\widetilde\Delta_k^\#,\widetilde\Delta_{k+1}^\#,\widetilde\Pi_k^\#$ 为相应的 (B) 型膜。更准确地说，$\widetilde\Delta_k^\#$ 的分次任意，但其余两个分次应作调整，使它们与 Dehn 扭转的典范分次一起，与\stref{引理}{18.3}中的同痕相容。重述该结果，现在有膜的同痕
\begin{equation}\label{eq:18.28}
\widetilde\Pi_k^\#\simeq\tau_{\widetilde\Delta_k}(\widetilde\Delta_{k+1}^\#).
\end{equation}
如\stref{笔记}{18.5}所述，这一同痕实际经过 $\iota_{\widetilde E}$-不变 Lagrangian 子流形。为使它成为等变膜的同痕，还需比较等变 Pin 结构。由\stref{引理}{14.2}，选择一个等变 Pin 结构，对应于选择不动点集的余定向。我们的约定是，为 $\widetilde\Delta_{k+1}$ 与 $\widetilde\Pi_k$ 选择的余定向，在 $D\widetilde\pi$ 下投影到 $\gamma'_{k+1}(0)$ 和 $\delta'_k(0)$ 的虚部为正的平方根。\znote{286}{这与 §18e 对 (B) 型膜的规定是同一条：等变 Pin 结构对应不动点集的余定向（引理 14.2），而对 $\widetilde\Delta_\gamma$ 有 $D\widetilde\pi(T\widetilde\Delta_\gamma|V_\gamma)=\R\sqrt{\gamma'(0)}$，故「取虚部为正的平方根」唯一确定余定向。要点在于两个球面的余定向分别用 $\gamma'_{k+1}$、$\delta'_k$ 定义，比较之前必须借\eref{18.28}的同痕把二者对齐；笔记 18.5 中相交数 $\widetilde\delta_k\cdot\widetilde\gamma_{k+1}=-1$ 的计算，正是说明这两个选择相容而非相差一个符号。[来源：原书 §18e、§18h、笔记 18.5（印刷页 279、282、286）]}\stref{笔记}{18.5}中的计算说明，这些选择确实与\eref{18.28}相容。现在可应用\secref{17k}的讨论，特别是\stref{引理}{17.23}。将由此得到的代数数据代入\stref{引理}{6.4}，便得到：

\StatementNumber{23}\begin{proposition}\label{st:18.23}
在 $D\F(\pi)$ 中，$\widetilde\Pi_k^\#$ 同构于 $T_{\widetilde\Delta_k^\#}\widetilde\Delta_{k+1}^\#$。
\end{proposition}

\subsection{最初的推论}\label{sec:18j}
\stref{命题}{18.23}最直接的应用，是将 Hurwitz 移动与变异联系起来。取消失路径的一个特选基
\SourcePage{295}
$\symbfit{\gamma}=(\gamma_1,\ldots,\gamma_m)$，考虑相应的 (B) 型等变膜集 $\widetilde{\symbfit{\Delta}}_{\symbfit{\gamma}}^\#$，其中分次任意选取；同时考虑带诱导膜结构的相应消失圈集 $\symbfit V_{\symbfit{\gamma}}^\#$。由\stref{命题}{18.14}，已知 $\widetilde{\symbfit{\Delta}}_{\symbfit{\gamma}}^\#$ 是例外集。\stref{命题}{18.23}说明，若对特选基施加\eref{16.7}中定义的一次初等 Hurwitz 移动，所得结果就是\eref{5.24}意义下该例外集的一次初等变异。选择相关记号时，我们已隐含地承认了这一对应，因此可写成
\begin{equation}\label{eq:18.29}
\widetilde{\symbfit{\Delta}}_{\mathcal L_k\symbfit{\gamma}}^\#
\cong\mathcal L_k\widetilde{\symbfit{\Delta}}_{\symbfit{\gamma}}^\#.
\end{equation}
严格地说，为使\eref{18.29}成为 $D\F(\pi)$ 中例外集的真正同构，分次必须按\stref{命题}{18.23}的方式专门选择；否则，它只在相差平移的意义下是同构，而这已完全足够。显然，只需改变所涉及对象集的记号，同样的结论就适用于逆移动 $\mathcal L_k^{-1}$，再通过复合，适用于任意变异。换言之，辫群 $\operatorname{Br}_m$ 在模去平移与同构的例外集上的代数作用，见\secref{5j}，可与它在消失路径特选基上的几何作用，见\secref{16d}，等同起来。

固定一个 $\symbfit{\gamma}$。暂以 $\mathcal T$ 表示 $\Tw\F(\pi)$ 中由对象 $\widetilde\Delta_{\gamma_1}^\#,\ldots,\widetilde\Delta_{\gamma_m}^\#$ 生成的三角子范畴。依定义，$\widetilde{\symbfit{\Delta}}_{\symbfit{\gamma}}^\#$ 是该子范畴中的满例外集，所有变异之后也仍如此。由\eref{18.29}，这意味着，若对基施加 Hurwitz 移动，再考虑与所得基相应的对象集，则其全部对象仍在 $\mathcal T$ 中，并生成该范畴。由于 Hurwitz 移动在特选基上可迁地作用，$\mathcal T$ 与 $\symbfit{\gamma}$ 无关。事实上，\stref{命题}{18.17}说明，任意 (U) 型等变膜都属于 $\mathcal T$，因此，由 $\F(\pi)$ 的定义，必有 $\mathcal T=\Tw\F(\pi)$。另一方面，\stref{引理}{5.22}说明，$\mathcal T$ 拟等价于与 $\widetilde{\symbfit{\Delta}}_{\symbfit{\gamma}}^\#$ 相应的有向 $\Ai$ 子范畴上的扭曲复形范畴；再由\stref{命题}{18.14}，该有向子范畴拟同构于 $\F^\to(\symbfit V_{\symbfit{\gamma}}^\#)$。\znote{287}{定理的证明是四步，第一步最容易在读时被跳过：(1) 引理 18.2 与命题 18.23 给出\eref{18.29}，即初等 Hurwitz 移动对应初等变异，于是「满例外集的变异」与「基的 Hurwitz 移动」同步；(2) Hurwitz 移动在特选基集合上可迁地作用，故 $\mathcal T$ 与所取基无关；(3) 命题 18.17 把任意 (U) 型膜收进 $\mathcal T$，而 $\F(\pi)$ 由 $\widetilde\Delta_{k}^\#$ 与这些膜生成，故 $\mathcal T=\Tw\F(\pi)$；(4) 引理 5.22 把 $\mathcal T$ 实现为 $\widetilde{\symbfit{\Delta}}_{\symbfit{\gamma}}^\#$ 上扭曲复形范畴，命题 18.14 再把后者换成 $\Tw\F^\to(\symbfit V_{\symbfit{\gamma}}^\#)$。注意定理本身不另列假设：它是在 §18f 起固定的 $\operatorname{char}\K\ne2$ 之下陈述的，代价在那里一次付清；四步中只有 (1) 真正用到二重分歧覆盖。[来源：原书 §18j（印刷页 287）]}因此：

\StatementNumber{24}\begin{theorem}\label{st:18.24}
$\Tw\F^\to(\symbfit V_{\symbfit{\gamma}}^\#)$ 拟等价于 $\Tw\F(\pi)$。
\end{theorem}

一个直接推论是得到\stref{推论}{17.21}的另一证明；事实上，也得到\stref{定理}{17.20}的另一证明，因为任意 Lagrangian 球面集都可实现为某个适当构造的 Lefschetz 纤维化的消失圈基。这比原来的证明有所改进：通过将 $\Tw\F^\to(\symbfit V_{\symbfit{\gamma}}^\#)$ 与显然不涉及基的选择的 $\Tw\F(\pi)$ 等同起来，证明了它与 $\symbfit{\gamma}$ 无关。一个具体收获是，可以为不同 $\symbfit{\gamma}$ 对应的范畴 $\Tw\F^\to(\symbfit V_{\symbfit{\gamma}}^\#)$ 构造典范的拟等价，至少在相应函子范畴 $H^0(\operatorname{fun})$ 中的同构意义下如此。
\SourcePage{296}
为此付出的唯一代价，是技术假设 $\operatorname{char}(\K)\ne2$；使用二重分歧覆盖这一技巧，迫使我们作出该假设。

其次，注意\stref{命题}{18.13}与\stref{引理}{3.23}给出一个在上同调上全忠实的函子 $\Tw\F(E)\to\Tw\F(\pi)$。结合\stref{定理}{18.24}，得到：
\StatementNumber{25}\begin{corollary}\label{st:18.25}
对任意 $\symbfit{\gamma}$，存在 $\F(E)$ 到 $\Tw\F^\to(\symbfit V_{\symbfit{\gamma}}^\#)$ 的在上同调上全忠实的嵌入。
\end{corollary}

关键在于，有向 Fukaya 范畴是在 $\F(M)$ 内部构造的，因此，从抽象角度说，总空间 $E$ 中的 Floer 上同调计算可以完全归约为纤维中的计算。下面尝试把这些讨论变得更加具体。\znote{288}{注意这里是「拟满忠实」而非「等价」：$\F(E)\hookrightarrow\Tw\F^\to(\symbfit V_{\symbfit{\gamma}}^\#)$ 只保证 $\Hom$ 上的同构，并不保证每个对象都来自 $E$ 中的膜，故 $E$ 的 Fukaya 范畴实现为消失圈所给有向 $\Ai$ 范畴的扭曲复形范畴的一个全子范畴。这句话的实际含义因此是：全空间的 Floer 数据被纤维内一列球面（连同它们的相交数与高阶运算）完全承载，这正是 Picard--Lefschetz 理论把整体问题化为组合数据的着力点。也正因如此，命题 18.13 的证明分两步：先说明 $E^{\mathrm{shrink}}$ 内的膜已足以算出 $\F(E)$（用定理 2.9 反演），再说明它们与 $\widetilde{\A}^{\mathrm{inv}}$ 的一个全子范畴一致。[来源：原书 §18j、命题 18.13、定理 18.24（印刷页 282、287、288）]}

\subsection{再论匹配圈}\label{sec:18k}
取\stref{定理}{18.24}中拟等价的一个逆拟等价，将其记为
\[
\G:\Tw\F(\pi)\longrightarrow\Tw\F^\to(\symbfit V_{\symbfit{\gamma}}^\#).
\]
首先注意，可以确定每个 (B) 型等变膜在 $\G$ 下映到哪里。由构造，对用于定义 $\G$ 的那个特定基，有
\[
\G(\widetilde\Delta_{\gamma_j}^\#)=V_{\gamma_j}^\#,
\]
其中右端是作为 $\F^\to(\symbfit V_{\symbfit{\gamma}}^\#)$ 本身的第 $j$ 个对象的消失圈。现考虑任意对象集 $\symbfit{\delta}$，及其第 $k$ 次初等 Hurwitz 移动 $\symbfit{\rho}=\mathcal L_k\symbfit{\delta}$。利用\stref{引理}{5.6}及\stref{命题}{18.23}，得到
\begin{equation}\label{eq:18.30}
\G(\widetilde\Delta_{\rho_j}^\#)\cong
\begin{cases}
\G(\widetilde\Delta_{\delta_j}^\#),&j\ne k,k+1,\\
\G(\widetilde\Delta_{\delta_k}^\#),&j=k+1,\\
T_{\G(\widetilde\Delta_{\delta_k}^\#)}\G(\widetilde\Delta_{\delta_{k+1}}^\#),&j=k.
\end{cases}
\end{equation}
这里的 $\cong$ 在 $D\F^\to(\symbfit V_{\symbfit{\gamma}}^\#)$ 中成立，而且实际可能只在相差平移的意义下成立，除非额外追踪正确的分次；我们不坚持这样做。\znote{289}{\eref{18.30} 最容易看反的是 $j=k$ 与 $j=k+1$ 两行的分工：初等 Hurwitz 移动把第 $k$ 条与第 $k+1$ 条路径换位，于是新基的第 $k+1$ 个对象就是原来的 $\widetilde\Delta_{\delta_k}^\#$（不加扭转），只有第 $k$ 个对象被扭转。逆移动公式中出现 $T^\vee$，方向相反，正说明 $\mathcal L_k$ 与 $\mathcal L_k^{-1}$ 互逆。另需留意扭转的下标是像对象 $\G(\widetilde\Delta_{\delta_k}^\#)$ 而不是原对象，因为变异是在目标范畴中进行的。「只到相差平移」并不削弱结论——后续比较都允许对象平移——但若把它当作严格同构去复合，分次上必出错。[来源：原书 §18k（印刷页 288、289）]}通过代入，可得到逆 Hurwitz 移动 $\symbfit{\sigma}=\mathcal L_k^{-1}\symbfit{\delta}$ 的类似公式，即
\[
\G(\widetilde\Delta_{\sigma_j}^\#)\cong
\begin{cases}
\G(\widetilde\Delta_{\delta_j}^\#),&j\ne k,k+1,\\
\G(\widetilde\Delta_{\delta_{k+1}}^\#),&j=k,\\
T^\vee_{\G(\widetilde\Delta_{\delta_{k+1}}^\#)}\bigl(\G(\widetilde\Delta_{\delta_k}^\#)\bigr),&j=k+1.
\end{cases}
\]
每条消失路径都是某个特选基的一部分，这就完成了第一步的说明。

\SourcePage{297}
其次，设 $\mu$ 为匹配路径，它相对于 $\symbfit{\delta}$ 的位置如\sr{fig:18.3}{图~18.3} 所示。由\stref{引理}{18.18}，$D\F(\pi)$ 中从 $\widetilde\Delta_{\delta_k}^\#$ 到 $\widetilde\Delta_{\delta_{k+1}}^\#$ 的态射空间有两个生成元，次数相差 $n$。准确的绝对次数
\SourceFigure{figure-18-3}{0.191076717}{18.3}
取决于分次的选择；我们平移这些分次，直到两个次数成为 $0$ 与 $n$。取匹配圈 $\Sigma_\mu$ 及相应的 (U) 型等变膜 $\widetilde\Sigma_\mu^\#$。由\stref{推论}{3.9}与\stref{命题}{18.21}，有
\[
\G(\widetilde\Sigma_\mu^\#)
\cong\G\bigl(\Cone(\widetilde\Delta_{\delta_k}^\#\to\widetilde\Delta_{\delta_{k+1}}^\#)\bigr)
\cong\Cone\bigl(\G(\widetilde\Delta_{\delta_k}^\#)\to\G(\widetilde\Delta_{\delta_{k+1}}^\#)\bigr),
\]
其中锥在任意非平凡零次态射上构造，同构的意义与\eref{18.30}相同。注意，对任意给定匹配路径 $\mu$，都能找到使这一做法适用的特选基 $\symbfit{\delta}$，这意味着全部匹配圈在 $\G$ 下的像都可显式确定。应用\stref{命题}{18.13}，得到：

\StatementNumber{26}\begin{corollary}\label{st:18.26}
设 $\mu_1,\ldots,\mu_l$ 为 $\pi$ 的有限条匹配路径。那么，只需在拟同构意义下知道一个特选基 $\symbfit{\gamma}$ 对应的 $\F^\to(\symbfit V_{\symbfit{\gamma}}^\#)$，就能计算 $\F(E)$ 中由匹配圈 $\Sigma_{\mu_j}^\#$ 构成的全 $\Ai$ 子范畴，计算结果确定到拟同构以及可能的对象平移。
\end{corollary}

\clearpage
\subsection{一个谱序列}\label{sec:18l}
设 $\symbfit{\gamma}$ 为消失路径的特选基，$\symbfit{\delta}$ 为对它施加与\eref{5.25}中的元素 $\Delta^{1/2}\in\operatorname{Br}_m$ 相应的 Hurwitz 移动后得到的基。注意，虽然这是辫群的正元素，但按\secref{16d}的讨论，在几何上，相应圆盘微分同胚是沿一条平行于边界的圆周所作的负 Dehn 扭转的平方根。\znote{289}{这里容易读反：$\Delta^{1/2}\in\operatorname{Br}_m$ 在辫群中是正元素，而它对应的圆盘微分同胚却是「负」Dehn 扭转的平方根。原因在于辫群与映射类的关系是按几何作用（对消失路径的 Hurwitz 移动）定义的（§16d），其符号约定与「正扭转」的字面直觉相反：整圈 $\Delta$ 对应沿边界平行圆的负 Dehn 扭转，故其一半才记为 $\Delta^{1/2}$。这一支在本书中的作用是：$\symbfit{\delta}$ 由 $\symbfit{\gamma}$ 经这一整族移动得到，于是可借\eref{18.30}逐次算出 $\G$ 在新基上的取值，最后再把基点移到 $-*$ 得到 $\symbfit{\gamma}^!$。[来源：原书 §16d、§18l（印刷页 289、290）]}此外，还要考虑将基点从 $*$ 沿顺时针半圆移到 $-*$，由 $\symbfit{\delta}$ 得到的路径集 $\symbfit{\gamma}^!$；见\sr{fig:18.4}{图~18.4}。

\SourcePage{298}
\SourceFigure{figure-18-4}{0.328153493}{18.4}
设 $L_j^\#$，$j=0,1$，为包含在 $E^{\mathrm{shrink}}$ 中的两个膜，$\widetilde L_j^\#$ 为相应的 (U) 型等变膜。将它们视为 $D\F(\pi)$ 的对象，对其应用\stref{命题}{5.17}。所得谱序列收敛到
\[
HF^*(\widetilde L_0^\#,\widetilde L_1^\#)^{\Z/2}\cong HF^*(L_0^\#,L_1^\#),
\]
其中的同构是\stref{命题}{18.13}的一部分。谱序列的 $E_1$ 页由两类 Floer 群的张量积构成。其中一类是
\begin{equation}\label{eq:18.31}
HF^*(\widetilde\Delta_k^\#,\widetilde L_1^\#)^{\Z/2}\cong HF^*(\Delta_k^\#,L_1^\#),
\end{equation}
其中 $\Delta_k$ 为 $\gamma_k$ 对应的 Lefschetz 顶针，$\widetilde\Delta_k$ 为其在 $\widetilde E$ 中的原像。照常将后者赋予 (B) 型等变膜结构，分次任意选取，使\eref{18.31}左端成为 $H\F(\pi)$ 中的态射空间；再应用\stref{笔记}{18.11}，将其等同于 $E$ 本身中适当的 Floer 上同调群，即右端。出现的另一类 Floer 群可类似地写成
\[
HF^*(\widetilde L_0^\#,\widetilde\Pi_k^\#)^{\Z/2}
\cong HF^*(L_0^\#,\Pi_k^\#)
\cong HF^*(L_0^\#,\Delta_k^{!,\#}),
\]
其中 $\Pi_k,\Delta_k^!$ 为 $\delta_k,\gamma_k^!$ 对应的 Lefschetz 顶针。第二个同构由这些顶针之间显然的同痕给出；该同痕移动边界，但仍保持 Floer 上同调群，因为 Floer 方程的相关解从不到达边界。准确地说，应为 $\Delta_k^!$ 使用与变异描述相容的分次：它可由 $\Delta_{m+1-k}\cap\Delta_k^!$ 的唯一交点具有零 Maslov 指标这一性质刻画，因为这是\secref{5k}中 Koszul 对偶集的一般性质之一。据此，有：

\SourcePage{299}
\StatementNumber{27}\begin{corollary}\label{st:18.27}
存在收敛到 $HF^*(L_0^\#,L_1^\#)$ 的谱序列，其起始页为
\begin{equation}\label{eq:18.32}
E_1^{rs}=\bigl(HF^*(L_0^\#,\Delta_{r+1}^{!,\#})\otimes HF^*(\Delta_{m-r}^\#,L_1^\#)\bigr)^{r+s}.
\end{equation}
\end{corollary}

即使 $L_0^\#=L_1^\#=L^\#$，这一结果也很有意义，此时谱序列的目标化为 $H^*(L;\K)$。注意，\stref{命题}{5.17}也描述了\eref{18.32}上第一个微分的形式：它由两个对象集中相继对象之间互相对偶的态射群中的元素作乘法给出。此外，由\stref{命题}{18.14}，这些群可用消失圈写成 $HF^*(V_k^\#,V_{k+1}^\#)$ 与 $HF^*(V_{m-k}^{!,\#},V_{m+1-k}^{!,\#})$。\znote{291}{关于\eref{18.32}的记号：$r,s$ 是双分次，外层上标 $r+s$ 不是矩阵元的下标，而是取该张量积的总次数 $r+s$ 分量（$HF^*$ 自带 $\Z$-分次），把齐次分量写成上标是全书惯用写法。两个因子的指标之和恒为 $m+1$：$r+1$ 与 $m-r$ 恰好配对，核对下标时很有用。又注意 $\Delta^{!,\#}$ 上标的 $!$ 标记与变异描述配套的那族路径 $\gamma_k^!$（见 §5k 的 Koszul 对偶集），不是求导，也不是取对偶空间。[来源：原书 §18l（印刷页 291）]}

\StatementNumber{28}\begin{note}\label{st:18.28}
最初是 Donaldson 向作者提出了谱序列\eref{18.32}应当存在。他的观点基于对角线的几何分解：为简便，假设 $\pi$ 的任意两个临界值的虚部不同。取对角线 $\Delta\subset E\times E$，其中第一个因子的辛形式变号；在第一个因子上施加 $\im(\pi)$ 的 Hamilton 流，使该对角线变形。对角线的一部分会流出 $E$，但其余部分会逼近稳定流形与不稳定流形的乘积，比较如\bcite{73}。这些流形是沿水平方向走向 $\pm\infty$ 的路径对应的 Lefschetz 顶针，应将其看成\sr{fig:18.4}{图~18.4} 中路径 $\gamma_k,\gamma_k^!$ 的变形版本。在我们的证明中，这一思想被 Beilinson 意义下对角线的代数分解所取代。
\end{note}

\section{代数簇}\label{sec:19}\label{section:19}
\subsection{多重次调和函数}\label{sec:19a}
Stein 流形是一个复流形 $(X,I_X)$，它容许一个穷竭的，也就是固有且有下界的多重次调和函数 $\phi:X\to\R$。多重次调和性是说，令 $\theta_\phi=-d^c\phi=-d\phi\circ I_X$ 时，$\omega_\phi=d\theta_\phi=-dd^c\phi$ 是 Kähler 形式。注意，$dd^c\phi$ 总是 $(1,1)$ 型的；非平凡的条件是，它在复切直线上为负，复切直线指 $TX$ 中带自然定向的实 $I_X$-不变二维子空间。一个直接推论是，给定任意全纯映射 $u:S\to X$，函数 $\rho=\phi\circ u$ 如\eref{7.4}那样满足 $\Delta\rho\ge0$，因此服从最大值原理。

由于多重次调和性条件具有凸性，比较容易对这类函数作调整。下面需要其中两个初等实例；有关这一方向的更多内容，见\bcite{42}或\bcite{22}。

\SourcePage{300}
\StatementNumber{1}\begin{lemma}\label{st:19.1}
设 $\phi$ 为穷竭的多重次调和函数，$h:\R\to\R$ 处处满足 $h'>0$、$h''\ge0$。则 $h(\phi)$ 仍为穷竭的多重次调和函数。
\end{lemma}
\begin{proof}
有 $dd^ch(\phi)=h'(\phi)dd^c\phi+h''(\phi)d\phi\wedge d^c\phi$。在复切直线上，第一个加项为负，第二个非正。\znote{292}{本书的 $d^c$ 不带 $\tfrac12$：由 $\theta_\phi=-d^c\phi=-d\phi\circ I_X$ 得 $d^c\phi=d\phi\circ I_X$，而不少文献取 $d^c=\tfrac i2(\bar\partial-\partial)$，换算 $\omega_\phi=-dd^c\phi$ 时会差一个因子 $2$。第二个加项非正可直接验证：在复切直线上取基 $v,Iv$，则 $d\phi\wedge d^c\phi(v,Iv)=-(d\phi(v)^2+d\phi(Iv)^2)\le0$，而 $h''\ge0$ 又使它无法抵消第一项的负性；若 $h$ 为仿射函数则第二项整个消失，结论同样成立（引理 19.2 证明中的 $l$ 在 $t\ll0$ 与 $t\gg0$ 两段就分别取仿射形式）。[来源：A，原书 p291 对 $\theta_\phi$、$d^c$ 的定义与引理 19.1、19.2 的证明。]}
\end{proof}

\StatementNumber{2}\begin{lemma}\label{st:19.2}
设 $\phi,\psi$ 为 $(X,I_X)$ 上穷竭的多重次调和函数。假设 $\psi$ 比 $\phi$ 增长得快，意思是 $\psi-\phi$ 仍为穷竭函数。固定 $X$ 的某个紧子集。那么，存在另一个穷竭的多重次调和函数 $\widetilde\psi$，它在该紧子集上等于 $\phi$，在无穷远处等于 $\psi+\text{常数}$。
\end{lemma}
\begin{proof}
选择函数 $l$，使 $t\ll0$ 时 $l(t)=0$，$t\gg0$ 时 $l(t)=t$，且处处 $l''(t)\ge0$。对某个常数 $b$，取 $\xi=\phi+l(\psi-\phi-b)$。由与前面相同的理由，它是多重次调和的：
\[
\begin{aligned}
dd^c\xi={}&\bigl(1-l'(\psi-\phi-b)\bigr)dd^c\phi
+l'(\psi-\phi-b)dd^c\psi\\
&+l''(\psi-\phi-b)d(\psi-\phi)\wedge d^c(\psi-\phi).
\end{aligned}
\]
将常数 $b$ 取充分大，即可保证在选定的紧子集上 $l(\psi-\phi-b)=0$。
\end{proof}

若 $\phi$ 的临界点集紧，就称它为\emph{有限型}。一般来说，判断一个 Stein 流形是否容许这样的函数可能很困难；对于四维流形，微分结构会带来微妙的障碍\bcite{63}。不过，在许多大类情形中，这很容易做到。

\StatementNumber{3}\begin{example}\label{st:19.3}
(i) 设 $\overline X$ 为光滑复射影簇，带有丰沛线丛 $\mathcal L$。丰沛性意味着该线丛上有度量 $|\cdot|$，其曲率是 $\overline X$ 上的 Kähler 形式。设 $s$ 为 $\mathcal L$ 的全纯截面，使 $s^{-1}(0)$ 是正规交叉除子。那么
\[
\phi(x)=-\log|s(x)|
\]
是 $X=\overline X\setminus s^{-1}(0)$ 上有限型的穷竭严格多重次调和函数。相应 Kähler 形式 $\omega_\phi$ 是前述 Kähler 形式在 $X$ 上的限制。有限型性质可由直接的局部计算证明，见\bcite{133}，引理~8。

(ii) 设 $X\subset\CC^N$ 为光滑仿射代数簇，$\phi:X\to\R$ 为穷竭的多重次调和函数，而且是一个实多项式的限制。那么 $\phi$ 自动是有限型的，因为其临界点集本身是实代数簇，因此只有有限多个连通分支，每个分支对应于唯一的临界值。

(iii) 类似地，设 $q:\CC^N\to\CC$ 为复多项式，原点是它的孤立临界点，且 $q(0)=0$、$Dq(0)=0$。那么，对
\SourcePage{301}
所有 $0<|\epsilon|\ll\delta$，Milnor 纤维 $X=\{x\in\CC^N:q(x)=\epsilon,\ |x|<\delta\}$ 是 Stein 流形，$\phi(x)=(\delta^2-|x|^2)^{-1}$ 是其上的有限型多重次调和函数。注意，Milnor 纤维一般不双全纯同构于代数簇。
\end{example}

此后固定 $(X,I_X)$，并要求它带有二次复体积形式 $\eta_X^2$ 及有限型的穷竭多重次调和函数 $\phi$。取 $c$ 大于 $\phi$ 的全部临界值，考虑次水平集
\[
M=X^{\phi,c}=\phi^{-1}(({-\infty},c])
\]
及 $\theta_X,\omega_X$ 的限制。这是带接触型边界的正合辛流形，$I_X$ 的限制也是边界附近的接触型殆复结构。再使用 $\eta_X^2$ 的限制，就具备了定义 Fukaya 范畴 $\F(M)$ 所需的全部数据。

\StatementNumber{4}\begin{proposition}\label{st:19.4}
在拟等价意义下，$\F(M)$ 与 $c$ 及 $\phi$ 无关。
\end{proposition}
\begin{proof}
整个讨论与\secref{10d}中的讨论有些类似。第一步说明，\emph{固定 $\phi$ 时，$c$ 的选择无关紧要}。另取 $d>c$。定义 $\F(X^{\phi,d})$ 时，对于包含在较小子集 $M=X^{\phi,c}$ 中的正合 Lagrangian 子流形，可以使 Floer 数据与扰动数据只使用在 $X^{\phi,d}\setminus\operatorname{int}(X^{\phi,c})$ 的某个邻域上等于 $I_X$ 的殆复结构，以及支撑于 $\operatorname{int}(X^{\phi,c})$ 中的 Hamilton 函数，即非齐次项。最大值原理说明，非齐次伪全纯多边形留在 $X^{\phi,c}$ 内，因此 $\F(X^{\phi,c})$ 作为全 $\Ai$ 子范畴包含在 $\F(X^{\phi,d})$ 中。另一方面，$X^{\phi,d}$ 中每个正合 Lagrangian 子流形，都能由负 Liouville 流推入 $X^{\phi,c}$；在 Stein 情形中，该流就是 $\phi$ 的负梯度流。这说明包含 $\F(X^{\phi,c})\to\F(X^{\phi,d})$ 是拟等价。

第二，\emph{若 $\phi=\psi$ 在某个紧子集之外成立，则相应 Fukaya 范畴相同}。可在两个函数之间插值，再应用标准的 Moser 型论证，说明充分大的次水平集辛同构，而且该辛同构正合，并在边界附近为恒等。

第三，\emph{若如\stref{引理}{19.1}那样，从 $\phi$ 换成某个 $h(\phi)$，则 Fukaya 范畴保持不变}。可选择具有同样性质的辅助函数 $\widetilde h$，使它在 $r\gg0$ 时与 $h$ 一致，且 $r\le c$ 时 $\widetilde h(r)=ar+b$，其中 $a>0$、$b\in\R$ 为常数。从 $\phi$ 换成 $\widetilde h(\phi)$，只使次水平集 $X^{\phi,c}$ 的辛几何发生常数倍缩放，这很容易通过调整 Floer 数据与扰动数据加以处理。由于两个函数在该次水平集之外都没有临界点，相应 Fukaya 范畴相同。另一方面，从 $\widetilde h(\phi)$ 换成 $h(\phi)$ 是紧支撑的改变，因此上面的论证适用。

\SourcePage{302}
作了这些准备后，进入主要论证。给定有限型的穷竭多重次调和函数 $\phi,\psi$。由\stref{引理}{19.1}，找到比 $\psi$ 增长得快的函数 $h(\phi)$。取 $d$ 大于 $\psi$ 的全部临界值。\stref{引理}{19.2}给出 $\widetilde\phi$，它在无穷远处等于 $h(\phi)+\text{常数}$，在 $X^{\psi,d}$ 上与 $\psi$ 一致。再重复同样的论证：取 $c$ 大于 $\widetilde\phi$ 的全部临界值，且使 $X^{\psi,d}\subset\operatorname{int}(X^{\widetilde\phi,c})$。再次由\stref{引理}{19.1}，找到 $k(\psi)$，它在 $X^{\psi,d}$ 上等于 $\psi$，并比 $\widetilde\phi$ 增长得快。再应用\stref{引理}{19.2}，找到另一个函数 $\widetilde\psi$，使它在无穷远处等于 $k(\psi)+\text{常数}$，在 $X^{\widetilde\phi,c}$ 上等于 $\widetilde\phi$。

由上述讨论，$\widetilde\phi$ 相应的 Fukaya 范畴与 $\phi$ 的相同。$\widetilde\psi$ 与 $\psi$ 也如此，但在后一情形中还可作进一步观察。考虑在两个函数之间线性插值并应用 Moser 技巧所得的 $(X,\omega_{\widetilde\psi},\theta_{\widetilde\psi})$ 与 $(X,\omega_{k(\psi)},\theta_{k(\psi)})$ 之间的辛同构。该同构在 $X^{\psi,d}$ 上为恒等，因为在那里 $\widetilde\psi,\widetilde\phi,\psi,k(\psi)$ 仅相差常数。$k(\psi)$ 的负 Liouville 流将 $(X,\omega_{k(\psi)},\theta_{k(\psi)})$ 的任意紧子集推入这个次水平集。因此，应用 Moser 辛同构后，$(X,\omega_{\widetilde\psi},\theta_{\widetilde\psi})$ 的每个正合 Lagrangian 子流形都同痕于一个位于 $X^{\psi,d}$ 中的子流形。

取 $\widetilde d\gg c$。如\stref{命题}{10.5}那样，包含
\[
(X^{\widetilde\phi,c},\omega_{\widetilde\phi},\theta_{\widetilde\phi})
\hookrightarrow(X^{\widetilde\psi,\widetilde d},\omega_{\widetilde\psi},\theta_{\widetilde\psi})
\]
在 Fukaya 范畴上诱导一个在上同调上全忠实的函子。另一方面，该包含的像含有 $X^{\psi,d}$。因此，由前面的观察，该函子是拟等价。
\end{proof}

\StatementNumber{5}\begin{note}\label{st:19.5}
还有其他将 Stein 流形截取成紧子集的可能方式。例如，若 $X\subset\CC^N$ 为仿射簇，可考虑它与多球
\[
P=\bigl\{|x_1|^2+\cdots+|x_{N_1}|^2\le c_1;\ \ldots;\
|x_{N_1+\cdots+N_{r-1}+1}|^2+\cdots+|x_N|^2\le c_r\bigr\}
\]
的交 $M=X\cap P$，其中 $N_1+\cdots+N_r=N$，$c_1,\ldots,c_r$ 充分大。选一个多重次调和函数，使其负 Liouville 流相对于每个变量 $x_k$ 都指向内部，例如标准函数 $\tfrac14|x|^2$，或 Fubini--Study Kähler 势 $\log(1+|x|^2)$。相应一形式与二形式的限制使 $M$ 成为带角的正合辛流形。由\stref{命题}{10.5}，Fukaya 范畴 $\F(M)$ 与先前构造所得的相同。

另一种等价方法是放弃紧性，直接在整个 $X$ 上工作，使用无穷远处等于 $I_X$ 的殆复结构。这不完全符合本书的框架，但也许是最自然的想法，而且还能推广到无限型情形。
\end{note}

\SourcePage{303}
\subsection{Lefschetz 束}\label{sec:19b}
设 $\overline X$ 为 $(n+1)$ 维光滑射影簇，带有丰沛线丛 $\mathcal L$。取两个线性无关截面 $s_0,s_\infty\in H^0(\mathcal L)$，考虑有理函数 $\overline p=s_0/s_\infty$，其纤维为 $z\in\CP^1$ 对应的超曲面 $\{s_0(y)=zs_\infty(y)\}$。$\overline p$ 的不定点集为基点集 $B$，由 $s_0$ 与 $s_\infty$ 同时为零定义。若 $B$ 光滑，且 $\overline p$ 只有非退化临界点，每个纤维至多含一个，就称 $\overline p$ 为 \emph{Lefschetz 束}。还假设 $s_\infty^{-1}(0)$ 本身光滑，并考虑 $\overline p$ 在其补集上的限制，记为
\begin{equation}\label{eq:19.1}
p:X=\overline X\setminus s_\infty^{-1}(0)\longrightarrow\CC.
\end{equation}
在 $\mathcal L$ 上选择一个给出 Kähler 形式的度量。如\stref{例题}{19.3}(i)，$s_\infty$ 的范数的负对数定义一个有限型多重次调和函数 $\phi$。所得设置与第\sr{sec:15}{15 节}引入的正合 Lefschetz 纤维化概念密切相关，但并不完全相同。有两种处理这一差别的方式：一是说明\eref{19.1}的几何性质足够好，可顺利推广所需论证；另一种是修改给定情形，使它完全符合前面建立的形式体系。

第一种方法的要点，是说明\eref{19.1}上由 $\omega_\phi$ 给出的联络具有良好的平行移动映射，尽管纤维并不紧。为此，沿基点集作爆破，得到 $\overline p$ 的图像
\[
\widehat X=\{(z,y)\in\CP^1\times\overline X:s_0(y)=zs_\infty(y)\}.
\]
它带有投影 $\widehat p:\widehat X\to\CP^1$ 与 $\widehat X\to\overline X$。通过第二个映射拉回给定 Kähler 形式，记所得形式为 $\omega_{\widehat X}$。它沿例外除子 $\widehat B=\CP^1\times B$ 退化，因而不再是 Kähler 形式，但在光滑纤维上的限制仍是 Kähler 形式，因此定义一个辛联络。该联络实际上在 $\widehat B$ 上平凡，所以在补集 $\widehat X\setminus\widehat B$ 上，其平行移动在奇点之外有良好定义。另一方面，从 $\widehat p:\widehat X\setminus\widehat B\to\CP^1$ 中去掉无穷远处的纤维，就恢复了 $p:X\to\CC$ 及由 $\omega_\phi$ 给出的联络。

这一论证还给出更多信息，在分次背景下尤其有意义。具体地，假设对某个 $d\in\Z$，有
\begin{equation}\label{eq:19.2}
\mathcal K_{\overline X}\cong\mathcal L^{\otimes-d}.
\end{equation}
在此同构下，$s_\infty^{\otimes-d}$ 的原像是一个有理全纯 $(n+1)$ 形式，它沿 $s_\infty^{-1}(0)$ 可能有零点，若 $d<0$，或极点，若 $d>0$。因此，它在 $X$ 上的限制为全纯体积形式 $\eta_X$。有了这一形式，准确地说，有了相关的相对二次形式 $\eta_X^2/dz^2$，全部平行移动映射都如\secref{15c}那样自然带有分次。考虑路径为充分大的圆周的情形。纤维化在无穷远处光滑延拓，但相对二次形式
\SourcePage{304}
不能如此延拓。因此，无穷远处的分次单值变换同痕于某个平移 $S^k$，其中 $k\in2\Z$；这里 $S^k$ 是带常值分次 $-k$ 的恒等映射。显然，$k$ 仅依赖于 $\eta_X/dz$ 沿无穷远纤维的极点阶数 $d-2$。为简便，改取平凡纤维化 $\CC^*\times Y\to\CC^*$，相对体积形式为 $z^{d-2}\eta_Y$。若 $\lambda(t)=e^{2\pi it}$ 为单位圆周，则沿 $\lambda|[0,t]$ 的平行移动是带分次 $(2d-4)t$ 的 $\Id_Y$。因此，沿圆周走完一圈后得到 $k=4-2d$。\znote{296}{极点阶 $d-2$ 的来历：(19.2) 给出 $\mathcal K_{\overline X}\cong\mathcal L^{\otimes-d}$，即 $\eta_X$ 沿除子 $s_\infty^{-1}(0)$ 有 $d$ 阶极点；而 $z=s_0/s_\infty$ 在该除子上恰有二阶极点，故相对形式 $\eta_X/dz$ 的极点阶为 $d-2$。取平凡化后相对形式写成 $z^{d-2}\eta_Y$，绕一圈时 $\arg z$ 增加 $2\pi$，分次增量即 $(d-2)\cdot2=2d-4$；又因 $S^k$ 携带的是常值分次 $-k$，故得 $k=4-2d$。这里的关键是：绕无穷远的单值变换只在分次意义下与恒等相差一个常值，因为决定分次的是相对二次形式而不是形式本身，故平移量必为 $2\Z$ 中的常数，且只能由沿无穷远纤维的极点阶读出。[来源：A，原书 (19.2) 及本段计算。]}

现在讨论第二种方法的一种具体实现，将\eref{19.1}变成第\sr{sec:15}{15 节}意义下的正合 Lefschetz 纤维化。需选择基点 $*\in\CC$；为简化记号，取 $*=0$，因此假设 $0$ 是 $p$ 的正则值。固定 $\CC$ 上的标准辛形式 $\omega_{\CC}$ 及其标准原始一形式 $\theta_{\CC}$。随后将对某个 $r\ge0$，把给定 $\omega_X$ 变形成 $\omega_X^r=\omega_X+rp^*\omega_{\CC}$，类似地把 $\theta_X$ 变形成 $\theta_X^r=\theta_X+rp^*\theta_{\CC}$。尽管缺少紧性，Moser 论证仍然适用，因为满足 $\omega^r(\cdot,Y^r)=\partial\theta_X^r/\partial r=p^*\theta_{\CC}$ 的向量场 $Y^r$，处处与底空间 $\CC$ 上向内的径向向量场 $-z\partial_z$ 的水平提升成正比例。所得结果是一族正合辛嵌入
\[
\rho^r:X\longrightarrow X,\qquad
(\rho^r)^*\omega_X^r=\omega_X,\qquad
(\rho^r)^*\theta_X^r=\theta_X+d(\text{某个函数}).
\]
记住这些嵌入的一个特殊性质很有用；它直接来自上述对 $Y^r$ 的描述：

\StatementNumber{6}\begin{lemma}\label{st:19.6}
取紧子集 $K\subset X$。则存在另一个紧子集 $\widetilde K\subset X$，使对所有 $r\ge0$，都有 $\rho^r(K)\subset\widetilde K$。
\end{lemma}

下一步是对给定辛联络，使用径向平行移动，为\eref{19.1}构造无穷远处的平凡化。它由开子集 $U\subset X$ 及 $U$ 与 $\CC\times U_0$ 之间的逐纤维微分同胚构成，其中 $p|(X\setminus U)$ 为固有映射，$U_0=U\cap X_0$。在此平凡化下，
\[
\theta_X=\theta_{X_0}+\kappa+dR,\qquad
\omega_X=\omega_{X_0}+d\kappa,
\]
其中 $\theta_{X_0},\omega_{X_0}$ 为 $\theta_X,\omega_X$ 在 $X_0$ 上的限制；$\kappa$ 是在纤维切向量上为零的一形式，$R$ 为函数。事实上，在所有 $z=0$ 的点 $(z,x)\in\CC\times U_0$ 上，二者都为零。现在希望修改平行移动，使它在给定平凡化下于无穷远处平凡。为此，取 $X_0$ 上紧支撑的函数 $l$，使它在 $U_0$ 之外等于 $1$，考虑修改后的形式
\begin{equation}\label{eq:19.3}
\widetilde\theta_X=\theta_{X_0}+l\kappa+d(lR)+rp^*\theta_{\CC},\qquad
\widetilde\omega_X=d\widetilde\theta_X,
\end{equation}
其中 $r\gg0$，并约定在平凡化适用的区域之外，将它们延拓为 $\theta_X^r,\omega_X^r$。取大圆盘 $S\subset\CC$，
\SourcePage{305}
赋予限制 $\theta_S=\theta_{\CC}|S$、$\omega_S=\omega_{\CC}|S$。“大”首先意味着 $S$ 的内部包含 $p$ 的全部临界值；稍后还会施加更多条件。取紧子集 $E_0\subset X_0$，使它为某个多重次调和函数的大次水平集，并考虑一块 $E\subset p^{-1}(S)$，使它在无穷远处形如 $S\times E_0$。这里约定，$E$ 包含 $p^{-1}(S)$ 中在无穷远平凡化下看不到的全部部分。此外，选择 $E_0$，使 $l$ 支撑于其内部。令 $\omega_E=\widetilde\omega_X|E$、$\theta_E=\widetilde\theta_X|E$，考虑限制 $\pi=p|E:E\to S$。最后，取与 $\omega_E$ 相容的殆复结构 $I_E$，使它在临界点附近与给定 $I_X$ 一致，使 $\pi$ 伪全纯，并在无穷远附近的平凡化下形如自然乘积结构 $i\times I_{X_0}$。我们断言，所得结果是正合 Lefschetz 纤维化。由于限制在纤维上时 $\widetilde\omega_X=\omega_X$，而且 $r$ 很大，显然 $\widetilde\omega_X$ 仍为辛形式，负 Liouville 向量场沿 $\partial^vE=\pi^{-1}(\partial S)$ 指向内部；考察\eref{15.6}时已见过这一论证。另一方面，已知负 Liouville 场沿 $E_0$ 的边界指向内部，因此由构造，它沿全部 $\partial^hE$ 也如此。事实上，已将纤维化安排成在 $\partial^hE$ 附近按\stref{笔记}{15.2}的意义平凡。$\partial^hE$ 的弱 $I_E$-凸性同样来自 $E_0$ 的相应性质；$\partial^vE$ 的对应性质则由将伪全纯曲线投影到底空间得到。最后，假设 $\overline X$ 带有同构\eref{19.2}。由构造，$I_E$ 作为殆复结构同伦于 $I_X|E$；沿此同伦传递 $\eta_X|E$，即可得到 $(E,I_E)$ 的复体积形式 $\eta_E$。

验证了所需性质后，还须考察\eref{19.3}对辛结构的改变及其对 Fukaya 范畴的影响。$X$ 上如\stref{例题}{19.3}(i)给出的自然多重次调和函数 $\phi$ 为有限型。取次水平集 $X^{\phi,c}$，其中 $c$ 大于全部临界值。将它作为\stref{引理}{19.6}中的 $K$，考虑相应的 $\widetilde K$。构造时，可使 $\widetilde\omega_X=\omega_X^r$ 与 $\widetilde\theta_X=\theta_X^r$ 在 $\widetilde K$ 上成立，于是存在 $X^{\phi,c}$ 到 $E$ 的正合辛嵌入。由\stref{命题}{10.5}，$\F(E)$ 因此含有一个拟等价于 $\F(X^{\phi,c})$ 的全 $\Ai$ 子范畴。这已完全满足我们的目的，即利用 Lefschetz 纤维化得出关于 $X$ 的结论；$\F(E)$ 中任何其他对象都可视为额外对象而忽略。若考察纤维化的几何，而不是总空间 $E$，事情会更简单。按定义，纤维 $E_0$ 是 $X_0$ 上某个多重次调和函数的大次水平集。$\pi:E\to S$ 的平行移动映射当然不是 $p:X\to\CC$ 的平行移动的限制，但在 $E\setminus\partial^hE$ 的一个大紧子集上，两者一致；这是无穷远平凡化没有改动的部分。可以将各项选择安排成使 $E_0$ 中每个 Lagrangian 子流形都能由负 Liouville 流推入该部分。特别地，这使我们能沿用前面关于整体单值变换的结果：若 $\lambda$ 为 $S\setminus\operatorname{Critv}(\pi)$ 中沿正向平行于边界的环路，$L$ 为纤维
\SourcePage{306}
$E_{\lambda(0)}$ 中的正合 Lagrangian 子流形，则单值变换像 $h_\lambda(L)$ 同痕于 $L$。存在\eref{19.2}时，还有更强的版本：任意正合 Lagrangian 膜 $L^\#$ 在 $h_\lambda$ 的自然分次版本下的像满足
\begin{equation}\label{eq:19.4}
h_\lambda(L^\#)\simeq S^{4-2d}L^\#.
\end{equation}

\StatementNumber{7}\begin{proposition}\label{st:19.7}
设 $\pi:E^{2n+2}\to S$ 为来自 Lefschetz 束的正合 Lefschetz 纤维化，带有由\eref{19.2}得到的相对复体积形式。固定基点 $*\in S$，令 $V^\#=(V_1^\#,\ldots,V_m^\#)$ 为纤维 $M=E_*$ 中的消失圈基，任意选择分次使它们成为 Lagrangian 膜；当 $n=1$ 且 $\operatorname{char}(\K)\ne2$ 时，取非平凡 Pin 结构。设 $\A=\Pi(\Tw\F(M))$ 为扭曲复形 $\Ai$ 范畴的分裂闭包，故按定义 $H^0(\A)=D^\pi\F(M)$。若 $d\ne2$，则 $V^\#$ 中的对象分裂生成 $\A$。
\end{proposition}
\begin{proof}
结合\eref{19.4}、\stref{推论}{17.17}及\stref{推论}{5.8}即可。\znote{298}{$d\ne2$ 正是 (19.4) 中平移量 $4-2d$ 非零的条件：$d=2$ 时平移为零，绕边界一圈的分次单值变换就是恒等映射，\eref{19.4} 不再提供任何信息；而“消失圈分裂生成”恰需这个非零平移来配合\stref{推论}{17.17}（非平凡平移给出沿环路单值变换的非平凡性，从而使相应对象在分裂闭包中生成）。下面的\stref{笔记}{19.8}用反例说明此时结论确实不成立，故该限制不是技术性的。[来源：A，原书 (19.4)、命题 19.7 的证明及\stref{笔记}{19.8}。]}
\end{proof}

\StatementNumber{8}\begin{note}\label{st:19.8}
为说明条件 $d\ne2$ 确实必需，取 $\overline X=\CP^1\times\CP^1$，$\mathcal L=\mathcal O_{\overline X}(1,1)$。此时 $X_0=\CC^*$，两个消失圈都同痕于同一圆周 $L=S^1\subset\CC^*$，带有非平凡 Pin 结构。将它记为 $L^\#$，将带平凡 Pin 结构的对应物记为 $L^\#\otimes\beta$。若系数域 $\K$ 的特征不为 $2$，则由\eref{12.14}，$HF^*(L^\#,L^\#\otimes\beta)=H^*(L;|\beta|_{\K})=0$。因此，$L^\#\otimes\beta$ 不在由 $L^\#$ 分裂生成的子范畴中。\znote{298}{这个例子可逐项核对：$\mathcal L=\mathcal O(1,1)$，而 $K_{\CP^1\times\CP^1}=\mathcal O(-2,-2)=\mathcal L^{\otimes-2}$，故 $d=2$。$D=s_\infty^{-1}(0)$ 是 $(1,1)$ 除子，$X=\overline X\setminus D$；纤维 $X_0$ 为 $s_0^{-1}(0)\setminus D$，其中 $s_0^{-1}(0)\cong\CP^1$ 与 $D$ 恰交于两点，故 $X_0=\CC^*$，两个消失圈都同痕于 $S^1\subset\CC^*$。$d=2$ 使 (19.4) 的平移为零，这正是结论失效的根源。再看后半句：$HF^*(L^\#,L^\#\otimes\beta)=H^*(L;|\beta|_{\K})$ 之所以为零，是因为 $\beta$ 使系数成为秩一的非平凡局部系统（沿 $S^1$ 的单值变换为 $-1$），此时 $H^0$ 与 $H^1$ 同时为零，故 $L^\#\otimes\beta$ 无法由 $L^\#$ 分裂生成。[来源：A，原书 (19.2)、\stref{例题}{19.3}(i)、\eref{12.14} 及本节。]}

尽管如此，仍可尝试将 $d=2$ 的一般理论再推进一些。写 $L_k=\tau_{V_{k+1}}\cdots\tau_{V_m}(L)$，赋予诱导的膜结构。在正合三角
\SourceDiagram{diagram-19-b-triangle}{0.384230355}
中，将所有 $\to$ 映射复合，得到数
\[
c_{L^\#}\in HF^0(L_m^\#,L_0^\#)\cong HF^0(L^\#,L^\#)\cong H^0(L;\K)=\K.
\]
考察\stref{推论}{5.8}的证明，可知 $c_{L^\#}=0$ 是 $L^\#$ 属于由消失圈分裂生成的子范畴的一个充分条件。另一方面，回顾\stref{推论}{17.18}中这些正合三角的构造，可知每个 $\to$ 映射都是计数给定 Lefschetz 纤维化一小块中的伪全纯截面的不变量，该小块仅含一个临界点。利用标准粘合定理，可将 $c_{L^\#}$ 写成整个纤维化的一个不变量 $\Phi_{E/S,*}$，其中移动的 Lagrangian 边界条件 $F$ 实现同痕
\SourcePage{307}
\eref{19.4}。等价地，从原来的代数几何数据看，可将 $F$ 视为 $X$ 中在 $p:X\to\CC$ 的底空间内一个大圆周上纤维化的 Lagrangian 子流形，因此它靠近无穷远纤维。于是，$c_{L^\#}$ 计数 $X$ 中边界位于 $F$ 的全纯圆盘。可尝试通过向法锥退化，或辛切割，进一步分析这一情形。该过程使 $X$ 退化为两块之并：$\overline X$ 与 $\mathcal L^{-1}\to s_\infty^{-1}(0)$ 的总空间，两者沿 $s_\infty^{-1}(0)$ 粘合。Lagrangian 子流形 $F$ 移入第二块，可用类似\bcite{79}的技术计数全纯圆盘。这里不再尝试作更具体的说明。
\end{note}

\clearpage
\subsection{Picard--Lefschetz 数据}\label{sec:19c}
再次考虑 $(n+1)$ 维光滑射影簇，这次明确将它嵌入射影空间 $\overline X\subset\CP^N$。与前面一样，对某个 $d\in\Z$，要求
\begin{equation}\label{eq:19.5}
\mathcal K_{\overline X}\cong\mathcal O_{\overline X}(-d).
\end{equation}
假设射影嵌入相对于标准齐次坐标 $(y_0:\cdots:y_N)$ 处于一般位置。准确的条件将在下面列出，但要点是，施加一个适当的 $\operatorname{PGL}_{N+1}$ 元素后，或者等价地，换用另一齐次坐标系后，这些条件总能成立。

取无穷远超平面为 $\CP^{N-1}=\{y_N=0\}$，将其补集等同于 $\CC^N$，坐标为 $(x_0=y_0/y_N,\ldots,x_{N-1}=y_{N-1}/y_N)$。对 $0\le k\le n+1$，取子空间 $H^k=\{x_0=\cdots=x_{k-1}=0\}\subset\CC^N$ 及其闭包 $\overline H^k\subset\CP^N$。定义
\[
\begin{aligned}
\overline X^k&=\overline X\cap\overline H^k, &X^k&=X\cap H^k,\\
\Psi^k&=x_k:X^k\longrightarrow\CC,&
\varpi^k&=(x_k,x_{k+1}):X^k\longrightarrow\CC^2.
\end{aligned}
\]
一般情形下，以下性质成立：
\begin{itemize}
\item 对 $0\le k\le n+1$，$\overline X$ 与全部 $\overline H^k$ 及 $\overline H^k\cap\CP^{N-1}$ 横截。这意味着每个 $X^k$ 光滑，而且在无穷远处也光滑。
\end{itemize}
注意，按定义，$X^{k+1}=(\Psi^k)^{-1}(0)$、$\Psi^{k+1}=\varpi^k|X^{k+1}$。此外，由\eref{19.5}，有 $\mathcal K_{\overline X^k}\cong\mathcal O_{\overline X^k}(k-d)$。\znote{299}{这是本节最容易错位的指标约定：$\Psi^k=x_k$ 的纤维是 $X^{k+1}$，所以由 $\Psi^k$ 的临界点集合 $C^k$（含于 $X^k$）标号的消失圈都住在 $X^{k+1}$ 里，$\A^{k+1}$ 才是由它们生成的有向范畴。\eref{19.6} 中 $\F(X^k)\to\F(\Psi^k)\to\Tw\A^{k+1,\to}$ 的指标比 $k$ 大一位，正是由此而来。原书 p302 有一处把 $\gamma_j^{k-1}$ 印作 $\gamma_j^{k+1}$，读时须按此约定校正。顺带一提，$\varpi^k=(x_k,x_{k+1})$ 的两个分量正是 $\Psi^k$ 与 $\Psi^{k+1}$，故双纤维化的指标同样是“错一位”的。[来源：A，原书 (19b)–(19c) 的定义与 (19.6)。]}继续列出一般性质：
\begin{itemize}
\item 对 $0\le k\le n$，每个函数 $\Psi^k$ 只有非退化临界点，且任意两点不在同一纤维中。
\SourcePage{308}
\item 对 $0\le k\le n-1$，$D\varpi^k$ 是一般的，即横截于 $\Hom(TX^k,\CC^2)$ 中由低秩线性映射构成的各子簇。若 $w\in\CC$ 为 $\Psi^k$ 的临界值，$x$ 为纤维 $(\Psi^k)^{-1}(w)$ 中的唯一临界点，则还满足：$(D^2\Psi^k)_x$ 在 $(D\varpi^k)_x$ 的核上的限制仍非退化；函数
\[
(\Psi^k)^{-1}(w)\setminus\{x\}\xrightarrow{\ \varpi^k\ }\{w\}\times\CC=\CC
\]
具有非退化临界点；任意两个这样的点的值不同，并且也都不同于 $\varpi^k(x)$ 本身。
\end{itemize}
这些性质在一般情形下成立是经典事实，但为方便起见，仍概述证明。第一部分是 Bertini--Sard 定理的直接应用。第二部分可用以下标准论证归约到同一定理。考虑拉回丛 $T^*X^k\to\CC^{N-k}\times X^k$ 的截面 $(\lambda,x)\mapsto d(\sum_{i\ge k}\lambda_i x_i)$。它总是与零截面横截相交，所以其零点集 $Z\subset\CC^{N-k}\times X^k$ 光滑。若取投影 $Z\to\CC^{N-k}$ 的一个正则值 $\lambda$，则函数 $\sum_i\lambda_i x_i$ 只有非退化临界点。显然，任意这样的函数都可通过线性坐标变换变成 $x_k$。因此，给定任意 $\overline X$，可用一个接近恒等的 $\operatorname{PGL}_{N+1}$ 元素改变其嵌入，使变换后的簇上函数 $\Psi^k=x_k$ 只有非退化临界点。再对嵌入作微扰，可使各临界点位于不同纤维。同一论证适用于 $\varpi^k$。读者想必已经注意到，这些条件是\secref{15e}中条件的代数几何对应物。其推论也类似：若 $x$ 为 $\Psi^k$ 的临界点，值为 $w$，则 $\operatorname{Critv}(\varpi^k)$ 在 $\varpi^k(x)$ 处与直线 $\{w\}\times\CC$ 相切，在其他各处都与之横截。按\secref{15e}，定义 $C^k=\operatorname{Critv}(\Psi^k)$，$0\le k\le n$，及 $F^k=\operatorname{Fakev}(\Psi^k,\varpi^k)$，$0\le k\le n-1$。为简化记号，令 $F^n=\varnothing$。注意，由假设，$0\notin C^k\cup F^k$。

固定 $0\le k\le n$，写 $C^k=\{c_j^k\}$。为每个 $j$ 选择嵌入路径 $\gamma_j^k:[0,1]\to\CC$，满足 $\gamma_j^k(0)=0$、$\gamma_j^k(1)=c_j^k$，$t<1$ 时 $\gamma_j^k(t)\notin C^k$，且对所有 $t$，$\gamma_j^k(t)\notin F^k$。这些路径应构成\secref{16d}意义下的基。在最低维情形 $k=n$，它们只是从 $0$ 走向分歧覆盖 $\Psi^n:X^n\to\CC$ 的某个分歧点的路径。可为每条路径配上无序点对 $\mu_j^{n+1}\subset X^{n+1}$，标记沿路径合并的覆盖的两层。更高维时，如\secref{16i}那样，对 $t\in[0,1)$，追踪子集 $\operatorname{Critv}(\varpi^k)\cap(\{\gamma_j^k(t)\}\times\CC)\subset\CC$ 的移动。由此得到嵌入路径 $\mu_j^{k+1}:[-1,1]\to\CC$，满足 $(\mu_j^{k+1})^{-1}(C^{k+1})=\{-1,1\}$，在同痕及反向意义下唯一。这一过程产生以下信息，称为 $X$ 的 \emph{Picard--Lefschetz 数据}：
\SourcePage{309}
\begin{itemize}
\item 有限集 $X^{n+1}$，以及子集族 $\{\mu_j^{n+1}\}$；
\item 对每个 $0\le k\le n$，在带标记点 $\{0\}\cup C^k\cup F^k$ 的平面 $\CC$ 中的路径集 $\{\gamma_j^k\}$，以及 $k>0$ 时的 $\{\mu_j^k\}$；
\item 当 $H^1(\overline X)\ne0$ 时，还需一个稍后说明的同伦数据。
\end{itemize}
注意，这些应被视为拓扑图形：可移动标记点，同时对相关路径作适当同痕，而不实质改变数据。因此，可通过指定一个穿孔平面微分同胚如何把标准路径变成给定路径，来编码每条路径。这样，剩下的主要是一组辫群中的字，这是 Moishezon 辫分解技巧的弱化版本。

考察这些数据包含多少关于 $X$ 的信息。首先，分歧数据容许在拓扑上重构 $\Psi^n:X^n\to\CC$。为了得到 $X^n$ 的良好定义的 Fukaya 范畴，还需知道 $\mathcal K_{X^n}$ 的平凡化，至少在拓扑上知道，即知道其在同伦意义下的 $C^\infty$ 平凡化。由构造，该平凡化来自一个处处非零的复一形式，在每个无穷远点处有 $d-n$ 阶极点。为确定其同伦类，应记录一组张成 $H_1(\overline X^n)$ 的浸入曲线的“旋转数”，即 Maslov 指标。若一条曲线是 $\Psi^{n-1}$ 的消失圈，已知其旋转数为零，因为该曲线容许分次。假设 $H^1(\overline X)=0$。由弱 Lefschetz 定理，$H^1(\overline X^{n-1})=0$，且消失圈张成 $H_1(\overline X^n)$，这意味着 $\mathcal K_{X^n}$ 的平凡化唯一确定。否则，可能还须记录一些额外旋转数，这就是上面提到的最后一项数据。

\StatementNumber{9}\begin{note}\label{st:19.9}
另一种做法有些偏离一般思路：假设精确知道分歧点集 $C^n$。此时，$X^n$ 在全纯同构意义下确定，并带有一个在非零常数倍意义下唯一的全纯一形式，它没有零点，在无穷远点附近具有所需阶数的极点。
\end{note}

在上一层，$X^n$ 中存在唯一的简单闭曲线，使它们在 $\Psi^n$ 下的像为 $\mu_j^n$。此外，这些曲线都不与零同调，因为它们来自消失圈，故对适当的一形式选择是正合的。按 Donaldson 与 Gompf 的方法，这组曲线连同相应路径 $\gamma_j^{n-1}$，在辛同构意义下确定 $\overline X^{n-1}$；这里辛同构相对于无穷远除子而言，因而诱导 $X^{n-1}$ 的正合辛同构。最后，$\mathcal K_{X^n}$ 的平凡化以在同伦意义下唯一的方式，诱导 $\mathcal K_{X^{n-1}}$ 的平凡化。

对 $k<n$，Picard--Lefschetz 数据本身是否足以确定 $X^k$ 或 $\overline X^k$ 的辛同构类型，似乎值得怀疑。因为我们记录了匹配路径，却没有记录构造实际
\SourcePage{310}
匹配圈所需的同痕。将 $X^k$ 紧化为 $\overline X^k$ 时，实际上还有一个潜在问题，即无穷远除子的附着方式。可尝试通过指定额外数据克服这些问题，但我们不采用这一方向。\znote{302}{为什么说数据不足，可与上文三级信息对照：$X^n$ 由分歧数据与旋转数确定；$X^{n-1}$ 由 $X^n$ 中的曲线连同路径 $\gamma_j^{n-1}$ 确定；再往下一层，则需把 $\gamma_j^{k-1}$ 同痕成匹配路径 $\mu_j^k$ 的具体同痕，而同痕不同可以使得到的匹配圈相差一次 Lagrangian 手术，$X^k$ 的辛结构随之改变。紧化还多一层歧义：无穷远除子如何附着并未被数据记录。[来源：C，译者按（上下文推断，无外部依据），据本节末三段。]}

\subsection{结论}\label{sec:19d}
我们将在 Fukaya 范畴框架内进行类似的代数构造，以取代从 Picard--Lefschetz 数据几何地恢复各 $X^k$ 的尝试。事实表明，在这一层次上，上述困难消失了，比较\stref{笔记}{16.11}。为简化叙述，将 $\Psi^k$ 与 $\varpi^k$ 分别当作真正的正合 Lefschetz 纤维化与双纤维化。这样做没有问题：$\Psi^k$ 的情形已在\secref{19b}解释，$\varpi^k$ 的讨论完全类似。因此，有两列 Fukaya 范畴 $\F(X^k)$ 与 $\F(\Psi^k)$。对每个 $0<k\le n$，令 $V_j^k\subset X^k$ 为 $\gamma_j^{k-1}$ 对应的消失圈；由构造，它也是 $\mu_j^k$ 对应的匹配圈，见\stref{引理}{16.15}。照常将它们赋予 Lagrangian 膜结构，将所得对象集记为 $V^{k,\#}$。写成
\[
\A^k=\F(V^{k,\#}),\qquad
\A^{k,\to}=\F^\to(V^{k,\#});
\]
也就是说，$\A^k$ 是 $\F(X^k)$ 中由这组消失圈构成的全 $\Ai$ 子范畴，$\A^{k,\to}$ 为相应有向 $\Ai$ 子范畴。

算法的起点与前面相同，即在拓扑上重构 $X^n$、$\mathcal K_{X^n}$ 的平凡化及消失圈 $V^n$。给定这些数据后，可按第\sr{section:13}{13 节}的方法组合地计算 $\A^{n,\to}$。迭代步骤如下。假设给定一个拟同构于 $\A^{k+1,\to}$ 的有向 $\Ai$ 范畴 $\B^{k+1,\to}$。由\stref{命题}{18.13}与\stref{定理}{18.24}，得到在上同调上全忠实的 $\Ai$ 函子
\begin{equation}\label{eq:19.6}
\F(X^k)\longrightarrow\F(\Psi^k)\longrightarrow\Tw\A^{k+1,\to}.
\end{equation}
如\stref{推论}{18.26}的证明中那样，可确定任意匹配圈在\eref{19.6}下的像，以及 $\Tw\B^{k+1,\to}$ 中的相应对象，可能相差一个分次平移；这种不确定性最终无关紧要，为简便，假设分次已正确选取。因此，$\A^k$ 拟同构于一个显式给出的全 $\Ai$ 子范畴 $\B^k\subset\Tw\B^{k+1,\to}$。若 $k>1$，就只保留拟同构于 $\A^{k,\to}$ 的 $\B^{k,\to}$，\emph{舍去部分信息}，再继续到高一维。这样到达 $k=1$ 时，得到拟同构于 $\A^1$ 的 $\Ai$ 范畴 $\B^1$。最后还可再前进半步，重新应用\eref{19.6}，说明 $\F(X)$ 本身拟等价于 $\Tw\B^{1,\to}$ 的某个事先未知的全 $\Ai$ 子范畴。\znote{302}{为什么要舍去：归纳的输入只有有向范畴 $\A^{k,\to}$，而 $\B^k$ 中非有向的态射与高阶运算在\eref{19.6}的函子里根本不出现；保留它们反而使下一步无法与 $\Tw\B^{k+1,\to}$ 比较。这也是上面 19c 末尾“数据不足以恢复 $X^k$”的代数对应物：几何上丢掉的是同痕，代数上丢掉的是非有向部分。[来源：A，原书本节与 19c 末段。]}

\SourcePage{311}
\begin{proof}[定理~\sr{st:B}{B} 在 $n\ge1$ 时的证明]
由上述论证，结论现已显然成立，所考虑的 $\Ai$ 范畴为 $\Tw\B^{1,\to}$。
\end{proof}

\begin{proof}[推论~\sr{st:C}{C} 在 $n\ge1$ 时的证明]
由\stref{推论}{5.27}，可将 $\B^{1,\to}$ 替换为一个与之拟同构、实际上是 $\dg$ 范畴的有向 $\Ai$ 范畴 $\C^\to$。取在上同调上全忠实的 $\Ai$ 函子 $\F(X)\to\Tw(\C^\to)$。必要时引入已有对象的重复副本后，可假设此函子在对象上单射，于是 $\F(X)$ 拟同构于其像，而该像是具有有限维态射空间的 $\dg$ 范畴。
\end{proof}

\begin{proof}[定理~\sr{st:A}{A} 的证明]
由\stref{命题}{19.7}，消失圈分裂生成 $\Pi\Tw\F(X^1)$。由\stref{推论}{4.9}与\stref{引理}{4.7}，$\Pi\Tw\F(X^1)$ 拟等价于 $\Pi\Tw\A^1$，因此也拟等价于 $\Pi\Tw\B^1$。
\end{proof}

细心的读者会注意到，该算法在最低维时未给出完整信息：可以很容易地确定 $\F(X^n)$ 的任意有向 $\Ai$ 子范畴，但这一做法并不先验地适用于 $\F(X^n)$ 本身。部分原因在于，没有从真正平凡的零维情形开始归纳；可如下弥补。取 $\Psi^n:X^n\to\CC$，按第\sr{section:18}{18 节}的过程定义 $\F(\Psi^n)$，即取沿光滑纤维分歧的二重覆盖 $\widetilde X^n\to X^n$，并考虑 $\Z/2$-不变 Lagrangian 子流形。特别地，$\widetilde X^n$ 中有 (B) 型等变膜，它们是消失路径 $\gamma_j^n$ 相应 Lefschetz 顶针的原像。考察\stref{命题}{18.14}的证明，可知它们在 $\F(\Psi^n)$ 中构成例外集。因此，相应全 $\Ai$ 子范畴与有向 $\Ai$ 子范畴拟同构。另一方面，该有向 $\Ai$ 子范畴可用第\sr{section:13}{13 节}的组合方法处理，因为定义相关 Floer 上链群与 $\mu^d$ 运算时，可以使用标准复结构及零非齐次项，\stref{假设}{14.5}仍会满足。在这一特殊情形中，所有伪全纯多边形实际上都是常值的；其正则性来自\stref{引理}{11.5}，如\stref{引理}{14.6}、\stref{引理}{14.7}的证明中那样。将所得有向 $\Ai$ 范畴记为 $\B^{n+1,\to}$。它带有在上同调上全忠实的函子
\begin{equation}\label{eq:19.7}
\F(X^n)\longrightarrow\Tw\B^{n+1,\to}.
\end{equation}
利用它，可将定理~\sr{st:B}{B} 与推论~\sr{st:C}{C} 的证明推广到 $n=0$ 的情形，此时 $X^n=X$。

\StatementNumber{10}\begin{note}\label{st:19.10}
为类似地推广定理~\sr{st:A}{A}，需确定 $\Psi^{n-1}$ 的消失圈在\eref{19.7}下的像。遗憾的是，这里有一个高维时未出现的微小但棘手之处：在\stref{命题}{18.21}的对应陈述中，零次态射空间现在为二维，
\SourcePage{312}
必须在该空间的正确元素上构造锥，它对应于消失圈上的非平凡 Pin 结构。
\end{note}

\section{\texorpdfstring{$(A_m)$}{(Am)} 型 Milnor 纤维}\label{sec:20}\label{section:20}
\subsection{生成元}\label{sec:20a}
对任意 $n\ge1$、$m\ge2$，考虑函数 $\Psi(x)=p(x_0)+x_1^2+\cdots+x_n^2:\CC^{n+1}\to\CC$，其中 $p$ 为某个 $m+1$ 次多项式。我们关心
\[
X_m^n=\Psi^{-1}(w),
\]
其中 $w$ 为 $\Psi$ 的正则值。这是光滑超曲面，特别地，是具有平凡典范丛的 $n$ 维仿射代数簇。遵循\stref{笔记}{19.5}的思路，我们直接在 $X_m^n$ 本身上工作，赋予它 $\CC^{n+1}$ 上平坦 Kähler 形式及其标准原始一形式的限制，以及满足 $\eta\wedge d\Psi=dx_0\wedge\cdots\wedge dx_n$ 的二次复体积形式 $\eta^2$。在正合辛同构意义下，$X_m^n$ 与 $p,w$ 的选择无关，因此 $\F(X_m^n)$ 也如此。

假设 $p$ 是一般的，即其临界点非退化。此时，经适当技术调整，可将 $\Psi$ 视为 Lefschetz 纤维化。适当选择消失路径基，得到消失圈基 $(V_1,\ldots,V_m)$，它们两两横截，其交点满足
\begin{equation}\label{eq:20.1}
|V_i\cap V_j|=\begin{cases}1,&|i-j|=1,\\0,&|i-j|\ge2.\end{cases}
\end{equation}
将其提升为 Lagrangian 膜 $(V_1^\#,\ldots,V_m^\#)$，取相应全 $\Ai$ 子范畴 $\F_m^n\subset\F(X_m^n)$。

只用\eref{20.1}及 Floer 上同调的一般性质，就容易算出上同调层次的 Donaldson--Fukaya 范畴 $H(\F_m^n)$。按\bcite{85}，用箭图代数来表述结果。考虑分次箭图
\SourceEquationDiagram{diagram-20-2}{0.336461176}{20.2}
一般地，可为任意这样的箭图配上分次结合代数，称为\emph{路径代数}。在本例中，路径代数的元素是基本表达式 $(i_l,\ldots,i_0)$ 的线性组合，其中 $l\ge0$、$i_j\in\{1,\ldots,m\}$、$|i_{j+1}-i_j|=1$。这一基元素的次数等于满足 $i_{j+1}=i_j-1$ 的 $j$ 的个数乘以 $n$。\znote{304}{这条分次规则的几何来源是 Maslov 指标：$V_i\cap V_{i+1}$ 的唯一交点作为 $V_i\to V_{i+1}$ 的态射次数为 $0$，作为 $V_{i+1}\to V_i$ 的态射次数则为 $n$（同一个横截交点，两种读法相差 $n$），于是每“下降”一步贡献 $n$。下文\eref{20.4}中 $|f|=n-1$ 与符号 $(-1)^{n-1}$ 都由此而来。另外，$|i_{j+1}-i_j|=1$ 表示路径只在相邻顶点之间走，故每条基本路径的次数是 $n$ 的非负整数倍；这与 (20.1) 中“下标相距两个顶点以上的消失圈不相交”是同一件事的两种说法。[来源：A，原书 (20.1) 与 (20.4)。]}乘法由路径复合给出，例如 $(i+2,i+1)(i+1,i)=(i+2,i+1,i)$。特别地，
\SourcePage{313}
长度为零的路径 $(i)$ 是两两正交的幂等元。它们使路径代数成为半单环 $R_m\cong\K^m$ 上的代数。

令 $A_m^n$ 为路径代数模去关系
\[
(i+2,i+1,i)=0,\qquad(i-2,i-1,i)=0,\qquad
(i,i+1,i)+(i,i-1,i)=0
\]
得到的商；关系在所写路径存在的指标范围内施加，当 $m=2$ 时还令所有长度为三的路径为零。\znote{305}{原书未单独讨论 $m=2$：此时上面三条关系全部落空（$(1,0,1)$、$(2,3,2)$ 都不存在），路径代数便含任意长的交替路径，与紧随其后“基由 $(i)$ 及 $\pm(i,i\pm1,i)$ 组成”这一有限维描述相矛盾。译文因此补出“长度三路径为零”一句，此时 $A_2^n$ 恰为 6 维：$e_1,e_2$、$(1,2),(2,1)$ 以及 $(1,2,1),(2,1,2)$。原书此处未作说明，删去该句会使同段的基描述不成立，故保留译文并以此注交代。[来源：A，原书 (20.1)–(20.2) 及该段对基的描述。]}这仍为分次 $R_m$-代数，故也可等价地视为具有 $m$ 个对象的分次线性范畴；其中态射空间 $(j)A_m^n(i)$ 由从第 $i$ 个到第 $j$ 个顶点的路径表示。据此，结论是：适当选择消失圈的分次后，有 $A_m^n\cong H(\F_m^n)$。更具体地，$A_m^n$ 作为 $\K$-向量空间的一组基由每个 $i$ 对应的路径 $(i)$ 及 $\pm(i,i\pm1,i)$ 组成，它们对应于 $H^*(V_i;\K)$ 的标准生成元；还包括 $(i,i+1)$ 与 $(i+1,i)$，它们是 $V_i\cap V_{i+1}$ 中唯一交点对应的 Floer 元素。

乘积 $\phi=\tau_{V_1}\cdots\tau_{V_m}$ 为 $X$ 的辛自同构，用经典术语说，是 $(A_m)$ 型超曲面奇点的单值变换映射。由于该奇点是加权齐次的，可以如\bcite{125}，小节~4c 那样，显式分析 $\phi$ 的几何。计入 Dehn 扭转的自然分次后，得到：对任意正合 Lagrangian 膜 $L^\#$，有
\[
\phi^{2m+2}(L^\#)\simeq S^{(4-2n)(m+1)-4}L^\#.
\]
平移量总是非零。由与\stref{命题}{19.7}相同的论证，各 $V_j^\#$ 分裂生成 Fukaya 范畴，因此存在拟等价
\begin{equation}\label{eq:20.3}
\Pi\Tw\bigl(\F(X_m^n)\bigr)\simeq\Pi\Tw(\F_m^n).
\end{equation}
为完全确定 $\Pi\Tw\bigl(\F(X_m^n)\bigr)$，还需在所有上同调为 $A_m^n$ 的 $\Ai$ 范畴的拟同构类中，刻画 $\F_m^n$。

\subsection{高维情形}\label{sec:20b}
假设 $n\ge2$。此时，可将\stref{引理}{1.9}与适当的 Hochschild 上同调群消失结果结合，证明任意上同调为 $A_m^n$ 的 $\Ai$ 范畴必为形式的，即拟同构于 $A_m^n$ 上的平凡 $\Ai$ 结构，后者的所有阶数大于 $2$ 的复合映射都为零。这已在\bcite{134}中完成。下面采用更具体的方法证明这一结果，避免使用抽象分类理论。在几何方面，这依赖于第\sr{section:18}{18 节}的材料，因此需要技术假设 $\operatorname{char}(\K)\ne2$。

取函数 $\pi:X\to\CC$，$\pi(x)=x_0$。再次经过技术调整后，可将它视为 Lefschetz 纤维化，将 $(\Psi,\pi)$ 视为双纤维化。此外，可显式分析 $\pi$ 的纤维与平行移动映射\bcite{85}：每个光滑纤维都辛同构于余切丛 $T^*S^{n-1}$；每条消失路径都给出同一个消失圈，即零截面
\SourcePage{314}
$S^{n-1}\subset T^*S^{n-1}$。因此，每条连接两个临界值的路径自然都是匹配路径。特别地，$V_j$ 是一链匹配路径 $\mu_j$ 对应的匹配圈，见\sr{fig:20.1}{图~20.1}。
\SourceFigure{figure-20-1}{0.350999622}{20.1}

此后固定 $m$ 及 $n\ge2$，写 $\F=\F_m^n$、$A=A_m^n$。若只考虑余切丛中的零截面 $S^{n-1}$，则相应 Fukaya $\Ai$ 代数拟同构于 $\K\cdot e\oplus\K\cdot f$，其中
\begin{equation}\label{eq:20.4}
\begin{gathered}
|e|=0,\qquad |f|=n-1,\\
\mu^2(e,e)=e,\qquad\mu^2(e,f)=(-1)^{n-1}f,\qquad\mu^2(f,e)=f,
\end{gathered}
\end{equation}
其余所有乘积都为零。\znote{306}{(20.4) 中 $\mu^2(e,f)=(-1)^{n-1}f$ 而 $\mu^2(f,e)=f$，并非笔误：它等价于 $\mu^2(e,f)=(-1)^{|f|}\mu^2(f,e)$，即单位元 $e$ 在左侧作用时带符号 $(-1)^{|x|}$。这是本书 $\Ai$ 符号约定（输入按圆盘边界顺序排列）的产物，不能按严格含幺代数的 $e\cdot f=f$ 去“改正”，否则\eref{20.5}的复合及后文各式的符号会整体出错。另一种看法：$e$ 是单位、$f$ 是 $S^{n-1}$ 的基本类，$\mu^2(f,e)=f$ 即杯积，而 $\mu^2(e,f)$ 上的符号正是把 $e$ 从右端移到左端所付的代价。[来源：A，原书 (20.4) 及其后的 (20.5)。]}在上同调层次上，这是显然的，因为考察的只是
\[
HF^*(S^{n-1},S^{n-1})=H^*(S^{n-1};\K).
\]
在上链层次上，所需消失结果来自次数考察，至少在把 $\Ai$ 结构做成严格含幺之后是如此。另一种方法是使用\bcite{58}的结果，以及 $S^{n-1}$ 在传统代数拓扑意义下为形式的这一事实。下一步是写出 $\pi$ 的一个消失圈基所对应的有向 $\Ai$ 范畴 $\B^\to$；该基当然只由 $m+1$ 个 $S^{n-1}$ 的副本构成。这很直接：将对象记为 $(Z_0,\ldots,Z_m)$，则态射为
\begin{equation}\label{eq:20.5}
\hom_{\B^\to}(Z_i,Z_j)=\begin{cases}
\K\cdot e_{ji}\oplus\K\cdot f_{ji},&i<j,\\
\K\cdot e_{ii},&i=j,\\
0,&i>j,
\end{cases}
\end{equation}
次数与复合均同\eref{20.4}。取 $\Tw\B^\to$ 的对象 $(C_1,\ldots,C_m)$，定义为
\begin{equation}\label{eq:20.6}
C_j=\Cone(e_{j,j-1}).
\end{equation}
\SourcePage{315}
\stref{命题}{18.21}说明，由这些对象构成的全 $\Ai$ 子范畴拟同构于 $\F$。容易验证，在上同调层次上，这恢复了 $A$。更重要的是，该过程给出显式的链层次模型。最后，这个模型还具有一个有用的额外对称性，它并不先验地来自几何。

为简化记号，使用\stref{笔记}{3.26}中的扭曲复形版本。给 $\B^\to$ 再赋予长度分次，使 $e_{ji}$ 的长度为零，$f_{ji}$ 的长度为二；这一名称稍后会得到解释。于是可引入\emph{齐次加性完备化} $\Sigma^{\mathrm{ho}}\B^\to$，其对象为有限和
\begin{equation}\label{eq:20.7}
\bigoplus_i S^{-\sigma_i}\Lambda^{-\tau_i}Y_i,
\end{equation}
其中 $\Lambda$ 为使长度分次降低一的形式运算。任意两个这样的对象之间的态射，便带有按次数与长度给出的双分次。同样，有齐次扭曲复形的 $\Ai$ 范畴 $\Tw^{\mathrm{ho}}\B^\to$，其对象是带保持长度分次的微分的表达式\eref{20.7}。\znote{307}{为什么要多引入一个长度分次：次数分次无法表达“$\Cone(e_{j,j-1})$ 应记作长度为 $j$ 的对象”这一信息。把 $e_{ji}$ 的长度定为 $0$、$f_{ji}$ 定为 $2$，再用形式运算 $\Lambda$ 降低长度，就能把 (20.6) 写成 $S\Lambda^jZ_{j-1}\oplus\Lambda^jZ_j$，使“关于长度分次齐次”的 $\Ai$ 结构有意义；这样用\stref{命题}{1.12}消去 $\mu^1$ 时才不会破坏双分次，最后才能断言 $d>2$ 的高阶 $\mu^d_{\A}$ 必为零。[来源：A，原书本节。]}例如，可将\eref{20.6}变成如下齐次扭曲复形：
\[
C_j=\Lambda^j\Cone(e_{j,j-1})=S\Lambda^jZ_{j-1}\oplus\Lambda^jZ_j,\qquad
\delta_{C_j}=\begin{pmatrix}0&0\\e_{j,j-1}&0\end{pmatrix}.
\]
如前，取 $\Tw^{\mathrm{ho}}\B^\to$ 中由这些对象构成的全 $\Ai$ 子范畴，它是双分次 $\Ai$ 范畴。在上同调层次上，这一额外分次确实恢复了 $A$ 作为\eref{20.2}的路径代数的商所继承的长度分次。

对前面构造的全 $\Ai$ 子范畴应用\stref{命题}{1.12}，以消去 $\mu^1$。可用保持双分次的方式做到这一点，结果是 $A$ 上关于长度分次齐次的 $\Ai$ 结构 $\A$。考察可知，此时高阶复合映射 $\mu_{\A}^d$，$d>2$，必为零。另一方面，由构造，这一 $\Ai$ 结构拟同构于 $\F$，所以 $\F$ 为形式的。

\subsection{曲面情形}\label{sec:20c}
剩下的 $n=1$ 情形更为棘手。我们将采用一种较为间接的方法，把显式计算与抽象分类理论结合起来。以下始终假设系数域满足 $\operatorname{char}(\K)=0$；这一假设会被反复用到，甚至在论证中纯代数的部分也是如此。例如，若取 $\operatorname{char}(\K)=2$，\stref{引理}{20.2}就不成立。

首先，我们考察一种可能导致上同调为 $A=A_m^1$ 的 $\Ai$ 范畴 $\A$ 不具形式性的情形。不失一般性，可设 $\mu_{\A}^1=0$，这样 $\A$ 就是 $A$ 本身上的一个 $\Ai$ 结构。还可假设 $\A$ 严格含幺。
\SourcePage{316}
\StatementNumber{1}\begin{lemma}\label{st:20.1}
再假设 $\mu_{\A}^3=0$，且 $m\ge5$。考虑
\begin{equation}\label{eq:20.8}
\begin{aligned}
&\mu_{\A}^4\bigl((3,2),(2,1),(1,2),(2,3)\bigr)
 -\mu_{\A}^4\bigl((3,2),(2,3),(3,2),(2,3)\bigr)\\
&\quad+\mu_{\A}^4\bigl((3,4),(4,3),(3,2),(2,3)\bigr)
 +\mu_{\A}^4\bigl((3,2),(2,3),(3,4),(4,3)\bigr)\\
&\quad-\mu_{\A}^4\bigl((3,4),(4,3),(3,4),(4,3)\bigr)
 +\mu_{\A}^4\bigl((3,4),(4,5),(5,4),(4,3)\bigr)=q(3),
\end{aligned}
\end{equation}
其中 $q\in\K$。若 $q\ne0$，则 $\A$ 不是形式的。
\end{lemma}
\begin{proof}
设 $\A,\widetilde{\A}$ 为两个这样的 $\Ai$ 结构，并以 $q,\widetilde q$ 表示它们按\eref{20.8}得到的系数。假设有一个严格含幺的形式微分同胚 $\Phi:\A\to\widetilde{\A}$，其一次项为 $\Phi^1=\Id$。直接计算可得
\begin{equation}\label{eq:20.9}
(\widetilde q-q)(3)=
\bigl(\Phi^2\bigl((3,2),(2,3)\bigr)-\Phi^2\bigl((3,4),(4,3)\bigr)\bigr)^2.
\end{equation}
再考虑其逆形式微分同胚 $\Psi$，它满足 $\Psi^2=-\Phi^2$，并考察 $q-\widetilde q$ 对应的等式。比较两式可见，\eref{20.9}的右端为零，因而 $q=\widetilde q$。\znote{308}{两处“显然”值得展开。其一，$m\ge5$ 是必需的：(20.8) 中出现 $(3,4),(4,5),(5,4),(4,3)$，用到顶点 $5$，故 $m\ge5$ 时六项才都有意义。其二，(20.9) 的右端是平方，因此把同一式用于逆微分同胚便得到 $x^2=-x^2$，特征不为 $2$ 时 $x=0$；即 $q$ 在形式同构下不变，$q\ne0$ 就意味着非形式。注意 (20.9) 只用到 $\mu^3=0$ 与 $\Phi^1=\Id$ 两个条件，故它对任意 $m\ge5$ 的 $\A$ 都成立；“$q\ne0$ 则不存在到平凡结构的形式微分同胚”正是非形式的定义。[来源：A，原书引理 20.1 及其证明。]}

现在将此结论应用于给定的 $\A$，并假设它是形式的。那么，可以找到上述类型的形式微分同胚 $\Phi$，使 $\widetilde{\A}$ 为平凡的 $\Ai$ 结构，因此 $q$ 必须为零。
\end{proof}

\subsection{一个 Hochschild 上同调计算}\label{sec:20d}
我们暂不直接对 $A$ 上的 $\Ai$ 结构进行分类，而先转向一个稍大的代数 $\overline A$。它是扩充箭图——比原来多一个顶点——
\SourceEquationDiagram{diagram-20-10}{0.415384168}{20.10}
模去关系
\[
\begin{gathered}
(i+2,i+1,i)=0,\qquad(i-2,i-1,i)=0,\\
(i,i+1,i)+(i,i-1,i)=0,\qquad(0,1,0)=0
\end{gathered}
\]
所得的商。由构造，$\overline A$ 是 $\overline R=\K^{m+1}$ 上的线性代数。或许令人意外的是，这个扩充代数比 $A$ 具有更好的性质。它的上同调维数有限，而且还是 Koszul 代数。下面解释后一性质的含义。

由于 $\overline A$ 是路径代数模去一个关于次数和路径长度都齐次的理想所得的商，它同样具有双分次。因此，Ext 代数
\begin{equation}\label{eq:20.11}
\operatorname{Ext}_{\overline A}(\overline R,\overline R)
\end{equation}
\SourcePage{317}
先验地带有三重分次：上同调次数、由 $\overline A$ 中元素的次数导出的内部分次，以及同样由 $\overline A$ 中相应分次导出的长度分次。考察计算\eref{20.11}所用的约化条形构造，显然，上链的上同调分次总不超过其长度分次。此外，可以显式描述上链复形中这两个分次相等的部分。结果如下。令 $\overline A^*$ 为 $\overline A$ 的二次对偶；其生成元 $(i,i\pm1)^*$ 与 $(i\pm1,i)$ 对偶，次数为
\[
|(i,i\pm1)^*|=1-|(i\pm1,i)|,
\]
并满足关系
\[
(i,i+1,i)^*+(i,i-1,i)^*=0,\qquad(m,m-1,m)^*=0.
\]
于是，将条形上链限制到生成元 $(i\pm1,i)$ 上，便定义了一个典范同态
\begin{equation}\label{eq:20.12}
\operatorname{Ext}_{\overline A}(\overline R,\overline R)\longrightarrow\overline A^*.
\end{equation}
在本例中，由\bcite{85}，命题~2.1 给出的单模的投射分解可知，Ext 代数的三重分次部分地集中在对角线上，即所有非零分次部分的上同调分次都等于长度分次。这恰好是\bcite{110}意义下的 Koszul 代数性质；这里稍作推广，以半单环 $\overline R$ 而非域作为基环。该性质等价于\eref{20.12}为同构。\znote{309}{“三重分次部分地对角”指：Ext 代数中每个非零分次部分都满足上同调次数等于长度分次。这与上文“上链的上同调次数总不超过长度分次”合起来，说明整个 Ext 代数由对角部分决定，于是 (20.12) 是同构就等价于 Koszul 性。此处基环取半单环 $\overline R=\K^{m+1}$ 而非域，所以作者说这是\bcite{110}中定义的轻微推广。[来源：A，原书本节与所引文献。]}

接着，考虑 $\overline A$ 的约化 Hochschild 上链复形，它由从 $\overline A/\overline R$ 到 $\overline A$ 的多线性映射组成。可按映射的目标中的路径长度过滤这个复形。所得谱序列的 $E_1$ 页可写成
\begin{equation}\label{eq:20.13}
\operatorname{Ext}_{\overline A}(\overline R,\overline R)
\otimes_{\overline R\otimes\overline R}\overline A
\cong\overline A\otimes_{\overline R\otimes\overline R}\overline A^*.
\end{equation}
由于 $\overline A$ 及其对偶均为有限维，这个谱序列显然收敛到 $\HH(\overline A,\overline A)$。从图形上看，$E_1$ 页的生成元可以视为由两段组成的闭路径，两段分别位于箭图\eref{20.10}及其对偶中。具体地，它们可写成
\begin{equation}\label{eq:20.14}
(i_l,\ldots,i_0)\otimes(j_0,\ldots,j_r)^*,\qquad j_r=i_l,\quad j_0=i_0.
\end{equation}
分次问题可以如下理清。假设我们关心 $\HH^{r+s}(\overline A,\overline A)^s$，它是由内部次数为 $s$ 的 $r$ 线性 Hochschild 上闭链表示的群；这里的记号与\secref{1f}相同。那么，\eref{20.13}中对它有贡献的项，是\eref{20.14}中具有给定 $r$ 并满足
\[
|(j_0,\ldots,j_r)^*|+|(i_l,\ldots,i_0)|=r+s
\]
的那些表达式。由箭图的分次方式，得到 $l=2s+r$。
\StatementNumber{2}\begin{lemma}\label{st:20.2}
当 $r\ge5$ 或 $r=3$ 时，$\HH^2(\overline A,\overline A)^{2-r}$ 为零；当 $r=4$ 时，它为 $\K$。
\end{lemma}
\SourcePage{318}
\begin{proof}[证明提要]
当 $r\ge5$ 时，结论直接成立，因为有贡献的项\eref{20.14}必须满足 $l=4-r<0$。$r=3,4$ 的情形略为复杂；这里略去细节，但记下一个稍后要用的推论。\eref{20.13}的边缘同态给出映射
\begin{equation}\label{eq:20.15}
\HH^2(\overline A,\overline A)^{-2}\longrightarrow
\overline R\otimes_{\overline R\otimes\overline R}(\overline A^*)^2.
\end{equation}
这里，目标空间由 $\overline A^*$ 中次数为二的元素组成，$\overline R$ 的幂等元在这些元素上对角地作用，即左、右作用相同。由计算可知，\eref{20.15}是单射，其像由下列元素的倍数组成；该元素属于 $\overline A^*$ 的中心：
\begin{equation}\label{eq:20.16}
\begin{aligned}
&(0,1,2,1,0)^*+(1,2,3,2,1)^*+\cdots\\
&\hspace{4em}+(m-2,m-1,m,m-1,m-2)^*.
\end{aligned}
\end{equation}
\end{proof}

结合一般分类理论，\stref{引理}{20.2}表明，在固定上同调的拟同构意义下，满足 $H(\overline{\A})\cong\overline A$ 的 $\Ai$ 范畴 $\overline{\A}$ 至多构成一个单参数族。将参数记为 $q\in\K$；$q=0$ 对应形式的 $\Ai$ 结构。单凭\stref{引理}{20.2}本身，尚不能证明对应于 $q\ne0$ 的非形式结构存在。实际上，不同的非零参数值可通过对各个分次部分重新缩放而联系起来，所以它们并无实质区别。不过，暂且假设这样的结构存在，那么可以安排它们的第一个非平凡高阶乘积为 $\mu^4$。此外，由\eref{20.16}可知，当 $m\ge5$ 时，这一乘积满足\eref{20.8}的类似等式，且 $q\ne0$。特别地，此时若删去额外顶点对应的对象，由\stref{引理}{20.1}，所得全 $\Ai$ 子范畴仍不是形式的。

\subsection{应用}\label{sec:20e}
我们沿用高维情形中的策略。为二重分歧覆盖 $\pi:X=X_m^1\to\CC$ 选取消失路径基，考察 $\F(\pi)$ 中相应的全 $\Ai$ 子范畴，将其记为 $\B^\to$。显式地，它的对象为 $(Z_0,\ldots,Z_m)$，态射空间为
\[
\hom_{\B^\to}(Z_i,Z_j)=\begin{cases}
\K\cdot c_{ji}\oplus\K\cdot d_{ji},&i<j,\\
\K\cdot e_{ii},&i=j,\\
0,&i>j,
\end{cases}
\]
其中所有次数均为零。$e_{ii}$ 是严格单位元，除此之外，有
\[
\begin{aligned}
\mu_{\B^\to}^2(c_{kj},c_{ji})&=c_{ki},&
\mu_{\B^\to}^2(d_{kj},d_{ji})&=d_{ki},\\
\mu_{\B^\to}^2(d_{kj},c_{ji})&=0,&
\mu_{\B^\to}^2(c_{kj},d_{ji})&=0.
\end{aligned}
\]
\SourcePage{319}
这与高维情形中的对应表达式\eref{20.5}存在差异，归根结底是因为 $S^0$ 的上同调环与高维球面的上同调环有不同的性质。取扭曲复形 $(C_0,\ldots,C_m)$，定义为
\begin{equation}\label{eq:20.17}
C_0=Z_0,\qquad C_j=\Cone(c_{j,j-1}\pm d_{j,j-1}),\quad j>0,
\end{equation}
其中对每个 $j$ 都可自由选择符号；这一歧义将在下面的\stref{笔记}{20.5}中讨论。令 $\overline{\A}$ 为 $\Tw\B^\to$ 中由这些对象构成的全 $\Ai$ 子范畴。容易算得 $H(\overline{\A})\cong\overline A$。
\StatementNumber{3}\begin{lemma}\label{st:20.3}
$\overline{\A}$ 不是形式的。
\end{lemma}
\begin{proof}
显然，只须证明 $m=2$ 的情形，因为其余所有 $\overline{\A}$ 都包含它作为全 $\Ai$ 子范畴。由于集合\eref{20.17}生成 $\Tw(\B^\to)$，$\overline{\A}$ 上与 $\B^\to$ 上的扭曲复形 $\Ai$ 范畴拟等价，进而它们的 Hochschild 上同调群同构。对于代数，这类结果可追溯到\bcite{71}；证明稍作调整便适用于 $\Ai$ 情形。直接计算得到
\[
\HH(\B^\to,\B^\to)=
\underbrace{\K}_{\text{次数为零}}\oplus
\underbrace{\K^2}_{\text{次数为一}}.
\]
这与\stref{引理}{20.2}中部分算出的 $\overline A$ 的 Hochschild 上同调不同，由此得到所需的非形式性结论。
\end{proof}

令 $\A$ 为 $\overline{\A}$ 中以 $(C_1,\ldots,C_m)$ 为对象的全 $\Ai$ 子范畴。由前面的讨论，它也不是形式的，至少当 $m\ge5$ 时如此。另一方面，将\stref{命题}{18.21}适当推广，见\stref{笔记}{20.5}，可知 $\A$ 拟同构于 $\F$。因此：
\StatementNumber{4}\begin{corollary}\label{st:20.4}
对任意 $m\ge5$，$\F=\F_m^1$ 都不是形式的。
\end{corollary}
\StatementNumber{5}\begin{note}\label{st:20.5}
我们在二维情形中使用了\stref{命题}{18.21}，此时 Lefschetz 纤维化退化为二重分歧覆盖。虽然原先的讨论并未涵盖这种情形，但很容易调整其证明，也可改用更初等的论证。不过，写出精确陈述时须稍加谨慎；\stref{笔记}{19.10}已简要提到这一点。态射 $c_{j,j-1}$ 与 $d_{j,j-1}$ 对应于 Lefschetz 顶针之间的交点。相应的陈述是，匹配圈为 $c_{j,j-1}\pm d_{j,j-1}$ 的映射锥，而符号的选择决定匹配圈所带的 Pin 结构。在我们的情形中，匹配圈是在 $X$ 上构成 $(A_m)$ 构型的曲线；原先假设它们带有非平凡 Pin 结构，因为我们曾应用\stref{推论}{17.17}证明它们是分裂生成元。所幸，这些曲线恰好在
\SourcePage{320}
$H_1(X;\Z/2)$ 中线性无关，因而任意两组 Pin 结构的选择，都可通过用 $X$ 上的实线丛作扭曲而相互转换；按\secref{12i}的语言，这就是膜自同构。在代数上，这反映为 $\B^\to$ 具有改变 $c_{j,j-1}\pm d_{j,j-1}$ 中符号的自同构。因此，\eref{20.17}中符号的选择无关紧要。
\end{note}

\subsection{辫群与映射类群}\label{sec:20f}
现在专门考察最简单的非形式情形 $n=1$、$m=5$，并换一个视角来理解其中的情形，从而说明高阶复合映射的几何意义。将等价\eref{20.3}与 Yoneda 嵌入结合，得到全忠实 $\Ai$ 函子
\begin{equation}\label{eq:20.18}
\F(X)\longrightarrow\Q=\modcat(\F).
\end{equation}
它把正合 Lagrangian 膜 $L^\#$ 映到 $\Ai$ 模 $\M$，其上同调为
\[
M=H(\M)=\bigoplus_i HF^*(V_i^\#,L^\#),
\]
将其视为 $A=\bigoplus_{i,j}HF^*(V_i^\#,V_j^\#)$ 上的模。我们特别关心\sr{fig:20.2}{图~20.2}所示的曲线 $L$。它的 Maslov 数为零，并且可使它正合，因此它成为一个正合 Lagrangian 膜。适当选择分次后，对应的模为
\[
M=U_3\oplus SU_3,
\]
其中 $U_i=(i)R\cong\K$ 表示初等单模，即将 $R$ 视为右 $A$ 模时的各直和项；$S$ 像往常一样表示平移。
\SourceFigure{figure-20-2}{0.458999506}{20.2}
接下来更仔细地考察单模 $U_3$ 的同调代数。以 $P_i=(i)A$ 表示初等投射模。存在无上同调复形
\[
\begin{aligned}
0\longrightarrow S^{-3}U_3\longrightarrow S^{-2}P_3
&\longrightarrow S^{-1}P_2\oplus S^{-2}P_4\\
&\longrightarrow P_1\oplus S^{-1}P_3\oplus S^{-2}P_5\\
&\longrightarrow P_2\oplus S^{-1}P_4\longrightarrow P_3\longrightarrow U_3\longrightarrow0,
\end{aligned}
\]
\SourcePage{321}
可将它首尾接续，得到 $U_3$ 的无限投射分解。由此计算得到下列结果；其中的分次记号与 Hochschild 上同调群的记号类似：
\begin{equation}\label{eq:20.19}
\begin{aligned}
&\operatorname{Ext}_{A}^{r+s}(U_3,U_3)^s=0,
&&r+s<0;\\
&\operatorname{Ext}_{A}^{0}(U_3,U_3)^0=\K,\qquad
\operatorname{Ext}_{A}^{0}(U_3,U_3)^s=0,
&&s\ne0;\\
&\operatorname{Ext}_{A}^{1}(U_3,U_3)^{-1}=\K,\qquad
\operatorname{Ext}_{A}^{1}(U_3,U_3)^s=0,
&&s\ne-1.
\end{aligned}
\end{equation}
最后，在已知这些结果的基础上，使用条形分解，或者与扩充代数 $\overline A$ 的对应结果比较，还可看出乘积
\begin{equation}\label{eq:20.20}
\operatorname{Ext}_{A}^{1}(U_3,U_3)^{-1}\otimes
\operatorname{Ext}_{A}^{1}(U_3,U_3)^{-1}
\longrightarrow\operatorname{Ext}_{A}^{2}(U_3,U_3)^{-2}
\end{equation}
非零。这些计算的意义如下。令 $\widetilde{\A}$ 为 $A$ 上的平凡 $\Ai$ 结构，并令 $\widetilde{\Q}=\modcat(\widetilde{\A})$ 为其 $\Ai$ 模范畴。对于上同调模 $M$ 固定的 $\widetilde{\A}$ 模 $\widetilde{\M}$，有一种一般的分类方法，在 Kodaira--Spencer 形变理论的意义下由 $\operatorname{Ext}_{A}(M,M)$ 控制。在此背景下，\eref{20.19}和\eref{20.20}描述了一种一阶形变空间为一维、且所有这些一阶形变在二阶均受阻的理论。\znote{313}{把 (20.19)(20.20) 读成形变论：$\operatorname{Ext}^1_A(U_3,U_3)^{-1}=\K$ 是一维的一阶形变空间，(20.20) 非零则说明两个一阶形变之积给出非零的二阶障碍，即每个一阶形变都不可积。于是 $U_3$ 上只存在平凡的 $\widetilde{\A}$ 模结构，这正是\stref{引理}{20.6}的出发点。符号的配位也相容：形变理论中一阶形变住在 $\operatorname{Ext}^1$，障碍住在 $\operatorname{Ext}^2$，而此处内部分次分别为 $-1$ 与 $-2$。[来源：A，原书本节。]}因此，在 $H^0(\widetilde{\Q})$ 的同构意义下，将 $U_3$ 赋予 $\widetilde{\A}$ 模结构 $\widetilde U_3$ 的方式只有一种，即平凡的方式。使用同一形式体系，通过略为复杂的论证，还可得到下列结论；这里略去其细节。
\StatementNumber{6}\begin{lemma}\label{st:20.6}
设 $\widetilde{\M}$ 为 $\widetilde{\A}$ 上的 $\Ai$ 模，其底层上同调模为 $H(\widetilde{\M})\cong U_3\oplus SU_3$。那么，在 $H^0(\widetilde{\Q})$ 中有
\[
\widetilde{\M}\cong\ACone(t),
\]
其中 $t\in\hom_{\widetilde{\Q}}^0(\ModShift^{-1}\widetilde U_3,\ModShift\widetilde U_3)$ 表示 $\operatorname{Ext}_{A}^{2}(U_3,U_3)$ 中的一个元素。因此，
\[
H^{-1}\bigl(\hom_{\widetilde{\Q}}(\widetilde{\M},\widetilde{\M})\bigr)\ne0.
\]
\end{lemma}

回到几何，暂且假设 $\F$ 是形式的，此时 $\Q$ 与 $\widetilde{\Q}$ 拟等价。于是，\sr{fig:20.2}{图~20.2}中的曲线 $L$ 会给出一个适用\stref{引理}{20.6}的 $\Ai$ 模。另一方面，由于\eref{20.18}是忠实的，这个模的自同态应恢复 $HF^*(L^\#,L^\#)\cong H^*(L;\K)$，而它在 $-1$ 次为零，产生矛盾。这给出了\stref{推论}{20.4}的另一种证明；这里证明了 $m=5$ 的情形，一般情形可直接由此推出。

$\Q$ 和 $\widetilde{\Q}$ 都带有辫群 $\operatorname{Br}_6$ 的作用，该作用由对应于基本球面对象的扭转函子给出。不过，在几何情形 $\Q\cong\modcat(\F)$ 中，这一作用可嵌入一个更大的群的作用，即模去同痕的 $X$ 的分次
\SourcePage{322}
正合辛自同构群。沿 $L$ 的 Dehn 扭转就是这个更大的群中不属于 $\operatorname{Br}_6$ 的一个元素；这是因为 $L$ 本身的同痕类在 $\pi:X\to\CC$ 的覆盖对合下不保持不变。直观地说，可以认为 $\F$ 的非形式性与这个更大的对称群相联系。
\clearpage
\StatementNumber{7}\begin{note}\label{st:20.7}
上面的讨论中出现的辫群与映射类群之间的关系，是经典低维拓扑中已得到充分理解的一个课题\bcite{23}。较近的研究表明，广义辫群的情形更为复杂\bcite{142}。值得从 Fukaya 范畴的角度简要考察这一问题。

考虑 $(E_6)$ 奇点的 Milnor 纤维，即 $\{x_0^3+x_1^4=1\}\subset\CC^2$。在拓扑上，它是具有一个穿孔的三亏格曲面，并带有构成 $(E_6)$ 构型的一组消失圈基。将它们记为 $(V_1,\ldots,V_6)$，编号采用\bcite{142}，图~2 所示的次序，尽管这不是它们在该基中出现的自然次序。令 $\A$ 为该 Milnor 纤维的 Fukaya 范畴中以 $V_i^\#$ 为对象的全 $\Ai$ 子范畴。在上同调层次，$A=H(\A)$ 可表示为适当分次的 $(E_6)$ Dynkin 箭图的路径代数的商，其含义与之前相同。令
\[
L=\tau_{V_4}\tau_{V_3}\tau_{V_5}\tau_{V_4}\tau_{V_2}\tau_{V_6}\tau_{V_3}\tau_{V_5}\tau_{V_4}(V_1).
\]
经过适当同痕后，$L$ 与 $V_1$ 恰有一个交点，由此得到两条曲线相应的 Dehn 扭转之间的辫关系。如\bcite{142}所指出，$(E_6)$ 型广义辫群中不存在对应的关系，因此，从该辫群到 Milnor 纤维的映射类群的自然同态不是单射。\znote{314}{这里的机制是：$L$ 由这一串 Dehn 扭转作用于 $V_1$ 得到，适当同痕后与 $V_1$ 恰交一点，于是 $\tau_L\tau_{V_1}$ 与 $\tau_{V_1}\tau_L$ 之间出现了通常的辫关系；而\bcite{142}指出 $(E_6)$ 型广义辫群中没有对应关系，故该同态不单射。后半段的计算把这一失效搬到 $\Tw\widetilde{\A}$ 上：若 $\A$ 形式，则 $H(\hom(Y_1,C))$ 应与 $HF^*(V_1,L)$ 同为一维，而实算为三维（次数 $-2,-1,0$）。[来源：A，原书\stref{笔记}{20.7}与所引文献。]}

令 $\widetilde{\A}$ 为 $A$ 上的平凡 $\Ai$ 结构。在 $\Tw\widetilde{\A}$ 中考虑扭曲复形
\[
C=T_{Y_4}T_{Y_3}T_{Y_5}T_{Y_4}T_{Y_2}T_{Y_6}T_{Y_3}T_{Y_5}T_{Y_4}(Y_1),
\]
其中 $Y_i$ 为 $\widetilde{\A}$ 本身的对象。直接计算可知，$H\bigl(\hom_{\Tw\widetilde{\A}}(Y_1,C)\bigr)$ 为三维，生成元的次数分别为 $-2,-1,0$。

若 $\A$ 拟同构于 $\widetilde{\A}$，则可用\stref{推论}{17.17}将这个 $H(\hom)$ 群与 $(V_1,L)$ 的 Floer 上同调等同。但我们知道，后者是一维的，因此 $\A$ 不可能是形式的。由此可见，$\A$ 中的高阶乘积会导致扭曲复形之间的同态发生抵消，而这些抵消似乎又是 $(E_6)$ 辫群作用不忠实的原因。我们说“似乎”，是因为尚不知道该群是否忠实地作用于 $\Tw\widetilde{\A}$。
\end{note}

\backmatter
\theendnotes
\chapter*{读书报告}
\addcontentsline{toc}{chapter}{读书报告}
\markboth{读书报告}{读书报告}
\begin{center}
\large 从辛几何到可计算的 $A_\infty$ 范畴
\end{center}

\begin{quote}
\small 本报告基于全书通读与本次校订工作写成。文中页码均为原书印刷页码；定理编号从原书。书末译者注共 187 条，可与本报告互参。
\end{quote}

\section*{一、这本书要解决什么问题}

Fukaya 范畴把辛流形的几何编码成一个 $A_\infty$ 范畴：对象是 Lagrangian 子流形，态射是 Floer 上链群，复合映射由伪全纯圆盘的计数给出。这个定义在概念上很自然，在计算上却几乎无法直接使用——对象本身就无法分类，更不用说数清楚所有圆盘。Seidel 这本书的目标，不是再讲一遍定义，而是回答一个更实际的问题：\emph{在什么条件下，Fukaya 范畴可以被约化为有限、可计算的数据？}

他的答案是：把辛流形限制到正合（exact）且 Calabi--Yau 的仿射代数簇，用 Picard--Lefschetz 理论把几何分解为消失圈与辫单值变换，再用 $A_\infty$ 范畴的代数形式（扭曲复形、分裂闭包、变异）把归纳计算组织起来。书中的定理 A、B 与推论 C（原书 pp.~4--5）是这条路线的主定理；全书三章分别准备代数、分析、几何三套工具。

\section*{二、三章的结构}

\textbf{第一章（§1--6）是纯同调代数。}它从 $A_\infty$ 范畴的定义讲起，重点不在铺陈理论，而在几件后文真正用得到的事：恒等态射在 $A_\infty$ 层次并不自动存在；映射锥与正合三角如何构造；幂等元如何形式地分裂出直和项；以及“沿一个对象的扭转”如何实现变异。作者明确说，这一章只详细讨论恒等态射，因为它在 Fukaya 范畴中恰恰最棘手。章末引入 $\Z/2$ 等变结构，为第三章的二重覆盖做准备。

\textbf{第二章（§7--14）是 Floer 理论。}它先以一点辛几何（§7）作玩具模型，再回顾经典 Floer 上同调（§8），然后逐步给出正合 Fukaya 范畴的“初步版本”（§9）与“完整版本”（§12）。中间三节分别处理基本性质（§10）、指标与行列式线（§11）、曲面上的多边形（§13）。分次（grading）与 Pin 结构是这一章的技术核心：前者把 Floer 群变成 $\Z$ 分次对象，后者决定符号，因而只在 $\operatorname{char}\K\ne2$ 时必需。§14 讨论带辛对合的流形及其等变 Fukaya 范畴，是第三章定义 $\F(\pi)$ 的直接工具。

\textbf{第三章（§15--20）是 Picard--Lefschetz 理论。}它从 Lefschetz 纤维化的初等几何开始（§15），引入消失圈、Lefschetz 顶针与匹配圈（§16），用伪全纯截面证明扭转函子与辛 Dehn 扭转的相容性（§17），再借二重覆盖定义 Lefschetz 纤维化的 Fukaya 范畴 $\F(\pi)$（§18）。§19 回到仿射簇，证明引言中的三个定理；§20 用 $(A_m)$ 型 Milnor 纤维作例子，展示这套机器如何真正算出东西。

\section*{三、第一章：把“算”变成“形式化”}

$A_\infty$ 范畴的定义是一组多重线性映射 $\mu^d$，满足带符号的二次关系。与 dg 范畴不同，它不要求 $\mu^1$ 与 $\mu^2$ 先满足结合律；结合律只在同调层次成立。这带来两个后果。其一，恒等态射通常不是严格存在的，只能在一个拟同构的模型里找到；Seidel 用同伦扰动（homological perturbation）处理这一点，并反复提醒读者：在 $A_\infty$ 层次做计算时，不能随意把恒等态射当成真的单位元。其二，映射锥不能按朴素方式定义，必须引入扭曲复形。

扭曲复形的对象是一族对象 $X_i$ 加上满足 Maurer--Cartan 方程的矩阵 $\delta=(\delta_{ij})$，$\delta_{ij}\in\hom^1(X_i,X_j)$。它的上同调 $H^0(\Tw\A)$ 是三角范畴；这一步把“映射锥”“正合三角”从链复形的直觉搬到了 $A_\infty$ 世界。接下来是分裂闭包：把每个幂等自同态的像形式地加进来，得到 $\Pi(\Tw\A)$。作者强调，分裂生成（split-generation）是很弱的要求——它甚至不保证 $K_0$ 由这些生成元的类张成——但正因为它弱，才可能在几何中实现。

扭转（twist）是本章与几何的接口。给定对象 $E$，沿 $E$ 的扭转 $T_E$ 在三角范畴中对应“变异”（mutation）；在 $A_\infty$ 层次，它由一串带符号的复合映射给出。§5 证明 $T_E$ 是 $A_\infty$ 函子，并在 $E$ 为球面对象时给出逆。这条代数线索在第三章会变成：沿消失圈的扭转等于沿匹配圈的 Dehn 扭转。

\section*{四、第二章：正合 Floer 理论的技术骨架}

第二章的出发点是正合辛流形 $(M,\omega=d\theta)$。正合性排除了闭流形，但作者施加凸性（convexity）条件，使有限面积的伪全纯曲线不会跑到“无穷远”。在这个框架里，正合 Lagrangian 子流形 $L$（即 $\theta|_L$ 恰当）没有障碍（unobstructed），每个都给出非零对象，且自同态代数就是经典上同调 $H^*(L)$，没有量子修正。这使 Fukaya 范畴摆脱 Novikov 参数，代价是失去闭流形上的许多例子。

分次是第二个技术支柱。Kontsevich 的分次 Lagrangian 是指一个实值函数 $\alpha_L:L\to\R$，使 $e^{2\pi i\alpha_L}$ 等于某个行列式线的相位；这要求 $2c_1(M)=0$。分次把 Floer 群变成 $\Z$ 分次向量空间，使映射锥、扭转、变异都能逐次进行。Pin 结构则负责符号：伪全纯圆盘模空间的定向由 $L$ 的 Pin 结构决定，若 $\operatorname{char}\K=2$，符号消失，Pin 结构可以省略。

完整定义（§12）用“带边界标记点的圆盘”统一 $\mu^d$：$d$ 个边界标记点把圆周分成 $d$ 段，每段落在一个 Lagrangian 上；模空间是这些圆盘的等价类。这里的技术难点是横截性（transversality）与紧性（compactness）；作者用零阶非齐次项扰动方程，避免虚基本链，但把分析细节压到最低，只给证明概要。§13 讨论曲面：此时伪全纯曲线可以用多边形组合地描述，Floer 理论退化为组合计数；这是全书唯一一处能把分析完全拿掉的地方。

\section*{五、第三章：从几何到辫群，再回到代数}

Picard--Lefschetz 理论的核心图像是：一个 Lefschetz 纤维化 $\pi:E\to\C$ 的临界值绕一圈，纤维上的消失圈就经历一次 Dehn 扭转。把这些扭转按临界值的顺序复合起来，得到辫单值变换（braid monodromy）。Seidel 的观察是：\emph{临界值的 Hurwitz 移动（交换两个临界值的顺序）对应消失圈变异}。这正是定理 17.20（原书 p.~263）的内容，也是“为什么 Fukaya 范畴可以按维数归纳计算”的几何解释；定理 18.24（原书 p.~287）与推论 18.27（原书 p.~291）把它组织成可用的代数形式。

为了把这个观察变成定理，需要两件工具。第一是伪全纯截面：Lefschetz 纤维化的全纯截面在纤维上留下一个点，绕临界值时这个点沿着消失圈移动；计数经过两个截面的伪全纯截面，就得到 Floer 上同调中的扭转映射。第二是 Lefschetz 纤维化的 Fukaya 范畴 $\F(\pi)$：取沿一个光滑纤维分歧的二重覆盖，用 $\Z/2$ 等变构造把纤维与总空间的 Fukaya 范畴联系起来。这个定义不自然，但可以让第二章的等变技术直接复用。

§19 把上述工具用于仿射簇 $X\subset\C^N$ 及其超平面截面 $X^1$。对 $X^1$，消失圈给出自然的一族对象；对 $X$，则用 Lefschetz 束的匹配圈。组合数据（临界值、辫、消失圈的交叠数）编码成一个有限的、有向的 $A_\infty$ 范畴 $\mathcal B$。定理说：$\Pi(\Tw\F(X^1))$ 与 $\Pi(\Tw\mathcal B)$ 拟等价（定理 A）；$\Tw\mathcal B$ 含有与 $\F(X)$ 拟等价的全子范畴（定理 B）；因此 $\F(X)$ 与一个 hom 空间有限维的 dg 范畴拟同构（推论 C）。条件 $d\ne2$、$n\ge1$、$\operatorname{char}\K\ne2$ 都是这条证明路线的技术限制，作者在引言中逐条说明。

这里有一处原书自身的不统一：定理 A 印作 $n\ge1$，但紧接着的解释段（p.~5）与定理 B 之后的说明都按 $n>1$ 使用，并说 $n=1$ 的情形“只是为方便起见”才排除。译文从原书定理陈述，取 $n\ge1$，并在译者注中登记；第 19 节的归纳证明本身也是按 $n\ge1$ 陈述的。

\section*{六、贡献、局限与后续}

这本书的贡献不在于发明 Fukaya 范畴，而在于\emph{把“计算”变成一个有结构的程序}。它给出了从辛几何到组合数据的清晰路线：正合性消去分析上的无穷远，分次与 Pin 结构消去符号歧义，扭曲复形与分裂闭包消去 $A_\infty$ 层次的技术麻烦，PL 理论则把几何归纳为辫群。Hurwitz 移动等于变异这一定理，后来成为镜像对称中“例外集合变异”与“球面对象”理论的基石之一。

局限也很明显。第一，全书只处理正合、Calabi--Yau 的仿射簇；闭流形、一般辛流形不在范围内。第二，分析细节被大幅压缩，读者若想真正补全横截性与紧性，必须回到 Fukaya--Oh--Ohta--Ono 的巨著。第三，条件 $\operatorname{char}\K\ne2$ 与 $d\ne2$ 是证明的技术产物，作者自己说前者“很可能”可以去掉，后者则更本质。第四，除了 $(A_m)$ Milnor 纤维，书中几乎没有例子；计算复杂度在实际中仍然很高。

从今天的视角看，这本书是“Fukaya 范畴的可计算性”这一方向的早期纲领。它之后的 wrapped Fukaya 范畴、family Floer 理论、microlocal sheaf 与 quiver-with-potential 的工作，都可以看作沿着它提出的问题继续走。读这本书的正确方式，不是把它当入门教材逐页啃，而是把它当一张施工图：先看引言与 §19 的定理陈述，再按需回到 §1--6（代数）、§8--12（分析）、§15--18（几何）取工具。

\section*{七、阅读建议}

\begin{itemize}
\item 第一次读：引言（原书 pp.~1--7）$\to$ §19 的定理与证明梗概 $\to$ §20 的例子。先建立“为什么可以归纳计算”的直觉。
\item 第二次读：§1--5（$A_\infty$ 形式化）$\to$ §8--12（Floer 理论）$\to$ §15--18（PL 理论）。每一步都对照书末的符号表与索引。
\item 公式核对：本书的指标与符号约定很密，建议把 §11（指标与行列式线）与 §17（伪全纯截面）的公式单独抄一遍，确认 $(-1)$ 符号与分次移位。
\item 配合阅读：作者在引言中反复引用的 [60]（Fukaya--Oh--Ohta--Ono）是分析细节的正式来源；[131]（Seidel 的 Lefschetz fibration 论文）是第三章的直接前身。
\end{itemize}

% 本批前向引用均已在累计正文中落地，无待译定位条目。

\chapter*{参考文献}
\markboth{参考文献}{参考文献}
\addcontentsline{toc}{chapter}{参考文献}
\begin{list}{}{\setlength{\leftmargin}{3em}\setlength{\labelwidth}{2.5em}\setlength{\itemsep}{.5em}\small}
\item[\hypertarget{bib:1}{[1]}] M. Akveld and D. Salamon, Loops of Lagrangian submanifolds and pseudo-holomorphic discs, Geom. Funct. Anal. 11 (2001), 609–650. \BibBackref{1}
\item[\hypertarget{bib:2}{[2]}] V. I. Arnol'd, On a characteristic class entering into conditions of quantization, Funct. Anal. Appl. 1 (1967), 1–13. \BibBackref{2}
\item[\hypertarget{bib:3}{[3]}] —–, Some remarks on symplectic monodromy of Milnor fibrations, in The Floer Memorial Volume (H. Hofer, C. Taubes, A. Weinstein, and E. Zehnder, eds.), Progr. Math. 133, Birkhäuser, 1995, pp. 99–104. \BibBackref{3}
\item[\hypertarget{bib:4}{[4]}] V. Arnol'd, V. Goryunov, O. Lyashko, and V. Vasiliev, Singularity theory. I, Encyclopaedia Math. Sci. 6, Springer-Verlag, 1998. \BibBackref{4}
\item[\hypertarget{bib:5}{[5]}] V. I. Arnol'd, S. M. Gusein-Zade, and A. N. Varchenko, Singularities of differentiable maps, Birkhäuser, 1988. \BibBackref{5}
\item[\hypertarget{bib:6}{[6]}] P. S. Aspinwall, Some navigation rules for D-brane monodromy, J. Math. Phys. 42 (2001), 5534–5552. \BibBackref{6}
\item[\hypertarget{bib:7}{[7]}] M. Atiyah, Topological quantum field theories, Inst. Hautes Études Sci. Publ. Math. 68 (1989), 175–186. \BibBackref{7}
\item[\hypertarget{bib:8}{[8]}] —–, On analytic surfaces with double points, Proc. Royal Soc. Ser. A 247 (1958), 237–244. \BibBackref{8}
\item[\hypertarget{bib:9}{[9]}] M. Atiyah and I. Singer, The index of elliptic operators. V, Ann. of Math. 93 (1971), 139–149. \BibBackref{9}
\item[\hypertarget{bib:10}{[10]}] M. Audin and J. Lafontaine (eds.), Holomorphic curves in symplectic geometry, Progr. Math. 117, Birkhäuser, 1994. \BibBackref{10}
\item[\hypertarget{bib:11}{[11]}] D. Auroux, Monodromy invariants in symplectic topology, Preprint math.SG/0304113, 2003. \BibBackref{11}
\item[\hypertarget{bib:12}{[12]}] —–, Symplectic 4-manifolds as branched coverings of $\mathbb{CP}^2$, Invent. Math. 139 (2000), 551–602. \BibBackref{12}
\item[\hypertarget{bib:13}{[13]}] —–, Symplectic maps to projective spaces and symplectic invariants, Turkish J. Math. 25 (2001), 1–42. \BibBackref{13}
\item[\hypertarget{bib:14}{[14]}] D. Auroux, L. Katzarkov, and D. Orlov, Mirror symmetry for weighted projective planes and their noncommutative deformations, Preprint math.AG/0404281, 2004. \BibBackref{14}
\item[\hypertarget{bib:15}{[15]}] —–, Mirror symmetry for Del Pezzo surfaces: Vanishing cycles and coherent sheaves, Invent. Math. 166 (2006), 537–582. \BibBackref{15}
\item[\hypertarget{bib:16}{[16]}] D. Auroux, V. Munoz, and F. Presas, Lagrangian submanifolds and Lefschetz pencils, J. Symplectic Geom. 3 (2005), 171–219. \BibBackref{16}
\item[\hypertarget{bib:17}{[17]}] D. Auroux and I. Smith, Lefschetz pencils, branched covers and symplectic invariants, Preprint math.SG/0401021, 2004. \BibBackref{17}
\item[\hypertarget{bib:18}{[18]}] P. Balmer and M. Schlichting, Idempotent completion of triangulated categories, J. Algebra 236 (2001), 819–834. \BibBackref{18}
\item[\hypertarget{bib:19}{[19]}] A. Beilinson, Coherent sheaves on $\mathbb{P}^n$ and problems in linear algebra, Functional Anal. Appl. 12 (1978), 214–216. \BibBackref{19}
\item[\hypertarget{bib:20}{[20]}] A. Beilinson, J. Bernstein, and P. Deligne, Faisceaux pervers, Analyse et topologie sur les espaces singuliers (I), Astérisque 100, Soc. Math. France, 1983, pp. 7–171. \BibBackref{20}
\item[\hypertarget{bib:21}{[21]}] J. Bernstein and V. Lunts, Equivariant sheaves and functors, Lecture Notes in Math. 1578, Springer-Verlag, 1994. \BibBackref{21}
\item[\hypertarget{bib:22}{[22]}] P. Biran and K. Cieliebak, Lagrangian embeddings into subcritical Stein manifolds, Israel J. Math. 127 (2002), 221–244. \BibBackref{22}
\item[\hypertarget{bib:23}{[23]}] J. Birman and H. Hilden, On the mapping class groups of closed surfaces as covering spaces, in Advances in the theory of Riemann surfaces (L. Ahlfors, L. Bers, H. Farkas, R. Gunning, I. Kra and H. Rauch, eds.), Ann. of Math Stud. 66 (1971), Princeton University Press, pp. 81–115. \BibBackref{23}
\item[\hypertarget{bib:24}{[24]}] M. Bökstedt and A. Neeman, Homotopy limits in triangulated categories, Compositio Math. 86 (1993), 209–234. \BibBackref{24}
\item[\hypertarget{bib:25}{[25]}] A. I. Bondal, Helices, representations of quivers, and Koszul algebras, \hyperlink{bib:117}{[117]}, pp. 75–95. \BibBackref{25}
\item[\hypertarget{bib:26}{[26]}] —–, Representations of associative algebras and coherent sheaves, Math. USSR Izv. 34 (1990), 23–42. \BibBackref{26}
\item[\hypertarget{bib:27}{[27]}] A. I. Bondal and M. M. Kapranov, Representable functors, Serre functors, and mutations, Math. USSR Izv. 35 (1990), 519–541. \BibBackref{27}
\item[\hypertarget{bib:28}{[28]}] —–, Enhanced triangulated categories, Math. USSR Sb. 70 (1991), 93–107. \BibBackref{28}
\item[\hypertarget{bib:29}{[29]}] B. Booss-Bavnbeck and K. Wojciechowski, Elliptic boundary value problems for Dirac operators, Birkhäuser, 1993. \BibBackref{29}
\item[\hypertarget{bib:30}{[30]}] Yu. Chekanov, Differential algebras of Legendrian links, Invent. Math. 150 (2002), 441–483. \BibBackref{30}
\item[\hypertarget{bib:31}{[31]}] C.-H. Cho, Holomorphic discs, spin structures, and the Floer cohomology of the Clifford torus, Internat. Math. Res. Notices 35 (2004), 1803–1843. \BibBackref{31}
\item[\hypertarget{bib:32}{[32]}] C.-H. Cho and Y.-G. Oh, Floer cohomology and disc instantons of Lagrangian torus fibers in toric Fano manifolds, Asian J. Math. 10 (2006), 773–814. \BibBackref{32}
\item[\hypertarget{bib:33}{[33]}] O. Cornea and F. Lalonde, Cluster homology, Electron. Res. Announc. Amer. Math. Soc. 12 (2006), 1–12. \BibBackref{33}
\item[\hypertarget{bib:34}{[34]}] V. de Silva, The index of a family of pseudo-holomorphic disks, Preprint, 2000. \BibBackref{34}
\item[\hypertarget{bib:35}{[35]}] V. de Silva, R. Gautschi, J. Robbin, and D. Salamon, Combinatorial Floer homology, Manuscript, 2001. \BibBackref{35}
\item[\hypertarget{bib:36}{[36]}] P. Deligne et al. (eds.), Quantum fields and strings: a course for mathematicians, Amer. Math. Soc., 1999. \BibBackref{36}
\item[\hypertarget{bib:37}{[37]}] S. K. Donaldson, Floer homology groups in Yang-Mills theory, Cambridge Tracts in Math. 147, Cambridge University Press, 2002. \BibBackref{37}
\item[\hypertarget{bib:38}{[38]}] —–, Symplectic submanifolds and almost-complex geometry, J. Differential Geom. 44 (1996), 666–705. \BibBackref{38}
\item[\hypertarget{bib:39}{[39]}] —–, Lefschetz pencils on symplectic manifolds, J. Differential Geom. 53 (1999), 205–236. \BibBackref{39}
\item[\hypertarget{bib:40}{[40]}] —–, Polynomials, vanishing cycles, and Floer homology, in Mathematics: frontiers and perspectives, Amer. Math. Soc., 2000, pp. 55–64. \BibBackref{40}
\item[\hypertarget{bib:41}{[41]}] Yu. Drozd, Tame and wild matrix problems, in Thirteen papers in algebra and number theory, Amer. Math. Soc. Transl. 128, Amer. Math. Soc., 1986, pp. 31–55. \BibBackref{41}
\item[\hypertarget{bib:42}{[42]}] Ya. Eliashberg and M. Gromov, Convex symplectic manifolds, in Several complex variables, Proc. Symposia Pure Math. 52, Amer. Math. Soc., 1991, pp. 135–162. \BibBackref{42}
\item[\hypertarget{bib:43}{[43]}] J. Etnyre, L. Ng, and J. Sabloff, Invariants of Legendrian knots and coherent orientations, J. Symplectic Geom. 1 (2002), 321–367. \BibBackref{43}
\item[\hypertarget{bib:44}{[44]}] A. Floer, Morse theory for Lagrangian intersections, J. Differential Geom. 28 (1988), 513–547. \BibBackref{44}
\item[\hypertarget{bib:45}{[45]}] —–, A relative Morse index for the symplectic action, Comm. Pure Appl. Math. 41 (1988), 393–407. \BibBackref{45}
\item[\hypertarget{bib:46}{[46]}] —–, The unregularized gradient flow of the symplectic action, Comm. Pure Appl. Math. 41 (1988), 775–813. \BibBackref{46}
\item[\hypertarget{bib:47}{[47]}] —–, Witten's complex and infinite dimensional Morse theory, J. Differential Geom. 30 (1989), 207–221. \BibBackref{47}
\item[\hypertarget{bib:48}{[48]}] A. Floer and H. Hofer, Coherent orientations for periodic orbit problems in symplectic geometry, Math. Z. 212 (1993), 13–38. \BibBackref{48}
\item[\hypertarget{bib:49}{[49]}] A. Floer, H. Hofer, and D. Salamon, Transversality in elliptic Morse theory for the symplectic action, Duke Math. J. 80 (1995), 251–292. \BibBackref{49}
\item[\hypertarget{bib:50}{[50]}] U. Frauenfelder, Gromov convergence of holomorphic discs, Diplomarbeit, ETH Zürich, 2001. \BibBackref{50}
\item[\hypertarget{bib:51}{[51]}] K. Fukaya, Floer homology for 3-manifolds with boundary (abstract), in Topology, Geometry, and Field theory, World Scientific, 1992, pp. 1–21. \BibBackref{51}
\item[\hypertarget{bib:52}{[52]}] —–, Morse homotopy, $A_\infty$-categories, and Floer homologies, in Proceedings of GARC workshop on Geometry and Topology (H. J. Kim, ed.), Seoul National University, 1993, pp. 1–102. \BibBackref{52}
\item[\hypertarget{bib:53}{[53]}] —–, Floer homology for three-manifolds with boundary I, Preprint, 1997. \BibBackref{53}
\item[\hypertarget{bib:54}{[54]}] —–, Morse homotopy and its quantization, in Geometric topology (Athens, GA, 1993) (W. H. Kazez, ed.), American Math. Society, 1997, pp. 409–440. \BibBackref{54}
\item[\hypertarget{bib:55}{[55]}] —–, Floer homology and mirror symmetry. II, in Minimal surfaces, geometric analysis and symplectic geometry (Baltimore, MD, 1999), Math. Soc. Japan, 2002, pp. 31–127. \BibBackref{55}
\item[\hypertarget{bib:56}{[56]}] —–, Deformation theory, homological algebra, and mirror symmetry, in Geometry and physics of branes (Como, 2001), IOP, 2003, pp. 121–209. \BibBackref{56}
\item[\hypertarget{bib:57}{[57]}] —–, Mirror symmetry of abelian varieties and multi-theta functions, J. Algebraic Geom. 11 (2002), 393–512. \BibBackref{57}
\item[\hypertarget{bib:58}{[58]}] K. Fukaya and Y.-G. Oh, Zero-loop open strings in the cotangent bundle and Morse homotopy, Asian J. Math. 1 (1998), 96–180. \BibBackref{58}
\item[\hypertarget{bib:59}{[59]}] K. Fukaya, P. Seidel, and I. Smith, Exact Lagrangian submanifolds in simply-connected cotangent bundles, Invent. Math., to appear. \BibBackref{59}
\item[\hypertarget{bib:60}{[60]}] K. Fukaya, Y.-G. Oh, H. Ohta, and K. Ono, Lagrangian intersection Floer theory–anomaly and obstruction, Preprint, first version, 2000; second expanded version, 2006. \BibBackref{60}
\item[\hypertarget{bib:61}{[61]}] S. Gelfand and Yu. Manin, Methods of homological algebra, Springer-Verlag, 1996. \BibBackref{61}
\item[\hypertarget{bib:62}{[62]}] E. Getzler and J. D. S. Jones, $A_\infty$-algebras and the cyclic bar complex, Illinois J. Math. 34 (1990), 256–283. \BibBackref{62}
\item[\hypertarget{bib:63}{[63]}] R. Gompf, Handlebody construction of Stein surfaces, Ann. of Math. 148 (1998), 619–693. \BibBackref{63}
\item[\hypertarget{bib:64}{[64]}] —–, Symplectic structures from Lefschetz pencils in high dimensions, in Proceedings of the Casson Fest, Geom. Topol. Publishers, 2004, pp. 267–290. \BibBackref{64}
\item[\hypertarget{bib:65}{[65]}] A. Gorodentsev, Exceptional objects and mutations in derived categories, in \hyperlink{bib:117}{[117]}, pp. 57–74. \BibBackref{65}
\item[\hypertarget{bib:66}{[66]}] A. Gorodentsev and S. Kuleshov, Helix theory, Moscow Math. J. 4 (2004), 377–440. \BibBackref{66}
\item[\hypertarget{bib:67}{[67]}] A. Gorodentsev and A. Rudakov, Exceptional vector bundles on projective spaces, Duke Math. J. 54 (1987), 115–130. \BibBackref{67}
\item[\hypertarget{bib:68}{[68]}] P. Griffiths and J. Morgan, Rational homotopy theory and differential forms, Progr. Math. 16, Birkhäuser, 1981. \BibBackref{68}
\item[\hypertarget{bib:69}{[69]}] V. Gugenheim, On a perturbation theory for the homology of the loop-space, J. Pure Appl. Algebra 25 (1982), 197–205. \BibBackref{69}
\item[\hypertarget{bib:70}{[70]}] V. Gugenheim, L. Lambe, and J. Stasheff, Algebraic aspects of Chen's twisting cochain, Illinois J. Math. 34 (1990), 485–502. \BibBackref{70}
\item[\hypertarget{bib:71}{[71]}] D. Happel, Hochschild cohomology of finite-dimensional algebras, in Séminaire d'algèbre Paul Dubreil et Marie-Paul Malliavin, 39ème Année (Paris, 1987/1988), Lecture Notes in Math. 1404, Springer-Verlag, 1989, pp. 108–126. \BibBackref{71}
\item[\hypertarget{bib:72}{[72]}] J. Harris and I. Morrison, Moduli of curves, Grad. Texts in Math. 187, Springer-Verlag, 1998. \BibBackref{72}
\item[\hypertarget{bib:73}{[73]}] F. Harvey and H. B. Lawson, Finite volume flows and Morse theory, Ann. of Math. 153 (2001), 1–25. \BibBackref{73}
\item[\hypertarget{bib:74}{[74]}] R. Hind, Lagrangian isotopies in Stein manifolds, Preprint math.SG/0311093, 2003. \BibBackref{74}
\item[\hypertarget{bib:75}{[75]}] —–, Lagrangian spheres in $S^2\times S^2$, Geom. Funct. Anal. 14 (2004), 303–318. \BibBackref{75}
\item[\hypertarget{bib:76}{[76]}] H. Hofer, Pseudoholomorphic curves in symplectizations with applications to the Weinstein conjecture in dimension three, Invent. Math. 114 (1993), 515–563. \BibBackref{76}
\item[\hypertarget{bib:77}{[77]}] Y.-Z. Huang, Two-dimensional conformal geometry and vertex operator algebras, Progr. Math. 148, Birkhäuser, 1997. \BibBackref{77}
\item[\hypertarget{bib:78}{[78]}] M. Hutchings, Reidemeister torsion in generalized Morse theory, Forum Math. 14 (2002), 209–244. \BibBackref{78}
\item[\hypertarget{bib:79}{[79]}] E.-N. Ionel and T. Parker, Gromov-Witten invariants of symplectic sums, Math. Res. Lett. 5 (1998), 563–576. \BibBackref{79}
\item[\hypertarget{bib:80}{[80]}] T. Kadeishvili, The algebraic structure in the homology of an $A(\infty)$-algebra, Soobshch. Akad. Nauk Gruzin. SSR 108 (2) (1982), 249–252. \BibBackref{80}
\item[\hypertarget{bib:81}{[81]}] —–, The category of differential coalgebras and the category of $A(\infty)$-algebras, Trudy Tbiliss. Mat. Inst. Razmadze Akad. Nauk Gruzin. SSR 77 (1985), 50–70. \BibBackref{81}
\item[\hypertarget{bib:82}{[82]}] B. Keller, Introduction to A-infinity algebras and modules, Homology Homotopy Appl. 3 (2001), 1–35. \BibBackref{82}
\item[\hypertarget{bib:83}{[83]}] —–, Bimodule complexes via strong homotopy actions, Algebr. Represent. Theory 3 (2000), 357–376. \BibBackref{83}
\item[\hypertarget{bib:84}{[84]}] —–, A-infinity algebras in representation theory, in Representations of algebras, vol. 1, Beijing Normal University Press, 2002, pp. 74–86. \BibBackref{84}
\item[\hypertarget{bib:85}{[85]}] M. Khovanov and P. Seidel, Quivers, Floer cohomology, and braid group actions, J. Amer. Math. Soc. 15 (2002), 203–271. \BibBackref{85}
\item[\hypertarget{bib:86}{[86]}] R. Kirby and L. Taylor, Pin structures on low-dimensional manifolds, in Geometry of low-dimensional manifolds (S. Donaldson and C. B. Thomas, eds.), vol. 2, London Math. Soc. Lecture Note Ser. 151, Cambridge University Press, 1990, pp. 177–242. \BibBackref{86}
\item[\hypertarget{bib:87}{[87]}] F. Knudsen and D. Mumford, The projectivity of the moduli space of stable curves I: preliminaries on ``det'' and ``div'', Math. Scand. 39 (1976), 19–55. \BibBackref{87}
\item[\hypertarget{bib:88}{[88]}] M. Kontsevich, Deformation quantization of Poisson manifolds I, Lett. Math. Phys. 66 (2003), 157–216. \BibBackref{88}
\item[\hypertarget{bib:89}{[89]}] —–, Lectures at ENS Paris, Spring 1998, set of notes taken by J. Bellaiche, J.-F. Dat, I. Marin, G. Racinet and H. Randriambololona. \BibBackref{89}
\item[\hypertarget{bib:90}{[90]}] —–, Homological algebra of mirror symmetry, in Proceedings of the International Congress of Mathematicians (Zürich, 1994), Birkhäuser, 1995, pp. 120–139. \BibBackref{90}
\item[\hypertarget{bib:91}{[91]}] D. Kwon and Y.-G. Oh, Structure of the image of (pseudo)-holomorphic discs with totally real boundary conditions, Comm. Anal. Geom. 8 (2000), 31–82. \BibBackref{91}
\item[\hypertarget{bib:92}{[92]}] K. Lamotke, The topology of complex projective varieties after S. Lefschetz, Topology 20 (1981), 15–51. \BibBackref{92}
\item[\hypertarget{bib:93}{[93]}] H. B. Lawson and M.-L. Michelsohn, Spin geometry, Princeton University Press, 1989. \BibBackref{93}
\item[\hypertarget{bib:94}{[94]}] Y.-J. Lee, Reidemeister torsion in symplectic Floer theory and counting pseudoholomorphic tori, I, J. Symplectic Geom. 3 (2005), 221–311. \BibBackref{94}
\item[\hypertarget{bib:95}{[95]}] K. Lefèvre, Sur les $A_\infty$-catégories, Ph.D. thesis, Université Paris 7, 2002. \BibBackref{95}
\item[\hypertarget{bib:96}{[96]}] J.-L. Loday, Cyclic homology, Springer-Verlag, 2nd edition, 1997. \BibBackref{96}
\item[\hypertarget{bib:97}{[97]}] V. Lyubashenko, Category of $A_\infty$-categories, Homology Homotopy Appl. 5 (2003), 1–48. \BibBackref{97}
\item[\hypertarget{bib:98}{[98]}] M. Markl, Transferring $A_\infty$ (strongly homotopy associative) structures, Rend. Circ. Mat. Palermo (2) Suppl. No. 79 (2006), 139–151. \BibBackref{98}
\item[\hypertarget{bib:99}{[99]}] —–, Homotopy algebras are homotopy algebras, Forum Math. 16 (2004), 129–160. \BibBackref{99}
\item[\hypertarget{bib:100}{[100]}] D. McDuff, Symplectic manifolds with contact type boundaries, Invent. Math. 103 (1991), 651–671. \BibBackref{100}
\item[\hypertarget{bib:101}{[101]}] D. Nadler, Microlocal branes are constructible sheaves, Preprint math.SG/0612399, 2007. \BibBackref{101}
\item[\hypertarget{bib:102}{[102]}] D. Nadler and E. Zaslow, Constructible Sheaves and the Fukaya Category, Preprint math.SG/0604379, 2007. \BibBackref{102}
\item[\hypertarget{bib:103}{[103]}] Y.-G. Oh, Floer cohomology of Lagrangian intersections and pseudo-holomorphic discs I, Comm. Pure Appl. Math. 46 (1993), 949–994. \BibBackref{103}
\item[\hypertarget{bib:104}{[104]}] C. Okonek, M. Schneider, and H. Spindler, Vector bundles on complex projective spaces, Progr. Math. 3, Birkhäuser, 1980. \BibBackref{104}
\item[\hypertarget{bib:105}{[105]}] P. Piazza, Determinant bundles, manifolds with boundary and surgery, Comm. Math. Phys. 178 (1996), 597–626. \BibBackref{105}
\item[\hypertarget{bib:106}{[106]}] S. Piunikhin, D. Salamon, and M. Schwarz, Symplectic Floer-Donaldson theory and quantum cohomology, in Contact and symplectic geometry (C. B. Thomas, ed.), Cambridge University Press, 1996, pp. 171–200. \BibBackref{106}
\item[\hypertarget{bib:107}{[107]}] A. Polishchuk, Massey and Fukaya products on elliptic curves, Adv. Theor. Math. Phys. 4 (2000), 1187–1207. \BibBackref{107}
\item[\hypertarget{bib:108}{[108]}] A. Polishchuk and E. Zaslow, Categorical mirror symmetry: the elliptic curve, Adv. Theor. Math. Phys. 2 (1998), 443–470. \BibBackref{108}
\item[\hypertarget{bib:109}{[109]}] L. Polterovich, Surgery of Lagrange submanifolds, Geom. Funct. Anal. 1 (1991), 198–210. \BibBackref{109}
\item[\hypertarget{bib:110}{[110]}] S. Priddy, Koszul resolutions, Trans. Amer. Math. Soc. 152 (1970), 39–60. \BibBackref{110}
\item[\hypertarget{bib:111}{[111]}] A. Prouté, Algèbres différentielles fortement homotopiquement associatives, Thèse d'État, Université Paris VII, 1985. \BibBackref{111}
\item[\hypertarget{bib:112}{[112]}] J. Robbin and D. Salamon, The Maslov index for paths, Topology 32 (1993), 827–844. \BibBackref{112}
\item[\hypertarget{bib:113}{[113]}] —–, The spectral flow and the Maslov index, Bull. London Math. Soc. 27 (1995), 1–33. \BibBackref{113}
\item[\hypertarget{bib:114}{[114]}] —–, Asymptotic behaviour of holomorphic strips, Ann. Inst. H. Poincaré Anal. Non Linéaire 18 (2001), 573–612. \BibBackref{114}
\item[\hypertarget{bib:115}{[115]}] Y. Ruan and G. Tian, A mathematical theory of Quantum Cohomology, J. Differential Geom. 42 (1995), 259–367. \BibBackref{115}
\item[\hypertarget{bib:116}{[116]}] —–, Higher genus symplectic invariants and sigma models coupled with gravity, Invent. Math. 130 (1997), 455–516. \BibBackref{116}
\item[\hypertarget{bib:117}{[117]}] A. Rudakov et al., Helices and vector bundles: Seminaire Rudakov, London Math. Soc. Lecture Note Ser. 148, Cambridge University Press, 1990. \BibBackref{117}
\item[\hypertarget{bib:118}{[118]}] D. Salamon and E. Zehnder, Morse theory for periodic solutions of Hamiltonian systems and the Maslov index, Comm. Pure Appl. Math. 45 (1992), 1303–1360. \BibBackref{118}
\item[\hypertarget{bib:119}{[119]}] M. Schwarz, Cohomology operations from $S^1$-cobordisms in Floer homology, Ph.D. thesis, ETH Zürich, 1995. \BibBackref{119}
\item[\hypertarget{bib:120}{[120]}] —–, A quantum cup-length estimate for symplectic fixed points, Invent. Math. 133 (1998), 353–397. \BibBackref{120}
\item[\hypertarget{bib:121}{[121]}] S. Scott, Splitting the curvature of the determinant line bundle, Proc. Amer. Math. Soc. 128 (2000), 2763–2775. \BibBackref{121}
\item[\hypertarget{bib:122}{[122]}] G. Segal, Lectures on topological field theory (Stanford notes), Manuscript, 1999. \BibBackref{122}
\item[\hypertarget{bib:123}{[123]}] P. Seidel, $\pi_1$ of symplectic automorphism groups and invertibles in quantum homology rings, Geom. Funct. Anal. 7 (1997), 1046–1095. \BibBackref{123}
\item[\hypertarget{bib:124}{[124]}] —–, Lagrangian two-spheres can be symplectically knotted, J. Differential Geom. 52 (1999), 145–171. \BibBackref{124}
\item[\hypertarget{bib:125}{[125]}] —–, Graded Lagrangian submanifolds, Bull. Soc. Math. France 128 (2000), 103–146. \BibBackref{125}
\item[\hypertarget{bib:126}{[126]}] —–, More about vanishing cycles and mutation, in Symplectic Geometry and Mirror Symmetry (Proceedings of the 4th KIAS Annual International Conference) (K. Fukaya, Y.-G. Oh, K. Ono, and G. Tian, eds.), World Scientific, 2001, pp. 429–465. \BibBackref{126}
\item[\hypertarget{bib:127}{[127]}] —–, Vanishing cycles and mutation, in Proceedings of the 3rd European Congress of Mathematics (Barcelona), Vol. 2, Birkhäuser, 2001, pp. 65–85. \BibBackref{127}
\item[\hypertarget{bib:128}{[128]}] —–, Braids and symplectic four-manifolds with abelian fundamental group, Turkish J. Math. 26 (2001 Gököva Geometry and Topology Conference Proceedings Special Volume) (2002), 93–100. \BibBackref{128}
\item[\hypertarget{bib:129}{[129]}] —–, Fukaya categories and deformations, in Proceedings of the International Congress of Mathematicians (Beijing, 2002), Vol. II, Higher Ed. Press, 2002, pp. 351–360. \BibBackref{129}
\item[\hypertarget{bib:130}{[130]}] —–, Homological mirror symmetry for the quartic surface, Preprint math.SG/0310414, 2003. \BibBackref{130}
\item[\hypertarget{bib:131}{[131]}] —–, A long exact sequence for symplectic Floer cohomology, Topology 42 (2003), 1003–1063. \BibBackref{131}
\item[\hypertarget{bib:132}{[132]}] —–, Exact Lagrangian submanifolds in $T^*S^n$ and the graded Kronecker quiver, in Different faces of geometry, Kluwer/Plenum, 2004, pp. 349–364. \BibBackref{132}
\item[\hypertarget{bib:133}{[133]}] P. Seidel and I. Smith, The symplectic topology of Ramanujam's surface, Comment. Math. Helv. 80 (2005), 859–881. \BibBackref{133}
\item[\hypertarget{bib:134}{[134]}] P. Seidel and R. Thomas, Braid group actions on derived categories of coherent sheaves, Duke Math. J. 108 (2001), 37–108. \BibBackref{134}
\item[\hypertarget{bib:135}{[135]}] V. De Silva, Products in the symplectic Floer homology of Lagrangian intersections, Ph.D. thesis, Oxford University, 1998. \BibBackref{135}
\item[\hypertarget{bib:136}{[136]}] N. Spaltenstein, Resolutions of unbounded complexes, Compositio Math. 65 (1988), 121–154. \BibBackref{136}
\item[\hypertarget{bib:137}{[137]}] J. Stasheff, Homotopy associativity of H-spaces I, Trans. Amer. Math. Soc. 108 (1963), 275–292. \BibBackref{137}
\item[\hypertarget{bib:138}{[138]}] R. Switzer, Algebraic topology - homotopy and homology, Grundlehren Math. Wiss. 212, Springer-Verlag, 1975. \BibBackref{138}
\item[\hypertarget{bib:139}{[139]}] R. Thomas, Stability conditions and the braid group, Comm. Anal. Geom. 14 (2006), 135–161. \BibBackref{139}
\item[\hypertarget{bib:140}{[140]}] R. Thomas and S.-T. Yau, Special Lagrangians, stable bundles and mean curvature flow, Comm. Anal. Geom. 10 (2002), 1075–1113. \BibBackref{140}
\item[\hypertarget{bib:141}{[141]}] J.-L. Verdier, Des catégories dérivées des catégories abéliennes, Astérisque 239, Soc. Math. France, 1996. \BibBackref{141}
\item[\hypertarget{bib:142}{[142]}] B. Wajnryb, Artin groups and geometric monodromy, Invent. Math. 138 (1999), 563–571. \BibBackref{142}
\item[\hypertarget{bib:143}{[143]}] C. A. Weibel, An introduction to homological algebra, Cambridge University Press, 1994. \BibBackref{143}
\end{list}

\chapter*{符号表}
\phantomsection\addcontentsline{toc}{chapter}{符号表}\label{symbols}
\markboth{符号表}{符号表}
\newcommand{\SymbolPage}[1]{\hyperref[source:\number\numexpr#1+8\relax]{\pageref*{source:\number\numexpr#1+8\relax}}}
\setlength{\LTpre}{0pt}
\begin{longtable}{@{}p{0.32\textwidth}p{0.54\textwidth}r@{}}
\toprule
符号 & 含义 & 页码\\
\midrule
\endfirsthead
\toprule
符号 & 含义 & 页码\\
\midrule
\endhead
\bottomrule
\endfoot
$\sgnmark_n$ & 符号约定中经常出现的项 & \SymbolPage{8}\\[4pt]
$\operatorname{Aut}(M)$ & 正合辛自同构群 & \SymbolPage{96}\\[4pt]
$\operatorname{Aut}(\widehat M)$ & 接上锥体后的辛自同构群 & \SymbolPage{98}\\[4pt]
$\operatorname{Aut}(M,\partial M)$ & 相对于边界的正合辛自同构群 & \SymbolPage{96}\\[4pt]
$\operatorname{Aut}^{\mathrm{gr}}(M,\partial M)$ & 分次辛自同构群 & \SymbolPage{188}\\[4pt]
$\operatorname{Aut}^{\#}(M,\partial M)$ & 膜自同构群 & \SymbolPage{188}\\[4pt]
$CF^{\mathrm{pr}}(L_0,L_1)$ & Floer 上链群，初步版本 & \SymbolPage{109}\\[4pt]
$CF^*(L_0^\#,L_1^\#)$ & Floer 上链群 & \SymbolPage{181}\\[4pt]
$\operatorname{Crit}(\pi)$ & 临界点集 & \SymbolPage{212}\\[4pt]
$\operatorname{Critv}(\pi)$ & 临界值集 & \SymbolPage{212}\\[4pt]
$D(\A),\ D^\pi(\A)$ & 导出范畴 & \SymbolPage{42}\\[4pt]
$D_{S,u}$ & 线性化算子 & \SymbolPage{107}\\[4pt]
$D_{\mathcal S,r,u}$ & 扩展线性化算子 & \SymbolPage{122}\\[4pt]
$\partial^h E$ & 水平边界 & \SymbolPage{212}\\[4pt]
$\partial^v E$ & 垂直边界 & \SymbolPage{212}\\[4pt]
$\operatorname{Fakev}(\Psi,\varpi)$ & 伪临界值集 & \SymbolPage{220}\\[4pt]
$\F(M)^{\mathrm{pr}}$ & Fukaya 范畴，初步版本 & \SymbolPage{126}\\[4pt]
$\F(\pi)$ & Lefschetz 纤维化的 Fukaya 范畴 & \SymbolPage{282}\\[4pt]
$H(\A),\ H^0(\A)$ & 上同调范畴 & \SymbolPage{8}\\[4pt]
$H\F(M)$ & Donaldson--Fukaya 范畴 & \SymbolPage{181}\\[4pt]
$H\F(M)^{\mathrm{pr}}$ & Donaldson--Fukaya 范畴，初步版本 & \SymbolPage{101}\\[4pt]
$HF^{\mathrm{pr}}$ & Floer 上同调，初步版本 & \SymbolPage{101}, \SymbolPage{109}\\[4pt]
$HF^*(L_0^\#,L_1^\#)$ & Floer 上同调 & \SymbolPage{182}\\[4pt]
$\M^{d+1}(y_0,\ldots,y_d)$ & 伪全纯多边形的模空间 & \SymbolPage{126}\\[4pt]
$\M_{E/S}(\{y_\zeta\})$ & 非齐次伪全纯截面的模空间 & \SymbolPage{240}\\[4pt]
$\M_S(\{y_\zeta\}),\ \M_{\mathcal S}(\{y_\zeta\})$ & 伪全纯映射的模空间 & \SymbolPage{105}, \SymbolPage{122}\\[4pt]
$\M_Z^*(y_0,y_1)$ & Floer 轨线的模空间 & \SymbolPage{105}\\[4pt]
$\nufun(\A,\B),\ \operatorname{fun}(\A,\B)$ & $\Ai$ 函子的范畴 & \SymbolPage{10}, \SymbolPage{28}\\[4pt]
$\operatorname{nu\text{-}mod}(\A),\ \modcat(\A)$ & $\Ai$ 模的范畴 & \SymbolPage{19}, \SymbolPage{30}\\[4pt]
$\mathcal R^{d+1},\ \overline{\mathcal R}^{d+1}$ & 带标点圆盘的模空间及其紧化 & \SymbolPage{115}, \SymbolPage{118}\\[4pt]
$TE^v,\ TE^h$ & 沿纤维的切空间及其补空间 & \SymbolPage{212}\\[4pt]
$T_{Y_0}(Y_1)$ & 对象的扭转 & \SymbolPage{63}\\[4pt]
$Z$ & 无限条带 & \SymbolPage{100}\\[4pt]
$Z^\pm$ & 半无限条带 & \SymbolPage{103}\\
\end{longtable}

% Complete translation of the original subject index; page numbers refer to this edition.
\makeatletter
\newcommand{\BookIndex}[2]{%
  \begingroup
  \edef\BookIndexPage{\getpagerefnumber{source:#2}}%
  \imki@wrindexentry{main}{#1|hyperpage}{\BookIndexPage}%
  \endgroup}
\makeatother
\BookIndex{juedui zhibiao@绝对指标}{172}
\BookIndex{juedui zhibiao@绝对指标}{183}
\BookIndex{a category@A∞ 范畴}{16}
\BookIndex{a functor@A∞ 函子}{17}
\BookIndex{a module@A∞ 模}{27}
\BookIndex{jiehe@结合律方程}{16}
\BookIndex{jiehe@结合律方程}{140}
\BookIndex{jiehe@结合律方程}{193}
\BookIndex{Beilinson puxulie@Beilinson 谱序列}{88}
\BookIndex{Beilinson puxulie@Beilinson 谱序列}{298}
\BookIndex{bian pingxing yidong@辫平行移动}{242}
\BookIndex{shangtongdiao youxian@上同调有限}{72}
\BookIndex{shangtongdiao fanti@上同调范畴}{16}
\BookIndex{shangtongdiao quanzhongshi@上同调全忠实}{17}
\BookIndex{fuhe lv@复合律}{109}
\BookIndex{yizhixing@一致性!tiaodai@条带状端}{128}
\BookIndex{yizhixing@一致性!raodong@扰动数据}{132}
\BookIndex{yantuo yingshe@延拓映射}{118}
\BookIndex{yantuo yingshe@延拓映射}{145}
\BookIndex{chuanyue xing@穿越型}{168}
\BookIndex{shangtongdiao hanyao@上同调含幺}{30}
\BookIndex{shangtongdiao hanyao@上同调含幺!hanzi@函子}{31}
\BookIndex{shangtongdiao hanyao@上同调含幺!mo@模}{38}
\BookIndex{Dehnniuzhuan@Dehn 扭转}{231}
\BookIndex{daochu@导出范畴}{50}
\BookIndex{daochu@导出范畴}{68}
\BookIndex{hanglieshi xian@行列式线}{157}
\BookIndex{youxiang@有向 A∞ 范畴}{88}
\BookIndex{texuanji@特选基!xiaoshiquan@消失圈}{235}
\BookIndex{texuanji@特选基!xiaoshilujing@消失路径}{233}
\BookIndex{Donaldson Fukaya fanchou@Donaldson–Fukaya 范畴}{109}
\BookIndex{Donaldson Fukaya fanchou@Donaldson–Fukaya 范畴}{189}
\BookIndex{nengliang@能量!weiquanchun yingshe@伪全纯映射}{113}
\BookIndex{nengliang@能量!weiquanchunjiemian@伪全纯截面}{244}
\BookIndex{Lagrangian@Lagrangian 子流形!zhenghe@正合}{104}
\BookIndex{zhenghe xinliuxing@正合辛流形!jiechuxing@带接触型边界}{105}
\BookIndex{zhenghe xinliuxing@正合辛流形!daijiao@带角}{104}
\BookIndex{zhenghe@正合三角}{44}
\BookIndex{liwai ji@例外集}{83}
\BookIndex{wei linjizhi@伪临界值}{228}
\BookIndex{Riemann qumian zu@Riemann 曲面族}{121}
\BookIndex{daijidianxianweihua@带奇点的纤维化}{220}
\BookIndex{Floer@Floer 上同调}{109}
\BookIndex{Floer@Floer 上同调}{117}
\BookIndex{Floer@Floer 上同调}{190}
\BookIndex{Floer shuju@Floer 数据}{112}
\BookIndex{xingshi@形式微分同胚}{18}
\BookIndex{xingshixing@形式性}{313}
\BookIndex{biaoji@标架（Lagrangian 球面）}{228}
\BookIndex{ziyou duixiang@自由对象!Z2@带 Z/2 作用的范畴}{96}
\BookIndex{Fukaya@Fukaya 范畴}{134}
\BookIndex{Fukaya@Fukaya 范畴!dengbian@等变}{217}
\BookIndex{Fukaya@Fukaya 范畴!Lefschetz@Lefschetz 纤维化的}{290}
\BookIndex{shengcheng@生成元}{50}
\BookIndex{fenci@分次!Lagrangian ziliuxing@Lagrangian 子流形的}{182}
\BookIndex{fenci@分次!Lagrangian zikongjian@Lagrangian 子空间的}{177}
\BookIndex{fenci@分次!xin zitonggou@辛自同构的}{195}
\BookIndex{Gromov jinhua@Gromov 紧化}{138}
\BookIndex{qun zuoyong@群作用!fanchou@范畴上的}{95}
\BookIndex{qun zuoyong@群作用!fanchou@范畴上的}{143}
\BookIndex{hochschild cohomology@Hochschild 上同调}{21}
\BookIndex{tonglun@同伦!a functor@A∞ 函子}{23}
\BookIndex{Hurwitz dengjia@Hurwitz 等价}{236}
\BookIndex{Hurwitzyidong@Hurwitz 移动!xiaoshilujing@消失路径基的}{234}
\BookIndex{Hurwitzyidong@Hurwitz 移动!xiaoshiquan@消失圈基的}{235}
\BookIndex{Lagrangianbianjietiaojian@Lagrangian 边界条件}{243}
\BookIndex{Lagrangianbianjietiaojian@Lagrangian 边界条件!biaozhun@标准}{246}
\BookIndex{Lagrangian mo@Lagrangian 膜}{182}
\BookIndex{Lagrangian mo@Lagrangian 膜!chouxiang@抽象}{171}
\BookIndex{Lagrangian mo@Lagrangian 膜!dengbian@等变}{207}
\BookIndex{Lagrangian mo@Lagrangian 膜!xianxing@线性}{177}
\BookIndex{Lagrangian mo@Lagrangian 膜!xiangdui@相对}{266}
\BookIndex{Lefschetzshuangxianweihua@Lefschetz 双纤维化}{226}
\BookIndex{Lefschetzxianweihua@Lefschetz 纤维化}{224}
\BookIndex{Lefschetzxianweihua@Lefschetz 纤维化!tiaodaizhuangduan@带条带状端}{247}
\BookIndex{Lefschetz shu@Lefschetz 束}{303}
\BookIndex{Lefschetzzhitao@Lefschetz 顶针}{229}
\BookIndex{xianxinghua suanzi@线性化算子}{115}
\BookIndex{xianxinghua suanzi@线性化算子!kuozhan@扩展}{130}
\BookIndex{yingshe@映射锥}{43}
\BookIndex{yingshe@映射锥}{56}
\BookIndex{massey@Massey 积}{17}
\BookIndex{pipeitonghen@匹配同痕}{238}
\BookIndex{pipeilujing@匹配路径}{238}
\BookIndex{bianyi@变异}{85}
\BookIndex{ziran@自然变换}{19}
\BookIndex{faxiang jiao@法向角}{211}
\BookIndex{bamian@八面体公理}{49}
\BookIndex{dingxiang suanzi@定向算子}{182}
\BookIndex{dingxiang kongjian@定向空间}{172}
\BookIndex{dingxiang kongjian@定向空间}{183}
\BookIndex{raodong shuju@扰动数据}{112}
\BookIndex{raodong shuju@扰动数据!xiangdui@相对}{244}
\BookIndex{raodong@扰动引理}{25}
\BookIndex{Picard Lefschetz shuju@Picard–Lefschetz 数据}{308}
\BookIndex{Picard Lefschetz dingli@Picard–Lefschetz 定理}{232}
\BookIndex{Pin jiegou@Pin 结构}{175}
\BookIndex{PSS tonggou@Piunikhin–Salamon–Schwarz（PSS）同构}{110}
\BookIndex{PSS tonggou@Piunikhin–Salamon–Schwarz（PSS）同构}{188}
\BookIndex{daibiaodian yuanpan@带标点圆盘}{122}
\BookIndex{Postnikov@Postnikov 分解}{84}
\BookIndex{weiquanchun yingshe@伪全纯映射}{113}
\BookIndex{weiquanchun duobianxing@伪全纯多边形}{134}
\BookIndex{weiquanchunjiemian@伪全纯截面}{244}
\BookIndex{erci futiji xingshi@二次复体积形式}{177}
\BookIndex{erci futiji xingshi@二次复体积形式!xiangdui@相对}{224}
\BookIndex{ni dengjia@拟等价}{36}
\BookIndex{ni biaoshi@拟表示}{40}
\BookIndex{pingyi@平移!daishu@代数}{42}
\BookIndex{pingyi@平移!daishu@代数}{56}
\BookIndex{pingyi@平移!Lagrangian mo@Lagrangian 膜的}{179}
\BookIndex{pingyi@平移!Lagrangian mo@Lagrangian 膜的}{182}
\BookIndex{pingyi@平移!Lagrangian mo@Lagrangian 膜的}{194}
\BookIndex{qiumian duixiang@球面对象}{80}
\BookIndex{fenlie bi fanchou@分裂闭范畴}{63}
\BookIndex{fenlie bi fanchou@分裂闭范畴}{67}
\BookIndex{fenlie shengcheng@分裂生成元}{68}
\BookIndex{Stein liuxing@Stein 流形}{299}
\BookIndex{yange@严格含幺}{30}
\BookIndex{tiaodai zhuangduan@条带状端}{111}
\BookIndex{xin zitonggou@辛自同构!zhenghe@正合}{104}
\BookIndex{xin zitonggou@辛自同构!fenci@分次}{195}
\BookIndex{shu@树}{123}
\BookIndex{sanjiao ji@三角形乘积}{109}
\BookIndex{sanjiao ji@三角形乘积}{188}
\BookIndex{sanjiao hua fanti@三角化的 A∞ 范畴}{47}
\BookIndex{niuqu@扭曲复形}{52}
\BookIndex{niuzhuan@扭转}{71}
\BookIndex{niuzhuan@扭转}{79}
\BookIndex{B xing@(B)型 Lagrangian 子流形或膜}{275}
\BookIndex{B xing@(B)型 Lagrangian 子流形或膜}{286}
\BookIndex{U xing@(U)型 Lagrangian 子流形或膜}{275}
\BookIndex{U xing@(U)型 Lagrangian 子流形或膜}{284}
\BookIndex{xiaoshiquan@消失圈}{229}
\BookIndex{xiaoshilujing@消失路径}{229}
\BookIndex{yoneda@Yoneda 嵌入}{29}
\BookIndex{yoneda@Yoneda 嵌入}{39}

\printindex
\end{document}
